\documentclass[12pt]{article}

% ===== 切換模式：只保留一行未註解 =====
% \def\docmode{chinese}
% \def\docmode{german}
\def\docmode{bilingual}

% ===== 雙語對照橫式紙張：只保留一行未註解 =====
\def\bilingualpaper{a4}
% \def\bilingualpaper{a3}

% bverfge_layout.tex
% 共用排版基礎設施，供 Schwangerschaftsabbruch I 與 II 共享。
% 由各文件以 \documentclass 後 % bverfge_layout.tex
% 共用排版基礎設施，供 Schwangerschaftsabbruch I 與 II 共享。
% 由各文件以 \documentclass 後 % bverfge_layout.tex
% 共用排版基礎設施，供 Schwangerschaftsabbruch I 與 II 共享。
% 由各文件以 \documentclass 後 \input{../../共用LaTeX/bverfge_layout} 載入。
% 不含：\documentclass、\documentversion、\ggversionde/zh、\tocde/\toczh。
% 日期與首頁版本列由本檔自動產生；各文件只需手動指定 \documentversion。

\usepackage{xeCJK}
\usepackage{fontspec}
\usepackage{geometry}
\usepackage{setspace}
\usepackage{hyperref}
\usepackage{xurl}
\usepackage{bookmark}
\usepackage{enumitem}
\usepackage{xcolor}
\usepackage{titlesec}
\usepackage{array}
\usepackage{longtable}
\usepackage{tcolorbox}
\usepackage{environ}
\usepackage{pdflscape}
\usepackage{etoolbox}
\usepackage{ragged2e}
\usepackage{paracol}
\usepackage{multicol}
\usepackage{needspace}
\usepackage{fancyhdr}
\usepackage{tikz}
\usepackage{zref-abspage}
\tcbuselibrary{breakable,skins}
\usetikzlibrary{calc}
\usepackage{pgfcalendar}

\hypersetup{
  unicode=true,
  bookmarksopen=true,
  bookmarksnumbered=false,
  hidelinks
}

% ===== 自動推導，不要手動改下面 =====
% （\docmode 與 \bilingualpaper 由各文件在 \input{bverfge_layout} 之前設定；
%   預設 chinese / a4）
\ifdefined\docmode\else\def\docmode{chinese}\fi
\ifdefined\bilingualpaper\else\def\bilingualpaper{a4}\fi
\newif\ifshoworiginal
\newif\ifshowtranslation
\newif\ifbilinguallandscape
\newif\ifbilingualcolumns
\newif\ifintranslation
\newif\iffirstbilingualsection

\showoriginalfalse
\showtranslationfalse
\bilinguallandscapefalse
\bilingualcolumnsfalse
\intranslationfalse
\firstbilingualsectionfalse

\def\modechinese{chinese}
\def\modegerman{german}
\def\modebilingual{bilingual}
\def\bilingualpaperaThree{a3}
\def\bilingualpaperaFour{a4}

\ifx\docmode\modechinese
  \showtranslationtrue
\fi

\ifx\docmode\modegerman
  \showoriginaltrue
\fi

\ifx\docmode\modebilingual
  \showoriginaltrue
  \showtranslationtrue
  \bilinguallandscapetrue
  \bilingualcolumnstrue
\fi

% 編排模式說明：
% chinese  → \showoriginalfalse  \showtranslationtrue  （A4 直式，純中文）
% german   → \showoriginaltrue   \showtranslationfalse （A4 直式，純德文）
% bilingual→ \showoriginaltrue   \showtranslationtrue  \bilingualcolumnstrue（A3/A4 橫式對照）

\ifbilingualcolumns
  \ifx\bilingualpaper\bilingualpaperaFour
    \geometry{
      a4paper,
      landscape,
      left=1.25cm,
      right=2.25cm,
      top=1.25cm,
      bottom=1.75cm,
      footskip=8.5mm,
      marginparwidth=0.75cm,
      marginparsep=0.25cm
    }
  \else
    \geometry{
      a3paper,
      landscape,
      left=1.55cm,
      right=2.75cm,
      top=1.55cm,
      bottom=2.05cm,
      footskip=9.5mm,
      marginparwidth=0.9cm,
      marginparsep=0.35cm
    }
  \fi
\else
\geometry{
  a4paper,
  portrait,
  left=2.2cm,
  right=2.6cm,
  top=2.0cm,
  bottom=2.0cm,
  marginparwidth=0.8cm,
  marginparsep=0.3cm
}
\fi
% ===== 時間戳基礎設施 =====
\newcount\stamptimeval
\newcount\stampminuteval
\newcommand{\stamphh}{%
  \stamptimeval=\time
  \divide\stamptimeval by 60
  \ifnum\stamptimeval<10 0\fi
  \the\stamptimeval}
\newcommand{\stampmm}{%
  \stamptimeval=\time
  \stampminuteval=\time
  \divide\stampminuteval by 60
  \multiply\stampminuteval by 60
  \advance\stamptimeval by -\stampminuteval
  \ifnum\stamptimeval<10 0\fi
  \the\stamptimeval}
\newcommand{\stampwkde}[1]{%
  \ifcase#1Montag\or Dienstag\or Mittwoch\or Donnerstag\or Freitag\or Samstag\or Sonntag\fi}
\newcommand{\stampwkzh}[1]{%
  \ifcase#1 一\or 二\or 三\or 四\or 五\or 六\or 日\fi}
\newcommand{\stampmonthde}[1]{%
  \ifcase#1\or Januar\or Februar\or März\or April\or Mai\or Juni\or
  Juli\or August\or September\or Oktober\or November\or Dezember\fi}
\newcount\stampjdval
\newcount\stampwdval
\pgfcalendardatetojulian{\the\year-\the\month-\the\day}{\stampjdval}
\pgfcalendarjuliantoweekday{\the\stampjdval}{\stampwdval}
\newcommand{\timestampzh}{%
  \songfont 最後修改：\the\year 年\the\month 月\the\day 日%
  % （星期\stampwkzh{\the\stampwdval}）%
  % \space\stamphh:\stampmm
  }

\newcommand{\documentdatezh}{\the\year 年 \the\month 月 \the\day 日}
\newcommand{\documentdatede}{\the\day. \stampmonthde{\the\month} \the\year}
\newcommand{\bverfgetranslationnote}{%
  {\songfont 翻譯說明：初譯使用 Claude Code 與 GPT 協助迭代；仍請以德文原文為準。}%
}
\newcommand{\bverfgecourseinfo}{%
  {\songfont 課程：114 學年度第 2 學期；憲法專題研究（四）；授課老師：湯德宗教授；導讀日期：115 年 6 月 1 日。}%
}
\newcommand{\bverfgeeditorinfo}{%
  {\songfont 編者：王逸帆（R10A21126），國立台灣大學法律系研究所碩士班。}%
}
\newcommand{\bverfgetitleversionline}{%
  \ifbilingualcolumns
    Fassung \documentversion\enspace\textperiodcentered\enspace Stand: \documentdatede
    \quad
    {\songfont 版本 \documentversion\enspace\textperiodcentered\enspace 修訂日期：\documentdatezh\par}%
  \else
    \ifshowtranslation
      {\songfont 版本 \documentversion\enspace\textperiodcentered\enspace 修訂日期：\documentdatezh\par}%
    \else
      Fassung \documentversion\enspace\textperiodcentered\enspace Stand: \documentdatede\par
    \fi
  \fi
}
\newcommand{\bverfgetitlecornerinfo}{%
  \ifshowtranslation
    \vspace{0.72em}%
    \noindent\hfill
    \begin{minipage}{0.66\textwidth}
      \raggedleft\tiny\color{pendingink}\songfont
      \bverfgecourseinfo\par
      \bverfgeeditorinfo
    \end{minipage}%
  \fi
}
\newcommand{\bverfgetitlemeta}{%
  \vfill
  \begin{center}%
    \footnotesize\color{pendingink}%
    \bverfgetitleversionline
    \ifshowtranslation
      \vspace{0.28em}{\scriptsize\bverfgetranslationnote\par}%
    \fi
  \end{center}%
  \bverfgetitlecornerinfo
}

\pagestyle{fancy}
\fancyhf{}
\renewcommand{\headrulewidth}{0pt}
\renewcommand{\footrulewidth}{0pt}
\fancyfoot[C]{\thepage}
\ifbilingualcolumns
  \fancyfoot[R]{\scriptsize\color{pendingink}\timestampzh}%
\else
  \ifshoworiginal
  \else
    \fancyfoot[R]{\scriptsize\color{pendingink}\timestampzh\hspace{0.45cm}}%
  \fi
\fi
\setmainfont{Times New Roman}
\setsansfont{Helvetica Neue}
\setmonofont{Menlo}
\IfFileExists{/Users/iw/Library/Fonts/kaiu.ttf}{
  \setCJKmainfont[Path=/Users/iw/Library/Fonts/,AutoFakeBold=4]{kaiu.ttf}
}{
  \IfFontExistsTF{標楷體}{
    \setCJKmainfont[AutoFakeBold=4]{標楷體}
  }{
    \setCJKmainfont{Songti TC}
  }
}
\IfFontExistsTF{Songti TC}{
  \setCJKfamilyfont{song}{Songti TC}
  \setCJKmonofont{Songti TC}
}{
  \setCJKfamilyfont{song}[Path=/Users/iw/Library/Fonts/,AutoFakeBold=4]{kaiu.ttf}
  \setCJKmonofont[Path=/Users/iw/Library/Fonts/,AutoFakeBold=4]{kaiu.ttf}
}
\newcommand{\songfont}{\CJKfamily{song}}

% ===== 封面／目錄專用打字機字體（僅限封面、目錄頁使用，不影響正文）=====
\IfFontExistsTF{Radio Newsman}{%
  \newfontfamily\bverfgetitlede{Radio Newsman}%
}{%
  \newfontfamily\bverfgetitlede{Courier New}%
}
\IfFontExistsTF{Erikas Farbband}{%
  \newfontfamily\bverfgebodyde{Erikas Farbband}%
}{%
  \let\bverfgebodyde\bverfgetitlede
}
\IfFontExistsTF{Maoken Heavy Labourer}{%
  \setCJKfamilyfont{bverfgetitlecjk}{Maoken Heavy Labourer}%
}{%
  \setCJKfamilyfont{bverfgetitlecjk}{Songti TC}%
}
\IfFontExistsTF{Huiwen-mincho}{%
  \setCJKfamilyfont{bverfgebodycjk}{Huiwen-mincho}%
}{%
  \setCJKfamilyfont{bverfgebodycjk}{Songti TC}%
}
\newcommand{\bverfgetitlezh}{\bverfgetitlede\CJKfamily{bverfgetitlecjk}}
\newcommand{\bverfgebodyzh}{\bverfgebodyde\CJKfamily{bverfgebodycjk}}

\setstretch{1.35}
\setlist{itemsep=0.2em, topsep=0.4em}
\setlength{\parindent}{0pt}
\setlength{\parskip}{0.45em}
\raggedbottom
\setcounter{secnumdepth}{0}
\setcounter{tocdepth}{2}

\newcommand{\lawboxsourcefont}{\footnotesize}
\newcommand{\lawboxtransfont}{\footnotesize}
\newcommand{\lawboxsourcestretch}{1.02}
\newcommand{\lawboxtransstretch}{0.92}

\titleformat{\section}{\centering\Large\bfseries\songfont}{\thesection}{1em}{}
\titleformat{\subsection}{\large\bfseries\songfont}{\thesubsection}{1em}{}
\titleformat{\subsubsection}{\normalsize\bfseries\songfont}{\thesubsubsection}{1em}{}
\titleformat{\paragraph}{\normalsize\bfseries\songfont}{\theparagraph}{1em}{}
\titlespacing*{\section}{0pt}{1.8em}{1.0em}
\titlespacing*{\subsection}{0pt}{1.2em}{0.6em}
\titlespacing*{\subsubsection}{0pt}{1.0em}{0.45em}
\titlespacing*{\paragraph}{0pt}{0.6em}{0.4em}

\definecolor{noteblue}{HTML}{A9C9F5}
\definecolor{noteteal}{HTML}{AEE1D6}
\definecolor{notegray}{HTML}{D7DADF}
\definecolor{caseink}{HTML}{000000}
\definecolor{casepaper}{HTML}{FCFEFD}
\definecolor{sourcepaper}{HTML}{FAFBF8}
\definecolor{sourceline}{HTML}{D7DED8}
\definecolor{sourceink}{HTML}{000000}
\definecolor{pendingink}{HTML}{000000}

\tcbset{
  caseboxbase/.style={
    enhanced,
    breakable,
    width=0.985\linewidth,
    colback=casepaper,
    coltitle=caseink,
    fonttitle=\bfseries\songfont,
    boxrule=0.28pt,
    arc=1.2mm,
    left=2.4mm,
    right=2.4mm,
    top=1.5mm,
    bottom=1.5mm,
    boxsep=1.5mm,
    before skip=0.65em,
    after skip=0.65em,
    before upper={\setlength{\parindent}{0pt}\setlength{\parskip}{0.35em}},
  }
}

\newcommand{\bilingualstart}{}
\newcommand{\bilingualstop}{}
\ifbilingualcolumns
  \renewcommand{\bilingualstart}{\small\setlength{\columnsep}{1.25cm}\setlength{\columnseprule}{0.12pt}\columnratio{0.5,0.5}\multicolumnfootnotes\flushbottom\sloppy\interlinepenalty=80\clubpenalty=800\widowpenalty=800\displaywidowpenalty=800\begin{paracol}{2}\global\intranslationfalse\global\firstbilingualsectiontrue{\footnotesize\color{sourceink}Deutsch\par}\switchcolumn{\footnotesize\songfont\color{sourceink}中文譯文\par}\switchcolumn*}
  \renewcommand{\bilingualstop}{\ifintranslation\switchcolumn*\global\intranslationfalse\fi\end{paracol}\clearpage\raggedbottom}
\fi
\newif\ifrandnumbers
\randnumbersfalse
\newcounter{randnumber}
\newcounter{randnumberlocal}
\newcounter{lastsourcerandstart}
\newif\iflastsourcehasrandnumbers
\lastsourcehasrandnumbersfalse
\newif\ifinlawbox
\inlawboxfalse
\newif\iflawboxhaspendingrn
\lawboxhaspendingrnfalse
\newcommand{\lawboxpendingrn}{}
\newcommand{\startrandnumbers}{\global\randnumberstrue\setcounter{randnumber}{1}\global\lastsourcehasrandnumbersfalse}
\newcommand{\stoprandnumbers}{\global\randnumbersfalse\global\lastsourcehasrandnumbersfalse}
\newlength{\randnumberwidth}
\setlength{\randnumberwidth}{0.95cm}
\newlength{\randnumberrightoffset}
\setlength{\randnumberrightoffset}{0.85cm}
\newlength{\randnumberlawboxrightoffset}
\setlength{\randnumberlawboxrightoffset}{1.40cm}
\newcommand{\randnumberbox}[1]{%
  \leavevmode
  \makebox[0pt][l]{%
    \smash{%
      \tikz[remember picture,overlay,baseline=(rnmark.base)]{%
        \coordinate (rnmark) at (0,0);%
        \ifinlawbox
          \path let \p1=(current page.east), \p2=(rnmark.base) in
            node[
              anchor=base east,
              inner sep=0pt,
              outer sep=0pt,
              text width=\randnumberwidth,
              align=right,
              text=pendingink,
              font=\normalfont\mdseries\scriptsize\ttfamily
            ] at (\x1-\randnumberlawboxrightoffset,\y2) {{\normalfont\mdseries\scriptsize\ttfamily #1}};%
        \else
          \path let \p1=(current page.east), \p2=(rnmark.base) in
            node[
              anchor=base east,
              inner sep=0pt,
              outer sep=0pt,
              text width=\randnumberwidth,
              align=right,
              text=pendingink,
              font=\normalfont\mdseries\scriptsize\ttfamily
            ] at (\x1-\randnumberrightoffset,\y2) {{\normalfont\mdseries\scriptsize\ttfamily #1}};%
        \fi
      }%
    }%
  }%
  \ignorespaces
}
\newcommand{\lawboxstorerandnumber}[1]{%
  \gdef\lawboxpendingrn{#1}%
  \global\lawboxhaspendingrntrue
  \ignorespaces
}
\newcommand{\lawboxuserandnumber}{%
  \iflawboxhaspendingrn
    \randnumberbox{\lawboxpendingrn}%
    \global\lawboxhaspendingrnfalse
  \fi
}
\newcommand{\sourcern}[1]{%
  \ifrandnumbers
    \ifshoworiginal
      \ifshowtranslation\else
        \ifinlawbox\lawboxstorerandnumber{#1}\else\randnumberbox{#1}\fi
      \fi
    \fi
  \fi
}
\newcommand{\transrn}[1]{%
  \ifrandnumbers
    \ifshowtranslation
      \ifinlawbox\lawboxstorerandnumber{#1}\else\randnumberbox{#1}\fi
    \fi
  \fi
}
% ===== 課堂模式：德文原文縮寫註解 =====
% 日常切換只改各判決檔的一行：
%   \classroommodeon        開啟課堂版；註解掉這行即回到一般版。
% 模板預設：
%   1. 課堂模式關閉。
%   2. 課堂模式開啟時，縮寫註解包含「德文全稱＋中文譯名」。
%   3. 課堂模式開啟時，GG/條文編者註仍會自動顯示。
% 進階微調才需要在單案檔覆寫：
%   \classroomabbrzhoff     縮寫註解僅顯示德文全稱。
%   \showeditornotestrue    一般版也顯示編者註。
% 註解策略：同一實際 PDF 頁面中，同一縮寫只在第一次出現處加註。
% 這裡用 zref-abspage 的絕對頁序，不用 \thepage，避免文件內頁碼重置時誤判。
\newif\ifclassroommode
\newif\ifclassroomabbrzh
\classroommodefalse
\classroomabbrzhtrue
\newcommand{\classroommodeon}{\classroommodetrue}
\newcommand{\classroommodeoff}{\classroommodefalse}
\newcommand{\classroomabbrzhon}{\classroomabbrzhtrue}
\newcommand{\classroomabbrzhoff}{\classroomabbrzhfalse}
\newcommand{\abbrnotelabel}{\emph{Abk.}: }
\newcommand{\abbrnote}[3]{%
  \ifclassroommode
    \ifshoworiginal
      \ifintranslation\else
        \footnote{\normalfont\footnotesize\abbrnotelabel #2\ifclassroomabbrzh；{\songfont #3}\fi}%
      \fi
    \fi
  \fi
}
\newcommand{\abbrnoteonce}[4]{%
  #2%
  \ifclassroommode
    \ifshoworiginal
      \ifintranslation\else
        \begingroup
          \edef\abbrcurrentabspage{\number\numexpr\value{abspage}+1\relax}%
          \ifcsname abbrnoteseen@#1@\abbrcurrentabspage\endcsname\else
            \global\expandafter\let\csname abbrnoteseen@#1@\abbrcurrentabspage\endcsname\empty
            \abbrnote{#1}{#3}{#4}%
          \fi
        \endgroup
      \fi
    \fi
  \fi
}
\newcommand{\abbrAbs}{\abbrnoteonce{abs}{Abs.}{Absatz}{項}}
\newcommand{\abbrAAO}{\abbrnoteonce{aao}{a.a.O.}{am angegebenen Ort}{前引處}}
\newcommand{\abbrAE}{\abbrnoteonce{ae}{AE}{Alternativ-Entwurf}{替代草案}}
\newcommand{\abbrArt}{\abbrnoteonce{art}{Art.}{Artikel}{條}}
\newcommand{\abbrAufl}{\abbrnoteonce{aufl}{Aufl.}{Auflage}{版}}
\newcommand{\abbrBd}{\abbrnoteonce{bd}{Bd.}{Band}{卷}}
\newcommand{\abbrBGBl}{\abbrnoteonce{bgbl}{BGBl.}{Bundesgesetzblatt}{聯邦法律公報}}
\newcommand{\abbrBGHSt}{\abbrnoteonce{bghst}{BGHSt}{Entscheidungen des Bundesgerichtshofes in Strafsachen}{聯邦普通法院刑事判決彙編}}
\newcommand{\abbrBRDrucks}{\abbrnoteonce{brdrucks}{BRDrucks.}{Bundesrats-Drucksache}{聯邦參議院文件}}
\newcommand{\abbrBTDrucks}{\abbrnoteonce{btdrucks}{BTDrucks.}{Bundestags-Drucksache}{聯邦議院文件}}
\newcommand{\abbrBundesgesetzbl}{\abbrnoteonce{bundesgesetzbl}{Bundesgesetzbl.}{Bundesgesetzblatt}{聯邦法律公報}}
\newcommand{\abbrBvF}{\abbrnoteonce{bvf}{BvF}{Bundesverfassungsgerichtsverfahren}{聯邦憲法法院程序案號}}
\newcommand{\abbrBvL}{\abbrnoteonce{bvl}{BvL}{Bundesverfassungsgerichtsverfahren}{聯邦憲法法院程序案號}}
\newcommand{\abbrBvR}{\abbrnoteonce{bvr}{BvR}{Bundesverfassungsgerichtsverfahren}{聯邦憲法法院程序案號}}
\newcommand{\abbrBVerfG}{\abbrnoteonce{bverfg}{BVerfG}{Bundesverfassungsgericht}{聯邦憲法法院}}
\newcommand{\abbrBVerfGE}{\abbrnoteonce{bverfge}{BVerfGE}{Entscheidungen des Bundesverfassungsgerichts}{聯邦憲法法院判決集}}
\newcommand{\abbrBVerfGG}{\abbrnoteonce{bverfgg}{BVerfGG}{Bundesverfassungsgerichtsgesetz}{聯邦憲法法院法}}
\newcommand{\abbrDDR}{\abbrnoteonce{ddr}{DDR}{Deutsche Demokratische Republik}{德意志民主共和國；東德}}
\newcommand{\abbrDJT}{\abbrnoteonce{djt}{DJT}{Deutscher Juristentag}{德國法學家大會}}
\newcommand{\abbrDr}{\abbrnoteonce{dr}{Dr.}{Doktor}{Dr.}}
\newcommand{\abbrEuGRZ}{\abbrnoteonce{eugrz}{EuGRZ}{Europäische Grundrechte-Zeitschrift}{歐洲基本權期刊}}
\newcommand{\abbrEV}{\abbrnoteonce{ev}{EV}{Einigungsvertrag}{統一條約}}
\newcommand{\abbrF}{\abbrnoteonce{f}{f.}{folgende}{以下一頁；下一段}}
\newcommand{\abbrFF}{\abbrnoteonce{ff}{ff.}{fortfolgende}{以下數頁；以下各段}}
\newcommand{\abbrGBlDDR}{\abbrnoteonce{gblddr}{GBl. DDR}{Gesetzblatt der Deutschen Demokratischen Republik}{德意志民主共和國法律公報}}
\newcommand{\abbrGG}{\abbrnoteonce{gg}{GG}{Grundgesetz}{基本法}}
\newcommand{\abbrGVBl}{\abbrnoteonce{gvbl}{GVBl.}{Gesetz- und Verordnungsblatt}{法律及命令公報}}
\newcommand{\abbrJR}{\abbrnoteonce{jr}{JR}{Juristische Rundschau}{法學評論}}
\newcommand{\abbrJZ}{\abbrnoteonce{jz}{JZ}{JuristenZeitung}{法學家報}}
\newcommand{\abbrIVm}{\abbrnoteonce{ivm}{i.V.m.}{in Verbindung mit}{結合；連同}}
\newcommand{\abbrMdB}{\abbrnoteonce{mdb}{MdB}{Mitglied des Bundestages}{聯邦議院議員}}
\newcommand{\abbrMwN}{\abbrnoteonce{mwn}{m.w.N.}{mit weiteren Nachweisen}{附進一步文獻／判例出處}}
\newcommand{\abbrNF}{\abbrnoteonce{nf}{n.F.}{neue Fassung}{新版本}}
\newcommand{\abbrAF}{\abbrnoteonce{af}{a.F.}{alte Fassung}{舊版本}}
\newcommand{\abbrNr}{\abbrnoteonce{nr}{Nr.}{Nummer}{款、目或號，依語境}}
\newcommand{\abbrProf}{\abbrnoteonce{prof}{Prof.}{Professor}{教授}}
\newcommand{\abbrRdnr}{\abbrnoteonce{rdnr}{Rdnr.}{Randnummer}{邊碼；段碼}}
\newcommand{\abbrRdnrn}{\abbrnoteonce{rdnrn}{Rdnrn.}{Randnummern}{邊碼；段碼}}
\newcommand{\abbrRGBl}{\abbrnoteonce{rgbl}{RGBl.}{Reichsgesetzblatt}{帝國法律公報}}
\newcommand{\abbrRGSt}{\abbrnoteonce{rgst}{RGSt}{Entscheidungen des Reichsgerichts in Strafsachen}{帝國法院刑事判決彙編}}
\newcommand{\abbrRspr}{\abbrnoteonce{rspr}{Rspr.}{Rechtsprechung}{判例；司法實務}}
\newcommand{\abbrRVO}{\abbrnoteonce{rvo}{RVO}{Reichsversicherungsordnung}{帝國保險法}}
\newcommand{\abbrS}{\abbrnoteonce{s}{S.}{Seite}{頁}}
\newcommand{\abbrSFHG}{\abbrnoteonce{sfhg}{SFHG}{Schwangeren- und Familienhilfegesetz}{孕婦及家庭扶助法}}
\newcommand{\abbrSGBV}{\abbrnoteonce{sgbv}{SGB V}{Fünftes Buch Sozialgesetzbuch}{社會法典第五編}}
\newcommand{\abbrStAG}{\abbrnoteonce{staeg}{StÄG}{Strafrechtsänderungsgesetz}{刑法修正法}}
\newcommand{\abbrStenBer}{\abbrnoteonce{stenber}{StenBer.}{Stenographischer Bericht}{速記紀錄}}
\newcommand{\abbrStGB}{\abbrnoteonce{stgb}{StGB}{Strafgesetzbuch}{刑法}}
\newcommand{\abbrStREG}{\abbrnoteonce{streg}{StREG}{Strafrechtsreform-Ergänzungsgesetz}{刑法改革補充法}}
\newcommand{\abbrStrRG}{\abbrnoteonce{strrg}{StrRG}{Strafrechtsreformgesetz}{刑法改革法}}
\newcommand{\abbrUS}{\abbrnoteonce{us}{U.S.}{United States Reports}{美國最高法院判決彙編}}
\newcommand{\abbrVgl}{\abbrnoteonce{vgl}{vgl.}{vergleiche}{參見}}
\newcommand{\abbrVVDStRL}{\abbrnoteonce{vvdstrl}{VVDStRL}{Veröffentlichungen der Vereinigung der Deutschen Staatsrechtslehrer}{德國公法教師協會論文集}}
\newcommand{\abbrWP}{\abbrnoteonce{wp}{WP.}{Wahlperiode}{屆期}}
\newcommand{\abbrWp}{\abbrnoteonce{wp}{Wp.}{Wahlperiode}{屆期}}
\newcommand{\abbrZStrW}{\abbrnoteonce{zstrw}{ZStrW}{Zeitschrift für die gesamte Strafrechtswissenschaft}{整體刑法學期刊}}
\newcommand{\recordsourcerandstart}{%
  \ifrandnumbers
    \setcounter{lastsourcerandstart}{\value{randnumber}}%
    \global\lastsourcehasrandnumberstrue
  \else
    \global\lastsourcehasrandnumbersfalse
  \fi
}
\newlength{\biparagraphskip}
\setlength{\biparagraphskip}{0.52em}
\newlength{\lawboxblockbeforeskip}
\setlength{\lawboxblockbeforeskip}{0.72em}
\newlength{\lawboxblockafterskip}
\setlength{\lawboxblockafterskip}{0.92em}
\newif\iflawboxpartendedblock
\lawboxpartendedblockfalse
\newcommand{\sourceparstyle}{\footnotesize\color{sourceink}\setstretch{1.12}\RaggedRight\emergencystretch=3em\setlength{\parskip}{0.42em}\obeylines\selectfont}
\newcommand{\transparstyle}{\songfont\color{caseink}\setstretch{1.12}\setlength{\parindent}{0pt}\setlength{\parskip}{0.42em}}
\newcommand{\bilingualpairskip}{%
  \par\addvspace{\biparagraphskip}%
  \switchcolumn
  \par\addvspace{\biparagraphskip}%
  \switchcolumn*%
  \global\intranslationfalse
}
\newcommand{\bilingualboxpairskip}{%
  \par
  \switchcolumn
  \par
  \switchcolumn*%
  \global\intranslationfalse
}
\newcommand{\bilinguallawboxblockskip}{%
  \switchcolumn*%
  \par\addvspace{\lawboxblockafterskip}%
  \switchcolumn
  \par\addvspace{\lawboxblockafterskip}%
  \switchcolumn*%
  \global\intranslationfalse
}
\newcommand{\bikeepwithnext}[1][4]{%
  \ifbilingualcolumns
    \syncleft
    \Needspace{#1\baselineskip}%
    \switchcolumn
    \Needspace{#1\baselineskip}%
    \switchcolumn*%
    \global\intranslationfalse
  \else
    \Needspace{#1\baselineskip}%
  \fi
}
\NewEnviron{sourcepara}[1][]{
  \recordsourcerandstart
  \ifshoworiginal
    \ifbilingualcolumns
      \ifintranslation\switchcolumn*\global\intranslationfalse\fi
    \else
      \Needspace{5\baselineskip}
    \fi
    \ifbilingualcolumns\else\par\addvspace{0.35em}\fi
    {\sourceparstyle
    \noindent\BODY\par}
    \ifbilingualcolumns\else\addvspace{0.25em}\fi
    \ifbilingualcolumns
      \switchcolumn
      \global\intranslationtrue
    \fi
  \else
    \ifrandnumbers
      \begingroup
        \setbox0=\vbox{\hsize=\linewidth\sourceparstyle\noindent\BODY\par}%
      \endgroup
    \fi
  \fi
}
\NewEnviron{transpara}{
  \ifshowtranslation
    \setcounter{randnumberlocal}{\value{lastsourcerandstart}}%
    \ifbilingualcolumns
      \ifintranslation\else\switchcolumn\global\intranslationtrue\fi
      {\transparstyle\setlength{\leftskip}{0.55em}\setlength{\rightskip}{0.15em plus 1em}\addtolength{\linewidth}{-0.55em}\BODY\par}
      \switchcolumn*
      \bilingualpairskip
    \else
      {\transparstyle\BODY\par}
    \fi
  \else
    \ifbilingualcolumns
      \ifintranslation\switchcolumn*\global\intranslationfalse\fi
    \fi
  \fi
}
\NewEnviron{sourceboxpara}[1][]{
  \global\lawboxpartendedblockfalse
  \recordsourcerandstart
  \ifshoworiginal
    \ifbilingualcolumns
      \ifintranslation\switchcolumn*\global\intranslationfalse\fi
    \else
      \Needspace{3\baselineskip}
    \fi
    {\sourceparstyle
    \BODY\par}
    \ifbilingualcolumns\else
      \iflawboxpartendedblock\addvspace{\lawboxblockafterskip}\fi
    \fi
    \ifbilingualcolumns
      \switchcolumn
      \global\intranslationtrue
    \fi
  \fi
}
\NewEnviron{transboxpara}{
  \global\lawboxpartendedblockfalse
  \ifshowtranslation
    \setcounter{randnumberlocal}{\value{lastsourcerandstart}}%
    \ifbilingualcolumns
      \ifintranslation\else\switchcolumn\global\intranslationtrue\fi
      {\transparstyle\BODY\par}
      \iflawboxpartendedblock
        \bilinguallawboxblockskip
      \else
        \switchcolumn*
        \bilingualboxpairskip
      \fi
    \else
      {\transparstyle\BODY\par}
      \iflawboxpartendedblock\addvspace{\lawboxblockafterskip}\fi
    \fi
  \else
    \ifbilingualcolumns
      \ifintranslation\switchcolumn*\global\intranslationfalse\fi
    \fi
  \fi
}
\newcommand{\leitsatznum}[1]{\noindent\hangindent=3.0em\hangafter=1\makebox[2.25em][r]{\bfseries #1.}\hspace{0.55em}\ignorespaces}
\newcommand{\syncleft}{\ifbilingualcolumns\ifintranslation\switchcolumn*\global\intranslationfalse\fi\fi}
\pretocmd{\section}{\syncleft\clearpage}{}{}
\pretocmd{\subsection}{\syncleft}{}{}
\pretocmd{\subsubsection}{\syncleft}{}{}
\newcommand{\biheadingfont}{\songfont}
\newcommand{\bitocentry}[2]{%
  \ifshoworiginal
    \ifshowtranslation
      \ifstrequal{#1}{#2}{#1}{#1\space #2}%
    \else
      #1%
    \fi
  \else
    #2%
  \fi
}
\pdfstringdefDisableCommands{%
  \def\bitocentry#1#2{#1}%
}
\newcommand{\biheading}[7]{%
  \syncleft#1\bikeepwithnext[#3]\phantomsection\addcontentsline{toc}{#2}{\bitocentry{#6}{#7}}%
  \ifbilingualcolumns
    {#4 #6\par}%
    \switchcolumn
    {#4\biheadingfont #7\par}%
    \switchcolumn*%
    \global\intranslationfalse
    \par\addvspace{#5}%
    \switchcolumn
    \par\addvspace{#5}%
    \switchcolumn*%
  \else
    \ifshoworiginal{#4 #6\par}\fi
    \ifshowtranslation{#4\biheadingfont #7\par}\fi
    \addvspace{#5}%
  \fi
}
\newcommand{\termsectionheading}[1]{%
  \par\Needspace{9\baselineskip}%
  \vspace{0.65em}%
  {\songfont\bfseries #1\par}%
  \vspace{0.2em}%
}
% 章節換頁：
%   雙語 paracol 欄內不要用 \clearpage；它會在同步兩欄時吐出只有欄線/頁碼的空頁。
%   \newpage 足以讓下一個 section 起新頁，又不會強制清空所有待排材料。
%   單語模式不在 paracol 內，仍可用 \clearpage。
\newcommand{\bisectionpagebreak}{\ifbilingualcolumns\newpage\else\clearpage\fi}
\newcommand{\bisection}[2]{%
  \biheading{\iffirstbilingualsection\global\firstbilingualsectionfalse\else\bisectionpagebreak\fi}{section}{6}{\centering\Large\bfseries}{0.8em}{#1}{#2}%
}
\newcommand{\bisubsection}[2]{%
  \biheading{}{subsection}{5}{\large\bfseries}{0.55em}{#1}{#2}%
}
\newcommand{\bisubsubsection}[2]{%
  \biheading{}{subsubsection}{4}{\normalsize\bfseries}{0.45em}{#1}{#2}%
}
\newcommand{\bicompositeheading}[4]{%
  \biheading{}{subsection}{5}{\large\bfseries}{0.16em}{#1}{#2}%
  \biheading{}{subsubsection}{3}{\normalsize\bfseries}{0.45em}{#3}{#4}%
}
\newif\iflawboxfirstitem
\lawboxfirstitemfalse
\newcommand{\lawboxprespace}[1]{%
  \par
  \iflawboxfirstitem
    \global\lawboxfirstitemfalse
  \else
    \addvspace{#1}%
  \fi
}
\newtcolorbox{lawboxinner}{
  caseboxbase,
  colframe=sourceline!85!black,
  colback=white,
  left=1.05em,
  right=1.05em,
  top=0.62em,
  bottom=0.62em,
  before upper={\global\inlawboxtrue\global\lawboxfirstitemtrue\global\lawboxhaspendingrnfalse\ifintranslation\lawboxtransfont\else\lawboxsourcefont\fi\RaggedRight\setlength{\parindent}{0pt}\setlength{\parskip}{0pt}\ifintranslation\setstretch{\lawboxtransstretch}\else\setstretch{\lawboxsourcestretch}\fi},
  after upper={\global\lawboxhaspendingrnfalse\global\inlawboxfalse}
}
\tcbset{
  lawboxpartbase/.style={
    enhanced,
    breakable=false,
    width=0.985\linewidth,
    colback=white,
    frame hidden,
    boxrule=0pt,
    arc=0pt,
    left=1.05em,
    right=1.05em,
    boxsep=1.5mm,
    before skip=0pt,
    after skip=0pt,
    before upper={\global\inlawboxtrue\global\lawboxfirstitemtrue\global\lawboxhaspendingrnfalse\ifintranslation\lawboxtransfont\else\lawboxsourcefont\fi\RaggedRight\setlength{\parindent}{0pt}\setlength{\parskip}{0pt}\ifintranslation\setstretch{\lawboxtransstretch}\else\setstretch{\lawboxsourcestretch}\fi},
    after upper={\global\lawboxhaspendingrnfalse\global\inlawboxfalse},
    borderline west={0.28pt}{0pt}{sourceline!85!black},
    borderline east={0.28pt}{0pt}{sourceline!85!black},
  },
  lawboxpartfirst/.style={
    lawboxpartbase,
    top=0.62em,
    bottom=0.04em,
    before skip=\lawboxblockbeforeskip,
    borderline north={0.28pt}{0pt}{sourceline!85!black},
  },
  lawboxpartmiddle/.style={
    lawboxpartbase,
    top=0.04em,
    bottom=0.04em,
  },
  lawboxpartlast/.style={
    lawboxpartbase,
    top=0.04em,
    bottom=0.62em,
    after skip=0pt,
    borderline south={0.28pt}{0pt}{sourceline!85!black},
  }
}
\newtcolorbox{lawboxpartinner}[1]{#1}
\newtcolorbox{termboxinner}{
  caseboxbase,
  colframe=noteblue!70!black,
  colback=casepaper,
  before upper={\songfont\setlength{\parindent}{0pt}\setlength{\parskip}{0.35em}}
}
\newtcolorbox{deboxinner}{
  caseboxbase,
  colframe=notegray!72!black,
  colback=white,
  left=1.05em,
  right=1.05em,
  top=0.62em,
  bottom=0.62em,
  before upper={\global\inlawboxtrue\global\lawboxfirstitemtrue\global\lawboxhaspendingrnfalse\small\setlength{\parindent}{0pt}\setlength{\parskip}{0.35em}},
  after upper={\global\lawboxhaspendingrnfalse\global\inlawboxfalse}
}
\NewEnviron{lawbox}{
  \begin{lawboxinner}\BODY\end{lawboxinner}
}
\NewEnviron{lawboxpart}[2]{
  \begin{lawboxpartinner}{lawboxpart#1,equal height group=#2}\BODY\end{lawboxpartinner}%
  \ifstrequal{#1}{last}{\global\lawboxpartendedblocktrue}{}%
}
\NewEnviron{termbox}{
  \par\addvspace{0.55em}
  {\color{noteblue!70!black}\rule{\linewidth}{0.28pt}}\par
  \begingroup
    \songfont\small\color{caseink}
    \setlength{\parindent}{0pt}
    \setlength{\parskip}{0.24em}
    \BODY
    \par
  \endgroup
  {\color{noteblue!70!black}\rule{\linewidth}{0.28pt}}\par
  \addvspace{0.55em}
}
\NewEnviron{debox}{
  \begin{deboxinner}\BODY\end{deboxinner}
}
\providecommand{\paragraphmark}[1]{}
\newcommand{\lawboxtitleafter}{\addvspace{0.06em}}
\newcommand{\lawtitle}[1]{\lawboxprespace{0.22em}{\lawboxtransfont\bfseries\lawboxuserandnumber #1}\par\lawboxtitleafter}
\newcommand{\absno}[2]{\lawboxprespace{0.04em}\noindent\lawboxuserandnumber\hangindent=2.6em\hangafter=1\makebox[2.6em][l]{(#1)}#2\par}
\newcommand{\nrno}[2]{\lawboxprespace{0pt}\noindent\lawboxuserandnumber\hspace*{2.6em}\hangindent=4.4em\hangafter=1\makebox[1.8em][l]{#1.}#2\par}
\newcommand{\letterno}[2]{\lawboxprespace{0pt}\noindent\lawboxuserandnumber\hspace*{4.4em}\hangindent=6.0em\hangafter=1\makebox[1.6em][l]{#1)}#2\par}
\newcommand{\sourcelawtitle}[1]{\lawboxprespace{0.22em}{\lawboxsourcefont\bfseries\lawboxuserandnumber #1}\par\lawboxtitleafter}
\newcommand{\sourcelawsubtitle}[1]{\lawboxprespace{0.04em}{\lawboxsourcefont\bfseries\lawboxuserandnumber #1}\par\lawboxtitleafter}
\newcommand{\sourceabsno}[2]{\lawboxprespace{0.04em}\noindent\lawboxuserandnumber\hangindent=2.45em\hangafter=1\makebox[2.45em][l]{(#1)}#2\par}
\newcommand{\sourcenrno}[2]{\lawboxprespace{0pt}\noindent\lawboxuserandnumber\hspace*{2.45em}\hangindent=4.25em\hangafter=1\makebox[1.8em][l]{#1.}#2\par}
\newcommand{\sourceletterno}[2]{\lawboxprespace{0pt}\noindent\lawboxuserandnumber\hspace*{4.25em}\hangindent=5.85em\hangafter=1\makebox[1.6em][l]{#1)}#2\par}
\newcommand{\sourcequotepara}[1]{\lawboxprespace{0.04em}\noindent\lawboxuserandnumber\hangindent=1.2em\hangafter=1 #1\par}
\newcommand{\sourcelawline}[1]{\lawboxprespace{0.04em}\noindent\lawboxuserandnumber #1\par}
\newcommand{\lawline}[1]{\lawboxprespace{0.04em}\noindent\lawboxuserandnumber #1\par}
\newcommand{\pendingtranslation}{\par{\songfont\footnotesize\color{pendingink}（中文譯文待補。）}\par}
\newcommand{\tocpart}[1]{\par\addvspace{0.52em}{\footnotesize\bfseries\color{pendingink}#1\par}\addvspace{0.16em}}
\newcommand{\tocdashline}{\par\addvspace{0.48em}\noindent{\color{notegray}\xleaders\hbox to 0.82em{\hfil\rule{0.42em}{0.28pt}\hfil}\hskip\linewidth}\par\addvspace{0.38em}}
% \tocpanel{font-cmd}{content}：將目錄置中於所在欄寬（雙語欄適用）。
% A3 橫式使用 0.58\linewidth，A4 橫式使用 0.84\linewidth，
% 使兩種紙張下目錄欄內容實際寬度相近（均約 100–105 mm）。
\newcommand{\tocpanel}[2]{%
  \makebox[\linewidth][c]{%
    \ifx\bilingualpaper\bilingualpaperaThree
      \begin{minipage}{0.58\linewidth}%
    \else
      \begin{minipage}{0.84\linewidth}%
    \fi
    \toccolstyle
    #1#2%
    \end{minipage}%
  }%
}
% ===== 首頁簡目錄的兩層 description list 縮排 =====
% enumitem 的 description list 可把每一列拆成：
%   [label]  labelsep  body text
% 這裡的「起點」都以所在 minipage/欄位的左緣為 0。
%
% 核心規則：
%   labelindent = label 欄位從左緣開始的位置。
%   labelwidth  = label 欄位本身寬度；標籤太長（如 Abw. Meinung）就加大它。
%   labelsep    = label 欄位右緣到正文的距離。
%   leftmargin  = 正文 body text 的起點。判斷層級視覺時主要看這個。
%
% 所以：
%   想讓「Begründung / 判決理由」整列正文往右或往左：改 tocitems 的 leftmargin。
%   想讓「A. Verfahren und Gegenstand / A. 程序與審查標的」正文往右或往左：
%     改 tocsubitems 的 leftmargin。
%   想讓 A./B./C. 這些標籤本身往右或往左，但正文不動：改 tocsubitems 的 labelindent。
%   想讓 A./B./C. 與正文間距變大或變小：改 tocsubitems 的 labelsep。
%
% 層級原則：
%   tocsubitems.leftmargin 必須大於 tocitems.leftmargin，
%   否則子層級正文會看起來比「Gründe / 理由」還靠左。
\newlist{tocitems}{description}{1}
\setlist[tocitems]{%
  labelindent=0pt,    % 第一層 label 欄位起點；0pt 表示貼齊目錄欄左緣。
  leftmargin=2.54cm,  % 第一層正文起點；例如 Begründung / 判決理由 的左緣。
  labelwidth=2.16cm,  % 第一層 label 欄位寬度；需容納 Leitsätze、Abw. Meinung、不同意見等。
  labelsep=0.32cm,    % 第一層 label 與正文的水平間距；只改間距，不改正文起點。
  itemsep=0.28em,     % 第一層各項之間的垂直距離。
  topsep=0.20em,      % 第一層 list 上下與前後內容的垂直距離。
  parsep=0pt,         % 同一 item 內段落之間的距離；目錄項通常不需要。
  partopsep=0pt,      % list 若緊接段落開始時的額外距離；關閉以免目錄高度飄動。
  align=left,         % label 欄內左對齊；長短標籤不靠右貼齊。
  font=\normalfont\bfseries % label 的字型；正文不受這行影響。
}
\newlist{tocsubitems}{description}{1}
\setlist[tocsubitems]{%
  labelindent=1cm,    % 第二層 label 起點；只控制 A./B./C. 本身的位置。
  leftmargin=3.00cm,  % 第二層正文起點；必須大於 2.54cm，才會比 Begründung / 判決理由更右。
  labelwidth=0.5cm,   % 第二層 label 欄寬；A./B./C. 很短，可小於第一層。
  labelsep=1.5cm,    % 第二層 label 與正文間距；A. 與 Verfahren / 程序 的距離看這裡。
  itemsep=0.12em,     % 第二層各項較密，避免 A.–E. 佔太高。
  topsep=0.04em,      % 第二層 list 與 Gründe / 理由之間的垂直距離。
  parsep=0pt,         % 同一第二層 item 內段落距離；目錄項通常不需要。
  partopsep=0pt,      % 關閉額外段落起始距離，避免版面不穩。
  align=left,         % A./B./C. label 欄內左對齊。
  font=\normalfont\bfseries, % A./B./C. label 加粗；正文不受這行影響。
  before={}           % 不在第二層 list 前自動插入額外內容。
}
\newlist{termitems}{description}{1}
\ifbilingualcolumns
  \setlist[termitems]{%
    leftmargin=5.2cm,
    labelwidth=4.8cm,
    labelsep=0.35cm,
    itemsep=0.16em,
    topsep=0.2em,
    parsep=0pt,
    partopsep=0pt,
    font=\normalfont\bfseries,
    style=nextline
  }
\else
  \setlist[termitems]{%
    leftmargin=0pt,
    labelwidth=0pt,
    labelsep=0pt,
    itemsep=0.24em,
    topsep=0.2em,
    parsep=0pt,
    partopsep=0pt,
    font=\normalfont\bfseries,
    style=nextline
  }
\fi
\newlist{abbritems}{description}{1}
\setlist[abbritems]{%
  leftmargin=2.38cm,
  labelwidth=2.02cm,
  labelsep=0.28cm,
  itemsep=0.16em,
  topsep=0.2em,
  parsep=0pt,
  partopsep=0pt,
  align=left,
  font=\normalfont\bfseries
}
\ifbilingualcolumns
  \newenvironment{termcolumns}{\begin{multicols}{2}}{\end{multicols}}
\else
  \newenvironment{termcolumns}{}{}
\fi

% ===== 編者註腳開關 =====
% 一般版預設不顯示編者註，避免正文腳註過密。
% 若開啟 \classroommodeon，GG/條文編者註仍會自動顯示，無需另開本開關。
% 若一般版也要顯示編者註，可在單案檔 \input 後加 \showeditornotestrue。
\newif\ifshoweditornotes
\showeditornotesfalse

% ===== 目錄排版輔助 =====
\newcommand{\toccolstyle}{\raggedright\arraybackslash}
\newcommand{\tocsingleindent}{\leftskip=1.45cm\rightskip=1.2cm}

% ===== 行內引文環境（德文原文與中文譯文各一，帶左側灰色邊線）=====
\newcommand{\quoteparbreak}{\par\addvspace{0.48em}}
\newcommand{\sourcequoteinner}[1]{%
  \begin{tcolorbox}[
    enhanced, breakable, width=\dimexpr\linewidth-0.8em\relax, frame hidden,
    colback=white, coltext=caseink, boxrule=0pt,
    borderline west={0.55pt}{0.88em}{notegray},
    left=1.78em, right=0.72em, top=0.18em, bottom=0.18em,
    boxsep=0pt, before skip=0.24em, after skip=0.24em,
    before upper={\normalfont\footnotesize\color{caseink}\setlength{\parindent}{0pt}\setlength{\parskip}{0.34em}\raggedright\setlength{\rightskip}{0pt plus 1em}}
  ]%
  #1\par
  \end{tcolorbox}%
}
\newcommand{\transquoteinner}[1]{%
  \begin{tcolorbox}[
    enhanced, breakable, width=\dimexpr\linewidth-0.8em\relax, frame hidden,
    colback=white, coltext=caseink, boxrule=0pt,
    borderline west={0.55pt}{0.88em}{notegray},
    left=1.78em, right=0.72em, top=0.18em, bottom=0.18em,
    boxsep=0pt, before skip=0.24em, after skip=0.24em,
    before upper={\songfont\small\color{caseink}\setlength{\parindent}{0pt}\setlength{\parskip}{0.34em}\raggedright\setlength{\rightskip}{0pt plus 1em}}
  ]%
  #1\par
  \end{tcolorbox}%
}

% ===== 大綱（Gliederung）Rn 輔助命令 =====
\newcommand{\outrn}[2]{\enspace{\normalfont\footnotesize\color{pendingink}(Rn\,#1–#2)}}
\newcommand{\outrnsingle}[1]{\enspace{\normalfont\footnotesize\color{pendingink}(Rn\,#1)}}
\newcommand{\outrnzh}[2]{\enspace{\normalfont\footnotesize\color{pendingink}（第\,#1–#2\,段）}}
\newcommand{\outrnzhs}[1]{\enspace{\normalfont\footnotesize\color{pendingink}（第\,#1\,段）}}
% ===== 大綱雙語對照表格輔助命令（三欄：段碼 │ 德文 │ 中文）=====
% 欄 1：段碼（由欄定義自動格式化：小字、等寬、右對齊、pendingink）
% 欄 2：德文標籤＋正文
% 欄 3：中文標籤＋正文
%
% 無段碼的列（\omainrow、\osecrow），欄 1 以 & 空置。
% 有段碼的列（\osecrowrn、\osubrow、\ointrow），#1 為預格式化字串
%   例：\osubrow{2–49}{I.}{...}{I.}{...}
%       \ointrow{107}{...}{...}
%
% \opartrow：段組標題（橫跨三欄）
\newcommand{\opartrow}[2]{%
  \noalign{\vspace{0.30em}}%
  \multicolumn{3}{@{}l@{}}{%
    \footnotesize\bfseries\color{pendingink}#1\hfill\songfont#2%
  }\\[0.06em]%
}
% \odashrow：虛線分隔（橫跨三欄）
\newcommand{\odashrow}{%
  \noalign{\vspace{0.28em}}%
  \multicolumn{3}{@{}l@{}}{\color{notegray}%
    \xleaders\hbox to 0.82em{\hfil\rule{0.42em}{0.28pt}\hfil}\hskip 0.88\linewidth}%
  \\[0.24em]%
}
% \odissentpart：不同意見書區段標頭。比一般 \opartrow 更像附錄性轉場。
\newcommand{\odissentpart}[2]{%
  \noalign{\vspace{0.42em}}%
  \multicolumn{3}{@{}p{0.88\textwidth}@{}}{%
    \color{notegray}\rule{\linewidth}{0.18pt}%
  }\\[-0.18em]%
  \multicolumn{3}{@{}p{0.88\textwidth}@{}}{%
    \makebox[\linewidth][s]{%
      \footnotesize\bfseries\color{pendingink}#1%
      \hfill{\songfont #2}%
    }%
  }\\[-0.02em]%
}
% \odissentopinion：不同意見書的署名與一行摘要，右欄保留段碼範圍。
\newcommand{\odissentopinion}[5]{%
  \footnotesize\hangindent=0.52cm\hangafter=1%
    {\bfseries\color{caseink}#2}\enspace{\color{notegray}\textperiodcentered}\enspace\color{caseink}#3%
  &\footnotesize\hangindent=0.52cm\hangafter=1%
    {\songfont\bfseries\color{caseink}#4}\enspace{\color{notegray}\textperiodcentered}\enspace\songfont\color{caseink}#5%
  &#1%
  \\[0.12em]%
}
% \omainrow：頂層項目，無段碼（Leitsätze、Urteil、Gründe …）
\newcommand{\omainrow}[4]{%
  \footnotesize\hangindent=1.92cm\hangafter=1%
    \makebox[1.92cm][l]{\bfseries\color{caseink}#1}\color{caseink}#2%
  &\footnotesize\hangindent=1.92cm\hangafter=1%
    \makebox[1.92cm][l]{\bfseries\color{caseink}#3}\songfont\color{caseink}#4%
  &%
  \\[0.15em]%
}
% \omainrowrn：頂層項目，有段碼（不同意見等標籤寬於 \osecrow 框者）
\newcommand{\omainrowrn}[5]{%
  \footnotesize\hangindent=1.92cm\hangafter=1%
    \makebox[1.92cm][l]{\bfseries\color{caseink}#2}\color{caseink}#3%
  &\footnotesize\hangindent=1.92cm\hangafter=1%
    \makebox[1.92cm][l]{\bfseries\color{caseink}#4}\songfont\color{caseink}#5%
  &#1%
  \\[0.15em]%
}
% \osecrow：一級子項，無段碼（A.、C.、D. 等有子項者）
\newcommand{\osecrow}[4]{%
  \footnotesize\hspace{1.10cm}\hangindent=1.98cm\hangafter=1%
    \makebox[0.86cm][l]{\bfseries\color{caseink}#1}\color{caseink}#2%
  &\footnotesize\hspace{1.10cm}\hangindent=1.98cm\hangafter=1%
    \makebox[0.86cm][l]{\bfseries\color{caseink}#3}\songfont\color{caseink}#4%
  &%
  \\[0.11em]%
}
% \osecrowrn：一級子項，有段碼（B.、E. 等完整段落範圍者）
% #1=段碼字串（如 101–106）, #2=DE標籤, #3=DE正文, #4=ZH標籤, #5=ZH正文
\newcommand{\osecrowrn}[5]{%
  \footnotesize\hspace{1.10cm}\hangindent=1.98cm\hangafter=1%
    \makebox[0.86cm][l]{\bfseries\color{caseink}#2}\color{caseink}#3%
  &\footnotesize\hspace{1.10cm}\hangindent=1.98cm\hangafter=1%
    \makebox[0.86cm][l]{\bfseries\color{caseink}#4}\songfont\color{caseink}#5%
  &#1%
  \\[0.11em]%
}
% \osubrow：二級子項，有段碼範圍（I.、II.、A.I. …）
% #1=段碼字串, #2=DE標籤, #3=DE正文, #4=ZH標籤, #5=ZH正文
\newcommand{\osubrow}[5]{%
  \footnotesize\hspace{1.40cm}\hangindent=2.30cm\hangafter=1%
    \makebox[0.90cm][l]{\bfseries\color{caseink}#2}\color{caseink}#3%
  &\footnotesize\hspace{1.40cm}\hangindent=2.30cm\hangafter=1%
    \makebox[0.90cm][l]{\bfseries\color{caseink}#4}\songfont\color{caseink}#5%
  &#1%
  \\[0.09em]%
}
% \ointrow：引入段落，有單一段碼，無標籤
% #1=段碼字串, #2=DE正文, #3=ZH正文
\newcommand{\ointrow}[3]{%
  \footnotesize\hspace{0.52cm}\hangindent=0.52cm\hangafter=1\color{caseink}#2%
  &\footnotesize\hspace{0.52cm}\hangindent=0.52cm\hangafter=1\songfont\color{caseink}#3%
  &#1%
  \\[0.09em]%
}
% \osubsubrow：三級子項（1.、2.、3.）
% #1=段碼字串, #2=DE標籤, #3=DE正文, #4=ZH標籤, #5=ZH正文
\newcommand{\osubsubrow}[5]{%
  \footnotesize\hspace{1.80cm}\hangindent=2.48cm\hangafter=1%
    \makebox[0.68cm][l]{\bfseries\color{caseink}#2}\color{caseink}#3%
  &\footnotesize\hspace{1.80cm}\hangindent=2.48cm\hangafter=1%
    \makebox[0.68cm][l]{\bfseries\color{caseink}#4}\songfont\color{caseink}#5%
  &#1%
  \\[0.08em]%
}
% \osubsubsubrow：四級子項（a）、b））
% #1=段碼字串, #2=DE標籤, #3=DE正文, #4=ZH標籤, #5=ZH正文
\newcommand{\osubsubsubrow}[5]{%
  \footnotesize\hspace{2.10cm}\hangindent=2.68cm\hangafter=1%
    \makebox[0.58cm][l]{\bfseries\color{caseink}#2}\color{caseink}#3%
  &\footnotesize\hspace{2.10cm}\hangindent=2.68cm\hangafter=1%
    \makebox[0.58cm][l]{\bfseries\color{caseink}#4}\songfont\color{caseink}#5%
  &#1%
  \\[0.07em]%
}
% \osubsubsubsubrow：五級子項（aa）、bb））
% 標籤框 0.62cm：足以容納「aa)」在 footnotesize bold 下（約 0.50cm）並留有間距
% #1=段碼字串, #2=DE標籤, #3=DE正文, #4=ZH標籤, #5=ZH正文
\newcommand{\osubsubsubsubrow}[5]{%
  \footnotesize\hspace{2.38cm}\hangindent=3.06cm\hangafter=1%
    \makebox[0.68cm][l]{\bfseries\color{caseink}#2}\color{caseink}#3%
  &\footnotesize\hspace{2.38cm}\hangindent=3.06cm\hangafter=1%
    \makebox[0.68cm][l]{\bfseries\color{caseink}#4}\songfont\color{caseink}#5%
  &#1%
  \\[0.06em]%
}
% \outlinepage：雙語對照大綱頁包裝環境（longtable 自動跨頁）
% 三欄：德文欄 + 中文欄 + 段碼欄（最右，不折行）
% 段碼欄寬度：1.80cm，足以容納最長段碼如「309–348」
% DE/ZH 欄寬：(0.87\textwidth - 2.08cm) / 2
%   驗算：2×欄 + 0.05\textwidth + 0.28cm gap + 1.80cm = 0.87tw - 2.08cm + 0.05tw + 2.08cm = 0.92tw ✓
\newenvironment{outlinepage}{%
  \vspace*{0.40cm}%
  \setlength{\LTleft}{0.04\textwidth}%
  \setlength{\LTright}{0.04\textwidth}%
  \setlength{\tabcolsep}{0pt}%
  % 大綱列距由 longtable 的 \arraystretch 與各 row macro 結尾的 \\[...em] 共同決定。
  % 這裡控制「每一列本身的表格 strut 高度」；數值越大，A./I./1. 各項之間越像有空行。
  % A3 原本用 0.62，視覺上沒有空行；A4 也沿用同值，避免切到 A4 後項目被撐開。
  % 若日後覺得大綱太擠，只微調這個值（例如 0.68），不要改各 \omainrow/\osubrow 的縮排。
  \renewcommand{\arraystretch}{0.62}%
  \begin{longtable}{%
    >{\vspace{0pt}\raggedright\arraybackslash}p{\dimexpr(0.87\textwidth-2.08cm)/2\relax}%
    @{\hspace{0.025\textwidth}}%
    !{\color{notegray}\vrule width 0.12pt}%
    @{\hspace{0.025\textwidth}}%
    >{\vspace{0pt}\raggedright\arraybackslash}p{\dimexpr(0.87\textwidth-2.08cm)/2\relax}%
    @{\hspace{0.28cm}}%
    >{\vspace{0pt}\footnotesize\normalfont\color{pendingink}\raggedright\arraybackslash}p{1.80cm}%
    @{}%
  }%
  % 首頁欄標籤：不要橫跨置中；分別放在德文欄與中文欄的實際欄位起點。
  \footnotesize\bfseries\color{sourceink}Gliederung%
  &\footnotesize\bfseries\songfont\color{sourceink}大綱%
  &%
  \\[0.36em]%
  \endfirsthead%
  % 續頁欄標籤：同樣分欄定位，以灰色細體區分；無需標注「續」。
  \footnotesize\color{pendingink}Gliederung%
  &\footnotesize\songfont\color{pendingink}大綱%
  &%
  \\[0.24em]%
  \endhead%
}{%
  \end{longtable}%
}
 載入。
% 不含：\documentclass、\documentversion、\ggversionde/zh、\tocde/\toczh。
% 日期與首頁版本列由本檔自動產生；各文件只需手動指定 \documentversion。

\usepackage{xeCJK}
\usepackage{fontspec}
\usepackage{geometry}
\usepackage{setspace}
\usepackage{hyperref}
\usepackage{xurl}
\usepackage{bookmark}
\usepackage{enumitem}
\usepackage{xcolor}
\usepackage{titlesec}
\usepackage{array}
\usepackage{longtable}
\usepackage{tcolorbox}
\usepackage{environ}
\usepackage{pdflscape}
\usepackage{etoolbox}
\usepackage{ragged2e}
\usepackage{paracol}
\usepackage{multicol}
\usepackage{needspace}
\usepackage{fancyhdr}
\usepackage{tikz}
\usepackage{zref-abspage}
\tcbuselibrary{breakable,skins}
\usetikzlibrary{calc}
\usepackage{pgfcalendar}

\hypersetup{
  unicode=true,
  bookmarksopen=true,
  bookmarksnumbered=false,
  hidelinks
}

% ===== 自動推導，不要手動改下面 =====
% （\docmode 與 \bilingualpaper 由各文件在 % bverfge_layout.tex
% 共用排版基礎設施，供 Schwangerschaftsabbruch I 與 II 共享。
% 由各文件以 \documentclass 後 \input{../../共用LaTeX/bverfge_layout} 載入。
% 不含：\documentclass、\documentversion、\ggversionde/zh、\tocde/\toczh。
% 日期與首頁版本列由本檔自動產生；各文件只需手動指定 \documentversion。

\usepackage{xeCJK}
\usepackage{fontspec}
\usepackage{geometry}
\usepackage{setspace}
\usepackage{hyperref}
\usepackage{xurl}
\usepackage{bookmark}
\usepackage{enumitem}
\usepackage{xcolor}
\usepackage{titlesec}
\usepackage{array}
\usepackage{longtable}
\usepackage{tcolorbox}
\usepackage{environ}
\usepackage{pdflscape}
\usepackage{etoolbox}
\usepackage{ragged2e}
\usepackage{paracol}
\usepackage{multicol}
\usepackage{needspace}
\usepackage{fancyhdr}
\usepackage{tikz}
\usepackage{zref-abspage}
\tcbuselibrary{breakable,skins}
\usetikzlibrary{calc}
\usepackage{pgfcalendar}

\hypersetup{
  unicode=true,
  bookmarksopen=true,
  bookmarksnumbered=false,
  hidelinks
}

% ===== 自動推導，不要手動改下面 =====
% （\docmode 與 \bilingualpaper 由各文件在 \input{bverfge_layout} 之前設定；
%   預設 chinese / a4）
\ifdefined\docmode\else\def\docmode{chinese}\fi
\ifdefined\bilingualpaper\else\def\bilingualpaper{a4}\fi
\newif\ifshoworiginal
\newif\ifshowtranslation
\newif\ifbilinguallandscape
\newif\ifbilingualcolumns
\newif\ifintranslation
\newif\iffirstbilingualsection

\showoriginalfalse
\showtranslationfalse
\bilinguallandscapefalse
\bilingualcolumnsfalse
\intranslationfalse
\firstbilingualsectionfalse

\def\modechinese{chinese}
\def\modegerman{german}
\def\modebilingual{bilingual}
\def\bilingualpaperaThree{a3}
\def\bilingualpaperaFour{a4}

\ifx\docmode\modechinese
  \showtranslationtrue
\fi

\ifx\docmode\modegerman
  \showoriginaltrue
\fi

\ifx\docmode\modebilingual
  \showoriginaltrue
  \showtranslationtrue
  \bilinguallandscapetrue
  \bilingualcolumnstrue
\fi

% 編排模式說明：
% chinese  → \showoriginalfalse  \showtranslationtrue  （A4 直式，純中文）
% german   → \showoriginaltrue   \showtranslationfalse （A4 直式，純德文）
% bilingual→ \showoriginaltrue   \showtranslationtrue  \bilingualcolumnstrue（A3/A4 橫式對照）

\ifbilingualcolumns
  \ifx\bilingualpaper\bilingualpaperaFour
    \geometry{
      a4paper,
      landscape,
      left=1.25cm,
      right=2.25cm,
      top=1.25cm,
      bottom=1.75cm,
      footskip=8.5mm,
      marginparwidth=0.75cm,
      marginparsep=0.25cm
    }
  \else
    \geometry{
      a3paper,
      landscape,
      left=1.55cm,
      right=2.75cm,
      top=1.55cm,
      bottom=2.05cm,
      footskip=9.5mm,
      marginparwidth=0.9cm,
      marginparsep=0.35cm
    }
  \fi
\else
\geometry{
  a4paper,
  portrait,
  left=2.2cm,
  right=2.6cm,
  top=2.0cm,
  bottom=2.0cm,
  marginparwidth=0.8cm,
  marginparsep=0.3cm
}
\fi
% ===== 時間戳基礎設施 =====
\newcount\stamptimeval
\newcount\stampminuteval
\newcommand{\stamphh}{%
  \stamptimeval=\time
  \divide\stamptimeval by 60
  \ifnum\stamptimeval<10 0\fi
  \the\stamptimeval}
\newcommand{\stampmm}{%
  \stamptimeval=\time
  \stampminuteval=\time
  \divide\stampminuteval by 60
  \multiply\stampminuteval by 60
  \advance\stamptimeval by -\stampminuteval
  \ifnum\stamptimeval<10 0\fi
  \the\stamptimeval}
\newcommand{\stampwkde}[1]{%
  \ifcase#1Montag\or Dienstag\or Mittwoch\or Donnerstag\or Freitag\or Samstag\or Sonntag\fi}
\newcommand{\stampwkzh}[1]{%
  \ifcase#1 一\or 二\or 三\or 四\or 五\or 六\or 日\fi}
\newcommand{\stampmonthde}[1]{%
  \ifcase#1\or Januar\or Februar\or März\or April\or Mai\or Juni\or
  Juli\or August\or September\or Oktober\or November\or Dezember\fi}
\newcount\stampjdval
\newcount\stampwdval
\pgfcalendardatetojulian{\the\year-\the\month-\the\day}{\stampjdval}
\pgfcalendarjuliantoweekday{\the\stampjdval}{\stampwdval}
\newcommand{\timestampzh}{%
  \songfont 最後修改：\the\year 年\the\month 月\the\day 日%
  % （星期\stampwkzh{\the\stampwdval}）%
  % \space\stamphh:\stampmm
  }

\newcommand{\documentdatezh}{\the\year 年 \the\month 月 \the\day 日}
\newcommand{\documentdatede}{\the\day. \stampmonthde{\the\month} \the\year}
\newcommand{\bverfgetranslationnote}{%
  {\songfont 翻譯說明：初譯使用 Claude Code 與 GPT 協助迭代；仍請以德文原文為準。}%
}
\newcommand{\bverfgecourseinfo}{%
  {\songfont 課程：114 學年度第 2 學期；憲法專題研究（四）；授課老師：湯德宗教授；導讀日期：115 年 6 月 1 日。}%
}
\newcommand{\bverfgeeditorinfo}{%
  {\songfont 編者：王逸帆（R10A21126），國立台灣大學法律系研究所碩士班。}%
}
\newcommand{\bverfgetitleversionline}{%
  \ifbilingualcolumns
    Fassung \documentversion\enspace\textperiodcentered\enspace Stand: \documentdatede
    \quad
    {\songfont 版本 \documentversion\enspace\textperiodcentered\enspace 修訂日期：\documentdatezh\par}%
  \else
    \ifshowtranslation
      {\songfont 版本 \documentversion\enspace\textperiodcentered\enspace 修訂日期：\documentdatezh\par}%
    \else
      Fassung \documentversion\enspace\textperiodcentered\enspace Stand: \documentdatede\par
    \fi
  \fi
}
\newcommand{\bverfgetitlecornerinfo}{%
  \ifshowtranslation
    \vspace{0.72em}%
    \noindent\hfill
    \begin{minipage}{0.66\textwidth}
      \raggedleft\tiny\color{pendingink}\songfont
      \bverfgecourseinfo\par
      \bverfgeeditorinfo
    \end{minipage}%
  \fi
}
\newcommand{\bverfgetitlemeta}{%
  \vfill
  \begin{center}%
    \footnotesize\color{pendingink}%
    \bverfgetitleversionline
    \ifshowtranslation
      \vspace{0.28em}{\scriptsize\bverfgetranslationnote\par}%
    \fi
  \end{center}%
  \bverfgetitlecornerinfo
}

\pagestyle{fancy}
\fancyhf{}
\renewcommand{\headrulewidth}{0pt}
\renewcommand{\footrulewidth}{0pt}
\fancyfoot[C]{\thepage}
\ifbilingualcolumns
  \fancyfoot[R]{\scriptsize\color{pendingink}\timestampzh}%
\else
  \ifshoworiginal
  \else
    \fancyfoot[R]{\scriptsize\color{pendingink}\timestampzh\hspace{0.45cm}}%
  \fi
\fi
\setmainfont{Times New Roman}
\setsansfont{Helvetica Neue}
\setmonofont{Menlo}
\IfFileExists{/Users/iw/Library/Fonts/kaiu.ttf}{
  \setCJKmainfont[Path=/Users/iw/Library/Fonts/,AutoFakeBold=4]{kaiu.ttf}
}{
  \IfFontExistsTF{標楷體}{
    \setCJKmainfont[AutoFakeBold=4]{標楷體}
  }{
    \setCJKmainfont{Songti TC}
  }
}
\IfFontExistsTF{Songti TC}{
  \setCJKfamilyfont{song}{Songti TC}
  \setCJKmonofont{Songti TC}
}{
  \setCJKfamilyfont{song}[Path=/Users/iw/Library/Fonts/,AutoFakeBold=4]{kaiu.ttf}
  \setCJKmonofont[Path=/Users/iw/Library/Fonts/,AutoFakeBold=4]{kaiu.ttf}
}
\newcommand{\songfont}{\CJKfamily{song}}

% ===== 封面／目錄專用打字機字體（僅限封面、目錄頁使用，不影響正文）=====
\IfFontExistsTF{Radio Newsman}{%
  \newfontfamily\bverfgetitlede{Radio Newsman}%
}{%
  \newfontfamily\bverfgetitlede{Courier New}%
}
\IfFontExistsTF{Erikas Farbband}{%
  \newfontfamily\bverfgebodyde{Erikas Farbband}%
}{%
  \let\bverfgebodyde\bverfgetitlede
}
\IfFontExistsTF{Maoken Heavy Labourer}{%
  \setCJKfamilyfont{bverfgetitlecjk}{Maoken Heavy Labourer}%
}{%
  \setCJKfamilyfont{bverfgetitlecjk}{Songti TC}%
}
\IfFontExistsTF{Huiwen-mincho}{%
  \setCJKfamilyfont{bverfgebodycjk}{Huiwen-mincho}%
}{%
  \setCJKfamilyfont{bverfgebodycjk}{Songti TC}%
}
\newcommand{\bverfgetitlezh}{\bverfgetitlede\CJKfamily{bverfgetitlecjk}}
\newcommand{\bverfgebodyzh}{\bverfgebodyde\CJKfamily{bverfgebodycjk}}

\setstretch{1.35}
\setlist{itemsep=0.2em, topsep=0.4em}
\setlength{\parindent}{0pt}
\setlength{\parskip}{0.45em}
\raggedbottom
\setcounter{secnumdepth}{0}
\setcounter{tocdepth}{2}

\newcommand{\lawboxsourcefont}{\footnotesize}
\newcommand{\lawboxtransfont}{\footnotesize}
\newcommand{\lawboxsourcestretch}{1.02}
\newcommand{\lawboxtransstretch}{0.92}

\titleformat{\section}{\centering\Large\bfseries\songfont}{\thesection}{1em}{}
\titleformat{\subsection}{\large\bfseries\songfont}{\thesubsection}{1em}{}
\titleformat{\subsubsection}{\normalsize\bfseries\songfont}{\thesubsubsection}{1em}{}
\titleformat{\paragraph}{\normalsize\bfseries\songfont}{\theparagraph}{1em}{}
\titlespacing*{\section}{0pt}{1.8em}{1.0em}
\titlespacing*{\subsection}{0pt}{1.2em}{0.6em}
\titlespacing*{\subsubsection}{0pt}{1.0em}{0.45em}
\titlespacing*{\paragraph}{0pt}{0.6em}{0.4em}

\definecolor{noteblue}{HTML}{A9C9F5}
\definecolor{noteteal}{HTML}{AEE1D6}
\definecolor{notegray}{HTML}{D7DADF}
\definecolor{caseink}{HTML}{000000}
\definecolor{casepaper}{HTML}{FCFEFD}
\definecolor{sourcepaper}{HTML}{FAFBF8}
\definecolor{sourceline}{HTML}{D7DED8}
\definecolor{sourceink}{HTML}{000000}
\definecolor{pendingink}{HTML}{000000}

\tcbset{
  caseboxbase/.style={
    enhanced,
    breakable,
    width=0.985\linewidth,
    colback=casepaper,
    coltitle=caseink,
    fonttitle=\bfseries\songfont,
    boxrule=0.28pt,
    arc=1.2mm,
    left=2.4mm,
    right=2.4mm,
    top=1.5mm,
    bottom=1.5mm,
    boxsep=1.5mm,
    before skip=0.65em,
    after skip=0.65em,
    before upper={\setlength{\parindent}{0pt}\setlength{\parskip}{0.35em}},
  }
}

\newcommand{\bilingualstart}{}
\newcommand{\bilingualstop}{}
\ifbilingualcolumns
  \renewcommand{\bilingualstart}{\small\setlength{\columnsep}{1.25cm}\setlength{\columnseprule}{0.12pt}\columnratio{0.5,0.5}\multicolumnfootnotes\flushbottom\sloppy\interlinepenalty=80\clubpenalty=800\widowpenalty=800\displaywidowpenalty=800\begin{paracol}{2}\global\intranslationfalse\global\firstbilingualsectiontrue{\footnotesize\color{sourceink}Deutsch\par}\switchcolumn{\footnotesize\songfont\color{sourceink}中文譯文\par}\switchcolumn*}
  \renewcommand{\bilingualstop}{\ifintranslation\switchcolumn*\global\intranslationfalse\fi\end{paracol}\clearpage\raggedbottom}
\fi
\newif\ifrandnumbers
\randnumbersfalse
\newcounter{randnumber}
\newcounter{randnumberlocal}
\newcounter{lastsourcerandstart}
\newif\iflastsourcehasrandnumbers
\lastsourcehasrandnumbersfalse
\newif\ifinlawbox
\inlawboxfalse
\newif\iflawboxhaspendingrn
\lawboxhaspendingrnfalse
\newcommand{\lawboxpendingrn}{}
\newcommand{\startrandnumbers}{\global\randnumberstrue\setcounter{randnumber}{1}\global\lastsourcehasrandnumbersfalse}
\newcommand{\stoprandnumbers}{\global\randnumbersfalse\global\lastsourcehasrandnumbersfalse}
\newlength{\randnumberwidth}
\setlength{\randnumberwidth}{0.95cm}
\newlength{\randnumberrightoffset}
\setlength{\randnumberrightoffset}{0.85cm}
\newlength{\randnumberlawboxrightoffset}
\setlength{\randnumberlawboxrightoffset}{1.40cm}
\newcommand{\randnumberbox}[1]{%
  \leavevmode
  \makebox[0pt][l]{%
    \smash{%
      \tikz[remember picture,overlay,baseline=(rnmark.base)]{%
        \coordinate (rnmark) at (0,0);%
        \ifinlawbox
          \path let \p1=(current page.east), \p2=(rnmark.base) in
            node[
              anchor=base east,
              inner sep=0pt,
              outer sep=0pt,
              text width=\randnumberwidth,
              align=right,
              text=pendingink,
              font=\normalfont\mdseries\scriptsize\ttfamily
            ] at (\x1-\randnumberlawboxrightoffset,\y2) {{\normalfont\mdseries\scriptsize\ttfamily #1}};%
        \else
          \path let \p1=(current page.east), \p2=(rnmark.base) in
            node[
              anchor=base east,
              inner sep=0pt,
              outer sep=0pt,
              text width=\randnumberwidth,
              align=right,
              text=pendingink,
              font=\normalfont\mdseries\scriptsize\ttfamily
            ] at (\x1-\randnumberrightoffset,\y2) {{\normalfont\mdseries\scriptsize\ttfamily #1}};%
        \fi
      }%
    }%
  }%
  \ignorespaces
}
\newcommand{\lawboxstorerandnumber}[1]{%
  \gdef\lawboxpendingrn{#1}%
  \global\lawboxhaspendingrntrue
  \ignorespaces
}
\newcommand{\lawboxuserandnumber}{%
  \iflawboxhaspendingrn
    \randnumberbox{\lawboxpendingrn}%
    \global\lawboxhaspendingrnfalse
  \fi
}
\newcommand{\sourcern}[1]{%
  \ifrandnumbers
    \ifshoworiginal
      \ifshowtranslation\else
        \ifinlawbox\lawboxstorerandnumber{#1}\else\randnumberbox{#1}\fi
      \fi
    \fi
  \fi
}
\newcommand{\transrn}[1]{%
  \ifrandnumbers
    \ifshowtranslation
      \ifinlawbox\lawboxstorerandnumber{#1}\else\randnumberbox{#1}\fi
    \fi
  \fi
}
% ===== 課堂模式：德文原文縮寫註解 =====
% 日常切換只改各判決檔的一行：
%   \classroommodeon        開啟課堂版；註解掉這行即回到一般版。
% 模板預設：
%   1. 課堂模式關閉。
%   2. 課堂模式開啟時，縮寫註解包含「德文全稱＋中文譯名」。
%   3. 課堂模式開啟時，GG/條文編者註仍會自動顯示。
% 進階微調才需要在單案檔覆寫：
%   \classroomabbrzhoff     縮寫註解僅顯示德文全稱。
%   \showeditornotestrue    一般版也顯示編者註。
% 註解策略：同一實際 PDF 頁面中，同一縮寫只在第一次出現處加註。
% 這裡用 zref-abspage 的絕對頁序，不用 \thepage，避免文件內頁碼重置時誤判。
\newif\ifclassroommode
\newif\ifclassroomabbrzh
\classroommodefalse
\classroomabbrzhtrue
\newcommand{\classroommodeon}{\classroommodetrue}
\newcommand{\classroommodeoff}{\classroommodefalse}
\newcommand{\classroomabbrzhon}{\classroomabbrzhtrue}
\newcommand{\classroomabbrzhoff}{\classroomabbrzhfalse}
\newcommand{\abbrnotelabel}{\emph{Abk.}: }
\newcommand{\abbrnote}[3]{%
  \ifclassroommode
    \ifshoworiginal
      \ifintranslation\else
        \footnote{\normalfont\footnotesize\abbrnotelabel #2\ifclassroomabbrzh；{\songfont #3}\fi}%
      \fi
    \fi
  \fi
}
\newcommand{\abbrnoteonce}[4]{%
  #2%
  \ifclassroommode
    \ifshoworiginal
      \ifintranslation\else
        \begingroup
          \edef\abbrcurrentabspage{\number\numexpr\value{abspage}+1\relax}%
          \ifcsname abbrnoteseen@#1@\abbrcurrentabspage\endcsname\else
            \global\expandafter\let\csname abbrnoteseen@#1@\abbrcurrentabspage\endcsname\empty
            \abbrnote{#1}{#3}{#4}%
          \fi
        \endgroup
      \fi
    \fi
  \fi
}
\newcommand{\abbrAbs}{\abbrnoteonce{abs}{Abs.}{Absatz}{項}}
\newcommand{\abbrAAO}{\abbrnoteonce{aao}{a.a.O.}{am angegebenen Ort}{前引處}}
\newcommand{\abbrAE}{\abbrnoteonce{ae}{AE}{Alternativ-Entwurf}{替代草案}}
\newcommand{\abbrArt}{\abbrnoteonce{art}{Art.}{Artikel}{條}}
\newcommand{\abbrAufl}{\abbrnoteonce{aufl}{Aufl.}{Auflage}{版}}
\newcommand{\abbrBd}{\abbrnoteonce{bd}{Bd.}{Band}{卷}}
\newcommand{\abbrBGBl}{\abbrnoteonce{bgbl}{BGBl.}{Bundesgesetzblatt}{聯邦法律公報}}
\newcommand{\abbrBGHSt}{\abbrnoteonce{bghst}{BGHSt}{Entscheidungen des Bundesgerichtshofes in Strafsachen}{聯邦普通法院刑事判決彙編}}
\newcommand{\abbrBRDrucks}{\abbrnoteonce{brdrucks}{BRDrucks.}{Bundesrats-Drucksache}{聯邦參議院文件}}
\newcommand{\abbrBTDrucks}{\abbrnoteonce{btdrucks}{BTDrucks.}{Bundestags-Drucksache}{聯邦議院文件}}
\newcommand{\abbrBundesgesetzbl}{\abbrnoteonce{bundesgesetzbl}{Bundesgesetzbl.}{Bundesgesetzblatt}{聯邦法律公報}}
\newcommand{\abbrBvF}{\abbrnoteonce{bvf}{BvF}{Bundesverfassungsgerichtsverfahren}{聯邦憲法法院程序案號}}
\newcommand{\abbrBvL}{\abbrnoteonce{bvl}{BvL}{Bundesverfassungsgerichtsverfahren}{聯邦憲法法院程序案號}}
\newcommand{\abbrBvR}{\abbrnoteonce{bvr}{BvR}{Bundesverfassungsgerichtsverfahren}{聯邦憲法法院程序案號}}
\newcommand{\abbrBVerfG}{\abbrnoteonce{bverfg}{BVerfG}{Bundesverfassungsgericht}{聯邦憲法法院}}
\newcommand{\abbrBVerfGE}{\abbrnoteonce{bverfge}{BVerfGE}{Entscheidungen des Bundesverfassungsgerichts}{聯邦憲法法院判決集}}
\newcommand{\abbrBVerfGG}{\abbrnoteonce{bverfgg}{BVerfGG}{Bundesverfassungsgerichtsgesetz}{聯邦憲法法院法}}
\newcommand{\abbrDDR}{\abbrnoteonce{ddr}{DDR}{Deutsche Demokratische Republik}{德意志民主共和國；東德}}
\newcommand{\abbrDJT}{\abbrnoteonce{djt}{DJT}{Deutscher Juristentag}{德國法學家大會}}
\newcommand{\abbrDr}{\abbrnoteonce{dr}{Dr.}{Doktor}{Dr.}}
\newcommand{\abbrEuGRZ}{\abbrnoteonce{eugrz}{EuGRZ}{Europäische Grundrechte-Zeitschrift}{歐洲基本權期刊}}
\newcommand{\abbrEV}{\abbrnoteonce{ev}{EV}{Einigungsvertrag}{統一條約}}
\newcommand{\abbrF}{\abbrnoteonce{f}{f.}{folgende}{以下一頁；下一段}}
\newcommand{\abbrFF}{\abbrnoteonce{ff}{ff.}{fortfolgende}{以下數頁；以下各段}}
\newcommand{\abbrGBlDDR}{\abbrnoteonce{gblddr}{GBl. DDR}{Gesetzblatt der Deutschen Demokratischen Republik}{德意志民主共和國法律公報}}
\newcommand{\abbrGG}{\abbrnoteonce{gg}{GG}{Grundgesetz}{基本法}}
\newcommand{\abbrGVBl}{\abbrnoteonce{gvbl}{GVBl.}{Gesetz- und Verordnungsblatt}{法律及命令公報}}
\newcommand{\abbrJR}{\abbrnoteonce{jr}{JR}{Juristische Rundschau}{法學評論}}
\newcommand{\abbrJZ}{\abbrnoteonce{jz}{JZ}{JuristenZeitung}{法學家報}}
\newcommand{\abbrIVm}{\abbrnoteonce{ivm}{i.V.m.}{in Verbindung mit}{結合；連同}}
\newcommand{\abbrMdB}{\abbrnoteonce{mdb}{MdB}{Mitglied des Bundestages}{聯邦議院議員}}
\newcommand{\abbrMwN}{\abbrnoteonce{mwn}{m.w.N.}{mit weiteren Nachweisen}{附進一步文獻／判例出處}}
\newcommand{\abbrNF}{\abbrnoteonce{nf}{n.F.}{neue Fassung}{新版本}}
\newcommand{\abbrAF}{\abbrnoteonce{af}{a.F.}{alte Fassung}{舊版本}}
\newcommand{\abbrNr}{\abbrnoteonce{nr}{Nr.}{Nummer}{款、目或號，依語境}}
\newcommand{\abbrProf}{\abbrnoteonce{prof}{Prof.}{Professor}{教授}}
\newcommand{\abbrRdnr}{\abbrnoteonce{rdnr}{Rdnr.}{Randnummer}{邊碼；段碼}}
\newcommand{\abbrRdnrn}{\abbrnoteonce{rdnrn}{Rdnrn.}{Randnummern}{邊碼；段碼}}
\newcommand{\abbrRGBl}{\abbrnoteonce{rgbl}{RGBl.}{Reichsgesetzblatt}{帝國法律公報}}
\newcommand{\abbrRGSt}{\abbrnoteonce{rgst}{RGSt}{Entscheidungen des Reichsgerichts in Strafsachen}{帝國法院刑事判決彙編}}
\newcommand{\abbrRspr}{\abbrnoteonce{rspr}{Rspr.}{Rechtsprechung}{判例；司法實務}}
\newcommand{\abbrRVO}{\abbrnoteonce{rvo}{RVO}{Reichsversicherungsordnung}{帝國保險法}}
\newcommand{\abbrS}{\abbrnoteonce{s}{S.}{Seite}{頁}}
\newcommand{\abbrSFHG}{\abbrnoteonce{sfhg}{SFHG}{Schwangeren- und Familienhilfegesetz}{孕婦及家庭扶助法}}
\newcommand{\abbrSGBV}{\abbrnoteonce{sgbv}{SGB V}{Fünftes Buch Sozialgesetzbuch}{社會法典第五編}}
\newcommand{\abbrStAG}{\abbrnoteonce{staeg}{StÄG}{Strafrechtsänderungsgesetz}{刑法修正法}}
\newcommand{\abbrStenBer}{\abbrnoteonce{stenber}{StenBer.}{Stenographischer Bericht}{速記紀錄}}
\newcommand{\abbrStGB}{\abbrnoteonce{stgb}{StGB}{Strafgesetzbuch}{刑法}}
\newcommand{\abbrStREG}{\abbrnoteonce{streg}{StREG}{Strafrechtsreform-Ergänzungsgesetz}{刑法改革補充法}}
\newcommand{\abbrStrRG}{\abbrnoteonce{strrg}{StrRG}{Strafrechtsreformgesetz}{刑法改革法}}
\newcommand{\abbrUS}{\abbrnoteonce{us}{U.S.}{United States Reports}{美國最高法院判決彙編}}
\newcommand{\abbrVgl}{\abbrnoteonce{vgl}{vgl.}{vergleiche}{參見}}
\newcommand{\abbrVVDStRL}{\abbrnoteonce{vvdstrl}{VVDStRL}{Veröffentlichungen der Vereinigung der Deutschen Staatsrechtslehrer}{德國公法教師協會論文集}}
\newcommand{\abbrWP}{\abbrnoteonce{wp}{WP.}{Wahlperiode}{屆期}}
\newcommand{\abbrWp}{\abbrnoteonce{wp}{Wp.}{Wahlperiode}{屆期}}
\newcommand{\abbrZStrW}{\abbrnoteonce{zstrw}{ZStrW}{Zeitschrift für die gesamte Strafrechtswissenschaft}{整體刑法學期刊}}
\newcommand{\recordsourcerandstart}{%
  \ifrandnumbers
    \setcounter{lastsourcerandstart}{\value{randnumber}}%
    \global\lastsourcehasrandnumberstrue
  \else
    \global\lastsourcehasrandnumbersfalse
  \fi
}
\newlength{\biparagraphskip}
\setlength{\biparagraphskip}{0.52em}
\newlength{\lawboxblockbeforeskip}
\setlength{\lawboxblockbeforeskip}{0.72em}
\newlength{\lawboxblockafterskip}
\setlength{\lawboxblockafterskip}{0.92em}
\newif\iflawboxpartendedblock
\lawboxpartendedblockfalse
\newcommand{\sourceparstyle}{\footnotesize\color{sourceink}\setstretch{1.12}\RaggedRight\emergencystretch=3em\setlength{\parskip}{0.42em}\obeylines\selectfont}
\newcommand{\transparstyle}{\songfont\color{caseink}\setstretch{1.12}\setlength{\parindent}{0pt}\setlength{\parskip}{0.42em}}
\newcommand{\bilingualpairskip}{%
  \par\addvspace{\biparagraphskip}%
  \switchcolumn
  \par\addvspace{\biparagraphskip}%
  \switchcolumn*%
  \global\intranslationfalse
}
\newcommand{\bilingualboxpairskip}{%
  \par
  \switchcolumn
  \par
  \switchcolumn*%
  \global\intranslationfalse
}
\newcommand{\bilinguallawboxblockskip}{%
  \switchcolumn*%
  \par\addvspace{\lawboxblockafterskip}%
  \switchcolumn
  \par\addvspace{\lawboxblockafterskip}%
  \switchcolumn*%
  \global\intranslationfalse
}
\newcommand{\bikeepwithnext}[1][4]{%
  \ifbilingualcolumns
    \syncleft
    \Needspace{#1\baselineskip}%
    \switchcolumn
    \Needspace{#1\baselineskip}%
    \switchcolumn*%
    \global\intranslationfalse
  \else
    \Needspace{#1\baselineskip}%
  \fi
}
\NewEnviron{sourcepara}[1][]{
  \recordsourcerandstart
  \ifshoworiginal
    \ifbilingualcolumns
      \ifintranslation\switchcolumn*\global\intranslationfalse\fi
    \else
      \Needspace{5\baselineskip}
    \fi
    \ifbilingualcolumns\else\par\addvspace{0.35em}\fi
    {\sourceparstyle
    \noindent\BODY\par}
    \ifbilingualcolumns\else\addvspace{0.25em}\fi
    \ifbilingualcolumns
      \switchcolumn
      \global\intranslationtrue
    \fi
  \else
    \ifrandnumbers
      \begingroup
        \setbox0=\vbox{\hsize=\linewidth\sourceparstyle\noindent\BODY\par}%
      \endgroup
    \fi
  \fi
}
\NewEnviron{transpara}{
  \ifshowtranslation
    \setcounter{randnumberlocal}{\value{lastsourcerandstart}}%
    \ifbilingualcolumns
      \ifintranslation\else\switchcolumn\global\intranslationtrue\fi
      {\transparstyle\setlength{\leftskip}{0.55em}\setlength{\rightskip}{0.15em plus 1em}\addtolength{\linewidth}{-0.55em}\BODY\par}
      \switchcolumn*
      \bilingualpairskip
    \else
      {\transparstyle\BODY\par}
    \fi
  \else
    \ifbilingualcolumns
      \ifintranslation\switchcolumn*\global\intranslationfalse\fi
    \fi
  \fi
}
\NewEnviron{sourceboxpara}[1][]{
  \global\lawboxpartendedblockfalse
  \recordsourcerandstart
  \ifshoworiginal
    \ifbilingualcolumns
      \ifintranslation\switchcolumn*\global\intranslationfalse\fi
    \else
      \Needspace{3\baselineskip}
    \fi
    {\sourceparstyle
    \BODY\par}
    \ifbilingualcolumns\else
      \iflawboxpartendedblock\addvspace{\lawboxblockafterskip}\fi
    \fi
    \ifbilingualcolumns
      \switchcolumn
      \global\intranslationtrue
    \fi
  \fi
}
\NewEnviron{transboxpara}{
  \global\lawboxpartendedblockfalse
  \ifshowtranslation
    \setcounter{randnumberlocal}{\value{lastsourcerandstart}}%
    \ifbilingualcolumns
      \ifintranslation\else\switchcolumn\global\intranslationtrue\fi
      {\transparstyle\BODY\par}
      \iflawboxpartendedblock
        \bilinguallawboxblockskip
      \else
        \switchcolumn*
        \bilingualboxpairskip
      \fi
    \else
      {\transparstyle\BODY\par}
      \iflawboxpartendedblock\addvspace{\lawboxblockafterskip}\fi
    \fi
  \else
    \ifbilingualcolumns
      \ifintranslation\switchcolumn*\global\intranslationfalse\fi
    \fi
  \fi
}
\newcommand{\leitsatznum}[1]{\noindent\hangindent=3.0em\hangafter=1\makebox[2.25em][r]{\bfseries #1.}\hspace{0.55em}\ignorespaces}
\newcommand{\syncleft}{\ifbilingualcolumns\ifintranslation\switchcolumn*\global\intranslationfalse\fi\fi}
\pretocmd{\section}{\syncleft\clearpage}{}{}
\pretocmd{\subsection}{\syncleft}{}{}
\pretocmd{\subsubsection}{\syncleft}{}{}
\newcommand{\biheadingfont}{\songfont}
\newcommand{\bitocentry}[2]{%
  \ifshoworiginal
    \ifshowtranslation
      \ifstrequal{#1}{#2}{#1}{#1\space #2}%
    \else
      #1%
    \fi
  \else
    #2%
  \fi
}
\pdfstringdefDisableCommands{%
  \def\bitocentry#1#2{#1}%
}
\newcommand{\biheading}[7]{%
  \syncleft#1\bikeepwithnext[#3]\phantomsection\addcontentsline{toc}{#2}{\bitocentry{#6}{#7}}%
  \ifbilingualcolumns
    {#4 #6\par}%
    \switchcolumn
    {#4\biheadingfont #7\par}%
    \switchcolumn*%
    \global\intranslationfalse
    \par\addvspace{#5}%
    \switchcolumn
    \par\addvspace{#5}%
    \switchcolumn*%
  \else
    \ifshoworiginal{#4 #6\par}\fi
    \ifshowtranslation{#4\biheadingfont #7\par}\fi
    \addvspace{#5}%
  \fi
}
\newcommand{\termsectionheading}[1]{%
  \par\Needspace{9\baselineskip}%
  \vspace{0.65em}%
  {\songfont\bfseries #1\par}%
  \vspace{0.2em}%
}
% 章節換頁：
%   雙語 paracol 欄內不要用 \clearpage；它會在同步兩欄時吐出只有欄線/頁碼的空頁。
%   \newpage 足以讓下一個 section 起新頁，又不會強制清空所有待排材料。
%   單語模式不在 paracol 內，仍可用 \clearpage。
\newcommand{\bisectionpagebreak}{\ifbilingualcolumns\newpage\else\clearpage\fi}
\newcommand{\bisection}[2]{%
  \biheading{\iffirstbilingualsection\global\firstbilingualsectionfalse\else\bisectionpagebreak\fi}{section}{6}{\centering\Large\bfseries}{0.8em}{#1}{#2}%
}
\newcommand{\bisubsection}[2]{%
  \biheading{}{subsection}{5}{\large\bfseries}{0.55em}{#1}{#2}%
}
\newcommand{\bisubsubsection}[2]{%
  \biheading{}{subsubsection}{4}{\normalsize\bfseries}{0.45em}{#1}{#2}%
}
\newcommand{\bicompositeheading}[4]{%
  \biheading{}{subsection}{5}{\large\bfseries}{0.16em}{#1}{#2}%
  \biheading{}{subsubsection}{3}{\normalsize\bfseries}{0.45em}{#3}{#4}%
}
\newif\iflawboxfirstitem
\lawboxfirstitemfalse
\newcommand{\lawboxprespace}[1]{%
  \par
  \iflawboxfirstitem
    \global\lawboxfirstitemfalse
  \else
    \addvspace{#1}%
  \fi
}
\newtcolorbox{lawboxinner}{
  caseboxbase,
  colframe=sourceline!85!black,
  colback=white,
  left=1.05em,
  right=1.05em,
  top=0.62em,
  bottom=0.62em,
  before upper={\global\inlawboxtrue\global\lawboxfirstitemtrue\global\lawboxhaspendingrnfalse\ifintranslation\lawboxtransfont\else\lawboxsourcefont\fi\RaggedRight\setlength{\parindent}{0pt}\setlength{\parskip}{0pt}\ifintranslation\setstretch{\lawboxtransstretch}\else\setstretch{\lawboxsourcestretch}\fi},
  after upper={\global\lawboxhaspendingrnfalse\global\inlawboxfalse}
}
\tcbset{
  lawboxpartbase/.style={
    enhanced,
    breakable=false,
    width=0.985\linewidth,
    colback=white,
    frame hidden,
    boxrule=0pt,
    arc=0pt,
    left=1.05em,
    right=1.05em,
    boxsep=1.5mm,
    before skip=0pt,
    after skip=0pt,
    before upper={\global\inlawboxtrue\global\lawboxfirstitemtrue\global\lawboxhaspendingrnfalse\ifintranslation\lawboxtransfont\else\lawboxsourcefont\fi\RaggedRight\setlength{\parindent}{0pt}\setlength{\parskip}{0pt}\ifintranslation\setstretch{\lawboxtransstretch}\else\setstretch{\lawboxsourcestretch}\fi},
    after upper={\global\lawboxhaspendingrnfalse\global\inlawboxfalse},
    borderline west={0.28pt}{0pt}{sourceline!85!black},
    borderline east={0.28pt}{0pt}{sourceline!85!black},
  },
  lawboxpartfirst/.style={
    lawboxpartbase,
    top=0.62em,
    bottom=0.04em,
    before skip=\lawboxblockbeforeskip,
    borderline north={0.28pt}{0pt}{sourceline!85!black},
  },
  lawboxpartmiddle/.style={
    lawboxpartbase,
    top=0.04em,
    bottom=0.04em,
  },
  lawboxpartlast/.style={
    lawboxpartbase,
    top=0.04em,
    bottom=0.62em,
    after skip=0pt,
    borderline south={0.28pt}{0pt}{sourceline!85!black},
  }
}
\newtcolorbox{lawboxpartinner}[1]{#1}
\newtcolorbox{termboxinner}{
  caseboxbase,
  colframe=noteblue!70!black,
  colback=casepaper,
  before upper={\songfont\setlength{\parindent}{0pt}\setlength{\parskip}{0.35em}}
}
\newtcolorbox{deboxinner}{
  caseboxbase,
  colframe=notegray!72!black,
  colback=white,
  left=1.05em,
  right=1.05em,
  top=0.62em,
  bottom=0.62em,
  before upper={\global\inlawboxtrue\global\lawboxfirstitemtrue\global\lawboxhaspendingrnfalse\small\setlength{\parindent}{0pt}\setlength{\parskip}{0.35em}},
  after upper={\global\lawboxhaspendingrnfalse\global\inlawboxfalse}
}
\NewEnviron{lawbox}{
  \begin{lawboxinner}\BODY\end{lawboxinner}
}
\NewEnviron{lawboxpart}[2]{
  \begin{lawboxpartinner}{lawboxpart#1,equal height group=#2}\BODY\end{lawboxpartinner}%
  \ifstrequal{#1}{last}{\global\lawboxpartendedblocktrue}{}%
}
\NewEnviron{termbox}{
  \par\addvspace{0.55em}
  {\color{noteblue!70!black}\rule{\linewidth}{0.28pt}}\par
  \begingroup
    \songfont\small\color{caseink}
    \setlength{\parindent}{0pt}
    \setlength{\parskip}{0.24em}
    \BODY
    \par
  \endgroup
  {\color{noteblue!70!black}\rule{\linewidth}{0.28pt}}\par
  \addvspace{0.55em}
}
\NewEnviron{debox}{
  \begin{deboxinner}\BODY\end{deboxinner}
}
\providecommand{\paragraphmark}[1]{}
\newcommand{\lawboxtitleafter}{\addvspace{0.06em}}
\newcommand{\lawtitle}[1]{\lawboxprespace{0.22em}{\lawboxtransfont\bfseries\lawboxuserandnumber #1}\par\lawboxtitleafter}
\newcommand{\absno}[2]{\lawboxprespace{0.04em}\noindent\lawboxuserandnumber\hangindent=2.6em\hangafter=1\makebox[2.6em][l]{(#1)}#2\par}
\newcommand{\nrno}[2]{\lawboxprespace{0pt}\noindent\lawboxuserandnumber\hspace*{2.6em}\hangindent=4.4em\hangafter=1\makebox[1.8em][l]{#1.}#2\par}
\newcommand{\letterno}[2]{\lawboxprespace{0pt}\noindent\lawboxuserandnumber\hspace*{4.4em}\hangindent=6.0em\hangafter=1\makebox[1.6em][l]{#1)}#2\par}
\newcommand{\sourcelawtitle}[1]{\lawboxprespace{0.22em}{\lawboxsourcefont\bfseries\lawboxuserandnumber #1}\par\lawboxtitleafter}
\newcommand{\sourcelawsubtitle}[1]{\lawboxprespace{0.04em}{\lawboxsourcefont\bfseries\lawboxuserandnumber #1}\par\lawboxtitleafter}
\newcommand{\sourceabsno}[2]{\lawboxprespace{0.04em}\noindent\lawboxuserandnumber\hangindent=2.45em\hangafter=1\makebox[2.45em][l]{(#1)}#2\par}
\newcommand{\sourcenrno}[2]{\lawboxprespace{0pt}\noindent\lawboxuserandnumber\hspace*{2.45em}\hangindent=4.25em\hangafter=1\makebox[1.8em][l]{#1.}#2\par}
\newcommand{\sourceletterno}[2]{\lawboxprespace{0pt}\noindent\lawboxuserandnumber\hspace*{4.25em}\hangindent=5.85em\hangafter=1\makebox[1.6em][l]{#1)}#2\par}
\newcommand{\sourcequotepara}[1]{\lawboxprespace{0.04em}\noindent\lawboxuserandnumber\hangindent=1.2em\hangafter=1 #1\par}
\newcommand{\sourcelawline}[1]{\lawboxprespace{0.04em}\noindent\lawboxuserandnumber #1\par}
\newcommand{\lawline}[1]{\lawboxprespace{0.04em}\noindent\lawboxuserandnumber #1\par}
\newcommand{\pendingtranslation}{\par{\songfont\footnotesize\color{pendingink}（中文譯文待補。）}\par}
\newcommand{\tocpart}[1]{\par\addvspace{0.52em}{\footnotesize\bfseries\color{pendingink}#1\par}\addvspace{0.16em}}
\newcommand{\tocdashline}{\par\addvspace{0.48em}\noindent{\color{notegray}\xleaders\hbox to 0.82em{\hfil\rule{0.42em}{0.28pt}\hfil}\hskip\linewidth}\par\addvspace{0.38em}}
% \tocpanel{font-cmd}{content}：將目錄置中於所在欄寬（雙語欄適用）。
% A3 橫式使用 0.58\linewidth，A4 橫式使用 0.84\linewidth，
% 使兩種紙張下目錄欄內容實際寬度相近（均約 100–105 mm）。
\newcommand{\tocpanel}[2]{%
  \makebox[\linewidth][c]{%
    \ifx\bilingualpaper\bilingualpaperaThree
      \begin{minipage}{0.58\linewidth}%
    \else
      \begin{minipage}{0.84\linewidth}%
    \fi
    \toccolstyle
    #1#2%
    \end{minipage}%
  }%
}
% ===== 首頁簡目錄的兩層 description list 縮排 =====
% enumitem 的 description list 可把每一列拆成：
%   [label]  labelsep  body text
% 這裡的「起點」都以所在 minipage/欄位的左緣為 0。
%
% 核心規則：
%   labelindent = label 欄位從左緣開始的位置。
%   labelwidth  = label 欄位本身寬度；標籤太長（如 Abw. Meinung）就加大它。
%   labelsep    = label 欄位右緣到正文的距離。
%   leftmargin  = 正文 body text 的起點。判斷層級視覺時主要看這個。
%
% 所以：
%   想讓「Begründung / 判決理由」整列正文往右或往左：改 tocitems 的 leftmargin。
%   想讓「A. Verfahren und Gegenstand / A. 程序與審查標的」正文往右或往左：
%     改 tocsubitems 的 leftmargin。
%   想讓 A./B./C. 這些標籤本身往右或往左，但正文不動：改 tocsubitems 的 labelindent。
%   想讓 A./B./C. 與正文間距變大或變小：改 tocsubitems 的 labelsep。
%
% 層級原則：
%   tocsubitems.leftmargin 必須大於 tocitems.leftmargin，
%   否則子層級正文會看起來比「Gründe / 理由」還靠左。
\newlist{tocitems}{description}{1}
\setlist[tocitems]{%
  labelindent=0pt,    % 第一層 label 欄位起點；0pt 表示貼齊目錄欄左緣。
  leftmargin=2.54cm,  % 第一層正文起點；例如 Begründung / 判決理由 的左緣。
  labelwidth=2.16cm,  % 第一層 label 欄位寬度；需容納 Leitsätze、Abw. Meinung、不同意見等。
  labelsep=0.32cm,    % 第一層 label 與正文的水平間距；只改間距，不改正文起點。
  itemsep=0.28em,     % 第一層各項之間的垂直距離。
  topsep=0.20em,      % 第一層 list 上下與前後內容的垂直距離。
  parsep=0pt,         % 同一 item 內段落之間的距離；目錄項通常不需要。
  partopsep=0pt,      % list 若緊接段落開始時的額外距離；關閉以免目錄高度飄動。
  align=left,         % label 欄內左對齊；長短標籤不靠右貼齊。
  font=\normalfont\bfseries % label 的字型；正文不受這行影響。
}
\newlist{tocsubitems}{description}{1}
\setlist[tocsubitems]{%
  labelindent=1cm,    % 第二層 label 起點；只控制 A./B./C. 本身的位置。
  leftmargin=3.00cm,  % 第二層正文起點；必須大於 2.54cm，才會比 Begründung / 判決理由更右。
  labelwidth=0.5cm,   % 第二層 label 欄寬；A./B./C. 很短，可小於第一層。
  labelsep=1.5cm,    % 第二層 label 與正文間距；A. 與 Verfahren / 程序 的距離看這裡。
  itemsep=0.12em,     % 第二層各項較密，避免 A.–E. 佔太高。
  topsep=0.04em,      % 第二層 list 與 Gründe / 理由之間的垂直距離。
  parsep=0pt,         % 同一第二層 item 內段落距離；目錄項通常不需要。
  partopsep=0pt,      % 關閉額外段落起始距離，避免版面不穩。
  align=left,         % A./B./C. label 欄內左對齊。
  font=\normalfont\bfseries, % A./B./C. label 加粗；正文不受這行影響。
  before={}           % 不在第二層 list 前自動插入額外內容。
}
\newlist{termitems}{description}{1}
\ifbilingualcolumns
  \setlist[termitems]{%
    leftmargin=5.2cm,
    labelwidth=4.8cm,
    labelsep=0.35cm,
    itemsep=0.16em,
    topsep=0.2em,
    parsep=0pt,
    partopsep=0pt,
    font=\normalfont\bfseries,
    style=nextline
  }
\else
  \setlist[termitems]{%
    leftmargin=0pt,
    labelwidth=0pt,
    labelsep=0pt,
    itemsep=0.24em,
    topsep=0.2em,
    parsep=0pt,
    partopsep=0pt,
    font=\normalfont\bfseries,
    style=nextline
  }
\fi
\newlist{abbritems}{description}{1}
\setlist[abbritems]{%
  leftmargin=2.38cm,
  labelwidth=2.02cm,
  labelsep=0.28cm,
  itemsep=0.16em,
  topsep=0.2em,
  parsep=0pt,
  partopsep=0pt,
  align=left,
  font=\normalfont\bfseries
}
\ifbilingualcolumns
  \newenvironment{termcolumns}{\begin{multicols}{2}}{\end{multicols}}
\else
  \newenvironment{termcolumns}{}{}
\fi

% ===== 編者註腳開關 =====
% 一般版預設不顯示編者註，避免正文腳註過密。
% 若開啟 \classroommodeon，GG/條文編者註仍會自動顯示，無需另開本開關。
% 若一般版也要顯示編者註，可在單案檔 \input 後加 \showeditornotestrue。
\newif\ifshoweditornotes
\showeditornotesfalse

% ===== 目錄排版輔助 =====
\newcommand{\toccolstyle}{\raggedright\arraybackslash}
\newcommand{\tocsingleindent}{\leftskip=1.45cm\rightskip=1.2cm}

% ===== 行內引文環境（德文原文與中文譯文各一，帶左側灰色邊線）=====
\newcommand{\quoteparbreak}{\par\addvspace{0.48em}}
\newcommand{\sourcequoteinner}[1]{%
  \begin{tcolorbox}[
    enhanced, breakable, width=\dimexpr\linewidth-0.8em\relax, frame hidden,
    colback=white, coltext=caseink, boxrule=0pt,
    borderline west={0.55pt}{0.88em}{notegray},
    left=1.78em, right=0.72em, top=0.18em, bottom=0.18em,
    boxsep=0pt, before skip=0.24em, after skip=0.24em,
    before upper={\normalfont\footnotesize\color{caseink}\setlength{\parindent}{0pt}\setlength{\parskip}{0.34em}\raggedright\setlength{\rightskip}{0pt plus 1em}}
  ]%
  #1\par
  \end{tcolorbox}%
}
\newcommand{\transquoteinner}[1]{%
  \begin{tcolorbox}[
    enhanced, breakable, width=\dimexpr\linewidth-0.8em\relax, frame hidden,
    colback=white, coltext=caseink, boxrule=0pt,
    borderline west={0.55pt}{0.88em}{notegray},
    left=1.78em, right=0.72em, top=0.18em, bottom=0.18em,
    boxsep=0pt, before skip=0.24em, after skip=0.24em,
    before upper={\songfont\small\color{caseink}\setlength{\parindent}{0pt}\setlength{\parskip}{0.34em}\raggedright\setlength{\rightskip}{0pt plus 1em}}
  ]%
  #1\par
  \end{tcolorbox}%
}

% ===== 大綱（Gliederung）Rn 輔助命令 =====
\newcommand{\outrn}[2]{\enspace{\normalfont\footnotesize\color{pendingink}(Rn\,#1–#2)}}
\newcommand{\outrnsingle}[1]{\enspace{\normalfont\footnotesize\color{pendingink}(Rn\,#1)}}
\newcommand{\outrnzh}[2]{\enspace{\normalfont\footnotesize\color{pendingink}（第\,#1–#2\,段）}}
\newcommand{\outrnzhs}[1]{\enspace{\normalfont\footnotesize\color{pendingink}（第\,#1\,段）}}
% ===== 大綱雙語對照表格輔助命令（三欄：段碼 │ 德文 │ 中文）=====
% 欄 1：段碼（由欄定義自動格式化：小字、等寬、右對齊、pendingink）
% 欄 2：德文標籤＋正文
% 欄 3：中文標籤＋正文
%
% 無段碼的列（\omainrow、\osecrow），欄 1 以 & 空置。
% 有段碼的列（\osecrowrn、\osubrow、\ointrow），#1 為預格式化字串
%   例：\osubrow{2–49}{I.}{...}{I.}{...}
%       \ointrow{107}{...}{...}
%
% \opartrow：段組標題（橫跨三欄）
\newcommand{\opartrow}[2]{%
  \noalign{\vspace{0.30em}}%
  \multicolumn{3}{@{}l@{}}{%
    \footnotesize\bfseries\color{pendingink}#1\hfill\songfont#2%
  }\\[0.06em]%
}
% \odashrow：虛線分隔（橫跨三欄）
\newcommand{\odashrow}{%
  \noalign{\vspace{0.28em}}%
  \multicolumn{3}{@{}l@{}}{\color{notegray}%
    \xleaders\hbox to 0.82em{\hfil\rule{0.42em}{0.28pt}\hfil}\hskip 0.88\linewidth}%
  \\[0.24em]%
}
% \odissentpart：不同意見書區段標頭。比一般 \opartrow 更像附錄性轉場。
\newcommand{\odissentpart}[2]{%
  \noalign{\vspace{0.42em}}%
  \multicolumn{3}{@{}p{0.88\textwidth}@{}}{%
    \color{notegray}\rule{\linewidth}{0.18pt}%
  }\\[-0.18em]%
  \multicolumn{3}{@{}p{0.88\textwidth}@{}}{%
    \makebox[\linewidth][s]{%
      \footnotesize\bfseries\color{pendingink}#1%
      \hfill{\songfont #2}%
    }%
  }\\[-0.02em]%
}
% \odissentopinion：不同意見書的署名與一行摘要，右欄保留段碼範圍。
\newcommand{\odissentopinion}[5]{%
  \footnotesize\hangindent=0.52cm\hangafter=1%
    {\bfseries\color{caseink}#2}\enspace{\color{notegray}\textperiodcentered}\enspace\color{caseink}#3%
  &\footnotesize\hangindent=0.52cm\hangafter=1%
    {\songfont\bfseries\color{caseink}#4}\enspace{\color{notegray}\textperiodcentered}\enspace\songfont\color{caseink}#5%
  &#1%
  \\[0.12em]%
}
% \omainrow：頂層項目，無段碼（Leitsätze、Urteil、Gründe …）
\newcommand{\omainrow}[4]{%
  \footnotesize\hangindent=1.92cm\hangafter=1%
    \makebox[1.92cm][l]{\bfseries\color{caseink}#1}\color{caseink}#2%
  &\footnotesize\hangindent=1.92cm\hangafter=1%
    \makebox[1.92cm][l]{\bfseries\color{caseink}#3}\songfont\color{caseink}#4%
  &%
  \\[0.15em]%
}
% \omainrowrn：頂層項目，有段碼（不同意見等標籤寬於 \osecrow 框者）
\newcommand{\omainrowrn}[5]{%
  \footnotesize\hangindent=1.92cm\hangafter=1%
    \makebox[1.92cm][l]{\bfseries\color{caseink}#2}\color{caseink}#3%
  &\footnotesize\hangindent=1.92cm\hangafter=1%
    \makebox[1.92cm][l]{\bfseries\color{caseink}#4}\songfont\color{caseink}#5%
  &#1%
  \\[0.15em]%
}
% \osecrow：一級子項，無段碼（A.、C.、D. 等有子項者）
\newcommand{\osecrow}[4]{%
  \footnotesize\hspace{1.10cm}\hangindent=1.98cm\hangafter=1%
    \makebox[0.86cm][l]{\bfseries\color{caseink}#1}\color{caseink}#2%
  &\footnotesize\hspace{1.10cm}\hangindent=1.98cm\hangafter=1%
    \makebox[0.86cm][l]{\bfseries\color{caseink}#3}\songfont\color{caseink}#4%
  &%
  \\[0.11em]%
}
% \osecrowrn：一級子項，有段碼（B.、E. 等完整段落範圍者）
% #1=段碼字串（如 101–106）, #2=DE標籤, #3=DE正文, #4=ZH標籤, #5=ZH正文
\newcommand{\osecrowrn}[5]{%
  \footnotesize\hspace{1.10cm}\hangindent=1.98cm\hangafter=1%
    \makebox[0.86cm][l]{\bfseries\color{caseink}#2}\color{caseink}#3%
  &\footnotesize\hspace{1.10cm}\hangindent=1.98cm\hangafter=1%
    \makebox[0.86cm][l]{\bfseries\color{caseink}#4}\songfont\color{caseink}#5%
  &#1%
  \\[0.11em]%
}
% \osubrow：二級子項，有段碼範圍（I.、II.、A.I. …）
% #1=段碼字串, #2=DE標籤, #3=DE正文, #4=ZH標籤, #5=ZH正文
\newcommand{\osubrow}[5]{%
  \footnotesize\hspace{1.40cm}\hangindent=2.30cm\hangafter=1%
    \makebox[0.90cm][l]{\bfseries\color{caseink}#2}\color{caseink}#3%
  &\footnotesize\hspace{1.40cm}\hangindent=2.30cm\hangafter=1%
    \makebox[0.90cm][l]{\bfseries\color{caseink}#4}\songfont\color{caseink}#5%
  &#1%
  \\[0.09em]%
}
% \ointrow：引入段落，有單一段碼，無標籤
% #1=段碼字串, #2=DE正文, #3=ZH正文
\newcommand{\ointrow}[3]{%
  \footnotesize\hspace{0.52cm}\hangindent=0.52cm\hangafter=1\color{caseink}#2%
  &\footnotesize\hspace{0.52cm}\hangindent=0.52cm\hangafter=1\songfont\color{caseink}#3%
  &#1%
  \\[0.09em]%
}
% \osubsubrow：三級子項（1.、2.、3.）
% #1=段碼字串, #2=DE標籤, #3=DE正文, #4=ZH標籤, #5=ZH正文
\newcommand{\osubsubrow}[5]{%
  \footnotesize\hspace{1.80cm}\hangindent=2.48cm\hangafter=1%
    \makebox[0.68cm][l]{\bfseries\color{caseink}#2}\color{caseink}#3%
  &\footnotesize\hspace{1.80cm}\hangindent=2.48cm\hangafter=1%
    \makebox[0.68cm][l]{\bfseries\color{caseink}#4}\songfont\color{caseink}#5%
  &#1%
  \\[0.08em]%
}
% \osubsubsubrow：四級子項（a）、b））
% #1=段碼字串, #2=DE標籤, #3=DE正文, #4=ZH標籤, #5=ZH正文
\newcommand{\osubsubsubrow}[5]{%
  \footnotesize\hspace{2.10cm}\hangindent=2.68cm\hangafter=1%
    \makebox[0.58cm][l]{\bfseries\color{caseink}#2}\color{caseink}#3%
  &\footnotesize\hspace{2.10cm}\hangindent=2.68cm\hangafter=1%
    \makebox[0.58cm][l]{\bfseries\color{caseink}#4}\songfont\color{caseink}#5%
  &#1%
  \\[0.07em]%
}
% \osubsubsubsubrow：五級子項（aa）、bb））
% 標籤框 0.62cm：足以容納「aa)」在 footnotesize bold 下（約 0.50cm）並留有間距
% #1=段碼字串, #2=DE標籤, #3=DE正文, #4=ZH標籤, #5=ZH正文
\newcommand{\osubsubsubsubrow}[5]{%
  \footnotesize\hspace{2.38cm}\hangindent=3.06cm\hangafter=1%
    \makebox[0.68cm][l]{\bfseries\color{caseink}#2}\color{caseink}#3%
  &\footnotesize\hspace{2.38cm}\hangindent=3.06cm\hangafter=1%
    \makebox[0.68cm][l]{\bfseries\color{caseink}#4}\songfont\color{caseink}#5%
  &#1%
  \\[0.06em]%
}
% \outlinepage：雙語對照大綱頁包裝環境（longtable 自動跨頁）
% 三欄：德文欄 + 中文欄 + 段碼欄（最右，不折行）
% 段碼欄寬度：1.80cm，足以容納最長段碼如「309–348」
% DE/ZH 欄寬：(0.87\textwidth - 2.08cm) / 2
%   驗算：2×欄 + 0.05\textwidth + 0.28cm gap + 1.80cm = 0.87tw - 2.08cm + 0.05tw + 2.08cm = 0.92tw ✓
\newenvironment{outlinepage}{%
  \vspace*{0.40cm}%
  \setlength{\LTleft}{0.04\textwidth}%
  \setlength{\LTright}{0.04\textwidth}%
  \setlength{\tabcolsep}{0pt}%
  % 大綱列距由 longtable 的 \arraystretch 與各 row macro 結尾的 \\[...em] 共同決定。
  % 這裡控制「每一列本身的表格 strut 高度」；數值越大，A./I./1. 各項之間越像有空行。
  % A3 原本用 0.62，視覺上沒有空行；A4 也沿用同值，避免切到 A4 後項目被撐開。
  % 若日後覺得大綱太擠，只微調這個值（例如 0.68），不要改各 \omainrow/\osubrow 的縮排。
  \renewcommand{\arraystretch}{0.62}%
  \begin{longtable}{%
    >{\vspace{0pt}\raggedright\arraybackslash}p{\dimexpr(0.87\textwidth-2.08cm)/2\relax}%
    @{\hspace{0.025\textwidth}}%
    !{\color{notegray}\vrule width 0.12pt}%
    @{\hspace{0.025\textwidth}}%
    >{\vspace{0pt}\raggedright\arraybackslash}p{\dimexpr(0.87\textwidth-2.08cm)/2\relax}%
    @{\hspace{0.28cm}}%
    >{\vspace{0pt}\footnotesize\normalfont\color{pendingink}\raggedright\arraybackslash}p{1.80cm}%
    @{}%
  }%
  % 首頁欄標籤：不要橫跨置中；分別放在德文欄與中文欄的實際欄位起點。
  \footnotesize\bfseries\color{sourceink}Gliederung%
  &\footnotesize\bfseries\songfont\color{sourceink}大綱%
  &%
  \\[0.36em]%
  \endfirsthead%
  % 續頁欄標籤：同樣分欄定位，以灰色細體區分；無需標注「續」。
  \footnotesize\color{pendingink}Gliederung%
  &\footnotesize\songfont\color{pendingink}大綱%
  &%
  \\[0.24em]%
  \endhead%
}{%
  \end{longtable}%
}
 之前設定；
%   預設 chinese / a4）
\ifdefined\docmode\else\def\docmode{chinese}\fi
\ifdefined\bilingualpaper\else\def\bilingualpaper{a4}\fi
\newif\ifshoworiginal
\newif\ifshowtranslation
\newif\ifbilinguallandscape
\newif\ifbilingualcolumns
\newif\ifintranslation
\newif\iffirstbilingualsection

\showoriginalfalse
\showtranslationfalse
\bilinguallandscapefalse
\bilingualcolumnsfalse
\intranslationfalse
\firstbilingualsectionfalse

\def\modechinese{chinese}
\def\modegerman{german}
\def\modebilingual{bilingual}
\def\bilingualpaperaThree{a3}
\def\bilingualpaperaFour{a4}

\ifx\docmode\modechinese
  \showtranslationtrue
\fi

\ifx\docmode\modegerman
  \showoriginaltrue
\fi

\ifx\docmode\modebilingual
  \showoriginaltrue
  \showtranslationtrue
  \bilinguallandscapetrue
  \bilingualcolumnstrue
\fi

% 編排模式說明：
% chinese  → \showoriginalfalse  \showtranslationtrue  （A4 直式，純中文）
% german   → \showoriginaltrue   \showtranslationfalse （A4 直式，純德文）
% bilingual→ \showoriginaltrue   \showtranslationtrue  \bilingualcolumnstrue（A3/A4 橫式對照）

\ifbilingualcolumns
  \ifx\bilingualpaper\bilingualpaperaFour
    \geometry{
      a4paper,
      landscape,
      left=1.25cm,
      right=2.25cm,
      top=1.25cm,
      bottom=1.75cm,
      footskip=8.5mm,
      marginparwidth=0.75cm,
      marginparsep=0.25cm
    }
  \else
    \geometry{
      a3paper,
      landscape,
      left=1.55cm,
      right=2.75cm,
      top=1.55cm,
      bottom=2.05cm,
      footskip=9.5mm,
      marginparwidth=0.9cm,
      marginparsep=0.35cm
    }
  \fi
\else
\geometry{
  a4paper,
  portrait,
  left=2.2cm,
  right=2.6cm,
  top=2.0cm,
  bottom=2.0cm,
  marginparwidth=0.8cm,
  marginparsep=0.3cm
}
\fi
% ===== 時間戳基礎設施 =====
\newcount\stamptimeval
\newcount\stampminuteval
\newcommand{\stamphh}{%
  \stamptimeval=\time
  \divide\stamptimeval by 60
  \ifnum\stamptimeval<10 0\fi
  \the\stamptimeval}
\newcommand{\stampmm}{%
  \stamptimeval=\time
  \stampminuteval=\time
  \divide\stampminuteval by 60
  \multiply\stampminuteval by 60
  \advance\stamptimeval by -\stampminuteval
  \ifnum\stamptimeval<10 0\fi
  \the\stamptimeval}
\newcommand{\stampwkde}[1]{%
  \ifcase#1Montag\or Dienstag\or Mittwoch\or Donnerstag\or Freitag\or Samstag\or Sonntag\fi}
\newcommand{\stampwkzh}[1]{%
  \ifcase#1 一\or 二\or 三\or 四\or 五\or 六\or 日\fi}
\newcommand{\stampmonthde}[1]{%
  \ifcase#1\or Januar\or Februar\or März\or April\or Mai\or Juni\or
  Juli\or August\or September\or Oktober\or November\or Dezember\fi}
\newcount\stampjdval
\newcount\stampwdval
\pgfcalendardatetojulian{\the\year-\the\month-\the\day}{\stampjdval}
\pgfcalendarjuliantoweekday{\the\stampjdval}{\stampwdval}
\newcommand{\timestampzh}{%
  \songfont 最後修改：\the\year 年\the\month 月\the\day 日%
  % （星期\stampwkzh{\the\stampwdval}）%
  % \space\stamphh:\stampmm
  }

\newcommand{\documentdatezh}{\the\year 年 \the\month 月 \the\day 日}
\newcommand{\documentdatede}{\the\day. \stampmonthde{\the\month} \the\year}
\newcommand{\bverfgetranslationnote}{%
  {\songfont 翻譯說明：初譯使用 Claude Code 與 GPT 協助迭代；仍請以德文原文為準。}%
}
\newcommand{\bverfgecourseinfo}{%
  {\songfont 課程：114 學年度第 2 學期；憲法專題研究（四）；授課老師：湯德宗教授；導讀日期：115 年 6 月 1 日。}%
}
\newcommand{\bverfgeeditorinfo}{%
  {\songfont 編者：王逸帆（R10A21126），國立台灣大學法律系研究所碩士班。}%
}
\newcommand{\bverfgetitleversionline}{%
  \ifbilingualcolumns
    Fassung \documentversion\enspace\textperiodcentered\enspace Stand: \documentdatede
    \quad
    {\songfont 版本 \documentversion\enspace\textperiodcentered\enspace 修訂日期：\documentdatezh\par}%
  \else
    \ifshowtranslation
      {\songfont 版本 \documentversion\enspace\textperiodcentered\enspace 修訂日期：\documentdatezh\par}%
    \else
      Fassung \documentversion\enspace\textperiodcentered\enspace Stand: \documentdatede\par
    \fi
  \fi
}
\newcommand{\bverfgetitlecornerinfo}{%
  \ifshowtranslation
    \vspace{0.72em}%
    \noindent\hfill
    \begin{minipage}{0.66\textwidth}
      \raggedleft\tiny\color{pendingink}\songfont
      \bverfgecourseinfo\par
      \bverfgeeditorinfo
    \end{minipage}%
  \fi
}
\newcommand{\bverfgetitlemeta}{%
  \vfill
  \begin{center}%
    \footnotesize\color{pendingink}%
    \bverfgetitleversionline
    \ifshowtranslation
      \vspace{0.28em}{\scriptsize\bverfgetranslationnote\par}%
    \fi
  \end{center}%
  \bverfgetitlecornerinfo
}

\pagestyle{fancy}
\fancyhf{}
\renewcommand{\headrulewidth}{0pt}
\renewcommand{\footrulewidth}{0pt}
\fancyfoot[C]{\thepage}
\ifbilingualcolumns
  \fancyfoot[R]{\scriptsize\color{pendingink}\timestampzh}%
\else
  \ifshoworiginal
  \else
    \fancyfoot[R]{\scriptsize\color{pendingink}\timestampzh\hspace{0.45cm}}%
  \fi
\fi
\setmainfont{Times New Roman}
\setsansfont{Helvetica Neue}
\setmonofont{Menlo}
\IfFileExists{/Users/iw/Library/Fonts/kaiu.ttf}{
  \setCJKmainfont[Path=/Users/iw/Library/Fonts/,AutoFakeBold=4]{kaiu.ttf}
}{
  \IfFontExistsTF{標楷體}{
    \setCJKmainfont[AutoFakeBold=4]{標楷體}
  }{
    \setCJKmainfont{Songti TC}
  }
}
\IfFontExistsTF{Songti TC}{
  \setCJKfamilyfont{song}{Songti TC}
  \setCJKmonofont{Songti TC}
}{
  \setCJKfamilyfont{song}[Path=/Users/iw/Library/Fonts/,AutoFakeBold=4]{kaiu.ttf}
  \setCJKmonofont[Path=/Users/iw/Library/Fonts/,AutoFakeBold=4]{kaiu.ttf}
}
\newcommand{\songfont}{\CJKfamily{song}}

% ===== 封面／目錄專用打字機字體（僅限封面、目錄頁使用，不影響正文）=====
\IfFontExistsTF{Radio Newsman}{%
  \newfontfamily\bverfgetitlede{Radio Newsman}%
}{%
  \newfontfamily\bverfgetitlede{Courier New}%
}
\IfFontExistsTF{Erikas Farbband}{%
  \newfontfamily\bverfgebodyde{Erikas Farbband}%
}{%
  \let\bverfgebodyde\bverfgetitlede
}
\IfFontExistsTF{Maoken Heavy Labourer}{%
  \setCJKfamilyfont{bverfgetitlecjk}{Maoken Heavy Labourer}%
}{%
  \setCJKfamilyfont{bverfgetitlecjk}{Songti TC}%
}
\IfFontExistsTF{Huiwen-mincho}{%
  \setCJKfamilyfont{bverfgebodycjk}{Huiwen-mincho}%
}{%
  \setCJKfamilyfont{bverfgebodycjk}{Songti TC}%
}
\newcommand{\bverfgetitlezh}{\bverfgetitlede\CJKfamily{bverfgetitlecjk}}
\newcommand{\bverfgebodyzh}{\bverfgebodyde\CJKfamily{bverfgebodycjk}}

\setstretch{1.35}
\setlist{itemsep=0.2em, topsep=0.4em}
\setlength{\parindent}{0pt}
\setlength{\parskip}{0.45em}
\raggedbottom
\setcounter{secnumdepth}{0}
\setcounter{tocdepth}{2}

\newcommand{\lawboxsourcefont}{\footnotesize}
\newcommand{\lawboxtransfont}{\footnotesize}
\newcommand{\lawboxsourcestretch}{1.02}
\newcommand{\lawboxtransstretch}{0.92}

\titleformat{\section}{\centering\Large\bfseries\songfont}{\thesection}{1em}{}
\titleformat{\subsection}{\large\bfseries\songfont}{\thesubsection}{1em}{}
\titleformat{\subsubsection}{\normalsize\bfseries\songfont}{\thesubsubsection}{1em}{}
\titleformat{\paragraph}{\normalsize\bfseries\songfont}{\theparagraph}{1em}{}
\titlespacing*{\section}{0pt}{1.8em}{1.0em}
\titlespacing*{\subsection}{0pt}{1.2em}{0.6em}
\titlespacing*{\subsubsection}{0pt}{1.0em}{0.45em}
\titlespacing*{\paragraph}{0pt}{0.6em}{0.4em}

\definecolor{noteblue}{HTML}{A9C9F5}
\definecolor{noteteal}{HTML}{AEE1D6}
\definecolor{notegray}{HTML}{D7DADF}
\definecolor{caseink}{HTML}{000000}
\definecolor{casepaper}{HTML}{FCFEFD}
\definecolor{sourcepaper}{HTML}{FAFBF8}
\definecolor{sourceline}{HTML}{D7DED8}
\definecolor{sourceink}{HTML}{000000}
\definecolor{pendingink}{HTML}{000000}

\tcbset{
  caseboxbase/.style={
    enhanced,
    breakable,
    width=0.985\linewidth,
    colback=casepaper,
    coltitle=caseink,
    fonttitle=\bfseries\songfont,
    boxrule=0.28pt,
    arc=1.2mm,
    left=2.4mm,
    right=2.4mm,
    top=1.5mm,
    bottom=1.5mm,
    boxsep=1.5mm,
    before skip=0.65em,
    after skip=0.65em,
    before upper={\setlength{\parindent}{0pt}\setlength{\parskip}{0.35em}},
  }
}

\newcommand{\bilingualstart}{}
\newcommand{\bilingualstop}{}
\ifbilingualcolumns
  \renewcommand{\bilingualstart}{\small\setlength{\columnsep}{1.25cm}\setlength{\columnseprule}{0.12pt}\columnratio{0.5,0.5}\multicolumnfootnotes\flushbottom\sloppy\interlinepenalty=80\clubpenalty=800\widowpenalty=800\displaywidowpenalty=800\begin{paracol}{2}\global\intranslationfalse\global\firstbilingualsectiontrue{\footnotesize\color{sourceink}Deutsch\par}\switchcolumn{\footnotesize\songfont\color{sourceink}中文譯文\par}\switchcolumn*}
  \renewcommand{\bilingualstop}{\ifintranslation\switchcolumn*\global\intranslationfalse\fi\end{paracol}\clearpage\raggedbottom}
\fi
\newif\ifrandnumbers
\randnumbersfalse
\newcounter{randnumber}
\newcounter{randnumberlocal}
\newcounter{lastsourcerandstart}
\newif\iflastsourcehasrandnumbers
\lastsourcehasrandnumbersfalse
\newif\ifinlawbox
\inlawboxfalse
\newif\iflawboxhaspendingrn
\lawboxhaspendingrnfalse
\newcommand{\lawboxpendingrn}{}
\newcommand{\startrandnumbers}{\global\randnumberstrue\setcounter{randnumber}{1}\global\lastsourcehasrandnumbersfalse}
\newcommand{\stoprandnumbers}{\global\randnumbersfalse\global\lastsourcehasrandnumbersfalse}
\newlength{\randnumberwidth}
\setlength{\randnumberwidth}{0.95cm}
\newlength{\randnumberrightoffset}
\setlength{\randnumberrightoffset}{0.85cm}
\newlength{\randnumberlawboxrightoffset}
\setlength{\randnumberlawboxrightoffset}{1.40cm}
\newcommand{\randnumberbox}[1]{%
  \leavevmode
  \makebox[0pt][l]{%
    \smash{%
      \tikz[remember picture,overlay,baseline=(rnmark.base)]{%
        \coordinate (rnmark) at (0,0);%
        \ifinlawbox
          \path let \p1=(current page.east), \p2=(rnmark.base) in
            node[
              anchor=base east,
              inner sep=0pt,
              outer sep=0pt,
              text width=\randnumberwidth,
              align=right,
              text=pendingink,
              font=\normalfont\mdseries\scriptsize\ttfamily
            ] at (\x1-\randnumberlawboxrightoffset,\y2) {{\normalfont\mdseries\scriptsize\ttfamily #1}};%
        \else
          \path let \p1=(current page.east), \p2=(rnmark.base) in
            node[
              anchor=base east,
              inner sep=0pt,
              outer sep=0pt,
              text width=\randnumberwidth,
              align=right,
              text=pendingink,
              font=\normalfont\mdseries\scriptsize\ttfamily
            ] at (\x1-\randnumberrightoffset,\y2) {{\normalfont\mdseries\scriptsize\ttfamily #1}};%
        \fi
      }%
    }%
  }%
  \ignorespaces
}
\newcommand{\lawboxstorerandnumber}[1]{%
  \gdef\lawboxpendingrn{#1}%
  \global\lawboxhaspendingrntrue
  \ignorespaces
}
\newcommand{\lawboxuserandnumber}{%
  \iflawboxhaspendingrn
    \randnumberbox{\lawboxpendingrn}%
    \global\lawboxhaspendingrnfalse
  \fi
}
\newcommand{\sourcern}[1]{%
  \ifrandnumbers
    \ifshoworiginal
      \ifshowtranslation\else
        \ifinlawbox\lawboxstorerandnumber{#1}\else\randnumberbox{#1}\fi
      \fi
    \fi
  \fi
}
\newcommand{\transrn}[1]{%
  \ifrandnumbers
    \ifshowtranslation
      \ifinlawbox\lawboxstorerandnumber{#1}\else\randnumberbox{#1}\fi
    \fi
  \fi
}
% ===== 課堂模式：德文原文縮寫註解 =====
% 日常切換只改各判決檔的一行：
%   \classroommodeon        開啟課堂版；註解掉這行即回到一般版。
% 模板預設：
%   1. 課堂模式關閉。
%   2. 課堂模式開啟時，縮寫註解包含「德文全稱＋中文譯名」。
%   3. 課堂模式開啟時，GG/條文編者註仍會自動顯示。
% 進階微調才需要在單案檔覆寫：
%   \classroomabbrzhoff     縮寫註解僅顯示德文全稱。
%   \showeditornotestrue    一般版也顯示編者註。
% 註解策略：同一實際 PDF 頁面中，同一縮寫只在第一次出現處加註。
% 這裡用 zref-abspage 的絕對頁序，不用 \thepage，避免文件內頁碼重置時誤判。
\newif\ifclassroommode
\newif\ifclassroomabbrzh
\classroommodefalse
\classroomabbrzhtrue
\newcommand{\classroommodeon}{\classroommodetrue}
\newcommand{\classroommodeoff}{\classroommodefalse}
\newcommand{\classroomabbrzhon}{\classroomabbrzhtrue}
\newcommand{\classroomabbrzhoff}{\classroomabbrzhfalse}
\newcommand{\abbrnotelabel}{\emph{Abk.}: }
\newcommand{\abbrnote}[3]{%
  \ifclassroommode
    \ifshoworiginal
      \ifintranslation\else
        \footnote{\normalfont\footnotesize\abbrnotelabel #2\ifclassroomabbrzh；{\songfont #3}\fi}%
      \fi
    \fi
  \fi
}
\newcommand{\abbrnoteonce}[4]{%
  #2%
  \ifclassroommode
    \ifshoworiginal
      \ifintranslation\else
        \begingroup
          \edef\abbrcurrentabspage{\number\numexpr\value{abspage}+1\relax}%
          \ifcsname abbrnoteseen@#1@\abbrcurrentabspage\endcsname\else
            \global\expandafter\let\csname abbrnoteseen@#1@\abbrcurrentabspage\endcsname\empty
            \abbrnote{#1}{#3}{#4}%
          \fi
        \endgroup
      \fi
    \fi
  \fi
}
\newcommand{\abbrAbs}{\abbrnoteonce{abs}{Abs.}{Absatz}{項}}
\newcommand{\abbrAAO}{\abbrnoteonce{aao}{a.a.O.}{am angegebenen Ort}{前引處}}
\newcommand{\abbrAE}{\abbrnoteonce{ae}{AE}{Alternativ-Entwurf}{替代草案}}
\newcommand{\abbrArt}{\abbrnoteonce{art}{Art.}{Artikel}{條}}
\newcommand{\abbrAufl}{\abbrnoteonce{aufl}{Aufl.}{Auflage}{版}}
\newcommand{\abbrBd}{\abbrnoteonce{bd}{Bd.}{Band}{卷}}
\newcommand{\abbrBGBl}{\abbrnoteonce{bgbl}{BGBl.}{Bundesgesetzblatt}{聯邦法律公報}}
\newcommand{\abbrBGHSt}{\abbrnoteonce{bghst}{BGHSt}{Entscheidungen des Bundesgerichtshofes in Strafsachen}{聯邦普通法院刑事判決彙編}}
\newcommand{\abbrBRDrucks}{\abbrnoteonce{brdrucks}{BRDrucks.}{Bundesrats-Drucksache}{聯邦參議院文件}}
\newcommand{\abbrBTDrucks}{\abbrnoteonce{btdrucks}{BTDrucks.}{Bundestags-Drucksache}{聯邦議院文件}}
\newcommand{\abbrBundesgesetzbl}{\abbrnoteonce{bundesgesetzbl}{Bundesgesetzbl.}{Bundesgesetzblatt}{聯邦法律公報}}
\newcommand{\abbrBvF}{\abbrnoteonce{bvf}{BvF}{Bundesverfassungsgerichtsverfahren}{聯邦憲法法院程序案號}}
\newcommand{\abbrBvL}{\abbrnoteonce{bvl}{BvL}{Bundesverfassungsgerichtsverfahren}{聯邦憲法法院程序案號}}
\newcommand{\abbrBvR}{\abbrnoteonce{bvr}{BvR}{Bundesverfassungsgerichtsverfahren}{聯邦憲法法院程序案號}}
\newcommand{\abbrBVerfG}{\abbrnoteonce{bverfg}{BVerfG}{Bundesverfassungsgericht}{聯邦憲法法院}}
\newcommand{\abbrBVerfGE}{\abbrnoteonce{bverfge}{BVerfGE}{Entscheidungen des Bundesverfassungsgerichts}{聯邦憲法法院判決集}}
\newcommand{\abbrBVerfGG}{\abbrnoteonce{bverfgg}{BVerfGG}{Bundesverfassungsgerichtsgesetz}{聯邦憲法法院法}}
\newcommand{\abbrDDR}{\abbrnoteonce{ddr}{DDR}{Deutsche Demokratische Republik}{德意志民主共和國；東德}}
\newcommand{\abbrDJT}{\abbrnoteonce{djt}{DJT}{Deutscher Juristentag}{德國法學家大會}}
\newcommand{\abbrDr}{\abbrnoteonce{dr}{Dr.}{Doktor}{Dr.}}
\newcommand{\abbrEuGRZ}{\abbrnoteonce{eugrz}{EuGRZ}{Europäische Grundrechte-Zeitschrift}{歐洲基本權期刊}}
\newcommand{\abbrEV}{\abbrnoteonce{ev}{EV}{Einigungsvertrag}{統一條約}}
\newcommand{\abbrF}{\abbrnoteonce{f}{f.}{folgende}{以下一頁；下一段}}
\newcommand{\abbrFF}{\abbrnoteonce{ff}{ff.}{fortfolgende}{以下數頁；以下各段}}
\newcommand{\abbrGBlDDR}{\abbrnoteonce{gblddr}{GBl. DDR}{Gesetzblatt der Deutschen Demokratischen Republik}{德意志民主共和國法律公報}}
\newcommand{\abbrGG}{\abbrnoteonce{gg}{GG}{Grundgesetz}{基本法}}
\newcommand{\abbrGVBl}{\abbrnoteonce{gvbl}{GVBl.}{Gesetz- und Verordnungsblatt}{法律及命令公報}}
\newcommand{\abbrJR}{\abbrnoteonce{jr}{JR}{Juristische Rundschau}{法學評論}}
\newcommand{\abbrJZ}{\abbrnoteonce{jz}{JZ}{JuristenZeitung}{法學家報}}
\newcommand{\abbrIVm}{\abbrnoteonce{ivm}{i.V.m.}{in Verbindung mit}{結合；連同}}
\newcommand{\abbrMdB}{\abbrnoteonce{mdb}{MdB}{Mitglied des Bundestages}{聯邦議院議員}}
\newcommand{\abbrMwN}{\abbrnoteonce{mwn}{m.w.N.}{mit weiteren Nachweisen}{附進一步文獻／判例出處}}
\newcommand{\abbrNF}{\abbrnoteonce{nf}{n.F.}{neue Fassung}{新版本}}
\newcommand{\abbrAF}{\abbrnoteonce{af}{a.F.}{alte Fassung}{舊版本}}
\newcommand{\abbrNr}{\abbrnoteonce{nr}{Nr.}{Nummer}{款、目或號，依語境}}
\newcommand{\abbrProf}{\abbrnoteonce{prof}{Prof.}{Professor}{教授}}
\newcommand{\abbrRdnr}{\abbrnoteonce{rdnr}{Rdnr.}{Randnummer}{邊碼；段碼}}
\newcommand{\abbrRdnrn}{\abbrnoteonce{rdnrn}{Rdnrn.}{Randnummern}{邊碼；段碼}}
\newcommand{\abbrRGBl}{\abbrnoteonce{rgbl}{RGBl.}{Reichsgesetzblatt}{帝國法律公報}}
\newcommand{\abbrRGSt}{\abbrnoteonce{rgst}{RGSt}{Entscheidungen des Reichsgerichts in Strafsachen}{帝國法院刑事判決彙編}}
\newcommand{\abbrRspr}{\abbrnoteonce{rspr}{Rspr.}{Rechtsprechung}{判例；司法實務}}
\newcommand{\abbrRVO}{\abbrnoteonce{rvo}{RVO}{Reichsversicherungsordnung}{帝國保險法}}
\newcommand{\abbrS}{\abbrnoteonce{s}{S.}{Seite}{頁}}
\newcommand{\abbrSFHG}{\abbrnoteonce{sfhg}{SFHG}{Schwangeren- und Familienhilfegesetz}{孕婦及家庭扶助法}}
\newcommand{\abbrSGBV}{\abbrnoteonce{sgbv}{SGB V}{Fünftes Buch Sozialgesetzbuch}{社會法典第五編}}
\newcommand{\abbrStAG}{\abbrnoteonce{staeg}{StÄG}{Strafrechtsänderungsgesetz}{刑法修正法}}
\newcommand{\abbrStenBer}{\abbrnoteonce{stenber}{StenBer.}{Stenographischer Bericht}{速記紀錄}}
\newcommand{\abbrStGB}{\abbrnoteonce{stgb}{StGB}{Strafgesetzbuch}{刑法}}
\newcommand{\abbrStREG}{\abbrnoteonce{streg}{StREG}{Strafrechtsreform-Ergänzungsgesetz}{刑法改革補充法}}
\newcommand{\abbrStrRG}{\abbrnoteonce{strrg}{StrRG}{Strafrechtsreformgesetz}{刑法改革法}}
\newcommand{\abbrUS}{\abbrnoteonce{us}{U.S.}{United States Reports}{美國最高法院判決彙編}}
\newcommand{\abbrVgl}{\abbrnoteonce{vgl}{vgl.}{vergleiche}{參見}}
\newcommand{\abbrVVDStRL}{\abbrnoteonce{vvdstrl}{VVDStRL}{Veröffentlichungen der Vereinigung der Deutschen Staatsrechtslehrer}{德國公法教師協會論文集}}
\newcommand{\abbrWP}{\abbrnoteonce{wp}{WP.}{Wahlperiode}{屆期}}
\newcommand{\abbrWp}{\abbrnoteonce{wp}{Wp.}{Wahlperiode}{屆期}}
\newcommand{\abbrZStrW}{\abbrnoteonce{zstrw}{ZStrW}{Zeitschrift für die gesamte Strafrechtswissenschaft}{整體刑法學期刊}}
\newcommand{\recordsourcerandstart}{%
  \ifrandnumbers
    \setcounter{lastsourcerandstart}{\value{randnumber}}%
    \global\lastsourcehasrandnumberstrue
  \else
    \global\lastsourcehasrandnumbersfalse
  \fi
}
\newlength{\biparagraphskip}
\setlength{\biparagraphskip}{0.52em}
\newlength{\lawboxblockbeforeskip}
\setlength{\lawboxblockbeforeskip}{0.72em}
\newlength{\lawboxblockafterskip}
\setlength{\lawboxblockafterskip}{0.92em}
\newif\iflawboxpartendedblock
\lawboxpartendedblockfalse
\newcommand{\sourceparstyle}{\footnotesize\color{sourceink}\setstretch{1.12}\RaggedRight\emergencystretch=3em\setlength{\parskip}{0.42em}\obeylines\selectfont}
\newcommand{\transparstyle}{\songfont\color{caseink}\setstretch{1.12}\setlength{\parindent}{0pt}\setlength{\parskip}{0.42em}}
\newcommand{\bilingualpairskip}{%
  \par\addvspace{\biparagraphskip}%
  \switchcolumn
  \par\addvspace{\biparagraphskip}%
  \switchcolumn*%
  \global\intranslationfalse
}
\newcommand{\bilingualboxpairskip}{%
  \par
  \switchcolumn
  \par
  \switchcolumn*%
  \global\intranslationfalse
}
\newcommand{\bilinguallawboxblockskip}{%
  \switchcolumn*%
  \par\addvspace{\lawboxblockafterskip}%
  \switchcolumn
  \par\addvspace{\lawboxblockafterskip}%
  \switchcolumn*%
  \global\intranslationfalse
}
\newcommand{\bikeepwithnext}[1][4]{%
  \ifbilingualcolumns
    \syncleft
    \Needspace{#1\baselineskip}%
    \switchcolumn
    \Needspace{#1\baselineskip}%
    \switchcolumn*%
    \global\intranslationfalse
  \else
    \Needspace{#1\baselineskip}%
  \fi
}
\NewEnviron{sourcepara}[1][]{
  \recordsourcerandstart
  \ifshoworiginal
    \ifbilingualcolumns
      \ifintranslation\switchcolumn*\global\intranslationfalse\fi
    \else
      \Needspace{5\baselineskip}
    \fi
    \ifbilingualcolumns\else\par\addvspace{0.35em}\fi
    {\sourceparstyle
    \noindent\BODY\par}
    \ifbilingualcolumns\else\addvspace{0.25em}\fi
    \ifbilingualcolumns
      \switchcolumn
      \global\intranslationtrue
    \fi
  \else
    \ifrandnumbers
      \begingroup
        \setbox0=\vbox{\hsize=\linewidth\sourceparstyle\noindent\BODY\par}%
      \endgroup
    \fi
  \fi
}
\NewEnviron{transpara}{
  \ifshowtranslation
    \setcounter{randnumberlocal}{\value{lastsourcerandstart}}%
    \ifbilingualcolumns
      \ifintranslation\else\switchcolumn\global\intranslationtrue\fi
      {\transparstyle\setlength{\leftskip}{0.55em}\setlength{\rightskip}{0.15em plus 1em}\addtolength{\linewidth}{-0.55em}\BODY\par}
      \switchcolumn*
      \bilingualpairskip
    \else
      {\transparstyle\BODY\par}
    \fi
  \else
    \ifbilingualcolumns
      \ifintranslation\switchcolumn*\global\intranslationfalse\fi
    \fi
  \fi
}
\NewEnviron{sourceboxpara}[1][]{
  \global\lawboxpartendedblockfalse
  \recordsourcerandstart
  \ifshoworiginal
    \ifbilingualcolumns
      \ifintranslation\switchcolumn*\global\intranslationfalse\fi
    \else
      \Needspace{3\baselineskip}
    \fi
    {\sourceparstyle
    \BODY\par}
    \ifbilingualcolumns\else
      \iflawboxpartendedblock\addvspace{\lawboxblockafterskip}\fi
    \fi
    \ifbilingualcolumns
      \switchcolumn
      \global\intranslationtrue
    \fi
  \fi
}
\NewEnviron{transboxpara}{
  \global\lawboxpartendedblockfalse
  \ifshowtranslation
    \setcounter{randnumberlocal}{\value{lastsourcerandstart}}%
    \ifbilingualcolumns
      \ifintranslation\else\switchcolumn\global\intranslationtrue\fi
      {\transparstyle\BODY\par}
      \iflawboxpartendedblock
        \bilinguallawboxblockskip
      \else
        \switchcolumn*
        \bilingualboxpairskip
      \fi
    \else
      {\transparstyle\BODY\par}
      \iflawboxpartendedblock\addvspace{\lawboxblockafterskip}\fi
    \fi
  \else
    \ifbilingualcolumns
      \ifintranslation\switchcolumn*\global\intranslationfalse\fi
    \fi
  \fi
}
\newcommand{\leitsatznum}[1]{\noindent\hangindent=3.0em\hangafter=1\makebox[2.25em][r]{\bfseries #1.}\hspace{0.55em}\ignorespaces}
\newcommand{\syncleft}{\ifbilingualcolumns\ifintranslation\switchcolumn*\global\intranslationfalse\fi\fi}
\pretocmd{\section}{\syncleft\clearpage}{}{}
\pretocmd{\subsection}{\syncleft}{}{}
\pretocmd{\subsubsection}{\syncleft}{}{}
\newcommand{\biheadingfont}{\songfont}
\newcommand{\bitocentry}[2]{%
  \ifshoworiginal
    \ifshowtranslation
      \ifstrequal{#1}{#2}{#1}{#1\space #2}%
    \else
      #1%
    \fi
  \else
    #2%
  \fi
}
\pdfstringdefDisableCommands{%
  \def\bitocentry#1#2{#1}%
}
\newcommand{\biheading}[7]{%
  \syncleft#1\bikeepwithnext[#3]\phantomsection\addcontentsline{toc}{#2}{\bitocentry{#6}{#7}}%
  \ifbilingualcolumns
    {#4 #6\par}%
    \switchcolumn
    {#4\biheadingfont #7\par}%
    \switchcolumn*%
    \global\intranslationfalse
    \par\addvspace{#5}%
    \switchcolumn
    \par\addvspace{#5}%
    \switchcolumn*%
  \else
    \ifshoworiginal{#4 #6\par}\fi
    \ifshowtranslation{#4\biheadingfont #7\par}\fi
    \addvspace{#5}%
  \fi
}
\newcommand{\termsectionheading}[1]{%
  \par\Needspace{9\baselineskip}%
  \vspace{0.65em}%
  {\songfont\bfseries #1\par}%
  \vspace{0.2em}%
}
% 章節換頁：
%   雙語 paracol 欄內不要用 \clearpage；它會在同步兩欄時吐出只有欄線/頁碼的空頁。
%   \newpage 足以讓下一個 section 起新頁，又不會強制清空所有待排材料。
%   單語模式不在 paracol 內，仍可用 \clearpage。
\newcommand{\bisectionpagebreak}{\ifbilingualcolumns\newpage\else\clearpage\fi}
\newcommand{\bisection}[2]{%
  \biheading{\iffirstbilingualsection\global\firstbilingualsectionfalse\else\bisectionpagebreak\fi}{section}{6}{\centering\Large\bfseries}{0.8em}{#1}{#2}%
}
\newcommand{\bisubsection}[2]{%
  \biheading{}{subsection}{5}{\large\bfseries}{0.55em}{#1}{#2}%
}
\newcommand{\bisubsubsection}[2]{%
  \biheading{}{subsubsection}{4}{\normalsize\bfseries}{0.45em}{#1}{#2}%
}
\newcommand{\bicompositeheading}[4]{%
  \biheading{}{subsection}{5}{\large\bfseries}{0.16em}{#1}{#2}%
  \biheading{}{subsubsection}{3}{\normalsize\bfseries}{0.45em}{#3}{#4}%
}
\newif\iflawboxfirstitem
\lawboxfirstitemfalse
\newcommand{\lawboxprespace}[1]{%
  \par
  \iflawboxfirstitem
    \global\lawboxfirstitemfalse
  \else
    \addvspace{#1}%
  \fi
}
\newtcolorbox{lawboxinner}{
  caseboxbase,
  colframe=sourceline!85!black,
  colback=white,
  left=1.05em,
  right=1.05em,
  top=0.62em,
  bottom=0.62em,
  before upper={\global\inlawboxtrue\global\lawboxfirstitemtrue\global\lawboxhaspendingrnfalse\ifintranslation\lawboxtransfont\else\lawboxsourcefont\fi\RaggedRight\setlength{\parindent}{0pt}\setlength{\parskip}{0pt}\ifintranslation\setstretch{\lawboxtransstretch}\else\setstretch{\lawboxsourcestretch}\fi},
  after upper={\global\lawboxhaspendingrnfalse\global\inlawboxfalse}
}
\tcbset{
  lawboxpartbase/.style={
    enhanced,
    breakable=false,
    width=0.985\linewidth,
    colback=white,
    frame hidden,
    boxrule=0pt,
    arc=0pt,
    left=1.05em,
    right=1.05em,
    boxsep=1.5mm,
    before skip=0pt,
    after skip=0pt,
    before upper={\global\inlawboxtrue\global\lawboxfirstitemtrue\global\lawboxhaspendingrnfalse\ifintranslation\lawboxtransfont\else\lawboxsourcefont\fi\RaggedRight\setlength{\parindent}{0pt}\setlength{\parskip}{0pt}\ifintranslation\setstretch{\lawboxtransstretch}\else\setstretch{\lawboxsourcestretch}\fi},
    after upper={\global\lawboxhaspendingrnfalse\global\inlawboxfalse},
    borderline west={0.28pt}{0pt}{sourceline!85!black},
    borderline east={0.28pt}{0pt}{sourceline!85!black},
  },
  lawboxpartfirst/.style={
    lawboxpartbase,
    top=0.62em,
    bottom=0.04em,
    before skip=\lawboxblockbeforeskip,
    borderline north={0.28pt}{0pt}{sourceline!85!black},
  },
  lawboxpartmiddle/.style={
    lawboxpartbase,
    top=0.04em,
    bottom=0.04em,
  },
  lawboxpartlast/.style={
    lawboxpartbase,
    top=0.04em,
    bottom=0.62em,
    after skip=0pt,
    borderline south={0.28pt}{0pt}{sourceline!85!black},
  }
}
\newtcolorbox{lawboxpartinner}[1]{#1}
\newtcolorbox{termboxinner}{
  caseboxbase,
  colframe=noteblue!70!black,
  colback=casepaper,
  before upper={\songfont\setlength{\parindent}{0pt}\setlength{\parskip}{0.35em}}
}
\newtcolorbox{deboxinner}{
  caseboxbase,
  colframe=notegray!72!black,
  colback=white,
  left=1.05em,
  right=1.05em,
  top=0.62em,
  bottom=0.62em,
  before upper={\global\inlawboxtrue\global\lawboxfirstitemtrue\global\lawboxhaspendingrnfalse\small\setlength{\parindent}{0pt}\setlength{\parskip}{0.35em}},
  after upper={\global\lawboxhaspendingrnfalse\global\inlawboxfalse}
}
\NewEnviron{lawbox}{
  \begin{lawboxinner}\BODY\end{lawboxinner}
}
\NewEnviron{lawboxpart}[2]{
  \begin{lawboxpartinner}{lawboxpart#1,equal height group=#2}\BODY\end{lawboxpartinner}%
  \ifstrequal{#1}{last}{\global\lawboxpartendedblocktrue}{}%
}
\NewEnviron{termbox}{
  \par\addvspace{0.55em}
  {\color{noteblue!70!black}\rule{\linewidth}{0.28pt}}\par
  \begingroup
    \songfont\small\color{caseink}
    \setlength{\parindent}{0pt}
    \setlength{\parskip}{0.24em}
    \BODY
    \par
  \endgroup
  {\color{noteblue!70!black}\rule{\linewidth}{0.28pt}}\par
  \addvspace{0.55em}
}
\NewEnviron{debox}{
  \begin{deboxinner}\BODY\end{deboxinner}
}
\providecommand{\paragraphmark}[1]{}
\newcommand{\lawboxtitleafter}{\addvspace{0.06em}}
\newcommand{\lawtitle}[1]{\lawboxprespace{0.22em}{\lawboxtransfont\bfseries\lawboxuserandnumber #1}\par\lawboxtitleafter}
\newcommand{\absno}[2]{\lawboxprespace{0.04em}\noindent\lawboxuserandnumber\hangindent=2.6em\hangafter=1\makebox[2.6em][l]{(#1)}#2\par}
\newcommand{\nrno}[2]{\lawboxprespace{0pt}\noindent\lawboxuserandnumber\hspace*{2.6em}\hangindent=4.4em\hangafter=1\makebox[1.8em][l]{#1.}#2\par}
\newcommand{\letterno}[2]{\lawboxprespace{0pt}\noindent\lawboxuserandnumber\hspace*{4.4em}\hangindent=6.0em\hangafter=1\makebox[1.6em][l]{#1)}#2\par}
\newcommand{\sourcelawtitle}[1]{\lawboxprespace{0.22em}{\lawboxsourcefont\bfseries\lawboxuserandnumber #1}\par\lawboxtitleafter}
\newcommand{\sourcelawsubtitle}[1]{\lawboxprespace{0.04em}{\lawboxsourcefont\bfseries\lawboxuserandnumber #1}\par\lawboxtitleafter}
\newcommand{\sourceabsno}[2]{\lawboxprespace{0.04em}\noindent\lawboxuserandnumber\hangindent=2.45em\hangafter=1\makebox[2.45em][l]{(#1)}#2\par}
\newcommand{\sourcenrno}[2]{\lawboxprespace{0pt}\noindent\lawboxuserandnumber\hspace*{2.45em}\hangindent=4.25em\hangafter=1\makebox[1.8em][l]{#1.}#2\par}
\newcommand{\sourceletterno}[2]{\lawboxprespace{0pt}\noindent\lawboxuserandnumber\hspace*{4.25em}\hangindent=5.85em\hangafter=1\makebox[1.6em][l]{#1)}#2\par}
\newcommand{\sourcequotepara}[1]{\lawboxprespace{0.04em}\noindent\lawboxuserandnumber\hangindent=1.2em\hangafter=1 #1\par}
\newcommand{\sourcelawline}[1]{\lawboxprespace{0.04em}\noindent\lawboxuserandnumber #1\par}
\newcommand{\lawline}[1]{\lawboxprespace{0.04em}\noindent\lawboxuserandnumber #1\par}
\newcommand{\pendingtranslation}{\par{\songfont\footnotesize\color{pendingink}（中文譯文待補。）}\par}
\newcommand{\tocpart}[1]{\par\addvspace{0.52em}{\footnotesize\bfseries\color{pendingink}#1\par}\addvspace{0.16em}}
\newcommand{\tocdashline}{\par\addvspace{0.48em}\noindent{\color{notegray}\xleaders\hbox to 0.82em{\hfil\rule{0.42em}{0.28pt}\hfil}\hskip\linewidth}\par\addvspace{0.38em}}
% \tocpanel{font-cmd}{content}：將目錄置中於所在欄寬（雙語欄適用）。
% A3 橫式使用 0.58\linewidth，A4 橫式使用 0.84\linewidth，
% 使兩種紙張下目錄欄內容實際寬度相近（均約 100–105 mm）。
\newcommand{\tocpanel}[2]{%
  \makebox[\linewidth][c]{%
    \ifx\bilingualpaper\bilingualpaperaThree
      \begin{minipage}{0.58\linewidth}%
    \else
      \begin{minipage}{0.84\linewidth}%
    \fi
    \toccolstyle
    #1#2%
    \end{minipage}%
  }%
}
% ===== 首頁簡目錄的兩層 description list 縮排 =====
% enumitem 的 description list 可把每一列拆成：
%   [label]  labelsep  body text
% 這裡的「起點」都以所在 minipage/欄位的左緣為 0。
%
% 核心規則：
%   labelindent = label 欄位從左緣開始的位置。
%   labelwidth  = label 欄位本身寬度；標籤太長（如 Abw. Meinung）就加大它。
%   labelsep    = label 欄位右緣到正文的距離。
%   leftmargin  = 正文 body text 的起點。判斷層級視覺時主要看這個。
%
% 所以：
%   想讓「Begründung / 判決理由」整列正文往右或往左：改 tocitems 的 leftmargin。
%   想讓「A. Verfahren und Gegenstand / A. 程序與審查標的」正文往右或往左：
%     改 tocsubitems 的 leftmargin。
%   想讓 A./B./C. 這些標籤本身往右或往左，但正文不動：改 tocsubitems 的 labelindent。
%   想讓 A./B./C. 與正文間距變大或變小：改 tocsubitems 的 labelsep。
%
% 層級原則：
%   tocsubitems.leftmargin 必須大於 tocitems.leftmargin，
%   否則子層級正文會看起來比「Gründe / 理由」還靠左。
\newlist{tocitems}{description}{1}
\setlist[tocitems]{%
  labelindent=0pt,    % 第一層 label 欄位起點；0pt 表示貼齊目錄欄左緣。
  leftmargin=2.54cm,  % 第一層正文起點；例如 Begründung / 判決理由 的左緣。
  labelwidth=2.16cm,  % 第一層 label 欄位寬度；需容納 Leitsätze、Abw. Meinung、不同意見等。
  labelsep=0.32cm,    % 第一層 label 與正文的水平間距；只改間距，不改正文起點。
  itemsep=0.28em,     % 第一層各項之間的垂直距離。
  topsep=0.20em,      % 第一層 list 上下與前後內容的垂直距離。
  parsep=0pt,         % 同一 item 內段落之間的距離；目錄項通常不需要。
  partopsep=0pt,      % list 若緊接段落開始時的額外距離；關閉以免目錄高度飄動。
  align=left,         % label 欄內左對齊；長短標籤不靠右貼齊。
  font=\normalfont\bfseries % label 的字型；正文不受這行影響。
}
\newlist{tocsubitems}{description}{1}
\setlist[tocsubitems]{%
  labelindent=1cm,    % 第二層 label 起點；只控制 A./B./C. 本身的位置。
  leftmargin=3.00cm,  % 第二層正文起點；必須大於 2.54cm，才會比 Begründung / 判決理由更右。
  labelwidth=0.5cm,   % 第二層 label 欄寬；A./B./C. 很短，可小於第一層。
  labelsep=1.5cm,    % 第二層 label 與正文間距；A. 與 Verfahren / 程序 的距離看這裡。
  itemsep=0.12em,     % 第二層各項較密，避免 A.–E. 佔太高。
  topsep=0.04em,      % 第二層 list 與 Gründe / 理由之間的垂直距離。
  parsep=0pt,         % 同一第二層 item 內段落距離；目錄項通常不需要。
  partopsep=0pt,      % 關閉額外段落起始距離，避免版面不穩。
  align=left,         % A./B./C. label 欄內左對齊。
  font=\normalfont\bfseries, % A./B./C. label 加粗；正文不受這行影響。
  before={}           % 不在第二層 list 前自動插入額外內容。
}
\newlist{termitems}{description}{1}
\ifbilingualcolumns
  \setlist[termitems]{%
    leftmargin=5.2cm,
    labelwidth=4.8cm,
    labelsep=0.35cm,
    itemsep=0.16em,
    topsep=0.2em,
    parsep=0pt,
    partopsep=0pt,
    font=\normalfont\bfseries,
    style=nextline
  }
\else
  \setlist[termitems]{%
    leftmargin=0pt,
    labelwidth=0pt,
    labelsep=0pt,
    itemsep=0.24em,
    topsep=0.2em,
    parsep=0pt,
    partopsep=0pt,
    font=\normalfont\bfseries,
    style=nextline
  }
\fi
\newlist{abbritems}{description}{1}
\setlist[abbritems]{%
  leftmargin=2.38cm,
  labelwidth=2.02cm,
  labelsep=0.28cm,
  itemsep=0.16em,
  topsep=0.2em,
  parsep=0pt,
  partopsep=0pt,
  align=left,
  font=\normalfont\bfseries
}
\ifbilingualcolumns
  \newenvironment{termcolumns}{\begin{multicols}{2}}{\end{multicols}}
\else
  \newenvironment{termcolumns}{}{}
\fi

% ===== 編者註腳開關 =====
% 一般版預設不顯示編者註，避免正文腳註過密。
% 若開啟 \classroommodeon，GG/條文編者註仍會自動顯示，無需另開本開關。
% 若一般版也要顯示編者註，可在單案檔 \input 後加 \showeditornotestrue。
\newif\ifshoweditornotes
\showeditornotesfalse

% ===== 目錄排版輔助 =====
\newcommand{\toccolstyle}{\raggedright\arraybackslash}
\newcommand{\tocsingleindent}{\leftskip=1.45cm\rightskip=1.2cm}

% ===== 行內引文環境（德文原文與中文譯文各一，帶左側灰色邊線）=====
\newcommand{\quoteparbreak}{\par\addvspace{0.48em}}
\newcommand{\sourcequoteinner}[1]{%
  \begin{tcolorbox}[
    enhanced, breakable, width=\dimexpr\linewidth-0.8em\relax, frame hidden,
    colback=white, coltext=caseink, boxrule=0pt,
    borderline west={0.55pt}{0.88em}{notegray},
    left=1.78em, right=0.72em, top=0.18em, bottom=0.18em,
    boxsep=0pt, before skip=0.24em, after skip=0.24em,
    before upper={\normalfont\footnotesize\color{caseink}\setlength{\parindent}{0pt}\setlength{\parskip}{0.34em}\raggedright\setlength{\rightskip}{0pt plus 1em}}
  ]%
  #1\par
  \end{tcolorbox}%
}
\newcommand{\transquoteinner}[1]{%
  \begin{tcolorbox}[
    enhanced, breakable, width=\dimexpr\linewidth-0.8em\relax, frame hidden,
    colback=white, coltext=caseink, boxrule=0pt,
    borderline west={0.55pt}{0.88em}{notegray},
    left=1.78em, right=0.72em, top=0.18em, bottom=0.18em,
    boxsep=0pt, before skip=0.24em, after skip=0.24em,
    before upper={\songfont\small\color{caseink}\setlength{\parindent}{0pt}\setlength{\parskip}{0.34em}\raggedright\setlength{\rightskip}{0pt plus 1em}}
  ]%
  #1\par
  \end{tcolorbox}%
}

% ===== 大綱（Gliederung）Rn 輔助命令 =====
\newcommand{\outrn}[2]{\enspace{\normalfont\footnotesize\color{pendingink}(Rn\,#1–#2)}}
\newcommand{\outrnsingle}[1]{\enspace{\normalfont\footnotesize\color{pendingink}(Rn\,#1)}}
\newcommand{\outrnzh}[2]{\enspace{\normalfont\footnotesize\color{pendingink}（第\,#1–#2\,段）}}
\newcommand{\outrnzhs}[1]{\enspace{\normalfont\footnotesize\color{pendingink}（第\,#1\,段）}}
% ===== 大綱雙語對照表格輔助命令（三欄：段碼 │ 德文 │ 中文）=====
% 欄 1：段碼（由欄定義自動格式化：小字、等寬、右對齊、pendingink）
% 欄 2：德文標籤＋正文
% 欄 3：中文標籤＋正文
%
% 無段碼的列（\omainrow、\osecrow），欄 1 以 & 空置。
% 有段碼的列（\osecrowrn、\osubrow、\ointrow），#1 為預格式化字串
%   例：\osubrow{2–49}{I.}{...}{I.}{...}
%       \ointrow{107}{...}{...}
%
% \opartrow：段組標題（橫跨三欄）
\newcommand{\opartrow}[2]{%
  \noalign{\vspace{0.30em}}%
  \multicolumn{3}{@{}l@{}}{%
    \footnotesize\bfseries\color{pendingink}#1\hfill\songfont#2%
  }\\[0.06em]%
}
% \odashrow：虛線分隔（橫跨三欄）
\newcommand{\odashrow}{%
  \noalign{\vspace{0.28em}}%
  \multicolumn{3}{@{}l@{}}{\color{notegray}%
    \xleaders\hbox to 0.82em{\hfil\rule{0.42em}{0.28pt}\hfil}\hskip 0.88\linewidth}%
  \\[0.24em]%
}
% \odissentpart：不同意見書區段標頭。比一般 \opartrow 更像附錄性轉場。
\newcommand{\odissentpart}[2]{%
  \noalign{\vspace{0.42em}}%
  \multicolumn{3}{@{}p{0.88\textwidth}@{}}{%
    \color{notegray}\rule{\linewidth}{0.18pt}%
  }\\[-0.18em]%
  \multicolumn{3}{@{}p{0.88\textwidth}@{}}{%
    \makebox[\linewidth][s]{%
      \footnotesize\bfseries\color{pendingink}#1%
      \hfill{\songfont #2}%
    }%
  }\\[-0.02em]%
}
% \odissentopinion：不同意見書的署名與一行摘要，右欄保留段碼範圍。
\newcommand{\odissentopinion}[5]{%
  \footnotesize\hangindent=0.52cm\hangafter=1%
    {\bfseries\color{caseink}#2}\enspace{\color{notegray}\textperiodcentered}\enspace\color{caseink}#3%
  &\footnotesize\hangindent=0.52cm\hangafter=1%
    {\songfont\bfseries\color{caseink}#4}\enspace{\color{notegray}\textperiodcentered}\enspace\songfont\color{caseink}#5%
  &#1%
  \\[0.12em]%
}
% \omainrow：頂層項目，無段碼（Leitsätze、Urteil、Gründe …）
\newcommand{\omainrow}[4]{%
  \footnotesize\hangindent=1.92cm\hangafter=1%
    \makebox[1.92cm][l]{\bfseries\color{caseink}#1}\color{caseink}#2%
  &\footnotesize\hangindent=1.92cm\hangafter=1%
    \makebox[1.92cm][l]{\bfseries\color{caseink}#3}\songfont\color{caseink}#4%
  &%
  \\[0.15em]%
}
% \omainrowrn：頂層項目，有段碼（不同意見等標籤寬於 \osecrow 框者）
\newcommand{\omainrowrn}[5]{%
  \footnotesize\hangindent=1.92cm\hangafter=1%
    \makebox[1.92cm][l]{\bfseries\color{caseink}#2}\color{caseink}#3%
  &\footnotesize\hangindent=1.92cm\hangafter=1%
    \makebox[1.92cm][l]{\bfseries\color{caseink}#4}\songfont\color{caseink}#5%
  &#1%
  \\[0.15em]%
}
% \osecrow：一級子項，無段碼（A.、C.、D. 等有子項者）
\newcommand{\osecrow}[4]{%
  \footnotesize\hspace{1.10cm}\hangindent=1.98cm\hangafter=1%
    \makebox[0.86cm][l]{\bfseries\color{caseink}#1}\color{caseink}#2%
  &\footnotesize\hspace{1.10cm}\hangindent=1.98cm\hangafter=1%
    \makebox[0.86cm][l]{\bfseries\color{caseink}#3}\songfont\color{caseink}#4%
  &%
  \\[0.11em]%
}
% \osecrowrn：一級子項，有段碼（B.、E. 等完整段落範圍者）
% #1=段碼字串（如 101–106）, #2=DE標籤, #3=DE正文, #4=ZH標籤, #5=ZH正文
\newcommand{\osecrowrn}[5]{%
  \footnotesize\hspace{1.10cm}\hangindent=1.98cm\hangafter=1%
    \makebox[0.86cm][l]{\bfseries\color{caseink}#2}\color{caseink}#3%
  &\footnotesize\hspace{1.10cm}\hangindent=1.98cm\hangafter=1%
    \makebox[0.86cm][l]{\bfseries\color{caseink}#4}\songfont\color{caseink}#5%
  &#1%
  \\[0.11em]%
}
% \osubrow：二級子項，有段碼範圍（I.、II.、A.I. …）
% #1=段碼字串, #2=DE標籤, #3=DE正文, #4=ZH標籤, #5=ZH正文
\newcommand{\osubrow}[5]{%
  \footnotesize\hspace{1.40cm}\hangindent=2.30cm\hangafter=1%
    \makebox[0.90cm][l]{\bfseries\color{caseink}#2}\color{caseink}#3%
  &\footnotesize\hspace{1.40cm}\hangindent=2.30cm\hangafter=1%
    \makebox[0.90cm][l]{\bfseries\color{caseink}#4}\songfont\color{caseink}#5%
  &#1%
  \\[0.09em]%
}
% \ointrow：引入段落，有單一段碼，無標籤
% #1=段碼字串, #2=DE正文, #3=ZH正文
\newcommand{\ointrow}[3]{%
  \footnotesize\hspace{0.52cm}\hangindent=0.52cm\hangafter=1\color{caseink}#2%
  &\footnotesize\hspace{0.52cm}\hangindent=0.52cm\hangafter=1\songfont\color{caseink}#3%
  &#1%
  \\[0.09em]%
}
% \osubsubrow：三級子項（1.、2.、3.）
% #1=段碼字串, #2=DE標籤, #3=DE正文, #4=ZH標籤, #5=ZH正文
\newcommand{\osubsubrow}[5]{%
  \footnotesize\hspace{1.80cm}\hangindent=2.48cm\hangafter=1%
    \makebox[0.68cm][l]{\bfseries\color{caseink}#2}\color{caseink}#3%
  &\footnotesize\hspace{1.80cm}\hangindent=2.48cm\hangafter=1%
    \makebox[0.68cm][l]{\bfseries\color{caseink}#4}\songfont\color{caseink}#5%
  &#1%
  \\[0.08em]%
}
% \osubsubsubrow：四級子項（a）、b））
% #1=段碼字串, #2=DE標籤, #3=DE正文, #4=ZH標籤, #5=ZH正文
\newcommand{\osubsubsubrow}[5]{%
  \footnotesize\hspace{2.10cm}\hangindent=2.68cm\hangafter=1%
    \makebox[0.58cm][l]{\bfseries\color{caseink}#2}\color{caseink}#3%
  &\footnotesize\hspace{2.10cm}\hangindent=2.68cm\hangafter=1%
    \makebox[0.58cm][l]{\bfseries\color{caseink}#4}\songfont\color{caseink}#5%
  &#1%
  \\[0.07em]%
}
% \osubsubsubsubrow：五級子項（aa）、bb））
% 標籤框 0.62cm：足以容納「aa)」在 footnotesize bold 下（約 0.50cm）並留有間距
% #1=段碼字串, #2=DE標籤, #3=DE正文, #4=ZH標籤, #5=ZH正文
\newcommand{\osubsubsubsubrow}[5]{%
  \footnotesize\hspace{2.38cm}\hangindent=3.06cm\hangafter=1%
    \makebox[0.68cm][l]{\bfseries\color{caseink}#2}\color{caseink}#3%
  &\footnotesize\hspace{2.38cm}\hangindent=3.06cm\hangafter=1%
    \makebox[0.68cm][l]{\bfseries\color{caseink}#4}\songfont\color{caseink}#5%
  &#1%
  \\[0.06em]%
}
% \outlinepage：雙語對照大綱頁包裝環境（longtable 自動跨頁）
% 三欄：德文欄 + 中文欄 + 段碼欄（最右，不折行）
% 段碼欄寬度：1.80cm，足以容納最長段碼如「309–348」
% DE/ZH 欄寬：(0.87\textwidth - 2.08cm) / 2
%   驗算：2×欄 + 0.05\textwidth + 0.28cm gap + 1.80cm = 0.87tw - 2.08cm + 0.05tw + 2.08cm = 0.92tw ✓
\newenvironment{outlinepage}{%
  \vspace*{0.40cm}%
  \setlength{\LTleft}{0.04\textwidth}%
  \setlength{\LTright}{0.04\textwidth}%
  \setlength{\tabcolsep}{0pt}%
  % 大綱列距由 longtable 的 \arraystretch 與各 row macro 結尾的 \\[...em] 共同決定。
  % 這裡控制「每一列本身的表格 strut 高度」；數值越大，A./I./1. 各項之間越像有空行。
  % A3 原本用 0.62，視覺上沒有空行；A4 也沿用同值，避免切到 A4 後項目被撐開。
  % 若日後覺得大綱太擠，只微調這個值（例如 0.68），不要改各 \omainrow/\osubrow 的縮排。
  \renewcommand{\arraystretch}{0.62}%
  \begin{longtable}{%
    >{\vspace{0pt}\raggedright\arraybackslash}p{\dimexpr(0.87\textwidth-2.08cm)/2\relax}%
    @{\hspace{0.025\textwidth}}%
    !{\color{notegray}\vrule width 0.12pt}%
    @{\hspace{0.025\textwidth}}%
    >{\vspace{0pt}\raggedright\arraybackslash}p{\dimexpr(0.87\textwidth-2.08cm)/2\relax}%
    @{\hspace{0.28cm}}%
    >{\vspace{0pt}\footnotesize\normalfont\color{pendingink}\raggedright\arraybackslash}p{1.80cm}%
    @{}%
  }%
  % 首頁欄標籤：不要橫跨置中；分別放在德文欄與中文欄的實際欄位起點。
  \footnotesize\bfseries\color{sourceink}Gliederung%
  &\footnotesize\bfseries\songfont\color{sourceink}大綱%
  &%
  \\[0.36em]%
  \endfirsthead%
  % 續頁欄標籤：同樣分欄定位，以灰色細體區分；無需標注「續」。
  \footnotesize\color{pendingink}Gliederung%
  &\footnotesize\songfont\color{pendingink}大綱%
  &%
  \\[0.24em]%
  \endhead%
}{%
  \end{longtable}%
}
 載入。
% 不含：\documentclass、\documentversion、\ggversionde/zh、\tocde/\toczh。
% 日期與首頁版本列由本檔自動產生；各文件只需手動指定 \documentversion。

\usepackage{xeCJK}
\usepackage{fontspec}
\usepackage{geometry}
\usepackage{setspace}
\usepackage{hyperref}
\usepackage{xurl}
\usepackage{bookmark}
\usepackage{enumitem}
\usepackage{xcolor}
\usepackage{titlesec}
\usepackage{array}
\usepackage{longtable}
\usepackage{tcolorbox}
\usepackage{environ}
\usepackage{pdflscape}
\usepackage{etoolbox}
\usepackage{ragged2e}
\usepackage{paracol}
\usepackage{multicol}
\usepackage{needspace}
\usepackage{fancyhdr}
\usepackage{tikz}
\usepackage{zref-abspage}
\tcbuselibrary{breakable,skins}
\usetikzlibrary{calc}
\usepackage{pgfcalendar}

\hypersetup{
  unicode=true,
  bookmarksopen=true,
  bookmarksnumbered=false,
  hidelinks
}

% ===== 自動推導，不要手動改下面 =====
% （\docmode 與 \bilingualpaper 由各文件在 % bverfge_layout.tex
% 共用排版基礎設施，供 Schwangerschaftsabbruch I 與 II 共享。
% 由各文件以 \documentclass 後 % bverfge_layout.tex
% 共用排版基礎設施，供 Schwangerschaftsabbruch I 與 II 共享。
% 由各文件以 \documentclass 後 \input{../../共用LaTeX/bverfge_layout} 載入。
% 不含：\documentclass、\documentversion、\ggversionde/zh、\tocde/\toczh。
% 日期與首頁版本列由本檔自動產生；各文件只需手動指定 \documentversion。

\usepackage{xeCJK}
\usepackage{fontspec}
\usepackage{geometry}
\usepackage{setspace}
\usepackage{hyperref}
\usepackage{xurl}
\usepackage{bookmark}
\usepackage{enumitem}
\usepackage{xcolor}
\usepackage{titlesec}
\usepackage{array}
\usepackage{longtable}
\usepackage{tcolorbox}
\usepackage{environ}
\usepackage{pdflscape}
\usepackage{etoolbox}
\usepackage{ragged2e}
\usepackage{paracol}
\usepackage{multicol}
\usepackage{needspace}
\usepackage{fancyhdr}
\usepackage{tikz}
\usepackage{zref-abspage}
\tcbuselibrary{breakable,skins}
\usetikzlibrary{calc}
\usepackage{pgfcalendar}

\hypersetup{
  unicode=true,
  bookmarksopen=true,
  bookmarksnumbered=false,
  hidelinks
}

% ===== 自動推導，不要手動改下面 =====
% （\docmode 與 \bilingualpaper 由各文件在 \input{bverfge_layout} 之前設定；
%   預設 chinese / a4）
\ifdefined\docmode\else\def\docmode{chinese}\fi
\ifdefined\bilingualpaper\else\def\bilingualpaper{a4}\fi
\newif\ifshoworiginal
\newif\ifshowtranslation
\newif\ifbilinguallandscape
\newif\ifbilingualcolumns
\newif\ifintranslation
\newif\iffirstbilingualsection

\showoriginalfalse
\showtranslationfalse
\bilinguallandscapefalse
\bilingualcolumnsfalse
\intranslationfalse
\firstbilingualsectionfalse

\def\modechinese{chinese}
\def\modegerman{german}
\def\modebilingual{bilingual}
\def\bilingualpaperaThree{a3}
\def\bilingualpaperaFour{a4}

\ifx\docmode\modechinese
  \showtranslationtrue
\fi

\ifx\docmode\modegerman
  \showoriginaltrue
\fi

\ifx\docmode\modebilingual
  \showoriginaltrue
  \showtranslationtrue
  \bilinguallandscapetrue
  \bilingualcolumnstrue
\fi

% 編排模式說明：
% chinese  → \showoriginalfalse  \showtranslationtrue  （A4 直式，純中文）
% german   → \showoriginaltrue   \showtranslationfalse （A4 直式，純德文）
% bilingual→ \showoriginaltrue   \showtranslationtrue  \bilingualcolumnstrue（A3/A4 橫式對照）

\ifbilingualcolumns
  \ifx\bilingualpaper\bilingualpaperaFour
    \geometry{
      a4paper,
      landscape,
      left=1.25cm,
      right=2.25cm,
      top=1.25cm,
      bottom=1.75cm,
      footskip=8.5mm,
      marginparwidth=0.75cm,
      marginparsep=0.25cm
    }
  \else
    \geometry{
      a3paper,
      landscape,
      left=1.55cm,
      right=2.75cm,
      top=1.55cm,
      bottom=2.05cm,
      footskip=9.5mm,
      marginparwidth=0.9cm,
      marginparsep=0.35cm
    }
  \fi
\else
\geometry{
  a4paper,
  portrait,
  left=2.2cm,
  right=2.6cm,
  top=2.0cm,
  bottom=2.0cm,
  marginparwidth=0.8cm,
  marginparsep=0.3cm
}
\fi
% ===== 時間戳基礎設施 =====
\newcount\stamptimeval
\newcount\stampminuteval
\newcommand{\stamphh}{%
  \stamptimeval=\time
  \divide\stamptimeval by 60
  \ifnum\stamptimeval<10 0\fi
  \the\stamptimeval}
\newcommand{\stampmm}{%
  \stamptimeval=\time
  \stampminuteval=\time
  \divide\stampminuteval by 60
  \multiply\stampminuteval by 60
  \advance\stamptimeval by -\stampminuteval
  \ifnum\stamptimeval<10 0\fi
  \the\stamptimeval}
\newcommand{\stampwkde}[1]{%
  \ifcase#1Montag\or Dienstag\or Mittwoch\or Donnerstag\or Freitag\or Samstag\or Sonntag\fi}
\newcommand{\stampwkzh}[1]{%
  \ifcase#1 一\or 二\or 三\or 四\or 五\or 六\or 日\fi}
\newcommand{\stampmonthde}[1]{%
  \ifcase#1\or Januar\or Februar\or März\or April\or Mai\or Juni\or
  Juli\or August\or September\or Oktober\or November\or Dezember\fi}
\newcount\stampjdval
\newcount\stampwdval
\pgfcalendardatetojulian{\the\year-\the\month-\the\day}{\stampjdval}
\pgfcalendarjuliantoweekday{\the\stampjdval}{\stampwdval}
\newcommand{\timestampzh}{%
  \songfont 最後修改：\the\year 年\the\month 月\the\day 日%
  % （星期\stampwkzh{\the\stampwdval}）%
  % \space\stamphh:\stampmm
  }

\newcommand{\documentdatezh}{\the\year 年 \the\month 月 \the\day 日}
\newcommand{\documentdatede}{\the\day. \stampmonthde{\the\month} \the\year}
\newcommand{\bverfgetranslationnote}{%
  {\songfont 翻譯說明：初譯使用 Claude Code 與 GPT 協助迭代；仍請以德文原文為準。}%
}
\newcommand{\bverfgecourseinfo}{%
  {\songfont 課程：114 學年度第 2 學期；憲法專題研究（四）；授課老師：湯德宗教授；導讀日期：115 年 6 月 1 日。}%
}
\newcommand{\bverfgeeditorinfo}{%
  {\songfont 編者：王逸帆（R10A21126），國立台灣大學法律系研究所碩士班。}%
}
\newcommand{\bverfgetitleversionline}{%
  \ifbilingualcolumns
    Fassung \documentversion\enspace\textperiodcentered\enspace Stand: \documentdatede
    \quad
    {\songfont 版本 \documentversion\enspace\textperiodcentered\enspace 修訂日期：\documentdatezh\par}%
  \else
    \ifshowtranslation
      {\songfont 版本 \documentversion\enspace\textperiodcentered\enspace 修訂日期：\documentdatezh\par}%
    \else
      Fassung \documentversion\enspace\textperiodcentered\enspace Stand: \documentdatede\par
    \fi
  \fi
}
\newcommand{\bverfgetitlecornerinfo}{%
  \ifshowtranslation
    \vspace{0.72em}%
    \noindent\hfill
    \begin{minipage}{0.66\textwidth}
      \raggedleft\tiny\color{pendingink}\songfont
      \bverfgecourseinfo\par
      \bverfgeeditorinfo
    \end{minipage}%
  \fi
}
\newcommand{\bverfgetitlemeta}{%
  \vfill
  \begin{center}%
    \footnotesize\color{pendingink}%
    \bverfgetitleversionline
    \ifshowtranslation
      \vspace{0.28em}{\scriptsize\bverfgetranslationnote\par}%
    \fi
  \end{center}%
  \bverfgetitlecornerinfo
}

\pagestyle{fancy}
\fancyhf{}
\renewcommand{\headrulewidth}{0pt}
\renewcommand{\footrulewidth}{0pt}
\fancyfoot[C]{\thepage}
\ifbilingualcolumns
  \fancyfoot[R]{\scriptsize\color{pendingink}\timestampzh}%
\else
  \ifshoworiginal
  \else
    \fancyfoot[R]{\scriptsize\color{pendingink}\timestampzh\hspace{0.45cm}}%
  \fi
\fi
\setmainfont{Times New Roman}
\setsansfont{Helvetica Neue}
\setmonofont{Menlo}
\IfFileExists{/Users/iw/Library/Fonts/kaiu.ttf}{
  \setCJKmainfont[Path=/Users/iw/Library/Fonts/,AutoFakeBold=4]{kaiu.ttf}
}{
  \IfFontExistsTF{標楷體}{
    \setCJKmainfont[AutoFakeBold=4]{標楷體}
  }{
    \setCJKmainfont{Songti TC}
  }
}
\IfFontExistsTF{Songti TC}{
  \setCJKfamilyfont{song}{Songti TC}
  \setCJKmonofont{Songti TC}
}{
  \setCJKfamilyfont{song}[Path=/Users/iw/Library/Fonts/,AutoFakeBold=4]{kaiu.ttf}
  \setCJKmonofont[Path=/Users/iw/Library/Fonts/,AutoFakeBold=4]{kaiu.ttf}
}
\newcommand{\songfont}{\CJKfamily{song}}

% ===== 封面／目錄專用打字機字體（僅限封面、目錄頁使用，不影響正文）=====
\IfFontExistsTF{Radio Newsman}{%
  \newfontfamily\bverfgetitlede{Radio Newsman}%
}{%
  \newfontfamily\bverfgetitlede{Courier New}%
}
\IfFontExistsTF{Erikas Farbband}{%
  \newfontfamily\bverfgebodyde{Erikas Farbband}%
}{%
  \let\bverfgebodyde\bverfgetitlede
}
\IfFontExistsTF{Maoken Heavy Labourer}{%
  \setCJKfamilyfont{bverfgetitlecjk}{Maoken Heavy Labourer}%
}{%
  \setCJKfamilyfont{bverfgetitlecjk}{Songti TC}%
}
\IfFontExistsTF{Huiwen-mincho}{%
  \setCJKfamilyfont{bverfgebodycjk}{Huiwen-mincho}%
}{%
  \setCJKfamilyfont{bverfgebodycjk}{Songti TC}%
}
\newcommand{\bverfgetitlezh}{\bverfgetitlede\CJKfamily{bverfgetitlecjk}}
\newcommand{\bverfgebodyzh}{\bverfgebodyde\CJKfamily{bverfgebodycjk}}

\setstretch{1.35}
\setlist{itemsep=0.2em, topsep=0.4em}
\setlength{\parindent}{0pt}
\setlength{\parskip}{0.45em}
\raggedbottom
\setcounter{secnumdepth}{0}
\setcounter{tocdepth}{2}

\newcommand{\lawboxsourcefont}{\footnotesize}
\newcommand{\lawboxtransfont}{\footnotesize}
\newcommand{\lawboxsourcestretch}{1.02}
\newcommand{\lawboxtransstretch}{0.92}

\titleformat{\section}{\centering\Large\bfseries\songfont}{\thesection}{1em}{}
\titleformat{\subsection}{\large\bfseries\songfont}{\thesubsection}{1em}{}
\titleformat{\subsubsection}{\normalsize\bfseries\songfont}{\thesubsubsection}{1em}{}
\titleformat{\paragraph}{\normalsize\bfseries\songfont}{\theparagraph}{1em}{}
\titlespacing*{\section}{0pt}{1.8em}{1.0em}
\titlespacing*{\subsection}{0pt}{1.2em}{0.6em}
\titlespacing*{\subsubsection}{0pt}{1.0em}{0.45em}
\titlespacing*{\paragraph}{0pt}{0.6em}{0.4em}

\definecolor{noteblue}{HTML}{A9C9F5}
\definecolor{noteteal}{HTML}{AEE1D6}
\definecolor{notegray}{HTML}{D7DADF}
\definecolor{caseink}{HTML}{000000}
\definecolor{casepaper}{HTML}{FCFEFD}
\definecolor{sourcepaper}{HTML}{FAFBF8}
\definecolor{sourceline}{HTML}{D7DED8}
\definecolor{sourceink}{HTML}{000000}
\definecolor{pendingink}{HTML}{000000}

\tcbset{
  caseboxbase/.style={
    enhanced,
    breakable,
    width=0.985\linewidth,
    colback=casepaper,
    coltitle=caseink,
    fonttitle=\bfseries\songfont,
    boxrule=0.28pt,
    arc=1.2mm,
    left=2.4mm,
    right=2.4mm,
    top=1.5mm,
    bottom=1.5mm,
    boxsep=1.5mm,
    before skip=0.65em,
    after skip=0.65em,
    before upper={\setlength{\parindent}{0pt}\setlength{\parskip}{0.35em}},
  }
}

\newcommand{\bilingualstart}{}
\newcommand{\bilingualstop}{}
\ifbilingualcolumns
  \renewcommand{\bilingualstart}{\small\setlength{\columnsep}{1.25cm}\setlength{\columnseprule}{0.12pt}\columnratio{0.5,0.5}\multicolumnfootnotes\flushbottom\sloppy\interlinepenalty=80\clubpenalty=800\widowpenalty=800\displaywidowpenalty=800\begin{paracol}{2}\global\intranslationfalse\global\firstbilingualsectiontrue{\footnotesize\color{sourceink}Deutsch\par}\switchcolumn{\footnotesize\songfont\color{sourceink}中文譯文\par}\switchcolumn*}
  \renewcommand{\bilingualstop}{\ifintranslation\switchcolumn*\global\intranslationfalse\fi\end{paracol}\clearpage\raggedbottom}
\fi
\newif\ifrandnumbers
\randnumbersfalse
\newcounter{randnumber}
\newcounter{randnumberlocal}
\newcounter{lastsourcerandstart}
\newif\iflastsourcehasrandnumbers
\lastsourcehasrandnumbersfalse
\newif\ifinlawbox
\inlawboxfalse
\newif\iflawboxhaspendingrn
\lawboxhaspendingrnfalse
\newcommand{\lawboxpendingrn}{}
\newcommand{\startrandnumbers}{\global\randnumberstrue\setcounter{randnumber}{1}\global\lastsourcehasrandnumbersfalse}
\newcommand{\stoprandnumbers}{\global\randnumbersfalse\global\lastsourcehasrandnumbersfalse}
\newlength{\randnumberwidth}
\setlength{\randnumberwidth}{0.95cm}
\newlength{\randnumberrightoffset}
\setlength{\randnumberrightoffset}{0.85cm}
\newlength{\randnumberlawboxrightoffset}
\setlength{\randnumberlawboxrightoffset}{1.40cm}
\newcommand{\randnumberbox}[1]{%
  \leavevmode
  \makebox[0pt][l]{%
    \smash{%
      \tikz[remember picture,overlay,baseline=(rnmark.base)]{%
        \coordinate (rnmark) at (0,0);%
        \ifinlawbox
          \path let \p1=(current page.east), \p2=(rnmark.base) in
            node[
              anchor=base east,
              inner sep=0pt,
              outer sep=0pt,
              text width=\randnumberwidth,
              align=right,
              text=pendingink,
              font=\normalfont\mdseries\scriptsize\ttfamily
            ] at (\x1-\randnumberlawboxrightoffset,\y2) {{\normalfont\mdseries\scriptsize\ttfamily #1}};%
        \else
          \path let \p1=(current page.east), \p2=(rnmark.base) in
            node[
              anchor=base east,
              inner sep=0pt,
              outer sep=0pt,
              text width=\randnumberwidth,
              align=right,
              text=pendingink,
              font=\normalfont\mdseries\scriptsize\ttfamily
            ] at (\x1-\randnumberrightoffset,\y2) {{\normalfont\mdseries\scriptsize\ttfamily #1}};%
        \fi
      }%
    }%
  }%
  \ignorespaces
}
\newcommand{\lawboxstorerandnumber}[1]{%
  \gdef\lawboxpendingrn{#1}%
  \global\lawboxhaspendingrntrue
  \ignorespaces
}
\newcommand{\lawboxuserandnumber}{%
  \iflawboxhaspendingrn
    \randnumberbox{\lawboxpendingrn}%
    \global\lawboxhaspendingrnfalse
  \fi
}
\newcommand{\sourcern}[1]{%
  \ifrandnumbers
    \ifshoworiginal
      \ifshowtranslation\else
        \ifinlawbox\lawboxstorerandnumber{#1}\else\randnumberbox{#1}\fi
      \fi
    \fi
  \fi
}
\newcommand{\transrn}[1]{%
  \ifrandnumbers
    \ifshowtranslation
      \ifinlawbox\lawboxstorerandnumber{#1}\else\randnumberbox{#1}\fi
    \fi
  \fi
}
% ===== 課堂模式：德文原文縮寫註解 =====
% 日常切換只改各判決檔的一行：
%   \classroommodeon        開啟課堂版；註解掉這行即回到一般版。
% 模板預設：
%   1. 課堂模式關閉。
%   2. 課堂模式開啟時，縮寫註解包含「德文全稱＋中文譯名」。
%   3. 課堂模式開啟時，GG/條文編者註仍會自動顯示。
% 進階微調才需要在單案檔覆寫：
%   \classroomabbrzhoff     縮寫註解僅顯示德文全稱。
%   \showeditornotestrue    一般版也顯示編者註。
% 註解策略：同一實際 PDF 頁面中，同一縮寫只在第一次出現處加註。
% 這裡用 zref-abspage 的絕對頁序，不用 \thepage，避免文件內頁碼重置時誤判。
\newif\ifclassroommode
\newif\ifclassroomabbrzh
\classroommodefalse
\classroomabbrzhtrue
\newcommand{\classroommodeon}{\classroommodetrue}
\newcommand{\classroommodeoff}{\classroommodefalse}
\newcommand{\classroomabbrzhon}{\classroomabbrzhtrue}
\newcommand{\classroomabbrzhoff}{\classroomabbrzhfalse}
\newcommand{\abbrnotelabel}{\emph{Abk.}: }
\newcommand{\abbrnote}[3]{%
  \ifclassroommode
    \ifshoworiginal
      \ifintranslation\else
        \footnote{\normalfont\footnotesize\abbrnotelabel #2\ifclassroomabbrzh；{\songfont #3}\fi}%
      \fi
    \fi
  \fi
}
\newcommand{\abbrnoteonce}[4]{%
  #2%
  \ifclassroommode
    \ifshoworiginal
      \ifintranslation\else
        \begingroup
          \edef\abbrcurrentabspage{\number\numexpr\value{abspage}+1\relax}%
          \ifcsname abbrnoteseen@#1@\abbrcurrentabspage\endcsname\else
            \global\expandafter\let\csname abbrnoteseen@#1@\abbrcurrentabspage\endcsname\empty
            \abbrnote{#1}{#3}{#4}%
          \fi
        \endgroup
      \fi
    \fi
  \fi
}
\newcommand{\abbrAbs}{\abbrnoteonce{abs}{Abs.}{Absatz}{項}}
\newcommand{\abbrAAO}{\abbrnoteonce{aao}{a.a.O.}{am angegebenen Ort}{前引處}}
\newcommand{\abbrAE}{\abbrnoteonce{ae}{AE}{Alternativ-Entwurf}{替代草案}}
\newcommand{\abbrArt}{\abbrnoteonce{art}{Art.}{Artikel}{條}}
\newcommand{\abbrAufl}{\abbrnoteonce{aufl}{Aufl.}{Auflage}{版}}
\newcommand{\abbrBd}{\abbrnoteonce{bd}{Bd.}{Band}{卷}}
\newcommand{\abbrBGBl}{\abbrnoteonce{bgbl}{BGBl.}{Bundesgesetzblatt}{聯邦法律公報}}
\newcommand{\abbrBGHSt}{\abbrnoteonce{bghst}{BGHSt}{Entscheidungen des Bundesgerichtshofes in Strafsachen}{聯邦普通法院刑事判決彙編}}
\newcommand{\abbrBRDrucks}{\abbrnoteonce{brdrucks}{BRDrucks.}{Bundesrats-Drucksache}{聯邦參議院文件}}
\newcommand{\abbrBTDrucks}{\abbrnoteonce{btdrucks}{BTDrucks.}{Bundestags-Drucksache}{聯邦議院文件}}
\newcommand{\abbrBundesgesetzbl}{\abbrnoteonce{bundesgesetzbl}{Bundesgesetzbl.}{Bundesgesetzblatt}{聯邦法律公報}}
\newcommand{\abbrBvF}{\abbrnoteonce{bvf}{BvF}{Bundesverfassungsgerichtsverfahren}{聯邦憲法法院程序案號}}
\newcommand{\abbrBvL}{\abbrnoteonce{bvl}{BvL}{Bundesverfassungsgerichtsverfahren}{聯邦憲法法院程序案號}}
\newcommand{\abbrBvR}{\abbrnoteonce{bvr}{BvR}{Bundesverfassungsgerichtsverfahren}{聯邦憲法法院程序案號}}
\newcommand{\abbrBVerfG}{\abbrnoteonce{bverfg}{BVerfG}{Bundesverfassungsgericht}{聯邦憲法法院}}
\newcommand{\abbrBVerfGE}{\abbrnoteonce{bverfge}{BVerfGE}{Entscheidungen des Bundesverfassungsgerichts}{聯邦憲法法院判決集}}
\newcommand{\abbrBVerfGG}{\abbrnoteonce{bverfgg}{BVerfGG}{Bundesverfassungsgerichtsgesetz}{聯邦憲法法院法}}
\newcommand{\abbrDDR}{\abbrnoteonce{ddr}{DDR}{Deutsche Demokratische Republik}{德意志民主共和國；東德}}
\newcommand{\abbrDJT}{\abbrnoteonce{djt}{DJT}{Deutscher Juristentag}{德國法學家大會}}
\newcommand{\abbrDr}{\abbrnoteonce{dr}{Dr.}{Doktor}{Dr.}}
\newcommand{\abbrEuGRZ}{\abbrnoteonce{eugrz}{EuGRZ}{Europäische Grundrechte-Zeitschrift}{歐洲基本權期刊}}
\newcommand{\abbrEV}{\abbrnoteonce{ev}{EV}{Einigungsvertrag}{統一條約}}
\newcommand{\abbrF}{\abbrnoteonce{f}{f.}{folgende}{以下一頁；下一段}}
\newcommand{\abbrFF}{\abbrnoteonce{ff}{ff.}{fortfolgende}{以下數頁；以下各段}}
\newcommand{\abbrGBlDDR}{\abbrnoteonce{gblddr}{GBl. DDR}{Gesetzblatt der Deutschen Demokratischen Republik}{德意志民主共和國法律公報}}
\newcommand{\abbrGG}{\abbrnoteonce{gg}{GG}{Grundgesetz}{基本法}}
\newcommand{\abbrGVBl}{\abbrnoteonce{gvbl}{GVBl.}{Gesetz- und Verordnungsblatt}{法律及命令公報}}
\newcommand{\abbrJR}{\abbrnoteonce{jr}{JR}{Juristische Rundschau}{法學評論}}
\newcommand{\abbrJZ}{\abbrnoteonce{jz}{JZ}{JuristenZeitung}{法學家報}}
\newcommand{\abbrIVm}{\abbrnoteonce{ivm}{i.V.m.}{in Verbindung mit}{結合；連同}}
\newcommand{\abbrMdB}{\abbrnoteonce{mdb}{MdB}{Mitglied des Bundestages}{聯邦議院議員}}
\newcommand{\abbrMwN}{\abbrnoteonce{mwn}{m.w.N.}{mit weiteren Nachweisen}{附進一步文獻／判例出處}}
\newcommand{\abbrNF}{\abbrnoteonce{nf}{n.F.}{neue Fassung}{新版本}}
\newcommand{\abbrAF}{\abbrnoteonce{af}{a.F.}{alte Fassung}{舊版本}}
\newcommand{\abbrNr}{\abbrnoteonce{nr}{Nr.}{Nummer}{款、目或號，依語境}}
\newcommand{\abbrProf}{\abbrnoteonce{prof}{Prof.}{Professor}{教授}}
\newcommand{\abbrRdnr}{\abbrnoteonce{rdnr}{Rdnr.}{Randnummer}{邊碼；段碼}}
\newcommand{\abbrRdnrn}{\abbrnoteonce{rdnrn}{Rdnrn.}{Randnummern}{邊碼；段碼}}
\newcommand{\abbrRGBl}{\abbrnoteonce{rgbl}{RGBl.}{Reichsgesetzblatt}{帝國法律公報}}
\newcommand{\abbrRGSt}{\abbrnoteonce{rgst}{RGSt}{Entscheidungen des Reichsgerichts in Strafsachen}{帝國法院刑事判決彙編}}
\newcommand{\abbrRspr}{\abbrnoteonce{rspr}{Rspr.}{Rechtsprechung}{判例；司法實務}}
\newcommand{\abbrRVO}{\abbrnoteonce{rvo}{RVO}{Reichsversicherungsordnung}{帝國保險法}}
\newcommand{\abbrS}{\abbrnoteonce{s}{S.}{Seite}{頁}}
\newcommand{\abbrSFHG}{\abbrnoteonce{sfhg}{SFHG}{Schwangeren- und Familienhilfegesetz}{孕婦及家庭扶助法}}
\newcommand{\abbrSGBV}{\abbrnoteonce{sgbv}{SGB V}{Fünftes Buch Sozialgesetzbuch}{社會法典第五編}}
\newcommand{\abbrStAG}{\abbrnoteonce{staeg}{StÄG}{Strafrechtsänderungsgesetz}{刑法修正法}}
\newcommand{\abbrStenBer}{\abbrnoteonce{stenber}{StenBer.}{Stenographischer Bericht}{速記紀錄}}
\newcommand{\abbrStGB}{\abbrnoteonce{stgb}{StGB}{Strafgesetzbuch}{刑法}}
\newcommand{\abbrStREG}{\abbrnoteonce{streg}{StREG}{Strafrechtsreform-Ergänzungsgesetz}{刑法改革補充法}}
\newcommand{\abbrStrRG}{\abbrnoteonce{strrg}{StrRG}{Strafrechtsreformgesetz}{刑法改革法}}
\newcommand{\abbrUS}{\abbrnoteonce{us}{U.S.}{United States Reports}{美國最高法院判決彙編}}
\newcommand{\abbrVgl}{\abbrnoteonce{vgl}{vgl.}{vergleiche}{參見}}
\newcommand{\abbrVVDStRL}{\abbrnoteonce{vvdstrl}{VVDStRL}{Veröffentlichungen der Vereinigung der Deutschen Staatsrechtslehrer}{德國公法教師協會論文集}}
\newcommand{\abbrWP}{\abbrnoteonce{wp}{WP.}{Wahlperiode}{屆期}}
\newcommand{\abbrWp}{\abbrnoteonce{wp}{Wp.}{Wahlperiode}{屆期}}
\newcommand{\abbrZStrW}{\abbrnoteonce{zstrw}{ZStrW}{Zeitschrift für die gesamte Strafrechtswissenschaft}{整體刑法學期刊}}
\newcommand{\recordsourcerandstart}{%
  \ifrandnumbers
    \setcounter{lastsourcerandstart}{\value{randnumber}}%
    \global\lastsourcehasrandnumberstrue
  \else
    \global\lastsourcehasrandnumbersfalse
  \fi
}
\newlength{\biparagraphskip}
\setlength{\biparagraphskip}{0.52em}
\newlength{\lawboxblockbeforeskip}
\setlength{\lawboxblockbeforeskip}{0.72em}
\newlength{\lawboxblockafterskip}
\setlength{\lawboxblockafterskip}{0.92em}
\newif\iflawboxpartendedblock
\lawboxpartendedblockfalse
\newcommand{\sourceparstyle}{\footnotesize\color{sourceink}\setstretch{1.12}\RaggedRight\emergencystretch=3em\setlength{\parskip}{0.42em}\obeylines\selectfont}
\newcommand{\transparstyle}{\songfont\color{caseink}\setstretch{1.12}\setlength{\parindent}{0pt}\setlength{\parskip}{0.42em}}
\newcommand{\bilingualpairskip}{%
  \par\addvspace{\biparagraphskip}%
  \switchcolumn
  \par\addvspace{\biparagraphskip}%
  \switchcolumn*%
  \global\intranslationfalse
}
\newcommand{\bilingualboxpairskip}{%
  \par
  \switchcolumn
  \par
  \switchcolumn*%
  \global\intranslationfalse
}
\newcommand{\bilinguallawboxblockskip}{%
  \switchcolumn*%
  \par\addvspace{\lawboxblockafterskip}%
  \switchcolumn
  \par\addvspace{\lawboxblockafterskip}%
  \switchcolumn*%
  \global\intranslationfalse
}
\newcommand{\bikeepwithnext}[1][4]{%
  \ifbilingualcolumns
    \syncleft
    \Needspace{#1\baselineskip}%
    \switchcolumn
    \Needspace{#1\baselineskip}%
    \switchcolumn*%
    \global\intranslationfalse
  \else
    \Needspace{#1\baselineskip}%
  \fi
}
\NewEnviron{sourcepara}[1][]{
  \recordsourcerandstart
  \ifshoworiginal
    \ifbilingualcolumns
      \ifintranslation\switchcolumn*\global\intranslationfalse\fi
    \else
      \Needspace{5\baselineskip}
    \fi
    \ifbilingualcolumns\else\par\addvspace{0.35em}\fi
    {\sourceparstyle
    \noindent\BODY\par}
    \ifbilingualcolumns\else\addvspace{0.25em}\fi
    \ifbilingualcolumns
      \switchcolumn
      \global\intranslationtrue
    \fi
  \else
    \ifrandnumbers
      \begingroup
        \setbox0=\vbox{\hsize=\linewidth\sourceparstyle\noindent\BODY\par}%
      \endgroup
    \fi
  \fi
}
\NewEnviron{transpara}{
  \ifshowtranslation
    \setcounter{randnumberlocal}{\value{lastsourcerandstart}}%
    \ifbilingualcolumns
      \ifintranslation\else\switchcolumn\global\intranslationtrue\fi
      {\transparstyle\setlength{\leftskip}{0.55em}\setlength{\rightskip}{0.15em plus 1em}\addtolength{\linewidth}{-0.55em}\BODY\par}
      \switchcolumn*
      \bilingualpairskip
    \else
      {\transparstyle\BODY\par}
    \fi
  \else
    \ifbilingualcolumns
      \ifintranslation\switchcolumn*\global\intranslationfalse\fi
    \fi
  \fi
}
\NewEnviron{sourceboxpara}[1][]{
  \global\lawboxpartendedblockfalse
  \recordsourcerandstart
  \ifshoworiginal
    \ifbilingualcolumns
      \ifintranslation\switchcolumn*\global\intranslationfalse\fi
    \else
      \Needspace{3\baselineskip}
    \fi
    {\sourceparstyle
    \BODY\par}
    \ifbilingualcolumns\else
      \iflawboxpartendedblock\addvspace{\lawboxblockafterskip}\fi
    \fi
    \ifbilingualcolumns
      \switchcolumn
      \global\intranslationtrue
    \fi
  \fi
}
\NewEnviron{transboxpara}{
  \global\lawboxpartendedblockfalse
  \ifshowtranslation
    \setcounter{randnumberlocal}{\value{lastsourcerandstart}}%
    \ifbilingualcolumns
      \ifintranslation\else\switchcolumn\global\intranslationtrue\fi
      {\transparstyle\BODY\par}
      \iflawboxpartendedblock
        \bilinguallawboxblockskip
      \else
        \switchcolumn*
        \bilingualboxpairskip
      \fi
    \else
      {\transparstyle\BODY\par}
      \iflawboxpartendedblock\addvspace{\lawboxblockafterskip}\fi
    \fi
  \else
    \ifbilingualcolumns
      \ifintranslation\switchcolumn*\global\intranslationfalse\fi
    \fi
  \fi
}
\newcommand{\leitsatznum}[1]{\noindent\hangindent=3.0em\hangafter=1\makebox[2.25em][r]{\bfseries #1.}\hspace{0.55em}\ignorespaces}
\newcommand{\syncleft}{\ifbilingualcolumns\ifintranslation\switchcolumn*\global\intranslationfalse\fi\fi}
\pretocmd{\section}{\syncleft\clearpage}{}{}
\pretocmd{\subsection}{\syncleft}{}{}
\pretocmd{\subsubsection}{\syncleft}{}{}
\newcommand{\biheadingfont}{\songfont}
\newcommand{\bitocentry}[2]{%
  \ifshoworiginal
    \ifshowtranslation
      \ifstrequal{#1}{#2}{#1}{#1\space #2}%
    \else
      #1%
    \fi
  \else
    #2%
  \fi
}
\pdfstringdefDisableCommands{%
  \def\bitocentry#1#2{#1}%
}
\newcommand{\biheading}[7]{%
  \syncleft#1\bikeepwithnext[#3]\phantomsection\addcontentsline{toc}{#2}{\bitocentry{#6}{#7}}%
  \ifbilingualcolumns
    {#4 #6\par}%
    \switchcolumn
    {#4\biheadingfont #7\par}%
    \switchcolumn*%
    \global\intranslationfalse
    \par\addvspace{#5}%
    \switchcolumn
    \par\addvspace{#5}%
    \switchcolumn*%
  \else
    \ifshoworiginal{#4 #6\par}\fi
    \ifshowtranslation{#4\biheadingfont #7\par}\fi
    \addvspace{#5}%
  \fi
}
\newcommand{\termsectionheading}[1]{%
  \par\Needspace{9\baselineskip}%
  \vspace{0.65em}%
  {\songfont\bfseries #1\par}%
  \vspace{0.2em}%
}
% 章節換頁：
%   雙語 paracol 欄內不要用 \clearpage；它會在同步兩欄時吐出只有欄線/頁碼的空頁。
%   \newpage 足以讓下一個 section 起新頁，又不會強制清空所有待排材料。
%   單語模式不在 paracol 內，仍可用 \clearpage。
\newcommand{\bisectionpagebreak}{\ifbilingualcolumns\newpage\else\clearpage\fi}
\newcommand{\bisection}[2]{%
  \biheading{\iffirstbilingualsection\global\firstbilingualsectionfalse\else\bisectionpagebreak\fi}{section}{6}{\centering\Large\bfseries}{0.8em}{#1}{#2}%
}
\newcommand{\bisubsection}[2]{%
  \biheading{}{subsection}{5}{\large\bfseries}{0.55em}{#1}{#2}%
}
\newcommand{\bisubsubsection}[2]{%
  \biheading{}{subsubsection}{4}{\normalsize\bfseries}{0.45em}{#1}{#2}%
}
\newcommand{\bicompositeheading}[4]{%
  \biheading{}{subsection}{5}{\large\bfseries}{0.16em}{#1}{#2}%
  \biheading{}{subsubsection}{3}{\normalsize\bfseries}{0.45em}{#3}{#4}%
}
\newif\iflawboxfirstitem
\lawboxfirstitemfalse
\newcommand{\lawboxprespace}[1]{%
  \par
  \iflawboxfirstitem
    \global\lawboxfirstitemfalse
  \else
    \addvspace{#1}%
  \fi
}
\newtcolorbox{lawboxinner}{
  caseboxbase,
  colframe=sourceline!85!black,
  colback=white,
  left=1.05em,
  right=1.05em,
  top=0.62em,
  bottom=0.62em,
  before upper={\global\inlawboxtrue\global\lawboxfirstitemtrue\global\lawboxhaspendingrnfalse\ifintranslation\lawboxtransfont\else\lawboxsourcefont\fi\RaggedRight\setlength{\parindent}{0pt}\setlength{\parskip}{0pt}\ifintranslation\setstretch{\lawboxtransstretch}\else\setstretch{\lawboxsourcestretch}\fi},
  after upper={\global\lawboxhaspendingrnfalse\global\inlawboxfalse}
}
\tcbset{
  lawboxpartbase/.style={
    enhanced,
    breakable=false,
    width=0.985\linewidth,
    colback=white,
    frame hidden,
    boxrule=0pt,
    arc=0pt,
    left=1.05em,
    right=1.05em,
    boxsep=1.5mm,
    before skip=0pt,
    after skip=0pt,
    before upper={\global\inlawboxtrue\global\lawboxfirstitemtrue\global\lawboxhaspendingrnfalse\ifintranslation\lawboxtransfont\else\lawboxsourcefont\fi\RaggedRight\setlength{\parindent}{0pt}\setlength{\parskip}{0pt}\ifintranslation\setstretch{\lawboxtransstretch}\else\setstretch{\lawboxsourcestretch}\fi},
    after upper={\global\lawboxhaspendingrnfalse\global\inlawboxfalse},
    borderline west={0.28pt}{0pt}{sourceline!85!black},
    borderline east={0.28pt}{0pt}{sourceline!85!black},
  },
  lawboxpartfirst/.style={
    lawboxpartbase,
    top=0.62em,
    bottom=0.04em,
    before skip=\lawboxblockbeforeskip,
    borderline north={0.28pt}{0pt}{sourceline!85!black},
  },
  lawboxpartmiddle/.style={
    lawboxpartbase,
    top=0.04em,
    bottom=0.04em,
  },
  lawboxpartlast/.style={
    lawboxpartbase,
    top=0.04em,
    bottom=0.62em,
    after skip=0pt,
    borderline south={0.28pt}{0pt}{sourceline!85!black},
  }
}
\newtcolorbox{lawboxpartinner}[1]{#1}
\newtcolorbox{termboxinner}{
  caseboxbase,
  colframe=noteblue!70!black,
  colback=casepaper,
  before upper={\songfont\setlength{\parindent}{0pt}\setlength{\parskip}{0.35em}}
}
\newtcolorbox{deboxinner}{
  caseboxbase,
  colframe=notegray!72!black,
  colback=white,
  left=1.05em,
  right=1.05em,
  top=0.62em,
  bottom=0.62em,
  before upper={\global\inlawboxtrue\global\lawboxfirstitemtrue\global\lawboxhaspendingrnfalse\small\setlength{\parindent}{0pt}\setlength{\parskip}{0.35em}},
  after upper={\global\lawboxhaspendingrnfalse\global\inlawboxfalse}
}
\NewEnviron{lawbox}{
  \begin{lawboxinner}\BODY\end{lawboxinner}
}
\NewEnviron{lawboxpart}[2]{
  \begin{lawboxpartinner}{lawboxpart#1,equal height group=#2}\BODY\end{lawboxpartinner}%
  \ifstrequal{#1}{last}{\global\lawboxpartendedblocktrue}{}%
}
\NewEnviron{termbox}{
  \par\addvspace{0.55em}
  {\color{noteblue!70!black}\rule{\linewidth}{0.28pt}}\par
  \begingroup
    \songfont\small\color{caseink}
    \setlength{\parindent}{0pt}
    \setlength{\parskip}{0.24em}
    \BODY
    \par
  \endgroup
  {\color{noteblue!70!black}\rule{\linewidth}{0.28pt}}\par
  \addvspace{0.55em}
}
\NewEnviron{debox}{
  \begin{deboxinner}\BODY\end{deboxinner}
}
\providecommand{\paragraphmark}[1]{}
\newcommand{\lawboxtitleafter}{\addvspace{0.06em}}
\newcommand{\lawtitle}[1]{\lawboxprespace{0.22em}{\lawboxtransfont\bfseries\lawboxuserandnumber #1}\par\lawboxtitleafter}
\newcommand{\absno}[2]{\lawboxprespace{0.04em}\noindent\lawboxuserandnumber\hangindent=2.6em\hangafter=1\makebox[2.6em][l]{(#1)}#2\par}
\newcommand{\nrno}[2]{\lawboxprespace{0pt}\noindent\lawboxuserandnumber\hspace*{2.6em}\hangindent=4.4em\hangafter=1\makebox[1.8em][l]{#1.}#2\par}
\newcommand{\letterno}[2]{\lawboxprespace{0pt}\noindent\lawboxuserandnumber\hspace*{4.4em}\hangindent=6.0em\hangafter=1\makebox[1.6em][l]{#1)}#2\par}
\newcommand{\sourcelawtitle}[1]{\lawboxprespace{0.22em}{\lawboxsourcefont\bfseries\lawboxuserandnumber #1}\par\lawboxtitleafter}
\newcommand{\sourcelawsubtitle}[1]{\lawboxprespace{0.04em}{\lawboxsourcefont\bfseries\lawboxuserandnumber #1}\par\lawboxtitleafter}
\newcommand{\sourceabsno}[2]{\lawboxprespace{0.04em}\noindent\lawboxuserandnumber\hangindent=2.45em\hangafter=1\makebox[2.45em][l]{(#1)}#2\par}
\newcommand{\sourcenrno}[2]{\lawboxprespace{0pt}\noindent\lawboxuserandnumber\hspace*{2.45em}\hangindent=4.25em\hangafter=1\makebox[1.8em][l]{#1.}#2\par}
\newcommand{\sourceletterno}[2]{\lawboxprespace{0pt}\noindent\lawboxuserandnumber\hspace*{4.25em}\hangindent=5.85em\hangafter=1\makebox[1.6em][l]{#1)}#2\par}
\newcommand{\sourcequotepara}[1]{\lawboxprespace{0.04em}\noindent\lawboxuserandnumber\hangindent=1.2em\hangafter=1 #1\par}
\newcommand{\sourcelawline}[1]{\lawboxprespace{0.04em}\noindent\lawboxuserandnumber #1\par}
\newcommand{\lawline}[1]{\lawboxprespace{0.04em}\noindent\lawboxuserandnumber #1\par}
\newcommand{\pendingtranslation}{\par{\songfont\footnotesize\color{pendingink}（中文譯文待補。）}\par}
\newcommand{\tocpart}[1]{\par\addvspace{0.52em}{\footnotesize\bfseries\color{pendingink}#1\par}\addvspace{0.16em}}
\newcommand{\tocdashline}{\par\addvspace{0.48em}\noindent{\color{notegray}\xleaders\hbox to 0.82em{\hfil\rule{0.42em}{0.28pt}\hfil}\hskip\linewidth}\par\addvspace{0.38em}}
% \tocpanel{font-cmd}{content}：將目錄置中於所在欄寬（雙語欄適用）。
% A3 橫式使用 0.58\linewidth，A4 橫式使用 0.84\linewidth，
% 使兩種紙張下目錄欄內容實際寬度相近（均約 100–105 mm）。
\newcommand{\tocpanel}[2]{%
  \makebox[\linewidth][c]{%
    \ifx\bilingualpaper\bilingualpaperaThree
      \begin{minipage}{0.58\linewidth}%
    \else
      \begin{minipage}{0.84\linewidth}%
    \fi
    \toccolstyle
    #1#2%
    \end{minipage}%
  }%
}
% ===== 首頁簡目錄的兩層 description list 縮排 =====
% enumitem 的 description list 可把每一列拆成：
%   [label]  labelsep  body text
% 這裡的「起點」都以所在 minipage/欄位的左緣為 0。
%
% 核心規則：
%   labelindent = label 欄位從左緣開始的位置。
%   labelwidth  = label 欄位本身寬度；標籤太長（如 Abw. Meinung）就加大它。
%   labelsep    = label 欄位右緣到正文的距離。
%   leftmargin  = 正文 body text 的起點。判斷層級視覺時主要看這個。
%
% 所以：
%   想讓「Begründung / 判決理由」整列正文往右或往左：改 tocitems 的 leftmargin。
%   想讓「A. Verfahren und Gegenstand / A. 程序與審查標的」正文往右或往左：
%     改 tocsubitems 的 leftmargin。
%   想讓 A./B./C. 這些標籤本身往右或往左，但正文不動：改 tocsubitems 的 labelindent。
%   想讓 A./B./C. 與正文間距變大或變小：改 tocsubitems 的 labelsep。
%
% 層級原則：
%   tocsubitems.leftmargin 必須大於 tocitems.leftmargin，
%   否則子層級正文會看起來比「Gründe / 理由」還靠左。
\newlist{tocitems}{description}{1}
\setlist[tocitems]{%
  labelindent=0pt,    % 第一層 label 欄位起點；0pt 表示貼齊目錄欄左緣。
  leftmargin=2.54cm,  % 第一層正文起點；例如 Begründung / 判決理由 的左緣。
  labelwidth=2.16cm,  % 第一層 label 欄位寬度；需容納 Leitsätze、Abw. Meinung、不同意見等。
  labelsep=0.32cm,    % 第一層 label 與正文的水平間距；只改間距，不改正文起點。
  itemsep=0.28em,     % 第一層各項之間的垂直距離。
  topsep=0.20em,      % 第一層 list 上下與前後內容的垂直距離。
  parsep=0pt,         % 同一 item 內段落之間的距離；目錄項通常不需要。
  partopsep=0pt,      % list 若緊接段落開始時的額外距離；關閉以免目錄高度飄動。
  align=left,         % label 欄內左對齊；長短標籤不靠右貼齊。
  font=\normalfont\bfseries % label 的字型；正文不受這行影響。
}
\newlist{tocsubitems}{description}{1}
\setlist[tocsubitems]{%
  labelindent=1cm,    % 第二層 label 起點；只控制 A./B./C. 本身的位置。
  leftmargin=3.00cm,  % 第二層正文起點；必須大於 2.54cm，才會比 Begründung / 判決理由更右。
  labelwidth=0.5cm,   % 第二層 label 欄寬；A./B./C. 很短，可小於第一層。
  labelsep=1.5cm,    % 第二層 label 與正文間距；A. 與 Verfahren / 程序 的距離看這裡。
  itemsep=0.12em,     % 第二層各項較密，避免 A.–E. 佔太高。
  topsep=0.04em,      % 第二層 list 與 Gründe / 理由之間的垂直距離。
  parsep=0pt,         % 同一第二層 item 內段落距離；目錄項通常不需要。
  partopsep=0pt,      % 關閉額外段落起始距離，避免版面不穩。
  align=left,         % A./B./C. label 欄內左對齊。
  font=\normalfont\bfseries, % A./B./C. label 加粗；正文不受這行影響。
  before={}           % 不在第二層 list 前自動插入額外內容。
}
\newlist{termitems}{description}{1}
\ifbilingualcolumns
  \setlist[termitems]{%
    leftmargin=5.2cm,
    labelwidth=4.8cm,
    labelsep=0.35cm,
    itemsep=0.16em,
    topsep=0.2em,
    parsep=0pt,
    partopsep=0pt,
    font=\normalfont\bfseries,
    style=nextline
  }
\else
  \setlist[termitems]{%
    leftmargin=0pt,
    labelwidth=0pt,
    labelsep=0pt,
    itemsep=0.24em,
    topsep=0.2em,
    parsep=0pt,
    partopsep=0pt,
    font=\normalfont\bfseries,
    style=nextline
  }
\fi
\newlist{abbritems}{description}{1}
\setlist[abbritems]{%
  leftmargin=2.38cm,
  labelwidth=2.02cm,
  labelsep=0.28cm,
  itemsep=0.16em,
  topsep=0.2em,
  parsep=0pt,
  partopsep=0pt,
  align=left,
  font=\normalfont\bfseries
}
\ifbilingualcolumns
  \newenvironment{termcolumns}{\begin{multicols}{2}}{\end{multicols}}
\else
  \newenvironment{termcolumns}{}{}
\fi

% ===== 編者註腳開關 =====
% 一般版預設不顯示編者註，避免正文腳註過密。
% 若開啟 \classroommodeon，GG/條文編者註仍會自動顯示，無需另開本開關。
% 若一般版也要顯示編者註，可在單案檔 \input 後加 \showeditornotestrue。
\newif\ifshoweditornotes
\showeditornotesfalse

% ===== 目錄排版輔助 =====
\newcommand{\toccolstyle}{\raggedright\arraybackslash}
\newcommand{\tocsingleindent}{\leftskip=1.45cm\rightskip=1.2cm}

% ===== 行內引文環境（德文原文與中文譯文各一，帶左側灰色邊線）=====
\newcommand{\quoteparbreak}{\par\addvspace{0.48em}}
\newcommand{\sourcequoteinner}[1]{%
  \begin{tcolorbox}[
    enhanced, breakable, width=\dimexpr\linewidth-0.8em\relax, frame hidden,
    colback=white, coltext=caseink, boxrule=0pt,
    borderline west={0.55pt}{0.88em}{notegray},
    left=1.78em, right=0.72em, top=0.18em, bottom=0.18em,
    boxsep=0pt, before skip=0.24em, after skip=0.24em,
    before upper={\normalfont\footnotesize\color{caseink}\setlength{\parindent}{0pt}\setlength{\parskip}{0.34em}\raggedright\setlength{\rightskip}{0pt plus 1em}}
  ]%
  #1\par
  \end{tcolorbox}%
}
\newcommand{\transquoteinner}[1]{%
  \begin{tcolorbox}[
    enhanced, breakable, width=\dimexpr\linewidth-0.8em\relax, frame hidden,
    colback=white, coltext=caseink, boxrule=0pt,
    borderline west={0.55pt}{0.88em}{notegray},
    left=1.78em, right=0.72em, top=0.18em, bottom=0.18em,
    boxsep=0pt, before skip=0.24em, after skip=0.24em,
    before upper={\songfont\small\color{caseink}\setlength{\parindent}{0pt}\setlength{\parskip}{0.34em}\raggedright\setlength{\rightskip}{0pt plus 1em}}
  ]%
  #1\par
  \end{tcolorbox}%
}

% ===== 大綱（Gliederung）Rn 輔助命令 =====
\newcommand{\outrn}[2]{\enspace{\normalfont\footnotesize\color{pendingink}(Rn\,#1–#2)}}
\newcommand{\outrnsingle}[1]{\enspace{\normalfont\footnotesize\color{pendingink}(Rn\,#1)}}
\newcommand{\outrnzh}[2]{\enspace{\normalfont\footnotesize\color{pendingink}（第\,#1–#2\,段）}}
\newcommand{\outrnzhs}[1]{\enspace{\normalfont\footnotesize\color{pendingink}（第\,#1\,段）}}
% ===== 大綱雙語對照表格輔助命令（三欄：段碼 │ 德文 │ 中文）=====
% 欄 1：段碼（由欄定義自動格式化：小字、等寬、右對齊、pendingink）
% 欄 2：德文標籤＋正文
% 欄 3：中文標籤＋正文
%
% 無段碼的列（\omainrow、\osecrow），欄 1 以 & 空置。
% 有段碼的列（\osecrowrn、\osubrow、\ointrow），#1 為預格式化字串
%   例：\osubrow{2–49}{I.}{...}{I.}{...}
%       \ointrow{107}{...}{...}
%
% \opartrow：段組標題（橫跨三欄）
\newcommand{\opartrow}[2]{%
  \noalign{\vspace{0.30em}}%
  \multicolumn{3}{@{}l@{}}{%
    \footnotesize\bfseries\color{pendingink}#1\hfill\songfont#2%
  }\\[0.06em]%
}
% \odashrow：虛線分隔（橫跨三欄）
\newcommand{\odashrow}{%
  \noalign{\vspace{0.28em}}%
  \multicolumn{3}{@{}l@{}}{\color{notegray}%
    \xleaders\hbox to 0.82em{\hfil\rule{0.42em}{0.28pt}\hfil}\hskip 0.88\linewidth}%
  \\[0.24em]%
}
% \odissentpart：不同意見書區段標頭。比一般 \opartrow 更像附錄性轉場。
\newcommand{\odissentpart}[2]{%
  \noalign{\vspace{0.42em}}%
  \multicolumn{3}{@{}p{0.88\textwidth}@{}}{%
    \color{notegray}\rule{\linewidth}{0.18pt}%
  }\\[-0.18em]%
  \multicolumn{3}{@{}p{0.88\textwidth}@{}}{%
    \makebox[\linewidth][s]{%
      \footnotesize\bfseries\color{pendingink}#1%
      \hfill{\songfont #2}%
    }%
  }\\[-0.02em]%
}
% \odissentopinion：不同意見書的署名與一行摘要，右欄保留段碼範圍。
\newcommand{\odissentopinion}[5]{%
  \footnotesize\hangindent=0.52cm\hangafter=1%
    {\bfseries\color{caseink}#2}\enspace{\color{notegray}\textperiodcentered}\enspace\color{caseink}#3%
  &\footnotesize\hangindent=0.52cm\hangafter=1%
    {\songfont\bfseries\color{caseink}#4}\enspace{\color{notegray}\textperiodcentered}\enspace\songfont\color{caseink}#5%
  &#1%
  \\[0.12em]%
}
% \omainrow：頂層項目，無段碼（Leitsätze、Urteil、Gründe …）
\newcommand{\omainrow}[4]{%
  \footnotesize\hangindent=1.92cm\hangafter=1%
    \makebox[1.92cm][l]{\bfseries\color{caseink}#1}\color{caseink}#2%
  &\footnotesize\hangindent=1.92cm\hangafter=1%
    \makebox[1.92cm][l]{\bfseries\color{caseink}#3}\songfont\color{caseink}#4%
  &%
  \\[0.15em]%
}
% \omainrowrn：頂層項目，有段碼（不同意見等標籤寬於 \osecrow 框者）
\newcommand{\omainrowrn}[5]{%
  \footnotesize\hangindent=1.92cm\hangafter=1%
    \makebox[1.92cm][l]{\bfseries\color{caseink}#2}\color{caseink}#3%
  &\footnotesize\hangindent=1.92cm\hangafter=1%
    \makebox[1.92cm][l]{\bfseries\color{caseink}#4}\songfont\color{caseink}#5%
  &#1%
  \\[0.15em]%
}
% \osecrow：一級子項，無段碼（A.、C.、D. 等有子項者）
\newcommand{\osecrow}[4]{%
  \footnotesize\hspace{1.10cm}\hangindent=1.98cm\hangafter=1%
    \makebox[0.86cm][l]{\bfseries\color{caseink}#1}\color{caseink}#2%
  &\footnotesize\hspace{1.10cm}\hangindent=1.98cm\hangafter=1%
    \makebox[0.86cm][l]{\bfseries\color{caseink}#3}\songfont\color{caseink}#4%
  &%
  \\[0.11em]%
}
% \osecrowrn：一級子項，有段碼（B.、E. 等完整段落範圍者）
% #1=段碼字串（如 101–106）, #2=DE標籤, #3=DE正文, #4=ZH標籤, #5=ZH正文
\newcommand{\osecrowrn}[5]{%
  \footnotesize\hspace{1.10cm}\hangindent=1.98cm\hangafter=1%
    \makebox[0.86cm][l]{\bfseries\color{caseink}#2}\color{caseink}#3%
  &\footnotesize\hspace{1.10cm}\hangindent=1.98cm\hangafter=1%
    \makebox[0.86cm][l]{\bfseries\color{caseink}#4}\songfont\color{caseink}#5%
  &#1%
  \\[0.11em]%
}
% \osubrow：二級子項，有段碼範圍（I.、II.、A.I. …）
% #1=段碼字串, #2=DE標籤, #3=DE正文, #4=ZH標籤, #5=ZH正文
\newcommand{\osubrow}[5]{%
  \footnotesize\hspace{1.40cm}\hangindent=2.30cm\hangafter=1%
    \makebox[0.90cm][l]{\bfseries\color{caseink}#2}\color{caseink}#3%
  &\footnotesize\hspace{1.40cm}\hangindent=2.30cm\hangafter=1%
    \makebox[0.90cm][l]{\bfseries\color{caseink}#4}\songfont\color{caseink}#5%
  &#1%
  \\[0.09em]%
}
% \ointrow：引入段落，有單一段碼，無標籤
% #1=段碼字串, #2=DE正文, #3=ZH正文
\newcommand{\ointrow}[3]{%
  \footnotesize\hspace{0.52cm}\hangindent=0.52cm\hangafter=1\color{caseink}#2%
  &\footnotesize\hspace{0.52cm}\hangindent=0.52cm\hangafter=1\songfont\color{caseink}#3%
  &#1%
  \\[0.09em]%
}
% \osubsubrow：三級子項（1.、2.、3.）
% #1=段碼字串, #2=DE標籤, #3=DE正文, #4=ZH標籤, #5=ZH正文
\newcommand{\osubsubrow}[5]{%
  \footnotesize\hspace{1.80cm}\hangindent=2.48cm\hangafter=1%
    \makebox[0.68cm][l]{\bfseries\color{caseink}#2}\color{caseink}#3%
  &\footnotesize\hspace{1.80cm}\hangindent=2.48cm\hangafter=1%
    \makebox[0.68cm][l]{\bfseries\color{caseink}#4}\songfont\color{caseink}#5%
  &#1%
  \\[0.08em]%
}
% \osubsubsubrow：四級子項（a）、b））
% #1=段碼字串, #2=DE標籤, #3=DE正文, #4=ZH標籤, #5=ZH正文
\newcommand{\osubsubsubrow}[5]{%
  \footnotesize\hspace{2.10cm}\hangindent=2.68cm\hangafter=1%
    \makebox[0.58cm][l]{\bfseries\color{caseink}#2}\color{caseink}#3%
  &\footnotesize\hspace{2.10cm}\hangindent=2.68cm\hangafter=1%
    \makebox[0.58cm][l]{\bfseries\color{caseink}#4}\songfont\color{caseink}#5%
  &#1%
  \\[0.07em]%
}
% \osubsubsubsubrow：五級子項（aa）、bb））
% 標籤框 0.62cm：足以容納「aa)」在 footnotesize bold 下（約 0.50cm）並留有間距
% #1=段碼字串, #2=DE標籤, #3=DE正文, #4=ZH標籤, #5=ZH正文
\newcommand{\osubsubsubsubrow}[5]{%
  \footnotesize\hspace{2.38cm}\hangindent=3.06cm\hangafter=1%
    \makebox[0.68cm][l]{\bfseries\color{caseink}#2}\color{caseink}#3%
  &\footnotesize\hspace{2.38cm}\hangindent=3.06cm\hangafter=1%
    \makebox[0.68cm][l]{\bfseries\color{caseink}#4}\songfont\color{caseink}#5%
  &#1%
  \\[0.06em]%
}
% \outlinepage：雙語對照大綱頁包裝環境（longtable 自動跨頁）
% 三欄：德文欄 + 中文欄 + 段碼欄（最右，不折行）
% 段碼欄寬度：1.80cm，足以容納最長段碼如「309–348」
% DE/ZH 欄寬：(0.87\textwidth - 2.08cm) / 2
%   驗算：2×欄 + 0.05\textwidth + 0.28cm gap + 1.80cm = 0.87tw - 2.08cm + 0.05tw + 2.08cm = 0.92tw ✓
\newenvironment{outlinepage}{%
  \vspace*{0.40cm}%
  \setlength{\LTleft}{0.04\textwidth}%
  \setlength{\LTright}{0.04\textwidth}%
  \setlength{\tabcolsep}{0pt}%
  % 大綱列距由 longtable 的 \arraystretch 與各 row macro 結尾的 \\[...em] 共同決定。
  % 這裡控制「每一列本身的表格 strut 高度」；數值越大，A./I./1. 各項之間越像有空行。
  % A3 原本用 0.62，視覺上沒有空行；A4 也沿用同值，避免切到 A4 後項目被撐開。
  % 若日後覺得大綱太擠，只微調這個值（例如 0.68），不要改各 \omainrow/\osubrow 的縮排。
  \renewcommand{\arraystretch}{0.62}%
  \begin{longtable}{%
    >{\vspace{0pt}\raggedright\arraybackslash}p{\dimexpr(0.87\textwidth-2.08cm)/2\relax}%
    @{\hspace{0.025\textwidth}}%
    !{\color{notegray}\vrule width 0.12pt}%
    @{\hspace{0.025\textwidth}}%
    >{\vspace{0pt}\raggedright\arraybackslash}p{\dimexpr(0.87\textwidth-2.08cm)/2\relax}%
    @{\hspace{0.28cm}}%
    >{\vspace{0pt}\footnotesize\normalfont\color{pendingink}\raggedright\arraybackslash}p{1.80cm}%
    @{}%
  }%
  % 首頁欄標籤：不要橫跨置中；分別放在德文欄與中文欄的實際欄位起點。
  \footnotesize\bfseries\color{sourceink}Gliederung%
  &\footnotesize\bfseries\songfont\color{sourceink}大綱%
  &%
  \\[0.36em]%
  \endfirsthead%
  % 續頁欄標籤：同樣分欄定位，以灰色細體區分；無需標注「續」。
  \footnotesize\color{pendingink}Gliederung%
  &\footnotesize\songfont\color{pendingink}大綱%
  &%
  \\[0.24em]%
  \endhead%
}{%
  \end{longtable}%
}
 載入。
% 不含：\documentclass、\documentversion、\ggversionde/zh、\tocde/\toczh。
% 日期與首頁版本列由本檔自動產生；各文件只需手動指定 \documentversion。

\usepackage{xeCJK}
\usepackage{fontspec}
\usepackage{geometry}
\usepackage{setspace}
\usepackage{hyperref}
\usepackage{xurl}
\usepackage{bookmark}
\usepackage{enumitem}
\usepackage{xcolor}
\usepackage{titlesec}
\usepackage{array}
\usepackage{longtable}
\usepackage{tcolorbox}
\usepackage{environ}
\usepackage{pdflscape}
\usepackage{etoolbox}
\usepackage{ragged2e}
\usepackage{paracol}
\usepackage{multicol}
\usepackage{needspace}
\usepackage{fancyhdr}
\usepackage{tikz}
\usepackage{zref-abspage}
\tcbuselibrary{breakable,skins}
\usetikzlibrary{calc}
\usepackage{pgfcalendar}

\hypersetup{
  unicode=true,
  bookmarksopen=true,
  bookmarksnumbered=false,
  hidelinks
}

% ===== 自動推導，不要手動改下面 =====
% （\docmode 與 \bilingualpaper 由各文件在 % bverfge_layout.tex
% 共用排版基礎設施，供 Schwangerschaftsabbruch I 與 II 共享。
% 由各文件以 \documentclass 後 \input{../../共用LaTeX/bverfge_layout} 載入。
% 不含：\documentclass、\documentversion、\ggversionde/zh、\tocde/\toczh。
% 日期與首頁版本列由本檔自動產生；各文件只需手動指定 \documentversion。

\usepackage{xeCJK}
\usepackage{fontspec}
\usepackage{geometry}
\usepackage{setspace}
\usepackage{hyperref}
\usepackage{xurl}
\usepackage{bookmark}
\usepackage{enumitem}
\usepackage{xcolor}
\usepackage{titlesec}
\usepackage{array}
\usepackage{longtable}
\usepackage{tcolorbox}
\usepackage{environ}
\usepackage{pdflscape}
\usepackage{etoolbox}
\usepackage{ragged2e}
\usepackage{paracol}
\usepackage{multicol}
\usepackage{needspace}
\usepackage{fancyhdr}
\usepackage{tikz}
\usepackage{zref-abspage}
\tcbuselibrary{breakable,skins}
\usetikzlibrary{calc}
\usepackage{pgfcalendar}

\hypersetup{
  unicode=true,
  bookmarksopen=true,
  bookmarksnumbered=false,
  hidelinks
}

% ===== 自動推導，不要手動改下面 =====
% （\docmode 與 \bilingualpaper 由各文件在 \input{bverfge_layout} 之前設定；
%   預設 chinese / a4）
\ifdefined\docmode\else\def\docmode{chinese}\fi
\ifdefined\bilingualpaper\else\def\bilingualpaper{a4}\fi
\newif\ifshoworiginal
\newif\ifshowtranslation
\newif\ifbilinguallandscape
\newif\ifbilingualcolumns
\newif\ifintranslation
\newif\iffirstbilingualsection

\showoriginalfalse
\showtranslationfalse
\bilinguallandscapefalse
\bilingualcolumnsfalse
\intranslationfalse
\firstbilingualsectionfalse

\def\modechinese{chinese}
\def\modegerman{german}
\def\modebilingual{bilingual}
\def\bilingualpaperaThree{a3}
\def\bilingualpaperaFour{a4}

\ifx\docmode\modechinese
  \showtranslationtrue
\fi

\ifx\docmode\modegerman
  \showoriginaltrue
\fi

\ifx\docmode\modebilingual
  \showoriginaltrue
  \showtranslationtrue
  \bilinguallandscapetrue
  \bilingualcolumnstrue
\fi

% 編排模式說明：
% chinese  → \showoriginalfalse  \showtranslationtrue  （A4 直式，純中文）
% german   → \showoriginaltrue   \showtranslationfalse （A4 直式，純德文）
% bilingual→ \showoriginaltrue   \showtranslationtrue  \bilingualcolumnstrue（A3/A4 橫式對照）

\ifbilingualcolumns
  \ifx\bilingualpaper\bilingualpaperaFour
    \geometry{
      a4paper,
      landscape,
      left=1.25cm,
      right=2.25cm,
      top=1.25cm,
      bottom=1.75cm,
      footskip=8.5mm,
      marginparwidth=0.75cm,
      marginparsep=0.25cm
    }
  \else
    \geometry{
      a3paper,
      landscape,
      left=1.55cm,
      right=2.75cm,
      top=1.55cm,
      bottom=2.05cm,
      footskip=9.5mm,
      marginparwidth=0.9cm,
      marginparsep=0.35cm
    }
  \fi
\else
\geometry{
  a4paper,
  portrait,
  left=2.2cm,
  right=2.6cm,
  top=2.0cm,
  bottom=2.0cm,
  marginparwidth=0.8cm,
  marginparsep=0.3cm
}
\fi
% ===== 時間戳基礎設施 =====
\newcount\stamptimeval
\newcount\stampminuteval
\newcommand{\stamphh}{%
  \stamptimeval=\time
  \divide\stamptimeval by 60
  \ifnum\stamptimeval<10 0\fi
  \the\stamptimeval}
\newcommand{\stampmm}{%
  \stamptimeval=\time
  \stampminuteval=\time
  \divide\stampminuteval by 60
  \multiply\stampminuteval by 60
  \advance\stamptimeval by -\stampminuteval
  \ifnum\stamptimeval<10 0\fi
  \the\stamptimeval}
\newcommand{\stampwkde}[1]{%
  \ifcase#1Montag\or Dienstag\or Mittwoch\or Donnerstag\or Freitag\or Samstag\or Sonntag\fi}
\newcommand{\stampwkzh}[1]{%
  \ifcase#1 一\or 二\or 三\or 四\or 五\or 六\or 日\fi}
\newcommand{\stampmonthde}[1]{%
  \ifcase#1\or Januar\or Februar\or März\or April\or Mai\or Juni\or
  Juli\or August\or September\or Oktober\or November\or Dezember\fi}
\newcount\stampjdval
\newcount\stampwdval
\pgfcalendardatetojulian{\the\year-\the\month-\the\day}{\stampjdval}
\pgfcalendarjuliantoweekday{\the\stampjdval}{\stampwdval}
\newcommand{\timestampzh}{%
  \songfont 最後修改：\the\year 年\the\month 月\the\day 日%
  % （星期\stampwkzh{\the\stampwdval}）%
  % \space\stamphh:\stampmm
  }

\newcommand{\documentdatezh}{\the\year 年 \the\month 月 \the\day 日}
\newcommand{\documentdatede}{\the\day. \stampmonthde{\the\month} \the\year}
\newcommand{\bverfgetranslationnote}{%
  {\songfont 翻譯說明：初譯使用 Claude Code 與 GPT 協助迭代；仍請以德文原文為準。}%
}
\newcommand{\bverfgecourseinfo}{%
  {\songfont 課程：114 學年度第 2 學期；憲法專題研究（四）；授課老師：湯德宗教授；導讀日期：115 年 6 月 1 日。}%
}
\newcommand{\bverfgeeditorinfo}{%
  {\songfont 編者：王逸帆（R10A21126），國立台灣大學法律系研究所碩士班。}%
}
\newcommand{\bverfgetitleversionline}{%
  \ifbilingualcolumns
    Fassung \documentversion\enspace\textperiodcentered\enspace Stand: \documentdatede
    \quad
    {\songfont 版本 \documentversion\enspace\textperiodcentered\enspace 修訂日期：\documentdatezh\par}%
  \else
    \ifshowtranslation
      {\songfont 版本 \documentversion\enspace\textperiodcentered\enspace 修訂日期：\documentdatezh\par}%
    \else
      Fassung \documentversion\enspace\textperiodcentered\enspace Stand: \documentdatede\par
    \fi
  \fi
}
\newcommand{\bverfgetitlecornerinfo}{%
  \ifshowtranslation
    \vspace{0.72em}%
    \noindent\hfill
    \begin{minipage}{0.66\textwidth}
      \raggedleft\tiny\color{pendingink}\songfont
      \bverfgecourseinfo\par
      \bverfgeeditorinfo
    \end{minipage}%
  \fi
}
\newcommand{\bverfgetitlemeta}{%
  \vfill
  \begin{center}%
    \footnotesize\color{pendingink}%
    \bverfgetitleversionline
    \ifshowtranslation
      \vspace{0.28em}{\scriptsize\bverfgetranslationnote\par}%
    \fi
  \end{center}%
  \bverfgetitlecornerinfo
}

\pagestyle{fancy}
\fancyhf{}
\renewcommand{\headrulewidth}{0pt}
\renewcommand{\footrulewidth}{0pt}
\fancyfoot[C]{\thepage}
\ifbilingualcolumns
  \fancyfoot[R]{\scriptsize\color{pendingink}\timestampzh}%
\else
  \ifshoworiginal
  \else
    \fancyfoot[R]{\scriptsize\color{pendingink}\timestampzh\hspace{0.45cm}}%
  \fi
\fi
\setmainfont{Times New Roman}
\setsansfont{Helvetica Neue}
\setmonofont{Menlo}
\IfFileExists{/Users/iw/Library/Fonts/kaiu.ttf}{
  \setCJKmainfont[Path=/Users/iw/Library/Fonts/,AutoFakeBold=4]{kaiu.ttf}
}{
  \IfFontExistsTF{標楷體}{
    \setCJKmainfont[AutoFakeBold=4]{標楷體}
  }{
    \setCJKmainfont{Songti TC}
  }
}
\IfFontExistsTF{Songti TC}{
  \setCJKfamilyfont{song}{Songti TC}
  \setCJKmonofont{Songti TC}
}{
  \setCJKfamilyfont{song}[Path=/Users/iw/Library/Fonts/,AutoFakeBold=4]{kaiu.ttf}
  \setCJKmonofont[Path=/Users/iw/Library/Fonts/,AutoFakeBold=4]{kaiu.ttf}
}
\newcommand{\songfont}{\CJKfamily{song}}

% ===== 封面／目錄專用打字機字體（僅限封面、目錄頁使用，不影響正文）=====
\IfFontExistsTF{Radio Newsman}{%
  \newfontfamily\bverfgetitlede{Radio Newsman}%
}{%
  \newfontfamily\bverfgetitlede{Courier New}%
}
\IfFontExistsTF{Erikas Farbband}{%
  \newfontfamily\bverfgebodyde{Erikas Farbband}%
}{%
  \let\bverfgebodyde\bverfgetitlede
}
\IfFontExistsTF{Maoken Heavy Labourer}{%
  \setCJKfamilyfont{bverfgetitlecjk}{Maoken Heavy Labourer}%
}{%
  \setCJKfamilyfont{bverfgetitlecjk}{Songti TC}%
}
\IfFontExistsTF{Huiwen-mincho}{%
  \setCJKfamilyfont{bverfgebodycjk}{Huiwen-mincho}%
}{%
  \setCJKfamilyfont{bverfgebodycjk}{Songti TC}%
}
\newcommand{\bverfgetitlezh}{\bverfgetitlede\CJKfamily{bverfgetitlecjk}}
\newcommand{\bverfgebodyzh}{\bverfgebodyde\CJKfamily{bverfgebodycjk}}

\setstretch{1.35}
\setlist{itemsep=0.2em, topsep=0.4em}
\setlength{\parindent}{0pt}
\setlength{\parskip}{0.45em}
\raggedbottom
\setcounter{secnumdepth}{0}
\setcounter{tocdepth}{2}

\newcommand{\lawboxsourcefont}{\footnotesize}
\newcommand{\lawboxtransfont}{\footnotesize}
\newcommand{\lawboxsourcestretch}{1.02}
\newcommand{\lawboxtransstretch}{0.92}

\titleformat{\section}{\centering\Large\bfseries\songfont}{\thesection}{1em}{}
\titleformat{\subsection}{\large\bfseries\songfont}{\thesubsection}{1em}{}
\titleformat{\subsubsection}{\normalsize\bfseries\songfont}{\thesubsubsection}{1em}{}
\titleformat{\paragraph}{\normalsize\bfseries\songfont}{\theparagraph}{1em}{}
\titlespacing*{\section}{0pt}{1.8em}{1.0em}
\titlespacing*{\subsection}{0pt}{1.2em}{0.6em}
\titlespacing*{\subsubsection}{0pt}{1.0em}{0.45em}
\titlespacing*{\paragraph}{0pt}{0.6em}{0.4em}

\definecolor{noteblue}{HTML}{A9C9F5}
\definecolor{noteteal}{HTML}{AEE1D6}
\definecolor{notegray}{HTML}{D7DADF}
\definecolor{caseink}{HTML}{000000}
\definecolor{casepaper}{HTML}{FCFEFD}
\definecolor{sourcepaper}{HTML}{FAFBF8}
\definecolor{sourceline}{HTML}{D7DED8}
\definecolor{sourceink}{HTML}{000000}
\definecolor{pendingink}{HTML}{000000}

\tcbset{
  caseboxbase/.style={
    enhanced,
    breakable,
    width=0.985\linewidth,
    colback=casepaper,
    coltitle=caseink,
    fonttitle=\bfseries\songfont,
    boxrule=0.28pt,
    arc=1.2mm,
    left=2.4mm,
    right=2.4mm,
    top=1.5mm,
    bottom=1.5mm,
    boxsep=1.5mm,
    before skip=0.65em,
    after skip=0.65em,
    before upper={\setlength{\parindent}{0pt}\setlength{\parskip}{0.35em}},
  }
}

\newcommand{\bilingualstart}{}
\newcommand{\bilingualstop}{}
\ifbilingualcolumns
  \renewcommand{\bilingualstart}{\small\setlength{\columnsep}{1.25cm}\setlength{\columnseprule}{0.12pt}\columnratio{0.5,0.5}\multicolumnfootnotes\flushbottom\sloppy\interlinepenalty=80\clubpenalty=800\widowpenalty=800\displaywidowpenalty=800\begin{paracol}{2}\global\intranslationfalse\global\firstbilingualsectiontrue{\footnotesize\color{sourceink}Deutsch\par}\switchcolumn{\footnotesize\songfont\color{sourceink}中文譯文\par}\switchcolumn*}
  \renewcommand{\bilingualstop}{\ifintranslation\switchcolumn*\global\intranslationfalse\fi\end{paracol}\clearpage\raggedbottom}
\fi
\newif\ifrandnumbers
\randnumbersfalse
\newcounter{randnumber}
\newcounter{randnumberlocal}
\newcounter{lastsourcerandstart}
\newif\iflastsourcehasrandnumbers
\lastsourcehasrandnumbersfalse
\newif\ifinlawbox
\inlawboxfalse
\newif\iflawboxhaspendingrn
\lawboxhaspendingrnfalse
\newcommand{\lawboxpendingrn}{}
\newcommand{\startrandnumbers}{\global\randnumberstrue\setcounter{randnumber}{1}\global\lastsourcehasrandnumbersfalse}
\newcommand{\stoprandnumbers}{\global\randnumbersfalse\global\lastsourcehasrandnumbersfalse}
\newlength{\randnumberwidth}
\setlength{\randnumberwidth}{0.95cm}
\newlength{\randnumberrightoffset}
\setlength{\randnumberrightoffset}{0.85cm}
\newlength{\randnumberlawboxrightoffset}
\setlength{\randnumberlawboxrightoffset}{1.40cm}
\newcommand{\randnumberbox}[1]{%
  \leavevmode
  \makebox[0pt][l]{%
    \smash{%
      \tikz[remember picture,overlay,baseline=(rnmark.base)]{%
        \coordinate (rnmark) at (0,0);%
        \ifinlawbox
          \path let \p1=(current page.east), \p2=(rnmark.base) in
            node[
              anchor=base east,
              inner sep=0pt,
              outer sep=0pt,
              text width=\randnumberwidth,
              align=right,
              text=pendingink,
              font=\normalfont\mdseries\scriptsize\ttfamily
            ] at (\x1-\randnumberlawboxrightoffset,\y2) {{\normalfont\mdseries\scriptsize\ttfamily #1}};%
        \else
          \path let \p1=(current page.east), \p2=(rnmark.base) in
            node[
              anchor=base east,
              inner sep=0pt,
              outer sep=0pt,
              text width=\randnumberwidth,
              align=right,
              text=pendingink,
              font=\normalfont\mdseries\scriptsize\ttfamily
            ] at (\x1-\randnumberrightoffset,\y2) {{\normalfont\mdseries\scriptsize\ttfamily #1}};%
        \fi
      }%
    }%
  }%
  \ignorespaces
}
\newcommand{\lawboxstorerandnumber}[1]{%
  \gdef\lawboxpendingrn{#1}%
  \global\lawboxhaspendingrntrue
  \ignorespaces
}
\newcommand{\lawboxuserandnumber}{%
  \iflawboxhaspendingrn
    \randnumberbox{\lawboxpendingrn}%
    \global\lawboxhaspendingrnfalse
  \fi
}
\newcommand{\sourcern}[1]{%
  \ifrandnumbers
    \ifshoworiginal
      \ifshowtranslation\else
        \ifinlawbox\lawboxstorerandnumber{#1}\else\randnumberbox{#1}\fi
      \fi
    \fi
  \fi
}
\newcommand{\transrn}[1]{%
  \ifrandnumbers
    \ifshowtranslation
      \ifinlawbox\lawboxstorerandnumber{#1}\else\randnumberbox{#1}\fi
    \fi
  \fi
}
% ===== 課堂模式：德文原文縮寫註解 =====
% 日常切換只改各判決檔的一行：
%   \classroommodeon        開啟課堂版；註解掉這行即回到一般版。
% 模板預設：
%   1. 課堂模式關閉。
%   2. 課堂模式開啟時，縮寫註解包含「德文全稱＋中文譯名」。
%   3. 課堂模式開啟時，GG/條文編者註仍會自動顯示。
% 進階微調才需要在單案檔覆寫：
%   \classroomabbrzhoff     縮寫註解僅顯示德文全稱。
%   \showeditornotestrue    一般版也顯示編者註。
% 註解策略：同一實際 PDF 頁面中，同一縮寫只在第一次出現處加註。
% 這裡用 zref-abspage 的絕對頁序，不用 \thepage，避免文件內頁碼重置時誤判。
\newif\ifclassroommode
\newif\ifclassroomabbrzh
\classroommodefalse
\classroomabbrzhtrue
\newcommand{\classroommodeon}{\classroommodetrue}
\newcommand{\classroommodeoff}{\classroommodefalse}
\newcommand{\classroomabbrzhon}{\classroomabbrzhtrue}
\newcommand{\classroomabbrzhoff}{\classroomabbrzhfalse}
\newcommand{\abbrnotelabel}{\emph{Abk.}: }
\newcommand{\abbrnote}[3]{%
  \ifclassroommode
    \ifshoworiginal
      \ifintranslation\else
        \footnote{\normalfont\footnotesize\abbrnotelabel #2\ifclassroomabbrzh；{\songfont #3}\fi}%
      \fi
    \fi
  \fi
}
\newcommand{\abbrnoteonce}[4]{%
  #2%
  \ifclassroommode
    \ifshoworiginal
      \ifintranslation\else
        \begingroup
          \edef\abbrcurrentabspage{\number\numexpr\value{abspage}+1\relax}%
          \ifcsname abbrnoteseen@#1@\abbrcurrentabspage\endcsname\else
            \global\expandafter\let\csname abbrnoteseen@#1@\abbrcurrentabspage\endcsname\empty
            \abbrnote{#1}{#3}{#4}%
          \fi
        \endgroup
      \fi
    \fi
  \fi
}
\newcommand{\abbrAbs}{\abbrnoteonce{abs}{Abs.}{Absatz}{項}}
\newcommand{\abbrAAO}{\abbrnoteonce{aao}{a.a.O.}{am angegebenen Ort}{前引處}}
\newcommand{\abbrAE}{\abbrnoteonce{ae}{AE}{Alternativ-Entwurf}{替代草案}}
\newcommand{\abbrArt}{\abbrnoteonce{art}{Art.}{Artikel}{條}}
\newcommand{\abbrAufl}{\abbrnoteonce{aufl}{Aufl.}{Auflage}{版}}
\newcommand{\abbrBd}{\abbrnoteonce{bd}{Bd.}{Band}{卷}}
\newcommand{\abbrBGBl}{\abbrnoteonce{bgbl}{BGBl.}{Bundesgesetzblatt}{聯邦法律公報}}
\newcommand{\abbrBGHSt}{\abbrnoteonce{bghst}{BGHSt}{Entscheidungen des Bundesgerichtshofes in Strafsachen}{聯邦普通法院刑事判決彙編}}
\newcommand{\abbrBRDrucks}{\abbrnoteonce{brdrucks}{BRDrucks.}{Bundesrats-Drucksache}{聯邦參議院文件}}
\newcommand{\abbrBTDrucks}{\abbrnoteonce{btdrucks}{BTDrucks.}{Bundestags-Drucksache}{聯邦議院文件}}
\newcommand{\abbrBundesgesetzbl}{\abbrnoteonce{bundesgesetzbl}{Bundesgesetzbl.}{Bundesgesetzblatt}{聯邦法律公報}}
\newcommand{\abbrBvF}{\abbrnoteonce{bvf}{BvF}{Bundesverfassungsgerichtsverfahren}{聯邦憲法法院程序案號}}
\newcommand{\abbrBvL}{\abbrnoteonce{bvl}{BvL}{Bundesverfassungsgerichtsverfahren}{聯邦憲法法院程序案號}}
\newcommand{\abbrBvR}{\abbrnoteonce{bvr}{BvR}{Bundesverfassungsgerichtsverfahren}{聯邦憲法法院程序案號}}
\newcommand{\abbrBVerfG}{\abbrnoteonce{bverfg}{BVerfG}{Bundesverfassungsgericht}{聯邦憲法法院}}
\newcommand{\abbrBVerfGE}{\abbrnoteonce{bverfge}{BVerfGE}{Entscheidungen des Bundesverfassungsgerichts}{聯邦憲法法院判決集}}
\newcommand{\abbrBVerfGG}{\abbrnoteonce{bverfgg}{BVerfGG}{Bundesverfassungsgerichtsgesetz}{聯邦憲法法院法}}
\newcommand{\abbrDDR}{\abbrnoteonce{ddr}{DDR}{Deutsche Demokratische Republik}{德意志民主共和國；東德}}
\newcommand{\abbrDJT}{\abbrnoteonce{djt}{DJT}{Deutscher Juristentag}{德國法學家大會}}
\newcommand{\abbrDr}{\abbrnoteonce{dr}{Dr.}{Doktor}{Dr.}}
\newcommand{\abbrEuGRZ}{\abbrnoteonce{eugrz}{EuGRZ}{Europäische Grundrechte-Zeitschrift}{歐洲基本權期刊}}
\newcommand{\abbrEV}{\abbrnoteonce{ev}{EV}{Einigungsvertrag}{統一條約}}
\newcommand{\abbrF}{\abbrnoteonce{f}{f.}{folgende}{以下一頁；下一段}}
\newcommand{\abbrFF}{\abbrnoteonce{ff}{ff.}{fortfolgende}{以下數頁；以下各段}}
\newcommand{\abbrGBlDDR}{\abbrnoteonce{gblddr}{GBl. DDR}{Gesetzblatt der Deutschen Demokratischen Republik}{德意志民主共和國法律公報}}
\newcommand{\abbrGG}{\abbrnoteonce{gg}{GG}{Grundgesetz}{基本法}}
\newcommand{\abbrGVBl}{\abbrnoteonce{gvbl}{GVBl.}{Gesetz- und Verordnungsblatt}{法律及命令公報}}
\newcommand{\abbrJR}{\abbrnoteonce{jr}{JR}{Juristische Rundschau}{法學評論}}
\newcommand{\abbrJZ}{\abbrnoteonce{jz}{JZ}{JuristenZeitung}{法學家報}}
\newcommand{\abbrIVm}{\abbrnoteonce{ivm}{i.V.m.}{in Verbindung mit}{結合；連同}}
\newcommand{\abbrMdB}{\abbrnoteonce{mdb}{MdB}{Mitglied des Bundestages}{聯邦議院議員}}
\newcommand{\abbrMwN}{\abbrnoteonce{mwn}{m.w.N.}{mit weiteren Nachweisen}{附進一步文獻／判例出處}}
\newcommand{\abbrNF}{\abbrnoteonce{nf}{n.F.}{neue Fassung}{新版本}}
\newcommand{\abbrAF}{\abbrnoteonce{af}{a.F.}{alte Fassung}{舊版本}}
\newcommand{\abbrNr}{\abbrnoteonce{nr}{Nr.}{Nummer}{款、目或號，依語境}}
\newcommand{\abbrProf}{\abbrnoteonce{prof}{Prof.}{Professor}{教授}}
\newcommand{\abbrRdnr}{\abbrnoteonce{rdnr}{Rdnr.}{Randnummer}{邊碼；段碼}}
\newcommand{\abbrRdnrn}{\abbrnoteonce{rdnrn}{Rdnrn.}{Randnummern}{邊碼；段碼}}
\newcommand{\abbrRGBl}{\abbrnoteonce{rgbl}{RGBl.}{Reichsgesetzblatt}{帝國法律公報}}
\newcommand{\abbrRGSt}{\abbrnoteonce{rgst}{RGSt}{Entscheidungen des Reichsgerichts in Strafsachen}{帝國法院刑事判決彙編}}
\newcommand{\abbrRspr}{\abbrnoteonce{rspr}{Rspr.}{Rechtsprechung}{判例；司法實務}}
\newcommand{\abbrRVO}{\abbrnoteonce{rvo}{RVO}{Reichsversicherungsordnung}{帝國保險法}}
\newcommand{\abbrS}{\abbrnoteonce{s}{S.}{Seite}{頁}}
\newcommand{\abbrSFHG}{\abbrnoteonce{sfhg}{SFHG}{Schwangeren- und Familienhilfegesetz}{孕婦及家庭扶助法}}
\newcommand{\abbrSGBV}{\abbrnoteonce{sgbv}{SGB V}{Fünftes Buch Sozialgesetzbuch}{社會法典第五編}}
\newcommand{\abbrStAG}{\abbrnoteonce{staeg}{StÄG}{Strafrechtsänderungsgesetz}{刑法修正法}}
\newcommand{\abbrStenBer}{\abbrnoteonce{stenber}{StenBer.}{Stenographischer Bericht}{速記紀錄}}
\newcommand{\abbrStGB}{\abbrnoteonce{stgb}{StGB}{Strafgesetzbuch}{刑法}}
\newcommand{\abbrStREG}{\abbrnoteonce{streg}{StREG}{Strafrechtsreform-Ergänzungsgesetz}{刑法改革補充法}}
\newcommand{\abbrStrRG}{\abbrnoteonce{strrg}{StrRG}{Strafrechtsreformgesetz}{刑法改革法}}
\newcommand{\abbrUS}{\abbrnoteonce{us}{U.S.}{United States Reports}{美國最高法院判決彙編}}
\newcommand{\abbrVgl}{\abbrnoteonce{vgl}{vgl.}{vergleiche}{參見}}
\newcommand{\abbrVVDStRL}{\abbrnoteonce{vvdstrl}{VVDStRL}{Veröffentlichungen der Vereinigung der Deutschen Staatsrechtslehrer}{德國公法教師協會論文集}}
\newcommand{\abbrWP}{\abbrnoteonce{wp}{WP.}{Wahlperiode}{屆期}}
\newcommand{\abbrWp}{\abbrnoteonce{wp}{Wp.}{Wahlperiode}{屆期}}
\newcommand{\abbrZStrW}{\abbrnoteonce{zstrw}{ZStrW}{Zeitschrift für die gesamte Strafrechtswissenschaft}{整體刑法學期刊}}
\newcommand{\recordsourcerandstart}{%
  \ifrandnumbers
    \setcounter{lastsourcerandstart}{\value{randnumber}}%
    \global\lastsourcehasrandnumberstrue
  \else
    \global\lastsourcehasrandnumbersfalse
  \fi
}
\newlength{\biparagraphskip}
\setlength{\biparagraphskip}{0.52em}
\newlength{\lawboxblockbeforeskip}
\setlength{\lawboxblockbeforeskip}{0.72em}
\newlength{\lawboxblockafterskip}
\setlength{\lawboxblockafterskip}{0.92em}
\newif\iflawboxpartendedblock
\lawboxpartendedblockfalse
\newcommand{\sourceparstyle}{\footnotesize\color{sourceink}\setstretch{1.12}\RaggedRight\emergencystretch=3em\setlength{\parskip}{0.42em}\obeylines\selectfont}
\newcommand{\transparstyle}{\songfont\color{caseink}\setstretch{1.12}\setlength{\parindent}{0pt}\setlength{\parskip}{0.42em}}
\newcommand{\bilingualpairskip}{%
  \par\addvspace{\biparagraphskip}%
  \switchcolumn
  \par\addvspace{\biparagraphskip}%
  \switchcolumn*%
  \global\intranslationfalse
}
\newcommand{\bilingualboxpairskip}{%
  \par
  \switchcolumn
  \par
  \switchcolumn*%
  \global\intranslationfalse
}
\newcommand{\bilinguallawboxblockskip}{%
  \switchcolumn*%
  \par\addvspace{\lawboxblockafterskip}%
  \switchcolumn
  \par\addvspace{\lawboxblockafterskip}%
  \switchcolumn*%
  \global\intranslationfalse
}
\newcommand{\bikeepwithnext}[1][4]{%
  \ifbilingualcolumns
    \syncleft
    \Needspace{#1\baselineskip}%
    \switchcolumn
    \Needspace{#1\baselineskip}%
    \switchcolumn*%
    \global\intranslationfalse
  \else
    \Needspace{#1\baselineskip}%
  \fi
}
\NewEnviron{sourcepara}[1][]{
  \recordsourcerandstart
  \ifshoworiginal
    \ifbilingualcolumns
      \ifintranslation\switchcolumn*\global\intranslationfalse\fi
    \else
      \Needspace{5\baselineskip}
    \fi
    \ifbilingualcolumns\else\par\addvspace{0.35em}\fi
    {\sourceparstyle
    \noindent\BODY\par}
    \ifbilingualcolumns\else\addvspace{0.25em}\fi
    \ifbilingualcolumns
      \switchcolumn
      \global\intranslationtrue
    \fi
  \else
    \ifrandnumbers
      \begingroup
        \setbox0=\vbox{\hsize=\linewidth\sourceparstyle\noindent\BODY\par}%
      \endgroup
    \fi
  \fi
}
\NewEnviron{transpara}{
  \ifshowtranslation
    \setcounter{randnumberlocal}{\value{lastsourcerandstart}}%
    \ifbilingualcolumns
      \ifintranslation\else\switchcolumn\global\intranslationtrue\fi
      {\transparstyle\setlength{\leftskip}{0.55em}\setlength{\rightskip}{0.15em plus 1em}\addtolength{\linewidth}{-0.55em}\BODY\par}
      \switchcolumn*
      \bilingualpairskip
    \else
      {\transparstyle\BODY\par}
    \fi
  \else
    \ifbilingualcolumns
      \ifintranslation\switchcolumn*\global\intranslationfalse\fi
    \fi
  \fi
}
\NewEnviron{sourceboxpara}[1][]{
  \global\lawboxpartendedblockfalse
  \recordsourcerandstart
  \ifshoworiginal
    \ifbilingualcolumns
      \ifintranslation\switchcolumn*\global\intranslationfalse\fi
    \else
      \Needspace{3\baselineskip}
    \fi
    {\sourceparstyle
    \BODY\par}
    \ifbilingualcolumns\else
      \iflawboxpartendedblock\addvspace{\lawboxblockafterskip}\fi
    \fi
    \ifbilingualcolumns
      \switchcolumn
      \global\intranslationtrue
    \fi
  \fi
}
\NewEnviron{transboxpara}{
  \global\lawboxpartendedblockfalse
  \ifshowtranslation
    \setcounter{randnumberlocal}{\value{lastsourcerandstart}}%
    \ifbilingualcolumns
      \ifintranslation\else\switchcolumn\global\intranslationtrue\fi
      {\transparstyle\BODY\par}
      \iflawboxpartendedblock
        \bilinguallawboxblockskip
      \else
        \switchcolumn*
        \bilingualboxpairskip
      \fi
    \else
      {\transparstyle\BODY\par}
      \iflawboxpartendedblock\addvspace{\lawboxblockafterskip}\fi
    \fi
  \else
    \ifbilingualcolumns
      \ifintranslation\switchcolumn*\global\intranslationfalse\fi
    \fi
  \fi
}
\newcommand{\leitsatznum}[1]{\noindent\hangindent=3.0em\hangafter=1\makebox[2.25em][r]{\bfseries #1.}\hspace{0.55em}\ignorespaces}
\newcommand{\syncleft}{\ifbilingualcolumns\ifintranslation\switchcolumn*\global\intranslationfalse\fi\fi}
\pretocmd{\section}{\syncleft\clearpage}{}{}
\pretocmd{\subsection}{\syncleft}{}{}
\pretocmd{\subsubsection}{\syncleft}{}{}
\newcommand{\biheadingfont}{\songfont}
\newcommand{\bitocentry}[2]{%
  \ifshoworiginal
    \ifshowtranslation
      \ifstrequal{#1}{#2}{#1}{#1\space #2}%
    \else
      #1%
    \fi
  \else
    #2%
  \fi
}
\pdfstringdefDisableCommands{%
  \def\bitocentry#1#2{#1}%
}
\newcommand{\biheading}[7]{%
  \syncleft#1\bikeepwithnext[#3]\phantomsection\addcontentsline{toc}{#2}{\bitocentry{#6}{#7}}%
  \ifbilingualcolumns
    {#4 #6\par}%
    \switchcolumn
    {#4\biheadingfont #7\par}%
    \switchcolumn*%
    \global\intranslationfalse
    \par\addvspace{#5}%
    \switchcolumn
    \par\addvspace{#5}%
    \switchcolumn*%
  \else
    \ifshoworiginal{#4 #6\par}\fi
    \ifshowtranslation{#4\biheadingfont #7\par}\fi
    \addvspace{#5}%
  \fi
}
\newcommand{\termsectionheading}[1]{%
  \par\Needspace{9\baselineskip}%
  \vspace{0.65em}%
  {\songfont\bfseries #1\par}%
  \vspace{0.2em}%
}
% 章節換頁：
%   雙語 paracol 欄內不要用 \clearpage；它會在同步兩欄時吐出只有欄線/頁碼的空頁。
%   \newpage 足以讓下一個 section 起新頁，又不會強制清空所有待排材料。
%   單語模式不在 paracol 內，仍可用 \clearpage。
\newcommand{\bisectionpagebreak}{\ifbilingualcolumns\newpage\else\clearpage\fi}
\newcommand{\bisection}[2]{%
  \biheading{\iffirstbilingualsection\global\firstbilingualsectionfalse\else\bisectionpagebreak\fi}{section}{6}{\centering\Large\bfseries}{0.8em}{#1}{#2}%
}
\newcommand{\bisubsection}[2]{%
  \biheading{}{subsection}{5}{\large\bfseries}{0.55em}{#1}{#2}%
}
\newcommand{\bisubsubsection}[2]{%
  \biheading{}{subsubsection}{4}{\normalsize\bfseries}{0.45em}{#1}{#2}%
}
\newcommand{\bicompositeheading}[4]{%
  \biheading{}{subsection}{5}{\large\bfseries}{0.16em}{#1}{#2}%
  \biheading{}{subsubsection}{3}{\normalsize\bfseries}{0.45em}{#3}{#4}%
}
\newif\iflawboxfirstitem
\lawboxfirstitemfalse
\newcommand{\lawboxprespace}[1]{%
  \par
  \iflawboxfirstitem
    \global\lawboxfirstitemfalse
  \else
    \addvspace{#1}%
  \fi
}
\newtcolorbox{lawboxinner}{
  caseboxbase,
  colframe=sourceline!85!black,
  colback=white,
  left=1.05em,
  right=1.05em,
  top=0.62em,
  bottom=0.62em,
  before upper={\global\inlawboxtrue\global\lawboxfirstitemtrue\global\lawboxhaspendingrnfalse\ifintranslation\lawboxtransfont\else\lawboxsourcefont\fi\RaggedRight\setlength{\parindent}{0pt}\setlength{\parskip}{0pt}\ifintranslation\setstretch{\lawboxtransstretch}\else\setstretch{\lawboxsourcestretch}\fi},
  after upper={\global\lawboxhaspendingrnfalse\global\inlawboxfalse}
}
\tcbset{
  lawboxpartbase/.style={
    enhanced,
    breakable=false,
    width=0.985\linewidth,
    colback=white,
    frame hidden,
    boxrule=0pt,
    arc=0pt,
    left=1.05em,
    right=1.05em,
    boxsep=1.5mm,
    before skip=0pt,
    after skip=0pt,
    before upper={\global\inlawboxtrue\global\lawboxfirstitemtrue\global\lawboxhaspendingrnfalse\ifintranslation\lawboxtransfont\else\lawboxsourcefont\fi\RaggedRight\setlength{\parindent}{0pt}\setlength{\parskip}{0pt}\ifintranslation\setstretch{\lawboxtransstretch}\else\setstretch{\lawboxsourcestretch}\fi},
    after upper={\global\lawboxhaspendingrnfalse\global\inlawboxfalse},
    borderline west={0.28pt}{0pt}{sourceline!85!black},
    borderline east={0.28pt}{0pt}{sourceline!85!black},
  },
  lawboxpartfirst/.style={
    lawboxpartbase,
    top=0.62em,
    bottom=0.04em,
    before skip=\lawboxblockbeforeskip,
    borderline north={0.28pt}{0pt}{sourceline!85!black},
  },
  lawboxpartmiddle/.style={
    lawboxpartbase,
    top=0.04em,
    bottom=0.04em,
  },
  lawboxpartlast/.style={
    lawboxpartbase,
    top=0.04em,
    bottom=0.62em,
    after skip=0pt,
    borderline south={0.28pt}{0pt}{sourceline!85!black},
  }
}
\newtcolorbox{lawboxpartinner}[1]{#1}
\newtcolorbox{termboxinner}{
  caseboxbase,
  colframe=noteblue!70!black,
  colback=casepaper,
  before upper={\songfont\setlength{\parindent}{0pt}\setlength{\parskip}{0.35em}}
}
\newtcolorbox{deboxinner}{
  caseboxbase,
  colframe=notegray!72!black,
  colback=white,
  left=1.05em,
  right=1.05em,
  top=0.62em,
  bottom=0.62em,
  before upper={\global\inlawboxtrue\global\lawboxfirstitemtrue\global\lawboxhaspendingrnfalse\small\setlength{\parindent}{0pt}\setlength{\parskip}{0.35em}},
  after upper={\global\lawboxhaspendingrnfalse\global\inlawboxfalse}
}
\NewEnviron{lawbox}{
  \begin{lawboxinner}\BODY\end{lawboxinner}
}
\NewEnviron{lawboxpart}[2]{
  \begin{lawboxpartinner}{lawboxpart#1,equal height group=#2}\BODY\end{lawboxpartinner}%
  \ifstrequal{#1}{last}{\global\lawboxpartendedblocktrue}{}%
}
\NewEnviron{termbox}{
  \par\addvspace{0.55em}
  {\color{noteblue!70!black}\rule{\linewidth}{0.28pt}}\par
  \begingroup
    \songfont\small\color{caseink}
    \setlength{\parindent}{0pt}
    \setlength{\parskip}{0.24em}
    \BODY
    \par
  \endgroup
  {\color{noteblue!70!black}\rule{\linewidth}{0.28pt}}\par
  \addvspace{0.55em}
}
\NewEnviron{debox}{
  \begin{deboxinner}\BODY\end{deboxinner}
}
\providecommand{\paragraphmark}[1]{}
\newcommand{\lawboxtitleafter}{\addvspace{0.06em}}
\newcommand{\lawtitle}[1]{\lawboxprespace{0.22em}{\lawboxtransfont\bfseries\lawboxuserandnumber #1}\par\lawboxtitleafter}
\newcommand{\absno}[2]{\lawboxprespace{0.04em}\noindent\lawboxuserandnumber\hangindent=2.6em\hangafter=1\makebox[2.6em][l]{(#1)}#2\par}
\newcommand{\nrno}[2]{\lawboxprespace{0pt}\noindent\lawboxuserandnumber\hspace*{2.6em}\hangindent=4.4em\hangafter=1\makebox[1.8em][l]{#1.}#2\par}
\newcommand{\letterno}[2]{\lawboxprespace{0pt}\noindent\lawboxuserandnumber\hspace*{4.4em}\hangindent=6.0em\hangafter=1\makebox[1.6em][l]{#1)}#2\par}
\newcommand{\sourcelawtitle}[1]{\lawboxprespace{0.22em}{\lawboxsourcefont\bfseries\lawboxuserandnumber #1}\par\lawboxtitleafter}
\newcommand{\sourcelawsubtitle}[1]{\lawboxprespace{0.04em}{\lawboxsourcefont\bfseries\lawboxuserandnumber #1}\par\lawboxtitleafter}
\newcommand{\sourceabsno}[2]{\lawboxprespace{0.04em}\noindent\lawboxuserandnumber\hangindent=2.45em\hangafter=1\makebox[2.45em][l]{(#1)}#2\par}
\newcommand{\sourcenrno}[2]{\lawboxprespace{0pt}\noindent\lawboxuserandnumber\hspace*{2.45em}\hangindent=4.25em\hangafter=1\makebox[1.8em][l]{#1.}#2\par}
\newcommand{\sourceletterno}[2]{\lawboxprespace{0pt}\noindent\lawboxuserandnumber\hspace*{4.25em}\hangindent=5.85em\hangafter=1\makebox[1.6em][l]{#1)}#2\par}
\newcommand{\sourcequotepara}[1]{\lawboxprespace{0.04em}\noindent\lawboxuserandnumber\hangindent=1.2em\hangafter=1 #1\par}
\newcommand{\sourcelawline}[1]{\lawboxprespace{0.04em}\noindent\lawboxuserandnumber #1\par}
\newcommand{\lawline}[1]{\lawboxprespace{0.04em}\noindent\lawboxuserandnumber #1\par}
\newcommand{\pendingtranslation}{\par{\songfont\footnotesize\color{pendingink}（中文譯文待補。）}\par}
\newcommand{\tocpart}[1]{\par\addvspace{0.52em}{\footnotesize\bfseries\color{pendingink}#1\par}\addvspace{0.16em}}
\newcommand{\tocdashline}{\par\addvspace{0.48em}\noindent{\color{notegray}\xleaders\hbox to 0.82em{\hfil\rule{0.42em}{0.28pt}\hfil}\hskip\linewidth}\par\addvspace{0.38em}}
% \tocpanel{font-cmd}{content}：將目錄置中於所在欄寬（雙語欄適用）。
% A3 橫式使用 0.58\linewidth，A4 橫式使用 0.84\linewidth，
% 使兩種紙張下目錄欄內容實際寬度相近（均約 100–105 mm）。
\newcommand{\tocpanel}[2]{%
  \makebox[\linewidth][c]{%
    \ifx\bilingualpaper\bilingualpaperaThree
      \begin{minipage}{0.58\linewidth}%
    \else
      \begin{minipage}{0.84\linewidth}%
    \fi
    \toccolstyle
    #1#2%
    \end{minipage}%
  }%
}
% ===== 首頁簡目錄的兩層 description list 縮排 =====
% enumitem 的 description list 可把每一列拆成：
%   [label]  labelsep  body text
% 這裡的「起點」都以所在 minipage/欄位的左緣為 0。
%
% 核心規則：
%   labelindent = label 欄位從左緣開始的位置。
%   labelwidth  = label 欄位本身寬度；標籤太長（如 Abw. Meinung）就加大它。
%   labelsep    = label 欄位右緣到正文的距離。
%   leftmargin  = 正文 body text 的起點。判斷層級視覺時主要看這個。
%
% 所以：
%   想讓「Begründung / 判決理由」整列正文往右或往左：改 tocitems 的 leftmargin。
%   想讓「A. Verfahren und Gegenstand / A. 程序與審查標的」正文往右或往左：
%     改 tocsubitems 的 leftmargin。
%   想讓 A./B./C. 這些標籤本身往右或往左，但正文不動：改 tocsubitems 的 labelindent。
%   想讓 A./B./C. 與正文間距變大或變小：改 tocsubitems 的 labelsep。
%
% 層級原則：
%   tocsubitems.leftmargin 必須大於 tocitems.leftmargin，
%   否則子層級正文會看起來比「Gründe / 理由」還靠左。
\newlist{tocitems}{description}{1}
\setlist[tocitems]{%
  labelindent=0pt,    % 第一層 label 欄位起點；0pt 表示貼齊目錄欄左緣。
  leftmargin=2.54cm,  % 第一層正文起點；例如 Begründung / 判決理由 的左緣。
  labelwidth=2.16cm,  % 第一層 label 欄位寬度；需容納 Leitsätze、Abw. Meinung、不同意見等。
  labelsep=0.32cm,    % 第一層 label 與正文的水平間距；只改間距，不改正文起點。
  itemsep=0.28em,     % 第一層各項之間的垂直距離。
  topsep=0.20em,      % 第一層 list 上下與前後內容的垂直距離。
  parsep=0pt,         % 同一 item 內段落之間的距離；目錄項通常不需要。
  partopsep=0pt,      % list 若緊接段落開始時的額外距離；關閉以免目錄高度飄動。
  align=left,         % label 欄內左對齊；長短標籤不靠右貼齊。
  font=\normalfont\bfseries % label 的字型；正文不受這行影響。
}
\newlist{tocsubitems}{description}{1}
\setlist[tocsubitems]{%
  labelindent=1cm,    % 第二層 label 起點；只控制 A./B./C. 本身的位置。
  leftmargin=3.00cm,  % 第二層正文起點；必須大於 2.54cm，才會比 Begründung / 判決理由更右。
  labelwidth=0.5cm,   % 第二層 label 欄寬；A./B./C. 很短，可小於第一層。
  labelsep=1.5cm,    % 第二層 label 與正文間距；A. 與 Verfahren / 程序 的距離看這裡。
  itemsep=0.12em,     % 第二層各項較密，避免 A.–E. 佔太高。
  topsep=0.04em,      % 第二層 list 與 Gründe / 理由之間的垂直距離。
  parsep=0pt,         % 同一第二層 item 內段落距離；目錄項通常不需要。
  partopsep=0pt,      % 關閉額外段落起始距離，避免版面不穩。
  align=left,         % A./B./C. label 欄內左對齊。
  font=\normalfont\bfseries, % A./B./C. label 加粗；正文不受這行影響。
  before={}           % 不在第二層 list 前自動插入額外內容。
}
\newlist{termitems}{description}{1}
\ifbilingualcolumns
  \setlist[termitems]{%
    leftmargin=5.2cm,
    labelwidth=4.8cm,
    labelsep=0.35cm,
    itemsep=0.16em,
    topsep=0.2em,
    parsep=0pt,
    partopsep=0pt,
    font=\normalfont\bfseries,
    style=nextline
  }
\else
  \setlist[termitems]{%
    leftmargin=0pt,
    labelwidth=0pt,
    labelsep=0pt,
    itemsep=0.24em,
    topsep=0.2em,
    parsep=0pt,
    partopsep=0pt,
    font=\normalfont\bfseries,
    style=nextline
  }
\fi
\newlist{abbritems}{description}{1}
\setlist[abbritems]{%
  leftmargin=2.38cm,
  labelwidth=2.02cm,
  labelsep=0.28cm,
  itemsep=0.16em,
  topsep=0.2em,
  parsep=0pt,
  partopsep=0pt,
  align=left,
  font=\normalfont\bfseries
}
\ifbilingualcolumns
  \newenvironment{termcolumns}{\begin{multicols}{2}}{\end{multicols}}
\else
  \newenvironment{termcolumns}{}{}
\fi

% ===== 編者註腳開關 =====
% 一般版預設不顯示編者註，避免正文腳註過密。
% 若開啟 \classroommodeon，GG/條文編者註仍會自動顯示，無需另開本開關。
% 若一般版也要顯示編者註，可在單案檔 \input 後加 \showeditornotestrue。
\newif\ifshoweditornotes
\showeditornotesfalse

% ===== 目錄排版輔助 =====
\newcommand{\toccolstyle}{\raggedright\arraybackslash}
\newcommand{\tocsingleindent}{\leftskip=1.45cm\rightskip=1.2cm}

% ===== 行內引文環境（德文原文與中文譯文各一，帶左側灰色邊線）=====
\newcommand{\quoteparbreak}{\par\addvspace{0.48em}}
\newcommand{\sourcequoteinner}[1]{%
  \begin{tcolorbox}[
    enhanced, breakable, width=\dimexpr\linewidth-0.8em\relax, frame hidden,
    colback=white, coltext=caseink, boxrule=0pt,
    borderline west={0.55pt}{0.88em}{notegray},
    left=1.78em, right=0.72em, top=0.18em, bottom=0.18em,
    boxsep=0pt, before skip=0.24em, after skip=0.24em,
    before upper={\normalfont\footnotesize\color{caseink}\setlength{\parindent}{0pt}\setlength{\parskip}{0.34em}\raggedright\setlength{\rightskip}{0pt plus 1em}}
  ]%
  #1\par
  \end{tcolorbox}%
}
\newcommand{\transquoteinner}[1]{%
  \begin{tcolorbox}[
    enhanced, breakable, width=\dimexpr\linewidth-0.8em\relax, frame hidden,
    colback=white, coltext=caseink, boxrule=0pt,
    borderline west={0.55pt}{0.88em}{notegray},
    left=1.78em, right=0.72em, top=0.18em, bottom=0.18em,
    boxsep=0pt, before skip=0.24em, after skip=0.24em,
    before upper={\songfont\small\color{caseink}\setlength{\parindent}{0pt}\setlength{\parskip}{0.34em}\raggedright\setlength{\rightskip}{0pt plus 1em}}
  ]%
  #1\par
  \end{tcolorbox}%
}

% ===== 大綱（Gliederung）Rn 輔助命令 =====
\newcommand{\outrn}[2]{\enspace{\normalfont\footnotesize\color{pendingink}(Rn\,#1–#2)}}
\newcommand{\outrnsingle}[1]{\enspace{\normalfont\footnotesize\color{pendingink}(Rn\,#1)}}
\newcommand{\outrnzh}[2]{\enspace{\normalfont\footnotesize\color{pendingink}（第\,#1–#2\,段）}}
\newcommand{\outrnzhs}[1]{\enspace{\normalfont\footnotesize\color{pendingink}（第\,#1\,段）}}
% ===== 大綱雙語對照表格輔助命令（三欄：段碼 │ 德文 │ 中文）=====
% 欄 1：段碼（由欄定義自動格式化：小字、等寬、右對齊、pendingink）
% 欄 2：德文標籤＋正文
% 欄 3：中文標籤＋正文
%
% 無段碼的列（\omainrow、\osecrow），欄 1 以 & 空置。
% 有段碼的列（\osecrowrn、\osubrow、\ointrow），#1 為預格式化字串
%   例：\osubrow{2–49}{I.}{...}{I.}{...}
%       \ointrow{107}{...}{...}
%
% \opartrow：段組標題（橫跨三欄）
\newcommand{\opartrow}[2]{%
  \noalign{\vspace{0.30em}}%
  \multicolumn{3}{@{}l@{}}{%
    \footnotesize\bfseries\color{pendingink}#1\hfill\songfont#2%
  }\\[0.06em]%
}
% \odashrow：虛線分隔（橫跨三欄）
\newcommand{\odashrow}{%
  \noalign{\vspace{0.28em}}%
  \multicolumn{3}{@{}l@{}}{\color{notegray}%
    \xleaders\hbox to 0.82em{\hfil\rule{0.42em}{0.28pt}\hfil}\hskip 0.88\linewidth}%
  \\[0.24em]%
}
% \odissentpart：不同意見書區段標頭。比一般 \opartrow 更像附錄性轉場。
\newcommand{\odissentpart}[2]{%
  \noalign{\vspace{0.42em}}%
  \multicolumn{3}{@{}p{0.88\textwidth}@{}}{%
    \color{notegray}\rule{\linewidth}{0.18pt}%
  }\\[-0.18em]%
  \multicolumn{3}{@{}p{0.88\textwidth}@{}}{%
    \makebox[\linewidth][s]{%
      \footnotesize\bfseries\color{pendingink}#1%
      \hfill{\songfont #2}%
    }%
  }\\[-0.02em]%
}
% \odissentopinion：不同意見書的署名與一行摘要，右欄保留段碼範圍。
\newcommand{\odissentopinion}[5]{%
  \footnotesize\hangindent=0.52cm\hangafter=1%
    {\bfseries\color{caseink}#2}\enspace{\color{notegray}\textperiodcentered}\enspace\color{caseink}#3%
  &\footnotesize\hangindent=0.52cm\hangafter=1%
    {\songfont\bfseries\color{caseink}#4}\enspace{\color{notegray}\textperiodcentered}\enspace\songfont\color{caseink}#5%
  &#1%
  \\[0.12em]%
}
% \omainrow：頂層項目，無段碼（Leitsätze、Urteil、Gründe …）
\newcommand{\omainrow}[4]{%
  \footnotesize\hangindent=1.92cm\hangafter=1%
    \makebox[1.92cm][l]{\bfseries\color{caseink}#1}\color{caseink}#2%
  &\footnotesize\hangindent=1.92cm\hangafter=1%
    \makebox[1.92cm][l]{\bfseries\color{caseink}#3}\songfont\color{caseink}#4%
  &%
  \\[0.15em]%
}
% \omainrowrn：頂層項目，有段碼（不同意見等標籤寬於 \osecrow 框者）
\newcommand{\omainrowrn}[5]{%
  \footnotesize\hangindent=1.92cm\hangafter=1%
    \makebox[1.92cm][l]{\bfseries\color{caseink}#2}\color{caseink}#3%
  &\footnotesize\hangindent=1.92cm\hangafter=1%
    \makebox[1.92cm][l]{\bfseries\color{caseink}#4}\songfont\color{caseink}#5%
  &#1%
  \\[0.15em]%
}
% \osecrow：一級子項，無段碼（A.、C.、D. 等有子項者）
\newcommand{\osecrow}[4]{%
  \footnotesize\hspace{1.10cm}\hangindent=1.98cm\hangafter=1%
    \makebox[0.86cm][l]{\bfseries\color{caseink}#1}\color{caseink}#2%
  &\footnotesize\hspace{1.10cm}\hangindent=1.98cm\hangafter=1%
    \makebox[0.86cm][l]{\bfseries\color{caseink}#3}\songfont\color{caseink}#4%
  &%
  \\[0.11em]%
}
% \osecrowrn：一級子項，有段碼（B.、E. 等完整段落範圍者）
% #1=段碼字串（如 101–106）, #2=DE標籤, #3=DE正文, #4=ZH標籤, #5=ZH正文
\newcommand{\osecrowrn}[5]{%
  \footnotesize\hspace{1.10cm}\hangindent=1.98cm\hangafter=1%
    \makebox[0.86cm][l]{\bfseries\color{caseink}#2}\color{caseink}#3%
  &\footnotesize\hspace{1.10cm}\hangindent=1.98cm\hangafter=1%
    \makebox[0.86cm][l]{\bfseries\color{caseink}#4}\songfont\color{caseink}#5%
  &#1%
  \\[0.11em]%
}
% \osubrow：二級子項，有段碼範圍（I.、II.、A.I. …）
% #1=段碼字串, #2=DE標籤, #3=DE正文, #4=ZH標籤, #5=ZH正文
\newcommand{\osubrow}[5]{%
  \footnotesize\hspace{1.40cm}\hangindent=2.30cm\hangafter=1%
    \makebox[0.90cm][l]{\bfseries\color{caseink}#2}\color{caseink}#3%
  &\footnotesize\hspace{1.40cm}\hangindent=2.30cm\hangafter=1%
    \makebox[0.90cm][l]{\bfseries\color{caseink}#4}\songfont\color{caseink}#5%
  &#1%
  \\[0.09em]%
}
% \ointrow：引入段落，有單一段碼，無標籤
% #1=段碼字串, #2=DE正文, #3=ZH正文
\newcommand{\ointrow}[3]{%
  \footnotesize\hspace{0.52cm}\hangindent=0.52cm\hangafter=1\color{caseink}#2%
  &\footnotesize\hspace{0.52cm}\hangindent=0.52cm\hangafter=1\songfont\color{caseink}#3%
  &#1%
  \\[0.09em]%
}
% \osubsubrow：三級子項（1.、2.、3.）
% #1=段碼字串, #2=DE標籤, #3=DE正文, #4=ZH標籤, #5=ZH正文
\newcommand{\osubsubrow}[5]{%
  \footnotesize\hspace{1.80cm}\hangindent=2.48cm\hangafter=1%
    \makebox[0.68cm][l]{\bfseries\color{caseink}#2}\color{caseink}#3%
  &\footnotesize\hspace{1.80cm}\hangindent=2.48cm\hangafter=1%
    \makebox[0.68cm][l]{\bfseries\color{caseink}#4}\songfont\color{caseink}#5%
  &#1%
  \\[0.08em]%
}
% \osubsubsubrow：四級子項（a）、b））
% #1=段碼字串, #2=DE標籤, #3=DE正文, #4=ZH標籤, #5=ZH正文
\newcommand{\osubsubsubrow}[5]{%
  \footnotesize\hspace{2.10cm}\hangindent=2.68cm\hangafter=1%
    \makebox[0.58cm][l]{\bfseries\color{caseink}#2}\color{caseink}#3%
  &\footnotesize\hspace{2.10cm}\hangindent=2.68cm\hangafter=1%
    \makebox[0.58cm][l]{\bfseries\color{caseink}#4}\songfont\color{caseink}#5%
  &#1%
  \\[0.07em]%
}
% \osubsubsubsubrow：五級子項（aa）、bb））
% 標籤框 0.62cm：足以容納「aa)」在 footnotesize bold 下（約 0.50cm）並留有間距
% #1=段碼字串, #2=DE標籤, #3=DE正文, #4=ZH標籤, #5=ZH正文
\newcommand{\osubsubsubsubrow}[5]{%
  \footnotesize\hspace{2.38cm}\hangindent=3.06cm\hangafter=1%
    \makebox[0.68cm][l]{\bfseries\color{caseink}#2}\color{caseink}#3%
  &\footnotesize\hspace{2.38cm}\hangindent=3.06cm\hangafter=1%
    \makebox[0.68cm][l]{\bfseries\color{caseink}#4}\songfont\color{caseink}#5%
  &#1%
  \\[0.06em]%
}
% \outlinepage：雙語對照大綱頁包裝環境（longtable 自動跨頁）
% 三欄：德文欄 + 中文欄 + 段碼欄（最右，不折行）
% 段碼欄寬度：1.80cm，足以容納最長段碼如「309–348」
% DE/ZH 欄寬：(0.87\textwidth - 2.08cm) / 2
%   驗算：2×欄 + 0.05\textwidth + 0.28cm gap + 1.80cm = 0.87tw - 2.08cm + 0.05tw + 2.08cm = 0.92tw ✓
\newenvironment{outlinepage}{%
  \vspace*{0.40cm}%
  \setlength{\LTleft}{0.04\textwidth}%
  \setlength{\LTright}{0.04\textwidth}%
  \setlength{\tabcolsep}{0pt}%
  % 大綱列距由 longtable 的 \arraystretch 與各 row macro 結尾的 \\[...em] 共同決定。
  % 這裡控制「每一列本身的表格 strut 高度」；數值越大，A./I./1. 各項之間越像有空行。
  % A3 原本用 0.62，視覺上沒有空行；A4 也沿用同值，避免切到 A4 後項目被撐開。
  % 若日後覺得大綱太擠，只微調這個值（例如 0.68），不要改各 \omainrow/\osubrow 的縮排。
  \renewcommand{\arraystretch}{0.62}%
  \begin{longtable}{%
    >{\vspace{0pt}\raggedright\arraybackslash}p{\dimexpr(0.87\textwidth-2.08cm)/2\relax}%
    @{\hspace{0.025\textwidth}}%
    !{\color{notegray}\vrule width 0.12pt}%
    @{\hspace{0.025\textwidth}}%
    >{\vspace{0pt}\raggedright\arraybackslash}p{\dimexpr(0.87\textwidth-2.08cm)/2\relax}%
    @{\hspace{0.28cm}}%
    >{\vspace{0pt}\footnotesize\normalfont\color{pendingink}\raggedright\arraybackslash}p{1.80cm}%
    @{}%
  }%
  % 首頁欄標籤：不要橫跨置中；分別放在德文欄與中文欄的實際欄位起點。
  \footnotesize\bfseries\color{sourceink}Gliederung%
  &\footnotesize\bfseries\songfont\color{sourceink}大綱%
  &%
  \\[0.36em]%
  \endfirsthead%
  % 續頁欄標籤：同樣分欄定位，以灰色細體區分；無需標注「續」。
  \footnotesize\color{pendingink}Gliederung%
  &\footnotesize\songfont\color{pendingink}大綱%
  &%
  \\[0.24em]%
  \endhead%
}{%
  \end{longtable}%
}
 之前設定；
%   預設 chinese / a4）
\ifdefined\docmode\else\def\docmode{chinese}\fi
\ifdefined\bilingualpaper\else\def\bilingualpaper{a4}\fi
\newif\ifshoworiginal
\newif\ifshowtranslation
\newif\ifbilinguallandscape
\newif\ifbilingualcolumns
\newif\ifintranslation
\newif\iffirstbilingualsection

\showoriginalfalse
\showtranslationfalse
\bilinguallandscapefalse
\bilingualcolumnsfalse
\intranslationfalse
\firstbilingualsectionfalse

\def\modechinese{chinese}
\def\modegerman{german}
\def\modebilingual{bilingual}
\def\bilingualpaperaThree{a3}
\def\bilingualpaperaFour{a4}

\ifx\docmode\modechinese
  \showtranslationtrue
\fi

\ifx\docmode\modegerman
  \showoriginaltrue
\fi

\ifx\docmode\modebilingual
  \showoriginaltrue
  \showtranslationtrue
  \bilinguallandscapetrue
  \bilingualcolumnstrue
\fi

% 編排模式說明：
% chinese  → \showoriginalfalse  \showtranslationtrue  （A4 直式，純中文）
% german   → \showoriginaltrue   \showtranslationfalse （A4 直式，純德文）
% bilingual→ \showoriginaltrue   \showtranslationtrue  \bilingualcolumnstrue（A3/A4 橫式對照）

\ifbilingualcolumns
  \ifx\bilingualpaper\bilingualpaperaFour
    \geometry{
      a4paper,
      landscape,
      left=1.25cm,
      right=2.25cm,
      top=1.25cm,
      bottom=1.75cm,
      footskip=8.5mm,
      marginparwidth=0.75cm,
      marginparsep=0.25cm
    }
  \else
    \geometry{
      a3paper,
      landscape,
      left=1.55cm,
      right=2.75cm,
      top=1.55cm,
      bottom=2.05cm,
      footskip=9.5mm,
      marginparwidth=0.9cm,
      marginparsep=0.35cm
    }
  \fi
\else
\geometry{
  a4paper,
  portrait,
  left=2.2cm,
  right=2.6cm,
  top=2.0cm,
  bottom=2.0cm,
  marginparwidth=0.8cm,
  marginparsep=0.3cm
}
\fi
% ===== 時間戳基礎設施 =====
\newcount\stamptimeval
\newcount\stampminuteval
\newcommand{\stamphh}{%
  \stamptimeval=\time
  \divide\stamptimeval by 60
  \ifnum\stamptimeval<10 0\fi
  \the\stamptimeval}
\newcommand{\stampmm}{%
  \stamptimeval=\time
  \stampminuteval=\time
  \divide\stampminuteval by 60
  \multiply\stampminuteval by 60
  \advance\stamptimeval by -\stampminuteval
  \ifnum\stamptimeval<10 0\fi
  \the\stamptimeval}
\newcommand{\stampwkde}[1]{%
  \ifcase#1Montag\or Dienstag\or Mittwoch\or Donnerstag\or Freitag\or Samstag\or Sonntag\fi}
\newcommand{\stampwkzh}[1]{%
  \ifcase#1 一\or 二\or 三\or 四\or 五\or 六\or 日\fi}
\newcommand{\stampmonthde}[1]{%
  \ifcase#1\or Januar\or Februar\or März\or April\or Mai\or Juni\or
  Juli\or August\or September\or Oktober\or November\or Dezember\fi}
\newcount\stampjdval
\newcount\stampwdval
\pgfcalendardatetojulian{\the\year-\the\month-\the\day}{\stampjdval}
\pgfcalendarjuliantoweekday{\the\stampjdval}{\stampwdval}
\newcommand{\timestampzh}{%
  \songfont 最後修改：\the\year 年\the\month 月\the\day 日%
  % （星期\stampwkzh{\the\stampwdval}）%
  % \space\stamphh:\stampmm
  }

\newcommand{\documentdatezh}{\the\year 年 \the\month 月 \the\day 日}
\newcommand{\documentdatede}{\the\day. \stampmonthde{\the\month} \the\year}
\newcommand{\bverfgetranslationnote}{%
  {\songfont 翻譯說明：初譯使用 Claude Code 與 GPT 協助迭代；仍請以德文原文為準。}%
}
\newcommand{\bverfgecourseinfo}{%
  {\songfont 課程：114 學年度第 2 學期；憲法專題研究（四）；授課老師：湯德宗教授；導讀日期：115 年 6 月 1 日。}%
}
\newcommand{\bverfgeeditorinfo}{%
  {\songfont 編者：王逸帆（R10A21126），國立台灣大學法律系研究所碩士班。}%
}
\newcommand{\bverfgetitleversionline}{%
  \ifbilingualcolumns
    Fassung \documentversion\enspace\textperiodcentered\enspace Stand: \documentdatede
    \quad
    {\songfont 版本 \documentversion\enspace\textperiodcentered\enspace 修訂日期：\documentdatezh\par}%
  \else
    \ifshowtranslation
      {\songfont 版本 \documentversion\enspace\textperiodcentered\enspace 修訂日期：\documentdatezh\par}%
    \else
      Fassung \documentversion\enspace\textperiodcentered\enspace Stand: \documentdatede\par
    \fi
  \fi
}
\newcommand{\bverfgetitlecornerinfo}{%
  \ifshowtranslation
    \vspace{0.72em}%
    \noindent\hfill
    \begin{minipage}{0.66\textwidth}
      \raggedleft\tiny\color{pendingink}\songfont
      \bverfgecourseinfo\par
      \bverfgeeditorinfo
    \end{minipage}%
  \fi
}
\newcommand{\bverfgetitlemeta}{%
  \vfill
  \begin{center}%
    \footnotesize\color{pendingink}%
    \bverfgetitleversionline
    \ifshowtranslation
      \vspace{0.28em}{\scriptsize\bverfgetranslationnote\par}%
    \fi
  \end{center}%
  \bverfgetitlecornerinfo
}

\pagestyle{fancy}
\fancyhf{}
\renewcommand{\headrulewidth}{0pt}
\renewcommand{\footrulewidth}{0pt}
\fancyfoot[C]{\thepage}
\ifbilingualcolumns
  \fancyfoot[R]{\scriptsize\color{pendingink}\timestampzh}%
\else
  \ifshoworiginal
  \else
    \fancyfoot[R]{\scriptsize\color{pendingink}\timestampzh\hspace{0.45cm}}%
  \fi
\fi
\setmainfont{Times New Roman}
\setsansfont{Helvetica Neue}
\setmonofont{Menlo}
\IfFileExists{/Users/iw/Library/Fonts/kaiu.ttf}{
  \setCJKmainfont[Path=/Users/iw/Library/Fonts/,AutoFakeBold=4]{kaiu.ttf}
}{
  \IfFontExistsTF{標楷體}{
    \setCJKmainfont[AutoFakeBold=4]{標楷體}
  }{
    \setCJKmainfont{Songti TC}
  }
}
\IfFontExistsTF{Songti TC}{
  \setCJKfamilyfont{song}{Songti TC}
  \setCJKmonofont{Songti TC}
}{
  \setCJKfamilyfont{song}[Path=/Users/iw/Library/Fonts/,AutoFakeBold=4]{kaiu.ttf}
  \setCJKmonofont[Path=/Users/iw/Library/Fonts/,AutoFakeBold=4]{kaiu.ttf}
}
\newcommand{\songfont}{\CJKfamily{song}}

% ===== 封面／目錄專用打字機字體（僅限封面、目錄頁使用，不影響正文）=====
\IfFontExistsTF{Radio Newsman}{%
  \newfontfamily\bverfgetitlede{Radio Newsman}%
}{%
  \newfontfamily\bverfgetitlede{Courier New}%
}
\IfFontExistsTF{Erikas Farbband}{%
  \newfontfamily\bverfgebodyde{Erikas Farbband}%
}{%
  \let\bverfgebodyde\bverfgetitlede
}
\IfFontExistsTF{Maoken Heavy Labourer}{%
  \setCJKfamilyfont{bverfgetitlecjk}{Maoken Heavy Labourer}%
}{%
  \setCJKfamilyfont{bverfgetitlecjk}{Songti TC}%
}
\IfFontExistsTF{Huiwen-mincho}{%
  \setCJKfamilyfont{bverfgebodycjk}{Huiwen-mincho}%
}{%
  \setCJKfamilyfont{bverfgebodycjk}{Songti TC}%
}
\newcommand{\bverfgetitlezh}{\bverfgetitlede\CJKfamily{bverfgetitlecjk}}
\newcommand{\bverfgebodyzh}{\bverfgebodyde\CJKfamily{bverfgebodycjk}}

\setstretch{1.35}
\setlist{itemsep=0.2em, topsep=0.4em}
\setlength{\parindent}{0pt}
\setlength{\parskip}{0.45em}
\raggedbottom
\setcounter{secnumdepth}{0}
\setcounter{tocdepth}{2}

\newcommand{\lawboxsourcefont}{\footnotesize}
\newcommand{\lawboxtransfont}{\footnotesize}
\newcommand{\lawboxsourcestretch}{1.02}
\newcommand{\lawboxtransstretch}{0.92}

\titleformat{\section}{\centering\Large\bfseries\songfont}{\thesection}{1em}{}
\titleformat{\subsection}{\large\bfseries\songfont}{\thesubsection}{1em}{}
\titleformat{\subsubsection}{\normalsize\bfseries\songfont}{\thesubsubsection}{1em}{}
\titleformat{\paragraph}{\normalsize\bfseries\songfont}{\theparagraph}{1em}{}
\titlespacing*{\section}{0pt}{1.8em}{1.0em}
\titlespacing*{\subsection}{0pt}{1.2em}{0.6em}
\titlespacing*{\subsubsection}{0pt}{1.0em}{0.45em}
\titlespacing*{\paragraph}{0pt}{0.6em}{0.4em}

\definecolor{noteblue}{HTML}{A9C9F5}
\definecolor{noteteal}{HTML}{AEE1D6}
\definecolor{notegray}{HTML}{D7DADF}
\definecolor{caseink}{HTML}{000000}
\definecolor{casepaper}{HTML}{FCFEFD}
\definecolor{sourcepaper}{HTML}{FAFBF8}
\definecolor{sourceline}{HTML}{D7DED8}
\definecolor{sourceink}{HTML}{000000}
\definecolor{pendingink}{HTML}{000000}

\tcbset{
  caseboxbase/.style={
    enhanced,
    breakable,
    width=0.985\linewidth,
    colback=casepaper,
    coltitle=caseink,
    fonttitle=\bfseries\songfont,
    boxrule=0.28pt,
    arc=1.2mm,
    left=2.4mm,
    right=2.4mm,
    top=1.5mm,
    bottom=1.5mm,
    boxsep=1.5mm,
    before skip=0.65em,
    after skip=0.65em,
    before upper={\setlength{\parindent}{0pt}\setlength{\parskip}{0.35em}},
  }
}

\newcommand{\bilingualstart}{}
\newcommand{\bilingualstop}{}
\ifbilingualcolumns
  \renewcommand{\bilingualstart}{\small\setlength{\columnsep}{1.25cm}\setlength{\columnseprule}{0.12pt}\columnratio{0.5,0.5}\multicolumnfootnotes\flushbottom\sloppy\interlinepenalty=80\clubpenalty=800\widowpenalty=800\displaywidowpenalty=800\begin{paracol}{2}\global\intranslationfalse\global\firstbilingualsectiontrue{\footnotesize\color{sourceink}Deutsch\par}\switchcolumn{\footnotesize\songfont\color{sourceink}中文譯文\par}\switchcolumn*}
  \renewcommand{\bilingualstop}{\ifintranslation\switchcolumn*\global\intranslationfalse\fi\end{paracol}\clearpage\raggedbottom}
\fi
\newif\ifrandnumbers
\randnumbersfalse
\newcounter{randnumber}
\newcounter{randnumberlocal}
\newcounter{lastsourcerandstart}
\newif\iflastsourcehasrandnumbers
\lastsourcehasrandnumbersfalse
\newif\ifinlawbox
\inlawboxfalse
\newif\iflawboxhaspendingrn
\lawboxhaspendingrnfalse
\newcommand{\lawboxpendingrn}{}
\newcommand{\startrandnumbers}{\global\randnumberstrue\setcounter{randnumber}{1}\global\lastsourcehasrandnumbersfalse}
\newcommand{\stoprandnumbers}{\global\randnumbersfalse\global\lastsourcehasrandnumbersfalse}
\newlength{\randnumberwidth}
\setlength{\randnumberwidth}{0.95cm}
\newlength{\randnumberrightoffset}
\setlength{\randnumberrightoffset}{0.85cm}
\newlength{\randnumberlawboxrightoffset}
\setlength{\randnumberlawboxrightoffset}{1.40cm}
\newcommand{\randnumberbox}[1]{%
  \leavevmode
  \makebox[0pt][l]{%
    \smash{%
      \tikz[remember picture,overlay,baseline=(rnmark.base)]{%
        \coordinate (rnmark) at (0,0);%
        \ifinlawbox
          \path let \p1=(current page.east), \p2=(rnmark.base) in
            node[
              anchor=base east,
              inner sep=0pt,
              outer sep=0pt,
              text width=\randnumberwidth,
              align=right,
              text=pendingink,
              font=\normalfont\mdseries\scriptsize\ttfamily
            ] at (\x1-\randnumberlawboxrightoffset,\y2) {{\normalfont\mdseries\scriptsize\ttfamily #1}};%
        \else
          \path let \p1=(current page.east), \p2=(rnmark.base) in
            node[
              anchor=base east,
              inner sep=0pt,
              outer sep=0pt,
              text width=\randnumberwidth,
              align=right,
              text=pendingink,
              font=\normalfont\mdseries\scriptsize\ttfamily
            ] at (\x1-\randnumberrightoffset,\y2) {{\normalfont\mdseries\scriptsize\ttfamily #1}};%
        \fi
      }%
    }%
  }%
  \ignorespaces
}
\newcommand{\lawboxstorerandnumber}[1]{%
  \gdef\lawboxpendingrn{#1}%
  \global\lawboxhaspendingrntrue
  \ignorespaces
}
\newcommand{\lawboxuserandnumber}{%
  \iflawboxhaspendingrn
    \randnumberbox{\lawboxpendingrn}%
    \global\lawboxhaspendingrnfalse
  \fi
}
\newcommand{\sourcern}[1]{%
  \ifrandnumbers
    \ifshoworiginal
      \ifshowtranslation\else
        \ifinlawbox\lawboxstorerandnumber{#1}\else\randnumberbox{#1}\fi
      \fi
    \fi
  \fi
}
\newcommand{\transrn}[1]{%
  \ifrandnumbers
    \ifshowtranslation
      \ifinlawbox\lawboxstorerandnumber{#1}\else\randnumberbox{#1}\fi
    \fi
  \fi
}
% ===== 課堂模式：德文原文縮寫註解 =====
% 日常切換只改各判決檔的一行：
%   \classroommodeon        開啟課堂版；註解掉這行即回到一般版。
% 模板預設：
%   1. 課堂模式關閉。
%   2. 課堂模式開啟時，縮寫註解包含「德文全稱＋中文譯名」。
%   3. 課堂模式開啟時，GG/條文編者註仍會自動顯示。
% 進階微調才需要在單案檔覆寫：
%   \classroomabbrzhoff     縮寫註解僅顯示德文全稱。
%   \showeditornotestrue    一般版也顯示編者註。
% 註解策略：同一實際 PDF 頁面中，同一縮寫只在第一次出現處加註。
% 這裡用 zref-abspage 的絕對頁序，不用 \thepage，避免文件內頁碼重置時誤判。
\newif\ifclassroommode
\newif\ifclassroomabbrzh
\classroommodefalse
\classroomabbrzhtrue
\newcommand{\classroommodeon}{\classroommodetrue}
\newcommand{\classroommodeoff}{\classroommodefalse}
\newcommand{\classroomabbrzhon}{\classroomabbrzhtrue}
\newcommand{\classroomabbrzhoff}{\classroomabbrzhfalse}
\newcommand{\abbrnotelabel}{\emph{Abk.}: }
\newcommand{\abbrnote}[3]{%
  \ifclassroommode
    \ifshoworiginal
      \ifintranslation\else
        \footnote{\normalfont\footnotesize\abbrnotelabel #2\ifclassroomabbrzh；{\songfont #3}\fi}%
      \fi
    \fi
  \fi
}
\newcommand{\abbrnoteonce}[4]{%
  #2%
  \ifclassroommode
    \ifshoworiginal
      \ifintranslation\else
        \begingroup
          \edef\abbrcurrentabspage{\number\numexpr\value{abspage}+1\relax}%
          \ifcsname abbrnoteseen@#1@\abbrcurrentabspage\endcsname\else
            \global\expandafter\let\csname abbrnoteseen@#1@\abbrcurrentabspage\endcsname\empty
            \abbrnote{#1}{#3}{#4}%
          \fi
        \endgroup
      \fi
    \fi
  \fi
}
\newcommand{\abbrAbs}{\abbrnoteonce{abs}{Abs.}{Absatz}{項}}
\newcommand{\abbrAAO}{\abbrnoteonce{aao}{a.a.O.}{am angegebenen Ort}{前引處}}
\newcommand{\abbrAE}{\abbrnoteonce{ae}{AE}{Alternativ-Entwurf}{替代草案}}
\newcommand{\abbrArt}{\abbrnoteonce{art}{Art.}{Artikel}{條}}
\newcommand{\abbrAufl}{\abbrnoteonce{aufl}{Aufl.}{Auflage}{版}}
\newcommand{\abbrBd}{\abbrnoteonce{bd}{Bd.}{Band}{卷}}
\newcommand{\abbrBGBl}{\abbrnoteonce{bgbl}{BGBl.}{Bundesgesetzblatt}{聯邦法律公報}}
\newcommand{\abbrBGHSt}{\abbrnoteonce{bghst}{BGHSt}{Entscheidungen des Bundesgerichtshofes in Strafsachen}{聯邦普通法院刑事判決彙編}}
\newcommand{\abbrBRDrucks}{\abbrnoteonce{brdrucks}{BRDrucks.}{Bundesrats-Drucksache}{聯邦參議院文件}}
\newcommand{\abbrBTDrucks}{\abbrnoteonce{btdrucks}{BTDrucks.}{Bundestags-Drucksache}{聯邦議院文件}}
\newcommand{\abbrBundesgesetzbl}{\abbrnoteonce{bundesgesetzbl}{Bundesgesetzbl.}{Bundesgesetzblatt}{聯邦法律公報}}
\newcommand{\abbrBvF}{\abbrnoteonce{bvf}{BvF}{Bundesverfassungsgerichtsverfahren}{聯邦憲法法院程序案號}}
\newcommand{\abbrBvL}{\abbrnoteonce{bvl}{BvL}{Bundesverfassungsgerichtsverfahren}{聯邦憲法法院程序案號}}
\newcommand{\abbrBvR}{\abbrnoteonce{bvr}{BvR}{Bundesverfassungsgerichtsverfahren}{聯邦憲法法院程序案號}}
\newcommand{\abbrBVerfG}{\abbrnoteonce{bverfg}{BVerfG}{Bundesverfassungsgericht}{聯邦憲法法院}}
\newcommand{\abbrBVerfGE}{\abbrnoteonce{bverfge}{BVerfGE}{Entscheidungen des Bundesverfassungsgerichts}{聯邦憲法法院判決集}}
\newcommand{\abbrBVerfGG}{\abbrnoteonce{bverfgg}{BVerfGG}{Bundesverfassungsgerichtsgesetz}{聯邦憲法法院法}}
\newcommand{\abbrDDR}{\abbrnoteonce{ddr}{DDR}{Deutsche Demokratische Republik}{德意志民主共和國；東德}}
\newcommand{\abbrDJT}{\abbrnoteonce{djt}{DJT}{Deutscher Juristentag}{德國法學家大會}}
\newcommand{\abbrDr}{\abbrnoteonce{dr}{Dr.}{Doktor}{Dr.}}
\newcommand{\abbrEuGRZ}{\abbrnoteonce{eugrz}{EuGRZ}{Europäische Grundrechte-Zeitschrift}{歐洲基本權期刊}}
\newcommand{\abbrEV}{\abbrnoteonce{ev}{EV}{Einigungsvertrag}{統一條約}}
\newcommand{\abbrF}{\abbrnoteonce{f}{f.}{folgende}{以下一頁；下一段}}
\newcommand{\abbrFF}{\abbrnoteonce{ff}{ff.}{fortfolgende}{以下數頁；以下各段}}
\newcommand{\abbrGBlDDR}{\abbrnoteonce{gblddr}{GBl. DDR}{Gesetzblatt der Deutschen Demokratischen Republik}{德意志民主共和國法律公報}}
\newcommand{\abbrGG}{\abbrnoteonce{gg}{GG}{Grundgesetz}{基本法}}
\newcommand{\abbrGVBl}{\abbrnoteonce{gvbl}{GVBl.}{Gesetz- und Verordnungsblatt}{法律及命令公報}}
\newcommand{\abbrJR}{\abbrnoteonce{jr}{JR}{Juristische Rundschau}{法學評論}}
\newcommand{\abbrJZ}{\abbrnoteonce{jz}{JZ}{JuristenZeitung}{法學家報}}
\newcommand{\abbrIVm}{\abbrnoteonce{ivm}{i.V.m.}{in Verbindung mit}{結合；連同}}
\newcommand{\abbrMdB}{\abbrnoteonce{mdb}{MdB}{Mitglied des Bundestages}{聯邦議院議員}}
\newcommand{\abbrMwN}{\abbrnoteonce{mwn}{m.w.N.}{mit weiteren Nachweisen}{附進一步文獻／判例出處}}
\newcommand{\abbrNF}{\abbrnoteonce{nf}{n.F.}{neue Fassung}{新版本}}
\newcommand{\abbrAF}{\abbrnoteonce{af}{a.F.}{alte Fassung}{舊版本}}
\newcommand{\abbrNr}{\abbrnoteonce{nr}{Nr.}{Nummer}{款、目或號，依語境}}
\newcommand{\abbrProf}{\abbrnoteonce{prof}{Prof.}{Professor}{教授}}
\newcommand{\abbrRdnr}{\abbrnoteonce{rdnr}{Rdnr.}{Randnummer}{邊碼；段碼}}
\newcommand{\abbrRdnrn}{\abbrnoteonce{rdnrn}{Rdnrn.}{Randnummern}{邊碼；段碼}}
\newcommand{\abbrRGBl}{\abbrnoteonce{rgbl}{RGBl.}{Reichsgesetzblatt}{帝國法律公報}}
\newcommand{\abbrRGSt}{\abbrnoteonce{rgst}{RGSt}{Entscheidungen des Reichsgerichts in Strafsachen}{帝國法院刑事判決彙編}}
\newcommand{\abbrRspr}{\abbrnoteonce{rspr}{Rspr.}{Rechtsprechung}{判例；司法實務}}
\newcommand{\abbrRVO}{\abbrnoteonce{rvo}{RVO}{Reichsversicherungsordnung}{帝國保險法}}
\newcommand{\abbrS}{\abbrnoteonce{s}{S.}{Seite}{頁}}
\newcommand{\abbrSFHG}{\abbrnoteonce{sfhg}{SFHG}{Schwangeren- und Familienhilfegesetz}{孕婦及家庭扶助法}}
\newcommand{\abbrSGBV}{\abbrnoteonce{sgbv}{SGB V}{Fünftes Buch Sozialgesetzbuch}{社會法典第五編}}
\newcommand{\abbrStAG}{\abbrnoteonce{staeg}{StÄG}{Strafrechtsänderungsgesetz}{刑法修正法}}
\newcommand{\abbrStenBer}{\abbrnoteonce{stenber}{StenBer.}{Stenographischer Bericht}{速記紀錄}}
\newcommand{\abbrStGB}{\abbrnoteonce{stgb}{StGB}{Strafgesetzbuch}{刑法}}
\newcommand{\abbrStREG}{\abbrnoteonce{streg}{StREG}{Strafrechtsreform-Ergänzungsgesetz}{刑法改革補充法}}
\newcommand{\abbrStrRG}{\abbrnoteonce{strrg}{StrRG}{Strafrechtsreformgesetz}{刑法改革法}}
\newcommand{\abbrUS}{\abbrnoteonce{us}{U.S.}{United States Reports}{美國最高法院判決彙編}}
\newcommand{\abbrVgl}{\abbrnoteonce{vgl}{vgl.}{vergleiche}{參見}}
\newcommand{\abbrVVDStRL}{\abbrnoteonce{vvdstrl}{VVDStRL}{Veröffentlichungen der Vereinigung der Deutschen Staatsrechtslehrer}{德國公法教師協會論文集}}
\newcommand{\abbrWP}{\abbrnoteonce{wp}{WP.}{Wahlperiode}{屆期}}
\newcommand{\abbrWp}{\abbrnoteonce{wp}{Wp.}{Wahlperiode}{屆期}}
\newcommand{\abbrZStrW}{\abbrnoteonce{zstrw}{ZStrW}{Zeitschrift für die gesamte Strafrechtswissenschaft}{整體刑法學期刊}}
\newcommand{\recordsourcerandstart}{%
  \ifrandnumbers
    \setcounter{lastsourcerandstart}{\value{randnumber}}%
    \global\lastsourcehasrandnumberstrue
  \else
    \global\lastsourcehasrandnumbersfalse
  \fi
}
\newlength{\biparagraphskip}
\setlength{\biparagraphskip}{0.52em}
\newlength{\lawboxblockbeforeskip}
\setlength{\lawboxblockbeforeskip}{0.72em}
\newlength{\lawboxblockafterskip}
\setlength{\lawboxblockafterskip}{0.92em}
\newif\iflawboxpartendedblock
\lawboxpartendedblockfalse
\newcommand{\sourceparstyle}{\footnotesize\color{sourceink}\setstretch{1.12}\RaggedRight\emergencystretch=3em\setlength{\parskip}{0.42em}\obeylines\selectfont}
\newcommand{\transparstyle}{\songfont\color{caseink}\setstretch{1.12}\setlength{\parindent}{0pt}\setlength{\parskip}{0.42em}}
\newcommand{\bilingualpairskip}{%
  \par\addvspace{\biparagraphskip}%
  \switchcolumn
  \par\addvspace{\biparagraphskip}%
  \switchcolumn*%
  \global\intranslationfalse
}
\newcommand{\bilingualboxpairskip}{%
  \par
  \switchcolumn
  \par
  \switchcolumn*%
  \global\intranslationfalse
}
\newcommand{\bilinguallawboxblockskip}{%
  \switchcolumn*%
  \par\addvspace{\lawboxblockafterskip}%
  \switchcolumn
  \par\addvspace{\lawboxblockafterskip}%
  \switchcolumn*%
  \global\intranslationfalse
}
\newcommand{\bikeepwithnext}[1][4]{%
  \ifbilingualcolumns
    \syncleft
    \Needspace{#1\baselineskip}%
    \switchcolumn
    \Needspace{#1\baselineskip}%
    \switchcolumn*%
    \global\intranslationfalse
  \else
    \Needspace{#1\baselineskip}%
  \fi
}
\NewEnviron{sourcepara}[1][]{
  \recordsourcerandstart
  \ifshoworiginal
    \ifbilingualcolumns
      \ifintranslation\switchcolumn*\global\intranslationfalse\fi
    \else
      \Needspace{5\baselineskip}
    \fi
    \ifbilingualcolumns\else\par\addvspace{0.35em}\fi
    {\sourceparstyle
    \noindent\BODY\par}
    \ifbilingualcolumns\else\addvspace{0.25em}\fi
    \ifbilingualcolumns
      \switchcolumn
      \global\intranslationtrue
    \fi
  \else
    \ifrandnumbers
      \begingroup
        \setbox0=\vbox{\hsize=\linewidth\sourceparstyle\noindent\BODY\par}%
      \endgroup
    \fi
  \fi
}
\NewEnviron{transpara}{
  \ifshowtranslation
    \setcounter{randnumberlocal}{\value{lastsourcerandstart}}%
    \ifbilingualcolumns
      \ifintranslation\else\switchcolumn\global\intranslationtrue\fi
      {\transparstyle\setlength{\leftskip}{0.55em}\setlength{\rightskip}{0.15em plus 1em}\addtolength{\linewidth}{-0.55em}\BODY\par}
      \switchcolumn*
      \bilingualpairskip
    \else
      {\transparstyle\BODY\par}
    \fi
  \else
    \ifbilingualcolumns
      \ifintranslation\switchcolumn*\global\intranslationfalse\fi
    \fi
  \fi
}
\NewEnviron{sourceboxpara}[1][]{
  \global\lawboxpartendedblockfalse
  \recordsourcerandstart
  \ifshoworiginal
    \ifbilingualcolumns
      \ifintranslation\switchcolumn*\global\intranslationfalse\fi
    \else
      \Needspace{3\baselineskip}
    \fi
    {\sourceparstyle
    \BODY\par}
    \ifbilingualcolumns\else
      \iflawboxpartendedblock\addvspace{\lawboxblockafterskip}\fi
    \fi
    \ifbilingualcolumns
      \switchcolumn
      \global\intranslationtrue
    \fi
  \fi
}
\NewEnviron{transboxpara}{
  \global\lawboxpartendedblockfalse
  \ifshowtranslation
    \setcounter{randnumberlocal}{\value{lastsourcerandstart}}%
    \ifbilingualcolumns
      \ifintranslation\else\switchcolumn\global\intranslationtrue\fi
      {\transparstyle\BODY\par}
      \iflawboxpartendedblock
        \bilinguallawboxblockskip
      \else
        \switchcolumn*
        \bilingualboxpairskip
      \fi
    \else
      {\transparstyle\BODY\par}
      \iflawboxpartendedblock\addvspace{\lawboxblockafterskip}\fi
    \fi
  \else
    \ifbilingualcolumns
      \ifintranslation\switchcolumn*\global\intranslationfalse\fi
    \fi
  \fi
}
\newcommand{\leitsatznum}[1]{\noindent\hangindent=3.0em\hangafter=1\makebox[2.25em][r]{\bfseries #1.}\hspace{0.55em}\ignorespaces}
\newcommand{\syncleft}{\ifbilingualcolumns\ifintranslation\switchcolumn*\global\intranslationfalse\fi\fi}
\pretocmd{\section}{\syncleft\clearpage}{}{}
\pretocmd{\subsection}{\syncleft}{}{}
\pretocmd{\subsubsection}{\syncleft}{}{}
\newcommand{\biheadingfont}{\songfont}
\newcommand{\bitocentry}[2]{%
  \ifshoworiginal
    \ifshowtranslation
      \ifstrequal{#1}{#2}{#1}{#1\space #2}%
    \else
      #1%
    \fi
  \else
    #2%
  \fi
}
\pdfstringdefDisableCommands{%
  \def\bitocentry#1#2{#1}%
}
\newcommand{\biheading}[7]{%
  \syncleft#1\bikeepwithnext[#3]\phantomsection\addcontentsline{toc}{#2}{\bitocentry{#6}{#7}}%
  \ifbilingualcolumns
    {#4 #6\par}%
    \switchcolumn
    {#4\biheadingfont #7\par}%
    \switchcolumn*%
    \global\intranslationfalse
    \par\addvspace{#5}%
    \switchcolumn
    \par\addvspace{#5}%
    \switchcolumn*%
  \else
    \ifshoworiginal{#4 #6\par}\fi
    \ifshowtranslation{#4\biheadingfont #7\par}\fi
    \addvspace{#5}%
  \fi
}
\newcommand{\termsectionheading}[1]{%
  \par\Needspace{9\baselineskip}%
  \vspace{0.65em}%
  {\songfont\bfseries #1\par}%
  \vspace{0.2em}%
}
% 章節換頁：
%   雙語 paracol 欄內不要用 \clearpage；它會在同步兩欄時吐出只有欄線/頁碼的空頁。
%   \newpage 足以讓下一個 section 起新頁，又不會強制清空所有待排材料。
%   單語模式不在 paracol 內，仍可用 \clearpage。
\newcommand{\bisectionpagebreak}{\ifbilingualcolumns\newpage\else\clearpage\fi}
\newcommand{\bisection}[2]{%
  \biheading{\iffirstbilingualsection\global\firstbilingualsectionfalse\else\bisectionpagebreak\fi}{section}{6}{\centering\Large\bfseries}{0.8em}{#1}{#2}%
}
\newcommand{\bisubsection}[2]{%
  \biheading{}{subsection}{5}{\large\bfseries}{0.55em}{#1}{#2}%
}
\newcommand{\bisubsubsection}[2]{%
  \biheading{}{subsubsection}{4}{\normalsize\bfseries}{0.45em}{#1}{#2}%
}
\newcommand{\bicompositeheading}[4]{%
  \biheading{}{subsection}{5}{\large\bfseries}{0.16em}{#1}{#2}%
  \biheading{}{subsubsection}{3}{\normalsize\bfseries}{0.45em}{#3}{#4}%
}
\newif\iflawboxfirstitem
\lawboxfirstitemfalse
\newcommand{\lawboxprespace}[1]{%
  \par
  \iflawboxfirstitem
    \global\lawboxfirstitemfalse
  \else
    \addvspace{#1}%
  \fi
}
\newtcolorbox{lawboxinner}{
  caseboxbase,
  colframe=sourceline!85!black,
  colback=white,
  left=1.05em,
  right=1.05em,
  top=0.62em,
  bottom=0.62em,
  before upper={\global\inlawboxtrue\global\lawboxfirstitemtrue\global\lawboxhaspendingrnfalse\ifintranslation\lawboxtransfont\else\lawboxsourcefont\fi\RaggedRight\setlength{\parindent}{0pt}\setlength{\parskip}{0pt}\ifintranslation\setstretch{\lawboxtransstretch}\else\setstretch{\lawboxsourcestretch}\fi},
  after upper={\global\lawboxhaspendingrnfalse\global\inlawboxfalse}
}
\tcbset{
  lawboxpartbase/.style={
    enhanced,
    breakable=false,
    width=0.985\linewidth,
    colback=white,
    frame hidden,
    boxrule=0pt,
    arc=0pt,
    left=1.05em,
    right=1.05em,
    boxsep=1.5mm,
    before skip=0pt,
    after skip=0pt,
    before upper={\global\inlawboxtrue\global\lawboxfirstitemtrue\global\lawboxhaspendingrnfalse\ifintranslation\lawboxtransfont\else\lawboxsourcefont\fi\RaggedRight\setlength{\parindent}{0pt}\setlength{\parskip}{0pt}\ifintranslation\setstretch{\lawboxtransstretch}\else\setstretch{\lawboxsourcestretch}\fi},
    after upper={\global\lawboxhaspendingrnfalse\global\inlawboxfalse},
    borderline west={0.28pt}{0pt}{sourceline!85!black},
    borderline east={0.28pt}{0pt}{sourceline!85!black},
  },
  lawboxpartfirst/.style={
    lawboxpartbase,
    top=0.62em,
    bottom=0.04em,
    before skip=\lawboxblockbeforeskip,
    borderline north={0.28pt}{0pt}{sourceline!85!black},
  },
  lawboxpartmiddle/.style={
    lawboxpartbase,
    top=0.04em,
    bottom=0.04em,
  },
  lawboxpartlast/.style={
    lawboxpartbase,
    top=0.04em,
    bottom=0.62em,
    after skip=0pt,
    borderline south={0.28pt}{0pt}{sourceline!85!black},
  }
}
\newtcolorbox{lawboxpartinner}[1]{#1}
\newtcolorbox{termboxinner}{
  caseboxbase,
  colframe=noteblue!70!black,
  colback=casepaper,
  before upper={\songfont\setlength{\parindent}{0pt}\setlength{\parskip}{0.35em}}
}
\newtcolorbox{deboxinner}{
  caseboxbase,
  colframe=notegray!72!black,
  colback=white,
  left=1.05em,
  right=1.05em,
  top=0.62em,
  bottom=0.62em,
  before upper={\global\inlawboxtrue\global\lawboxfirstitemtrue\global\lawboxhaspendingrnfalse\small\setlength{\parindent}{0pt}\setlength{\parskip}{0.35em}},
  after upper={\global\lawboxhaspendingrnfalse\global\inlawboxfalse}
}
\NewEnviron{lawbox}{
  \begin{lawboxinner}\BODY\end{lawboxinner}
}
\NewEnviron{lawboxpart}[2]{
  \begin{lawboxpartinner}{lawboxpart#1,equal height group=#2}\BODY\end{lawboxpartinner}%
  \ifstrequal{#1}{last}{\global\lawboxpartendedblocktrue}{}%
}
\NewEnviron{termbox}{
  \par\addvspace{0.55em}
  {\color{noteblue!70!black}\rule{\linewidth}{0.28pt}}\par
  \begingroup
    \songfont\small\color{caseink}
    \setlength{\parindent}{0pt}
    \setlength{\parskip}{0.24em}
    \BODY
    \par
  \endgroup
  {\color{noteblue!70!black}\rule{\linewidth}{0.28pt}}\par
  \addvspace{0.55em}
}
\NewEnviron{debox}{
  \begin{deboxinner}\BODY\end{deboxinner}
}
\providecommand{\paragraphmark}[1]{}
\newcommand{\lawboxtitleafter}{\addvspace{0.06em}}
\newcommand{\lawtitle}[1]{\lawboxprespace{0.22em}{\lawboxtransfont\bfseries\lawboxuserandnumber #1}\par\lawboxtitleafter}
\newcommand{\absno}[2]{\lawboxprespace{0.04em}\noindent\lawboxuserandnumber\hangindent=2.6em\hangafter=1\makebox[2.6em][l]{(#1)}#2\par}
\newcommand{\nrno}[2]{\lawboxprespace{0pt}\noindent\lawboxuserandnumber\hspace*{2.6em}\hangindent=4.4em\hangafter=1\makebox[1.8em][l]{#1.}#2\par}
\newcommand{\letterno}[2]{\lawboxprespace{0pt}\noindent\lawboxuserandnumber\hspace*{4.4em}\hangindent=6.0em\hangafter=1\makebox[1.6em][l]{#1)}#2\par}
\newcommand{\sourcelawtitle}[1]{\lawboxprespace{0.22em}{\lawboxsourcefont\bfseries\lawboxuserandnumber #1}\par\lawboxtitleafter}
\newcommand{\sourcelawsubtitle}[1]{\lawboxprespace{0.04em}{\lawboxsourcefont\bfseries\lawboxuserandnumber #1}\par\lawboxtitleafter}
\newcommand{\sourceabsno}[2]{\lawboxprespace{0.04em}\noindent\lawboxuserandnumber\hangindent=2.45em\hangafter=1\makebox[2.45em][l]{(#1)}#2\par}
\newcommand{\sourcenrno}[2]{\lawboxprespace{0pt}\noindent\lawboxuserandnumber\hspace*{2.45em}\hangindent=4.25em\hangafter=1\makebox[1.8em][l]{#1.}#2\par}
\newcommand{\sourceletterno}[2]{\lawboxprespace{0pt}\noindent\lawboxuserandnumber\hspace*{4.25em}\hangindent=5.85em\hangafter=1\makebox[1.6em][l]{#1)}#2\par}
\newcommand{\sourcequotepara}[1]{\lawboxprespace{0.04em}\noindent\lawboxuserandnumber\hangindent=1.2em\hangafter=1 #1\par}
\newcommand{\sourcelawline}[1]{\lawboxprespace{0.04em}\noindent\lawboxuserandnumber #1\par}
\newcommand{\lawline}[1]{\lawboxprespace{0.04em}\noindent\lawboxuserandnumber #1\par}
\newcommand{\pendingtranslation}{\par{\songfont\footnotesize\color{pendingink}（中文譯文待補。）}\par}
\newcommand{\tocpart}[1]{\par\addvspace{0.52em}{\footnotesize\bfseries\color{pendingink}#1\par}\addvspace{0.16em}}
\newcommand{\tocdashline}{\par\addvspace{0.48em}\noindent{\color{notegray}\xleaders\hbox to 0.82em{\hfil\rule{0.42em}{0.28pt}\hfil}\hskip\linewidth}\par\addvspace{0.38em}}
% \tocpanel{font-cmd}{content}：將目錄置中於所在欄寬（雙語欄適用）。
% A3 橫式使用 0.58\linewidth，A4 橫式使用 0.84\linewidth，
% 使兩種紙張下目錄欄內容實際寬度相近（均約 100–105 mm）。
\newcommand{\tocpanel}[2]{%
  \makebox[\linewidth][c]{%
    \ifx\bilingualpaper\bilingualpaperaThree
      \begin{minipage}{0.58\linewidth}%
    \else
      \begin{minipage}{0.84\linewidth}%
    \fi
    \toccolstyle
    #1#2%
    \end{minipage}%
  }%
}
% ===== 首頁簡目錄的兩層 description list 縮排 =====
% enumitem 的 description list 可把每一列拆成：
%   [label]  labelsep  body text
% 這裡的「起點」都以所在 minipage/欄位的左緣為 0。
%
% 核心規則：
%   labelindent = label 欄位從左緣開始的位置。
%   labelwidth  = label 欄位本身寬度；標籤太長（如 Abw. Meinung）就加大它。
%   labelsep    = label 欄位右緣到正文的距離。
%   leftmargin  = 正文 body text 的起點。判斷層級視覺時主要看這個。
%
% 所以：
%   想讓「Begründung / 判決理由」整列正文往右或往左：改 tocitems 的 leftmargin。
%   想讓「A. Verfahren und Gegenstand / A. 程序與審查標的」正文往右或往左：
%     改 tocsubitems 的 leftmargin。
%   想讓 A./B./C. 這些標籤本身往右或往左，但正文不動：改 tocsubitems 的 labelindent。
%   想讓 A./B./C. 與正文間距變大或變小：改 tocsubitems 的 labelsep。
%
% 層級原則：
%   tocsubitems.leftmargin 必須大於 tocitems.leftmargin，
%   否則子層級正文會看起來比「Gründe / 理由」還靠左。
\newlist{tocitems}{description}{1}
\setlist[tocitems]{%
  labelindent=0pt,    % 第一層 label 欄位起點；0pt 表示貼齊目錄欄左緣。
  leftmargin=2.54cm,  % 第一層正文起點；例如 Begründung / 判決理由 的左緣。
  labelwidth=2.16cm,  % 第一層 label 欄位寬度；需容納 Leitsätze、Abw. Meinung、不同意見等。
  labelsep=0.32cm,    % 第一層 label 與正文的水平間距；只改間距，不改正文起點。
  itemsep=0.28em,     % 第一層各項之間的垂直距離。
  topsep=0.20em,      % 第一層 list 上下與前後內容的垂直距離。
  parsep=0pt,         % 同一 item 內段落之間的距離；目錄項通常不需要。
  partopsep=0pt,      % list 若緊接段落開始時的額外距離；關閉以免目錄高度飄動。
  align=left,         % label 欄內左對齊；長短標籤不靠右貼齊。
  font=\normalfont\bfseries % label 的字型；正文不受這行影響。
}
\newlist{tocsubitems}{description}{1}
\setlist[tocsubitems]{%
  labelindent=1cm,    % 第二層 label 起點；只控制 A./B./C. 本身的位置。
  leftmargin=3.00cm,  % 第二層正文起點；必須大於 2.54cm，才會比 Begründung / 判決理由更右。
  labelwidth=0.5cm,   % 第二層 label 欄寬；A./B./C. 很短，可小於第一層。
  labelsep=1.5cm,    % 第二層 label 與正文間距；A. 與 Verfahren / 程序 的距離看這裡。
  itemsep=0.12em,     % 第二層各項較密，避免 A.–E. 佔太高。
  topsep=0.04em,      % 第二層 list 與 Gründe / 理由之間的垂直距離。
  parsep=0pt,         % 同一第二層 item 內段落距離；目錄項通常不需要。
  partopsep=0pt,      % 關閉額外段落起始距離，避免版面不穩。
  align=left,         % A./B./C. label 欄內左對齊。
  font=\normalfont\bfseries, % A./B./C. label 加粗；正文不受這行影響。
  before={}           % 不在第二層 list 前自動插入額外內容。
}
\newlist{termitems}{description}{1}
\ifbilingualcolumns
  \setlist[termitems]{%
    leftmargin=5.2cm,
    labelwidth=4.8cm,
    labelsep=0.35cm,
    itemsep=0.16em,
    topsep=0.2em,
    parsep=0pt,
    partopsep=0pt,
    font=\normalfont\bfseries,
    style=nextline
  }
\else
  \setlist[termitems]{%
    leftmargin=0pt,
    labelwidth=0pt,
    labelsep=0pt,
    itemsep=0.24em,
    topsep=0.2em,
    parsep=0pt,
    partopsep=0pt,
    font=\normalfont\bfseries,
    style=nextline
  }
\fi
\newlist{abbritems}{description}{1}
\setlist[abbritems]{%
  leftmargin=2.38cm,
  labelwidth=2.02cm,
  labelsep=0.28cm,
  itemsep=0.16em,
  topsep=0.2em,
  parsep=0pt,
  partopsep=0pt,
  align=left,
  font=\normalfont\bfseries
}
\ifbilingualcolumns
  \newenvironment{termcolumns}{\begin{multicols}{2}}{\end{multicols}}
\else
  \newenvironment{termcolumns}{}{}
\fi

% ===== 編者註腳開關 =====
% 一般版預設不顯示編者註，避免正文腳註過密。
% 若開啟 \classroommodeon，GG/條文編者註仍會自動顯示，無需另開本開關。
% 若一般版也要顯示編者註，可在單案檔 \input 後加 \showeditornotestrue。
\newif\ifshoweditornotes
\showeditornotesfalse

% ===== 目錄排版輔助 =====
\newcommand{\toccolstyle}{\raggedright\arraybackslash}
\newcommand{\tocsingleindent}{\leftskip=1.45cm\rightskip=1.2cm}

% ===== 行內引文環境（德文原文與中文譯文各一，帶左側灰色邊線）=====
\newcommand{\quoteparbreak}{\par\addvspace{0.48em}}
\newcommand{\sourcequoteinner}[1]{%
  \begin{tcolorbox}[
    enhanced, breakable, width=\dimexpr\linewidth-0.8em\relax, frame hidden,
    colback=white, coltext=caseink, boxrule=0pt,
    borderline west={0.55pt}{0.88em}{notegray},
    left=1.78em, right=0.72em, top=0.18em, bottom=0.18em,
    boxsep=0pt, before skip=0.24em, after skip=0.24em,
    before upper={\normalfont\footnotesize\color{caseink}\setlength{\parindent}{0pt}\setlength{\parskip}{0.34em}\raggedright\setlength{\rightskip}{0pt plus 1em}}
  ]%
  #1\par
  \end{tcolorbox}%
}
\newcommand{\transquoteinner}[1]{%
  \begin{tcolorbox}[
    enhanced, breakable, width=\dimexpr\linewidth-0.8em\relax, frame hidden,
    colback=white, coltext=caseink, boxrule=0pt,
    borderline west={0.55pt}{0.88em}{notegray},
    left=1.78em, right=0.72em, top=0.18em, bottom=0.18em,
    boxsep=0pt, before skip=0.24em, after skip=0.24em,
    before upper={\songfont\small\color{caseink}\setlength{\parindent}{0pt}\setlength{\parskip}{0.34em}\raggedright\setlength{\rightskip}{0pt plus 1em}}
  ]%
  #1\par
  \end{tcolorbox}%
}

% ===== 大綱（Gliederung）Rn 輔助命令 =====
\newcommand{\outrn}[2]{\enspace{\normalfont\footnotesize\color{pendingink}(Rn\,#1–#2)}}
\newcommand{\outrnsingle}[1]{\enspace{\normalfont\footnotesize\color{pendingink}(Rn\,#1)}}
\newcommand{\outrnzh}[2]{\enspace{\normalfont\footnotesize\color{pendingink}（第\,#1–#2\,段）}}
\newcommand{\outrnzhs}[1]{\enspace{\normalfont\footnotesize\color{pendingink}（第\,#1\,段）}}
% ===== 大綱雙語對照表格輔助命令（三欄：段碼 │ 德文 │ 中文）=====
% 欄 1：段碼（由欄定義自動格式化：小字、等寬、右對齊、pendingink）
% 欄 2：德文標籤＋正文
% 欄 3：中文標籤＋正文
%
% 無段碼的列（\omainrow、\osecrow），欄 1 以 & 空置。
% 有段碼的列（\osecrowrn、\osubrow、\ointrow），#1 為預格式化字串
%   例：\osubrow{2–49}{I.}{...}{I.}{...}
%       \ointrow{107}{...}{...}
%
% \opartrow：段組標題（橫跨三欄）
\newcommand{\opartrow}[2]{%
  \noalign{\vspace{0.30em}}%
  \multicolumn{3}{@{}l@{}}{%
    \footnotesize\bfseries\color{pendingink}#1\hfill\songfont#2%
  }\\[0.06em]%
}
% \odashrow：虛線分隔（橫跨三欄）
\newcommand{\odashrow}{%
  \noalign{\vspace{0.28em}}%
  \multicolumn{3}{@{}l@{}}{\color{notegray}%
    \xleaders\hbox to 0.82em{\hfil\rule{0.42em}{0.28pt}\hfil}\hskip 0.88\linewidth}%
  \\[0.24em]%
}
% \odissentpart：不同意見書區段標頭。比一般 \opartrow 更像附錄性轉場。
\newcommand{\odissentpart}[2]{%
  \noalign{\vspace{0.42em}}%
  \multicolumn{3}{@{}p{0.88\textwidth}@{}}{%
    \color{notegray}\rule{\linewidth}{0.18pt}%
  }\\[-0.18em]%
  \multicolumn{3}{@{}p{0.88\textwidth}@{}}{%
    \makebox[\linewidth][s]{%
      \footnotesize\bfseries\color{pendingink}#1%
      \hfill{\songfont #2}%
    }%
  }\\[-0.02em]%
}
% \odissentopinion：不同意見書的署名與一行摘要，右欄保留段碼範圍。
\newcommand{\odissentopinion}[5]{%
  \footnotesize\hangindent=0.52cm\hangafter=1%
    {\bfseries\color{caseink}#2}\enspace{\color{notegray}\textperiodcentered}\enspace\color{caseink}#3%
  &\footnotesize\hangindent=0.52cm\hangafter=1%
    {\songfont\bfseries\color{caseink}#4}\enspace{\color{notegray}\textperiodcentered}\enspace\songfont\color{caseink}#5%
  &#1%
  \\[0.12em]%
}
% \omainrow：頂層項目，無段碼（Leitsätze、Urteil、Gründe …）
\newcommand{\omainrow}[4]{%
  \footnotesize\hangindent=1.92cm\hangafter=1%
    \makebox[1.92cm][l]{\bfseries\color{caseink}#1}\color{caseink}#2%
  &\footnotesize\hangindent=1.92cm\hangafter=1%
    \makebox[1.92cm][l]{\bfseries\color{caseink}#3}\songfont\color{caseink}#4%
  &%
  \\[0.15em]%
}
% \omainrowrn：頂層項目，有段碼（不同意見等標籤寬於 \osecrow 框者）
\newcommand{\omainrowrn}[5]{%
  \footnotesize\hangindent=1.92cm\hangafter=1%
    \makebox[1.92cm][l]{\bfseries\color{caseink}#2}\color{caseink}#3%
  &\footnotesize\hangindent=1.92cm\hangafter=1%
    \makebox[1.92cm][l]{\bfseries\color{caseink}#4}\songfont\color{caseink}#5%
  &#1%
  \\[0.15em]%
}
% \osecrow：一級子項，無段碼（A.、C.、D. 等有子項者）
\newcommand{\osecrow}[4]{%
  \footnotesize\hspace{1.10cm}\hangindent=1.98cm\hangafter=1%
    \makebox[0.86cm][l]{\bfseries\color{caseink}#1}\color{caseink}#2%
  &\footnotesize\hspace{1.10cm}\hangindent=1.98cm\hangafter=1%
    \makebox[0.86cm][l]{\bfseries\color{caseink}#3}\songfont\color{caseink}#4%
  &%
  \\[0.11em]%
}
% \osecrowrn：一級子項，有段碼（B.、E. 等完整段落範圍者）
% #1=段碼字串（如 101–106）, #2=DE標籤, #3=DE正文, #4=ZH標籤, #5=ZH正文
\newcommand{\osecrowrn}[5]{%
  \footnotesize\hspace{1.10cm}\hangindent=1.98cm\hangafter=1%
    \makebox[0.86cm][l]{\bfseries\color{caseink}#2}\color{caseink}#3%
  &\footnotesize\hspace{1.10cm}\hangindent=1.98cm\hangafter=1%
    \makebox[0.86cm][l]{\bfseries\color{caseink}#4}\songfont\color{caseink}#5%
  &#1%
  \\[0.11em]%
}
% \osubrow：二級子項，有段碼範圍（I.、II.、A.I. …）
% #1=段碼字串, #2=DE標籤, #3=DE正文, #4=ZH標籤, #5=ZH正文
\newcommand{\osubrow}[5]{%
  \footnotesize\hspace{1.40cm}\hangindent=2.30cm\hangafter=1%
    \makebox[0.90cm][l]{\bfseries\color{caseink}#2}\color{caseink}#3%
  &\footnotesize\hspace{1.40cm}\hangindent=2.30cm\hangafter=1%
    \makebox[0.90cm][l]{\bfseries\color{caseink}#4}\songfont\color{caseink}#5%
  &#1%
  \\[0.09em]%
}
% \ointrow：引入段落，有單一段碼，無標籤
% #1=段碼字串, #2=DE正文, #3=ZH正文
\newcommand{\ointrow}[3]{%
  \footnotesize\hspace{0.52cm}\hangindent=0.52cm\hangafter=1\color{caseink}#2%
  &\footnotesize\hspace{0.52cm}\hangindent=0.52cm\hangafter=1\songfont\color{caseink}#3%
  &#1%
  \\[0.09em]%
}
% \osubsubrow：三級子項（1.、2.、3.）
% #1=段碼字串, #2=DE標籤, #3=DE正文, #4=ZH標籤, #5=ZH正文
\newcommand{\osubsubrow}[5]{%
  \footnotesize\hspace{1.80cm}\hangindent=2.48cm\hangafter=1%
    \makebox[0.68cm][l]{\bfseries\color{caseink}#2}\color{caseink}#3%
  &\footnotesize\hspace{1.80cm}\hangindent=2.48cm\hangafter=1%
    \makebox[0.68cm][l]{\bfseries\color{caseink}#4}\songfont\color{caseink}#5%
  &#1%
  \\[0.08em]%
}
% \osubsubsubrow：四級子項（a）、b））
% #1=段碼字串, #2=DE標籤, #3=DE正文, #4=ZH標籤, #5=ZH正文
\newcommand{\osubsubsubrow}[5]{%
  \footnotesize\hspace{2.10cm}\hangindent=2.68cm\hangafter=1%
    \makebox[0.58cm][l]{\bfseries\color{caseink}#2}\color{caseink}#3%
  &\footnotesize\hspace{2.10cm}\hangindent=2.68cm\hangafter=1%
    \makebox[0.58cm][l]{\bfseries\color{caseink}#4}\songfont\color{caseink}#5%
  &#1%
  \\[0.07em]%
}
% \osubsubsubsubrow：五級子項（aa）、bb））
% 標籤框 0.62cm：足以容納「aa)」在 footnotesize bold 下（約 0.50cm）並留有間距
% #1=段碼字串, #2=DE標籤, #3=DE正文, #4=ZH標籤, #5=ZH正文
\newcommand{\osubsubsubsubrow}[5]{%
  \footnotesize\hspace{2.38cm}\hangindent=3.06cm\hangafter=1%
    \makebox[0.68cm][l]{\bfseries\color{caseink}#2}\color{caseink}#3%
  &\footnotesize\hspace{2.38cm}\hangindent=3.06cm\hangafter=1%
    \makebox[0.68cm][l]{\bfseries\color{caseink}#4}\songfont\color{caseink}#5%
  &#1%
  \\[0.06em]%
}
% \outlinepage：雙語對照大綱頁包裝環境（longtable 自動跨頁）
% 三欄：德文欄 + 中文欄 + 段碼欄（最右，不折行）
% 段碼欄寬度：1.80cm，足以容納最長段碼如「309–348」
% DE/ZH 欄寬：(0.87\textwidth - 2.08cm) / 2
%   驗算：2×欄 + 0.05\textwidth + 0.28cm gap + 1.80cm = 0.87tw - 2.08cm + 0.05tw + 2.08cm = 0.92tw ✓
\newenvironment{outlinepage}{%
  \vspace*{0.40cm}%
  \setlength{\LTleft}{0.04\textwidth}%
  \setlength{\LTright}{0.04\textwidth}%
  \setlength{\tabcolsep}{0pt}%
  % 大綱列距由 longtable 的 \arraystretch 與各 row macro 結尾的 \\[...em] 共同決定。
  % 這裡控制「每一列本身的表格 strut 高度」；數值越大，A./I./1. 各項之間越像有空行。
  % A3 原本用 0.62，視覺上沒有空行；A4 也沿用同值，避免切到 A4 後項目被撐開。
  % 若日後覺得大綱太擠，只微調這個值（例如 0.68），不要改各 \omainrow/\osubrow 的縮排。
  \renewcommand{\arraystretch}{0.62}%
  \begin{longtable}{%
    >{\vspace{0pt}\raggedright\arraybackslash}p{\dimexpr(0.87\textwidth-2.08cm)/2\relax}%
    @{\hspace{0.025\textwidth}}%
    !{\color{notegray}\vrule width 0.12pt}%
    @{\hspace{0.025\textwidth}}%
    >{\vspace{0pt}\raggedright\arraybackslash}p{\dimexpr(0.87\textwidth-2.08cm)/2\relax}%
    @{\hspace{0.28cm}}%
    >{\vspace{0pt}\footnotesize\normalfont\color{pendingink}\raggedright\arraybackslash}p{1.80cm}%
    @{}%
  }%
  % 首頁欄標籤：不要橫跨置中；分別放在德文欄與中文欄的實際欄位起點。
  \footnotesize\bfseries\color{sourceink}Gliederung%
  &\footnotesize\bfseries\songfont\color{sourceink}大綱%
  &%
  \\[0.36em]%
  \endfirsthead%
  % 續頁欄標籤：同樣分欄定位，以灰色細體區分；無需標注「續」。
  \footnotesize\color{pendingink}Gliederung%
  &\footnotesize\songfont\color{pendingink}大綱%
  &%
  \\[0.24em]%
  \endhead%
}{%
  \end{longtable}%
}
 之前設定；
%   預設 chinese / a4）
\ifdefined\docmode\else\def\docmode{chinese}\fi
\ifdefined\bilingualpaper\else\def\bilingualpaper{a4}\fi
\newif\ifshoworiginal
\newif\ifshowtranslation
\newif\ifbilinguallandscape
\newif\ifbilingualcolumns
\newif\ifintranslation
\newif\iffirstbilingualsection

\showoriginalfalse
\showtranslationfalse
\bilinguallandscapefalse
\bilingualcolumnsfalse
\intranslationfalse
\firstbilingualsectionfalse

\def\modechinese{chinese}
\def\modegerman{german}
\def\modebilingual{bilingual}
\def\bilingualpaperaThree{a3}
\def\bilingualpaperaFour{a4}

\ifx\docmode\modechinese
  \showtranslationtrue
\fi

\ifx\docmode\modegerman
  \showoriginaltrue
\fi

\ifx\docmode\modebilingual
  \showoriginaltrue
  \showtranslationtrue
  \bilinguallandscapetrue
  \bilingualcolumnstrue
\fi

% 編排模式說明：
% chinese  → \showoriginalfalse  \showtranslationtrue  （A4 直式，純中文）
% german   → \showoriginaltrue   \showtranslationfalse （A4 直式，純德文）
% bilingual→ \showoriginaltrue   \showtranslationtrue  \bilingualcolumnstrue（A3/A4 橫式對照）

\ifbilingualcolumns
  \ifx\bilingualpaper\bilingualpaperaFour
    \geometry{
      a4paper,
      landscape,
      left=1.25cm,
      right=2.25cm,
      top=1.25cm,
      bottom=1.75cm,
      footskip=8.5mm,
      marginparwidth=0.75cm,
      marginparsep=0.25cm
    }
  \else
    \geometry{
      a3paper,
      landscape,
      left=1.55cm,
      right=2.75cm,
      top=1.55cm,
      bottom=2.05cm,
      footskip=9.5mm,
      marginparwidth=0.9cm,
      marginparsep=0.35cm
    }
  \fi
\else
\geometry{
  a4paper,
  portrait,
  left=2.2cm,
  right=2.6cm,
  top=2.0cm,
  bottom=2.0cm,
  marginparwidth=0.8cm,
  marginparsep=0.3cm
}
\fi
% ===== 時間戳基礎設施 =====
\newcount\stamptimeval
\newcount\stampminuteval
\newcommand{\stamphh}{%
  \stamptimeval=\time
  \divide\stamptimeval by 60
  \ifnum\stamptimeval<10 0\fi
  \the\stamptimeval}
\newcommand{\stampmm}{%
  \stamptimeval=\time
  \stampminuteval=\time
  \divide\stampminuteval by 60
  \multiply\stampminuteval by 60
  \advance\stamptimeval by -\stampminuteval
  \ifnum\stamptimeval<10 0\fi
  \the\stamptimeval}
\newcommand{\stampwkde}[1]{%
  \ifcase#1Montag\or Dienstag\or Mittwoch\or Donnerstag\or Freitag\or Samstag\or Sonntag\fi}
\newcommand{\stampwkzh}[1]{%
  \ifcase#1 一\or 二\or 三\or 四\or 五\or 六\or 日\fi}
\newcommand{\stampmonthde}[1]{%
  \ifcase#1\or Januar\or Februar\or März\or April\or Mai\or Juni\or
  Juli\or August\or September\or Oktober\or November\or Dezember\fi}
\newcount\stampjdval
\newcount\stampwdval
\pgfcalendardatetojulian{\the\year-\the\month-\the\day}{\stampjdval}
\pgfcalendarjuliantoweekday{\the\stampjdval}{\stampwdval}
\newcommand{\timestampzh}{%
  \songfont 最後修改：\the\year 年\the\month 月\the\day 日%
  % （星期\stampwkzh{\the\stampwdval}）%
  % \space\stamphh:\stampmm
  }

\newcommand{\documentdatezh}{\the\year 年 \the\month 月 \the\day 日}
\newcommand{\documentdatede}{\the\day. \stampmonthde{\the\month} \the\year}
\newcommand{\bverfgetranslationnote}{%
  {\songfont 翻譯說明：初譯使用 Claude Code 與 GPT 協助迭代；仍請以德文原文為準。}%
}
\newcommand{\bverfgecourseinfo}{%
  {\songfont 課程：114 學年度第 2 學期；憲法專題研究（四）；授課老師：湯德宗教授；導讀日期：115 年 6 月 1 日。}%
}
\newcommand{\bverfgeeditorinfo}{%
  {\songfont 編者：王逸帆（R10A21126），國立台灣大學法律系研究所碩士班。}%
}
\newcommand{\bverfgetitleversionline}{%
  \ifbilingualcolumns
    Fassung \documentversion\enspace\textperiodcentered\enspace Stand: \documentdatede
    \quad
    {\songfont 版本 \documentversion\enspace\textperiodcentered\enspace 修訂日期：\documentdatezh\par}%
  \else
    \ifshowtranslation
      {\songfont 版本 \documentversion\enspace\textperiodcentered\enspace 修訂日期：\documentdatezh\par}%
    \else
      Fassung \documentversion\enspace\textperiodcentered\enspace Stand: \documentdatede\par
    \fi
  \fi
}
\newcommand{\bverfgetitlecornerinfo}{%
  \ifshowtranslation
    \vspace{0.72em}%
    \noindent\hfill
    \begin{minipage}{0.66\textwidth}
      \raggedleft\tiny\color{pendingink}\songfont
      \bverfgecourseinfo\par
      \bverfgeeditorinfo
    \end{minipage}%
  \fi
}
\newcommand{\bverfgetitlemeta}{%
  \vfill
  \begin{center}%
    \footnotesize\color{pendingink}%
    \bverfgetitleversionline
    \ifshowtranslation
      \vspace{0.28em}{\scriptsize\bverfgetranslationnote\par}%
    \fi
  \end{center}%
  \bverfgetitlecornerinfo
}

\pagestyle{fancy}
\fancyhf{}
\renewcommand{\headrulewidth}{0pt}
\renewcommand{\footrulewidth}{0pt}
\fancyfoot[C]{\thepage}
\ifbilingualcolumns
  \fancyfoot[R]{\scriptsize\color{pendingink}\timestampzh}%
\else
  \ifshoworiginal
  \else
    \fancyfoot[R]{\scriptsize\color{pendingink}\timestampzh\hspace{0.45cm}}%
  \fi
\fi
\setmainfont{Times New Roman}
\setsansfont{Helvetica Neue}
\setmonofont{Menlo}
\IfFileExists{/Users/iw/Library/Fonts/kaiu.ttf}{
  \setCJKmainfont[Path=/Users/iw/Library/Fonts/,AutoFakeBold=4]{kaiu.ttf}
}{
  \IfFontExistsTF{標楷體}{
    \setCJKmainfont[AutoFakeBold=4]{標楷體}
  }{
    \setCJKmainfont{Songti TC}
  }
}
\IfFontExistsTF{Songti TC}{
  \setCJKfamilyfont{song}{Songti TC}
  \setCJKmonofont{Songti TC}
}{
  \setCJKfamilyfont{song}[Path=/Users/iw/Library/Fonts/,AutoFakeBold=4]{kaiu.ttf}
  \setCJKmonofont[Path=/Users/iw/Library/Fonts/,AutoFakeBold=4]{kaiu.ttf}
}
\newcommand{\songfont}{\CJKfamily{song}}

% ===== 封面／目錄專用打字機字體（僅限封面、目錄頁使用，不影響正文）=====
\IfFontExistsTF{Radio Newsman}{%
  \newfontfamily\bverfgetitlede{Radio Newsman}%
}{%
  \newfontfamily\bverfgetitlede{Courier New}%
}
\IfFontExistsTF{Erikas Farbband}{%
  \newfontfamily\bverfgebodyde{Erikas Farbband}%
}{%
  \let\bverfgebodyde\bverfgetitlede
}
\IfFontExistsTF{Maoken Heavy Labourer}{%
  \setCJKfamilyfont{bverfgetitlecjk}{Maoken Heavy Labourer}%
}{%
  \setCJKfamilyfont{bverfgetitlecjk}{Songti TC}%
}
\IfFontExistsTF{Huiwen-mincho}{%
  \setCJKfamilyfont{bverfgebodycjk}{Huiwen-mincho}%
}{%
  \setCJKfamilyfont{bverfgebodycjk}{Songti TC}%
}
\newcommand{\bverfgetitlezh}{\bverfgetitlede\CJKfamily{bverfgetitlecjk}}
\newcommand{\bverfgebodyzh}{\bverfgebodyde\CJKfamily{bverfgebodycjk}}

\setstretch{1.35}
\setlist{itemsep=0.2em, topsep=0.4em}
\setlength{\parindent}{0pt}
\setlength{\parskip}{0.45em}
\raggedbottom
\setcounter{secnumdepth}{0}
\setcounter{tocdepth}{2}

\newcommand{\lawboxsourcefont}{\footnotesize}
\newcommand{\lawboxtransfont}{\footnotesize}
\newcommand{\lawboxsourcestretch}{1.02}
\newcommand{\lawboxtransstretch}{0.92}

\titleformat{\section}{\centering\Large\bfseries\songfont}{\thesection}{1em}{}
\titleformat{\subsection}{\large\bfseries\songfont}{\thesubsection}{1em}{}
\titleformat{\subsubsection}{\normalsize\bfseries\songfont}{\thesubsubsection}{1em}{}
\titleformat{\paragraph}{\normalsize\bfseries\songfont}{\theparagraph}{1em}{}
\titlespacing*{\section}{0pt}{1.8em}{1.0em}
\titlespacing*{\subsection}{0pt}{1.2em}{0.6em}
\titlespacing*{\subsubsection}{0pt}{1.0em}{0.45em}
\titlespacing*{\paragraph}{0pt}{0.6em}{0.4em}

\definecolor{noteblue}{HTML}{A9C9F5}
\definecolor{noteteal}{HTML}{AEE1D6}
\definecolor{notegray}{HTML}{D7DADF}
\definecolor{caseink}{HTML}{000000}
\definecolor{casepaper}{HTML}{FCFEFD}
\definecolor{sourcepaper}{HTML}{FAFBF8}
\definecolor{sourceline}{HTML}{D7DED8}
\definecolor{sourceink}{HTML}{000000}
\definecolor{pendingink}{HTML}{000000}

\tcbset{
  caseboxbase/.style={
    enhanced,
    breakable,
    width=0.985\linewidth,
    colback=casepaper,
    coltitle=caseink,
    fonttitle=\bfseries\songfont,
    boxrule=0.28pt,
    arc=1.2mm,
    left=2.4mm,
    right=2.4mm,
    top=1.5mm,
    bottom=1.5mm,
    boxsep=1.5mm,
    before skip=0.65em,
    after skip=0.65em,
    before upper={\setlength{\parindent}{0pt}\setlength{\parskip}{0.35em}},
  }
}

\newcommand{\bilingualstart}{}
\newcommand{\bilingualstop}{}
\ifbilingualcolumns
  \renewcommand{\bilingualstart}{\small\setlength{\columnsep}{1.25cm}\setlength{\columnseprule}{0.12pt}\columnratio{0.5,0.5}\multicolumnfootnotes\flushbottom\sloppy\interlinepenalty=80\clubpenalty=800\widowpenalty=800\displaywidowpenalty=800\begin{paracol}{2}\global\intranslationfalse\global\firstbilingualsectiontrue{\footnotesize\color{sourceink}Deutsch\par}\switchcolumn{\footnotesize\songfont\color{sourceink}中文譯文\par}\switchcolumn*}
  \renewcommand{\bilingualstop}{\ifintranslation\switchcolumn*\global\intranslationfalse\fi\end{paracol}\clearpage\raggedbottom}
\fi
\newif\ifrandnumbers
\randnumbersfalse
\newcounter{randnumber}
\newcounter{randnumberlocal}
\newcounter{lastsourcerandstart}
\newif\iflastsourcehasrandnumbers
\lastsourcehasrandnumbersfalse
\newif\ifinlawbox
\inlawboxfalse
\newif\iflawboxhaspendingrn
\lawboxhaspendingrnfalse
\newcommand{\lawboxpendingrn}{}
\newcommand{\startrandnumbers}{\global\randnumberstrue\setcounter{randnumber}{1}\global\lastsourcehasrandnumbersfalse}
\newcommand{\stoprandnumbers}{\global\randnumbersfalse\global\lastsourcehasrandnumbersfalse}
\newlength{\randnumberwidth}
\setlength{\randnumberwidth}{0.95cm}
\newlength{\randnumberrightoffset}
\setlength{\randnumberrightoffset}{0.85cm}
\newlength{\randnumberlawboxrightoffset}
\setlength{\randnumberlawboxrightoffset}{1.40cm}
\newcommand{\randnumberbox}[1]{%
  \leavevmode
  \makebox[0pt][l]{%
    \smash{%
      \tikz[remember picture,overlay,baseline=(rnmark.base)]{%
        \coordinate (rnmark) at (0,0);%
        \ifinlawbox
          \path let \p1=(current page.east), \p2=(rnmark.base) in
            node[
              anchor=base east,
              inner sep=0pt,
              outer sep=0pt,
              text width=\randnumberwidth,
              align=right,
              text=pendingink,
              font=\normalfont\mdseries\scriptsize\ttfamily
            ] at (\x1-\randnumberlawboxrightoffset,\y2) {{\normalfont\mdseries\scriptsize\ttfamily #1}};%
        \else
          \path let \p1=(current page.east), \p2=(rnmark.base) in
            node[
              anchor=base east,
              inner sep=0pt,
              outer sep=0pt,
              text width=\randnumberwidth,
              align=right,
              text=pendingink,
              font=\normalfont\mdseries\scriptsize\ttfamily
            ] at (\x1-\randnumberrightoffset,\y2) {{\normalfont\mdseries\scriptsize\ttfamily #1}};%
        \fi
      }%
    }%
  }%
  \ignorespaces
}
\newcommand{\lawboxstorerandnumber}[1]{%
  \gdef\lawboxpendingrn{#1}%
  \global\lawboxhaspendingrntrue
  \ignorespaces
}
\newcommand{\lawboxuserandnumber}{%
  \iflawboxhaspendingrn
    \randnumberbox{\lawboxpendingrn}%
    \global\lawboxhaspendingrnfalse
  \fi
}
\newcommand{\sourcern}[1]{%
  \ifrandnumbers
    \ifshoworiginal
      \ifshowtranslation\else
        \ifinlawbox\lawboxstorerandnumber{#1}\else\randnumberbox{#1}\fi
      \fi
    \fi
  \fi
}
\newcommand{\transrn}[1]{%
  \ifrandnumbers
    \ifshowtranslation
      \ifinlawbox\lawboxstorerandnumber{#1}\else\randnumberbox{#1}\fi
    \fi
  \fi
}
% ===== 課堂模式：德文原文縮寫註解 =====
% 日常切換只改各判決檔的一行：
%   \classroommodeon        開啟課堂版；註解掉這行即回到一般版。
% 模板預設：
%   1. 課堂模式關閉。
%   2. 課堂模式開啟時，縮寫註解包含「德文全稱＋中文譯名」。
%   3. 課堂模式開啟時，GG/條文編者註仍會自動顯示。
% 進階微調才需要在單案檔覆寫：
%   \classroomabbrzhoff     縮寫註解僅顯示德文全稱。
%   \showeditornotestrue    一般版也顯示編者註。
% 註解策略：同一實際 PDF 頁面中，同一縮寫只在第一次出現處加註。
% 這裡用 zref-abspage 的絕對頁序，不用 \thepage，避免文件內頁碼重置時誤判。
\newif\ifclassroommode
\newif\ifclassroomabbrzh
\classroommodefalse
\classroomabbrzhtrue
\newcommand{\classroommodeon}{\classroommodetrue}
\newcommand{\classroommodeoff}{\classroommodefalse}
\newcommand{\classroomabbrzhon}{\classroomabbrzhtrue}
\newcommand{\classroomabbrzhoff}{\classroomabbrzhfalse}
\newcommand{\abbrnotelabel}{\emph{Abk.}: }
\newcommand{\abbrnote}[3]{%
  \ifclassroommode
    \ifshoworiginal
      \ifintranslation\else
        \footnote{\normalfont\footnotesize\abbrnotelabel #2\ifclassroomabbrzh；{\songfont #3}\fi}%
      \fi
    \fi
  \fi
}
\newcommand{\abbrnoteonce}[4]{%
  #2%
  \ifclassroommode
    \ifshoworiginal
      \ifintranslation\else
        \begingroup
          \edef\abbrcurrentabspage{\number\numexpr\value{abspage}+1\relax}%
          \ifcsname abbrnoteseen@#1@\abbrcurrentabspage\endcsname\else
            \global\expandafter\let\csname abbrnoteseen@#1@\abbrcurrentabspage\endcsname\empty
            \abbrnote{#1}{#3}{#4}%
          \fi
        \endgroup
      \fi
    \fi
  \fi
}
\newcommand{\abbrAbs}{\abbrnoteonce{abs}{Abs.}{Absatz}{項}}
\newcommand{\abbrAAO}{\abbrnoteonce{aao}{a.a.O.}{am angegebenen Ort}{前引處}}
\newcommand{\abbrAE}{\abbrnoteonce{ae}{AE}{Alternativ-Entwurf}{替代草案}}
\newcommand{\abbrArt}{\abbrnoteonce{art}{Art.}{Artikel}{條}}
\newcommand{\abbrAufl}{\abbrnoteonce{aufl}{Aufl.}{Auflage}{版}}
\newcommand{\abbrBd}{\abbrnoteonce{bd}{Bd.}{Band}{卷}}
\newcommand{\abbrBGBl}{\abbrnoteonce{bgbl}{BGBl.}{Bundesgesetzblatt}{聯邦法律公報}}
\newcommand{\abbrBGHSt}{\abbrnoteonce{bghst}{BGHSt}{Entscheidungen des Bundesgerichtshofes in Strafsachen}{聯邦普通法院刑事判決彙編}}
\newcommand{\abbrBRDrucks}{\abbrnoteonce{brdrucks}{BRDrucks.}{Bundesrats-Drucksache}{聯邦參議院文件}}
\newcommand{\abbrBTDrucks}{\abbrnoteonce{btdrucks}{BTDrucks.}{Bundestags-Drucksache}{聯邦議院文件}}
\newcommand{\abbrBundesgesetzbl}{\abbrnoteonce{bundesgesetzbl}{Bundesgesetzbl.}{Bundesgesetzblatt}{聯邦法律公報}}
\newcommand{\abbrBvF}{\abbrnoteonce{bvf}{BvF}{Bundesverfassungsgerichtsverfahren}{聯邦憲法法院程序案號}}
\newcommand{\abbrBvL}{\abbrnoteonce{bvl}{BvL}{Bundesverfassungsgerichtsverfahren}{聯邦憲法法院程序案號}}
\newcommand{\abbrBvR}{\abbrnoteonce{bvr}{BvR}{Bundesverfassungsgerichtsverfahren}{聯邦憲法法院程序案號}}
\newcommand{\abbrBVerfG}{\abbrnoteonce{bverfg}{BVerfG}{Bundesverfassungsgericht}{聯邦憲法法院}}
\newcommand{\abbrBVerfGE}{\abbrnoteonce{bverfge}{BVerfGE}{Entscheidungen des Bundesverfassungsgerichts}{聯邦憲法法院判決集}}
\newcommand{\abbrBVerfGG}{\abbrnoteonce{bverfgg}{BVerfGG}{Bundesverfassungsgerichtsgesetz}{聯邦憲法法院法}}
\newcommand{\abbrDDR}{\abbrnoteonce{ddr}{DDR}{Deutsche Demokratische Republik}{德意志民主共和國；東德}}
\newcommand{\abbrDJT}{\abbrnoteonce{djt}{DJT}{Deutscher Juristentag}{德國法學家大會}}
\newcommand{\abbrDr}{\abbrnoteonce{dr}{Dr.}{Doktor}{Dr.}}
\newcommand{\abbrEuGRZ}{\abbrnoteonce{eugrz}{EuGRZ}{Europäische Grundrechte-Zeitschrift}{歐洲基本權期刊}}
\newcommand{\abbrEV}{\abbrnoteonce{ev}{EV}{Einigungsvertrag}{統一條約}}
\newcommand{\abbrF}{\abbrnoteonce{f}{f.}{folgende}{以下一頁；下一段}}
\newcommand{\abbrFF}{\abbrnoteonce{ff}{ff.}{fortfolgende}{以下數頁；以下各段}}
\newcommand{\abbrGBlDDR}{\abbrnoteonce{gblddr}{GBl. DDR}{Gesetzblatt der Deutschen Demokratischen Republik}{德意志民主共和國法律公報}}
\newcommand{\abbrGG}{\abbrnoteonce{gg}{GG}{Grundgesetz}{基本法}}
\newcommand{\abbrGVBl}{\abbrnoteonce{gvbl}{GVBl.}{Gesetz- und Verordnungsblatt}{法律及命令公報}}
\newcommand{\abbrJR}{\abbrnoteonce{jr}{JR}{Juristische Rundschau}{法學評論}}
\newcommand{\abbrJZ}{\abbrnoteonce{jz}{JZ}{JuristenZeitung}{法學家報}}
\newcommand{\abbrIVm}{\abbrnoteonce{ivm}{i.V.m.}{in Verbindung mit}{結合；連同}}
\newcommand{\abbrMdB}{\abbrnoteonce{mdb}{MdB}{Mitglied des Bundestages}{聯邦議院議員}}
\newcommand{\abbrMwN}{\abbrnoteonce{mwn}{m.w.N.}{mit weiteren Nachweisen}{附進一步文獻／判例出處}}
\newcommand{\abbrNF}{\abbrnoteonce{nf}{n.F.}{neue Fassung}{新版本}}
\newcommand{\abbrAF}{\abbrnoteonce{af}{a.F.}{alte Fassung}{舊版本}}
\newcommand{\abbrNr}{\abbrnoteonce{nr}{Nr.}{Nummer}{款、目或號，依語境}}
\newcommand{\abbrProf}{\abbrnoteonce{prof}{Prof.}{Professor}{教授}}
\newcommand{\abbrRdnr}{\abbrnoteonce{rdnr}{Rdnr.}{Randnummer}{邊碼；段碼}}
\newcommand{\abbrRdnrn}{\abbrnoteonce{rdnrn}{Rdnrn.}{Randnummern}{邊碼；段碼}}
\newcommand{\abbrRGBl}{\abbrnoteonce{rgbl}{RGBl.}{Reichsgesetzblatt}{帝國法律公報}}
\newcommand{\abbrRGSt}{\abbrnoteonce{rgst}{RGSt}{Entscheidungen des Reichsgerichts in Strafsachen}{帝國法院刑事判決彙編}}
\newcommand{\abbrRspr}{\abbrnoteonce{rspr}{Rspr.}{Rechtsprechung}{判例；司法實務}}
\newcommand{\abbrRVO}{\abbrnoteonce{rvo}{RVO}{Reichsversicherungsordnung}{帝國保險法}}
\newcommand{\abbrS}{\abbrnoteonce{s}{S.}{Seite}{頁}}
\newcommand{\abbrSFHG}{\abbrnoteonce{sfhg}{SFHG}{Schwangeren- und Familienhilfegesetz}{孕婦及家庭扶助法}}
\newcommand{\abbrSGBV}{\abbrnoteonce{sgbv}{SGB V}{Fünftes Buch Sozialgesetzbuch}{社會法典第五編}}
\newcommand{\abbrStAG}{\abbrnoteonce{staeg}{StÄG}{Strafrechtsänderungsgesetz}{刑法修正法}}
\newcommand{\abbrStenBer}{\abbrnoteonce{stenber}{StenBer.}{Stenographischer Bericht}{速記紀錄}}
\newcommand{\abbrStGB}{\abbrnoteonce{stgb}{StGB}{Strafgesetzbuch}{刑法}}
\newcommand{\abbrStREG}{\abbrnoteonce{streg}{StREG}{Strafrechtsreform-Ergänzungsgesetz}{刑法改革補充法}}
\newcommand{\abbrStrRG}{\abbrnoteonce{strrg}{StrRG}{Strafrechtsreformgesetz}{刑法改革法}}
\newcommand{\abbrUS}{\abbrnoteonce{us}{U.S.}{United States Reports}{美國最高法院判決彙編}}
\newcommand{\abbrVgl}{\abbrnoteonce{vgl}{vgl.}{vergleiche}{參見}}
\newcommand{\abbrVVDStRL}{\abbrnoteonce{vvdstrl}{VVDStRL}{Veröffentlichungen der Vereinigung der Deutschen Staatsrechtslehrer}{德國公法教師協會論文集}}
\newcommand{\abbrWP}{\abbrnoteonce{wp}{WP.}{Wahlperiode}{屆期}}
\newcommand{\abbrWp}{\abbrnoteonce{wp}{Wp.}{Wahlperiode}{屆期}}
\newcommand{\abbrZStrW}{\abbrnoteonce{zstrw}{ZStrW}{Zeitschrift für die gesamte Strafrechtswissenschaft}{整體刑法學期刊}}
\newcommand{\recordsourcerandstart}{%
  \ifrandnumbers
    \setcounter{lastsourcerandstart}{\value{randnumber}}%
    \global\lastsourcehasrandnumberstrue
  \else
    \global\lastsourcehasrandnumbersfalse
  \fi
}
\newlength{\biparagraphskip}
\setlength{\biparagraphskip}{0.52em}
\newlength{\lawboxblockbeforeskip}
\setlength{\lawboxblockbeforeskip}{0.72em}
\newlength{\lawboxblockafterskip}
\setlength{\lawboxblockafterskip}{0.92em}
\newif\iflawboxpartendedblock
\lawboxpartendedblockfalse
\newcommand{\sourceparstyle}{\footnotesize\color{sourceink}\setstretch{1.12}\RaggedRight\emergencystretch=3em\setlength{\parskip}{0.42em}\obeylines\selectfont}
\newcommand{\transparstyle}{\songfont\color{caseink}\setstretch{1.12}\setlength{\parindent}{0pt}\setlength{\parskip}{0.42em}}
\newcommand{\bilingualpairskip}{%
  \par\addvspace{\biparagraphskip}%
  \switchcolumn
  \par\addvspace{\biparagraphskip}%
  \switchcolumn*%
  \global\intranslationfalse
}
\newcommand{\bilingualboxpairskip}{%
  \par
  \switchcolumn
  \par
  \switchcolumn*%
  \global\intranslationfalse
}
\newcommand{\bilinguallawboxblockskip}{%
  \switchcolumn*%
  \par\addvspace{\lawboxblockafterskip}%
  \switchcolumn
  \par\addvspace{\lawboxblockafterskip}%
  \switchcolumn*%
  \global\intranslationfalse
}
\newcommand{\bikeepwithnext}[1][4]{%
  \ifbilingualcolumns
    \syncleft
    \Needspace{#1\baselineskip}%
    \switchcolumn
    \Needspace{#1\baselineskip}%
    \switchcolumn*%
    \global\intranslationfalse
  \else
    \Needspace{#1\baselineskip}%
  \fi
}
\NewEnviron{sourcepara}[1][]{
  \recordsourcerandstart
  \ifshoworiginal
    \ifbilingualcolumns
      \ifintranslation\switchcolumn*\global\intranslationfalse\fi
    \else
      \Needspace{5\baselineskip}
    \fi
    \ifbilingualcolumns\else\par\addvspace{0.35em}\fi
    {\sourceparstyle
    \noindent\BODY\par}
    \ifbilingualcolumns\else\addvspace{0.25em}\fi
    \ifbilingualcolumns
      \switchcolumn
      \global\intranslationtrue
    \fi
  \else
    \ifrandnumbers
      \begingroup
        \setbox0=\vbox{\hsize=\linewidth\sourceparstyle\noindent\BODY\par}%
      \endgroup
    \fi
  \fi
}
\NewEnviron{transpara}{
  \ifshowtranslation
    \setcounter{randnumberlocal}{\value{lastsourcerandstart}}%
    \ifbilingualcolumns
      \ifintranslation\else\switchcolumn\global\intranslationtrue\fi
      {\transparstyle\setlength{\leftskip}{0.55em}\setlength{\rightskip}{0.15em plus 1em}\addtolength{\linewidth}{-0.55em}\BODY\par}
      \switchcolumn*
      \bilingualpairskip
    \else
      {\transparstyle\BODY\par}
    \fi
  \else
    \ifbilingualcolumns
      \ifintranslation\switchcolumn*\global\intranslationfalse\fi
    \fi
  \fi
}
\NewEnviron{sourceboxpara}[1][]{
  \global\lawboxpartendedblockfalse
  \recordsourcerandstart
  \ifshoworiginal
    \ifbilingualcolumns
      \ifintranslation\switchcolumn*\global\intranslationfalse\fi
    \else
      \Needspace{3\baselineskip}
    \fi
    {\sourceparstyle
    \BODY\par}
    \ifbilingualcolumns\else
      \iflawboxpartendedblock\addvspace{\lawboxblockafterskip}\fi
    \fi
    \ifbilingualcolumns
      \switchcolumn
      \global\intranslationtrue
    \fi
  \fi
}
\NewEnviron{transboxpara}{
  \global\lawboxpartendedblockfalse
  \ifshowtranslation
    \setcounter{randnumberlocal}{\value{lastsourcerandstart}}%
    \ifbilingualcolumns
      \ifintranslation\else\switchcolumn\global\intranslationtrue\fi
      {\transparstyle\BODY\par}
      \iflawboxpartendedblock
        \bilinguallawboxblockskip
      \else
        \switchcolumn*
        \bilingualboxpairskip
      \fi
    \else
      {\transparstyle\BODY\par}
      \iflawboxpartendedblock\addvspace{\lawboxblockafterskip}\fi
    \fi
  \else
    \ifbilingualcolumns
      \ifintranslation\switchcolumn*\global\intranslationfalse\fi
    \fi
  \fi
}
\newcommand{\leitsatznum}[1]{\noindent\hangindent=3.0em\hangafter=1\makebox[2.25em][r]{\bfseries #1.}\hspace{0.55em}\ignorespaces}
\newcommand{\syncleft}{\ifbilingualcolumns\ifintranslation\switchcolumn*\global\intranslationfalse\fi\fi}
\pretocmd{\section}{\syncleft\clearpage}{}{}
\pretocmd{\subsection}{\syncleft}{}{}
\pretocmd{\subsubsection}{\syncleft}{}{}
\newcommand{\biheadingfont}{\songfont}
\newcommand{\bitocentry}[2]{%
  \ifshoworiginal
    \ifshowtranslation
      \ifstrequal{#1}{#2}{#1}{#1\space #2}%
    \else
      #1%
    \fi
  \else
    #2%
  \fi
}
\pdfstringdefDisableCommands{%
  \def\bitocentry#1#2{#1}%
}
\newcommand{\biheading}[7]{%
  \syncleft#1\bikeepwithnext[#3]\phantomsection\addcontentsline{toc}{#2}{\bitocentry{#6}{#7}}%
  \ifbilingualcolumns
    {#4 #6\par}%
    \switchcolumn
    {#4\biheadingfont #7\par}%
    \switchcolumn*%
    \global\intranslationfalse
    \par\addvspace{#5}%
    \switchcolumn
    \par\addvspace{#5}%
    \switchcolumn*%
  \else
    \ifshoworiginal{#4 #6\par}\fi
    \ifshowtranslation{#4\biheadingfont #7\par}\fi
    \addvspace{#5}%
  \fi
}
\newcommand{\termsectionheading}[1]{%
  \par\Needspace{9\baselineskip}%
  \vspace{0.65em}%
  {\songfont\bfseries #1\par}%
  \vspace{0.2em}%
}
% 章節換頁：
%   雙語 paracol 欄內不要用 \clearpage；它會在同步兩欄時吐出只有欄線/頁碼的空頁。
%   \newpage 足以讓下一個 section 起新頁，又不會強制清空所有待排材料。
%   單語模式不在 paracol 內，仍可用 \clearpage。
\newcommand{\bisectionpagebreak}{\ifbilingualcolumns\newpage\else\clearpage\fi}
\newcommand{\bisection}[2]{%
  \biheading{\iffirstbilingualsection\global\firstbilingualsectionfalse\else\bisectionpagebreak\fi}{section}{6}{\centering\Large\bfseries}{0.8em}{#1}{#2}%
}
\newcommand{\bisubsection}[2]{%
  \biheading{}{subsection}{5}{\large\bfseries}{0.55em}{#1}{#2}%
}
\newcommand{\bisubsubsection}[2]{%
  \biheading{}{subsubsection}{4}{\normalsize\bfseries}{0.45em}{#1}{#2}%
}
\newcommand{\bicompositeheading}[4]{%
  \biheading{}{subsection}{5}{\large\bfseries}{0.16em}{#1}{#2}%
  \biheading{}{subsubsection}{3}{\normalsize\bfseries}{0.45em}{#3}{#4}%
}
\newif\iflawboxfirstitem
\lawboxfirstitemfalse
\newcommand{\lawboxprespace}[1]{%
  \par
  \iflawboxfirstitem
    \global\lawboxfirstitemfalse
  \else
    \addvspace{#1}%
  \fi
}
\newtcolorbox{lawboxinner}{
  caseboxbase,
  colframe=sourceline!85!black,
  colback=white,
  left=1.05em,
  right=1.05em,
  top=0.62em,
  bottom=0.62em,
  before upper={\global\inlawboxtrue\global\lawboxfirstitemtrue\global\lawboxhaspendingrnfalse\ifintranslation\lawboxtransfont\else\lawboxsourcefont\fi\RaggedRight\setlength{\parindent}{0pt}\setlength{\parskip}{0pt}\ifintranslation\setstretch{\lawboxtransstretch}\else\setstretch{\lawboxsourcestretch}\fi},
  after upper={\global\lawboxhaspendingrnfalse\global\inlawboxfalse}
}
\tcbset{
  lawboxpartbase/.style={
    enhanced,
    breakable=false,
    width=0.985\linewidth,
    colback=white,
    frame hidden,
    boxrule=0pt,
    arc=0pt,
    left=1.05em,
    right=1.05em,
    boxsep=1.5mm,
    before skip=0pt,
    after skip=0pt,
    before upper={\global\inlawboxtrue\global\lawboxfirstitemtrue\global\lawboxhaspendingrnfalse\ifintranslation\lawboxtransfont\else\lawboxsourcefont\fi\RaggedRight\setlength{\parindent}{0pt}\setlength{\parskip}{0pt}\ifintranslation\setstretch{\lawboxtransstretch}\else\setstretch{\lawboxsourcestretch}\fi},
    after upper={\global\lawboxhaspendingrnfalse\global\inlawboxfalse},
    borderline west={0.28pt}{0pt}{sourceline!85!black},
    borderline east={0.28pt}{0pt}{sourceline!85!black},
  },
  lawboxpartfirst/.style={
    lawboxpartbase,
    top=0.62em,
    bottom=0.04em,
    before skip=\lawboxblockbeforeskip,
    borderline north={0.28pt}{0pt}{sourceline!85!black},
  },
  lawboxpartmiddle/.style={
    lawboxpartbase,
    top=0.04em,
    bottom=0.04em,
  },
  lawboxpartlast/.style={
    lawboxpartbase,
    top=0.04em,
    bottom=0.62em,
    after skip=0pt,
    borderline south={0.28pt}{0pt}{sourceline!85!black},
  }
}
\newtcolorbox{lawboxpartinner}[1]{#1}
\newtcolorbox{termboxinner}{
  caseboxbase,
  colframe=noteblue!70!black,
  colback=casepaper,
  before upper={\songfont\setlength{\parindent}{0pt}\setlength{\parskip}{0.35em}}
}
\newtcolorbox{deboxinner}{
  caseboxbase,
  colframe=notegray!72!black,
  colback=white,
  left=1.05em,
  right=1.05em,
  top=0.62em,
  bottom=0.62em,
  before upper={\global\inlawboxtrue\global\lawboxfirstitemtrue\global\lawboxhaspendingrnfalse\small\setlength{\parindent}{0pt}\setlength{\parskip}{0.35em}},
  after upper={\global\lawboxhaspendingrnfalse\global\inlawboxfalse}
}
\NewEnviron{lawbox}{
  \begin{lawboxinner}\BODY\end{lawboxinner}
}
\NewEnviron{lawboxpart}[2]{
  \begin{lawboxpartinner}{lawboxpart#1,equal height group=#2}\BODY\end{lawboxpartinner}%
  \ifstrequal{#1}{last}{\global\lawboxpartendedblocktrue}{}%
}
\NewEnviron{termbox}{
  \par\addvspace{0.55em}
  {\color{noteblue!70!black}\rule{\linewidth}{0.28pt}}\par
  \begingroup
    \songfont\small\color{caseink}
    \setlength{\parindent}{0pt}
    \setlength{\parskip}{0.24em}
    \BODY
    \par
  \endgroup
  {\color{noteblue!70!black}\rule{\linewidth}{0.28pt}}\par
  \addvspace{0.55em}
}
\NewEnviron{debox}{
  \begin{deboxinner}\BODY\end{deboxinner}
}
\providecommand{\paragraphmark}[1]{}
\newcommand{\lawboxtitleafter}{\addvspace{0.06em}}
\newcommand{\lawtitle}[1]{\lawboxprespace{0.22em}{\lawboxtransfont\bfseries\lawboxuserandnumber #1}\par\lawboxtitleafter}
\newcommand{\absno}[2]{\lawboxprespace{0.04em}\noindent\lawboxuserandnumber\hangindent=2.6em\hangafter=1\makebox[2.6em][l]{(#1)}#2\par}
\newcommand{\nrno}[2]{\lawboxprespace{0pt}\noindent\lawboxuserandnumber\hspace*{2.6em}\hangindent=4.4em\hangafter=1\makebox[1.8em][l]{#1.}#2\par}
\newcommand{\letterno}[2]{\lawboxprespace{0pt}\noindent\lawboxuserandnumber\hspace*{4.4em}\hangindent=6.0em\hangafter=1\makebox[1.6em][l]{#1)}#2\par}
\newcommand{\sourcelawtitle}[1]{\lawboxprespace{0.22em}{\lawboxsourcefont\bfseries\lawboxuserandnumber #1}\par\lawboxtitleafter}
\newcommand{\sourcelawsubtitle}[1]{\lawboxprespace{0.04em}{\lawboxsourcefont\bfseries\lawboxuserandnumber #1}\par\lawboxtitleafter}
\newcommand{\sourceabsno}[2]{\lawboxprespace{0.04em}\noindent\lawboxuserandnumber\hangindent=2.45em\hangafter=1\makebox[2.45em][l]{(#1)}#2\par}
\newcommand{\sourcenrno}[2]{\lawboxprespace{0pt}\noindent\lawboxuserandnumber\hspace*{2.45em}\hangindent=4.25em\hangafter=1\makebox[1.8em][l]{#1.}#2\par}
\newcommand{\sourceletterno}[2]{\lawboxprespace{0pt}\noindent\lawboxuserandnumber\hspace*{4.25em}\hangindent=5.85em\hangafter=1\makebox[1.6em][l]{#1)}#2\par}
\newcommand{\sourcequotepara}[1]{\lawboxprespace{0.04em}\noindent\lawboxuserandnumber\hangindent=1.2em\hangafter=1 #1\par}
\newcommand{\sourcelawline}[1]{\lawboxprespace{0.04em}\noindent\lawboxuserandnumber #1\par}
\newcommand{\lawline}[1]{\lawboxprespace{0.04em}\noindent\lawboxuserandnumber #1\par}
\newcommand{\pendingtranslation}{\par{\songfont\footnotesize\color{pendingink}（中文譯文待補。）}\par}
\newcommand{\tocpart}[1]{\par\addvspace{0.52em}{\footnotesize\bfseries\color{pendingink}#1\par}\addvspace{0.16em}}
\newcommand{\tocdashline}{\par\addvspace{0.48em}\noindent{\color{notegray}\xleaders\hbox to 0.82em{\hfil\rule{0.42em}{0.28pt}\hfil}\hskip\linewidth}\par\addvspace{0.38em}}
% \tocpanel{font-cmd}{content}：將目錄置中於所在欄寬（雙語欄適用）。
% A3 橫式使用 0.58\linewidth，A4 橫式使用 0.84\linewidth，
% 使兩種紙張下目錄欄內容實際寬度相近（均約 100–105 mm）。
\newcommand{\tocpanel}[2]{%
  \makebox[\linewidth][c]{%
    \ifx\bilingualpaper\bilingualpaperaThree
      \begin{minipage}{0.58\linewidth}%
    \else
      \begin{minipage}{0.84\linewidth}%
    \fi
    \toccolstyle
    #1#2%
    \end{minipage}%
  }%
}
% ===== 首頁簡目錄的兩層 description list 縮排 =====
% enumitem 的 description list 可把每一列拆成：
%   [label]  labelsep  body text
% 這裡的「起點」都以所在 minipage/欄位的左緣為 0。
%
% 核心規則：
%   labelindent = label 欄位從左緣開始的位置。
%   labelwidth  = label 欄位本身寬度；標籤太長（如 Abw. Meinung）就加大它。
%   labelsep    = label 欄位右緣到正文的距離。
%   leftmargin  = 正文 body text 的起點。判斷層級視覺時主要看這個。
%
% 所以：
%   想讓「Begründung / 判決理由」整列正文往右或往左：改 tocitems 的 leftmargin。
%   想讓「A. Verfahren und Gegenstand / A. 程序與審查標的」正文往右或往左：
%     改 tocsubitems 的 leftmargin。
%   想讓 A./B./C. 這些標籤本身往右或往左，但正文不動：改 tocsubitems 的 labelindent。
%   想讓 A./B./C. 與正文間距變大或變小：改 tocsubitems 的 labelsep。
%
% 層級原則：
%   tocsubitems.leftmargin 必須大於 tocitems.leftmargin，
%   否則子層級正文會看起來比「Gründe / 理由」還靠左。
\newlist{tocitems}{description}{1}
\setlist[tocitems]{%
  labelindent=0pt,    % 第一層 label 欄位起點；0pt 表示貼齊目錄欄左緣。
  leftmargin=2.54cm,  % 第一層正文起點；例如 Begründung / 判決理由 的左緣。
  labelwidth=2.16cm,  % 第一層 label 欄位寬度；需容納 Leitsätze、Abw. Meinung、不同意見等。
  labelsep=0.32cm,    % 第一層 label 與正文的水平間距；只改間距，不改正文起點。
  itemsep=0.28em,     % 第一層各項之間的垂直距離。
  topsep=0.20em,      % 第一層 list 上下與前後內容的垂直距離。
  parsep=0pt,         % 同一 item 內段落之間的距離；目錄項通常不需要。
  partopsep=0pt,      % list 若緊接段落開始時的額外距離；關閉以免目錄高度飄動。
  align=left,         % label 欄內左對齊；長短標籤不靠右貼齊。
  font=\normalfont\bfseries % label 的字型；正文不受這行影響。
}
\newlist{tocsubitems}{description}{1}
\setlist[tocsubitems]{%
  labelindent=1cm,    % 第二層 label 起點；只控制 A./B./C. 本身的位置。
  leftmargin=3.00cm,  % 第二層正文起點；必須大於 2.54cm，才會比 Begründung / 判決理由更右。
  labelwidth=0.5cm,   % 第二層 label 欄寬；A./B./C. 很短，可小於第一層。
  labelsep=1.5cm,    % 第二層 label 與正文間距；A. 與 Verfahren / 程序 的距離看這裡。
  itemsep=0.12em,     % 第二層各項較密，避免 A.–E. 佔太高。
  topsep=0.04em,      % 第二層 list 與 Gründe / 理由之間的垂直距離。
  parsep=0pt,         % 同一第二層 item 內段落距離；目錄項通常不需要。
  partopsep=0pt,      % 關閉額外段落起始距離，避免版面不穩。
  align=left,         % A./B./C. label 欄內左對齊。
  font=\normalfont\bfseries, % A./B./C. label 加粗；正文不受這行影響。
  before={}           % 不在第二層 list 前自動插入額外內容。
}
\newlist{termitems}{description}{1}
\ifbilingualcolumns
  \setlist[termitems]{%
    leftmargin=5.2cm,
    labelwidth=4.8cm,
    labelsep=0.35cm,
    itemsep=0.16em,
    topsep=0.2em,
    parsep=0pt,
    partopsep=0pt,
    font=\normalfont\bfseries,
    style=nextline
  }
\else
  \setlist[termitems]{%
    leftmargin=0pt,
    labelwidth=0pt,
    labelsep=0pt,
    itemsep=0.24em,
    topsep=0.2em,
    parsep=0pt,
    partopsep=0pt,
    font=\normalfont\bfseries,
    style=nextline
  }
\fi
\newlist{abbritems}{description}{1}
\setlist[abbritems]{%
  leftmargin=2.38cm,
  labelwidth=2.02cm,
  labelsep=0.28cm,
  itemsep=0.16em,
  topsep=0.2em,
  parsep=0pt,
  partopsep=0pt,
  align=left,
  font=\normalfont\bfseries
}
\ifbilingualcolumns
  \newenvironment{termcolumns}{\begin{multicols}{2}}{\end{multicols}}
\else
  \newenvironment{termcolumns}{}{}
\fi

% ===== 編者註腳開關 =====
% 一般版預設不顯示編者註，避免正文腳註過密。
% 若開啟 \classroommodeon，GG/條文編者註仍會自動顯示，無需另開本開關。
% 若一般版也要顯示編者註，可在單案檔 \input 後加 \showeditornotestrue。
\newif\ifshoweditornotes
\showeditornotesfalse

% ===== 目錄排版輔助 =====
\newcommand{\toccolstyle}{\raggedright\arraybackslash}
\newcommand{\tocsingleindent}{\leftskip=1.45cm\rightskip=1.2cm}

% ===== 行內引文環境（德文原文與中文譯文各一，帶左側灰色邊線）=====
\newcommand{\quoteparbreak}{\par\addvspace{0.48em}}
\newcommand{\sourcequoteinner}[1]{%
  \begin{tcolorbox}[
    enhanced, breakable, width=\dimexpr\linewidth-0.8em\relax, frame hidden,
    colback=white, coltext=caseink, boxrule=0pt,
    borderline west={0.55pt}{0.88em}{notegray},
    left=1.78em, right=0.72em, top=0.18em, bottom=0.18em,
    boxsep=0pt, before skip=0.24em, after skip=0.24em,
    before upper={\normalfont\footnotesize\color{caseink}\setlength{\parindent}{0pt}\setlength{\parskip}{0.34em}\raggedright\setlength{\rightskip}{0pt plus 1em}}
  ]%
  #1\par
  \end{tcolorbox}%
}
\newcommand{\transquoteinner}[1]{%
  \begin{tcolorbox}[
    enhanced, breakable, width=\dimexpr\linewidth-0.8em\relax, frame hidden,
    colback=white, coltext=caseink, boxrule=0pt,
    borderline west={0.55pt}{0.88em}{notegray},
    left=1.78em, right=0.72em, top=0.18em, bottom=0.18em,
    boxsep=0pt, before skip=0.24em, after skip=0.24em,
    before upper={\songfont\small\color{caseink}\setlength{\parindent}{0pt}\setlength{\parskip}{0.34em}\raggedright\setlength{\rightskip}{0pt plus 1em}}
  ]%
  #1\par
  \end{tcolorbox}%
}

% ===== 大綱（Gliederung）Rn 輔助命令 =====
\newcommand{\outrn}[2]{\enspace{\normalfont\footnotesize\color{pendingink}(Rn\,#1–#2)}}
\newcommand{\outrnsingle}[1]{\enspace{\normalfont\footnotesize\color{pendingink}(Rn\,#1)}}
\newcommand{\outrnzh}[2]{\enspace{\normalfont\footnotesize\color{pendingink}（第\,#1–#2\,段）}}
\newcommand{\outrnzhs}[1]{\enspace{\normalfont\footnotesize\color{pendingink}（第\,#1\,段）}}
% ===== 大綱雙語對照表格輔助命令（三欄：段碼 │ 德文 │ 中文）=====
% 欄 1：段碼（由欄定義自動格式化：小字、等寬、右對齊、pendingink）
% 欄 2：德文標籤＋正文
% 欄 3：中文標籤＋正文
%
% 無段碼的列（\omainrow、\osecrow），欄 1 以 & 空置。
% 有段碼的列（\osecrowrn、\osubrow、\ointrow），#1 為預格式化字串
%   例：\osubrow{2–49}{I.}{...}{I.}{...}
%       \ointrow{107}{...}{...}
%
% \opartrow：段組標題（橫跨三欄）
\newcommand{\opartrow}[2]{%
  \noalign{\vspace{0.30em}}%
  \multicolumn{3}{@{}l@{}}{%
    \footnotesize\bfseries\color{pendingink}#1\hfill\songfont#2%
  }\\[0.06em]%
}
% \odashrow：虛線分隔（橫跨三欄）
\newcommand{\odashrow}{%
  \noalign{\vspace{0.28em}}%
  \multicolumn{3}{@{}l@{}}{\color{notegray}%
    \xleaders\hbox to 0.82em{\hfil\rule{0.42em}{0.28pt}\hfil}\hskip 0.88\linewidth}%
  \\[0.24em]%
}
% \odissentpart：不同意見書區段標頭。比一般 \opartrow 更像附錄性轉場。
\newcommand{\odissentpart}[2]{%
  \noalign{\vspace{0.42em}}%
  \multicolumn{3}{@{}p{0.88\textwidth}@{}}{%
    \color{notegray}\rule{\linewidth}{0.18pt}%
  }\\[-0.18em]%
  \multicolumn{3}{@{}p{0.88\textwidth}@{}}{%
    \makebox[\linewidth][s]{%
      \footnotesize\bfseries\color{pendingink}#1%
      \hfill{\songfont #2}%
    }%
  }\\[-0.02em]%
}
% \odissentopinion：不同意見書的署名與一行摘要，右欄保留段碼範圍。
\newcommand{\odissentopinion}[5]{%
  \footnotesize\hangindent=0.52cm\hangafter=1%
    {\bfseries\color{caseink}#2}\enspace{\color{notegray}\textperiodcentered}\enspace\color{caseink}#3%
  &\footnotesize\hangindent=0.52cm\hangafter=1%
    {\songfont\bfseries\color{caseink}#4}\enspace{\color{notegray}\textperiodcentered}\enspace\songfont\color{caseink}#5%
  &#1%
  \\[0.12em]%
}
% \omainrow：頂層項目，無段碼（Leitsätze、Urteil、Gründe …）
\newcommand{\omainrow}[4]{%
  \footnotesize\hangindent=1.92cm\hangafter=1%
    \makebox[1.92cm][l]{\bfseries\color{caseink}#1}\color{caseink}#2%
  &\footnotesize\hangindent=1.92cm\hangafter=1%
    \makebox[1.92cm][l]{\bfseries\color{caseink}#3}\songfont\color{caseink}#4%
  &%
  \\[0.15em]%
}
% \omainrowrn：頂層項目，有段碼（不同意見等標籤寬於 \osecrow 框者）
\newcommand{\omainrowrn}[5]{%
  \footnotesize\hangindent=1.92cm\hangafter=1%
    \makebox[1.92cm][l]{\bfseries\color{caseink}#2}\color{caseink}#3%
  &\footnotesize\hangindent=1.92cm\hangafter=1%
    \makebox[1.92cm][l]{\bfseries\color{caseink}#4}\songfont\color{caseink}#5%
  &#1%
  \\[0.15em]%
}
% \osecrow：一級子項，無段碼（A.、C.、D. 等有子項者）
\newcommand{\osecrow}[4]{%
  \footnotesize\hspace{1.10cm}\hangindent=1.98cm\hangafter=1%
    \makebox[0.86cm][l]{\bfseries\color{caseink}#1}\color{caseink}#2%
  &\footnotesize\hspace{1.10cm}\hangindent=1.98cm\hangafter=1%
    \makebox[0.86cm][l]{\bfseries\color{caseink}#3}\songfont\color{caseink}#4%
  &%
  \\[0.11em]%
}
% \osecrowrn：一級子項，有段碼（B.、E. 等完整段落範圍者）
% #1=段碼字串（如 101–106）, #2=DE標籤, #3=DE正文, #4=ZH標籤, #5=ZH正文
\newcommand{\osecrowrn}[5]{%
  \footnotesize\hspace{1.10cm}\hangindent=1.98cm\hangafter=1%
    \makebox[0.86cm][l]{\bfseries\color{caseink}#2}\color{caseink}#3%
  &\footnotesize\hspace{1.10cm}\hangindent=1.98cm\hangafter=1%
    \makebox[0.86cm][l]{\bfseries\color{caseink}#4}\songfont\color{caseink}#5%
  &#1%
  \\[0.11em]%
}
% \osubrow：二級子項，有段碼範圍（I.、II.、A.I. …）
% #1=段碼字串, #2=DE標籤, #3=DE正文, #4=ZH標籤, #5=ZH正文
\newcommand{\osubrow}[5]{%
  \footnotesize\hspace{1.40cm}\hangindent=2.30cm\hangafter=1%
    \makebox[0.90cm][l]{\bfseries\color{caseink}#2}\color{caseink}#3%
  &\footnotesize\hspace{1.40cm}\hangindent=2.30cm\hangafter=1%
    \makebox[0.90cm][l]{\bfseries\color{caseink}#4}\songfont\color{caseink}#5%
  &#1%
  \\[0.09em]%
}
% \ointrow：引入段落，有單一段碼，無標籤
% #1=段碼字串, #2=DE正文, #3=ZH正文
\newcommand{\ointrow}[3]{%
  \footnotesize\hspace{0.52cm}\hangindent=0.52cm\hangafter=1\color{caseink}#2%
  &\footnotesize\hspace{0.52cm}\hangindent=0.52cm\hangafter=1\songfont\color{caseink}#3%
  &#1%
  \\[0.09em]%
}
% \osubsubrow：三級子項（1.、2.、3.）
% #1=段碼字串, #2=DE標籤, #3=DE正文, #4=ZH標籤, #5=ZH正文
\newcommand{\osubsubrow}[5]{%
  \footnotesize\hspace{1.80cm}\hangindent=2.48cm\hangafter=1%
    \makebox[0.68cm][l]{\bfseries\color{caseink}#2}\color{caseink}#3%
  &\footnotesize\hspace{1.80cm}\hangindent=2.48cm\hangafter=1%
    \makebox[0.68cm][l]{\bfseries\color{caseink}#4}\songfont\color{caseink}#5%
  &#1%
  \\[0.08em]%
}
% \osubsubsubrow：四級子項（a）、b））
% #1=段碼字串, #2=DE標籤, #3=DE正文, #4=ZH標籤, #5=ZH正文
\newcommand{\osubsubsubrow}[5]{%
  \footnotesize\hspace{2.10cm}\hangindent=2.68cm\hangafter=1%
    \makebox[0.58cm][l]{\bfseries\color{caseink}#2}\color{caseink}#3%
  &\footnotesize\hspace{2.10cm}\hangindent=2.68cm\hangafter=1%
    \makebox[0.58cm][l]{\bfseries\color{caseink}#4}\songfont\color{caseink}#5%
  &#1%
  \\[0.07em]%
}
% \osubsubsubsubrow：五級子項（aa）、bb））
% 標籤框 0.62cm：足以容納「aa)」在 footnotesize bold 下（約 0.50cm）並留有間距
% #1=段碼字串, #2=DE標籤, #3=DE正文, #4=ZH標籤, #5=ZH正文
\newcommand{\osubsubsubsubrow}[5]{%
  \footnotesize\hspace{2.38cm}\hangindent=3.06cm\hangafter=1%
    \makebox[0.68cm][l]{\bfseries\color{caseink}#2}\color{caseink}#3%
  &\footnotesize\hspace{2.38cm}\hangindent=3.06cm\hangafter=1%
    \makebox[0.68cm][l]{\bfseries\color{caseink}#4}\songfont\color{caseink}#5%
  &#1%
  \\[0.06em]%
}
% \outlinepage：雙語對照大綱頁包裝環境（longtable 自動跨頁）
% 三欄：德文欄 + 中文欄 + 段碼欄（最右，不折行）
% 段碼欄寬度：1.80cm，足以容納最長段碼如「309–348」
% DE/ZH 欄寬：(0.87\textwidth - 2.08cm) / 2
%   驗算：2×欄 + 0.05\textwidth + 0.28cm gap + 1.80cm = 0.87tw - 2.08cm + 0.05tw + 2.08cm = 0.92tw ✓
\newenvironment{outlinepage}{%
  \vspace*{0.40cm}%
  \setlength{\LTleft}{0.04\textwidth}%
  \setlength{\LTright}{0.04\textwidth}%
  \setlength{\tabcolsep}{0pt}%
  % 大綱列距由 longtable 的 \arraystretch 與各 row macro 結尾的 \\[...em] 共同決定。
  % 這裡控制「每一列本身的表格 strut 高度」；數值越大，A./I./1. 各項之間越像有空行。
  % A3 原本用 0.62，視覺上沒有空行；A4 也沿用同值，避免切到 A4 後項目被撐開。
  % 若日後覺得大綱太擠，只微調這個值（例如 0.68），不要改各 \omainrow/\osubrow 的縮排。
  \renewcommand{\arraystretch}{0.62}%
  \begin{longtable}{%
    >{\vspace{0pt}\raggedright\arraybackslash}p{\dimexpr(0.87\textwidth-2.08cm)/2\relax}%
    @{\hspace{0.025\textwidth}}%
    !{\color{notegray}\vrule width 0.12pt}%
    @{\hspace{0.025\textwidth}}%
    >{\vspace{0pt}\raggedright\arraybackslash}p{\dimexpr(0.87\textwidth-2.08cm)/2\relax}%
    @{\hspace{0.28cm}}%
    >{\vspace{0pt}\footnotesize\normalfont\color{pendingink}\raggedright\arraybackslash}p{1.80cm}%
    @{}%
  }%
  % 首頁欄標籤：不要橫跨置中；分別放在德文欄與中文欄的實際欄位起點。
  \footnotesize\bfseries\color{sourceink}Gliederung%
  &\footnotesize\bfseries\songfont\color{sourceink}大綱%
  &%
  \\[0.36em]%
  \endfirsthead%
  % 續頁欄標籤：同樣分欄定位，以灰色細體區分；無需標注「續」。
  \footnotesize\color{pendingink}Gliederung%
  &\footnotesize\songfont\color{pendingink}大綱%
  &%
  \\[0.24em]%
  \endhead%
}{%
  \end{longtable}%
}


% 課堂模式：取消下行註解即開啟；註解掉即回到一般版。
% 開啟時會在德文原文縮寫後加註，並自動顯示 GG/條文編者註。
% \classroommodeon

\newcommand{\documentversion}{v1.1}
\newcommand{\ggversionde}{\ifggversionnotede\else GG-Fassung: 28. Mai 1993, nach der 38. GG-Aenderung (BGBl. I 1992, S. 2086) und vor der 39. GG-Aenderung (BGBl. I 1993, S. 1002). \global\ggversionnotedetrue\fi}
\newcommand{\ggversionzh}{\ifggversionnotezh\else 以下引用《基本法》1993年5月28日判決時有效版本；即第38次《基本法》修正（BGBl. I 1992, S. 2086）後、第39次《基本法》修正（BGBl. I 1993, S. 1002）前之版本。 \global\ggversionnotezhtrue\fi}
% bverfge_gg.tex — 德國墮胎案 GG 條文腳註共用定義
% 使用方式：在各案 .tex 中先定義 \ggversionde/\ggversionzh，再 % bverfge_gg.tex — 德國墮胎案 GG 條文腳註共用定義
% 使用方式：在各案 .tex 中先定義 \ggversionde/\ggversionzh，再 % bverfge_gg.tex — 德國墮胎案 GG 條文腳註共用定義
% 使用方式：在各案 .tex 中先定義 \ggversionde/\ggversionzh，再 \input{../../共用LaTeX/bverfge_gg}

% ===== GG 版本旗標（\ggversionde/zh 由各案在 \input{../../共用LaTeX/bverfge_gg} 之前定義）=====
\newif\ifggversionnotede
\newif\ifggversionnotezh

% ===== 編者註腳基礎設施（首次標記模式）=====
\newif\ifeditornotelabelde
\newif\ifeditornotelabelzh
\newcommand{\editornotelabelde}{\ifeditornotelabelde\else\emph{Hrsg.-Anm. (nachfolgend ohne Kennzeichnung)}: \global\editornotelabeldetrue\fi}
\newcommand{\editornotelabelzh}{\ifeditornotelabelzh\else 編者註（下同）：\global\editornotelabelzhtrue\fi}
\newcommand{\editornotede}[1]{\ifshoweditornotes\footnote{\normalfont\footnotesize\editornotelabelde #1}\else\ifclassroommode\footnote{\normalfont\footnotesize\editornotelabelde #1}\fi\fi}
\newcommand{\editornotezh}[1]{\ifshoweditornotes\footnote{\normalfont\footnotesize\editornotelabelzh #1}\else\ifclassroommode\footnote{\normalfont\footnotesize\editornotelabelzh #1}\fi\fi}

% ===== GG 引用系統（\ggversionde/zh 由各案在 \input 前定義）=====
\newcommand{\ggmarknotede}[1]{\global\expandafter\let\csname ggnoteseen@de@#1\endcsname\empty}
\newcommand{\ggmarknotezh}[1]{\global\expandafter\let\csname ggnoteseen@zh@#1\endcsname\empty}
\newcommand{\ggnotedeonce}[2]{\ifcsname ggnoteseen@de@#1\endcsname\else\global\expandafter\let\csname ggnoteseen@de@#1\endcsname\empty\ggnotede{#2}\fi}
\newcommand{\ggnotezhonce}[2]{\ifcsname ggnoteseen@zh@#1\endcsname\else\global\expandafter\let\csname ggnoteseen@zh@#1\endcsname\empty\ggnotezh{#2}\fi}
\newcommand{\ggnotede}[1]{\editornotede{\ggversionde #1}}
\newcommand{\ggnotezh}[1]{\editornotezh{\ggversionzh #1}}

% ===== Art. 1 =====
\newcommand{\ggfnArtOneOneDE}{\ggnotedeonce{art-1-1}{Art. 1 Abs. 1: ``Die Wuerde des Menschen ist unantastbar. Sie zu achten und zu schuetzen ist Verpflichtung aller staatlichen Gewalt.''}}
\newcommand{\ggfnArtOneOneZH}{\ggnotezhonce{art-1-1}{第1條第1項：「人之尊嚴不可侵犯。尊重及保護之，為一切國家權力之義務。」}}
\newcommand{\ggfnArtOneOneTwoDE}{\ggnotedeonce{art-1-1-2}{Art. 1 Abs. 1 Satz 2: ``Sie zu achten und zu schuetzen ist Verpflichtung aller staatlichen Gewalt.''}}
\newcommand{\ggfnArtOneOneTwoZH}{\ggnotezhonce{art-1-1-2}{第1條第1項第2句：「尊重及保護之，為一切國家權力之義務。」}}

% ===== Art. 2 =====
\newcommand{\ggfnArtTwoOneDE}{\ggnotedeonce{art-2-1}{Art. 2 Abs. 1: ``Jeder hat das Recht auf die freie Entfaltung seiner Persoenlichkeit, soweit er nicht die Rechte anderer verletzt und nicht gegen die verfassungsmaessige Ordnung oder das Sittengesetz verstoesst.''}}
\newcommand{\ggfnArtTwoOneZH}{\ggnotezhonce{art-2-1}{第2條第1項：「人人有自由發展其人格之權利，但不得侵害他人權利，亦不得違反合憲秩序或道德律。」}}
% 簡易預設版（墮胎一案以 \renewcommand 覆寫為含條件邏輯之版本）
\newcommand{\ggfnArtTwoTwoDE}{\ggnotedeonce{art-2-2}{Art. 2 Abs. 2: ``Jeder hat das Recht auf Leben und koerperliche Unversehrtheit. Die Freiheit der Person ist unverletzlich. In diese Rechte darf nur auf Grund eines Gesetzes eingegriffen werden.''}\ggmarknotede{art-2-2-1}}
\newcommand{\ggfnArtTwoTwoZH}{\ggnotezhonce{art-2-2}{第2條第2項：「人人有生命與身體完整性之權利。人身自由不可侵犯。此等權利僅得依法律加以干預。」}\ggmarknotezh{art-2-2-1}}
\newcommand{\ggfnArtTwoTwoOneDE}{\ggnotedeonce{art-2-2-1}{Art. 2 Abs. 2 Satz 1: ``Jeder hat das Recht auf Leben und koerperliche Unversehrtheit.''}}
\newcommand{\ggfnArtTwoTwoOneZH}{\ggnotezhonce{art-2-2-1}{第2條第2項第1句：「人人有生命與身體完整性之權利。」}}
\newcommand{\ggfnArtTwoTwoThreeDE}{\ggnotedeonce{art-2-2-3}{Art. 2 Abs. 2 Satz 3: ``In diese Rechte darf nur auf Grund eines Gesetzes eingegriffen werden.''}}
\newcommand{\ggfnArtTwoTwoThreeZH}{\ggnotezhonce{art-2-2-3}{第2條第2項第3句：「此等權利僅得依法律加以干預。」}}
% 墮胎一案特有：第2條第2項第2–3句（Art. 2(2) 去掉第1句後的「餘文」）
\newcommand{\ggfnArtTwoTwoRestDE}{\ggnotedeonce{art-2-2-rest}{Art. 2 Abs. 2 Saetze 2-3: ``Die Freiheit der Person ist unverletzlich. In diese Rechte darf nur auf Grund eines Gesetzes eingegriffen werden.''}\ggmarknotede{art-2-2}}
\newcommand{\ggfnArtTwoTwoRestZH}{\ggnotezhonce{art-2-2-rest}{第2條第2項第2、3句：「人身自由不可侵犯。此等權利僅得依法律加以干預。」}\ggmarknotezh{art-2-2}}

% ===== Art. 3 =====
\newcommand{\ggfnArtThreeDE}{\ggnotedeonce{art-3}{Art. 3: ``Alle Menschen sind vor dem Gesetz gleich. Maenner und Frauen sind gleichberechtigt. Niemand darf wegen seines Geschlechtes, seiner Abstammung, seiner Rasse, seiner Sprache, seiner Heimat und Herkunft, seines Glaubens, seiner religioesen oder politischen Anschauungen benachteiligt oder bevorzugt werden.''}}
\newcommand{\ggfnArtThreeZH}{\ggnotezhonce{art-3}{第3條：「所有人在法律前一律平等。男女享有平等權利。任何人不得因其性別、血統、種族、語言、故鄉及出身、信仰、宗教或政治見解而受不利益或優惠。」}}
\newcommand{\ggfnArtThreeTwoDE}{\ggnotedeonce{art-3-2}{Art. 3 Abs. 2: ``Maenner und Frauen sind gleichberechtigt.''}}
\newcommand{\ggfnArtThreeTwoZH}{\ggnotezhonce{art-3-2}{第3條第2項：「男女享有平等權利。」}}

% ===== Art. 4 =====
\newcommand{\ggfnArtFourOneDE}{\ggnotedeonce{art-4-1}{Art. 4 Abs. 1: ``Die Freiheit des Glaubens, des Gewissens und die Freiheit des religioesen und weltanschaulichen Bekenntnisses sind unverletzlich.''}}
\newcommand{\ggfnArtFourOneZH}{\ggnotezhonce{art-4-1}{第4條第1項：「信仰自由、良心自由，以及宗教與世界觀信念自由，不可侵犯。」}}

% ===== Art. 5 =====
\newcommand{\ggfnArtFiveOneDE}{\ggnotedeonce{art-5-1}{Art. 5 Abs. 1: ``Jeder hat das Recht, seine Meinung in Wort, Schrift und Bild frei zu aeussern und zu verbreiten und sich aus allgemein zugaenglichen Quellen ungehindert zu unterrichten. Die Pressefreiheit und die Freiheit der Berichterstattung durch Rundfunk und Film werden gewaehrleistet. Eine Zensur findet nicht statt.''}}
\newcommand{\ggfnArtFiveOneZH}{\ggnotezhonce{art-5-1}{第5條第1項：「人人有以言語、文字及圖像自由表達及傳播其意見，並由一般可接近來源不受阻礙獲取資訊之權利。新聞自由及廣播、電影報導自由受保障。不設審查制度。」}}

% ===== Art. 6 =====
\newcommand{\ggfnArtSixDE}{\ggnotedeonce{art-6}{Art. 6 Abs. 1: ``Ehe und Familie stehen unter dem besonderen Schutze der staatlichen Ordnung.'' Abs. 2 Satz 1: ``Pflege und Erziehung der Kinder sind das natuerliche Recht der Eltern und die zuvoerderst ihnen obliegende Pflicht.'' Abs. 4: ``Jede Mutter hat Anspruch auf den Schutz und die Fuersorge der Gemeinschaft.''}}
\newcommand{\ggfnArtSixZH}{\ggnotezhonce{art-6}{第6條第1項：「婚姻與家庭受國家秩序之特別保護。」第2項第1句：「照護及教育子女，為父母之自然權利，並為首先由其承擔之義務。」第4項：「每一母親均有請求共同體保護與照顧之權利。」}}
\newcommand{\ggfnArtSixOneDE}{\ggnotedeonce{art-6-1}{Art. 6 Abs. 1: ``Ehe und Familie stehen unter dem besonderen Schutze der staatlichen Ordnung.''}}
\newcommand{\ggfnArtSixOneZH}{\ggnotezhonce{art-6-1}{第6條第1項：「婚姻與家庭受國家秩序之特別保護。」}}
\newcommand{\ggfnArtSixTwoDE}{\ggnotedeonce{art-6-2}{Art. 6 Abs. 2: ``Pflege und Erziehung der Kinder sind das natuerliche Recht der Eltern und die zuvoerderst ihnen obliegende Pflicht. Ueber ihre Betaetigung wacht die staatliche Gemeinschaft.''}}
\newcommand{\ggfnArtSixTwoZH}{\ggnotezhonce{art-6-2}{第6條第2項：「照護及教育子女，為父母之自然權利，並為首先由其承擔之義務。國家共同體監督其行使。」}}
\newcommand{\ggfnArtSixTwoOneDE}{\ggnotedeonce{art-6-2-1}{Art. 6 Abs. 2 Satz 1: ``Pflege und Erziehung der Kinder sind das natuerliche Recht der Eltern und die zuvoerderst ihnen obliegende Pflicht.''}}
\newcommand{\ggfnArtSixTwoOneZH}{\ggnotezhonce{art-6-2-1}{第6條第2項第1句：「照護及教育子女，為父母之自然權利，並為首先由其承擔之義務。」}}
\newcommand{\ggfnArtSixFourDE}{\ggnotedeonce{art-6-4}{Art. 6 Abs. 4: ``Jede Mutter hat Anspruch auf den Schutz und die Fuersorge der Gemeinschaft.''}}
\newcommand{\ggfnArtSixFourZH}{\ggnotezhonce{art-6-4}{第6條第4項：「每一母親均有請求共同體保護與照顧之權利。」}}

% ===== Art. 12 =====
\newcommand{\ggfnArtTwelveOneDE}{\ggnotedeonce{art-12-1}{Art. 12 Abs. 1: ``Alle Deutschen haben das Recht, Beruf, Arbeitsplatz und Ausbildungsstaette frei zu waehlen. Die Berufsausuebung kann durch Gesetz oder auf Grund eines Gesetzes geregelt werden.''}}
\newcommand{\ggfnArtTwelveOneZH}{\ggnotezhonce{art-12-1}{第12條第1項：「所有德國人均有自由選擇職業、工作場所及訓練場所之權利。職業行使得以法律或基於法律加以規範。」}}

% ===== Art. 19 =====
\newcommand{\ggfnArtNineteenTwoDE}{\ggnotedeonce{art-19-2}{Art. 19 Abs. 2: ``In keinem Falle darf ein Grundrecht in seinem Wesensgehalt angetastet werden.''}}
\newcommand{\ggfnArtNineteenTwoZH}{\ggnotezhonce{art-19-2}{第19條第2項：「任何情形下，基本權利之本質內容均不得受到侵害。」}}

% ===== Art. 20 =====
\newcommand{\ggfnArtTwentyOneDE}{\ggnotedeonce{art-20-1}{Art. 20 Abs. 1: ``Die Bundesrepublik Deutschland ist ein demokratischer und sozialer Bundesstaat.''}}
\newcommand{\ggfnArtTwentyOneZH}{\ggnotezhonce{art-20-1}{第20條第1項：「德意志聯邦共和國為民主及社會之聯邦國。」}}
\newcommand{\ggfnArtTwentyThreeDE}{\ggnotedeonce{art-20-3}{Art. 20 Abs. 3: ``Die Gesetzgebung ist an die verfassungsmaessige Ordnung, die vollziehende Gewalt und die Rechtsprechung sind an Gesetz und Recht gebunden.''}}
\newcommand{\ggfnArtTwentyThreeZH}{\ggnotezhonce{art-20-3}{第20條第3項：「立法受合憲秩序拘束；行政權與司法權受法律與法拘束。」}}

% ===== Art. 26 + Art. 143（墮胎一案特有）=====
\newcommand{\ggfnArtTwoSixAndOneFourThreeDE}{\ggnotedeonce{art-26-143}{Art. 26 Abs. 1 Satz 2: ``Sie sind unter Strafe zu stellen.'' Art. 143 in der urspruenglichen Fassung (24. Mai 1949 bis 1. September 1951) enthielt Strafvorschriften zum Hochverrat; 1975 war Art. 143 bereits weggefallen.}}
\newcommand{\ggfnArtTwoSixAndOneFourThreeZH}{\ggnotezhonce{art-26-143}{第26條第1項第2句：「此等行為應規定為犯罪處罰。」第143條原始版本（1949年5月24日至1951年9月1日）含有關於叛國罪之刑罰規定；至1975年時，第143條已廢止。}}

% ===== Art. 28 =====
\newcommand{\ggfnArtTwentyEightOneDE}{\ggnotedeonce{art-28-1-1}{Art. 28 Abs. 1 Satz 1: ``Die verfassungsmaessige Ordnung in den Laendern muss den Grundsaetzen des republikanischen, demokratischen und sozialen Rechtsstaates im Sinne dieses Grundgesetzes entsprechen.''}}
\newcommand{\ggfnArtTwentyEightOneZH}{\ggnotezhonce{art-28-1-1}{第28條第1項第1句：「各邦之合憲秩序，須符合本基本法意義下共和、民主及社會法治國之原則。」}}

% ===== Art. 30 =====
\newcommand{\ggfnArtThirtyDE}{\ggnotedeonce{art-30}{Art. 30: ``Die Ausuebung der staatlichen Befugnisse und die Erfuellung der staatlichen Aufgaben ist Sache der Laender, soweit dieses Grundgesetz keine andere Regelung trifft oder zulaesst.''}}
\newcommand{\ggfnArtThirtyZH}{\ggnotezhonce{art-30}{第30條：「國家權限之行使與國家任務之履行，屬於各邦，但本基本法另有規定或容許者，不在此限。」}}

% ===== Art. 74 =====
\newcommand{\ggfnArtSeventyFourOneDE}{\ggnotedeonce{art-74-1}{Art. 74 Nr. 1: Die konkurrierende Gesetzgebung erstreckt sich auf ``das buergerliche Recht, das Strafrecht und den Strafvollzug, die Gerichtsverfassung, das gerichtliche Verfahren, die Rechtsanwaltschaft, das Notariat und die Rechtsberatung''.}}
\newcommand{\ggfnArtSeventyFourOneZH}{\ggnotezhonce{art-74-1}{第74條第1款：競合立法及於「民法、刑法及刑罰執行、法院組織、訴訟程序、律師制度、公證制度及法律諮詢」。}}
\newcommand{\ggfnArtSeventyFourSevenDE}{\ggnotedeonce{art-74-7}{Art. 74 Nr. 7: Die konkurrierende Gesetzgebung erstreckt sich auf ``die oeffentliche Fuersorge''.}}
\newcommand{\ggfnArtSeventyFourSevenZH}{\ggnotezhonce{art-74-7}{第74條第7款：競合立法及於「公共扶助」。}}
\newcommand{\ggfnArtSeventyFourTwelveDE}{\ggnotedeonce{art-74-12}{Art. 74 Nr. 12: Die konkurrierende Gesetzgebung erstreckt sich auf ``das Arbeitsrecht einschliesslich der Betriebsverfassung, des Arbeitsschutzes und der Arbeitsvermittlung sowie die Sozialversicherung einschliesslich der Arbeitslosenversicherung''.}}
\newcommand{\ggfnArtSeventyFourTwelveZH}{\ggnotezhonce{art-74-12}{第74條第12款：競合立法及於「勞動法，包括企業組織、勞動保護及就業媒合，以及社會保險，包括失業保險」。}}

% ===== Art. 77（墮胎一案特有）=====
\newcommand{\ggfnArtSevenSevenThreeDE}{\ggnotedeonce{art-77-3}{Art. 77 Abs. 3 Satz 1: ``Soweit zu einem Gesetze die Zustimmung des Bundesrates nicht erforderlich ist, kann der Bundesrat, wenn das Verfahren nach Absatz 2 beendigt ist, gegen ein vom Bundestage beschlossenes Gesetz binnen zwei Wochen Einspruch einlegen.''}}
\newcommand{\ggfnArtSevenSevenThreeZH}{\ggnotezhonce{art-77-3}{第77條第3項第1句：「法律無須聯邦參議院同意者，於第2項程序終了後，聯邦參議院得於二週內對聯邦議院通過之法律提出異議。」}}

% ===== Art. 80 =====
\newcommand{\ggfnArtEightyOneOneDE}{\ggnotedeonce{art-80-1-1}{Art. 80 Abs. 1 Satz 1: ``Durch Gesetz koennen die Bundesregierung, ein Bundesminister oder die Landesregierungen ermaechtigt werden, Rechtsverordnungen zu erlassen.''}}
\newcommand{\ggfnArtEightyOneOneZH}{\ggnotezhonce{art-80-1-1}{第80條第1項第1句：「得以法律授權聯邦政府、聯邦部長或邦政府發布法規命令。」}}

% ===== Art. 83 =====
\newcommand{\ggfnArtEightyThreeDE}{\ggnotedeonce{art-83}{Art. 83: ``Die Laender fuehren die Bundesgesetze als eigene Angelegenheit aus, soweit dieses Grundgesetz nichts anderes bestimmt oder zulaesst.''}}
\newcommand{\ggfnArtEightyThreeZH}{\ggnotezhonce{art-83}{第83條：「各邦以自己事務執行聯邦法律，但本基本法另有規定或容許者，不在此限。」}}

% ===== Art. 84 =====
\newcommand{\ggfnArtEightyFourOneDE}{\ggnotedeonce{art-84-1}{Art. 84 Abs. 1: ``Fuehren die Laender die Bundesgesetze als eigene Angelegenheit aus, so regeln sie die Einrichtung der Behoerden und das Verwaltungsverfahren, soweit nicht Bundesgesetze mit Zustimmung des Bundesrates etwas anderes bestimmen.''}}
\newcommand{\ggfnArtEightyFourOneZH}{\ggnotezhonce{art-84-1}{第84條第1項：「各邦以自己事務執行聯邦法律時，除聯邦法律經聯邦參議院同意另有規定外，由各邦規範機關之設置及行政程序。」}}
% 別名（墮胎一案使用之縮寫命名）
\let\ggfnArtEightFourOneDE\ggfnArtEightyFourOneDE
\let\ggfnArtEightFourOneZH\ggfnArtEightyFourOneZH

% ===== Art. 93 =====
\newcommand{\ggfnArtNinetyThreeOneTwoDE}{\ggnotedeonce{art-93-1-2}{Art. 93 Abs. 1 Nr. 2: Das Bundesverfassungsgericht entscheidet ``bei Meinungsverschiedenheiten oder Zweifeln ueber die foermliche und sachliche Vereinbarkeit von Bundesrecht oder Landesrecht mit diesem Grundgesetze ... auf Antrag der Bundesregierung, einer Landesregierung oder eines Drittels der Mitglieder des Bundestages''.}}
\newcommand{\ggfnArtNinetyThreeOneTwoZH}{\ggnotezhonce{art-93-1-2}{第93條第1項第2款：聯邦憲法法院就「聯邦法或邦法在形式及實質上是否與本基本法相容之意見分歧或疑義……依聯邦政府、邦政府或聯邦議院三分之一議員之聲請」作成裁判。}}
% 別名（墮胎一案使用之縮寫命名）
\let\ggfnArtNineThreeOneTwoDE\ggfnArtNinetyThreeOneTwoDE
\let\ggfnArtNineThreeOneTwoZH\ggfnArtNinetyThreeOneTwoZH

% ===== Art. 102 =====
\newcommand{\ggfnArtOneZeroTwoDE}{\ggnotedeonce{art-102}{Art. 102: ``Die Todesstrafe ist abgeschafft.''}}
\newcommand{\ggfnArtOneZeroTwoZH}{\ggnotezhonce{art-102}{第102條：「死刑廢止。」}}

% ===== Art. 103 =====
\newcommand{\ggfnArtOneZeroThreeTwoDE}{\ggnotedeonce{art-103-2}{Art. 103 Abs. 2: ``Eine Tat kann nur bestraft werden, wenn die Strafbarkeit gesetzlich bestimmt war, bevor die Tat begangen wurde.''}}
\newcommand{\ggfnArtOneZeroThreeTwoZH}{\ggnotezhonce{art-103-2}{第103條第2項：「行為之可罰性，須於行為前以法律定之，始得處罰。」}}

% ===== Art. 104 =====
\newcommand{\ggfnArtOneZeroFourOneDE}{\ggnotedeonce{art-104-1}{Art. 104 Abs. 1: ``Die Freiheit der Person kann nur auf Grund eines foermlichen Gesetzes und nur unter Beachtung der darin vorgeschriebenen Formen beschraenkt werden.''}}
\newcommand{\ggfnArtOneZeroFourOneZH}{\ggnotezhonce{art-104-1}{第104條第1項：「人身自由僅得基於形式法律，並遵守其中規定之形式加以限制。」}}

% ===== 組合腳註：Art. 3 + Art. 6（墮胎一案特有）=====
\newcommand{\ggfnArtThreeSixDE}{\ggnotedeonce{art-3-6}{Art. 3: ``Alle Menschen sind vor dem Gesetz gleich. Maenner und Frauen sind gleichberechtigt. Niemand darf wegen seines Geschlechtes, seiner Abstammung, seiner Rasse, seiner Sprache, seiner Heimat und Herkunft, seines Glaubens, seiner religioesen oder politischen Anschauungen benachteiligt oder bevorzugt werden.'' Art. 6 Abs. 1: ``Ehe und Familie stehen unter dem besonderen Schutze der staatlichen Ordnung.'' Abs. 2 Satz 1: ``Pflege und Erziehung der Kinder sind das natuerliche Recht der Eltern und die zuvoerderst ihnen obliegende Pflicht.'' Abs. 4: ``Jede Mutter hat Anspruch auf den Schutz und die Fuersorge der Gemeinschaft.''}\ggmarknotede{art-3}\ggmarknotede{art-6}\ggmarknotede{art-6-1-4}}
\newcommand{\ggfnArtThreeSixZH}{\ggnotezhonce{art-3-6}{第3條：「所有人在法律前一律平等。男女享有平等權利。任何人不得因其性別、血統、種族、語言、故鄉及出身、信仰、宗教或政治見解而受不利益或優惠。」第6條第1項：「婚姻與家庭受國家秩序之特別保護。」第2項第1句：「照護及教育子女，為父母之自然權利，並為首先由其承擔之義務。」第4項：「每一母親均有請求共同體保護與照顧之權利。」}\ggmarknotezh{art-3}\ggmarknotezh{art-6}\ggmarknotezh{art-6-1-4}}

% ===== 組合腳註：Art. 6(1) + Art. 6(4)（墮胎一案特有）=====
\newcommand{\ggfnArtSixOneFourDE}{\ggnotedeonce{art-6-1-4}{Art. 6 Abs. 1: ``Ehe und Familie stehen unter dem besonderen Schutze der staatlichen Ordnung.'' Art. 6 Abs. 4: ``Jede Mutter hat Anspruch auf den Schutz und die Fuersorge der Gemeinschaft.''}\ggmarknotede{art-6}}
\newcommand{\ggfnArtSixOneFourZH}{\ggnotezhonce{art-6-1-4}{第6條第1項：「婚姻與家庭受國家秩序之特別保護。」第4項：「每一母親均有請求共同體保護與照顧之權利。」}\ggmarknotezh{art-6}}

% ===== 組合腳註：Art. 1(1) + Art. 2(2)(1)（墮胎二案特有）=====
\newcommand{\ggfnArtOneOneTwoTwoOneDE}{\ggnotedeonce{art-1-1+2-2-1}{Art. 1 Abs. 1: ``Die Wuerde des Menschen ist unantastbar. Sie zu achten und zu schuetzen ist Verpflichtung aller staatlichen Gewalt.'' Art. 2 Abs. 2 Satz 1: ``Jeder hat das Recht auf Leben und koerperliche Unversehrtheit.''}\ggmarknotede{art-1-1}\ggmarknotede{art-2-2-1}}
\newcommand{\ggfnArtOneOneTwoTwoOneZH}{\ggnotezhonce{art-1-1+2-2-1}{第1條第1項：「人之尊嚴不可侵犯。尊重及保護之，為一切國家權力之義務。」第2條第2項第1句：「人人有生命與身體完整性之權利。」}\ggmarknotezh{art-1-1}\ggmarknotezh{art-2-2-1}}

% ===== 組合腳註：Art. 1(1) + Art. 2(2)（墮胎二案特有）=====
\newcommand{\ggfnArtOneOneTwoTwoDE}{\ggnotedeonce{art-1-1+2-2}{Art. 1 Abs. 1: ``Die Wuerde des Menschen ist unantastbar. Sie zu achten und zu schuetzen ist Verpflichtung aller staatlichen Gewalt.'' Art. 2 Abs. 2: ``Jeder hat das Recht auf Leben und koerperliche Unversehrtheit. Die Freiheit der Person ist unverletzlich. In diese Rechte darf nur auf Grund eines Gesetzes eingegriffen werden.''}\ggmarknotede{art-1-1}\ggmarknotede{art-2-2}\ggmarknotede{art-2-2-1}}
\newcommand{\ggfnArtOneOneTwoTwoZH}{\ggnotezhonce{art-1-1+2-2}{第1條第1項：「人之尊嚴不可侵犯。尊重及保護之，為一切國家權力之義務。」第2條第2項：「人人有生命與身體完整性之權利。人身自由不可侵犯。此等權利僅得依法律加以干預。」}\ggmarknotezh{art-1-1}\ggmarknotezh{art-2-2}\ggmarknotezh{art-2-2-1}}

% ===== 組合腳註：Art. 1(1) + Art. 2(1)（墮胎二案特有）=====
\newcommand{\ggfnArtOneOneTwoOneDE}{\ggnotedeonce{art-1-1+2-1}{Art. 1 Abs. 1: ``Die Wuerde des Menschen ist unantastbar. Sie zu achten und zu schuetzen ist Verpflichtung aller staatlichen Gewalt.'' Art. 2 Abs. 1: ``Jeder hat das Recht auf die freie Entfaltung seiner Persoenlichkeit, soweit er nicht die Rechte anderer verletzt und nicht gegen die verfassungsmaessige Ordnung oder das Sittengesetz verstoesst.''}\ggmarknotede{art-1-1}\ggmarknotede{art-2-1}}
\newcommand{\ggfnArtOneOneTwoOneZH}{\ggnotezhonce{art-1-1+2-1}{第1條第1項：「人之尊嚴不可侵犯。尊重及保護之，為一切國家權力之義務。」第2條第1項：「人人有自由發展其人格之權利，但不得侵害他人權利，亦不得違反合憲秩序或道德律。」}\ggmarknotezh{art-1-1}\ggmarknotezh{art-2-1}}

% ===== 組合腳註：Art. 2(2)(1) + Art. 1(1)（墮胎一案特有，含交叉標記）=====
\newcommand{\ggfnLifeDignityDE}{\ggnotedeonce{life-dignity}{Art. 2 Abs. 2 Satz 1: ``Jeder hat das Recht auf Leben und koerperliche Unversehrtheit.'' Art. 1 Abs. 1: ``Die Wuerde des Menschen ist unantastbar. Sie zu achten und zu schuetzen ist Verpflichtung aller staatlichen Gewalt.''}\ggmarknotede{art-2-2-1}\ggmarknotede{art-1-1}\ggmarknotede{art-1-1-2}}
\newcommand{\ggfnLifeDignityZH}{\ggnotezhonce{life-dignity}{第2條第2項第1句：「人人有生命與身體完整性之權利。」第1條第1項：「人之尊嚴不可侵犯。尊重及保護之，為一切國家權力之義務。」}\ggmarknotezh{art-2-2-1}\ggmarknotezh{art-1-1}\ggmarknotezh{art-1-1-2}}

% ===== 組合腳註：Art. 1(1) + Art. 2(2)（墮胎一案特有，條件判斷版）=====
\newcommand{\ggfnArtOneTwoTwoDE}{\ifcsname ggnoteseen@de@life-dignity\endcsname\ggfnArtTwoTwoRestDE{}\else\ggnotedeonce{art-1-2-2}{Art. 1 Abs. 1: ``Die Wuerde des Menschen ist unantastbar. Sie zu achten und zu schuetzen ist Verpflichtung aller staatlichen Gewalt.'' Art. 2 Abs. 2: ``Jeder hat das Recht auf Leben und koerperliche Unversehrtheit. Die Freiheit der Person ist unverletzlich. In diese Rechte darf nur auf Grund eines Gesetzes eingegriffen werden.''}\ggmarknotede{art-1-1}\ggmarknotede{art-1-1-2}\ggmarknotede{art-2-2}\ggmarknotede{art-2-2-1}\fi}
\newcommand{\ggfnArtOneTwoTwoZH}{\ifcsname ggnoteseen@zh@life-dignity\endcsname\ggfnArtTwoTwoRestZH{}\else\ggnotezhonce{art-1-2-2}{第1條第1項：「人之尊嚴不可侵犯。尊重及保護之，為一切國家權力之義務。」第2條第2項：「人人有生命與身體完整性之權利。人身自由不可侵犯。此等權利僅得依法律加以干預。」}\ggmarknotezh{art-1-1}\ggmarknotezh{art-1-1-2}\ggmarknotezh{art-2-2}\ggmarknotezh{art-2-2-1}\fi}

% ===== 組合腳註：Art. 1(1) + Art. 19(2)（墮胎一案特有，條件判斷版）=====
\newcommand{\ggfnArtOneNineteenDE}{\ifcsname ggnoteseen@de@art-1-1\endcsname\ggfnArtNineteenTwoDE{}\else\ggnotedeonce{art-1-19}{Art. 1 Abs. 1: ``Die Wuerde des Menschen ist unantastbar. Sie zu achten und zu schuetzen ist Verpflichtung aller staatlichen Gewalt.'' Art. 19 Abs. 2: ``In keinem Falle darf ein Grundrecht in seinem Wesensgehalt angetastet werden.''}\ggmarknotede{art-1-1}\ggmarknotede{art-1-1-2}\ggmarknotede{art-19-2}\fi}
\newcommand{\ggfnArtOneNineteenZH}{\ifcsname ggnoteseen@zh@art-1-1\endcsname\ggfnArtNineteenTwoZH{}\else\ggnotezhonce{art-1-19}{第1條第1項：「人之尊嚴不可侵犯。尊重及保護之，為一切國家權力之義務。」第19條第2項：「任何情形下，基本權利之本質內容均不得受到侵害。」}\ggmarknotezh{art-1-1}\ggmarknotezh{art-1-1-2}\ggmarknotezh{art-19-2}\fi}


% ===== GG 版本旗標（\ggversionde/zh 由各案在 % bverfge_gg.tex — 德國墮胎案 GG 條文腳註共用定義
% 使用方式：在各案 .tex 中先定義 \ggversionde/\ggversionzh，再 \input{../../共用LaTeX/bverfge_gg}

% ===== GG 版本旗標（\ggversionde/zh 由各案在 \input{../../共用LaTeX/bverfge_gg} 之前定義）=====
\newif\ifggversionnotede
\newif\ifggversionnotezh

% ===== 編者註腳基礎設施（首次標記模式）=====
\newif\ifeditornotelabelde
\newif\ifeditornotelabelzh
\newcommand{\editornotelabelde}{\ifeditornotelabelde\else\emph{Hrsg.-Anm. (nachfolgend ohne Kennzeichnung)}: \global\editornotelabeldetrue\fi}
\newcommand{\editornotelabelzh}{\ifeditornotelabelzh\else 編者註（下同）：\global\editornotelabelzhtrue\fi}
\newcommand{\editornotede}[1]{\ifshoweditornotes\footnote{\normalfont\footnotesize\editornotelabelde #1}\else\ifclassroommode\footnote{\normalfont\footnotesize\editornotelabelde #1}\fi\fi}
\newcommand{\editornotezh}[1]{\ifshoweditornotes\footnote{\normalfont\footnotesize\editornotelabelzh #1}\else\ifclassroommode\footnote{\normalfont\footnotesize\editornotelabelzh #1}\fi\fi}

% ===== GG 引用系統（\ggversionde/zh 由各案在 \input 前定義）=====
\newcommand{\ggmarknotede}[1]{\global\expandafter\let\csname ggnoteseen@de@#1\endcsname\empty}
\newcommand{\ggmarknotezh}[1]{\global\expandafter\let\csname ggnoteseen@zh@#1\endcsname\empty}
\newcommand{\ggnotedeonce}[2]{\ifcsname ggnoteseen@de@#1\endcsname\else\global\expandafter\let\csname ggnoteseen@de@#1\endcsname\empty\ggnotede{#2}\fi}
\newcommand{\ggnotezhonce}[2]{\ifcsname ggnoteseen@zh@#1\endcsname\else\global\expandafter\let\csname ggnoteseen@zh@#1\endcsname\empty\ggnotezh{#2}\fi}
\newcommand{\ggnotede}[1]{\editornotede{\ggversionde #1}}
\newcommand{\ggnotezh}[1]{\editornotezh{\ggversionzh #1}}

% ===== Art. 1 =====
\newcommand{\ggfnArtOneOneDE}{\ggnotedeonce{art-1-1}{Art. 1 Abs. 1: ``Die Wuerde des Menschen ist unantastbar. Sie zu achten und zu schuetzen ist Verpflichtung aller staatlichen Gewalt.''}}
\newcommand{\ggfnArtOneOneZH}{\ggnotezhonce{art-1-1}{第1條第1項：「人之尊嚴不可侵犯。尊重及保護之，為一切國家權力之義務。」}}
\newcommand{\ggfnArtOneOneTwoDE}{\ggnotedeonce{art-1-1-2}{Art. 1 Abs. 1 Satz 2: ``Sie zu achten und zu schuetzen ist Verpflichtung aller staatlichen Gewalt.''}}
\newcommand{\ggfnArtOneOneTwoZH}{\ggnotezhonce{art-1-1-2}{第1條第1項第2句：「尊重及保護之，為一切國家權力之義務。」}}

% ===== Art. 2 =====
\newcommand{\ggfnArtTwoOneDE}{\ggnotedeonce{art-2-1}{Art. 2 Abs. 1: ``Jeder hat das Recht auf die freie Entfaltung seiner Persoenlichkeit, soweit er nicht die Rechte anderer verletzt und nicht gegen die verfassungsmaessige Ordnung oder das Sittengesetz verstoesst.''}}
\newcommand{\ggfnArtTwoOneZH}{\ggnotezhonce{art-2-1}{第2條第1項：「人人有自由發展其人格之權利，但不得侵害他人權利，亦不得違反合憲秩序或道德律。」}}
% 簡易預設版（墮胎一案以 \renewcommand 覆寫為含條件邏輯之版本）
\newcommand{\ggfnArtTwoTwoDE}{\ggnotedeonce{art-2-2}{Art. 2 Abs. 2: ``Jeder hat das Recht auf Leben und koerperliche Unversehrtheit. Die Freiheit der Person ist unverletzlich. In diese Rechte darf nur auf Grund eines Gesetzes eingegriffen werden.''}\ggmarknotede{art-2-2-1}}
\newcommand{\ggfnArtTwoTwoZH}{\ggnotezhonce{art-2-2}{第2條第2項：「人人有生命與身體完整性之權利。人身自由不可侵犯。此等權利僅得依法律加以干預。」}\ggmarknotezh{art-2-2-1}}
\newcommand{\ggfnArtTwoTwoOneDE}{\ggnotedeonce{art-2-2-1}{Art. 2 Abs. 2 Satz 1: ``Jeder hat das Recht auf Leben und koerperliche Unversehrtheit.''}}
\newcommand{\ggfnArtTwoTwoOneZH}{\ggnotezhonce{art-2-2-1}{第2條第2項第1句：「人人有生命與身體完整性之權利。」}}
\newcommand{\ggfnArtTwoTwoThreeDE}{\ggnotedeonce{art-2-2-3}{Art. 2 Abs. 2 Satz 3: ``In diese Rechte darf nur auf Grund eines Gesetzes eingegriffen werden.''}}
\newcommand{\ggfnArtTwoTwoThreeZH}{\ggnotezhonce{art-2-2-3}{第2條第2項第3句：「此等權利僅得依法律加以干預。」}}
% 墮胎一案特有：第2條第2項第2–3句（Art. 2(2) 去掉第1句後的「餘文」）
\newcommand{\ggfnArtTwoTwoRestDE}{\ggnotedeonce{art-2-2-rest}{Art. 2 Abs. 2 Saetze 2-3: ``Die Freiheit der Person ist unverletzlich. In diese Rechte darf nur auf Grund eines Gesetzes eingegriffen werden.''}\ggmarknotede{art-2-2}}
\newcommand{\ggfnArtTwoTwoRestZH}{\ggnotezhonce{art-2-2-rest}{第2條第2項第2、3句：「人身自由不可侵犯。此等權利僅得依法律加以干預。」}\ggmarknotezh{art-2-2}}

% ===== Art. 3 =====
\newcommand{\ggfnArtThreeDE}{\ggnotedeonce{art-3}{Art. 3: ``Alle Menschen sind vor dem Gesetz gleich. Maenner und Frauen sind gleichberechtigt. Niemand darf wegen seines Geschlechtes, seiner Abstammung, seiner Rasse, seiner Sprache, seiner Heimat und Herkunft, seines Glaubens, seiner religioesen oder politischen Anschauungen benachteiligt oder bevorzugt werden.''}}
\newcommand{\ggfnArtThreeZH}{\ggnotezhonce{art-3}{第3條：「所有人在法律前一律平等。男女享有平等權利。任何人不得因其性別、血統、種族、語言、故鄉及出身、信仰、宗教或政治見解而受不利益或優惠。」}}
\newcommand{\ggfnArtThreeTwoDE}{\ggnotedeonce{art-3-2}{Art. 3 Abs. 2: ``Maenner und Frauen sind gleichberechtigt.''}}
\newcommand{\ggfnArtThreeTwoZH}{\ggnotezhonce{art-3-2}{第3條第2項：「男女享有平等權利。」}}

% ===== Art. 4 =====
\newcommand{\ggfnArtFourOneDE}{\ggnotedeonce{art-4-1}{Art. 4 Abs. 1: ``Die Freiheit des Glaubens, des Gewissens und die Freiheit des religioesen und weltanschaulichen Bekenntnisses sind unverletzlich.''}}
\newcommand{\ggfnArtFourOneZH}{\ggnotezhonce{art-4-1}{第4條第1項：「信仰自由、良心自由，以及宗教與世界觀信念自由，不可侵犯。」}}

% ===== Art. 5 =====
\newcommand{\ggfnArtFiveOneDE}{\ggnotedeonce{art-5-1}{Art. 5 Abs. 1: ``Jeder hat das Recht, seine Meinung in Wort, Schrift und Bild frei zu aeussern und zu verbreiten und sich aus allgemein zugaenglichen Quellen ungehindert zu unterrichten. Die Pressefreiheit und die Freiheit der Berichterstattung durch Rundfunk und Film werden gewaehrleistet. Eine Zensur findet nicht statt.''}}
\newcommand{\ggfnArtFiveOneZH}{\ggnotezhonce{art-5-1}{第5條第1項：「人人有以言語、文字及圖像自由表達及傳播其意見，並由一般可接近來源不受阻礙獲取資訊之權利。新聞自由及廣播、電影報導自由受保障。不設審查制度。」}}

% ===== Art. 6 =====
\newcommand{\ggfnArtSixDE}{\ggnotedeonce{art-6}{Art. 6 Abs. 1: ``Ehe und Familie stehen unter dem besonderen Schutze der staatlichen Ordnung.'' Abs. 2 Satz 1: ``Pflege und Erziehung der Kinder sind das natuerliche Recht der Eltern und die zuvoerderst ihnen obliegende Pflicht.'' Abs. 4: ``Jede Mutter hat Anspruch auf den Schutz und die Fuersorge der Gemeinschaft.''}}
\newcommand{\ggfnArtSixZH}{\ggnotezhonce{art-6}{第6條第1項：「婚姻與家庭受國家秩序之特別保護。」第2項第1句：「照護及教育子女，為父母之自然權利，並為首先由其承擔之義務。」第4項：「每一母親均有請求共同體保護與照顧之權利。」}}
\newcommand{\ggfnArtSixOneDE}{\ggnotedeonce{art-6-1}{Art. 6 Abs. 1: ``Ehe und Familie stehen unter dem besonderen Schutze der staatlichen Ordnung.''}}
\newcommand{\ggfnArtSixOneZH}{\ggnotezhonce{art-6-1}{第6條第1項：「婚姻與家庭受國家秩序之特別保護。」}}
\newcommand{\ggfnArtSixTwoDE}{\ggnotedeonce{art-6-2}{Art. 6 Abs. 2: ``Pflege und Erziehung der Kinder sind das natuerliche Recht der Eltern und die zuvoerderst ihnen obliegende Pflicht. Ueber ihre Betaetigung wacht die staatliche Gemeinschaft.''}}
\newcommand{\ggfnArtSixTwoZH}{\ggnotezhonce{art-6-2}{第6條第2項：「照護及教育子女，為父母之自然權利，並為首先由其承擔之義務。國家共同體監督其行使。」}}
\newcommand{\ggfnArtSixTwoOneDE}{\ggnotedeonce{art-6-2-1}{Art. 6 Abs. 2 Satz 1: ``Pflege und Erziehung der Kinder sind das natuerliche Recht der Eltern und die zuvoerderst ihnen obliegende Pflicht.''}}
\newcommand{\ggfnArtSixTwoOneZH}{\ggnotezhonce{art-6-2-1}{第6條第2項第1句：「照護及教育子女，為父母之自然權利，並為首先由其承擔之義務。」}}
\newcommand{\ggfnArtSixFourDE}{\ggnotedeonce{art-6-4}{Art. 6 Abs. 4: ``Jede Mutter hat Anspruch auf den Schutz und die Fuersorge der Gemeinschaft.''}}
\newcommand{\ggfnArtSixFourZH}{\ggnotezhonce{art-6-4}{第6條第4項：「每一母親均有請求共同體保護與照顧之權利。」}}

% ===== Art. 12 =====
\newcommand{\ggfnArtTwelveOneDE}{\ggnotedeonce{art-12-1}{Art. 12 Abs. 1: ``Alle Deutschen haben das Recht, Beruf, Arbeitsplatz und Ausbildungsstaette frei zu waehlen. Die Berufsausuebung kann durch Gesetz oder auf Grund eines Gesetzes geregelt werden.''}}
\newcommand{\ggfnArtTwelveOneZH}{\ggnotezhonce{art-12-1}{第12條第1項：「所有德國人均有自由選擇職業、工作場所及訓練場所之權利。職業行使得以法律或基於法律加以規範。」}}

% ===== Art. 19 =====
\newcommand{\ggfnArtNineteenTwoDE}{\ggnotedeonce{art-19-2}{Art. 19 Abs. 2: ``In keinem Falle darf ein Grundrecht in seinem Wesensgehalt angetastet werden.''}}
\newcommand{\ggfnArtNineteenTwoZH}{\ggnotezhonce{art-19-2}{第19條第2項：「任何情形下，基本權利之本質內容均不得受到侵害。」}}

% ===== Art. 20 =====
\newcommand{\ggfnArtTwentyOneDE}{\ggnotedeonce{art-20-1}{Art. 20 Abs. 1: ``Die Bundesrepublik Deutschland ist ein demokratischer und sozialer Bundesstaat.''}}
\newcommand{\ggfnArtTwentyOneZH}{\ggnotezhonce{art-20-1}{第20條第1項：「德意志聯邦共和國為民主及社會之聯邦國。」}}
\newcommand{\ggfnArtTwentyThreeDE}{\ggnotedeonce{art-20-3}{Art. 20 Abs. 3: ``Die Gesetzgebung ist an die verfassungsmaessige Ordnung, die vollziehende Gewalt und die Rechtsprechung sind an Gesetz und Recht gebunden.''}}
\newcommand{\ggfnArtTwentyThreeZH}{\ggnotezhonce{art-20-3}{第20條第3項：「立法受合憲秩序拘束；行政權與司法權受法律與法拘束。」}}

% ===== Art. 26 + Art. 143（墮胎一案特有）=====
\newcommand{\ggfnArtTwoSixAndOneFourThreeDE}{\ggnotedeonce{art-26-143}{Art. 26 Abs. 1 Satz 2: ``Sie sind unter Strafe zu stellen.'' Art. 143 in der urspruenglichen Fassung (24. Mai 1949 bis 1. September 1951) enthielt Strafvorschriften zum Hochverrat; 1975 war Art. 143 bereits weggefallen.}}
\newcommand{\ggfnArtTwoSixAndOneFourThreeZH}{\ggnotezhonce{art-26-143}{第26條第1項第2句：「此等行為應規定為犯罪處罰。」第143條原始版本（1949年5月24日至1951年9月1日）含有關於叛國罪之刑罰規定；至1975年時，第143條已廢止。}}

% ===== Art. 28 =====
\newcommand{\ggfnArtTwentyEightOneDE}{\ggnotedeonce{art-28-1-1}{Art. 28 Abs. 1 Satz 1: ``Die verfassungsmaessige Ordnung in den Laendern muss den Grundsaetzen des republikanischen, demokratischen und sozialen Rechtsstaates im Sinne dieses Grundgesetzes entsprechen.''}}
\newcommand{\ggfnArtTwentyEightOneZH}{\ggnotezhonce{art-28-1-1}{第28條第1項第1句：「各邦之合憲秩序，須符合本基本法意義下共和、民主及社會法治國之原則。」}}

% ===== Art. 30 =====
\newcommand{\ggfnArtThirtyDE}{\ggnotedeonce{art-30}{Art. 30: ``Die Ausuebung der staatlichen Befugnisse und die Erfuellung der staatlichen Aufgaben ist Sache der Laender, soweit dieses Grundgesetz keine andere Regelung trifft oder zulaesst.''}}
\newcommand{\ggfnArtThirtyZH}{\ggnotezhonce{art-30}{第30條：「國家權限之行使與國家任務之履行，屬於各邦，但本基本法另有規定或容許者，不在此限。」}}

% ===== Art. 74 =====
\newcommand{\ggfnArtSeventyFourOneDE}{\ggnotedeonce{art-74-1}{Art. 74 Nr. 1: Die konkurrierende Gesetzgebung erstreckt sich auf ``das buergerliche Recht, das Strafrecht und den Strafvollzug, die Gerichtsverfassung, das gerichtliche Verfahren, die Rechtsanwaltschaft, das Notariat und die Rechtsberatung''.}}
\newcommand{\ggfnArtSeventyFourOneZH}{\ggnotezhonce{art-74-1}{第74條第1款：競合立法及於「民法、刑法及刑罰執行、法院組織、訴訟程序、律師制度、公證制度及法律諮詢」。}}
\newcommand{\ggfnArtSeventyFourSevenDE}{\ggnotedeonce{art-74-7}{Art. 74 Nr. 7: Die konkurrierende Gesetzgebung erstreckt sich auf ``die oeffentliche Fuersorge''.}}
\newcommand{\ggfnArtSeventyFourSevenZH}{\ggnotezhonce{art-74-7}{第74條第7款：競合立法及於「公共扶助」。}}
\newcommand{\ggfnArtSeventyFourTwelveDE}{\ggnotedeonce{art-74-12}{Art. 74 Nr. 12: Die konkurrierende Gesetzgebung erstreckt sich auf ``das Arbeitsrecht einschliesslich der Betriebsverfassung, des Arbeitsschutzes und der Arbeitsvermittlung sowie die Sozialversicherung einschliesslich der Arbeitslosenversicherung''.}}
\newcommand{\ggfnArtSeventyFourTwelveZH}{\ggnotezhonce{art-74-12}{第74條第12款：競合立法及於「勞動法，包括企業組織、勞動保護及就業媒合，以及社會保險，包括失業保險」。}}

% ===== Art. 77（墮胎一案特有）=====
\newcommand{\ggfnArtSevenSevenThreeDE}{\ggnotedeonce{art-77-3}{Art. 77 Abs. 3 Satz 1: ``Soweit zu einem Gesetze die Zustimmung des Bundesrates nicht erforderlich ist, kann der Bundesrat, wenn das Verfahren nach Absatz 2 beendigt ist, gegen ein vom Bundestage beschlossenes Gesetz binnen zwei Wochen Einspruch einlegen.''}}
\newcommand{\ggfnArtSevenSevenThreeZH}{\ggnotezhonce{art-77-3}{第77條第3項第1句：「法律無須聯邦參議院同意者，於第2項程序終了後，聯邦參議院得於二週內對聯邦議院通過之法律提出異議。」}}

% ===== Art. 80 =====
\newcommand{\ggfnArtEightyOneOneDE}{\ggnotedeonce{art-80-1-1}{Art. 80 Abs. 1 Satz 1: ``Durch Gesetz koennen die Bundesregierung, ein Bundesminister oder die Landesregierungen ermaechtigt werden, Rechtsverordnungen zu erlassen.''}}
\newcommand{\ggfnArtEightyOneOneZH}{\ggnotezhonce{art-80-1-1}{第80條第1項第1句：「得以法律授權聯邦政府、聯邦部長或邦政府發布法規命令。」}}

% ===== Art. 83 =====
\newcommand{\ggfnArtEightyThreeDE}{\ggnotedeonce{art-83}{Art. 83: ``Die Laender fuehren die Bundesgesetze als eigene Angelegenheit aus, soweit dieses Grundgesetz nichts anderes bestimmt oder zulaesst.''}}
\newcommand{\ggfnArtEightyThreeZH}{\ggnotezhonce{art-83}{第83條：「各邦以自己事務執行聯邦法律，但本基本法另有規定或容許者，不在此限。」}}

% ===== Art. 84 =====
\newcommand{\ggfnArtEightyFourOneDE}{\ggnotedeonce{art-84-1}{Art. 84 Abs. 1: ``Fuehren die Laender die Bundesgesetze als eigene Angelegenheit aus, so regeln sie die Einrichtung der Behoerden und das Verwaltungsverfahren, soweit nicht Bundesgesetze mit Zustimmung des Bundesrates etwas anderes bestimmen.''}}
\newcommand{\ggfnArtEightyFourOneZH}{\ggnotezhonce{art-84-1}{第84條第1項：「各邦以自己事務執行聯邦法律時，除聯邦法律經聯邦參議院同意另有規定外，由各邦規範機關之設置及行政程序。」}}
% 別名（墮胎一案使用之縮寫命名）
\let\ggfnArtEightFourOneDE\ggfnArtEightyFourOneDE
\let\ggfnArtEightFourOneZH\ggfnArtEightyFourOneZH

% ===== Art. 93 =====
\newcommand{\ggfnArtNinetyThreeOneTwoDE}{\ggnotedeonce{art-93-1-2}{Art. 93 Abs. 1 Nr. 2: Das Bundesverfassungsgericht entscheidet ``bei Meinungsverschiedenheiten oder Zweifeln ueber die foermliche und sachliche Vereinbarkeit von Bundesrecht oder Landesrecht mit diesem Grundgesetze ... auf Antrag der Bundesregierung, einer Landesregierung oder eines Drittels der Mitglieder des Bundestages''.}}
\newcommand{\ggfnArtNinetyThreeOneTwoZH}{\ggnotezhonce{art-93-1-2}{第93條第1項第2款：聯邦憲法法院就「聯邦法或邦法在形式及實質上是否與本基本法相容之意見分歧或疑義……依聯邦政府、邦政府或聯邦議院三分之一議員之聲請」作成裁判。}}
% 別名（墮胎一案使用之縮寫命名）
\let\ggfnArtNineThreeOneTwoDE\ggfnArtNinetyThreeOneTwoDE
\let\ggfnArtNineThreeOneTwoZH\ggfnArtNinetyThreeOneTwoZH

% ===== Art. 102 =====
\newcommand{\ggfnArtOneZeroTwoDE}{\ggnotedeonce{art-102}{Art. 102: ``Die Todesstrafe ist abgeschafft.''}}
\newcommand{\ggfnArtOneZeroTwoZH}{\ggnotezhonce{art-102}{第102條：「死刑廢止。」}}

% ===== Art. 103 =====
\newcommand{\ggfnArtOneZeroThreeTwoDE}{\ggnotedeonce{art-103-2}{Art. 103 Abs. 2: ``Eine Tat kann nur bestraft werden, wenn die Strafbarkeit gesetzlich bestimmt war, bevor die Tat begangen wurde.''}}
\newcommand{\ggfnArtOneZeroThreeTwoZH}{\ggnotezhonce{art-103-2}{第103條第2項：「行為之可罰性，須於行為前以法律定之，始得處罰。」}}

% ===== Art. 104 =====
\newcommand{\ggfnArtOneZeroFourOneDE}{\ggnotedeonce{art-104-1}{Art. 104 Abs. 1: ``Die Freiheit der Person kann nur auf Grund eines foermlichen Gesetzes und nur unter Beachtung der darin vorgeschriebenen Formen beschraenkt werden.''}}
\newcommand{\ggfnArtOneZeroFourOneZH}{\ggnotezhonce{art-104-1}{第104條第1項：「人身自由僅得基於形式法律，並遵守其中規定之形式加以限制。」}}

% ===== 組合腳註：Art. 3 + Art. 6（墮胎一案特有）=====
\newcommand{\ggfnArtThreeSixDE}{\ggnotedeonce{art-3-6}{Art. 3: ``Alle Menschen sind vor dem Gesetz gleich. Maenner und Frauen sind gleichberechtigt. Niemand darf wegen seines Geschlechtes, seiner Abstammung, seiner Rasse, seiner Sprache, seiner Heimat und Herkunft, seines Glaubens, seiner religioesen oder politischen Anschauungen benachteiligt oder bevorzugt werden.'' Art. 6 Abs. 1: ``Ehe und Familie stehen unter dem besonderen Schutze der staatlichen Ordnung.'' Abs. 2 Satz 1: ``Pflege und Erziehung der Kinder sind das natuerliche Recht der Eltern und die zuvoerderst ihnen obliegende Pflicht.'' Abs. 4: ``Jede Mutter hat Anspruch auf den Schutz und die Fuersorge der Gemeinschaft.''}\ggmarknotede{art-3}\ggmarknotede{art-6}\ggmarknotede{art-6-1-4}}
\newcommand{\ggfnArtThreeSixZH}{\ggnotezhonce{art-3-6}{第3條：「所有人在法律前一律平等。男女享有平等權利。任何人不得因其性別、血統、種族、語言、故鄉及出身、信仰、宗教或政治見解而受不利益或優惠。」第6條第1項：「婚姻與家庭受國家秩序之特別保護。」第2項第1句：「照護及教育子女，為父母之自然權利，並為首先由其承擔之義務。」第4項：「每一母親均有請求共同體保護與照顧之權利。」}\ggmarknotezh{art-3}\ggmarknotezh{art-6}\ggmarknotezh{art-6-1-4}}

% ===== 組合腳註：Art. 6(1) + Art. 6(4)（墮胎一案特有）=====
\newcommand{\ggfnArtSixOneFourDE}{\ggnotedeonce{art-6-1-4}{Art. 6 Abs. 1: ``Ehe und Familie stehen unter dem besonderen Schutze der staatlichen Ordnung.'' Art. 6 Abs. 4: ``Jede Mutter hat Anspruch auf den Schutz und die Fuersorge der Gemeinschaft.''}\ggmarknotede{art-6}}
\newcommand{\ggfnArtSixOneFourZH}{\ggnotezhonce{art-6-1-4}{第6條第1項：「婚姻與家庭受國家秩序之特別保護。」第4項：「每一母親均有請求共同體保護與照顧之權利。」}\ggmarknotezh{art-6}}

% ===== 組合腳註：Art. 1(1) + Art. 2(2)(1)（墮胎二案特有）=====
\newcommand{\ggfnArtOneOneTwoTwoOneDE}{\ggnotedeonce{art-1-1+2-2-1}{Art. 1 Abs. 1: ``Die Wuerde des Menschen ist unantastbar. Sie zu achten und zu schuetzen ist Verpflichtung aller staatlichen Gewalt.'' Art. 2 Abs. 2 Satz 1: ``Jeder hat das Recht auf Leben und koerperliche Unversehrtheit.''}\ggmarknotede{art-1-1}\ggmarknotede{art-2-2-1}}
\newcommand{\ggfnArtOneOneTwoTwoOneZH}{\ggnotezhonce{art-1-1+2-2-1}{第1條第1項：「人之尊嚴不可侵犯。尊重及保護之，為一切國家權力之義務。」第2條第2項第1句：「人人有生命與身體完整性之權利。」}\ggmarknotezh{art-1-1}\ggmarknotezh{art-2-2-1}}

% ===== 組合腳註：Art. 1(1) + Art. 2(2)（墮胎二案特有）=====
\newcommand{\ggfnArtOneOneTwoTwoDE}{\ggnotedeonce{art-1-1+2-2}{Art. 1 Abs. 1: ``Die Wuerde des Menschen ist unantastbar. Sie zu achten und zu schuetzen ist Verpflichtung aller staatlichen Gewalt.'' Art. 2 Abs. 2: ``Jeder hat das Recht auf Leben und koerperliche Unversehrtheit. Die Freiheit der Person ist unverletzlich. In diese Rechte darf nur auf Grund eines Gesetzes eingegriffen werden.''}\ggmarknotede{art-1-1}\ggmarknotede{art-2-2}\ggmarknotede{art-2-2-1}}
\newcommand{\ggfnArtOneOneTwoTwoZH}{\ggnotezhonce{art-1-1+2-2}{第1條第1項：「人之尊嚴不可侵犯。尊重及保護之，為一切國家權力之義務。」第2條第2項：「人人有生命與身體完整性之權利。人身自由不可侵犯。此等權利僅得依法律加以干預。」}\ggmarknotezh{art-1-1}\ggmarknotezh{art-2-2}\ggmarknotezh{art-2-2-1}}

% ===== 組合腳註：Art. 1(1) + Art. 2(1)（墮胎二案特有）=====
\newcommand{\ggfnArtOneOneTwoOneDE}{\ggnotedeonce{art-1-1+2-1}{Art. 1 Abs. 1: ``Die Wuerde des Menschen ist unantastbar. Sie zu achten und zu schuetzen ist Verpflichtung aller staatlichen Gewalt.'' Art. 2 Abs. 1: ``Jeder hat das Recht auf die freie Entfaltung seiner Persoenlichkeit, soweit er nicht die Rechte anderer verletzt und nicht gegen die verfassungsmaessige Ordnung oder das Sittengesetz verstoesst.''}\ggmarknotede{art-1-1}\ggmarknotede{art-2-1}}
\newcommand{\ggfnArtOneOneTwoOneZH}{\ggnotezhonce{art-1-1+2-1}{第1條第1項：「人之尊嚴不可侵犯。尊重及保護之，為一切國家權力之義務。」第2條第1項：「人人有自由發展其人格之權利，但不得侵害他人權利，亦不得違反合憲秩序或道德律。」}\ggmarknotezh{art-1-1}\ggmarknotezh{art-2-1}}

% ===== 組合腳註：Art. 2(2)(1) + Art. 1(1)（墮胎一案特有，含交叉標記）=====
\newcommand{\ggfnLifeDignityDE}{\ggnotedeonce{life-dignity}{Art. 2 Abs. 2 Satz 1: ``Jeder hat das Recht auf Leben und koerperliche Unversehrtheit.'' Art. 1 Abs. 1: ``Die Wuerde des Menschen ist unantastbar. Sie zu achten und zu schuetzen ist Verpflichtung aller staatlichen Gewalt.''}\ggmarknotede{art-2-2-1}\ggmarknotede{art-1-1}\ggmarknotede{art-1-1-2}}
\newcommand{\ggfnLifeDignityZH}{\ggnotezhonce{life-dignity}{第2條第2項第1句：「人人有生命與身體完整性之權利。」第1條第1項：「人之尊嚴不可侵犯。尊重及保護之，為一切國家權力之義務。」}\ggmarknotezh{art-2-2-1}\ggmarknotezh{art-1-1}\ggmarknotezh{art-1-1-2}}

% ===== 組合腳註：Art. 1(1) + Art. 2(2)（墮胎一案特有，條件判斷版）=====
\newcommand{\ggfnArtOneTwoTwoDE}{\ifcsname ggnoteseen@de@life-dignity\endcsname\ggfnArtTwoTwoRestDE{}\else\ggnotedeonce{art-1-2-2}{Art. 1 Abs. 1: ``Die Wuerde des Menschen ist unantastbar. Sie zu achten und zu schuetzen ist Verpflichtung aller staatlichen Gewalt.'' Art. 2 Abs. 2: ``Jeder hat das Recht auf Leben und koerperliche Unversehrtheit. Die Freiheit der Person ist unverletzlich. In diese Rechte darf nur auf Grund eines Gesetzes eingegriffen werden.''}\ggmarknotede{art-1-1}\ggmarknotede{art-1-1-2}\ggmarknotede{art-2-2}\ggmarknotede{art-2-2-1}\fi}
\newcommand{\ggfnArtOneTwoTwoZH}{\ifcsname ggnoteseen@zh@life-dignity\endcsname\ggfnArtTwoTwoRestZH{}\else\ggnotezhonce{art-1-2-2}{第1條第1項：「人之尊嚴不可侵犯。尊重及保護之，為一切國家權力之義務。」第2條第2項：「人人有生命與身體完整性之權利。人身自由不可侵犯。此等權利僅得依法律加以干預。」}\ggmarknotezh{art-1-1}\ggmarknotezh{art-1-1-2}\ggmarknotezh{art-2-2}\ggmarknotezh{art-2-2-1}\fi}

% ===== 組合腳註：Art. 1(1) + Art. 19(2)（墮胎一案特有，條件判斷版）=====
\newcommand{\ggfnArtOneNineteenDE}{\ifcsname ggnoteseen@de@art-1-1\endcsname\ggfnArtNineteenTwoDE{}\else\ggnotedeonce{art-1-19}{Art. 1 Abs. 1: ``Die Wuerde des Menschen ist unantastbar. Sie zu achten und zu schuetzen ist Verpflichtung aller staatlichen Gewalt.'' Art. 19 Abs. 2: ``In keinem Falle darf ein Grundrecht in seinem Wesensgehalt angetastet werden.''}\ggmarknotede{art-1-1}\ggmarknotede{art-1-1-2}\ggmarknotede{art-19-2}\fi}
\newcommand{\ggfnArtOneNineteenZH}{\ifcsname ggnoteseen@zh@art-1-1\endcsname\ggfnArtNineteenTwoZH{}\else\ggnotezhonce{art-1-19}{第1條第1項：「人之尊嚴不可侵犯。尊重及保護之，為一切國家權力之義務。」第19條第2項：「任何情形下，基本權利之本質內容均不得受到侵害。」}\ggmarknotezh{art-1-1}\ggmarknotezh{art-1-1-2}\ggmarknotezh{art-19-2}\fi}
 之前定義）=====
\newif\ifggversionnotede
\newif\ifggversionnotezh

% ===== 編者註腳基礎設施（首次標記模式）=====
\newif\ifeditornotelabelde
\newif\ifeditornotelabelzh
\newcommand{\editornotelabelde}{\ifeditornotelabelde\else\emph{Hrsg.-Anm. (nachfolgend ohne Kennzeichnung)}: \global\editornotelabeldetrue\fi}
\newcommand{\editornotelabelzh}{\ifeditornotelabelzh\else 編者註（下同）：\global\editornotelabelzhtrue\fi}
\newcommand{\editornotede}[1]{\ifshoweditornotes\footnote{\normalfont\footnotesize\editornotelabelde #1}\else\ifclassroommode\footnote{\normalfont\footnotesize\editornotelabelde #1}\fi\fi}
\newcommand{\editornotezh}[1]{\ifshoweditornotes\footnote{\normalfont\footnotesize\editornotelabelzh #1}\else\ifclassroommode\footnote{\normalfont\footnotesize\editornotelabelzh #1}\fi\fi}

% ===== GG 引用系統（\ggversionde/zh 由各案在 \input 前定義）=====
\newcommand{\ggmarknotede}[1]{\global\expandafter\let\csname ggnoteseen@de@#1\endcsname\empty}
\newcommand{\ggmarknotezh}[1]{\global\expandafter\let\csname ggnoteseen@zh@#1\endcsname\empty}
\newcommand{\ggnotedeonce}[2]{\ifcsname ggnoteseen@de@#1\endcsname\else\global\expandafter\let\csname ggnoteseen@de@#1\endcsname\empty\ggnotede{#2}\fi}
\newcommand{\ggnotezhonce}[2]{\ifcsname ggnoteseen@zh@#1\endcsname\else\global\expandafter\let\csname ggnoteseen@zh@#1\endcsname\empty\ggnotezh{#2}\fi}
\newcommand{\ggnotede}[1]{\editornotede{\ggversionde #1}}
\newcommand{\ggnotezh}[1]{\editornotezh{\ggversionzh #1}}

% ===== Art. 1 =====
\newcommand{\ggfnArtOneOneDE}{\ggnotedeonce{art-1-1}{Art. 1 Abs. 1: ``Die Wuerde des Menschen ist unantastbar. Sie zu achten und zu schuetzen ist Verpflichtung aller staatlichen Gewalt.''}}
\newcommand{\ggfnArtOneOneZH}{\ggnotezhonce{art-1-1}{第1條第1項：「人之尊嚴不可侵犯。尊重及保護之，為一切國家權力之義務。」}}
\newcommand{\ggfnArtOneOneTwoDE}{\ggnotedeonce{art-1-1-2}{Art. 1 Abs. 1 Satz 2: ``Sie zu achten und zu schuetzen ist Verpflichtung aller staatlichen Gewalt.''}}
\newcommand{\ggfnArtOneOneTwoZH}{\ggnotezhonce{art-1-1-2}{第1條第1項第2句：「尊重及保護之，為一切國家權力之義務。」}}

% ===== Art. 2 =====
\newcommand{\ggfnArtTwoOneDE}{\ggnotedeonce{art-2-1}{Art. 2 Abs. 1: ``Jeder hat das Recht auf die freie Entfaltung seiner Persoenlichkeit, soweit er nicht die Rechte anderer verletzt und nicht gegen die verfassungsmaessige Ordnung oder das Sittengesetz verstoesst.''}}
\newcommand{\ggfnArtTwoOneZH}{\ggnotezhonce{art-2-1}{第2條第1項：「人人有自由發展其人格之權利，但不得侵害他人權利，亦不得違反合憲秩序或道德律。」}}
% 簡易預設版（墮胎一案以 \renewcommand 覆寫為含條件邏輯之版本）
\newcommand{\ggfnArtTwoTwoDE}{\ggnotedeonce{art-2-2}{Art. 2 Abs. 2: ``Jeder hat das Recht auf Leben und koerperliche Unversehrtheit. Die Freiheit der Person ist unverletzlich. In diese Rechte darf nur auf Grund eines Gesetzes eingegriffen werden.''}\ggmarknotede{art-2-2-1}}
\newcommand{\ggfnArtTwoTwoZH}{\ggnotezhonce{art-2-2}{第2條第2項：「人人有生命與身體完整性之權利。人身自由不可侵犯。此等權利僅得依法律加以干預。」}\ggmarknotezh{art-2-2-1}}
\newcommand{\ggfnArtTwoTwoOneDE}{\ggnotedeonce{art-2-2-1}{Art. 2 Abs. 2 Satz 1: ``Jeder hat das Recht auf Leben und koerperliche Unversehrtheit.''}}
\newcommand{\ggfnArtTwoTwoOneZH}{\ggnotezhonce{art-2-2-1}{第2條第2項第1句：「人人有生命與身體完整性之權利。」}}
\newcommand{\ggfnArtTwoTwoThreeDE}{\ggnotedeonce{art-2-2-3}{Art. 2 Abs. 2 Satz 3: ``In diese Rechte darf nur auf Grund eines Gesetzes eingegriffen werden.''}}
\newcommand{\ggfnArtTwoTwoThreeZH}{\ggnotezhonce{art-2-2-3}{第2條第2項第3句：「此等權利僅得依法律加以干預。」}}
% 墮胎一案特有：第2條第2項第2–3句（Art. 2(2) 去掉第1句後的「餘文」）
\newcommand{\ggfnArtTwoTwoRestDE}{\ggnotedeonce{art-2-2-rest}{Art. 2 Abs. 2 Saetze 2-3: ``Die Freiheit der Person ist unverletzlich. In diese Rechte darf nur auf Grund eines Gesetzes eingegriffen werden.''}\ggmarknotede{art-2-2}}
\newcommand{\ggfnArtTwoTwoRestZH}{\ggnotezhonce{art-2-2-rest}{第2條第2項第2、3句：「人身自由不可侵犯。此等權利僅得依法律加以干預。」}\ggmarknotezh{art-2-2}}

% ===== Art. 3 =====
\newcommand{\ggfnArtThreeDE}{\ggnotedeonce{art-3}{Art. 3: ``Alle Menschen sind vor dem Gesetz gleich. Maenner und Frauen sind gleichberechtigt. Niemand darf wegen seines Geschlechtes, seiner Abstammung, seiner Rasse, seiner Sprache, seiner Heimat und Herkunft, seines Glaubens, seiner religioesen oder politischen Anschauungen benachteiligt oder bevorzugt werden.''}}
\newcommand{\ggfnArtThreeZH}{\ggnotezhonce{art-3}{第3條：「所有人在法律前一律平等。男女享有平等權利。任何人不得因其性別、血統、種族、語言、故鄉及出身、信仰、宗教或政治見解而受不利益或優惠。」}}
\newcommand{\ggfnArtThreeTwoDE}{\ggnotedeonce{art-3-2}{Art. 3 Abs. 2: ``Maenner und Frauen sind gleichberechtigt.''}}
\newcommand{\ggfnArtThreeTwoZH}{\ggnotezhonce{art-3-2}{第3條第2項：「男女享有平等權利。」}}

% ===== Art. 4 =====
\newcommand{\ggfnArtFourOneDE}{\ggnotedeonce{art-4-1}{Art. 4 Abs. 1: ``Die Freiheit des Glaubens, des Gewissens und die Freiheit des religioesen und weltanschaulichen Bekenntnisses sind unverletzlich.''}}
\newcommand{\ggfnArtFourOneZH}{\ggnotezhonce{art-4-1}{第4條第1項：「信仰自由、良心自由，以及宗教與世界觀信念自由，不可侵犯。」}}

% ===== Art. 5 =====
\newcommand{\ggfnArtFiveOneDE}{\ggnotedeonce{art-5-1}{Art. 5 Abs. 1: ``Jeder hat das Recht, seine Meinung in Wort, Schrift und Bild frei zu aeussern und zu verbreiten und sich aus allgemein zugaenglichen Quellen ungehindert zu unterrichten. Die Pressefreiheit und die Freiheit der Berichterstattung durch Rundfunk und Film werden gewaehrleistet. Eine Zensur findet nicht statt.''}}
\newcommand{\ggfnArtFiveOneZH}{\ggnotezhonce{art-5-1}{第5條第1項：「人人有以言語、文字及圖像自由表達及傳播其意見，並由一般可接近來源不受阻礙獲取資訊之權利。新聞自由及廣播、電影報導自由受保障。不設審查制度。」}}

% ===== Art. 6 =====
\newcommand{\ggfnArtSixDE}{\ggnotedeonce{art-6}{Art. 6 Abs. 1: ``Ehe und Familie stehen unter dem besonderen Schutze der staatlichen Ordnung.'' Abs. 2 Satz 1: ``Pflege und Erziehung der Kinder sind das natuerliche Recht der Eltern und die zuvoerderst ihnen obliegende Pflicht.'' Abs. 4: ``Jede Mutter hat Anspruch auf den Schutz und die Fuersorge der Gemeinschaft.''}}
\newcommand{\ggfnArtSixZH}{\ggnotezhonce{art-6}{第6條第1項：「婚姻與家庭受國家秩序之特別保護。」第2項第1句：「照護及教育子女，為父母之自然權利，並為首先由其承擔之義務。」第4項：「每一母親均有請求共同體保護與照顧之權利。」}}
\newcommand{\ggfnArtSixOneDE}{\ggnotedeonce{art-6-1}{Art. 6 Abs. 1: ``Ehe und Familie stehen unter dem besonderen Schutze der staatlichen Ordnung.''}}
\newcommand{\ggfnArtSixOneZH}{\ggnotezhonce{art-6-1}{第6條第1項：「婚姻與家庭受國家秩序之特別保護。」}}
\newcommand{\ggfnArtSixTwoDE}{\ggnotedeonce{art-6-2}{Art. 6 Abs. 2: ``Pflege und Erziehung der Kinder sind das natuerliche Recht der Eltern und die zuvoerderst ihnen obliegende Pflicht. Ueber ihre Betaetigung wacht die staatliche Gemeinschaft.''}}
\newcommand{\ggfnArtSixTwoZH}{\ggnotezhonce{art-6-2}{第6條第2項：「照護及教育子女，為父母之自然權利，並為首先由其承擔之義務。國家共同體監督其行使。」}}
\newcommand{\ggfnArtSixTwoOneDE}{\ggnotedeonce{art-6-2-1}{Art. 6 Abs. 2 Satz 1: ``Pflege und Erziehung der Kinder sind das natuerliche Recht der Eltern und die zuvoerderst ihnen obliegende Pflicht.''}}
\newcommand{\ggfnArtSixTwoOneZH}{\ggnotezhonce{art-6-2-1}{第6條第2項第1句：「照護及教育子女，為父母之自然權利，並為首先由其承擔之義務。」}}
\newcommand{\ggfnArtSixFourDE}{\ggnotedeonce{art-6-4}{Art. 6 Abs. 4: ``Jede Mutter hat Anspruch auf den Schutz und die Fuersorge der Gemeinschaft.''}}
\newcommand{\ggfnArtSixFourZH}{\ggnotezhonce{art-6-4}{第6條第4項：「每一母親均有請求共同體保護與照顧之權利。」}}

% ===== Art. 12 =====
\newcommand{\ggfnArtTwelveOneDE}{\ggnotedeonce{art-12-1}{Art. 12 Abs. 1: ``Alle Deutschen haben das Recht, Beruf, Arbeitsplatz und Ausbildungsstaette frei zu waehlen. Die Berufsausuebung kann durch Gesetz oder auf Grund eines Gesetzes geregelt werden.''}}
\newcommand{\ggfnArtTwelveOneZH}{\ggnotezhonce{art-12-1}{第12條第1項：「所有德國人均有自由選擇職業、工作場所及訓練場所之權利。職業行使得以法律或基於法律加以規範。」}}

% ===== Art. 19 =====
\newcommand{\ggfnArtNineteenTwoDE}{\ggnotedeonce{art-19-2}{Art. 19 Abs. 2: ``In keinem Falle darf ein Grundrecht in seinem Wesensgehalt angetastet werden.''}}
\newcommand{\ggfnArtNineteenTwoZH}{\ggnotezhonce{art-19-2}{第19條第2項：「任何情形下，基本權利之本質內容均不得受到侵害。」}}

% ===== Art. 20 =====
\newcommand{\ggfnArtTwentyOneDE}{\ggnotedeonce{art-20-1}{Art. 20 Abs. 1: ``Die Bundesrepublik Deutschland ist ein demokratischer und sozialer Bundesstaat.''}}
\newcommand{\ggfnArtTwentyOneZH}{\ggnotezhonce{art-20-1}{第20條第1項：「德意志聯邦共和國為民主及社會之聯邦國。」}}
\newcommand{\ggfnArtTwentyThreeDE}{\ggnotedeonce{art-20-3}{Art. 20 Abs. 3: ``Die Gesetzgebung ist an die verfassungsmaessige Ordnung, die vollziehende Gewalt und die Rechtsprechung sind an Gesetz und Recht gebunden.''}}
\newcommand{\ggfnArtTwentyThreeZH}{\ggnotezhonce{art-20-3}{第20條第3項：「立法受合憲秩序拘束；行政權與司法權受法律與法拘束。」}}

% ===== Art. 26 + Art. 143（墮胎一案特有）=====
\newcommand{\ggfnArtTwoSixAndOneFourThreeDE}{\ggnotedeonce{art-26-143}{Art. 26 Abs. 1 Satz 2: ``Sie sind unter Strafe zu stellen.'' Art. 143 in der urspruenglichen Fassung (24. Mai 1949 bis 1. September 1951) enthielt Strafvorschriften zum Hochverrat; 1975 war Art. 143 bereits weggefallen.}}
\newcommand{\ggfnArtTwoSixAndOneFourThreeZH}{\ggnotezhonce{art-26-143}{第26條第1項第2句：「此等行為應規定為犯罪處罰。」第143條原始版本（1949年5月24日至1951年9月1日）含有關於叛國罪之刑罰規定；至1975年時，第143條已廢止。}}

% ===== Art. 28 =====
\newcommand{\ggfnArtTwentyEightOneDE}{\ggnotedeonce{art-28-1-1}{Art. 28 Abs. 1 Satz 1: ``Die verfassungsmaessige Ordnung in den Laendern muss den Grundsaetzen des republikanischen, demokratischen und sozialen Rechtsstaates im Sinne dieses Grundgesetzes entsprechen.''}}
\newcommand{\ggfnArtTwentyEightOneZH}{\ggnotezhonce{art-28-1-1}{第28條第1項第1句：「各邦之合憲秩序，須符合本基本法意義下共和、民主及社會法治國之原則。」}}

% ===== Art. 30 =====
\newcommand{\ggfnArtThirtyDE}{\ggnotedeonce{art-30}{Art. 30: ``Die Ausuebung der staatlichen Befugnisse und die Erfuellung der staatlichen Aufgaben ist Sache der Laender, soweit dieses Grundgesetz keine andere Regelung trifft oder zulaesst.''}}
\newcommand{\ggfnArtThirtyZH}{\ggnotezhonce{art-30}{第30條：「國家權限之行使與國家任務之履行，屬於各邦，但本基本法另有規定或容許者，不在此限。」}}

% ===== Art. 74 =====
\newcommand{\ggfnArtSeventyFourOneDE}{\ggnotedeonce{art-74-1}{Art. 74 Nr. 1: Die konkurrierende Gesetzgebung erstreckt sich auf ``das buergerliche Recht, das Strafrecht und den Strafvollzug, die Gerichtsverfassung, das gerichtliche Verfahren, die Rechtsanwaltschaft, das Notariat und die Rechtsberatung''.}}
\newcommand{\ggfnArtSeventyFourOneZH}{\ggnotezhonce{art-74-1}{第74條第1款：競合立法及於「民法、刑法及刑罰執行、法院組織、訴訟程序、律師制度、公證制度及法律諮詢」。}}
\newcommand{\ggfnArtSeventyFourSevenDE}{\ggnotedeonce{art-74-7}{Art. 74 Nr. 7: Die konkurrierende Gesetzgebung erstreckt sich auf ``die oeffentliche Fuersorge''.}}
\newcommand{\ggfnArtSeventyFourSevenZH}{\ggnotezhonce{art-74-7}{第74條第7款：競合立法及於「公共扶助」。}}
\newcommand{\ggfnArtSeventyFourTwelveDE}{\ggnotedeonce{art-74-12}{Art. 74 Nr. 12: Die konkurrierende Gesetzgebung erstreckt sich auf ``das Arbeitsrecht einschliesslich der Betriebsverfassung, des Arbeitsschutzes und der Arbeitsvermittlung sowie die Sozialversicherung einschliesslich der Arbeitslosenversicherung''.}}
\newcommand{\ggfnArtSeventyFourTwelveZH}{\ggnotezhonce{art-74-12}{第74條第12款：競合立法及於「勞動法，包括企業組織、勞動保護及就業媒合，以及社會保險，包括失業保險」。}}

% ===== Art. 77（墮胎一案特有）=====
\newcommand{\ggfnArtSevenSevenThreeDE}{\ggnotedeonce{art-77-3}{Art. 77 Abs. 3 Satz 1: ``Soweit zu einem Gesetze die Zustimmung des Bundesrates nicht erforderlich ist, kann der Bundesrat, wenn das Verfahren nach Absatz 2 beendigt ist, gegen ein vom Bundestage beschlossenes Gesetz binnen zwei Wochen Einspruch einlegen.''}}
\newcommand{\ggfnArtSevenSevenThreeZH}{\ggnotezhonce{art-77-3}{第77條第3項第1句：「法律無須聯邦參議院同意者，於第2項程序終了後，聯邦參議院得於二週內對聯邦議院通過之法律提出異議。」}}

% ===== Art. 80 =====
\newcommand{\ggfnArtEightyOneOneDE}{\ggnotedeonce{art-80-1-1}{Art. 80 Abs. 1 Satz 1: ``Durch Gesetz koennen die Bundesregierung, ein Bundesminister oder die Landesregierungen ermaechtigt werden, Rechtsverordnungen zu erlassen.''}}
\newcommand{\ggfnArtEightyOneOneZH}{\ggnotezhonce{art-80-1-1}{第80條第1項第1句：「得以法律授權聯邦政府、聯邦部長或邦政府發布法規命令。」}}

% ===== Art. 83 =====
\newcommand{\ggfnArtEightyThreeDE}{\ggnotedeonce{art-83}{Art. 83: ``Die Laender fuehren die Bundesgesetze als eigene Angelegenheit aus, soweit dieses Grundgesetz nichts anderes bestimmt oder zulaesst.''}}
\newcommand{\ggfnArtEightyThreeZH}{\ggnotezhonce{art-83}{第83條：「各邦以自己事務執行聯邦法律，但本基本法另有規定或容許者，不在此限。」}}

% ===== Art. 84 =====
\newcommand{\ggfnArtEightyFourOneDE}{\ggnotedeonce{art-84-1}{Art. 84 Abs. 1: ``Fuehren die Laender die Bundesgesetze als eigene Angelegenheit aus, so regeln sie die Einrichtung der Behoerden und das Verwaltungsverfahren, soweit nicht Bundesgesetze mit Zustimmung des Bundesrates etwas anderes bestimmen.''}}
\newcommand{\ggfnArtEightyFourOneZH}{\ggnotezhonce{art-84-1}{第84條第1項：「各邦以自己事務執行聯邦法律時，除聯邦法律經聯邦參議院同意另有規定外，由各邦規範機關之設置及行政程序。」}}
% 別名（墮胎一案使用之縮寫命名）
\let\ggfnArtEightFourOneDE\ggfnArtEightyFourOneDE
\let\ggfnArtEightFourOneZH\ggfnArtEightyFourOneZH

% ===== Art. 93 =====
\newcommand{\ggfnArtNinetyThreeOneTwoDE}{\ggnotedeonce{art-93-1-2}{Art. 93 Abs. 1 Nr. 2: Das Bundesverfassungsgericht entscheidet ``bei Meinungsverschiedenheiten oder Zweifeln ueber die foermliche und sachliche Vereinbarkeit von Bundesrecht oder Landesrecht mit diesem Grundgesetze ... auf Antrag der Bundesregierung, einer Landesregierung oder eines Drittels der Mitglieder des Bundestages''.}}
\newcommand{\ggfnArtNinetyThreeOneTwoZH}{\ggnotezhonce{art-93-1-2}{第93條第1項第2款：聯邦憲法法院就「聯邦法或邦法在形式及實質上是否與本基本法相容之意見分歧或疑義……依聯邦政府、邦政府或聯邦議院三分之一議員之聲請」作成裁判。}}
% 別名（墮胎一案使用之縮寫命名）
\let\ggfnArtNineThreeOneTwoDE\ggfnArtNinetyThreeOneTwoDE
\let\ggfnArtNineThreeOneTwoZH\ggfnArtNinetyThreeOneTwoZH

% ===== Art. 102 =====
\newcommand{\ggfnArtOneZeroTwoDE}{\ggnotedeonce{art-102}{Art. 102: ``Die Todesstrafe ist abgeschafft.''}}
\newcommand{\ggfnArtOneZeroTwoZH}{\ggnotezhonce{art-102}{第102條：「死刑廢止。」}}

% ===== Art. 103 =====
\newcommand{\ggfnArtOneZeroThreeTwoDE}{\ggnotedeonce{art-103-2}{Art. 103 Abs. 2: ``Eine Tat kann nur bestraft werden, wenn die Strafbarkeit gesetzlich bestimmt war, bevor die Tat begangen wurde.''}}
\newcommand{\ggfnArtOneZeroThreeTwoZH}{\ggnotezhonce{art-103-2}{第103條第2項：「行為之可罰性，須於行為前以法律定之，始得處罰。」}}

% ===== Art. 104 =====
\newcommand{\ggfnArtOneZeroFourOneDE}{\ggnotedeonce{art-104-1}{Art. 104 Abs. 1: ``Die Freiheit der Person kann nur auf Grund eines foermlichen Gesetzes und nur unter Beachtung der darin vorgeschriebenen Formen beschraenkt werden.''}}
\newcommand{\ggfnArtOneZeroFourOneZH}{\ggnotezhonce{art-104-1}{第104條第1項：「人身自由僅得基於形式法律，並遵守其中規定之形式加以限制。」}}

% ===== 組合腳註：Art. 3 + Art. 6（墮胎一案特有）=====
\newcommand{\ggfnArtThreeSixDE}{\ggnotedeonce{art-3-6}{Art. 3: ``Alle Menschen sind vor dem Gesetz gleich. Maenner und Frauen sind gleichberechtigt. Niemand darf wegen seines Geschlechtes, seiner Abstammung, seiner Rasse, seiner Sprache, seiner Heimat und Herkunft, seines Glaubens, seiner religioesen oder politischen Anschauungen benachteiligt oder bevorzugt werden.'' Art. 6 Abs. 1: ``Ehe und Familie stehen unter dem besonderen Schutze der staatlichen Ordnung.'' Abs. 2 Satz 1: ``Pflege und Erziehung der Kinder sind das natuerliche Recht der Eltern und die zuvoerderst ihnen obliegende Pflicht.'' Abs. 4: ``Jede Mutter hat Anspruch auf den Schutz und die Fuersorge der Gemeinschaft.''}\ggmarknotede{art-3}\ggmarknotede{art-6}\ggmarknotede{art-6-1-4}}
\newcommand{\ggfnArtThreeSixZH}{\ggnotezhonce{art-3-6}{第3條：「所有人在法律前一律平等。男女享有平等權利。任何人不得因其性別、血統、種族、語言、故鄉及出身、信仰、宗教或政治見解而受不利益或優惠。」第6條第1項：「婚姻與家庭受國家秩序之特別保護。」第2項第1句：「照護及教育子女，為父母之自然權利，並為首先由其承擔之義務。」第4項：「每一母親均有請求共同體保護與照顧之權利。」}\ggmarknotezh{art-3}\ggmarknotezh{art-6}\ggmarknotezh{art-6-1-4}}

% ===== 組合腳註：Art. 6(1) + Art. 6(4)（墮胎一案特有）=====
\newcommand{\ggfnArtSixOneFourDE}{\ggnotedeonce{art-6-1-4}{Art. 6 Abs. 1: ``Ehe und Familie stehen unter dem besonderen Schutze der staatlichen Ordnung.'' Art. 6 Abs. 4: ``Jede Mutter hat Anspruch auf den Schutz und die Fuersorge der Gemeinschaft.''}\ggmarknotede{art-6}}
\newcommand{\ggfnArtSixOneFourZH}{\ggnotezhonce{art-6-1-4}{第6條第1項：「婚姻與家庭受國家秩序之特別保護。」第4項：「每一母親均有請求共同體保護與照顧之權利。」}\ggmarknotezh{art-6}}

% ===== 組合腳註：Art. 1(1) + Art. 2(2)(1)（墮胎二案特有）=====
\newcommand{\ggfnArtOneOneTwoTwoOneDE}{\ggnotedeonce{art-1-1+2-2-1}{Art. 1 Abs. 1: ``Die Wuerde des Menschen ist unantastbar. Sie zu achten und zu schuetzen ist Verpflichtung aller staatlichen Gewalt.'' Art. 2 Abs. 2 Satz 1: ``Jeder hat das Recht auf Leben und koerperliche Unversehrtheit.''}\ggmarknotede{art-1-1}\ggmarknotede{art-2-2-1}}
\newcommand{\ggfnArtOneOneTwoTwoOneZH}{\ggnotezhonce{art-1-1+2-2-1}{第1條第1項：「人之尊嚴不可侵犯。尊重及保護之，為一切國家權力之義務。」第2條第2項第1句：「人人有生命與身體完整性之權利。」}\ggmarknotezh{art-1-1}\ggmarknotezh{art-2-2-1}}

% ===== 組合腳註：Art. 1(1) + Art. 2(2)（墮胎二案特有）=====
\newcommand{\ggfnArtOneOneTwoTwoDE}{\ggnotedeonce{art-1-1+2-2}{Art. 1 Abs. 1: ``Die Wuerde des Menschen ist unantastbar. Sie zu achten und zu schuetzen ist Verpflichtung aller staatlichen Gewalt.'' Art. 2 Abs. 2: ``Jeder hat das Recht auf Leben und koerperliche Unversehrtheit. Die Freiheit der Person ist unverletzlich. In diese Rechte darf nur auf Grund eines Gesetzes eingegriffen werden.''}\ggmarknotede{art-1-1}\ggmarknotede{art-2-2}\ggmarknotede{art-2-2-1}}
\newcommand{\ggfnArtOneOneTwoTwoZH}{\ggnotezhonce{art-1-1+2-2}{第1條第1項：「人之尊嚴不可侵犯。尊重及保護之，為一切國家權力之義務。」第2條第2項：「人人有生命與身體完整性之權利。人身自由不可侵犯。此等權利僅得依法律加以干預。」}\ggmarknotezh{art-1-1}\ggmarknotezh{art-2-2}\ggmarknotezh{art-2-2-1}}

% ===== 組合腳註：Art. 1(1) + Art. 2(1)（墮胎二案特有）=====
\newcommand{\ggfnArtOneOneTwoOneDE}{\ggnotedeonce{art-1-1+2-1}{Art. 1 Abs. 1: ``Die Wuerde des Menschen ist unantastbar. Sie zu achten und zu schuetzen ist Verpflichtung aller staatlichen Gewalt.'' Art. 2 Abs. 1: ``Jeder hat das Recht auf die freie Entfaltung seiner Persoenlichkeit, soweit er nicht die Rechte anderer verletzt und nicht gegen die verfassungsmaessige Ordnung oder das Sittengesetz verstoesst.''}\ggmarknotede{art-1-1}\ggmarknotede{art-2-1}}
\newcommand{\ggfnArtOneOneTwoOneZH}{\ggnotezhonce{art-1-1+2-1}{第1條第1項：「人之尊嚴不可侵犯。尊重及保護之，為一切國家權力之義務。」第2條第1項：「人人有自由發展其人格之權利，但不得侵害他人權利，亦不得違反合憲秩序或道德律。」}\ggmarknotezh{art-1-1}\ggmarknotezh{art-2-1}}

% ===== 組合腳註：Art. 2(2)(1) + Art. 1(1)（墮胎一案特有，含交叉標記）=====
\newcommand{\ggfnLifeDignityDE}{\ggnotedeonce{life-dignity}{Art. 2 Abs. 2 Satz 1: ``Jeder hat das Recht auf Leben und koerperliche Unversehrtheit.'' Art. 1 Abs. 1: ``Die Wuerde des Menschen ist unantastbar. Sie zu achten und zu schuetzen ist Verpflichtung aller staatlichen Gewalt.''}\ggmarknotede{art-2-2-1}\ggmarknotede{art-1-1}\ggmarknotede{art-1-1-2}}
\newcommand{\ggfnLifeDignityZH}{\ggnotezhonce{life-dignity}{第2條第2項第1句：「人人有生命與身體完整性之權利。」第1條第1項：「人之尊嚴不可侵犯。尊重及保護之，為一切國家權力之義務。」}\ggmarknotezh{art-2-2-1}\ggmarknotezh{art-1-1}\ggmarknotezh{art-1-1-2}}

% ===== 組合腳註：Art. 1(1) + Art. 2(2)（墮胎一案特有，條件判斷版）=====
\newcommand{\ggfnArtOneTwoTwoDE}{\ifcsname ggnoteseen@de@life-dignity\endcsname\ggfnArtTwoTwoRestDE{}\else\ggnotedeonce{art-1-2-2}{Art. 1 Abs. 1: ``Die Wuerde des Menschen ist unantastbar. Sie zu achten und zu schuetzen ist Verpflichtung aller staatlichen Gewalt.'' Art. 2 Abs. 2: ``Jeder hat das Recht auf Leben und koerperliche Unversehrtheit. Die Freiheit der Person ist unverletzlich. In diese Rechte darf nur auf Grund eines Gesetzes eingegriffen werden.''}\ggmarknotede{art-1-1}\ggmarknotede{art-1-1-2}\ggmarknotede{art-2-2}\ggmarknotede{art-2-2-1}\fi}
\newcommand{\ggfnArtOneTwoTwoZH}{\ifcsname ggnoteseen@zh@life-dignity\endcsname\ggfnArtTwoTwoRestZH{}\else\ggnotezhonce{art-1-2-2}{第1條第1項：「人之尊嚴不可侵犯。尊重及保護之，為一切國家權力之義務。」第2條第2項：「人人有生命與身體完整性之權利。人身自由不可侵犯。此等權利僅得依法律加以干預。」}\ggmarknotezh{art-1-1}\ggmarknotezh{art-1-1-2}\ggmarknotezh{art-2-2}\ggmarknotezh{art-2-2-1}\fi}

% ===== 組合腳註：Art. 1(1) + Art. 19(2)（墮胎一案特有，條件判斷版）=====
\newcommand{\ggfnArtOneNineteenDE}{\ifcsname ggnoteseen@de@art-1-1\endcsname\ggfnArtNineteenTwoDE{}\else\ggnotedeonce{art-1-19}{Art. 1 Abs. 1: ``Die Wuerde des Menschen ist unantastbar. Sie zu achten und zu schuetzen ist Verpflichtung aller staatlichen Gewalt.'' Art. 19 Abs. 2: ``In keinem Falle darf ein Grundrecht in seinem Wesensgehalt angetastet werden.''}\ggmarknotede{art-1-1}\ggmarknotede{art-1-1-2}\ggmarknotede{art-19-2}\fi}
\newcommand{\ggfnArtOneNineteenZH}{\ifcsname ggnoteseen@zh@art-1-1\endcsname\ggfnArtNineteenTwoZH{}\else\ggnotezhonce{art-1-19}{第1條第1項：「人之尊嚴不可侵犯。尊重及保護之，為一切國家權力之義務。」第19條第2項：「任何情形下，基本權利之本質內容均不得受到侵害。」}\ggmarknotezh{art-1-1}\ggmarknotezh{art-1-1-2}\ggmarknotezh{art-19-2}\fi}


% ===== GG 版本旗標（\ggversionde/zh 由各案在 % bverfge_gg.tex — 德國墮胎案 GG 條文腳註共用定義
% 使用方式：在各案 .tex 中先定義 \ggversionde/\ggversionzh，再 % bverfge_gg.tex — 德國墮胎案 GG 條文腳註共用定義
% 使用方式：在各案 .tex 中先定義 \ggversionde/\ggversionzh，再 \input{../../共用LaTeX/bverfge_gg}

% ===== GG 版本旗標（\ggversionde/zh 由各案在 \input{../../共用LaTeX/bverfge_gg} 之前定義）=====
\newif\ifggversionnotede
\newif\ifggversionnotezh

% ===== 編者註腳基礎設施（首次標記模式）=====
\newif\ifeditornotelabelde
\newif\ifeditornotelabelzh
\newcommand{\editornotelabelde}{\ifeditornotelabelde\else\emph{Hrsg.-Anm. (nachfolgend ohne Kennzeichnung)}: \global\editornotelabeldetrue\fi}
\newcommand{\editornotelabelzh}{\ifeditornotelabelzh\else 編者註（下同）：\global\editornotelabelzhtrue\fi}
\newcommand{\editornotede}[1]{\ifshoweditornotes\footnote{\normalfont\footnotesize\editornotelabelde #1}\else\ifclassroommode\footnote{\normalfont\footnotesize\editornotelabelde #1}\fi\fi}
\newcommand{\editornotezh}[1]{\ifshoweditornotes\footnote{\normalfont\footnotesize\editornotelabelzh #1}\else\ifclassroommode\footnote{\normalfont\footnotesize\editornotelabelzh #1}\fi\fi}

% ===== GG 引用系統（\ggversionde/zh 由各案在 \input 前定義）=====
\newcommand{\ggmarknotede}[1]{\global\expandafter\let\csname ggnoteseen@de@#1\endcsname\empty}
\newcommand{\ggmarknotezh}[1]{\global\expandafter\let\csname ggnoteseen@zh@#1\endcsname\empty}
\newcommand{\ggnotedeonce}[2]{\ifcsname ggnoteseen@de@#1\endcsname\else\global\expandafter\let\csname ggnoteseen@de@#1\endcsname\empty\ggnotede{#2}\fi}
\newcommand{\ggnotezhonce}[2]{\ifcsname ggnoteseen@zh@#1\endcsname\else\global\expandafter\let\csname ggnoteseen@zh@#1\endcsname\empty\ggnotezh{#2}\fi}
\newcommand{\ggnotede}[1]{\editornotede{\ggversionde #1}}
\newcommand{\ggnotezh}[1]{\editornotezh{\ggversionzh #1}}

% ===== Art. 1 =====
\newcommand{\ggfnArtOneOneDE}{\ggnotedeonce{art-1-1}{Art. 1 Abs. 1: ``Die Wuerde des Menschen ist unantastbar. Sie zu achten und zu schuetzen ist Verpflichtung aller staatlichen Gewalt.''}}
\newcommand{\ggfnArtOneOneZH}{\ggnotezhonce{art-1-1}{第1條第1項：「人之尊嚴不可侵犯。尊重及保護之，為一切國家權力之義務。」}}
\newcommand{\ggfnArtOneOneTwoDE}{\ggnotedeonce{art-1-1-2}{Art. 1 Abs. 1 Satz 2: ``Sie zu achten und zu schuetzen ist Verpflichtung aller staatlichen Gewalt.''}}
\newcommand{\ggfnArtOneOneTwoZH}{\ggnotezhonce{art-1-1-2}{第1條第1項第2句：「尊重及保護之，為一切國家權力之義務。」}}

% ===== Art. 2 =====
\newcommand{\ggfnArtTwoOneDE}{\ggnotedeonce{art-2-1}{Art. 2 Abs. 1: ``Jeder hat das Recht auf die freie Entfaltung seiner Persoenlichkeit, soweit er nicht die Rechte anderer verletzt und nicht gegen die verfassungsmaessige Ordnung oder das Sittengesetz verstoesst.''}}
\newcommand{\ggfnArtTwoOneZH}{\ggnotezhonce{art-2-1}{第2條第1項：「人人有自由發展其人格之權利，但不得侵害他人權利，亦不得違反合憲秩序或道德律。」}}
% 簡易預設版（墮胎一案以 \renewcommand 覆寫為含條件邏輯之版本）
\newcommand{\ggfnArtTwoTwoDE}{\ggnotedeonce{art-2-2}{Art. 2 Abs. 2: ``Jeder hat das Recht auf Leben und koerperliche Unversehrtheit. Die Freiheit der Person ist unverletzlich. In diese Rechte darf nur auf Grund eines Gesetzes eingegriffen werden.''}\ggmarknotede{art-2-2-1}}
\newcommand{\ggfnArtTwoTwoZH}{\ggnotezhonce{art-2-2}{第2條第2項：「人人有生命與身體完整性之權利。人身自由不可侵犯。此等權利僅得依法律加以干預。」}\ggmarknotezh{art-2-2-1}}
\newcommand{\ggfnArtTwoTwoOneDE}{\ggnotedeonce{art-2-2-1}{Art. 2 Abs. 2 Satz 1: ``Jeder hat das Recht auf Leben und koerperliche Unversehrtheit.''}}
\newcommand{\ggfnArtTwoTwoOneZH}{\ggnotezhonce{art-2-2-1}{第2條第2項第1句：「人人有生命與身體完整性之權利。」}}
\newcommand{\ggfnArtTwoTwoThreeDE}{\ggnotedeonce{art-2-2-3}{Art. 2 Abs. 2 Satz 3: ``In diese Rechte darf nur auf Grund eines Gesetzes eingegriffen werden.''}}
\newcommand{\ggfnArtTwoTwoThreeZH}{\ggnotezhonce{art-2-2-3}{第2條第2項第3句：「此等權利僅得依法律加以干預。」}}
% 墮胎一案特有：第2條第2項第2–3句（Art. 2(2) 去掉第1句後的「餘文」）
\newcommand{\ggfnArtTwoTwoRestDE}{\ggnotedeonce{art-2-2-rest}{Art. 2 Abs. 2 Saetze 2-3: ``Die Freiheit der Person ist unverletzlich. In diese Rechte darf nur auf Grund eines Gesetzes eingegriffen werden.''}\ggmarknotede{art-2-2}}
\newcommand{\ggfnArtTwoTwoRestZH}{\ggnotezhonce{art-2-2-rest}{第2條第2項第2、3句：「人身自由不可侵犯。此等權利僅得依法律加以干預。」}\ggmarknotezh{art-2-2}}

% ===== Art. 3 =====
\newcommand{\ggfnArtThreeDE}{\ggnotedeonce{art-3}{Art. 3: ``Alle Menschen sind vor dem Gesetz gleich. Maenner und Frauen sind gleichberechtigt. Niemand darf wegen seines Geschlechtes, seiner Abstammung, seiner Rasse, seiner Sprache, seiner Heimat und Herkunft, seines Glaubens, seiner religioesen oder politischen Anschauungen benachteiligt oder bevorzugt werden.''}}
\newcommand{\ggfnArtThreeZH}{\ggnotezhonce{art-3}{第3條：「所有人在法律前一律平等。男女享有平等權利。任何人不得因其性別、血統、種族、語言、故鄉及出身、信仰、宗教或政治見解而受不利益或優惠。」}}
\newcommand{\ggfnArtThreeTwoDE}{\ggnotedeonce{art-3-2}{Art. 3 Abs. 2: ``Maenner und Frauen sind gleichberechtigt.''}}
\newcommand{\ggfnArtThreeTwoZH}{\ggnotezhonce{art-3-2}{第3條第2項：「男女享有平等權利。」}}

% ===== Art. 4 =====
\newcommand{\ggfnArtFourOneDE}{\ggnotedeonce{art-4-1}{Art. 4 Abs. 1: ``Die Freiheit des Glaubens, des Gewissens und die Freiheit des religioesen und weltanschaulichen Bekenntnisses sind unverletzlich.''}}
\newcommand{\ggfnArtFourOneZH}{\ggnotezhonce{art-4-1}{第4條第1項：「信仰自由、良心自由，以及宗教與世界觀信念自由，不可侵犯。」}}

% ===== Art. 5 =====
\newcommand{\ggfnArtFiveOneDE}{\ggnotedeonce{art-5-1}{Art. 5 Abs. 1: ``Jeder hat das Recht, seine Meinung in Wort, Schrift und Bild frei zu aeussern und zu verbreiten und sich aus allgemein zugaenglichen Quellen ungehindert zu unterrichten. Die Pressefreiheit und die Freiheit der Berichterstattung durch Rundfunk und Film werden gewaehrleistet. Eine Zensur findet nicht statt.''}}
\newcommand{\ggfnArtFiveOneZH}{\ggnotezhonce{art-5-1}{第5條第1項：「人人有以言語、文字及圖像自由表達及傳播其意見，並由一般可接近來源不受阻礙獲取資訊之權利。新聞自由及廣播、電影報導自由受保障。不設審查制度。」}}

% ===== Art. 6 =====
\newcommand{\ggfnArtSixDE}{\ggnotedeonce{art-6}{Art. 6 Abs. 1: ``Ehe und Familie stehen unter dem besonderen Schutze der staatlichen Ordnung.'' Abs. 2 Satz 1: ``Pflege und Erziehung der Kinder sind das natuerliche Recht der Eltern und die zuvoerderst ihnen obliegende Pflicht.'' Abs. 4: ``Jede Mutter hat Anspruch auf den Schutz und die Fuersorge der Gemeinschaft.''}}
\newcommand{\ggfnArtSixZH}{\ggnotezhonce{art-6}{第6條第1項：「婚姻與家庭受國家秩序之特別保護。」第2項第1句：「照護及教育子女，為父母之自然權利，並為首先由其承擔之義務。」第4項：「每一母親均有請求共同體保護與照顧之權利。」}}
\newcommand{\ggfnArtSixOneDE}{\ggnotedeonce{art-6-1}{Art. 6 Abs. 1: ``Ehe und Familie stehen unter dem besonderen Schutze der staatlichen Ordnung.''}}
\newcommand{\ggfnArtSixOneZH}{\ggnotezhonce{art-6-1}{第6條第1項：「婚姻與家庭受國家秩序之特別保護。」}}
\newcommand{\ggfnArtSixTwoDE}{\ggnotedeonce{art-6-2}{Art. 6 Abs. 2: ``Pflege und Erziehung der Kinder sind das natuerliche Recht der Eltern und die zuvoerderst ihnen obliegende Pflicht. Ueber ihre Betaetigung wacht die staatliche Gemeinschaft.''}}
\newcommand{\ggfnArtSixTwoZH}{\ggnotezhonce{art-6-2}{第6條第2項：「照護及教育子女，為父母之自然權利，並為首先由其承擔之義務。國家共同體監督其行使。」}}
\newcommand{\ggfnArtSixTwoOneDE}{\ggnotedeonce{art-6-2-1}{Art. 6 Abs. 2 Satz 1: ``Pflege und Erziehung der Kinder sind das natuerliche Recht der Eltern und die zuvoerderst ihnen obliegende Pflicht.''}}
\newcommand{\ggfnArtSixTwoOneZH}{\ggnotezhonce{art-6-2-1}{第6條第2項第1句：「照護及教育子女，為父母之自然權利，並為首先由其承擔之義務。」}}
\newcommand{\ggfnArtSixFourDE}{\ggnotedeonce{art-6-4}{Art. 6 Abs. 4: ``Jede Mutter hat Anspruch auf den Schutz und die Fuersorge der Gemeinschaft.''}}
\newcommand{\ggfnArtSixFourZH}{\ggnotezhonce{art-6-4}{第6條第4項：「每一母親均有請求共同體保護與照顧之權利。」}}

% ===== Art. 12 =====
\newcommand{\ggfnArtTwelveOneDE}{\ggnotedeonce{art-12-1}{Art. 12 Abs. 1: ``Alle Deutschen haben das Recht, Beruf, Arbeitsplatz und Ausbildungsstaette frei zu waehlen. Die Berufsausuebung kann durch Gesetz oder auf Grund eines Gesetzes geregelt werden.''}}
\newcommand{\ggfnArtTwelveOneZH}{\ggnotezhonce{art-12-1}{第12條第1項：「所有德國人均有自由選擇職業、工作場所及訓練場所之權利。職業行使得以法律或基於法律加以規範。」}}

% ===== Art. 19 =====
\newcommand{\ggfnArtNineteenTwoDE}{\ggnotedeonce{art-19-2}{Art. 19 Abs. 2: ``In keinem Falle darf ein Grundrecht in seinem Wesensgehalt angetastet werden.''}}
\newcommand{\ggfnArtNineteenTwoZH}{\ggnotezhonce{art-19-2}{第19條第2項：「任何情形下，基本權利之本質內容均不得受到侵害。」}}

% ===== Art. 20 =====
\newcommand{\ggfnArtTwentyOneDE}{\ggnotedeonce{art-20-1}{Art. 20 Abs. 1: ``Die Bundesrepublik Deutschland ist ein demokratischer und sozialer Bundesstaat.''}}
\newcommand{\ggfnArtTwentyOneZH}{\ggnotezhonce{art-20-1}{第20條第1項：「德意志聯邦共和國為民主及社會之聯邦國。」}}
\newcommand{\ggfnArtTwentyThreeDE}{\ggnotedeonce{art-20-3}{Art. 20 Abs. 3: ``Die Gesetzgebung ist an die verfassungsmaessige Ordnung, die vollziehende Gewalt und die Rechtsprechung sind an Gesetz und Recht gebunden.''}}
\newcommand{\ggfnArtTwentyThreeZH}{\ggnotezhonce{art-20-3}{第20條第3項：「立法受合憲秩序拘束；行政權與司法權受法律與法拘束。」}}

% ===== Art. 26 + Art. 143（墮胎一案特有）=====
\newcommand{\ggfnArtTwoSixAndOneFourThreeDE}{\ggnotedeonce{art-26-143}{Art. 26 Abs. 1 Satz 2: ``Sie sind unter Strafe zu stellen.'' Art. 143 in der urspruenglichen Fassung (24. Mai 1949 bis 1. September 1951) enthielt Strafvorschriften zum Hochverrat; 1975 war Art. 143 bereits weggefallen.}}
\newcommand{\ggfnArtTwoSixAndOneFourThreeZH}{\ggnotezhonce{art-26-143}{第26條第1項第2句：「此等行為應規定為犯罪處罰。」第143條原始版本（1949年5月24日至1951年9月1日）含有關於叛國罪之刑罰規定；至1975年時，第143條已廢止。}}

% ===== Art. 28 =====
\newcommand{\ggfnArtTwentyEightOneDE}{\ggnotedeonce{art-28-1-1}{Art. 28 Abs. 1 Satz 1: ``Die verfassungsmaessige Ordnung in den Laendern muss den Grundsaetzen des republikanischen, demokratischen und sozialen Rechtsstaates im Sinne dieses Grundgesetzes entsprechen.''}}
\newcommand{\ggfnArtTwentyEightOneZH}{\ggnotezhonce{art-28-1-1}{第28條第1項第1句：「各邦之合憲秩序，須符合本基本法意義下共和、民主及社會法治國之原則。」}}

% ===== Art. 30 =====
\newcommand{\ggfnArtThirtyDE}{\ggnotedeonce{art-30}{Art. 30: ``Die Ausuebung der staatlichen Befugnisse und die Erfuellung der staatlichen Aufgaben ist Sache der Laender, soweit dieses Grundgesetz keine andere Regelung trifft oder zulaesst.''}}
\newcommand{\ggfnArtThirtyZH}{\ggnotezhonce{art-30}{第30條：「國家權限之行使與國家任務之履行，屬於各邦，但本基本法另有規定或容許者，不在此限。」}}

% ===== Art. 74 =====
\newcommand{\ggfnArtSeventyFourOneDE}{\ggnotedeonce{art-74-1}{Art. 74 Nr. 1: Die konkurrierende Gesetzgebung erstreckt sich auf ``das buergerliche Recht, das Strafrecht und den Strafvollzug, die Gerichtsverfassung, das gerichtliche Verfahren, die Rechtsanwaltschaft, das Notariat und die Rechtsberatung''.}}
\newcommand{\ggfnArtSeventyFourOneZH}{\ggnotezhonce{art-74-1}{第74條第1款：競合立法及於「民法、刑法及刑罰執行、法院組織、訴訟程序、律師制度、公證制度及法律諮詢」。}}
\newcommand{\ggfnArtSeventyFourSevenDE}{\ggnotedeonce{art-74-7}{Art. 74 Nr. 7: Die konkurrierende Gesetzgebung erstreckt sich auf ``die oeffentliche Fuersorge''.}}
\newcommand{\ggfnArtSeventyFourSevenZH}{\ggnotezhonce{art-74-7}{第74條第7款：競合立法及於「公共扶助」。}}
\newcommand{\ggfnArtSeventyFourTwelveDE}{\ggnotedeonce{art-74-12}{Art. 74 Nr. 12: Die konkurrierende Gesetzgebung erstreckt sich auf ``das Arbeitsrecht einschliesslich der Betriebsverfassung, des Arbeitsschutzes und der Arbeitsvermittlung sowie die Sozialversicherung einschliesslich der Arbeitslosenversicherung''.}}
\newcommand{\ggfnArtSeventyFourTwelveZH}{\ggnotezhonce{art-74-12}{第74條第12款：競合立法及於「勞動法，包括企業組織、勞動保護及就業媒合，以及社會保險，包括失業保險」。}}

% ===== Art. 77（墮胎一案特有）=====
\newcommand{\ggfnArtSevenSevenThreeDE}{\ggnotedeonce{art-77-3}{Art. 77 Abs. 3 Satz 1: ``Soweit zu einem Gesetze die Zustimmung des Bundesrates nicht erforderlich ist, kann der Bundesrat, wenn das Verfahren nach Absatz 2 beendigt ist, gegen ein vom Bundestage beschlossenes Gesetz binnen zwei Wochen Einspruch einlegen.''}}
\newcommand{\ggfnArtSevenSevenThreeZH}{\ggnotezhonce{art-77-3}{第77條第3項第1句：「法律無須聯邦參議院同意者，於第2項程序終了後，聯邦參議院得於二週內對聯邦議院通過之法律提出異議。」}}

% ===== Art. 80 =====
\newcommand{\ggfnArtEightyOneOneDE}{\ggnotedeonce{art-80-1-1}{Art. 80 Abs. 1 Satz 1: ``Durch Gesetz koennen die Bundesregierung, ein Bundesminister oder die Landesregierungen ermaechtigt werden, Rechtsverordnungen zu erlassen.''}}
\newcommand{\ggfnArtEightyOneOneZH}{\ggnotezhonce{art-80-1-1}{第80條第1項第1句：「得以法律授權聯邦政府、聯邦部長或邦政府發布法規命令。」}}

% ===== Art. 83 =====
\newcommand{\ggfnArtEightyThreeDE}{\ggnotedeonce{art-83}{Art. 83: ``Die Laender fuehren die Bundesgesetze als eigene Angelegenheit aus, soweit dieses Grundgesetz nichts anderes bestimmt oder zulaesst.''}}
\newcommand{\ggfnArtEightyThreeZH}{\ggnotezhonce{art-83}{第83條：「各邦以自己事務執行聯邦法律，但本基本法另有規定或容許者，不在此限。」}}

% ===== Art. 84 =====
\newcommand{\ggfnArtEightyFourOneDE}{\ggnotedeonce{art-84-1}{Art. 84 Abs. 1: ``Fuehren die Laender die Bundesgesetze als eigene Angelegenheit aus, so regeln sie die Einrichtung der Behoerden und das Verwaltungsverfahren, soweit nicht Bundesgesetze mit Zustimmung des Bundesrates etwas anderes bestimmen.''}}
\newcommand{\ggfnArtEightyFourOneZH}{\ggnotezhonce{art-84-1}{第84條第1項：「各邦以自己事務執行聯邦法律時，除聯邦法律經聯邦參議院同意另有規定外，由各邦規範機關之設置及行政程序。」}}
% 別名（墮胎一案使用之縮寫命名）
\let\ggfnArtEightFourOneDE\ggfnArtEightyFourOneDE
\let\ggfnArtEightFourOneZH\ggfnArtEightyFourOneZH

% ===== Art. 93 =====
\newcommand{\ggfnArtNinetyThreeOneTwoDE}{\ggnotedeonce{art-93-1-2}{Art. 93 Abs. 1 Nr. 2: Das Bundesverfassungsgericht entscheidet ``bei Meinungsverschiedenheiten oder Zweifeln ueber die foermliche und sachliche Vereinbarkeit von Bundesrecht oder Landesrecht mit diesem Grundgesetze ... auf Antrag der Bundesregierung, einer Landesregierung oder eines Drittels der Mitglieder des Bundestages''.}}
\newcommand{\ggfnArtNinetyThreeOneTwoZH}{\ggnotezhonce{art-93-1-2}{第93條第1項第2款：聯邦憲法法院就「聯邦法或邦法在形式及實質上是否與本基本法相容之意見分歧或疑義……依聯邦政府、邦政府或聯邦議院三分之一議員之聲請」作成裁判。}}
% 別名（墮胎一案使用之縮寫命名）
\let\ggfnArtNineThreeOneTwoDE\ggfnArtNinetyThreeOneTwoDE
\let\ggfnArtNineThreeOneTwoZH\ggfnArtNinetyThreeOneTwoZH

% ===== Art. 102 =====
\newcommand{\ggfnArtOneZeroTwoDE}{\ggnotedeonce{art-102}{Art. 102: ``Die Todesstrafe ist abgeschafft.''}}
\newcommand{\ggfnArtOneZeroTwoZH}{\ggnotezhonce{art-102}{第102條：「死刑廢止。」}}

% ===== Art. 103 =====
\newcommand{\ggfnArtOneZeroThreeTwoDE}{\ggnotedeonce{art-103-2}{Art. 103 Abs. 2: ``Eine Tat kann nur bestraft werden, wenn die Strafbarkeit gesetzlich bestimmt war, bevor die Tat begangen wurde.''}}
\newcommand{\ggfnArtOneZeroThreeTwoZH}{\ggnotezhonce{art-103-2}{第103條第2項：「行為之可罰性，須於行為前以法律定之，始得處罰。」}}

% ===== Art. 104 =====
\newcommand{\ggfnArtOneZeroFourOneDE}{\ggnotedeonce{art-104-1}{Art. 104 Abs. 1: ``Die Freiheit der Person kann nur auf Grund eines foermlichen Gesetzes und nur unter Beachtung der darin vorgeschriebenen Formen beschraenkt werden.''}}
\newcommand{\ggfnArtOneZeroFourOneZH}{\ggnotezhonce{art-104-1}{第104條第1項：「人身自由僅得基於形式法律，並遵守其中規定之形式加以限制。」}}

% ===== 組合腳註：Art. 3 + Art. 6（墮胎一案特有）=====
\newcommand{\ggfnArtThreeSixDE}{\ggnotedeonce{art-3-6}{Art. 3: ``Alle Menschen sind vor dem Gesetz gleich. Maenner und Frauen sind gleichberechtigt. Niemand darf wegen seines Geschlechtes, seiner Abstammung, seiner Rasse, seiner Sprache, seiner Heimat und Herkunft, seines Glaubens, seiner religioesen oder politischen Anschauungen benachteiligt oder bevorzugt werden.'' Art. 6 Abs. 1: ``Ehe und Familie stehen unter dem besonderen Schutze der staatlichen Ordnung.'' Abs. 2 Satz 1: ``Pflege und Erziehung der Kinder sind das natuerliche Recht der Eltern und die zuvoerderst ihnen obliegende Pflicht.'' Abs. 4: ``Jede Mutter hat Anspruch auf den Schutz und die Fuersorge der Gemeinschaft.''}\ggmarknotede{art-3}\ggmarknotede{art-6}\ggmarknotede{art-6-1-4}}
\newcommand{\ggfnArtThreeSixZH}{\ggnotezhonce{art-3-6}{第3條：「所有人在法律前一律平等。男女享有平等權利。任何人不得因其性別、血統、種族、語言、故鄉及出身、信仰、宗教或政治見解而受不利益或優惠。」第6條第1項：「婚姻與家庭受國家秩序之特別保護。」第2項第1句：「照護及教育子女，為父母之自然權利，並為首先由其承擔之義務。」第4項：「每一母親均有請求共同體保護與照顧之權利。」}\ggmarknotezh{art-3}\ggmarknotezh{art-6}\ggmarknotezh{art-6-1-4}}

% ===== 組合腳註：Art. 6(1) + Art. 6(4)（墮胎一案特有）=====
\newcommand{\ggfnArtSixOneFourDE}{\ggnotedeonce{art-6-1-4}{Art. 6 Abs. 1: ``Ehe und Familie stehen unter dem besonderen Schutze der staatlichen Ordnung.'' Art. 6 Abs. 4: ``Jede Mutter hat Anspruch auf den Schutz und die Fuersorge der Gemeinschaft.''}\ggmarknotede{art-6}}
\newcommand{\ggfnArtSixOneFourZH}{\ggnotezhonce{art-6-1-4}{第6條第1項：「婚姻與家庭受國家秩序之特別保護。」第4項：「每一母親均有請求共同體保護與照顧之權利。」}\ggmarknotezh{art-6}}

% ===== 組合腳註：Art. 1(1) + Art. 2(2)(1)（墮胎二案特有）=====
\newcommand{\ggfnArtOneOneTwoTwoOneDE}{\ggnotedeonce{art-1-1+2-2-1}{Art. 1 Abs. 1: ``Die Wuerde des Menschen ist unantastbar. Sie zu achten und zu schuetzen ist Verpflichtung aller staatlichen Gewalt.'' Art. 2 Abs. 2 Satz 1: ``Jeder hat das Recht auf Leben und koerperliche Unversehrtheit.''}\ggmarknotede{art-1-1}\ggmarknotede{art-2-2-1}}
\newcommand{\ggfnArtOneOneTwoTwoOneZH}{\ggnotezhonce{art-1-1+2-2-1}{第1條第1項：「人之尊嚴不可侵犯。尊重及保護之，為一切國家權力之義務。」第2條第2項第1句：「人人有生命與身體完整性之權利。」}\ggmarknotezh{art-1-1}\ggmarknotezh{art-2-2-1}}

% ===== 組合腳註：Art. 1(1) + Art. 2(2)（墮胎二案特有）=====
\newcommand{\ggfnArtOneOneTwoTwoDE}{\ggnotedeonce{art-1-1+2-2}{Art. 1 Abs. 1: ``Die Wuerde des Menschen ist unantastbar. Sie zu achten und zu schuetzen ist Verpflichtung aller staatlichen Gewalt.'' Art. 2 Abs. 2: ``Jeder hat das Recht auf Leben und koerperliche Unversehrtheit. Die Freiheit der Person ist unverletzlich. In diese Rechte darf nur auf Grund eines Gesetzes eingegriffen werden.''}\ggmarknotede{art-1-1}\ggmarknotede{art-2-2}\ggmarknotede{art-2-2-1}}
\newcommand{\ggfnArtOneOneTwoTwoZH}{\ggnotezhonce{art-1-1+2-2}{第1條第1項：「人之尊嚴不可侵犯。尊重及保護之，為一切國家權力之義務。」第2條第2項：「人人有生命與身體完整性之權利。人身自由不可侵犯。此等權利僅得依法律加以干預。」}\ggmarknotezh{art-1-1}\ggmarknotezh{art-2-2}\ggmarknotezh{art-2-2-1}}

% ===== 組合腳註：Art. 1(1) + Art. 2(1)（墮胎二案特有）=====
\newcommand{\ggfnArtOneOneTwoOneDE}{\ggnotedeonce{art-1-1+2-1}{Art. 1 Abs. 1: ``Die Wuerde des Menschen ist unantastbar. Sie zu achten und zu schuetzen ist Verpflichtung aller staatlichen Gewalt.'' Art. 2 Abs. 1: ``Jeder hat das Recht auf die freie Entfaltung seiner Persoenlichkeit, soweit er nicht die Rechte anderer verletzt und nicht gegen die verfassungsmaessige Ordnung oder das Sittengesetz verstoesst.''}\ggmarknotede{art-1-1}\ggmarknotede{art-2-1}}
\newcommand{\ggfnArtOneOneTwoOneZH}{\ggnotezhonce{art-1-1+2-1}{第1條第1項：「人之尊嚴不可侵犯。尊重及保護之，為一切國家權力之義務。」第2條第1項：「人人有自由發展其人格之權利，但不得侵害他人權利，亦不得違反合憲秩序或道德律。」}\ggmarknotezh{art-1-1}\ggmarknotezh{art-2-1}}

% ===== 組合腳註：Art. 2(2)(1) + Art. 1(1)（墮胎一案特有，含交叉標記）=====
\newcommand{\ggfnLifeDignityDE}{\ggnotedeonce{life-dignity}{Art. 2 Abs. 2 Satz 1: ``Jeder hat das Recht auf Leben und koerperliche Unversehrtheit.'' Art. 1 Abs. 1: ``Die Wuerde des Menschen ist unantastbar. Sie zu achten und zu schuetzen ist Verpflichtung aller staatlichen Gewalt.''}\ggmarknotede{art-2-2-1}\ggmarknotede{art-1-1}\ggmarknotede{art-1-1-2}}
\newcommand{\ggfnLifeDignityZH}{\ggnotezhonce{life-dignity}{第2條第2項第1句：「人人有生命與身體完整性之權利。」第1條第1項：「人之尊嚴不可侵犯。尊重及保護之，為一切國家權力之義務。」}\ggmarknotezh{art-2-2-1}\ggmarknotezh{art-1-1}\ggmarknotezh{art-1-1-2}}

% ===== 組合腳註：Art. 1(1) + Art. 2(2)（墮胎一案特有，條件判斷版）=====
\newcommand{\ggfnArtOneTwoTwoDE}{\ifcsname ggnoteseen@de@life-dignity\endcsname\ggfnArtTwoTwoRestDE{}\else\ggnotedeonce{art-1-2-2}{Art. 1 Abs. 1: ``Die Wuerde des Menschen ist unantastbar. Sie zu achten und zu schuetzen ist Verpflichtung aller staatlichen Gewalt.'' Art. 2 Abs. 2: ``Jeder hat das Recht auf Leben und koerperliche Unversehrtheit. Die Freiheit der Person ist unverletzlich. In diese Rechte darf nur auf Grund eines Gesetzes eingegriffen werden.''}\ggmarknotede{art-1-1}\ggmarknotede{art-1-1-2}\ggmarknotede{art-2-2}\ggmarknotede{art-2-2-1}\fi}
\newcommand{\ggfnArtOneTwoTwoZH}{\ifcsname ggnoteseen@zh@life-dignity\endcsname\ggfnArtTwoTwoRestZH{}\else\ggnotezhonce{art-1-2-2}{第1條第1項：「人之尊嚴不可侵犯。尊重及保護之，為一切國家權力之義務。」第2條第2項：「人人有生命與身體完整性之權利。人身自由不可侵犯。此等權利僅得依法律加以干預。」}\ggmarknotezh{art-1-1}\ggmarknotezh{art-1-1-2}\ggmarknotezh{art-2-2}\ggmarknotezh{art-2-2-1}\fi}

% ===== 組合腳註：Art. 1(1) + Art. 19(2)（墮胎一案特有，條件判斷版）=====
\newcommand{\ggfnArtOneNineteenDE}{\ifcsname ggnoteseen@de@art-1-1\endcsname\ggfnArtNineteenTwoDE{}\else\ggnotedeonce{art-1-19}{Art. 1 Abs. 1: ``Die Wuerde des Menschen ist unantastbar. Sie zu achten und zu schuetzen ist Verpflichtung aller staatlichen Gewalt.'' Art. 19 Abs. 2: ``In keinem Falle darf ein Grundrecht in seinem Wesensgehalt angetastet werden.''}\ggmarknotede{art-1-1}\ggmarknotede{art-1-1-2}\ggmarknotede{art-19-2}\fi}
\newcommand{\ggfnArtOneNineteenZH}{\ifcsname ggnoteseen@zh@art-1-1\endcsname\ggfnArtNineteenTwoZH{}\else\ggnotezhonce{art-1-19}{第1條第1項：「人之尊嚴不可侵犯。尊重及保護之，為一切國家權力之義務。」第19條第2項：「任何情形下，基本權利之本質內容均不得受到侵害。」}\ggmarknotezh{art-1-1}\ggmarknotezh{art-1-1-2}\ggmarknotezh{art-19-2}\fi}


% ===== GG 版本旗標（\ggversionde/zh 由各案在 % bverfge_gg.tex — 德國墮胎案 GG 條文腳註共用定義
% 使用方式：在各案 .tex 中先定義 \ggversionde/\ggversionzh，再 \input{../../共用LaTeX/bverfge_gg}

% ===== GG 版本旗標（\ggversionde/zh 由各案在 \input{../../共用LaTeX/bverfge_gg} 之前定義）=====
\newif\ifggversionnotede
\newif\ifggversionnotezh

% ===== 編者註腳基礎設施（首次標記模式）=====
\newif\ifeditornotelabelde
\newif\ifeditornotelabelzh
\newcommand{\editornotelabelde}{\ifeditornotelabelde\else\emph{Hrsg.-Anm. (nachfolgend ohne Kennzeichnung)}: \global\editornotelabeldetrue\fi}
\newcommand{\editornotelabelzh}{\ifeditornotelabelzh\else 編者註（下同）：\global\editornotelabelzhtrue\fi}
\newcommand{\editornotede}[1]{\ifshoweditornotes\footnote{\normalfont\footnotesize\editornotelabelde #1}\else\ifclassroommode\footnote{\normalfont\footnotesize\editornotelabelde #1}\fi\fi}
\newcommand{\editornotezh}[1]{\ifshoweditornotes\footnote{\normalfont\footnotesize\editornotelabelzh #1}\else\ifclassroommode\footnote{\normalfont\footnotesize\editornotelabelzh #1}\fi\fi}

% ===== GG 引用系統（\ggversionde/zh 由各案在 \input 前定義）=====
\newcommand{\ggmarknotede}[1]{\global\expandafter\let\csname ggnoteseen@de@#1\endcsname\empty}
\newcommand{\ggmarknotezh}[1]{\global\expandafter\let\csname ggnoteseen@zh@#1\endcsname\empty}
\newcommand{\ggnotedeonce}[2]{\ifcsname ggnoteseen@de@#1\endcsname\else\global\expandafter\let\csname ggnoteseen@de@#1\endcsname\empty\ggnotede{#2}\fi}
\newcommand{\ggnotezhonce}[2]{\ifcsname ggnoteseen@zh@#1\endcsname\else\global\expandafter\let\csname ggnoteseen@zh@#1\endcsname\empty\ggnotezh{#2}\fi}
\newcommand{\ggnotede}[1]{\editornotede{\ggversionde #1}}
\newcommand{\ggnotezh}[1]{\editornotezh{\ggversionzh #1}}

% ===== Art. 1 =====
\newcommand{\ggfnArtOneOneDE}{\ggnotedeonce{art-1-1}{Art. 1 Abs. 1: ``Die Wuerde des Menschen ist unantastbar. Sie zu achten und zu schuetzen ist Verpflichtung aller staatlichen Gewalt.''}}
\newcommand{\ggfnArtOneOneZH}{\ggnotezhonce{art-1-1}{第1條第1項：「人之尊嚴不可侵犯。尊重及保護之，為一切國家權力之義務。」}}
\newcommand{\ggfnArtOneOneTwoDE}{\ggnotedeonce{art-1-1-2}{Art. 1 Abs. 1 Satz 2: ``Sie zu achten und zu schuetzen ist Verpflichtung aller staatlichen Gewalt.''}}
\newcommand{\ggfnArtOneOneTwoZH}{\ggnotezhonce{art-1-1-2}{第1條第1項第2句：「尊重及保護之，為一切國家權力之義務。」}}

% ===== Art. 2 =====
\newcommand{\ggfnArtTwoOneDE}{\ggnotedeonce{art-2-1}{Art. 2 Abs. 1: ``Jeder hat das Recht auf die freie Entfaltung seiner Persoenlichkeit, soweit er nicht die Rechte anderer verletzt und nicht gegen die verfassungsmaessige Ordnung oder das Sittengesetz verstoesst.''}}
\newcommand{\ggfnArtTwoOneZH}{\ggnotezhonce{art-2-1}{第2條第1項：「人人有自由發展其人格之權利，但不得侵害他人權利，亦不得違反合憲秩序或道德律。」}}
% 簡易預設版（墮胎一案以 \renewcommand 覆寫為含條件邏輯之版本）
\newcommand{\ggfnArtTwoTwoDE}{\ggnotedeonce{art-2-2}{Art. 2 Abs. 2: ``Jeder hat das Recht auf Leben und koerperliche Unversehrtheit. Die Freiheit der Person ist unverletzlich. In diese Rechte darf nur auf Grund eines Gesetzes eingegriffen werden.''}\ggmarknotede{art-2-2-1}}
\newcommand{\ggfnArtTwoTwoZH}{\ggnotezhonce{art-2-2}{第2條第2項：「人人有生命與身體完整性之權利。人身自由不可侵犯。此等權利僅得依法律加以干預。」}\ggmarknotezh{art-2-2-1}}
\newcommand{\ggfnArtTwoTwoOneDE}{\ggnotedeonce{art-2-2-1}{Art. 2 Abs. 2 Satz 1: ``Jeder hat das Recht auf Leben und koerperliche Unversehrtheit.''}}
\newcommand{\ggfnArtTwoTwoOneZH}{\ggnotezhonce{art-2-2-1}{第2條第2項第1句：「人人有生命與身體完整性之權利。」}}
\newcommand{\ggfnArtTwoTwoThreeDE}{\ggnotedeonce{art-2-2-3}{Art. 2 Abs. 2 Satz 3: ``In diese Rechte darf nur auf Grund eines Gesetzes eingegriffen werden.''}}
\newcommand{\ggfnArtTwoTwoThreeZH}{\ggnotezhonce{art-2-2-3}{第2條第2項第3句：「此等權利僅得依法律加以干預。」}}
% 墮胎一案特有：第2條第2項第2–3句（Art. 2(2) 去掉第1句後的「餘文」）
\newcommand{\ggfnArtTwoTwoRestDE}{\ggnotedeonce{art-2-2-rest}{Art. 2 Abs. 2 Saetze 2-3: ``Die Freiheit der Person ist unverletzlich. In diese Rechte darf nur auf Grund eines Gesetzes eingegriffen werden.''}\ggmarknotede{art-2-2}}
\newcommand{\ggfnArtTwoTwoRestZH}{\ggnotezhonce{art-2-2-rest}{第2條第2項第2、3句：「人身自由不可侵犯。此等權利僅得依法律加以干預。」}\ggmarknotezh{art-2-2}}

% ===== Art. 3 =====
\newcommand{\ggfnArtThreeDE}{\ggnotedeonce{art-3}{Art. 3: ``Alle Menschen sind vor dem Gesetz gleich. Maenner und Frauen sind gleichberechtigt. Niemand darf wegen seines Geschlechtes, seiner Abstammung, seiner Rasse, seiner Sprache, seiner Heimat und Herkunft, seines Glaubens, seiner religioesen oder politischen Anschauungen benachteiligt oder bevorzugt werden.''}}
\newcommand{\ggfnArtThreeZH}{\ggnotezhonce{art-3}{第3條：「所有人在法律前一律平等。男女享有平等權利。任何人不得因其性別、血統、種族、語言、故鄉及出身、信仰、宗教或政治見解而受不利益或優惠。」}}
\newcommand{\ggfnArtThreeTwoDE}{\ggnotedeonce{art-3-2}{Art. 3 Abs. 2: ``Maenner und Frauen sind gleichberechtigt.''}}
\newcommand{\ggfnArtThreeTwoZH}{\ggnotezhonce{art-3-2}{第3條第2項：「男女享有平等權利。」}}

% ===== Art. 4 =====
\newcommand{\ggfnArtFourOneDE}{\ggnotedeonce{art-4-1}{Art. 4 Abs. 1: ``Die Freiheit des Glaubens, des Gewissens und die Freiheit des religioesen und weltanschaulichen Bekenntnisses sind unverletzlich.''}}
\newcommand{\ggfnArtFourOneZH}{\ggnotezhonce{art-4-1}{第4條第1項：「信仰自由、良心自由，以及宗教與世界觀信念自由，不可侵犯。」}}

% ===== Art. 5 =====
\newcommand{\ggfnArtFiveOneDE}{\ggnotedeonce{art-5-1}{Art. 5 Abs. 1: ``Jeder hat das Recht, seine Meinung in Wort, Schrift und Bild frei zu aeussern und zu verbreiten und sich aus allgemein zugaenglichen Quellen ungehindert zu unterrichten. Die Pressefreiheit und die Freiheit der Berichterstattung durch Rundfunk und Film werden gewaehrleistet. Eine Zensur findet nicht statt.''}}
\newcommand{\ggfnArtFiveOneZH}{\ggnotezhonce{art-5-1}{第5條第1項：「人人有以言語、文字及圖像自由表達及傳播其意見，並由一般可接近來源不受阻礙獲取資訊之權利。新聞自由及廣播、電影報導自由受保障。不設審查制度。」}}

% ===== Art. 6 =====
\newcommand{\ggfnArtSixDE}{\ggnotedeonce{art-6}{Art. 6 Abs. 1: ``Ehe und Familie stehen unter dem besonderen Schutze der staatlichen Ordnung.'' Abs. 2 Satz 1: ``Pflege und Erziehung der Kinder sind das natuerliche Recht der Eltern und die zuvoerderst ihnen obliegende Pflicht.'' Abs. 4: ``Jede Mutter hat Anspruch auf den Schutz und die Fuersorge der Gemeinschaft.''}}
\newcommand{\ggfnArtSixZH}{\ggnotezhonce{art-6}{第6條第1項：「婚姻與家庭受國家秩序之特別保護。」第2項第1句：「照護及教育子女，為父母之自然權利，並為首先由其承擔之義務。」第4項：「每一母親均有請求共同體保護與照顧之權利。」}}
\newcommand{\ggfnArtSixOneDE}{\ggnotedeonce{art-6-1}{Art. 6 Abs. 1: ``Ehe und Familie stehen unter dem besonderen Schutze der staatlichen Ordnung.''}}
\newcommand{\ggfnArtSixOneZH}{\ggnotezhonce{art-6-1}{第6條第1項：「婚姻與家庭受國家秩序之特別保護。」}}
\newcommand{\ggfnArtSixTwoDE}{\ggnotedeonce{art-6-2}{Art. 6 Abs. 2: ``Pflege und Erziehung der Kinder sind das natuerliche Recht der Eltern und die zuvoerderst ihnen obliegende Pflicht. Ueber ihre Betaetigung wacht die staatliche Gemeinschaft.''}}
\newcommand{\ggfnArtSixTwoZH}{\ggnotezhonce{art-6-2}{第6條第2項：「照護及教育子女，為父母之自然權利，並為首先由其承擔之義務。國家共同體監督其行使。」}}
\newcommand{\ggfnArtSixTwoOneDE}{\ggnotedeonce{art-6-2-1}{Art. 6 Abs. 2 Satz 1: ``Pflege und Erziehung der Kinder sind das natuerliche Recht der Eltern und die zuvoerderst ihnen obliegende Pflicht.''}}
\newcommand{\ggfnArtSixTwoOneZH}{\ggnotezhonce{art-6-2-1}{第6條第2項第1句：「照護及教育子女，為父母之自然權利，並為首先由其承擔之義務。」}}
\newcommand{\ggfnArtSixFourDE}{\ggnotedeonce{art-6-4}{Art. 6 Abs. 4: ``Jede Mutter hat Anspruch auf den Schutz und die Fuersorge der Gemeinschaft.''}}
\newcommand{\ggfnArtSixFourZH}{\ggnotezhonce{art-6-4}{第6條第4項：「每一母親均有請求共同體保護與照顧之權利。」}}

% ===== Art. 12 =====
\newcommand{\ggfnArtTwelveOneDE}{\ggnotedeonce{art-12-1}{Art. 12 Abs. 1: ``Alle Deutschen haben das Recht, Beruf, Arbeitsplatz und Ausbildungsstaette frei zu waehlen. Die Berufsausuebung kann durch Gesetz oder auf Grund eines Gesetzes geregelt werden.''}}
\newcommand{\ggfnArtTwelveOneZH}{\ggnotezhonce{art-12-1}{第12條第1項：「所有德國人均有自由選擇職業、工作場所及訓練場所之權利。職業行使得以法律或基於法律加以規範。」}}

% ===== Art. 19 =====
\newcommand{\ggfnArtNineteenTwoDE}{\ggnotedeonce{art-19-2}{Art. 19 Abs. 2: ``In keinem Falle darf ein Grundrecht in seinem Wesensgehalt angetastet werden.''}}
\newcommand{\ggfnArtNineteenTwoZH}{\ggnotezhonce{art-19-2}{第19條第2項：「任何情形下，基本權利之本質內容均不得受到侵害。」}}

% ===== Art. 20 =====
\newcommand{\ggfnArtTwentyOneDE}{\ggnotedeonce{art-20-1}{Art. 20 Abs. 1: ``Die Bundesrepublik Deutschland ist ein demokratischer und sozialer Bundesstaat.''}}
\newcommand{\ggfnArtTwentyOneZH}{\ggnotezhonce{art-20-1}{第20條第1項：「德意志聯邦共和國為民主及社會之聯邦國。」}}
\newcommand{\ggfnArtTwentyThreeDE}{\ggnotedeonce{art-20-3}{Art. 20 Abs. 3: ``Die Gesetzgebung ist an die verfassungsmaessige Ordnung, die vollziehende Gewalt und die Rechtsprechung sind an Gesetz und Recht gebunden.''}}
\newcommand{\ggfnArtTwentyThreeZH}{\ggnotezhonce{art-20-3}{第20條第3項：「立法受合憲秩序拘束；行政權與司法權受法律與法拘束。」}}

% ===== Art. 26 + Art. 143（墮胎一案特有）=====
\newcommand{\ggfnArtTwoSixAndOneFourThreeDE}{\ggnotedeonce{art-26-143}{Art. 26 Abs. 1 Satz 2: ``Sie sind unter Strafe zu stellen.'' Art. 143 in der urspruenglichen Fassung (24. Mai 1949 bis 1. September 1951) enthielt Strafvorschriften zum Hochverrat; 1975 war Art. 143 bereits weggefallen.}}
\newcommand{\ggfnArtTwoSixAndOneFourThreeZH}{\ggnotezhonce{art-26-143}{第26條第1項第2句：「此等行為應規定為犯罪處罰。」第143條原始版本（1949年5月24日至1951年9月1日）含有關於叛國罪之刑罰規定；至1975年時，第143條已廢止。}}

% ===== Art. 28 =====
\newcommand{\ggfnArtTwentyEightOneDE}{\ggnotedeonce{art-28-1-1}{Art. 28 Abs. 1 Satz 1: ``Die verfassungsmaessige Ordnung in den Laendern muss den Grundsaetzen des republikanischen, demokratischen und sozialen Rechtsstaates im Sinne dieses Grundgesetzes entsprechen.''}}
\newcommand{\ggfnArtTwentyEightOneZH}{\ggnotezhonce{art-28-1-1}{第28條第1項第1句：「各邦之合憲秩序，須符合本基本法意義下共和、民主及社會法治國之原則。」}}

% ===== Art. 30 =====
\newcommand{\ggfnArtThirtyDE}{\ggnotedeonce{art-30}{Art. 30: ``Die Ausuebung der staatlichen Befugnisse und die Erfuellung der staatlichen Aufgaben ist Sache der Laender, soweit dieses Grundgesetz keine andere Regelung trifft oder zulaesst.''}}
\newcommand{\ggfnArtThirtyZH}{\ggnotezhonce{art-30}{第30條：「國家權限之行使與國家任務之履行，屬於各邦，但本基本法另有規定或容許者，不在此限。」}}

% ===== Art. 74 =====
\newcommand{\ggfnArtSeventyFourOneDE}{\ggnotedeonce{art-74-1}{Art. 74 Nr. 1: Die konkurrierende Gesetzgebung erstreckt sich auf ``das buergerliche Recht, das Strafrecht und den Strafvollzug, die Gerichtsverfassung, das gerichtliche Verfahren, die Rechtsanwaltschaft, das Notariat und die Rechtsberatung''.}}
\newcommand{\ggfnArtSeventyFourOneZH}{\ggnotezhonce{art-74-1}{第74條第1款：競合立法及於「民法、刑法及刑罰執行、法院組織、訴訟程序、律師制度、公證制度及法律諮詢」。}}
\newcommand{\ggfnArtSeventyFourSevenDE}{\ggnotedeonce{art-74-7}{Art. 74 Nr. 7: Die konkurrierende Gesetzgebung erstreckt sich auf ``die oeffentliche Fuersorge''.}}
\newcommand{\ggfnArtSeventyFourSevenZH}{\ggnotezhonce{art-74-7}{第74條第7款：競合立法及於「公共扶助」。}}
\newcommand{\ggfnArtSeventyFourTwelveDE}{\ggnotedeonce{art-74-12}{Art. 74 Nr. 12: Die konkurrierende Gesetzgebung erstreckt sich auf ``das Arbeitsrecht einschliesslich der Betriebsverfassung, des Arbeitsschutzes und der Arbeitsvermittlung sowie die Sozialversicherung einschliesslich der Arbeitslosenversicherung''.}}
\newcommand{\ggfnArtSeventyFourTwelveZH}{\ggnotezhonce{art-74-12}{第74條第12款：競合立法及於「勞動法，包括企業組織、勞動保護及就業媒合，以及社會保險，包括失業保險」。}}

% ===== Art. 77（墮胎一案特有）=====
\newcommand{\ggfnArtSevenSevenThreeDE}{\ggnotedeonce{art-77-3}{Art. 77 Abs. 3 Satz 1: ``Soweit zu einem Gesetze die Zustimmung des Bundesrates nicht erforderlich ist, kann der Bundesrat, wenn das Verfahren nach Absatz 2 beendigt ist, gegen ein vom Bundestage beschlossenes Gesetz binnen zwei Wochen Einspruch einlegen.''}}
\newcommand{\ggfnArtSevenSevenThreeZH}{\ggnotezhonce{art-77-3}{第77條第3項第1句：「法律無須聯邦參議院同意者，於第2項程序終了後，聯邦參議院得於二週內對聯邦議院通過之法律提出異議。」}}

% ===== Art. 80 =====
\newcommand{\ggfnArtEightyOneOneDE}{\ggnotedeonce{art-80-1-1}{Art. 80 Abs. 1 Satz 1: ``Durch Gesetz koennen die Bundesregierung, ein Bundesminister oder die Landesregierungen ermaechtigt werden, Rechtsverordnungen zu erlassen.''}}
\newcommand{\ggfnArtEightyOneOneZH}{\ggnotezhonce{art-80-1-1}{第80條第1項第1句：「得以法律授權聯邦政府、聯邦部長或邦政府發布法規命令。」}}

% ===== Art. 83 =====
\newcommand{\ggfnArtEightyThreeDE}{\ggnotedeonce{art-83}{Art. 83: ``Die Laender fuehren die Bundesgesetze als eigene Angelegenheit aus, soweit dieses Grundgesetz nichts anderes bestimmt oder zulaesst.''}}
\newcommand{\ggfnArtEightyThreeZH}{\ggnotezhonce{art-83}{第83條：「各邦以自己事務執行聯邦法律，但本基本法另有規定或容許者，不在此限。」}}

% ===== Art. 84 =====
\newcommand{\ggfnArtEightyFourOneDE}{\ggnotedeonce{art-84-1}{Art. 84 Abs. 1: ``Fuehren die Laender die Bundesgesetze als eigene Angelegenheit aus, so regeln sie die Einrichtung der Behoerden und das Verwaltungsverfahren, soweit nicht Bundesgesetze mit Zustimmung des Bundesrates etwas anderes bestimmen.''}}
\newcommand{\ggfnArtEightyFourOneZH}{\ggnotezhonce{art-84-1}{第84條第1項：「各邦以自己事務執行聯邦法律時，除聯邦法律經聯邦參議院同意另有規定外，由各邦規範機關之設置及行政程序。」}}
% 別名（墮胎一案使用之縮寫命名）
\let\ggfnArtEightFourOneDE\ggfnArtEightyFourOneDE
\let\ggfnArtEightFourOneZH\ggfnArtEightyFourOneZH

% ===== Art. 93 =====
\newcommand{\ggfnArtNinetyThreeOneTwoDE}{\ggnotedeonce{art-93-1-2}{Art. 93 Abs. 1 Nr. 2: Das Bundesverfassungsgericht entscheidet ``bei Meinungsverschiedenheiten oder Zweifeln ueber die foermliche und sachliche Vereinbarkeit von Bundesrecht oder Landesrecht mit diesem Grundgesetze ... auf Antrag der Bundesregierung, einer Landesregierung oder eines Drittels der Mitglieder des Bundestages''.}}
\newcommand{\ggfnArtNinetyThreeOneTwoZH}{\ggnotezhonce{art-93-1-2}{第93條第1項第2款：聯邦憲法法院就「聯邦法或邦法在形式及實質上是否與本基本法相容之意見分歧或疑義……依聯邦政府、邦政府或聯邦議院三分之一議員之聲請」作成裁判。}}
% 別名（墮胎一案使用之縮寫命名）
\let\ggfnArtNineThreeOneTwoDE\ggfnArtNinetyThreeOneTwoDE
\let\ggfnArtNineThreeOneTwoZH\ggfnArtNinetyThreeOneTwoZH

% ===== Art. 102 =====
\newcommand{\ggfnArtOneZeroTwoDE}{\ggnotedeonce{art-102}{Art. 102: ``Die Todesstrafe ist abgeschafft.''}}
\newcommand{\ggfnArtOneZeroTwoZH}{\ggnotezhonce{art-102}{第102條：「死刑廢止。」}}

% ===== Art. 103 =====
\newcommand{\ggfnArtOneZeroThreeTwoDE}{\ggnotedeonce{art-103-2}{Art. 103 Abs. 2: ``Eine Tat kann nur bestraft werden, wenn die Strafbarkeit gesetzlich bestimmt war, bevor die Tat begangen wurde.''}}
\newcommand{\ggfnArtOneZeroThreeTwoZH}{\ggnotezhonce{art-103-2}{第103條第2項：「行為之可罰性，須於行為前以法律定之，始得處罰。」}}

% ===== Art. 104 =====
\newcommand{\ggfnArtOneZeroFourOneDE}{\ggnotedeonce{art-104-1}{Art. 104 Abs. 1: ``Die Freiheit der Person kann nur auf Grund eines foermlichen Gesetzes und nur unter Beachtung der darin vorgeschriebenen Formen beschraenkt werden.''}}
\newcommand{\ggfnArtOneZeroFourOneZH}{\ggnotezhonce{art-104-1}{第104條第1項：「人身自由僅得基於形式法律，並遵守其中規定之形式加以限制。」}}

% ===== 組合腳註：Art. 3 + Art. 6（墮胎一案特有）=====
\newcommand{\ggfnArtThreeSixDE}{\ggnotedeonce{art-3-6}{Art. 3: ``Alle Menschen sind vor dem Gesetz gleich. Maenner und Frauen sind gleichberechtigt. Niemand darf wegen seines Geschlechtes, seiner Abstammung, seiner Rasse, seiner Sprache, seiner Heimat und Herkunft, seines Glaubens, seiner religioesen oder politischen Anschauungen benachteiligt oder bevorzugt werden.'' Art. 6 Abs. 1: ``Ehe und Familie stehen unter dem besonderen Schutze der staatlichen Ordnung.'' Abs. 2 Satz 1: ``Pflege und Erziehung der Kinder sind das natuerliche Recht der Eltern und die zuvoerderst ihnen obliegende Pflicht.'' Abs. 4: ``Jede Mutter hat Anspruch auf den Schutz und die Fuersorge der Gemeinschaft.''}\ggmarknotede{art-3}\ggmarknotede{art-6}\ggmarknotede{art-6-1-4}}
\newcommand{\ggfnArtThreeSixZH}{\ggnotezhonce{art-3-6}{第3條：「所有人在法律前一律平等。男女享有平等權利。任何人不得因其性別、血統、種族、語言、故鄉及出身、信仰、宗教或政治見解而受不利益或優惠。」第6條第1項：「婚姻與家庭受國家秩序之特別保護。」第2項第1句：「照護及教育子女，為父母之自然權利，並為首先由其承擔之義務。」第4項：「每一母親均有請求共同體保護與照顧之權利。」}\ggmarknotezh{art-3}\ggmarknotezh{art-6}\ggmarknotezh{art-6-1-4}}

% ===== 組合腳註：Art. 6(1) + Art. 6(4)（墮胎一案特有）=====
\newcommand{\ggfnArtSixOneFourDE}{\ggnotedeonce{art-6-1-4}{Art. 6 Abs. 1: ``Ehe und Familie stehen unter dem besonderen Schutze der staatlichen Ordnung.'' Art. 6 Abs. 4: ``Jede Mutter hat Anspruch auf den Schutz und die Fuersorge der Gemeinschaft.''}\ggmarknotede{art-6}}
\newcommand{\ggfnArtSixOneFourZH}{\ggnotezhonce{art-6-1-4}{第6條第1項：「婚姻與家庭受國家秩序之特別保護。」第4項：「每一母親均有請求共同體保護與照顧之權利。」}\ggmarknotezh{art-6}}

% ===== 組合腳註：Art. 1(1) + Art. 2(2)(1)（墮胎二案特有）=====
\newcommand{\ggfnArtOneOneTwoTwoOneDE}{\ggnotedeonce{art-1-1+2-2-1}{Art. 1 Abs. 1: ``Die Wuerde des Menschen ist unantastbar. Sie zu achten und zu schuetzen ist Verpflichtung aller staatlichen Gewalt.'' Art. 2 Abs. 2 Satz 1: ``Jeder hat das Recht auf Leben und koerperliche Unversehrtheit.''}\ggmarknotede{art-1-1}\ggmarknotede{art-2-2-1}}
\newcommand{\ggfnArtOneOneTwoTwoOneZH}{\ggnotezhonce{art-1-1+2-2-1}{第1條第1項：「人之尊嚴不可侵犯。尊重及保護之，為一切國家權力之義務。」第2條第2項第1句：「人人有生命與身體完整性之權利。」}\ggmarknotezh{art-1-1}\ggmarknotezh{art-2-2-1}}

% ===== 組合腳註：Art. 1(1) + Art. 2(2)（墮胎二案特有）=====
\newcommand{\ggfnArtOneOneTwoTwoDE}{\ggnotedeonce{art-1-1+2-2}{Art. 1 Abs. 1: ``Die Wuerde des Menschen ist unantastbar. Sie zu achten und zu schuetzen ist Verpflichtung aller staatlichen Gewalt.'' Art. 2 Abs. 2: ``Jeder hat das Recht auf Leben und koerperliche Unversehrtheit. Die Freiheit der Person ist unverletzlich. In diese Rechte darf nur auf Grund eines Gesetzes eingegriffen werden.''}\ggmarknotede{art-1-1}\ggmarknotede{art-2-2}\ggmarknotede{art-2-2-1}}
\newcommand{\ggfnArtOneOneTwoTwoZH}{\ggnotezhonce{art-1-1+2-2}{第1條第1項：「人之尊嚴不可侵犯。尊重及保護之，為一切國家權力之義務。」第2條第2項：「人人有生命與身體完整性之權利。人身自由不可侵犯。此等權利僅得依法律加以干預。」}\ggmarknotezh{art-1-1}\ggmarknotezh{art-2-2}\ggmarknotezh{art-2-2-1}}

% ===== 組合腳註：Art. 1(1) + Art. 2(1)（墮胎二案特有）=====
\newcommand{\ggfnArtOneOneTwoOneDE}{\ggnotedeonce{art-1-1+2-1}{Art. 1 Abs. 1: ``Die Wuerde des Menschen ist unantastbar. Sie zu achten und zu schuetzen ist Verpflichtung aller staatlichen Gewalt.'' Art. 2 Abs. 1: ``Jeder hat das Recht auf die freie Entfaltung seiner Persoenlichkeit, soweit er nicht die Rechte anderer verletzt und nicht gegen die verfassungsmaessige Ordnung oder das Sittengesetz verstoesst.''}\ggmarknotede{art-1-1}\ggmarknotede{art-2-1}}
\newcommand{\ggfnArtOneOneTwoOneZH}{\ggnotezhonce{art-1-1+2-1}{第1條第1項：「人之尊嚴不可侵犯。尊重及保護之，為一切國家權力之義務。」第2條第1項：「人人有自由發展其人格之權利，但不得侵害他人權利，亦不得違反合憲秩序或道德律。」}\ggmarknotezh{art-1-1}\ggmarknotezh{art-2-1}}

% ===== 組合腳註：Art. 2(2)(1) + Art. 1(1)（墮胎一案特有，含交叉標記）=====
\newcommand{\ggfnLifeDignityDE}{\ggnotedeonce{life-dignity}{Art. 2 Abs. 2 Satz 1: ``Jeder hat das Recht auf Leben und koerperliche Unversehrtheit.'' Art. 1 Abs. 1: ``Die Wuerde des Menschen ist unantastbar. Sie zu achten und zu schuetzen ist Verpflichtung aller staatlichen Gewalt.''}\ggmarknotede{art-2-2-1}\ggmarknotede{art-1-1}\ggmarknotede{art-1-1-2}}
\newcommand{\ggfnLifeDignityZH}{\ggnotezhonce{life-dignity}{第2條第2項第1句：「人人有生命與身體完整性之權利。」第1條第1項：「人之尊嚴不可侵犯。尊重及保護之，為一切國家權力之義務。」}\ggmarknotezh{art-2-2-1}\ggmarknotezh{art-1-1}\ggmarknotezh{art-1-1-2}}

% ===== 組合腳註：Art. 1(1) + Art. 2(2)（墮胎一案特有，條件判斷版）=====
\newcommand{\ggfnArtOneTwoTwoDE}{\ifcsname ggnoteseen@de@life-dignity\endcsname\ggfnArtTwoTwoRestDE{}\else\ggnotedeonce{art-1-2-2}{Art. 1 Abs. 1: ``Die Wuerde des Menschen ist unantastbar. Sie zu achten und zu schuetzen ist Verpflichtung aller staatlichen Gewalt.'' Art. 2 Abs. 2: ``Jeder hat das Recht auf Leben und koerperliche Unversehrtheit. Die Freiheit der Person ist unverletzlich. In diese Rechte darf nur auf Grund eines Gesetzes eingegriffen werden.''}\ggmarknotede{art-1-1}\ggmarknotede{art-1-1-2}\ggmarknotede{art-2-2}\ggmarknotede{art-2-2-1}\fi}
\newcommand{\ggfnArtOneTwoTwoZH}{\ifcsname ggnoteseen@zh@life-dignity\endcsname\ggfnArtTwoTwoRestZH{}\else\ggnotezhonce{art-1-2-2}{第1條第1項：「人之尊嚴不可侵犯。尊重及保護之，為一切國家權力之義務。」第2條第2項：「人人有生命與身體完整性之權利。人身自由不可侵犯。此等權利僅得依法律加以干預。」}\ggmarknotezh{art-1-1}\ggmarknotezh{art-1-1-2}\ggmarknotezh{art-2-2}\ggmarknotezh{art-2-2-1}\fi}

% ===== 組合腳註：Art. 1(1) + Art. 19(2)（墮胎一案特有，條件判斷版）=====
\newcommand{\ggfnArtOneNineteenDE}{\ifcsname ggnoteseen@de@art-1-1\endcsname\ggfnArtNineteenTwoDE{}\else\ggnotedeonce{art-1-19}{Art. 1 Abs. 1: ``Die Wuerde des Menschen ist unantastbar. Sie zu achten und zu schuetzen ist Verpflichtung aller staatlichen Gewalt.'' Art. 19 Abs. 2: ``In keinem Falle darf ein Grundrecht in seinem Wesensgehalt angetastet werden.''}\ggmarknotede{art-1-1}\ggmarknotede{art-1-1-2}\ggmarknotede{art-19-2}\fi}
\newcommand{\ggfnArtOneNineteenZH}{\ifcsname ggnoteseen@zh@art-1-1\endcsname\ggfnArtNineteenTwoZH{}\else\ggnotezhonce{art-1-19}{第1條第1項：「人之尊嚴不可侵犯。尊重及保護之，為一切國家權力之義務。」第19條第2項：「任何情形下，基本權利之本質內容均不得受到侵害。」}\ggmarknotezh{art-1-1}\ggmarknotezh{art-1-1-2}\ggmarknotezh{art-19-2}\fi}
 之前定義）=====
\newif\ifggversionnotede
\newif\ifggversionnotezh

% ===== 編者註腳基礎設施（首次標記模式）=====
\newif\ifeditornotelabelde
\newif\ifeditornotelabelzh
\newcommand{\editornotelabelde}{\ifeditornotelabelde\else\emph{Hrsg.-Anm. (nachfolgend ohne Kennzeichnung)}: \global\editornotelabeldetrue\fi}
\newcommand{\editornotelabelzh}{\ifeditornotelabelzh\else 編者註（下同）：\global\editornotelabelzhtrue\fi}
\newcommand{\editornotede}[1]{\ifshoweditornotes\footnote{\normalfont\footnotesize\editornotelabelde #1}\else\ifclassroommode\footnote{\normalfont\footnotesize\editornotelabelde #1}\fi\fi}
\newcommand{\editornotezh}[1]{\ifshoweditornotes\footnote{\normalfont\footnotesize\editornotelabelzh #1}\else\ifclassroommode\footnote{\normalfont\footnotesize\editornotelabelzh #1}\fi\fi}

% ===== GG 引用系統（\ggversionde/zh 由各案在 \input 前定義）=====
\newcommand{\ggmarknotede}[1]{\global\expandafter\let\csname ggnoteseen@de@#1\endcsname\empty}
\newcommand{\ggmarknotezh}[1]{\global\expandafter\let\csname ggnoteseen@zh@#1\endcsname\empty}
\newcommand{\ggnotedeonce}[2]{\ifcsname ggnoteseen@de@#1\endcsname\else\global\expandafter\let\csname ggnoteseen@de@#1\endcsname\empty\ggnotede{#2}\fi}
\newcommand{\ggnotezhonce}[2]{\ifcsname ggnoteseen@zh@#1\endcsname\else\global\expandafter\let\csname ggnoteseen@zh@#1\endcsname\empty\ggnotezh{#2}\fi}
\newcommand{\ggnotede}[1]{\editornotede{\ggversionde #1}}
\newcommand{\ggnotezh}[1]{\editornotezh{\ggversionzh #1}}

% ===== Art. 1 =====
\newcommand{\ggfnArtOneOneDE}{\ggnotedeonce{art-1-1}{Art. 1 Abs. 1: ``Die Wuerde des Menschen ist unantastbar. Sie zu achten und zu schuetzen ist Verpflichtung aller staatlichen Gewalt.''}}
\newcommand{\ggfnArtOneOneZH}{\ggnotezhonce{art-1-1}{第1條第1項：「人之尊嚴不可侵犯。尊重及保護之，為一切國家權力之義務。」}}
\newcommand{\ggfnArtOneOneTwoDE}{\ggnotedeonce{art-1-1-2}{Art. 1 Abs. 1 Satz 2: ``Sie zu achten und zu schuetzen ist Verpflichtung aller staatlichen Gewalt.''}}
\newcommand{\ggfnArtOneOneTwoZH}{\ggnotezhonce{art-1-1-2}{第1條第1項第2句：「尊重及保護之，為一切國家權力之義務。」}}

% ===== Art. 2 =====
\newcommand{\ggfnArtTwoOneDE}{\ggnotedeonce{art-2-1}{Art. 2 Abs. 1: ``Jeder hat das Recht auf die freie Entfaltung seiner Persoenlichkeit, soweit er nicht die Rechte anderer verletzt und nicht gegen die verfassungsmaessige Ordnung oder das Sittengesetz verstoesst.''}}
\newcommand{\ggfnArtTwoOneZH}{\ggnotezhonce{art-2-1}{第2條第1項：「人人有自由發展其人格之權利，但不得侵害他人權利，亦不得違反合憲秩序或道德律。」}}
% 簡易預設版（墮胎一案以 \renewcommand 覆寫為含條件邏輯之版本）
\newcommand{\ggfnArtTwoTwoDE}{\ggnotedeonce{art-2-2}{Art. 2 Abs. 2: ``Jeder hat das Recht auf Leben und koerperliche Unversehrtheit. Die Freiheit der Person ist unverletzlich. In diese Rechte darf nur auf Grund eines Gesetzes eingegriffen werden.''}\ggmarknotede{art-2-2-1}}
\newcommand{\ggfnArtTwoTwoZH}{\ggnotezhonce{art-2-2}{第2條第2項：「人人有生命與身體完整性之權利。人身自由不可侵犯。此等權利僅得依法律加以干預。」}\ggmarknotezh{art-2-2-1}}
\newcommand{\ggfnArtTwoTwoOneDE}{\ggnotedeonce{art-2-2-1}{Art. 2 Abs. 2 Satz 1: ``Jeder hat das Recht auf Leben und koerperliche Unversehrtheit.''}}
\newcommand{\ggfnArtTwoTwoOneZH}{\ggnotezhonce{art-2-2-1}{第2條第2項第1句：「人人有生命與身體完整性之權利。」}}
\newcommand{\ggfnArtTwoTwoThreeDE}{\ggnotedeonce{art-2-2-3}{Art. 2 Abs. 2 Satz 3: ``In diese Rechte darf nur auf Grund eines Gesetzes eingegriffen werden.''}}
\newcommand{\ggfnArtTwoTwoThreeZH}{\ggnotezhonce{art-2-2-3}{第2條第2項第3句：「此等權利僅得依法律加以干預。」}}
% 墮胎一案特有：第2條第2項第2–3句（Art. 2(2) 去掉第1句後的「餘文」）
\newcommand{\ggfnArtTwoTwoRestDE}{\ggnotedeonce{art-2-2-rest}{Art. 2 Abs. 2 Saetze 2-3: ``Die Freiheit der Person ist unverletzlich. In diese Rechte darf nur auf Grund eines Gesetzes eingegriffen werden.''}\ggmarknotede{art-2-2}}
\newcommand{\ggfnArtTwoTwoRestZH}{\ggnotezhonce{art-2-2-rest}{第2條第2項第2、3句：「人身自由不可侵犯。此等權利僅得依法律加以干預。」}\ggmarknotezh{art-2-2}}

% ===== Art. 3 =====
\newcommand{\ggfnArtThreeDE}{\ggnotedeonce{art-3}{Art. 3: ``Alle Menschen sind vor dem Gesetz gleich. Maenner und Frauen sind gleichberechtigt. Niemand darf wegen seines Geschlechtes, seiner Abstammung, seiner Rasse, seiner Sprache, seiner Heimat und Herkunft, seines Glaubens, seiner religioesen oder politischen Anschauungen benachteiligt oder bevorzugt werden.''}}
\newcommand{\ggfnArtThreeZH}{\ggnotezhonce{art-3}{第3條：「所有人在法律前一律平等。男女享有平等權利。任何人不得因其性別、血統、種族、語言、故鄉及出身、信仰、宗教或政治見解而受不利益或優惠。」}}
\newcommand{\ggfnArtThreeTwoDE}{\ggnotedeonce{art-3-2}{Art. 3 Abs. 2: ``Maenner und Frauen sind gleichberechtigt.''}}
\newcommand{\ggfnArtThreeTwoZH}{\ggnotezhonce{art-3-2}{第3條第2項：「男女享有平等權利。」}}

% ===== Art. 4 =====
\newcommand{\ggfnArtFourOneDE}{\ggnotedeonce{art-4-1}{Art. 4 Abs. 1: ``Die Freiheit des Glaubens, des Gewissens und die Freiheit des religioesen und weltanschaulichen Bekenntnisses sind unverletzlich.''}}
\newcommand{\ggfnArtFourOneZH}{\ggnotezhonce{art-4-1}{第4條第1項：「信仰自由、良心自由，以及宗教與世界觀信念自由，不可侵犯。」}}

% ===== Art. 5 =====
\newcommand{\ggfnArtFiveOneDE}{\ggnotedeonce{art-5-1}{Art. 5 Abs. 1: ``Jeder hat das Recht, seine Meinung in Wort, Schrift und Bild frei zu aeussern und zu verbreiten und sich aus allgemein zugaenglichen Quellen ungehindert zu unterrichten. Die Pressefreiheit und die Freiheit der Berichterstattung durch Rundfunk und Film werden gewaehrleistet. Eine Zensur findet nicht statt.''}}
\newcommand{\ggfnArtFiveOneZH}{\ggnotezhonce{art-5-1}{第5條第1項：「人人有以言語、文字及圖像自由表達及傳播其意見，並由一般可接近來源不受阻礙獲取資訊之權利。新聞自由及廣播、電影報導自由受保障。不設審查制度。」}}

% ===== Art. 6 =====
\newcommand{\ggfnArtSixDE}{\ggnotedeonce{art-6}{Art. 6 Abs. 1: ``Ehe und Familie stehen unter dem besonderen Schutze der staatlichen Ordnung.'' Abs. 2 Satz 1: ``Pflege und Erziehung der Kinder sind das natuerliche Recht der Eltern und die zuvoerderst ihnen obliegende Pflicht.'' Abs. 4: ``Jede Mutter hat Anspruch auf den Schutz und die Fuersorge der Gemeinschaft.''}}
\newcommand{\ggfnArtSixZH}{\ggnotezhonce{art-6}{第6條第1項：「婚姻與家庭受國家秩序之特別保護。」第2項第1句：「照護及教育子女，為父母之自然權利，並為首先由其承擔之義務。」第4項：「每一母親均有請求共同體保護與照顧之權利。」}}
\newcommand{\ggfnArtSixOneDE}{\ggnotedeonce{art-6-1}{Art. 6 Abs. 1: ``Ehe und Familie stehen unter dem besonderen Schutze der staatlichen Ordnung.''}}
\newcommand{\ggfnArtSixOneZH}{\ggnotezhonce{art-6-1}{第6條第1項：「婚姻與家庭受國家秩序之特別保護。」}}
\newcommand{\ggfnArtSixTwoDE}{\ggnotedeonce{art-6-2}{Art. 6 Abs. 2: ``Pflege und Erziehung der Kinder sind das natuerliche Recht der Eltern und die zuvoerderst ihnen obliegende Pflicht. Ueber ihre Betaetigung wacht die staatliche Gemeinschaft.''}}
\newcommand{\ggfnArtSixTwoZH}{\ggnotezhonce{art-6-2}{第6條第2項：「照護及教育子女，為父母之自然權利，並為首先由其承擔之義務。國家共同體監督其行使。」}}
\newcommand{\ggfnArtSixTwoOneDE}{\ggnotedeonce{art-6-2-1}{Art. 6 Abs. 2 Satz 1: ``Pflege und Erziehung der Kinder sind das natuerliche Recht der Eltern und die zuvoerderst ihnen obliegende Pflicht.''}}
\newcommand{\ggfnArtSixTwoOneZH}{\ggnotezhonce{art-6-2-1}{第6條第2項第1句：「照護及教育子女，為父母之自然權利，並為首先由其承擔之義務。」}}
\newcommand{\ggfnArtSixFourDE}{\ggnotedeonce{art-6-4}{Art. 6 Abs. 4: ``Jede Mutter hat Anspruch auf den Schutz und die Fuersorge der Gemeinschaft.''}}
\newcommand{\ggfnArtSixFourZH}{\ggnotezhonce{art-6-4}{第6條第4項：「每一母親均有請求共同體保護與照顧之權利。」}}

% ===== Art. 12 =====
\newcommand{\ggfnArtTwelveOneDE}{\ggnotedeonce{art-12-1}{Art. 12 Abs. 1: ``Alle Deutschen haben das Recht, Beruf, Arbeitsplatz und Ausbildungsstaette frei zu waehlen. Die Berufsausuebung kann durch Gesetz oder auf Grund eines Gesetzes geregelt werden.''}}
\newcommand{\ggfnArtTwelveOneZH}{\ggnotezhonce{art-12-1}{第12條第1項：「所有德國人均有自由選擇職業、工作場所及訓練場所之權利。職業行使得以法律或基於法律加以規範。」}}

% ===== Art. 19 =====
\newcommand{\ggfnArtNineteenTwoDE}{\ggnotedeonce{art-19-2}{Art. 19 Abs. 2: ``In keinem Falle darf ein Grundrecht in seinem Wesensgehalt angetastet werden.''}}
\newcommand{\ggfnArtNineteenTwoZH}{\ggnotezhonce{art-19-2}{第19條第2項：「任何情形下，基本權利之本質內容均不得受到侵害。」}}

% ===== Art. 20 =====
\newcommand{\ggfnArtTwentyOneDE}{\ggnotedeonce{art-20-1}{Art. 20 Abs. 1: ``Die Bundesrepublik Deutschland ist ein demokratischer und sozialer Bundesstaat.''}}
\newcommand{\ggfnArtTwentyOneZH}{\ggnotezhonce{art-20-1}{第20條第1項：「德意志聯邦共和國為民主及社會之聯邦國。」}}
\newcommand{\ggfnArtTwentyThreeDE}{\ggnotedeonce{art-20-3}{Art. 20 Abs. 3: ``Die Gesetzgebung ist an die verfassungsmaessige Ordnung, die vollziehende Gewalt und die Rechtsprechung sind an Gesetz und Recht gebunden.''}}
\newcommand{\ggfnArtTwentyThreeZH}{\ggnotezhonce{art-20-3}{第20條第3項：「立法受合憲秩序拘束；行政權與司法權受法律與法拘束。」}}

% ===== Art. 26 + Art. 143（墮胎一案特有）=====
\newcommand{\ggfnArtTwoSixAndOneFourThreeDE}{\ggnotedeonce{art-26-143}{Art. 26 Abs. 1 Satz 2: ``Sie sind unter Strafe zu stellen.'' Art. 143 in der urspruenglichen Fassung (24. Mai 1949 bis 1. September 1951) enthielt Strafvorschriften zum Hochverrat; 1975 war Art. 143 bereits weggefallen.}}
\newcommand{\ggfnArtTwoSixAndOneFourThreeZH}{\ggnotezhonce{art-26-143}{第26條第1項第2句：「此等行為應規定為犯罪處罰。」第143條原始版本（1949年5月24日至1951年9月1日）含有關於叛國罪之刑罰規定；至1975年時，第143條已廢止。}}

% ===== Art. 28 =====
\newcommand{\ggfnArtTwentyEightOneDE}{\ggnotedeonce{art-28-1-1}{Art. 28 Abs. 1 Satz 1: ``Die verfassungsmaessige Ordnung in den Laendern muss den Grundsaetzen des republikanischen, demokratischen und sozialen Rechtsstaates im Sinne dieses Grundgesetzes entsprechen.''}}
\newcommand{\ggfnArtTwentyEightOneZH}{\ggnotezhonce{art-28-1-1}{第28條第1項第1句：「各邦之合憲秩序，須符合本基本法意義下共和、民主及社會法治國之原則。」}}

% ===== Art. 30 =====
\newcommand{\ggfnArtThirtyDE}{\ggnotedeonce{art-30}{Art. 30: ``Die Ausuebung der staatlichen Befugnisse und die Erfuellung der staatlichen Aufgaben ist Sache der Laender, soweit dieses Grundgesetz keine andere Regelung trifft oder zulaesst.''}}
\newcommand{\ggfnArtThirtyZH}{\ggnotezhonce{art-30}{第30條：「國家權限之行使與國家任務之履行，屬於各邦，但本基本法另有規定或容許者，不在此限。」}}

% ===== Art. 74 =====
\newcommand{\ggfnArtSeventyFourOneDE}{\ggnotedeonce{art-74-1}{Art. 74 Nr. 1: Die konkurrierende Gesetzgebung erstreckt sich auf ``das buergerliche Recht, das Strafrecht und den Strafvollzug, die Gerichtsverfassung, das gerichtliche Verfahren, die Rechtsanwaltschaft, das Notariat und die Rechtsberatung''.}}
\newcommand{\ggfnArtSeventyFourOneZH}{\ggnotezhonce{art-74-1}{第74條第1款：競合立法及於「民法、刑法及刑罰執行、法院組織、訴訟程序、律師制度、公證制度及法律諮詢」。}}
\newcommand{\ggfnArtSeventyFourSevenDE}{\ggnotedeonce{art-74-7}{Art. 74 Nr. 7: Die konkurrierende Gesetzgebung erstreckt sich auf ``die oeffentliche Fuersorge''.}}
\newcommand{\ggfnArtSeventyFourSevenZH}{\ggnotezhonce{art-74-7}{第74條第7款：競合立法及於「公共扶助」。}}
\newcommand{\ggfnArtSeventyFourTwelveDE}{\ggnotedeonce{art-74-12}{Art. 74 Nr. 12: Die konkurrierende Gesetzgebung erstreckt sich auf ``das Arbeitsrecht einschliesslich der Betriebsverfassung, des Arbeitsschutzes und der Arbeitsvermittlung sowie die Sozialversicherung einschliesslich der Arbeitslosenversicherung''.}}
\newcommand{\ggfnArtSeventyFourTwelveZH}{\ggnotezhonce{art-74-12}{第74條第12款：競合立法及於「勞動法，包括企業組織、勞動保護及就業媒合，以及社會保險，包括失業保險」。}}

% ===== Art. 77（墮胎一案特有）=====
\newcommand{\ggfnArtSevenSevenThreeDE}{\ggnotedeonce{art-77-3}{Art. 77 Abs. 3 Satz 1: ``Soweit zu einem Gesetze die Zustimmung des Bundesrates nicht erforderlich ist, kann der Bundesrat, wenn das Verfahren nach Absatz 2 beendigt ist, gegen ein vom Bundestage beschlossenes Gesetz binnen zwei Wochen Einspruch einlegen.''}}
\newcommand{\ggfnArtSevenSevenThreeZH}{\ggnotezhonce{art-77-3}{第77條第3項第1句：「法律無須聯邦參議院同意者，於第2項程序終了後，聯邦參議院得於二週內對聯邦議院通過之法律提出異議。」}}

% ===== Art. 80 =====
\newcommand{\ggfnArtEightyOneOneDE}{\ggnotedeonce{art-80-1-1}{Art. 80 Abs. 1 Satz 1: ``Durch Gesetz koennen die Bundesregierung, ein Bundesminister oder die Landesregierungen ermaechtigt werden, Rechtsverordnungen zu erlassen.''}}
\newcommand{\ggfnArtEightyOneOneZH}{\ggnotezhonce{art-80-1-1}{第80條第1項第1句：「得以法律授權聯邦政府、聯邦部長或邦政府發布法規命令。」}}

% ===== Art. 83 =====
\newcommand{\ggfnArtEightyThreeDE}{\ggnotedeonce{art-83}{Art. 83: ``Die Laender fuehren die Bundesgesetze als eigene Angelegenheit aus, soweit dieses Grundgesetz nichts anderes bestimmt oder zulaesst.''}}
\newcommand{\ggfnArtEightyThreeZH}{\ggnotezhonce{art-83}{第83條：「各邦以自己事務執行聯邦法律，但本基本法另有規定或容許者，不在此限。」}}

% ===== Art. 84 =====
\newcommand{\ggfnArtEightyFourOneDE}{\ggnotedeonce{art-84-1}{Art. 84 Abs. 1: ``Fuehren die Laender die Bundesgesetze als eigene Angelegenheit aus, so regeln sie die Einrichtung der Behoerden und das Verwaltungsverfahren, soweit nicht Bundesgesetze mit Zustimmung des Bundesrates etwas anderes bestimmen.''}}
\newcommand{\ggfnArtEightyFourOneZH}{\ggnotezhonce{art-84-1}{第84條第1項：「各邦以自己事務執行聯邦法律時，除聯邦法律經聯邦參議院同意另有規定外，由各邦規範機關之設置及行政程序。」}}
% 別名（墮胎一案使用之縮寫命名）
\let\ggfnArtEightFourOneDE\ggfnArtEightyFourOneDE
\let\ggfnArtEightFourOneZH\ggfnArtEightyFourOneZH

% ===== Art. 93 =====
\newcommand{\ggfnArtNinetyThreeOneTwoDE}{\ggnotedeonce{art-93-1-2}{Art. 93 Abs. 1 Nr. 2: Das Bundesverfassungsgericht entscheidet ``bei Meinungsverschiedenheiten oder Zweifeln ueber die foermliche und sachliche Vereinbarkeit von Bundesrecht oder Landesrecht mit diesem Grundgesetze ... auf Antrag der Bundesregierung, einer Landesregierung oder eines Drittels der Mitglieder des Bundestages''.}}
\newcommand{\ggfnArtNinetyThreeOneTwoZH}{\ggnotezhonce{art-93-1-2}{第93條第1項第2款：聯邦憲法法院就「聯邦法或邦法在形式及實質上是否與本基本法相容之意見分歧或疑義……依聯邦政府、邦政府或聯邦議院三分之一議員之聲請」作成裁判。}}
% 別名（墮胎一案使用之縮寫命名）
\let\ggfnArtNineThreeOneTwoDE\ggfnArtNinetyThreeOneTwoDE
\let\ggfnArtNineThreeOneTwoZH\ggfnArtNinetyThreeOneTwoZH

% ===== Art. 102 =====
\newcommand{\ggfnArtOneZeroTwoDE}{\ggnotedeonce{art-102}{Art. 102: ``Die Todesstrafe ist abgeschafft.''}}
\newcommand{\ggfnArtOneZeroTwoZH}{\ggnotezhonce{art-102}{第102條：「死刑廢止。」}}

% ===== Art. 103 =====
\newcommand{\ggfnArtOneZeroThreeTwoDE}{\ggnotedeonce{art-103-2}{Art. 103 Abs. 2: ``Eine Tat kann nur bestraft werden, wenn die Strafbarkeit gesetzlich bestimmt war, bevor die Tat begangen wurde.''}}
\newcommand{\ggfnArtOneZeroThreeTwoZH}{\ggnotezhonce{art-103-2}{第103條第2項：「行為之可罰性，須於行為前以法律定之，始得處罰。」}}

% ===== Art. 104 =====
\newcommand{\ggfnArtOneZeroFourOneDE}{\ggnotedeonce{art-104-1}{Art. 104 Abs. 1: ``Die Freiheit der Person kann nur auf Grund eines foermlichen Gesetzes und nur unter Beachtung der darin vorgeschriebenen Formen beschraenkt werden.''}}
\newcommand{\ggfnArtOneZeroFourOneZH}{\ggnotezhonce{art-104-1}{第104條第1項：「人身自由僅得基於形式法律，並遵守其中規定之形式加以限制。」}}

% ===== 組合腳註：Art. 3 + Art. 6（墮胎一案特有）=====
\newcommand{\ggfnArtThreeSixDE}{\ggnotedeonce{art-3-6}{Art. 3: ``Alle Menschen sind vor dem Gesetz gleich. Maenner und Frauen sind gleichberechtigt. Niemand darf wegen seines Geschlechtes, seiner Abstammung, seiner Rasse, seiner Sprache, seiner Heimat und Herkunft, seines Glaubens, seiner religioesen oder politischen Anschauungen benachteiligt oder bevorzugt werden.'' Art. 6 Abs. 1: ``Ehe und Familie stehen unter dem besonderen Schutze der staatlichen Ordnung.'' Abs. 2 Satz 1: ``Pflege und Erziehung der Kinder sind das natuerliche Recht der Eltern und die zuvoerderst ihnen obliegende Pflicht.'' Abs. 4: ``Jede Mutter hat Anspruch auf den Schutz und die Fuersorge der Gemeinschaft.''}\ggmarknotede{art-3}\ggmarknotede{art-6}\ggmarknotede{art-6-1-4}}
\newcommand{\ggfnArtThreeSixZH}{\ggnotezhonce{art-3-6}{第3條：「所有人在法律前一律平等。男女享有平等權利。任何人不得因其性別、血統、種族、語言、故鄉及出身、信仰、宗教或政治見解而受不利益或優惠。」第6條第1項：「婚姻與家庭受國家秩序之特別保護。」第2項第1句：「照護及教育子女，為父母之自然權利，並為首先由其承擔之義務。」第4項：「每一母親均有請求共同體保護與照顧之權利。」}\ggmarknotezh{art-3}\ggmarknotezh{art-6}\ggmarknotezh{art-6-1-4}}

% ===== 組合腳註：Art. 6(1) + Art. 6(4)（墮胎一案特有）=====
\newcommand{\ggfnArtSixOneFourDE}{\ggnotedeonce{art-6-1-4}{Art. 6 Abs. 1: ``Ehe und Familie stehen unter dem besonderen Schutze der staatlichen Ordnung.'' Art. 6 Abs. 4: ``Jede Mutter hat Anspruch auf den Schutz und die Fuersorge der Gemeinschaft.''}\ggmarknotede{art-6}}
\newcommand{\ggfnArtSixOneFourZH}{\ggnotezhonce{art-6-1-4}{第6條第1項：「婚姻與家庭受國家秩序之特別保護。」第4項：「每一母親均有請求共同體保護與照顧之權利。」}\ggmarknotezh{art-6}}

% ===== 組合腳註：Art. 1(1) + Art. 2(2)(1)（墮胎二案特有）=====
\newcommand{\ggfnArtOneOneTwoTwoOneDE}{\ggnotedeonce{art-1-1+2-2-1}{Art. 1 Abs. 1: ``Die Wuerde des Menschen ist unantastbar. Sie zu achten und zu schuetzen ist Verpflichtung aller staatlichen Gewalt.'' Art. 2 Abs. 2 Satz 1: ``Jeder hat das Recht auf Leben und koerperliche Unversehrtheit.''}\ggmarknotede{art-1-1}\ggmarknotede{art-2-2-1}}
\newcommand{\ggfnArtOneOneTwoTwoOneZH}{\ggnotezhonce{art-1-1+2-2-1}{第1條第1項：「人之尊嚴不可侵犯。尊重及保護之，為一切國家權力之義務。」第2條第2項第1句：「人人有生命與身體完整性之權利。」}\ggmarknotezh{art-1-1}\ggmarknotezh{art-2-2-1}}

% ===== 組合腳註：Art. 1(1) + Art. 2(2)（墮胎二案特有）=====
\newcommand{\ggfnArtOneOneTwoTwoDE}{\ggnotedeonce{art-1-1+2-2}{Art. 1 Abs. 1: ``Die Wuerde des Menschen ist unantastbar. Sie zu achten und zu schuetzen ist Verpflichtung aller staatlichen Gewalt.'' Art. 2 Abs. 2: ``Jeder hat das Recht auf Leben und koerperliche Unversehrtheit. Die Freiheit der Person ist unverletzlich. In diese Rechte darf nur auf Grund eines Gesetzes eingegriffen werden.''}\ggmarknotede{art-1-1}\ggmarknotede{art-2-2}\ggmarknotede{art-2-2-1}}
\newcommand{\ggfnArtOneOneTwoTwoZH}{\ggnotezhonce{art-1-1+2-2}{第1條第1項：「人之尊嚴不可侵犯。尊重及保護之，為一切國家權力之義務。」第2條第2項：「人人有生命與身體完整性之權利。人身自由不可侵犯。此等權利僅得依法律加以干預。」}\ggmarknotezh{art-1-1}\ggmarknotezh{art-2-2}\ggmarknotezh{art-2-2-1}}

% ===== 組合腳註：Art. 1(1) + Art. 2(1)（墮胎二案特有）=====
\newcommand{\ggfnArtOneOneTwoOneDE}{\ggnotedeonce{art-1-1+2-1}{Art. 1 Abs. 1: ``Die Wuerde des Menschen ist unantastbar. Sie zu achten und zu schuetzen ist Verpflichtung aller staatlichen Gewalt.'' Art. 2 Abs. 1: ``Jeder hat das Recht auf die freie Entfaltung seiner Persoenlichkeit, soweit er nicht die Rechte anderer verletzt und nicht gegen die verfassungsmaessige Ordnung oder das Sittengesetz verstoesst.''}\ggmarknotede{art-1-1}\ggmarknotede{art-2-1}}
\newcommand{\ggfnArtOneOneTwoOneZH}{\ggnotezhonce{art-1-1+2-1}{第1條第1項：「人之尊嚴不可侵犯。尊重及保護之，為一切國家權力之義務。」第2條第1項：「人人有自由發展其人格之權利，但不得侵害他人權利，亦不得違反合憲秩序或道德律。」}\ggmarknotezh{art-1-1}\ggmarknotezh{art-2-1}}

% ===== 組合腳註：Art. 2(2)(1) + Art. 1(1)（墮胎一案特有，含交叉標記）=====
\newcommand{\ggfnLifeDignityDE}{\ggnotedeonce{life-dignity}{Art. 2 Abs. 2 Satz 1: ``Jeder hat das Recht auf Leben und koerperliche Unversehrtheit.'' Art. 1 Abs. 1: ``Die Wuerde des Menschen ist unantastbar. Sie zu achten und zu schuetzen ist Verpflichtung aller staatlichen Gewalt.''}\ggmarknotede{art-2-2-1}\ggmarknotede{art-1-1}\ggmarknotede{art-1-1-2}}
\newcommand{\ggfnLifeDignityZH}{\ggnotezhonce{life-dignity}{第2條第2項第1句：「人人有生命與身體完整性之權利。」第1條第1項：「人之尊嚴不可侵犯。尊重及保護之，為一切國家權力之義務。」}\ggmarknotezh{art-2-2-1}\ggmarknotezh{art-1-1}\ggmarknotezh{art-1-1-2}}

% ===== 組合腳註：Art. 1(1) + Art. 2(2)（墮胎一案特有，條件判斷版）=====
\newcommand{\ggfnArtOneTwoTwoDE}{\ifcsname ggnoteseen@de@life-dignity\endcsname\ggfnArtTwoTwoRestDE{}\else\ggnotedeonce{art-1-2-2}{Art. 1 Abs. 1: ``Die Wuerde des Menschen ist unantastbar. Sie zu achten und zu schuetzen ist Verpflichtung aller staatlichen Gewalt.'' Art. 2 Abs. 2: ``Jeder hat das Recht auf Leben und koerperliche Unversehrtheit. Die Freiheit der Person ist unverletzlich. In diese Rechte darf nur auf Grund eines Gesetzes eingegriffen werden.''}\ggmarknotede{art-1-1}\ggmarknotede{art-1-1-2}\ggmarknotede{art-2-2}\ggmarknotede{art-2-2-1}\fi}
\newcommand{\ggfnArtOneTwoTwoZH}{\ifcsname ggnoteseen@zh@life-dignity\endcsname\ggfnArtTwoTwoRestZH{}\else\ggnotezhonce{art-1-2-2}{第1條第1項：「人之尊嚴不可侵犯。尊重及保護之，為一切國家權力之義務。」第2條第2項：「人人有生命與身體完整性之權利。人身自由不可侵犯。此等權利僅得依法律加以干預。」}\ggmarknotezh{art-1-1}\ggmarknotezh{art-1-1-2}\ggmarknotezh{art-2-2}\ggmarknotezh{art-2-2-1}\fi}

% ===== 組合腳註：Art. 1(1) + Art. 19(2)（墮胎一案特有，條件判斷版）=====
\newcommand{\ggfnArtOneNineteenDE}{\ifcsname ggnoteseen@de@art-1-1\endcsname\ggfnArtNineteenTwoDE{}\else\ggnotedeonce{art-1-19}{Art. 1 Abs. 1: ``Die Wuerde des Menschen ist unantastbar. Sie zu achten und zu schuetzen ist Verpflichtung aller staatlichen Gewalt.'' Art. 19 Abs. 2: ``In keinem Falle darf ein Grundrecht in seinem Wesensgehalt angetastet werden.''}\ggmarknotede{art-1-1}\ggmarknotede{art-1-1-2}\ggmarknotede{art-19-2}\fi}
\newcommand{\ggfnArtOneNineteenZH}{\ifcsname ggnoteseen@zh@art-1-1\endcsname\ggfnArtNineteenTwoZH{}\else\ggnotezhonce{art-1-19}{第1條第1項：「人之尊嚴不可侵犯。尊重及保護之，為一切國家權力之義務。」第19條第2項：「任何情形下，基本權利之本質內容均不得受到侵害。」}\ggmarknotezh{art-1-1}\ggmarknotezh{art-1-1-2}\ggmarknotezh{art-19-2}\fi}
 之前定義）=====
\newif\ifggversionnotede
\newif\ifggversionnotezh

% ===== 編者註腳基礎設施（首次標記模式）=====
\newif\ifeditornotelabelde
\newif\ifeditornotelabelzh
\newcommand{\editornotelabelde}{\ifeditornotelabelde\else\emph{Hrsg.-Anm. (nachfolgend ohne Kennzeichnung)}: \global\editornotelabeldetrue\fi}
\newcommand{\editornotelabelzh}{\ifeditornotelabelzh\else 編者註（下同）：\global\editornotelabelzhtrue\fi}
\newcommand{\editornotede}[1]{\ifshoweditornotes\footnote{\normalfont\footnotesize\editornotelabelde #1}\else\ifclassroommode\footnote{\normalfont\footnotesize\editornotelabelde #1}\fi\fi}
\newcommand{\editornotezh}[1]{\ifshoweditornotes\footnote{\normalfont\footnotesize\editornotelabelzh #1}\else\ifclassroommode\footnote{\normalfont\footnotesize\editornotelabelzh #1}\fi\fi}

% ===== GG 引用系統（\ggversionde/zh 由各案在 \input 前定義）=====
\newcommand{\ggmarknotede}[1]{\global\expandafter\let\csname ggnoteseen@de@#1\endcsname\empty}
\newcommand{\ggmarknotezh}[1]{\global\expandafter\let\csname ggnoteseen@zh@#1\endcsname\empty}
\newcommand{\ggnotedeonce}[2]{\ifcsname ggnoteseen@de@#1\endcsname\else\global\expandafter\let\csname ggnoteseen@de@#1\endcsname\empty\ggnotede{#2}\fi}
\newcommand{\ggnotezhonce}[2]{\ifcsname ggnoteseen@zh@#1\endcsname\else\global\expandafter\let\csname ggnoteseen@zh@#1\endcsname\empty\ggnotezh{#2}\fi}
\newcommand{\ggnotede}[1]{\editornotede{\ggversionde #1}}
\newcommand{\ggnotezh}[1]{\editornotezh{\ggversionzh #1}}

% ===== Art. 1 =====
\newcommand{\ggfnArtOneOneDE}{\ggnotedeonce{art-1-1}{Art. 1 Abs. 1: ``Die Wuerde des Menschen ist unantastbar. Sie zu achten und zu schuetzen ist Verpflichtung aller staatlichen Gewalt.''}}
\newcommand{\ggfnArtOneOneZH}{\ggnotezhonce{art-1-1}{第1條第1項：「人之尊嚴不可侵犯。尊重及保護之，為一切國家權力之義務。」}}
\newcommand{\ggfnArtOneOneTwoDE}{\ggnotedeonce{art-1-1-2}{Art. 1 Abs. 1 Satz 2: ``Sie zu achten und zu schuetzen ist Verpflichtung aller staatlichen Gewalt.''}}
\newcommand{\ggfnArtOneOneTwoZH}{\ggnotezhonce{art-1-1-2}{第1條第1項第2句：「尊重及保護之，為一切國家權力之義務。」}}

% ===== Art. 2 =====
\newcommand{\ggfnArtTwoOneDE}{\ggnotedeonce{art-2-1}{Art. 2 Abs. 1: ``Jeder hat das Recht auf die freie Entfaltung seiner Persoenlichkeit, soweit er nicht die Rechte anderer verletzt und nicht gegen die verfassungsmaessige Ordnung oder das Sittengesetz verstoesst.''}}
\newcommand{\ggfnArtTwoOneZH}{\ggnotezhonce{art-2-1}{第2條第1項：「人人有自由發展其人格之權利，但不得侵害他人權利，亦不得違反合憲秩序或道德律。」}}
% 簡易預設版（墮胎一案以 \renewcommand 覆寫為含條件邏輯之版本）
\newcommand{\ggfnArtTwoTwoDE}{\ggnotedeonce{art-2-2}{Art. 2 Abs. 2: ``Jeder hat das Recht auf Leben und koerperliche Unversehrtheit. Die Freiheit der Person ist unverletzlich. In diese Rechte darf nur auf Grund eines Gesetzes eingegriffen werden.''}\ggmarknotede{art-2-2-1}}
\newcommand{\ggfnArtTwoTwoZH}{\ggnotezhonce{art-2-2}{第2條第2項：「人人有生命與身體完整性之權利。人身自由不可侵犯。此等權利僅得依法律加以干預。」}\ggmarknotezh{art-2-2-1}}
\newcommand{\ggfnArtTwoTwoOneDE}{\ggnotedeonce{art-2-2-1}{Art. 2 Abs. 2 Satz 1: ``Jeder hat das Recht auf Leben und koerperliche Unversehrtheit.''}}
\newcommand{\ggfnArtTwoTwoOneZH}{\ggnotezhonce{art-2-2-1}{第2條第2項第1句：「人人有生命與身體完整性之權利。」}}
\newcommand{\ggfnArtTwoTwoThreeDE}{\ggnotedeonce{art-2-2-3}{Art. 2 Abs. 2 Satz 3: ``In diese Rechte darf nur auf Grund eines Gesetzes eingegriffen werden.''}}
\newcommand{\ggfnArtTwoTwoThreeZH}{\ggnotezhonce{art-2-2-3}{第2條第2項第3句：「此等權利僅得依法律加以干預。」}}
% 墮胎一案特有：第2條第2項第2–3句（Art. 2(2) 去掉第1句後的「餘文」）
\newcommand{\ggfnArtTwoTwoRestDE}{\ggnotedeonce{art-2-2-rest}{Art. 2 Abs. 2 Saetze 2-3: ``Die Freiheit der Person ist unverletzlich. In diese Rechte darf nur auf Grund eines Gesetzes eingegriffen werden.''}\ggmarknotede{art-2-2}}
\newcommand{\ggfnArtTwoTwoRestZH}{\ggnotezhonce{art-2-2-rest}{第2條第2項第2、3句：「人身自由不可侵犯。此等權利僅得依法律加以干預。」}\ggmarknotezh{art-2-2}}

% ===== Art. 3 =====
\newcommand{\ggfnArtThreeDE}{\ggnotedeonce{art-3}{Art. 3: ``Alle Menschen sind vor dem Gesetz gleich. Maenner und Frauen sind gleichberechtigt. Niemand darf wegen seines Geschlechtes, seiner Abstammung, seiner Rasse, seiner Sprache, seiner Heimat und Herkunft, seines Glaubens, seiner religioesen oder politischen Anschauungen benachteiligt oder bevorzugt werden.''}}
\newcommand{\ggfnArtThreeZH}{\ggnotezhonce{art-3}{第3條：「所有人在法律前一律平等。男女享有平等權利。任何人不得因其性別、血統、種族、語言、故鄉及出身、信仰、宗教或政治見解而受不利益或優惠。」}}
\newcommand{\ggfnArtThreeTwoDE}{\ggnotedeonce{art-3-2}{Art. 3 Abs. 2: ``Maenner und Frauen sind gleichberechtigt.''}}
\newcommand{\ggfnArtThreeTwoZH}{\ggnotezhonce{art-3-2}{第3條第2項：「男女享有平等權利。」}}

% ===== Art. 4 =====
\newcommand{\ggfnArtFourOneDE}{\ggnotedeonce{art-4-1}{Art. 4 Abs. 1: ``Die Freiheit des Glaubens, des Gewissens und die Freiheit des religioesen und weltanschaulichen Bekenntnisses sind unverletzlich.''}}
\newcommand{\ggfnArtFourOneZH}{\ggnotezhonce{art-4-1}{第4條第1項：「信仰自由、良心自由，以及宗教與世界觀信念自由，不可侵犯。」}}

% ===== Art. 5 =====
\newcommand{\ggfnArtFiveOneDE}{\ggnotedeonce{art-5-1}{Art. 5 Abs. 1: ``Jeder hat das Recht, seine Meinung in Wort, Schrift und Bild frei zu aeussern und zu verbreiten und sich aus allgemein zugaenglichen Quellen ungehindert zu unterrichten. Die Pressefreiheit und die Freiheit der Berichterstattung durch Rundfunk und Film werden gewaehrleistet. Eine Zensur findet nicht statt.''}}
\newcommand{\ggfnArtFiveOneZH}{\ggnotezhonce{art-5-1}{第5條第1項：「人人有以言語、文字及圖像自由表達及傳播其意見，並由一般可接近來源不受阻礙獲取資訊之權利。新聞自由及廣播、電影報導自由受保障。不設審查制度。」}}

% ===== Art. 6 =====
\newcommand{\ggfnArtSixDE}{\ggnotedeonce{art-6}{Art. 6 Abs. 1: ``Ehe und Familie stehen unter dem besonderen Schutze der staatlichen Ordnung.'' Abs. 2 Satz 1: ``Pflege und Erziehung der Kinder sind das natuerliche Recht der Eltern und die zuvoerderst ihnen obliegende Pflicht.'' Abs. 4: ``Jede Mutter hat Anspruch auf den Schutz und die Fuersorge der Gemeinschaft.''}}
\newcommand{\ggfnArtSixZH}{\ggnotezhonce{art-6}{第6條第1項：「婚姻與家庭受國家秩序之特別保護。」第2項第1句：「照護及教育子女，為父母之自然權利，並為首先由其承擔之義務。」第4項：「每一母親均有請求共同體保護與照顧之權利。」}}
\newcommand{\ggfnArtSixOneDE}{\ggnotedeonce{art-6-1}{Art. 6 Abs. 1: ``Ehe und Familie stehen unter dem besonderen Schutze der staatlichen Ordnung.''}}
\newcommand{\ggfnArtSixOneZH}{\ggnotezhonce{art-6-1}{第6條第1項：「婚姻與家庭受國家秩序之特別保護。」}}
\newcommand{\ggfnArtSixTwoDE}{\ggnotedeonce{art-6-2}{Art. 6 Abs. 2: ``Pflege und Erziehung der Kinder sind das natuerliche Recht der Eltern und die zuvoerderst ihnen obliegende Pflicht. Ueber ihre Betaetigung wacht die staatliche Gemeinschaft.''}}
\newcommand{\ggfnArtSixTwoZH}{\ggnotezhonce{art-6-2}{第6條第2項：「照護及教育子女，為父母之自然權利，並為首先由其承擔之義務。國家共同體監督其行使。」}}
\newcommand{\ggfnArtSixTwoOneDE}{\ggnotedeonce{art-6-2-1}{Art. 6 Abs. 2 Satz 1: ``Pflege und Erziehung der Kinder sind das natuerliche Recht der Eltern und die zuvoerderst ihnen obliegende Pflicht.''}}
\newcommand{\ggfnArtSixTwoOneZH}{\ggnotezhonce{art-6-2-1}{第6條第2項第1句：「照護及教育子女，為父母之自然權利，並為首先由其承擔之義務。」}}
\newcommand{\ggfnArtSixFourDE}{\ggnotedeonce{art-6-4}{Art. 6 Abs. 4: ``Jede Mutter hat Anspruch auf den Schutz und die Fuersorge der Gemeinschaft.''}}
\newcommand{\ggfnArtSixFourZH}{\ggnotezhonce{art-6-4}{第6條第4項：「每一母親均有請求共同體保護與照顧之權利。」}}

% ===== Art. 12 =====
\newcommand{\ggfnArtTwelveOneDE}{\ggnotedeonce{art-12-1}{Art. 12 Abs. 1: ``Alle Deutschen haben das Recht, Beruf, Arbeitsplatz und Ausbildungsstaette frei zu waehlen. Die Berufsausuebung kann durch Gesetz oder auf Grund eines Gesetzes geregelt werden.''}}
\newcommand{\ggfnArtTwelveOneZH}{\ggnotezhonce{art-12-1}{第12條第1項：「所有德國人均有自由選擇職業、工作場所及訓練場所之權利。職業行使得以法律或基於法律加以規範。」}}

% ===== Art. 19 =====
\newcommand{\ggfnArtNineteenTwoDE}{\ggnotedeonce{art-19-2}{Art. 19 Abs. 2: ``In keinem Falle darf ein Grundrecht in seinem Wesensgehalt angetastet werden.''}}
\newcommand{\ggfnArtNineteenTwoZH}{\ggnotezhonce{art-19-2}{第19條第2項：「任何情形下，基本權利之本質內容均不得受到侵害。」}}

% ===== Art. 20 =====
\newcommand{\ggfnArtTwentyOneDE}{\ggnotedeonce{art-20-1}{Art. 20 Abs. 1: ``Die Bundesrepublik Deutschland ist ein demokratischer und sozialer Bundesstaat.''}}
\newcommand{\ggfnArtTwentyOneZH}{\ggnotezhonce{art-20-1}{第20條第1項：「德意志聯邦共和國為民主及社會之聯邦國。」}}
\newcommand{\ggfnArtTwentyThreeDE}{\ggnotedeonce{art-20-3}{Art. 20 Abs. 3: ``Die Gesetzgebung ist an die verfassungsmaessige Ordnung, die vollziehende Gewalt und die Rechtsprechung sind an Gesetz und Recht gebunden.''}}
\newcommand{\ggfnArtTwentyThreeZH}{\ggnotezhonce{art-20-3}{第20條第3項：「立法受合憲秩序拘束；行政權與司法權受法律與法拘束。」}}

% ===== Art. 26 + Art. 143（墮胎一案特有）=====
\newcommand{\ggfnArtTwoSixAndOneFourThreeDE}{\ggnotedeonce{art-26-143}{Art. 26 Abs. 1 Satz 2: ``Sie sind unter Strafe zu stellen.'' Art. 143 in der urspruenglichen Fassung (24. Mai 1949 bis 1. September 1951) enthielt Strafvorschriften zum Hochverrat; 1975 war Art. 143 bereits weggefallen.}}
\newcommand{\ggfnArtTwoSixAndOneFourThreeZH}{\ggnotezhonce{art-26-143}{第26條第1項第2句：「此等行為應規定為犯罪處罰。」第143條原始版本（1949年5月24日至1951年9月1日）含有關於叛國罪之刑罰規定；至1975年時，第143條已廢止。}}

% ===== Art. 28 =====
\newcommand{\ggfnArtTwentyEightOneDE}{\ggnotedeonce{art-28-1-1}{Art. 28 Abs. 1 Satz 1: ``Die verfassungsmaessige Ordnung in den Laendern muss den Grundsaetzen des republikanischen, demokratischen und sozialen Rechtsstaates im Sinne dieses Grundgesetzes entsprechen.''}}
\newcommand{\ggfnArtTwentyEightOneZH}{\ggnotezhonce{art-28-1-1}{第28條第1項第1句：「各邦之合憲秩序，須符合本基本法意義下共和、民主及社會法治國之原則。」}}

% ===== Art. 30 =====
\newcommand{\ggfnArtThirtyDE}{\ggnotedeonce{art-30}{Art. 30: ``Die Ausuebung der staatlichen Befugnisse und die Erfuellung der staatlichen Aufgaben ist Sache der Laender, soweit dieses Grundgesetz keine andere Regelung trifft oder zulaesst.''}}
\newcommand{\ggfnArtThirtyZH}{\ggnotezhonce{art-30}{第30條：「國家權限之行使與國家任務之履行，屬於各邦，但本基本法另有規定或容許者，不在此限。」}}

% ===== Art. 74 =====
\newcommand{\ggfnArtSeventyFourOneDE}{\ggnotedeonce{art-74-1}{Art. 74 Nr. 1: Die konkurrierende Gesetzgebung erstreckt sich auf ``das buergerliche Recht, das Strafrecht und den Strafvollzug, die Gerichtsverfassung, das gerichtliche Verfahren, die Rechtsanwaltschaft, das Notariat und die Rechtsberatung''.}}
\newcommand{\ggfnArtSeventyFourOneZH}{\ggnotezhonce{art-74-1}{第74條第1款：競合立法及於「民法、刑法及刑罰執行、法院組織、訴訟程序、律師制度、公證制度及法律諮詢」。}}
\newcommand{\ggfnArtSeventyFourSevenDE}{\ggnotedeonce{art-74-7}{Art. 74 Nr. 7: Die konkurrierende Gesetzgebung erstreckt sich auf ``die oeffentliche Fuersorge''.}}
\newcommand{\ggfnArtSeventyFourSevenZH}{\ggnotezhonce{art-74-7}{第74條第7款：競合立法及於「公共扶助」。}}
\newcommand{\ggfnArtSeventyFourTwelveDE}{\ggnotedeonce{art-74-12}{Art. 74 Nr. 12: Die konkurrierende Gesetzgebung erstreckt sich auf ``das Arbeitsrecht einschliesslich der Betriebsverfassung, des Arbeitsschutzes und der Arbeitsvermittlung sowie die Sozialversicherung einschliesslich der Arbeitslosenversicherung''.}}
\newcommand{\ggfnArtSeventyFourTwelveZH}{\ggnotezhonce{art-74-12}{第74條第12款：競合立法及於「勞動法，包括企業組織、勞動保護及就業媒合，以及社會保險，包括失業保險」。}}

% ===== Art. 77（墮胎一案特有）=====
\newcommand{\ggfnArtSevenSevenThreeDE}{\ggnotedeonce{art-77-3}{Art. 77 Abs. 3 Satz 1: ``Soweit zu einem Gesetze die Zustimmung des Bundesrates nicht erforderlich ist, kann der Bundesrat, wenn das Verfahren nach Absatz 2 beendigt ist, gegen ein vom Bundestage beschlossenes Gesetz binnen zwei Wochen Einspruch einlegen.''}}
\newcommand{\ggfnArtSevenSevenThreeZH}{\ggnotezhonce{art-77-3}{第77條第3項第1句：「法律無須聯邦參議院同意者，於第2項程序終了後，聯邦參議院得於二週內對聯邦議院通過之法律提出異議。」}}

% ===== Art. 80 =====
\newcommand{\ggfnArtEightyOneOneDE}{\ggnotedeonce{art-80-1-1}{Art. 80 Abs. 1 Satz 1: ``Durch Gesetz koennen die Bundesregierung, ein Bundesminister oder die Landesregierungen ermaechtigt werden, Rechtsverordnungen zu erlassen.''}}
\newcommand{\ggfnArtEightyOneOneZH}{\ggnotezhonce{art-80-1-1}{第80條第1項第1句：「得以法律授權聯邦政府、聯邦部長或邦政府發布法規命令。」}}

% ===== Art. 83 =====
\newcommand{\ggfnArtEightyThreeDE}{\ggnotedeonce{art-83}{Art. 83: ``Die Laender fuehren die Bundesgesetze als eigene Angelegenheit aus, soweit dieses Grundgesetz nichts anderes bestimmt oder zulaesst.''}}
\newcommand{\ggfnArtEightyThreeZH}{\ggnotezhonce{art-83}{第83條：「各邦以自己事務執行聯邦法律，但本基本法另有規定或容許者，不在此限。」}}

% ===== Art. 84 =====
\newcommand{\ggfnArtEightyFourOneDE}{\ggnotedeonce{art-84-1}{Art. 84 Abs. 1: ``Fuehren die Laender die Bundesgesetze als eigene Angelegenheit aus, so regeln sie die Einrichtung der Behoerden und das Verwaltungsverfahren, soweit nicht Bundesgesetze mit Zustimmung des Bundesrates etwas anderes bestimmen.''}}
\newcommand{\ggfnArtEightyFourOneZH}{\ggnotezhonce{art-84-1}{第84條第1項：「各邦以自己事務執行聯邦法律時，除聯邦法律經聯邦參議院同意另有規定外，由各邦規範機關之設置及行政程序。」}}
% 別名（墮胎一案使用之縮寫命名）
\let\ggfnArtEightFourOneDE\ggfnArtEightyFourOneDE
\let\ggfnArtEightFourOneZH\ggfnArtEightyFourOneZH

% ===== Art. 93 =====
\newcommand{\ggfnArtNinetyThreeOneTwoDE}{\ggnotedeonce{art-93-1-2}{Art. 93 Abs. 1 Nr. 2: Das Bundesverfassungsgericht entscheidet ``bei Meinungsverschiedenheiten oder Zweifeln ueber die foermliche und sachliche Vereinbarkeit von Bundesrecht oder Landesrecht mit diesem Grundgesetze ... auf Antrag der Bundesregierung, einer Landesregierung oder eines Drittels der Mitglieder des Bundestages''.}}
\newcommand{\ggfnArtNinetyThreeOneTwoZH}{\ggnotezhonce{art-93-1-2}{第93條第1項第2款：聯邦憲法法院就「聯邦法或邦法在形式及實質上是否與本基本法相容之意見分歧或疑義……依聯邦政府、邦政府或聯邦議院三分之一議員之聲請」作成裁判。}}
% 別名（墮胎一案使用之縮寫命名）
\let\ggfnArtNineThreeOneTwoDE\ggfnArtNinetyThreeOneTwoDE
\let\ggfnArtNineThreeOneTwoZH\ggfnArtNinetyThreeOneTwoZH

% ===== Art. 102 =====
\newcommand{\ggfnArtOneZeroTwoDE}{\ggnotedeonce{art-102}{Art. 102: ``Die Todesstrafe ist abgeschafft.''}}
\newcommand{\ggfnArtOneZeroTwoZH}{\ggnotezhonce{art-102}{第102條：「死刑廢止。」}}

% ===== Art. 103 =====
\newcommand{\ggfnArtOneZeroThreeTwoDE}{\ggnotedeonce{art-103-2}{Art. 103 Abs. 2: ``Eine Tat kann nur bestraft werden, wenn die Strafbarkeit gesetzlich bestimmt war, bevor die Tat begangen wurde.''}}
\newcommand{\ggfnArtOneZeroThreeTwoZH}{\ggnotezhonce{art-103-2}{第103條第2項：「行為之可罰性，須於行為前以法律定之，始得處罰。」}}

% ===== Art. 104 =====
\newcommand{\ggfnArtOneZeroFourOneDE}{\ggnotedeonce{art-104-1}{Art. 104 Abs. 1: ``Die Freiheit der Person kann nur auf Grund eines foermlichen Gesetzes und nur unter Beachtung der darin vorgeschriebenen Formen beschraenkt werden.''}}
\newcommand{\ggfnArtOneZeroFourOneZH}{\ggnotezhonce{art-104-1}{第104條第1項：「人身自由僅得基於形式法律，並遵守其中規定之形式加以限制。」}}

% ===== 組合腳註：Art. 3 + Art. 6（墮胎一案特有）=====
\newcommand{\ggfnArtThreeSixDE}{\ggnotedeonce{art-3-6}{Art. 3: ``Alle Menschen sind vor dem Gesetz gleich. Maenner und Frauen sind gleichberechtigt. Niemand darf wegen seines Geschlechtes, seiner Abstammung, seiner Rasse, seiner Sprache, seiner Heimat und Herkunft, seines Glaubens, seiner religioesen oder politischen Anschauungen benachteiligt oder bevorzugt werden.'' Art. 6 Abs. 1: ``Ehe und Familie stehen unter dem besonderen Schutze der staatlichen Ordnung.'' Abs. 2 Satz 1: ``Pflege und Erziehung der Kinder sind das natuerliche Recht der Eltern und die zuvoerderst ihnen obliegende Pflicht.'' Abs. 4: ``Jede Mutter hat Anspruch auf den Schutz und die Fuersorge der Gemeinschaft.''}\ggmarknotede{art-3}\ggmarknotede{art-6}\ggmarknotede{art-6-1-4}}
\newcommand{\ggfnArtThreeSixZH}{\ggnotezhonce{art-3-6}{第3條：「所有人在法律前一律平等。男女享有平等權利。任何人不得因其性別、血統、種族、語言、故鄉及出身、信仰、宗教或政治見解而受不利益或優惠。」第6條第1項：「婚姻與家庭受國家秩序之特別保護。」第2項第1句：「照護及教育子女，為父母之自然權利，並為首先由其承擔之義務。」第4項：「每一母親均有請求共同體保護與照顧之權利。」}\ggmarknotezh{art-3}\ggmarknotezh{art-6}\ggmarknotezh{art-6-1-4}}

% ===== 組合腳註：Art. 6(1) + Art. 6(4)（墮胎一案特有）=====
\newcommand{\ggfnArtSixOneFourDE}{\ggnotedeonce{art-6-1-4}{Art. 6 Abs. 1: ``Ehe und Familie stehen unter dem besonderen Schutze der staatlichen Ordnung.'' Art. 6 Abs. 4: ``Jede Mutter hat Anspruch auf den Schutz und die Fuersorge der Gemeinschaft.''}\ggmarknotede{art-6}}
\newcommand{\ggfnArtSixOneFourZH}{\ggnotezhonce{art-6-1-4}{第6條第1項：「婚姻與家庭受國家秩序之特別保護。」第4項：「每一母親均有請求共同體保護與照顧之權利。」}\ggmarknotezh{art-6}}

% ===== 組合腳註：Art. 1(1) + Art. 2(2)(1)（墮胎二案特有）=====
\newcommand{\ggfnArtOneOneTwoTwoOneDE}{\ggnotedeonce{art-1-1+2-2-1}{Art. 1 Abs. 1: ``Die Wuerde des Menschen ist unantastbar. Sie zu achten und zu schuetzen ist Verpflichtung aller staatlichen Gewalt.'' Art. 2 Abs. 2 Satz 1: ``Jeder hat das Recht auf Leben und koerperliche Unversehrtheit.''}\ggmarknotede{art-1-1}\ggmarknotede{art-2-2-1}}
\newcommand{\ggfnArtOneOneTwoTwoOneZH}{\ggnotezhonce{art-1-1+2-2-1}{第1條第1項：「人之尊嚴不可侵犯。尊重及保護之，為一切國家權力之義務。」第2條第2項第1句：「人人有生命與身體完整性之權利。」}\ggmarknotezh{art-1-1}\ggmarknotezh{art-2-2-1}}

% ===== 組合腳註：Art. 1(1) + Art. 2(2)（墮胎二案特有）=====
\newcommand{\ggfnArtOneOneTwoTwoDE}{\ggnotedeonce{art-1-1+2-2}{Art. 1 Abs. 1: ``Die Wuerde des Menschen ist unantastbar. Sie zu achten und zu schuetzen ist Verpflichtung aller staatlichen Gewalt.'' Art. 2 Abs. 2: ``Jeder hat das Recht auf Leben und koerperliche Unversehrtheit. Die Freiheit der Person ist unverletzlich. In diese Rechte darf nur auf Grund eines Gesetzes eingegriffen werden.''}\ggmarknotede{art-1-1}\ggmarknotede{art-2-2}\ggmarknotede{art-2-2-1}}
\newcommand{\ggfnArtOneOneTwoTwoZH}{\ggnotezhonce{art-1-1+2-2}{第1條第1項：「人之尊嚴不可侵犯。尊重及保護之，為一切國家權力之義務。」第2條第2項：「人人有生命與身體完整性之權利。人身自由不可侵犯。此等權利僅得依法律加以干預。」}\ggmarknotezh{art-1-1}\ggmarknotezh{art-2-2}\ggmarknotezh{art-2-2-1}}

% ===== 組合腳註：Art. 1(1) + Art. 2(1)（墮胎二案特有）=====
\newcommand{\ggfnArtOneOneTwoOneDE}{\ggnotedeonce{art-1-1+2-1}{Art. 1 Abs. 1: ``Die Wuerde des Menschen ist unantastbar. Sie zu achten und zu schuetzen ist Verpflichtung aller staatlichen Gewalt.'' Art. 2 Abs. 1: ``Jeder hat das Recht auf die freie Entfaltung seiner Persoenlichkeit, soweit er nicht die Rechte anderer verletzt und nicht gegen die verfassungsmaessige Ordnung oder das Sittengesetz verstoesst.''}\ggmarknotede{art-1-1}\ggmarknotede{art-2-1}}
\newcommand{\ggfnArtOneOneTwoOneZH}{\ggnotezhonce{art-1-1+2-1}{第1條第1項：「人之尊嚴不可侵犯。尊重及保護之，為一切國家權力之義務。」第2條第1項：「人人有自由發展其人格之權利，但不得侵害他人權利，亦不得違反合憲秩序或道德律。」}\ggmarknotezh{art-1-1}\ggmarknotezh{art-2-1}}

% ===== 組合腳註：Art. 2(2)(1) + Art. 1(1)（墮胎一案特有，含交叉標記）=====
\newcommand{\ggfnLifeDignityDE}{\ggnotedeonce{life-dignity}{Art. 2 Abs. 2 Satz 1: ``Jeder hat das Recht auf Leben und koerperliche Unversehrtheit.'' Art. 1 Abs. 1: ``Die Wuerde des Menschen ist unantastbar. Sie zu achten und zu schuetzen ist Verpflichtung aller staatlichen Gewalt.''}\ggmarknotede{art-2-2-1}\ggmarknotede{art-1-1}\ggmarknotede{art-1-1-2}}
\newcommand{\ggfnLifeDignityZH}{\ggnotezhonce{life-dignity}{第2條第2項第1句：「人人有生命與身體完整性之權利。」第1條第1項：「人之尊嚴不可侵犯。尊重及保護之，為一切國家權力之義務。」}\ggmarknotezh{art-2-2-1}\ggmarknotezh{art-1-1}\ggmarknotezh{art-1-1-2}}

% ===== 組合腳註：Art. 1(1) + Art. 2(2)（墮胎一案特有，條件判斷版）=====
\newcommand{\ggfnArtOneTwoTwoDE}{\ifcsname ggnoteseen@de@life-dignity\endcsname\ggfnArtTwoTwoRestDE{}\else\ggnotedeonce{art-1-2-2}{Art. 1 Abs. 1: ``Die Wuerde des Menschen ist unantastbar. Sie zu achten und zu schuetzen ist Verpflichtung aller staatlichen Gewalt.'' Art. 2 Abs. 2: ``Jeder hat das Recht auf Leben und koerperliche Unversehrtheit. Die Freiheit der Person ist unverletzlich. In diese Rechte darf nur auf Grund eines Gesetzes eingegriffen werden.''}\ggmarknotede{art-1-1}\ggmarknotede{art-1-1-2}\ggmarknotede{art-2-2}\ggmarknotede{art-2-2-1}\fi}
\newcommand{\ggfnArtOneTwoTwoZH}{\ifcsname ggnoteseen@zh@life-dignity\endcsname\ggfnArtTwoTwoRestZH{}\else\ggnotezhonce{art-1-2-2}{第1條第1項：「人之尊嚴不可侵犯。尊重及保護之，為一切國家權力之義務。」第2條第2項：「人人有生命與身體完整性之權利。人身自由不可侵犯。此等權利僅得依法律加以干預。」}\ggmarknotezh{art-1-1}\ggmarknotezh{art-1-1-2}\ggmarknotezh{art-2-2}\ggmarknotezh{art-2-2-1}\fi}

% ===== 組合腳註：Art. 1(1) + Art. 19(2)（墮胎一案特有，條件判斷版）=====
\newcommand{\ggfnArtOneNineteenDE}{\ifcsname ggnoteseen@de@art-1-1\endcsname\ggfnArtNineteenTwoDE{}\else\ggnotedeonce{art-1-19}{Art. 1 Abs. 1: ``Die Wuerde des Menschen ist unantastbar. Sie zu achten und zu schuetzen ist Verpflichtung aller staatlichen Gewalt.'' Art. 19 Abs. 2: ``In keinem Falle darf ein Grundrecht in seinem Wesensgehalt angetastet werden.''}\ggmarknotede{art-1-1}\ggmarknotede{art-1-1-2}\ggmarknotede{art-19-2}\fi}
\newcommand{\ggfnArtOneNineteenZH}{\ifcsname ggnoteseen@zh@art-1-1\endcsname\ggfnArtNineteenTwoZH{}\else\ggnotezhonce{art-1-19}{第1條第1項：「人之尊嚴不可侵犯。尊重及保護之，為一切國家權力之義務。」第19條第2項：「任何情形下，基本權利之本質內容均不得受到侵害。」}\ggmarknotezh{art-1-1}\ggmarknotezh{art-1-1-2}\ggmarknotezh{art-19-2}\fi}


% 本案 Art. 1(1) + Art. 2(2)(1) 組合腳註條件邏輯：
% \ggfnArtTwoTwoDE 在標記 art-2-2 時同步標記 art-2-2-1；
% 若兩者均已引用，組合腳註不重複。若僅 Art. 1(1) 未引，輸出全文；
% 若僅 Art. 2(2)(1) 未引，輸出 Art. 2(2)(1)；否則不輸出。
\renewcommand{\ggfnArtOneOneTwoTwoOneDE}{%
  \ifcsname ggnoteseen@de@art-1-1+2-2-1\endcsname\else
    \global\expandafter\let\csname ggnoteseen@de@art-1-1+2-2-1\endcsname\empty
    \ifcsname ggnoteseen@de@art-1-1\endcsname
      \ifcsname ggnoteseen@de@art-2-2-1\endcsname\else
        \ggnotede{Art. 2 Abs. 2 Satz 1: ``Jeder hat das Recht auf Leben und koerperliche Unversehrtheit.''}%
      \fi
    \else
      \ggnotede{Art. 1 Abs. 1: ``Die Wuerde des Menschen ist unantastbar. Sie zu achten und zu schuetzen ist Verpflichtung aller staatlichen Gewalt.'' Art. 2 Abs. 2 Satz 1: ``Jeder hat das Recht auf Leben und koerperliche Unversehrtheit.''}%
    \fi
    \ggmarknotede{art-1-1}\ggmarknotede{art-2-2-1}%
  \fi
}
\renewcommand{\ggfnArtOneOneTwoTwoOneZH}{%
  \ifcsname ggnoteseen@zh@art-1-1+2-2-1\endcsname\else
    \global\expandafter\let\csname ggnoteseen@zh@art-1-1+2-2-1\endcsname\empty
    \ifcsname ggnoteseen@zh@art-1-1\endcsname
      \ifcsname ggnoteseen@zh@art-2-2-1\endcsname\else
        \ggnotezh{第2條第2項第1句：「人人有生命與身體完整性之權利。」}%
      \fi
    \else
      \ggnotezh{第1條第1項：「人之尊嚴不可侵犯。尊重及保護之，為一切國家權力之義務。」第2條第2項第1句：「人人有生命與身體完整性之權利。」}%
    \fi
    \ggmarknotezh{art-1-1}\ggmarknotezh{art-2-2-1}%
  \fi
}
% 本案 Art. 1(1) + Art. 2(2) 組合腳註條件邏輯：
% 若兩條均已引用，不重複；若僅一條已引，僅顯示未引之條文。
\renewcommand{\ggfnArtOneOneTwoTwoDE}{%
  \ifcsname ggnoteseen@de@art-1-1+2-2\endcsname\else
    \global\expandafter\let\csname ggnoteseen@de@art-1-1+2-2\endcsname\empty
    \ifcsname ggnoteseen@de@art-1-1\endcsname
      \ifcsname ggnoteseen@de@art-2-2\endcsname\else
        \ggnotede{Art. 2 Abs. 2: ``Jeder hat das Recht auf Leben und koerperliche Unversehrtheit. Die Freiheit der Person ist unverletzlich. In diese Rechte darf nur auf Grund eines Gesetzes eingegriffen werden.''}%
      \fi
    \else
      \ggnotede{Art. 1 Abs. 1: ``Die Wuerde des Menschen ist unantastbar. Sie zu achten und zu schuetzen ist Verpflichtung aller staatlichen Gewalt.'' Art. 2 Abs. 2: ``Jeder hat das Recht auf Leben und koerperliche Unversehrtheit. Die Freiheit der Person ist unverletzlich. In diese Rechte darf nur auf Grund eines Gesetzes eingegriffen werden.''}%
    \fi
    \ggmarknotede{art-1-1}\ggmarknotede{art-2-2}\ggmarknotede{art-2-2-1}%
  \fi
}
\renewcommand{\ggfnArtOneOneTwoTwoZH}{%
  \ifcsname ggnoteseen@zh@art-1-1+2-2\endcsname\else
    \global\expandafter\let\csname ggnoteseen@zh@art-1-1+2-2\endcsname\empty
    \ifcsname ggnoteseen@zh@art-1-1\endcsname
      \ifcsname ggnoteseen@zh@art-2-2\endcsname\else
        \ggnotezh{第2條第2項：「人人有生命與身體完整性之權利。人身自由不可侵犯。此等權利僅得依法律加以干預。」}%
      \fi
    \else
      \ggnotezh{第1條第1項：「人之尊嚴不可侵犯。尊重及保護之，為一切國家權力之義務。」第2條第2項：「人人有生命與身體完整性之權利。人身自由不可侵犯。此等權利僅得依法律加以干預。」}%
    \fi
    \ggmarknotezh{art-1-1}\ggmarknotezh{art-2-2}\ggmarknotezh{art-2-2-1}%
  \fi
}

\newcommand{\tocde}{%
  \section*{\mdseries\bverfgetitlede Inhaltsverzeichnis}
  \tocpart{Entscheidungstext}
  \begin{tocitems}
  \item[Leitsätze] Amtliche Leitsätze
  \item[Urteil] Tenor
  \item[Gründe] Begründung
  \end{tocitems}
  \begin{tocsubitems}
  \item[A.] Verfahren und Gegenstand
  \item[B.] Zulässigkeit
  \item[C.] Verfassungsrechtliche Maßstäbe
  \item[D.] Prüfung der Neuregelung
  \item[E.] Ergebnis
  \end{tocsubitems}
  \begin{tocitems}
  \item[Abw. Meinung] Mahrenholz/Sommer
  \item[Abw. Meinung] Böckenförde
  \end{tocitems}
  \tocdashline
  \begin{tocitems}
  \item[Glossar] Terminologie
  \end{tocitems}
}
\newcommand{\toczh}{%
  \section*{\mdseries\bverfgetitlezh 目錄}
  \tocpart{判決原文部分}
  \begin{tocitems}
  \item[裁判要旨] 官方裁判要旨
  \item[判決] 主文
  \item[理由] 判決理由
  \end{tocitems}
  \begin{tocsubitems}
  \item[A.] 程序與審查標的
  \item[B.] 合法性
  \item[C.] 憲法審查標準
  \item[D.] 新法規範之審查
  \item[E.] 結論
  \end{tocsubitems}
  \begin{tocitems}
  \item[不同意見] Mahrenholz/Sommer
  \item[不同意見] Böckenförde
  \end{tocitems}
  \tocdashline
  \begin{tocitems}
  \item[譯語表] 術語對照
  \end{tocitems}
}
% 大綱雙語對照表格內容（由 outlinepage 環境包裝後插入文件）
\newcommand{\outlinepaired}{%
\opartrow{Entscheidungstext}{判決原文部分}%
\omainrow{Leitsätze}{Amtliche Leitsätze: staatliche Schutzpflicht für ungeborenes Leben; Beratungskonzept als verfassungsrechtliche Option; Anforderungen an Beratung, Arztpflichten und Drittschutz; spezifische Folgerungen für SFHG-Regelungen.}{裁判要旨}{裁判要旨：國家對未出生生命之保護義務；諮詢構想為合憲保護方案；對諮詢內容、醫師義務、第三人保護之要求；對《孕婦及家庭扶助法》各規定之具體結論。}%
\omainrow{Urteil}{Zweiter Senat, 28.\,Mai\,1993; Parteien der verbundenen Normenkontrollverfahren.}{判決}{第二庭，1993年5月28日；合併審理之各規範審查程序當事人。}%
\omainrow{Gründe}{Begründung der Entscheidung (A–E).}{理由}{判決理由（A–E）。}%
\opartrow{Gliederung der Gründe}{理由之大綱}%
\osecrow{A.}{\textit{Verfahren und Gegenstand}}{A.}{程序與審查標的}%
\ointrow{1}{Gegenstand: Genügen von 15.\,StÄG und SFHG der staatlichen Schutzpflicht für ungeborenes Leben.}{標的：《第十五刑法修正法》及《孕婦及家庭扶助法》是否符合國家保護義務。}%
\osubrow{2–31}{I.}{Inhalt des 15.\,StÄG und SFHG; Gesetzgebungsgeschichte seit BVerfGE 39,\,1.}{I.}{《第十五刑法修正法》及《孕婦及家庭扶助法》之內容；BVerfGE 39,\,1 後之立法沿革。}%
\osubrow{32–90}{II.}{Anträge und Verfassungsrügen aller Antragsteller; Stellungnahmen Dritter.}{II.}{所有聲請人之聲請狀及合憲性指摘；第三人意見陳述。}%
\osubrow{91}{III.}{Äußerungen weiterer Beteiligter.}{III.}{其他參與者之意見。}%
\osecrow{B.}{\textit{Zulässigkeit}}{B.}{合法性}%
\osubrow{92–100}{I.}{Erster Antrag (2\,BvF\,2/90): Bayern greift §§ 218b, 219 StGB\,a.F. und §§ 200f, 200g RVO an.}{I.}{第一件聲請（2\,BvF\,2/90）：巴伐利亞就§§\,218b、219 StGB 舊法及§§\,200f、200g RVO 提出聲請。}%
\osubrow{101–108}{II.}{Stellungnahmen der Bundesregierung und der Länder zum ersten Antrag.}{II.}{聯邦政府及各邦政府就第一件聲請之意見陳述。}%
\osecrow{C.}{\textit{Verfassungsrechtliche Maßstäbe}}{C.}{憲法審查標準}%
\osubrow{109–110}{I.}{Zweiter und dritter Antrag (2\,BvF\,4, 5/92): Bayern und 249 MdB greifen § 218a Abs.\,1 n.F. und § 219 n.F. an.}{I.}{第二、三件聲請（2\,BvF\,4, 5/92）：巴伐利亞及249名聯邦議員就新版§\,218a(1)及§\,219提出聲請。}%
\osubrow{111–127}{II.}{Antragsbegründungen: Verstoß gegen Art.\,1 Abs.\,1 i.\,V.\,m.\ Art.\,2 Abs.\,2 S.\,1 GG.}{II.}{聲請理由：違反《基本法》第1條第1項結合第2條第2項第1句。}%
\osubrow{128–144}{III.}{Äußerungen der Verfassungsorgane und sonstiger Beteiligter.}{III.}{憲法機構及其他參與者之意見陳述。}%
\osubrow{145–148}{IV.}{Sachverständigengutachten.}{IV.}{鑑定意見。}%
\osubrow{149}{V.}{Zulässigkeit der Anträge.}{V.}{聲請合法性裁定。}%
\osecrow{D.}{\textit{Prüfung der Neuregelung}}{D.}{新法規範之審查}%
\osubrow{150–183}{I.}{Verfassungsrechtliche Maßstäbe: staatliche Schutzpflicht (Art.\,1 Abs.\,1 i.\,V.\,m.\ Art.\,2 Abs.\,2 GG); Menschenwürde des Ungeborenen; Abwägung mit Grundrechten der Frau.}{I.}{憲法審查標準：國家保護義務（第1條第1項結合第2條第2項）；未出生者之人性尊嚴；與婦女基本權之衡量。}%
\osubsubrow{150–153}{1.}{Schutzpflicht und Menschenwürde des Ungeborenen als verfassungsrechtliche Grundlage.}{1.}{保護義務與未出生者之人性尊嚴作為憲法基礎。}%
\osubsubsubrow{151–152}{a)}{Menschenwürde kommt schon dem ungeborenen Leben zu; Schutz ab Nidation.}{a)}{人性尊嚴已歸屬於未出生生命；自著床起受保護。}%
\osubsubsubrow{153}{b)}{Schutzpflicht bezogen auf das einzelne Leben, nicht nur auf menschliches Leben allgemein.}{b)}{保護義務指向個別生命，而非僅一般意義之人類生命。}%
\osubsubrow{154–171}{2.}{Verhaltensanforderungen durch gesetzliche Gebote; Untermaßverbot als Mindeststandard.}{2.}{透過法律命令提出行為要求；不足禁止作為最低標準。}%
\osubsubsubrow{156–157}{a)}{Verhaltensgebote bewirken Prävention und Normstabilisierung.}{a)}{行為命令產生預防效果並穩定規範。}%
\osubsubsubrow{158–159}{b)}{Schutz des Lebens nicht absolut; Abwägung mit Grundrechten der Frau erforderlich.}{b)}{生命保護並非絕對；須與婦女基本權進行衡量。}%
\osubsubsubrow{160–171}{c)}{Untermaßverbot: Mindestanforderungen, die der Staat nicht unterschreiten darf.}{c)}{不足禁止：國家不得低於之最低要求。}%
\osubsubsubsubrow{161–162}{aa)}{Grundsätzliches Verbot des Schwangerschaftsabbruchs für gesamte Dauer; Bestätigung BVerfGE 39,\,1.}{aa)}{整個懷孕期間之原則禁止；確認 BVerfGE 39,\,1 之見解。}%
\osubsubsubsubrow{163–166}{bb)}{Ausnahmetatbestände; Unzumutbarkeit der Frau als Grenze der Austragungspflicht.}{bb)}{例外事由；不合理負擔作為懷孕義務之界限。}%
\osubsubsubsubrow{167}{cc)}{Unzumutbarkeit begrenzt Austragungspflicht, nicht aber staatliche Schutzpflicht.}{cc)}{不合理負擔限制懷孕義務，但不限制國家之保護義務。}%
\osubsubsubsubrow{168–171}{dd)}{Strafrecht als unverzichtbarer Mindestschutz; Untermaßverbot gebietet strafrechtliche Reaktion.}{dd)}{刑法作為不可或缺之最低保護；不足禁止命令刑法回應。}%
\osubsubrow{172–183}{3.}{Schutz nicht nur durch Verbot: soziale Hilfen und positive Maßnahmen als Ergänzung.}{3.}{保護不限於禁止：社會扶助及積極措施作為補充。}%
\osubsubsubrow{174–176}{a)}{Fürsorgepflicht: soziale und wirtschaftliche Hilfen für Mutter; Beteiligung Dritter.}{a)}{扶助義務：對母親之社會及經濟協助；第三人之參與。}%
\osubsubsubrow{177}{b)}{Schutzauftrag für Ehe und Familie (Art.\,6 GG).}{b)}{婚姻與家庭保護委託（第6條）。}%
\osubsubsubrow{178}{c)}{Berufliche Absicherung bei Erwerbsverzicht zugunsten des Kindes.}{c)}{為子女利益放棄就業者之職業保障。}%
\osubsubsubrow{179}{d)}{Rechtlicher Schutzanspruch des Ungeborenen auch gegenüber Dritten.}{d)}{未出生者對第三人亦享有法律保護請求。}%
\osubsubrow{180–183}{4.}{Gesamtanforderungen an das Schutzkonzept; Spielraum und Grenzen des Gesetzgebers.}{4.}{對保護構想之整體要求；立法者之形成空間與界限。}%
\osubsubsubrow{182}{a)}{Einschätzungs- und Gestaltungsspielraum bei Wahl und Ausgestaltung des Schutzkonzepts.}{a)}{選擇及設計保護構想時之評估與形成空間。}%
\osubsubsubrow{183}{b)}{Verfassungsgerichtliche Kontrolle: Untermaßverbot als Prüfungsmaßstab.}{b)}{憲法法院之審查：不足禁止作為審查標準。}%
\osubrow{184–196}{II.}{Beratungskonzept grundsätzlich verfassungsrechtlich zulässig; Übergang als grundsätzlich zulässige Entscheidung; Einschätzungsprärogative des SFHG-Gesetzgebers.}{II.}{諮詢構想在憲法上原則可行；轉向諮詢構想作為原則允許之立法選擇；《孕婦及家庭扶助法》立法者之評估可支持。}%
\osubsubrow{186}{1.}{Neuregelung durch Einigungsvertrag veranlasst; gesamtdeutsche Regelung erforderlich.}{1.}{新規制由《統一條約》所促成；需要適用全德之統一規定。}%
\osubsubrow{187–188}{2.}{Gesetzgeber berechtigt und verpflichtet zu Neukonzept; Beratung als wirksames Schutzmittel.}{2.}{立法者有權且有義務採用新構想；諮詢作為有效保護手段。}%
\osubsubrow{189–190}{3.}{Einschätzung des Gesetzgebers verfassungsrechtlich unbedenklich; Erfahrungen der Beratungspraxis.}{3.}{立法者之評估在憲法上無疑義；諮詢實務之經驗支持。}%
\osubsubrow{191–192}{4.}{Letztverantwortung der Frau vereinbar mit ihr geschuldeter Achtung; nicht schrankenlose Entscheidungsfreiheit.}{4.}{婦女最後責任與應給予其之尊重相容；非無限制之決定自由。}%
\osubsubrow{193–196}{5.}{Letztverantwortung nur für allgemeine Notlagenindikation; spezifische Indikationen bleiben unberührt.}{5.}{最後責任僅適用於一般困境指標事由；特定指標事由不受影響。}%
\osubsubsubrow{194}{a)}{Erfahrungen der Beratungspraxis: Beratung bewirkt Fortsetzung der Schwangerschaft.}{a)}{諮詢實務經驗：諮詢使繼續懷孕成為可能。}%
\osubsubsubrow{195}{b)}{Gilt nicht für medizinische, eugenische, kriminologische Indikation.}{b)}{不適用於醫學、優生及犯罪學指標事由。}%
\osubsubsubrow{196}{c)}{Wissenschaftlich und rechtspolitisch umstritten; Einschätzungsprärogative bleibt beim Gesetzgeber.}{c)}{學術及法政策上爭議持續；評估優先權仍屬立法者。}%
\osubrow{197–220}{III.}{Anforderungen an verfassungskonformes Beratungskonzept; Beratungskonzept bedeutet nicht Aufgabe des Schutzziels; § 218a Abs.\,1 n.F. darf keine Rechtfertigung begründen.}{III.}{合憲諮詢構想之要求；諮詢構想不意味放棄保護目標；新版§\,218a(1)不得構成正當化事由。}%
\osubsubrow{198–202}{1.}{Notwendige Rahmenbedingungen eines verfassungskonformen Beratungskonzepts.}{1.}{合憲諮詢構想之必要框架條件。}%
\osubsubsubrow{198}{a)}{Zielorientierung der Beratung auf Schutz des ungeborenen Lebens muss erkennbar sein.}{a)}{諮詢之目標導向——保護未出生生命——必須清楚可辨。}%
\osubsubsubrow{199}{b)}{Einbeziehung aller Personen, die positiv oder negativ auf Entscheidung einwirken können.}{b)}{納入所有可能正面或負面影響決定之相關人員。}%
\osubsubsubrow{200–201}{c)}{§\,218a Abs.\,1 n.F. darf keinen Rechtfertigungsgrund begründen; Unzumutbarkeit muss vorliegen.}{c)}{新版§\,218a(1)不得構成正當化事由；不合理負擔須確實存在。}%
\osubsubsubrow{202}{d)}{Einmaligkeit jeder Schwangerschaft und individuelle Lebenssituation berücksichtigen.}{d)}{每次懷孕之唯一性及個別生活處境須納入考量。}%
\osubsubrow{203–216}{2.}{§\,218a Abs.\,1 n.F. als Rechtfertigungsgrund verfassungswidrig.}{2.}{新版§\,218a(1)作為正當化事由：違憲。}%
\osubsubsubrow{204–206}{a)}{Tatbestandsausschluss vs. Rechtfertigungsgrund: verfassungsrechtlich unterschiedliche Wirkungen.}{a)}{構成要件排除與正當化事由：在憲法上產生不同效果。}%
\osubsubsubrow{207–216}{b)}{Rechtsstaatliche Grundsätze verbieten Rechtfertigungsgrund ohne Nachweis der Unzumutbarkeit.}{b)}{法治國原則禁止在未證明不合理負擔之情況下成立正當化事由。}%
\osubsubsubsubrow{208–209}{aa)}{Unzumutbarkeitsnachweis nicht durch Generalfiktion aus Beratungsteilnahme ersetzbar.}{aa)}{不合理負擔之證明不能由參加諮詢之一般推定所取代。}%
\osubsubsubsubrow{210–213}{bb)}{Nachweis auch nicht durch Erfahrungssätze oder Indizien allgemein fingierbar.}{bb)}{證明亦不能透過經驗法則或跡象概括推定。}%
\osubsubsubsubrow{214–216}{cc)}{Schutzwirkung der Beratungsregelung hängt nicht von Rechtswidrigkeitserklärung ab.}{cc)}{諮詢規定之保護效果不取決於違法性之宣告。}%
\osubsubrow{217–219}{3.}{Herausnahme aus Straftatbestand lässt Raum für wirksame Schutzregelung; Anforderungen.}{3.}{排除於刑事構成要件之外仍為有效保護規定留有空間；相應要求。}%
\osubsubrow{220}{4.}{Beratungskonzept-Ergebnis: Frau trägt rechtliche Verantwortung; Abbruch bleibt rechtlich Unrecht.}{4.}{諮詢構想之結論：婦女承擔法律責任；終止妊娠在法律上仍屬不法。}%
\osubrow{221–244}{IV.}{Inhalt, Durchführung und Organisation der Beratung; Untermaßverbot.}{IV.}{諮詢之內容、實施及組織；不足禁止之約束。}%
\osubsubrow{222–233}{1.}{Inhalt der Beratung: Zielorientierung, Methodik, Information über Hilfen, familiäres Umfeld.}{1.}{諮詢內容：目標導向、方法論、關於協助之資訊、家庭環境。}%
\osubsubrow{234–236}{2.}{Durchführung: Möglichkeit mehrerer Gespräche; zeitlicher Abstand zum Abbruch geboten.}{2.}{實施：多次談話之可能性；與終止妊娠之間應有時間間距。}%
\osubsubrow{237–244}{3.}{Organisation: Anforderungen an Träger, Berater, Trennung von Abbruch und Kontrollpflicht.}{3.}{組織：對機構、諮詢人員之要求，與終止妊娠之分離及監督義務。}%
\osubrow{245–264}{V.}{Arztpflichten im Beratungskonzept; verfassungsrechtliche Mindestanforderungen.}{V.}{諮詢構想中醫師之義務；憲法最低要求。}%
\osubsubrow{247–255}{1.}{Inhalt der ärztlichen Pflichten: Untersuchung, Beratungsgespräch, Entscheidung über Mitwirkung.}{1.}{醫師義務之內容：檢查、諮詢談話、對參與之決定。}%
\osubsubrow{256–259}{2.}{Durchführung und strafrechtliche Absicherung; Mitwirkungsrecht und -pflicht des Arztes.}{2.}{實施與刑法保障；醫師之參與權利與義務。}%
\osubsubrow{260–264}{3.}{Informations-, Schweige-, Melde- und zivilrechtliche Haftungspflichten des Arztes (Nrn.\,3–6).}{3.}{醫師之告知、保密、通報義務及民事責任（第3–6項）。}%
\osubrow{265–271}{VI.}{Schutz vor Druckausübung durch das familiäre und soziale Umfeld.}{VI.}{防止家庭及社會環境施壓之保護義務。}%
\osubsubrow{266–267}{1.}{Wirksamkeit des Schutzkonzepts erfordert Schutz der Frau vor Zumutungen ihres Umfeldes.}{1.}{保護構想之有效性要求保護婦女免受其周遭環境之不當要求。}%
\osubsubrow{268–271}{2.}{Rechtsordnung muss entgegentreten; strafbewehrte Verhaltensgebote für familiäres Umfeld.}{2.}{法律秩序必須加以對抗；對家庭環境課予刑事保障之行為義務。}%
\osecrow{E.}{\textit{Ergebnis}}{E.}{結論}%
\ointrow{272}{SFHG verletzt die Schutzpflicht trotz zulässigen Übergangs zum Beratungskonzept.}{《孕婦及家庭扶助法》雖得轉向諮詢構想，仍未充分履行保護義務。}%
\osubrow{273–279}{I.}{§ 218a Abs.\,1 StGB n.F. mit Schutzpflicht unvereinbar und nichtig.}{I.}{新版§\,218a(1) StGB 與保護義務不相容，宣告無效。}%
\osubrow{280–299}{II.}{§ 219 StGB n.F. (Beratung) unzureichend; Nachbesserungspflicht des Gesetzgebers.}{II.}{新版§\,219 StGB（諮詢）不足；立法者之補正義務。}%
\osubrow{300}{III.}{Arzt- und Drittpflichten nicht hinreichend festgelegt; Nachholungspflicht.}{III.}{醫師及第三人義務規定不足；立法者之補足義務。}%
\osubrow{301–308}{IV.}{Art.\,15 Nr.\,2 SFHG (Aufhebung der Statistikpflicht) verfassungswidrig; Beobachtungs- und Korrekturpflicht.}{IV.}{SFHG 第15條第2款（廢止統計義務）違憲無效；持續觀察及修正義務。}%
\osubrow{309–348}{V.}{§ 24b SGB\,V: teils verfassungsgemäß, teils nicht (kein Anspruch für Abbrüche ohne verfassungskonforme Rechtmäßigkeitsfeststellung).}{V.}{SGB\,V 第24b條：部分合憲，部分不合憲（未經合法確認之終止妊娠不得獲給付）。}%
\osubrow{349–378}{VI.}{Art.\,4 5.\,StrRG n.F. (Einrichtungspflicht): teils nichtig wegen Verstoßes gegen das bundesstaatliche Prinzip.}{VI.}{新版《第五刑法改革法》第4條（確保機構義務）：部分因違反聯邦國原則而無效。}%
\odissentpart{Abweichende Meinungen}{不同意見}%
\odissentopinion{379–423}{Mahrenholz / Sommer}{Schutzpflicht durch Grundrechtsstatus der Frau begrenzt; § 218a Abs.\,1 n.F. verfassungsgemäß; § 24b SGB\,V zulässig.}{Mahrenholz / Sommer}{保護義務受婦女基本權地位限制；新版§\,218a(1)合憲；SGB\,V 第24b條給付允許。}%
\ointrow{379}{Letztverantwortung der Frau verfassungsrechtlich geboten; Schutzpflicht insoweit begrenzt; Beratungskonzept tragfähig.}{婦女最後責任為基本權所要求；保護義務就此受限；諮詢構想具備憲法基礎。}%
\osubrow{382–400}{I.}{Grundrechtsstatus der Frau begrenzt die staatliche Schutzpflicht; Letztverantwortung verfassungsrechtlich geboten.}{I.}{婦女基本權地位限制國家保護義務；最後責任在憲法上受到要求。}%
\osubrow{401–416}{II.}{§ 218a Abs.\,1 n.F. lässt sich als Rechtfertigungsgrund verfassungsrechtlich begründen.}{II.}{新版§\,218a(1)得以正當化事由在憲法上加以論證。}%
\osubrow{417–420}{III.}{§ 218a Abs.\,1 n.F. verfassungsgemäß; § 24b SGB\,V zulässig.}{III.}{新版§\,218a(1)合憲；SGB\,V 第24b條給付允許。}%
\osubrow{421–423}{IV.}{Schlußfolgerung.}{IV.}{綜合結論。}%
\odissentopinion{424–440}{Böckenförde}{Stimmt Untermaßverbot-Maßstäben zu; § 218a Abs.\,1 n.F. verfassungswidrig; § 24b SGB\,V zulässig.}{Böckenförde}{贊同不足禁止審查標準；§\,218a(1)違憲；但 SGB\,V 第24b條保險給付應屬允許。}%
}
\begin{document}
\pagenumbering{roman}
\hypersetup{pageanchor=false}

\begin{titlepage}
\thispagestyle{empty}
\vspace*{2.2cm}
\ifbilingualcolumns
  \noindent
  \begin{tabular}{@{}>{\centering\arraybackslash}p{0.485\textwidth}|>{\centering\arraybackslash}p{0.485\textwidth}@{}}
  {\bverfgetitlede\LARGE BVerfGE 88, 203\par}
  \vspace{0.35cm}
  {\bverfgetitlede\Large Schwangerschaftsabbruch II\par}
  \vspace{0.8cm}
  {\bverfgetitlede\Large Synoptische Ausgabe\par}
  \vspace{1.4cm}
  {\bverfgebodyde Bundesverfassungsgericht, Zweiter Senat\par}
  \vspace{0.35cm}
  {\bverfgebodyde Urteil vom 28. Mai 1993\par}
  \vspace{0.35cm}
  {\bverfgebodyde Aktenzeichen: 2 BvF 2/90 und 4, 5/92\par}
  &
  {\bverfgetitlezh\LARGE BVerfGE 88, 203\par}
  \vspace{0.35cm}
  {\bverfgetitlezh\Large Schwangerschaftsabbruch II\par}
  \vspace{0.8cm}
  {\bverfgetitlezh\Large 德中對照本\par}
  \vspace{1.4cm}
  {\bverfgebodyzh 德國聯邦憲法法院第二庭\par}
  \vspace{0.35cm}
  {\bverfgebodyzh 判決，1993年5月28日\par}
  \vspace{0.35cm}
  {\bverfgebodyzh 案號：2 BvF 2/90 und 4, 5/92\par}
  \end{tabular}
\else
  \begin{center}
  {\bverfgetitlede\LARGE BVerfGE 88, 203\par}
  \vspace{0.35cm}
  {\bverfgetitlede\Large Schwangerschaftsabbruch II\par}
  \vspace{0.8cm}
  \ifshowtranslation
    {\bverfgetitlezh\Large 中文譯本\par}
    \vspace{1.4cm}
    {\bverfgebodyzh 德國聯邦憲法法院第二庭\par}
    \vspace{0.35cm}
    {\bverfgebodyzh 判決，1993年5月28日\par}
    \vspace{0.35cm}
    {\bverfgebodyzh 案號：2 BvF 2/90 und 4, 5/92\par}
  \else
    {\bverfgetitlede\Large Amtliche Textfassung\par}
    \vspace{1.4cm}
    {\bverfgebodyde Bundesverfassungsgericht, Zweiter Senat\par}
    \vspace{0.35cm}
    {\bverfgebodyde Urteil vom 28. Mai 1993\par}
    \vspace{0.35cm}
    {\bverfgebodyde Aktenzeichen: 2 BvF 2/90 und 4, 5/92\par}
  \fi
  \end{center}
\fi
\bverfgetitlemeta
\end{titlepage}

\hypersetup{pageanchor=true}
\vspace*{1.0cm}
\begin{center}
\begin{minipage}{0.92\textwidth}
\ifbilingualcolumns
  \noindent
  \begin{tabular}{@{}p{0.47\textwidth}@{\hspace{0.025\textwidth}}!{\color{notegray}\vrule width 0.12pt}@{\hspace{0.025\textwidth}}p{0.47\textwidth}@{}}
  \tocpanel{\bverfgebodyde}{\tocde}
  &
  \tocpanel{\bverfgebodyzh}{\toczh}
  \end{tabular}
\else
  \tocsingleindent
  \ifshowtranslation
    \bverfgebodyzh
    \toczh
  \else
    \bverfgebodyde
    \tocde
  \fi
\fi
\end{minipage}
\end{center}
\clearpage
\ifbilingualcolumns
  \begin{outlinepage}\outlinepaired\end{outlinepage}
\fi
\clearpage
\pagenumbering{arabic}
\bilingualstart

\bisection{Leitsätze}{裁判要旨}
\syncleft
\begin{sourcepara}[title={原文要旨 1}]
\leitsatznum{1}Das Grundgesetz verpflichtet den Staat, menschliches Leben, auch das ungeborene, zu schützen. Diese Schutzpflicht hat ihren Grund in \abbrArt{} 1 \abbrAbs{} 1 \abbrGG{}\ggfnArtOneOneDE{}; ihr Gegenstand und -- von ihm her -- ihr Maß werden durch \abbrArt{} 2 \abbrAbs{} 2 \abbrGG{}\ggfnArtTwoTwoDE{} näher bestimmt. Menschenwürde kommt schon dem ungeborenen menschlichen Leben zu. Die Rechtsordnung muß die rechtlichen Voraussetzungen seiner Entfaltung im Sinne eines eigenen Lebensrechts des Ungeborenen gewährleisten. Dieses Lebensrecht wird nicht erst durch die Annahme seitens der Mutter begründet.
\end{sourcepara}
\begin{transpara}
\leitsatznum{1}《基本法》課予國家保護人類生命，亦即保護未出生生命（\textit{ungeborenes Leben}）之義務。此一保護義務（\textit{Schutzpflicht}）之根據在於《基本法》第1條第1項\ggfnArtOneOneZH{}；其客體及由客體所決定之尺度，則由第2條第2項進一步具體化。未出生的人類生命人性尊嚴。法律秩序必須保障其發展所需之法律前提，使未出生者享有自身生命權（\textit{Lebensrecht des Ungeborenen}）。此生命權並非因母親接納而始成立。
\end{transpara}

\syncleft
\begin{sourcepara}[title={原文要旨 2}]
\leitsatznum{2}Die Schutzpflicht für das ungeborene Leben ist bezogen auf das einzelne Leben, nicht nur auf menschliches Leben allgemein.
\end{sourcepara}
\begin{transpara}
\leitsatznum{2}對未出生生命之保護義務，係指向每一個別生命，而不僅是指向一般意義上之人類生命。
\end{transpara}

\syncleft
\begin{sourcepara}[title={原文要旨 3}]
\leitsatznum{3}Rechtlicher Schutz gebührt dem Ungeborenen auch gegenüber seiner Mutter. Ein solcher Schutz ist nur möglich, wenn der Gesetzgeber ihr einen Schwangerschaftsabbruch grundsätzlich verbietet und ihr damit die grundsätzliche Rechtspflicht auferlegt, das Kind auszutragen. Das grundsätzliche Verbot des Schwangerschaftsabbruchs und die grundsätzliche Pflicht zum Austragen des Kindes sind zwei untrennbar verbundene Elemente des verfassungsrechtlich gebotenen Schutzes.
\end{sourcepara}
\begin{transpara}
\leitsatznum{3}未出生者即使相對於其母親，亦應享有法律保護。此種保護只有在立法者原則上禁止母親終止妊娠，並由此課予其原則上繼續懷孕之法律義務時，始有可能。終止妊娠之原則禁止，與繼續懷孕之原則義務，是憲法所要求保護之兩個不可分割要素。
\end{transpara}

\syncleft
\begin{sourcepara}[title={原文要旨 4}]
\leitsatznum{4}Der Schwangerschaftsabbruch muß für die ganze Dauer der Schwangerschaft grundsätzlich als Unrecht angesehen und demgemäß rechtlich verboten sein (Bestätigung von \abbrBVerfGE{} 39, 1 [44]). Das Lebensrecht des Ungeborenen darf nicht, wenn auch nur für eine begrenzte Zeit, der freien, rechtlich nicht gebundenen Entscheidung eines Dritten, und sei es selbst der Mutter, überantwortet werden.
\end{sourcepara}
\begin{transpara}
\leitsatznum{4}終止妊娠於整個懷孕期間，原則上均必須被視為不法（\textit{Unrecht}），並相應地在法律上被禁止（確認 BVerfGE 39, 1 [44]之見解）。未出生者之生命權不得，即使僅於有限期間內，交由第三人自由且不受法律拘束之決定支配；即使該第三人是母親本人亦同。
\end{transpara}

\syncleft
\begin{sourcepara}[title={原文要旨 5}]
\leitsatznum{5}Die Reichweite der Schutzpflicht für das ungeborene menschliche Leben ist im Blick auf die Bedeutung und Schutzbedürftigkeit des zu schützenden Rechtsguts einerseits und damit kollidierender Rechtsgüter andererseits zu bestimmen. Als vom Lebensrecht des Ungeborenen berührte Rechtsgüter kommen dabei -- ausgehend vom Anspruch der schwangeren Frau auf Schutz und Achtung ihrer Menschenwürde (\abbrArt{} 1 \abbrAbs{} 1 \abbrGG{}\ggfnArtOneOneDE{}) -- vor allem ihr Recht auf Leben und körperliche Unversehrtheit (\abbrArt{} 2 \abbrAbs{} 2 \abbrGG{}\ggfnArtTwoTwoDE{}) sowie ihr Persönlichkeitsrecht (\abbrArt{} 2 \abbrAbs{} 1 \abbrGG{}\ggfnArtTwoOneDE{}) in Betracht. Dagegen kann die Frau für die mit dem Schwangerschaftsabbruch einhergehende Tötung des Ungeborenen nicht eine grundrechtlich in \abbrArt{} 4 \abbrAbs{} 1 \abbrGG{}\ggfnArtFourOneDE{} geschützte Rechtsposition in Anspruch nehmen.
\end{sourcepara}
\begin{transpara}
\leitsatznum{5}對未出生人類生命之保護義務範圍，須一方面著眼於受保護法益之重要性與保護需求，另一方面著眼於與之衝突之法益而定。就受未出生者生命權觸及之法益而言，從懷孕婦女受保護並受尊重其人性尊嚴之請求（《基本法》第1條第1項\ggfnArtOneOneZH{}）出發，尤其包括其生命與身體完整權（第2條第2項\ggfnArtTwoTwoZH{}）以及其人格權（第2條第1項\ggfnArtTwoOneZH{}）。相對地，婦女不得就伴隨終止妊娠而發生之殺害未出生者，主張一項受第4條第1項\ggfnArtFourOneZH{}基本權保護之法律地位。
\end{transpara}

\syncleft
\begin{sourcepara}[title={原文要旨 6}]
\leitsatznum{6}Der Staat muß zur Erfüllung seiner Schutzpflicht ausreichende Maßnahmen normativer und tatsächlicher Art ergreifen, die dazu führen, daß ein -- unter Berücksichtigung entgegenstehender Rechtsgüter -- angemessener und als solcher wirksamer Schutz erreicht wird (Untermaßverbot). Dazu bedarf es eines Schutzkonzepts, das Elemente des präventiven wie des repressiven Schutzes miteinander verbindet.
\end{sourcepara}
\begin{transpara}
\leitsatznum{6}國家為履行其保護義務，必須採取充分之規範性與事實性措施，使在考量相反法益後仍屬相當且作為相當保護確實有效之保護得以實現（不足禁止，\textit{Untermaßverbot}）。為此，需要一套結合預防性保護與壓制性保護要素之保護構想。
\end{transpara}

\syncleft
\begin{sourcepara}[title={原文要旨 7}]
\leitsatznum{7}Grundrechte der Frau tragen nicht so weit, daß die Rechtspflicht zum Austragen des Kindes -- auch nur für eine bestimmte Zeit -- generell aufgehoben wäre. Die Grundrechtspositionen der Frau führen allerdings dazu, daß es in Ausnahmelagen zulässig, in manchen dieser Fälle womöglich geboten ist, eine solche Rechtspflicht nicht aufzuerlegen. Es ist Sache des Gesetzgebers, solche Ausnahmetatbestände im einzelnen nach dem Kriterium der Unzumutbarkeit zu bestimmen. Dafür müssen Belastungen gegeben sein, die ein solches Maß an Aufopferung eigener Lebenswerte verlangen, daß dies von der Frau nicht erwartet werden kann (Bestätigung von \abbrBVerfGE{} 39, 1 [48 \abbrFF{}]).
\end{sourcepara}
\begin{transpara}
\leitsatznum{7}婦女之基本權尚不足以使繼續懷孕之法律義務，即使僅於特定期間內，得被一般性取消。然而，婦女之基本權地位確實導致：在例外處境中，不課予此種法律義務是容許的；在其中若干案件中，甚至可能是必須的。此等例外構成要件應由立法者依不可期待性（\textit{Unzumutbarkeit}）標準具體決定。其前提必須是存在某些負擔，要求婦女犧牲自身生命價值達到無法期待其承受之程度（確認 BVerfGE 39, 1 [48 ff.]之見解）。
\end{transpara}

\syncleft
\begin{sourcepara}[title={原文要旨 8}]
\leitsatznum{8}Das Untermaßverbot läßt es nicht zu, auf den Einsatz auch des Strafrechts und die davon ausgehende Schutzwirkung für das menschliche Leben frei zu verzichten.
\end{sourcepara}
\begin{transpara}
\leitsatznum{8}不足禁止不允許國家自由放棄使用刑法及其對人類生命所產生之保護效果。
\end{transpara}

\syncleft
\begin{sourcepara}[title={原文要旨 9}]
\leitsatznum{9}Die staatliche Schutzpflicht umfaßt auch den Schutz vor Gefahren, die für das ungeborene menschliche Leben von Einflüssen aus dem familiären oder weiteren sozialen Umfeld der Schwangeren oder von gegenwärtigen und absehbaren realen Lebensverhältnissen der Frau und der Familie ausgehen und der Bereitschaft zum Austragen des Kindes entgegenwirken.
\end{sourcepara}
\begin{transpara}
\leitsatznum{9}國家保護義務亦包括防範下列危險：由懷孕婦女家庭或更廣泛社會環境之影響，或由婦女及其家庭目前與可預見之實際生活關係所生、並對繼續懷孕之意願產生反作用之危險。
\end{transpara}

\syncleft
\begin{sourcepara}[title={原文要旨 10}]
\leitsatznum{10}Der Schutzauftrag verpflichtet den Staat ferner, den rechtlichen Schutzanspruch des ungeborenen Lebens im allgemeinen Bewußtsein zu erhalten und zu beleben.
\end{sourcepara}
\begin{transpara}
\leitsatznum{10}保護任務並進一步要求國家，在一般意識中維持並活化未出生生命之法律保護請求。
\end{transpara}

\syncleft
\begin{sourcepara}[title={原文要旨 11}]
\leitsatznum{11}Dem Gesetzgeber ist es verfassungsrechtlich grundsätzlich nicht verwehrt, zu einem Konzept für den Schutz des ungeborenen Lebens überzugehen, das in der Frühphase der Schwangerschaft in Schwangerschaftskonflikten den Schwerpunkt auf die Beratung der schwangeren Frau legt, um sie für das Austragen des Kindes zu gewinnen, und dabei auf eine indikationsbestimmte Strafdrohung und die Feststellung von Indikationstatbeständen durch einen Dritten verzichtet.
\end{sourcepara}
\begin{transpara}
\leitsatznum{11}憲法上原則上並不禁止立法者轉向一種保護未出生生命之構想：在懷孕初期之懷孕衝突中，將重點置於對懷孕婦女之諮詢，以爭取其繼續懷孕；並在此同時放棄由指標事由決定之刑罰威嚇，以及由第三人確認指標事由構成要件。
\end{transpara}

\syncleft
\begin{sourcepara}[title={原文要旨 12}]
\leitsatznum{12}Ein solches Beratungskonzept erfordert Rahmenbedingungen, die positive Voraussetzungen für ein Handeln der Frau zugunsten des ungeborenen Lebens schaffen. Der Staat trägt für die Durchführung des Beratungsverfahrens die volle Verantwortung.
\end{sourcepara}
\begin{transpara}
\leitsatznum{12}此種諮詢構想需要一套框架條件，以創造婦女為未出生生命而行動之正面前提。國家對諮詢程序之實施負全部責任。
\end{transpara}

\syncleft
\begin{sourcepara}[title={原文要旨 13}]
\leitsatznum{13}Die staatliche Schutzpflicht erfordert es, daß die im Interesse der Frau notwendige Beteiligung des Arztes zugleich Schutz für das ungeborene Leben bewirkt.
\end{sourcepara}
\begin{transpara}
\leitsatznum{13}國家保護義務要求：基於婦女利益而必要之醫師參與，必須同時對未出生生命產生保護效果。
\end{transpara}

\syncleft
\begin{sourcepara}[title={原文要旨 14}]
\leitsatznum{14}Eine rechtliche Qualifikation des Daseins eines Kindes als Schadensquelle kommt von Verfassungs wegen (\abbrArt{} 1 \abbrAbs{} 1 \abbrGG{}\ggfnArtOneOneDE{}) nicht in Betracht. Deshalb verbietet es sich, die Unterhaltspflicht für ein Kind als Schaden zu begreifen.
\end{sourcepara}
\begin{transpara}
\leitsatznum{14}基於憲法（《基本法》第1條第1項\ggfnArtOneOneZH{}），不得在法律上將一名兒童之存在評價為損害來源。因此，亦不得將對子女之扶養義務理解為損害。
\end{transpara}

\syncleft
\begin{sourcepara}[title={原文要旨 15}]
\leitsatznum{15}Schwangerschaftsabbrüche, die ohne Feststellung einer Indikation nach der Beratungsregelung vorgenommen werden, dürfen nicht für gerechtfertigt (nicht rechtswidrig) erklärt werden. Es entspricht unverzichtbaren rechtsstaatlichen Grundsätzen, daß einem Ausnahmetatbestand rechtfertigende Wirkung nur dann zukommen kann, wenn das Vorliegen seiner Voraussetzungen unter staatlicher Verantwortung festgestellt werden muß.
\end{sourcepara}
\begin{transpara}
\leitsatznum{15}依諮詢規制而未經指標事由確認所實施之終止妊娠，不得被宣告為正當化（不具違法性）。不可或缺之法治國原則要求：例外構成要件只有在其要件是否存在必須由國家負責確認時，始得具有正當化效果。
\end{transpara}

\syncleft
\begin{sourcepara}[title={原文要旨 16}]
\leitsatznum{16}Das Grundgesetz läßt es nicht zu, für die Vornahme eines Schwangerschaftsabbruchs, dessen Rechtmäßigkeit nicht festgestellt wird, einen Anspruch auf Leistungen der gesetzlichen Krankenversicherung zu gewähren. Die Gewährung von Sozialhilfe für nicht mit Strafe bedrohte Schwangerschaftsabbrüche nach der Beratungsregelung in Fällen wirtschaftlicher Bedürftigkeit ist demgegenüber ebensowenig verfassungsrechtlich zu beanstanden wie die Fortzahlung des Arbeitsentgelts.
\end{sourcepara}
\begin{transpara}
\leitsatznum{16}《基本法》不容許就其合法性未經確認之終止妊娠，給予法定醫療保險給付請求權。相對地，於經濟上有需要之情形，就依諮詢規制不受刑罰威嚇之終止妊娠給予社會扶助，與繼續給付工作報酬一樣，均不受憲法上非難。
\end{transpara}

\syncleft
\begin{sourcepara}[title={原文要旨 17}]
\leitsatznum{17}Der Grundsatz der Organisationsgewalt der Länder gilt uneingeschränkt, wenn eine bundesgesetzliche Regelung lediglich eine von den Ländern zu erfüllende Staatsaufgabe vorsieht, nicht jedoch Einzelregelungen trifft, die behördlich-administrativ vollzogen werden könnten.
\end{sourcepara}
\begin{transpara}
\leitsatznum{17}若聯邦法律規定僅設置一項應由各邦履行之國家任務，而未作成可由行政機關具體執行之個別規定，則各邦組織權原則不受限制地適用。
\end{transpara}

\bisection{Urteil}{判決}

\begin{sourcepara}[title={原文判決標示}]
Urteil des Zweiten Senats vom 28. Mai 1993 aufgrund der mündlichen Verhandlung vom 8. und 9. Dezember 1992 -- 2 \abbrBvF{} 2/90 und 4, 5/92 -- in den Verfahren [...]
\end{sourcepara}
\begin{transpara}
第二庭依1992年12月8日及9日言詞辯論，於1993年5月28日作成判決；案號為 2 BvF 2/90 及 4, 5/92；各程序當事人略。
\end{transpara}

\bisection{Tenor}{主文}

\begin{sourcepara}[title={原文主文 I.1}]
I. 1. § 218a Absatz 1 des Strafgesetzbuches in der Fassung des Gesetzes zum Schutz des vorgeburtlichen/werdenden Lebens, zur Förderung einer kinderfreundlicheren Gesellschaft, für Hilfen im Schwangerschaftskonflikt und zur Regelung des Schwangerschaftsabbruchs (Schwangeren- und Familienhilfegesetz) vom 27. Juli 1992 (\abbrBundesgesetzbl{} I Seite 1398) ist insoweit mit Artikel 1 Absatz 1 in Verbindung mit Artikel 2 Absatz 2 Satz 1 des Grundgesetzes\ggfnArtOneOneTwoTwoOneDE{} unvereinbar, als die Vorschrift den unter den dort genannten Voraussetzungen vorgenommenen Schwangerschaftsabbruch für nicht rechtswidrig erklärt und in Nummer 1 auf eine Beratung Bezug nimmt, die ihrerseits den verfassungsrechtlichen Anforderungen aus Artikel 1 Absatz 1 in Verbindung mit Artikel 2 Absatz 2 Satz 1 des Grundgesetzes\ggfnArtOneOneTwoTwoOneDE{} nicht genügt. Die Bestimmung ist insgesamt nichtig.
\end{sourcepara}
\begin{transpara}
I. 1. 《刑法》第218a條第1項，即1992年7月27日《保護出生前／形成中生命、促進更友善兒童之社會、提供懷孕衝突協助及規制終止妊娠之法律（孕婦及家庭扶助法）》（聯邦法律公報第一部第1398頁）版本，在下列範圍內與《基本法》第1條第1項結合第2條第2項第1句\ggfnArtOneOneTwoTwoOneZH{}不相容：該規定將於其所列要件下實施之終止妊娠宣告為不具違法性，且其第1款所指之諮詢本身不符合《基本法》第1條第1項結合第2條第2項第1句\ggfnArtOneOneTwoTwoOneZH{}之憲法要求。該規定整體無效。
\end{transpara}

\begin{sourcepara}[title={原文主文 I.2}]
2. § 219 des Strafgesetzbuches in der Fassung des genannten Gesetzes ist mit Artikel 1 Absatz 1 in Verbindung mit Artikel 2 Absatz 2 Satz 1 des Grundgesetzes\ggfnArtOneOneTwoTwoOneDE{} unvereinbar und nichtig.
\end{sourcepara}
\begin{transpara}
2. 上述法律版本之《刑法》第219條，與《基本法》第1條第1項結合第2條第2項第1句\ggfnArtOneOneTwoTwoOneZH{}不相容並無效。
\end{transpara}

\begin{sourcepara}[title={原文主文 I.3}]
3. § 24b des Fünften Buches Sozialgesetzbuch ist nach Maßgabe der Urteilsgründe mit Artikel 1 Absatz 1 in Verbindung mit Artikel 2 Absatz 2 Satz 1 des Grundgesetzes\ggfnArtOneOneTwoTwoOneDE{} vereinbar.
\end{sourcepara}
\begin{transpara}
3. 《社會法典》第五編第24b條，依本判決理由所定尺度，與《基本法》第1條第1項結合第2條第2項第1句\ggfnArtOneOneTwoTwoOneZH{}相容。
\end{transpara}

\begin{sourcepara}[title={原文主文 I.4}]
4. §§ 200f, 200g der Reichsversicherungsordnung in der Fassung des Gesetzes über ergänzende Maßnahmen zum Fünften Strafrechtsreformgesetz (Strafrechtsreform-Ergänzungsgesetz -- \abbrStREG{}) vom 28. August 1975 (\abbrBundesgesetzbl{} I Seite 2289) waren, soweit sie Leistungen der gesetzlichen Krankenversicherung bei Schwangerschaftsabbrüchen nach § 218a Absatz 2 Nummer 3 des Strafgesetzbuches in der Fassung des Fünfzehnten Strafrechtsänderungsgesetzes vom 18. Mai 1976 (\abbrBundesgesetzbl{} I Seite 1213) vorsahen, nach Maßgabe der Gründe mit Artikel 1 Absatz 1 in Verbindung mit Artikel 2 Absatz 2 Satz 1 des Grundgesetzes\ggfnArtOneOneTwoTwoOneDE{} vereinbar.
\end{sourcepara}
\begin{transpara}
4. 1975年8月28日《關於第五刑法改革法補充措施之法律（刑法改革補充法，StREG）》（聯邦法律公報第一部第2289頁）版本之《帝國保險法》第200f條、第200g條，在其就1976年5月18日《第十五刑法修正法》（聯邦法律公報第一部第1213頁）版本《刑法》第218a條第2項第3款之終止妊娠規定法定醫療保險給付之範圍內，依本判決理由，曾與《基本法》第1條第1項結合第2條第2項第1句\ggfnArtOneOneTwoTwoOneZH{}相容。
\end{transpara}

\begin{sourcepara}[title={原文主文 I.5}]
5. Artikel 15 Nummer 2 des Schwangeren- und Familienhilfegesetzes ist mit Artikel 1 Absatz 1 in Verbindung mit Artikel 2 Absatz 2 Satz 1 des Grundgesetzes\ggfnArtOneOneTwoTwoOneDE{} unvereinbar und nichtig, soweit dadurch die bisher in Artikel 4 des Fünften Gesetzes zur Reform des Strafrechts (5. \abbrStrRG{}) vom 18. Juni 1974 (\abbrBundesgesetzbl{} I Seite 1297), geändert durch Artikel 3 und Artikel 4 des Fünfzehnten Strafrechtsänderungsgesetzes vom 18. Mai 1976 (\abbrBundesgesetzbl{} I Seite 1213), enthaltene Vorschrift betreffend die Bundesstatistik über Schwangerschaftsabbrüche aufgehoben wird.
\end{sourcepara}
\begin{transpara}
5. 《孕婦及家庭扶助法》第15條第2款，在其廢止原載於1974年6月18日《第五刑法改革法（5. StrRG）》（聯邦法律公報第一部第1297頁）、並經1976年5月18日《第十五刑法修正法》（聯邦法律公報第一部第1213頁）第3條及第4條修正之第4條關於終止妊娠聯邦統計之規定的範圍內，與《基本法》第1條第1項結合第2條第2項第1句\ggfnArtOneOneTwoTwoOneZH{}不相容並無效。
\end{transpara}

\begin{sourcepara}[title={原文主文 I.6}]
6. Artikel 4 des Fünften Gesetzes zur Reform des Strafrechts in der Fassung des Artikels 15 Nummer 2 des Schwangeren- und Familienhilfegesetzes ist mit dem bundesstaatlichen Prinzip (Artikel 20 Absatz 1, Artikel 28 Absatz 1 des Grundgesetzes) unvereinbar und nichtig, soweit die Bestimmung die zuständigen obersten Landesbehörden verpflichtet; sie ist im übrigen nach Maßgabe der Urteilsgründe mit dem Grundgesetz vereinbar.
\end{sourcepara}
\begin{transpara}
6. 《第五刑法改革法》第4條於《孕婦及家庭扶助法》第15條第2款版本中，在該規定課予各邦主管最高機關義務之範圍內，與聯邦國原則（《基本法》第20條第1項\ggfnArtTwentyOneZH{}、第28條第1項\ggfnArtTwentyEightOneZH{}）不相容並無效；其餘部分依本判決理由與《基本法》相容。
\end{transpara}

\begin{sourcepara}[title={原文主文 I.7}]
7. Die Anträge im Verfahren 2 \abbrBvF{} 2/90, betreffend die verfassungsrechtliche Prüfung des § 218b Absatz 1 Satz 1 und Absatz 2 sowie des § 219 Absatz 1 Satz 1 des Strafgesetzbuches in der Fassung des Fünfzehnten Strafrechtsänderungsgesetzes vom 18. Mai 1976 (\abbrBundesgesetzbl{} I Seite 1213), sind erledigt.
\end{sourcepara}
\begin{transpara}
7. 程序 2 BvF 2/90 中，關於審查1976年5月18日《第十五刑法修正法》（聯邦法律公報第一部第1213頁）版本《刑法》第218b條第1項第1句及第2項、以及第219條第1項第1句是否合憲之聲請，已經終結。
\end{transpara}

\begin{sourcepara}[title={原文主文 II. 開頭}]
II. Gemäß § 35 des Gesetzes über das Bundesverfassungsgericht wird angeordnet:
\end{sourcepara}
\begin{transpara}
II. 依《聯邦憲法法院法》第35條，命令如下：
\end{transpara}

\begin{sourcepara}[title={原文主文 II.1}]
1. Das bisher nach Maßgabe des Urteils vom 4. August 1992 geltende Recht bleibt bis zum 15. Juni 1993 anwendbar. Für die Zeit danach bis zum Inkrafttreten einer gesetzlichen Neuregelung gelten in Ergänzung zu den Vorschriften des Schwangeren- und Familienhilfegesetzes, soweit diese nicht durch Nummer I. der Urteilsformel für nichtig erklärt worden sind, die Nummern 2 bis 9 dieser Anordnung.
\end{sourcepara}
\begin{transpara}
1. 迄今依1992年8月4日判決尺度適用之法律，至1993年6月15日止仍得適用。此後至新法律規定生效前，除《孕婦及家庭扶助法》規定中已由本判決主文第I項宣告無效者外，並補充適用本命令第2款至第9款。
\end{transpara}

\begin{sourcepara}[title={原文主文 II.2}]
2. § 218 des Strafgesetzbuches in der Fassung des Schwangeren- und Familienhilfegesetzes findet keine Anwendung, wenn die Schwangerschaft innerhalb von zwölf Wochen nach der Empfängnis durch einen Arzt abgebrochen wird, die schwangere Frau den Abbruch verlangt und dem Arzt durch eine Bescheinigung nachgewiesen hat, daß sie sich mindestens drei Tage vor dem Eingriff von einer anerkannten Beratungsstelle (vgl. Nummer 4 dieser Anordnung) hat beraten lassen. Das grundsätzliche Verbot des Schwangerschaftsabbruchs bleibt auch in diesen Fällen unberührt.
\end{sourcepara}
\begin{transpara}
2. 《孕婦及家庭扶助法》版本之《刑法》第218條於下列情形不適用：懷孕婦女於受孕後十二週內要求終止妊娠並由醫師施行，且以證明向醫師證明其至少於介入前三日曾由經承認之諮詢機構接受諮詢（參本命令第4款）。即使在此等情形中，終止妊娠之原則禁止仍不受影響。
\end{transpara}

\bikeepwithnext[6]
\begin{sourcepara}[title={原文主文 II.3 開頭}]
3.
\end{sourcepara}
\begin{transpara}
3.
\end{transpara}

\begin{sourcepara}[title={原文主文 II.3(1)}]
(1) Die Beratung dient dem Schutz des ungeborenen Lebens. Sie hat sich von dem Bemühen leiten zu lassen, die Frau zur Fortsetzung der Schwangerschaft zu ermutigen und ihr Perspektiven für ein Leben mit dem Kind zu eröffnen; sie soll ihr helfen, eine verantwortliche und gewissenhafte Entscheidung zu treffen. Dabei muß der Frau bewußt sein, daß das Ungeborene in jedem Stadium der Schwangerschaft auch ihr gegenüber ein eigenes Recht auf Leben hat und daß deshalb nach der Rechtsordnung ein Schwangerschaftsabbruch nur in Ausnahmesituationen in Betracht kommen kann, wenn der Frau durch das Austragen des Kindes eine Belastung erwächst, die -- vergleichbar den Fällen des § 218a Absatz 2 und 3 des Strafgesetzbuches in der Fassung des Schwangeren- und Familienhilfegesetzes -- so schwer und außergewöhnlich ist, daß sie die zumutbare Opfergrenze übersteigt.
\end{sourcepara}
\begin{transpara}
(1) 諮詢服務於未出生生命之保護。諮詢應以努力鼓勵婦女繼續懷孕，並為其開啟與子女共同生活之前景為指導；諮詢應協助婦女作成負責且憑良心之決定。在此，婦女必須意識到：未出生者在懷孕任何階段，即使相對於她，也享有自身生命權；因此，依法律秩序，終止妊娠只有在例外情境下始得考慮，也就是繼續懷孕會使婦女承受一項負擔，而該負擔如同《孕婦及家庭扶助法》版本《刑法》第218a條第2項及第3項之情形般嚴重且異常，以致超越可期待之犧牲界限。
\end{transpara}

\begin{sourcepara}[title={原文主文 II.3(2)}]
(2) Die Beratung bietet der schwangeren Frau Rat und Hilfe. Sie trägt dazu bei, die im Zusammenhang mit der Schwangerschaft bestehende Konfliktlage zu bewältigen und einer Notlage abzuhelfen. Hierzu umfaßt die Beratung
\end{sourcepara}
\begin{transpara}
(2) 諮詢為懷孕婦女提供建議與協助。諮詢有助於處理與懷孕相關之衝突處境，並排除困境。為此，諮詢包括：
\end{transpara}

\begin{sourcepara}[title={原文主文 II.3(2)a}]
a) das Eintreten in eine Konfliktberatung; dazu wird erwartet, daß die schwangere Frau der sie beratenden Person die Tatsachen mitteilt, deretwegen sie einen Abbruch der Schwangerschaft erwägt;
\end{sourcepara}
\begin{transpara}
a) 進入衝突諮詢；在此期待懷孕婦女向諮詢者說明其考慮終止妊娠所依據之事實；
\end{transpara}

\begin{sourcepara}[title={原文主文 II.3(2)b}]
b) jede nach Sachlage erforderliche medizinische, soziale und juristische Information, die Darlegung der Rechtsansprüche von Mutter und Kind und der möglichen praktischen Hilfen, insbesondere solcher, die die Fortsetzung der Schwangerschaft und die Lage von Mutter und Kind erleichtern;
\end{sourcepara}
\begin{transpara}
b) 依事況所需之一切醫療、社會與法律資訊，說明母親與子女之法律請求權，以及可能之實際協助，尤其是有助於繼續懷孕並改善母子處境之協助；
\end{transpara}

\begin{sourcepara}[title={原文主文 II.3(2)c}]
c) das Angebot, die schwangere Frau bei der Geltendmachung von Ansprüchen, bei der Wohnungssuche, bei der Suche nach einer Betreuungsmöglichkeit für das Kind und bei der Fortsetzung ihrer Ausbildung zu unterstützen, sowie das Angebot einer Nachbetreuung.
\end{sourcepara}
\begin{transpara}
c) 提供協助懷孕婦女主張請求權、尋找住居、尋找子女照顧可能性、繼續其教育訓練之協助，以及提供後續照護。
\end{transpara}

\begin{sourcepara}[title={原文主文 II.3 避孕資訊}]
Die Beratung unterrichtet auch über Möglichkeiten, ungewollte Schwangerschaften zu vermeiden.
\end{sourcepara}
\begin{transpara}
諮詢亦告知避免非自願懷孕之可能方式。
\end{transpara}

\begin{sourcepara}[title={原文主文 II.3(3)}]
(3) Erforderlichenfalls sind ärztlich, psychologisch oder juristisch ausgebildete Fachkräfte oder andere Personen zu der Beratung hinzuzuziehen. Bei jeder Beratung ist zu prüfen, ob es angezeigt ist, im Einvernehmen mit der schwangeren Frau Dritte, insbesondere den Vater sowie nahe Angehörige beider Eltern des Ungeborenen hinzuzuziehen.
\end{sourcepara}
\begin{transpara}
(3) 必要時，應使受過醫學、心理學或法律訓練之專業人員或其他人員參與諮詢。每一次諮詢均應審查，在懷孕婦女同意下使第三人，尤其是父親及未出生者雙親之近親參與，是否適當。
\end{transpara}

\begin{sourcepara}[title={原文主文 II.3(4)}]
(4) Die schwangere Frau kann auf ihren Wunsch gegenüber der sie beratenden Person anonym bleiben.
\end{sourcepara}
\begin{transpara}
(4) 懷孕婦女得依其意願，相對於諮詢者保持匿名。
\end{transpara}

\begin{sourcepara}[title={原文主文 II.3(5)}]
(5) Ist es nach dem Inhalt des Beratungsgesprächs dem Ziel der Beratung (Absatz 1 [Satz 1]) dienlich, ist das Beratungsgespräch alsbald fortzusetzen. Sieht die beratende Person die Beratung als abgeschlossen an, hat die Beratungsstelle der Frau auf Antrag über die Tatsache, daß eine Beratung nach den Absätzen 1 bis 4 stattgefunden hat, eine auf ihren Namen lautende und mit dem Datum des letzten Beratungsgesprächs versehene Bescheinigung auszustellen.
\end{sourcepara}
\begin{transpara}
(5) 若依諮詢談話內容，為達成諮詢目的（第1項〔第1句〕）有助益，諮詢談話應儘速繼續。若諮詢者認為諮詢已完成，諮詢機構應依婦女聲請，就已依第1項至第4項進行諮詢之事實，出具以其姓名為名義並載有最後一次諮詢談話日期之證明。
\end{transpara}

\begin{sourcepara}[title={原文主文 II.3(6)}]
(6) Die beratende Person hat in einer Weise, die keine Rückschlüsse auf die Identität der Beratenen erlaubt, in einem Protokoll das Alter, den Familienstand und die Staatsangehörigkeit der Beratenen, die Zahl ihrer Schwangerschaften, ihrer Kinder und früherer Schwangerschaftsabbrüche festzuhalten. Sie hat ferner die für den Abbruch genannten wesentlichen Gründe, die Dauer des Beratungsgesprächs und gegebenenfalls die zu ihm hinzugezogenen weiteren Personen zu vermerken. Das Protokoll muß auch ausweisen, welche Informationen der Schwangeren vermittelt und welche Hilfen ihr angeboten worden sind.
\end{sourcepara}
\begin{transpara}
(6) 諮詢者應以不能推知受諮詢者身分之方式，在紀錄中記載受諮詢者之年齡、婚姻狀態與國籍，其懷孕次數、子女人數及既往終止妊娠次數。諮詢者並應記載終止妊娠所提出之主要理由、諮詢談話期間，以及必要時參與談話之其他人員。該紀錄亦必須顯示，已向懷孕婦女傳達何等資訊，並向其提供何等協助。
\end{transpara}

\bikeepwithnext[6]
\begin{sourcepara}[title={原文主文 II.4 開頭}]
4.
\end{sourcepara}
\begin{transpara}
4.
\end{transpara}

\begin{sourcepara}[title={原文主文 II.4(1)}]
(1) Stellen, die eine Beratung nach Nummer 3 vornehmen, bedürfen -- unabhängig von einer Anerkennung nach Artikel 1 § 3 des Schwangeren- und Familienhilfegesetzes -- besonderer staatlicher Anerkennung. Als Beratungsstellen können auch Einrichtungen freier Träger und Ärzte anerkannt werden.
\end{sourcepara}
\begin{transpara}
(1) 實施第3款諮詢之機構，不論是否已依《孕婦及家庭扶助法》第1條第3條獲承認，均須有特別之國家承認。自由團體之設施及醫師，亦得被承認為諮詢機構。
\end{transpara}

\begin{sourcepara}[title={原文主文 II.4(2)}]
(2) Beratungsstellen dürfen mit Einrichtungen, in denen Schwangerschaftsabbrüche vorgenommen werden, nicht derart organisatorisch oder durch wirtschaftliche Interessen verbunden sein, daß hiernach ein materielles Interesse der Beratungseinrichtung an der Durchführung von Schwangerschaftsabbrüchen nicht auszuschließen ist. Der Arzt, der den Schwangerschaftsabbruch vornimmt, ist als Berater ausgeschlossen; er darf auch nicht der Beratungsstelle angehören, die die Beratung durchgeführt hat.
\end{sourcepara}
\begin{transpara}
(2) 諮詢機構不得與實施終止妊娠之設施，在組織上或經濟利益上具有如下連結：依此不能排除諮詢機構對實施終止妊娠具有物質利益。實施終止妊娠之醫師不得擔任諮詢者；其亦不得隸屬於實施諮詢之諮詢機構。
\end{transpara}

\begin{sourcepara}[title={原文主文 II.4(3)}]
(3) Als Beratungsstelle kann nur anerkannt werden, wer für eine Beratung nach Maßgabe der Nummer 3 Gewähr bietet, über für eine solche Beratung in persönlicher und fachlicher Hinsicht qualifiziertes und der Zahl nach ausreichendes Personal verfügt und mit allen Stellen zusammenarbeitet, die öffentliche und private Hilfen für Mutter und Kind gewähren. Die Beratungsstellen sind verpflichtet, die ihrer Beratungstätigkeit zugrundeliegenden Maßstäbe und die dabei gesammelten Erfahrungen jährlich schriftlich niederzulegen.
\end{sourcepara}
\begin{transpara}
(3) 只有能保證依第3款實施諮詢、具備個人與專業上足以實施此種諮詢且數量充分之人員，並與所有提供母親與兒童公私協助之機構合作者，始得被承認為諮詢機構。諮詢機構有義務每年以書面記錄其諮詢活動所依據之尺度及在此過程中累積之經驗。
\end{transpara}

\begin{sourcepara}[title={原文主文 II.4(4)}]
(4) Die Anerkennung darf nur mit der Maßgabe erteilt werden, daß sie nach einer gesetzlich zu bestimmenden Frist jeweils der Bestätigung durch die zuständige Behörde bedarf.
\end{sourcepara}
\begin{transpara}
(4) 承認僅得附帶下列條件授予：於法律所定期間屆滿後，該承認均須經主管機關確認。
\end{transpara}

\begin{sourcepara}[title={原文主文 II.4(5)}]
(5) Die Länder stellen ein ausreichendes Angebot wohnortnaher Beratungsstellen sicher.
\end{sourcepara}
\begin{transpara}
(5) 各邦確保提供足夠且接近居住地之諮詢機構。
\end{transpara}

\begin{sourcepara}[title={原文主文 II.5}]
5. Dem Arzt, von dem die Frau den Abbruch der Schwangerschaft verlangt, obliegen die sich aus den Urteilsgründen ergebenden Pflichten (D.V.1. und 2.).
\end{sourcepara}
\begin{transpara}
5. 婦女向其要求終止妊娠之醫師，負有本判決理由所生之義務（D.V.1. 及 2.）。
\end{transpara}

\begin{sourcepara}[title={原文主文 II.6}]
6. Das in Nummer 4 vorgesehene Anerkennungsverfahren ist auch für bestehende Beratungsstellen durchzuführen. Bis zu dessen Abschluß, längstens bis zum 31. Dezember 1994, sind sie befugt, gemäß Nummer 3 dieser Anordnung zu beraten.
\end{sourcepara}
\begin{transpara}
6. 第4款所規定之承認程序，亦應對既有諮詢機構實施。於該程序完成以前，最遲至1994年12月31日止，該等機構有權依本命令第3款提供諮詢。
\end{transpara}

\begin{sourcepara}[title={原文主文 II.7}]
7. Die Pflicht zur Führung einer Bundesstatistik und die Meldepflicht nach Artikel 4 des Fünften Gesetzes zur Reform des Strafrechts (5. \abbrStrRG{}) vom 18. Juni 1974 (\abbrBundesgesetzbl{} I Seite 1297), geändert durch Artikel 3 und Artikel 4 des Fünfzehnten Strafrechtsänderungsgesetzes vom 18. Mai 1976 (\abbrBundesgesetzbl{} I Seite 1213), gelten auch in dem in Artikel 3 des Einigungsvertrages genannten Gebiet.
\end{sourcepara}
\begin{transpara}
7. 1974年6月18日《第五刑法改革法（5. StrRG）》（聯邦法律公報第一部第1297頁）第4條所定之聯邦統計製作義務及申報義務，經1976年5月18日《第十五刑法修正法》（聯邦法律公報第一部第1213頁）第3條及第4條修正者，亦適用於《統一條約》第3條所稱地區。
\end{transpara}

\begin{sourcepara}[title={原文主文 II.8}]
8. Die Regelung des § 37a des Bundessozialhilfegesetzes findet auch Anwendung bei Abbrüchen der Schwangerschaft nach Nummer 2 dieser Anordnung.
\end{sourcepara}
\begin{transpara}
8. 《聯邦社會扶助法》第37a條之規定，亦適用於依本命令第2款所為之終止妊娠。
\end{transpara}

\begin{sourcepara}[title={原文主文 II.9}]
9. Bis zu einer Entscheidung des Gesetzgebers über eine etwaige Einführung einer kriminologischen Indikation und deren Feststellung können Versicherte der gesetzlichen Krankenversicherung und nach Beihilfevorschriften Anspruchsberechtigte bei einem Abbruch der Schwangerschaft auf Antrag Leistungen erhalten, wenn die Voraussetzungen der Nummer 2 dieser Anordnung vorliegen und der zuständige Amtsarzt oder ein Vertrauensarzt der gesetzlichen Krankenkasse bescheinigt hat, daß nach seiner ärztlichen Erkenntnis an der Schwangeren eine rechtswidrige Tat nach den §§ 176 bis 179 des Strafgesetzbuches begangen worden ist und dringende Gründe für die Annahme sprechen, daß die Schwangerschaft auf der Tat beruht. Der Arzt kann mit Einwilligung der Frau eine Auskunft bei der Staatsanwaltschaft einholen und etwa vorhandene Ermittlungsakten einsehen; die hierbei gewonnenen Erkenntnisse unterliegen seiner ärztlichen Schweigepflicht.
\end{sourcepara}
\begin{transpara}
9. 在立法者就可能引入犯罪學指標事由（\textit{kriminologische Indikation}）及其確認作成決定以前，法定醫療保險之被保險人及依補助規定享有請求權者，於終止妊娠時，若符合本命令第2款要件，且主管公醫或法定醫療保險基金之信任醫師證明，依其醫學認識，懷孕婦女曾遭受《刑法》第176條至第179條之違法行為，且有迫切理由足以認為懷孕係基於該行為而生，得依聲請取得給付。醫師得經婦女同意，向檢察署取得資訊並閱覽既有偵查卷宗；由此取得之認識，受其醫師保密義務拘束。
\end{transpara}

\bisection{Gründe}{理由}
\startrandnumbers

\bisubsection{A.}{A.}

\begin{sourcepara}[title={原文理由 A. 開頭}]
\sourcern{1}Gegenstand der zu gemeinsamer Verhandlung und Entscheidung verbundenen Normenkontrollverfahren ist vor allem die Frage, ob verschiedene strafrechtliche, sozialversicherungsrechtliche und organisationsrechtliche Vorschriften über den Schwangerschaftsabbruch, die Teil der durch das Urteil des Bundesverfassungsgerichts vom 25. Februar 1975 (\abbrBVerfGE{} 39, 1 \abbrFF{}) veranlaßten Neuregelungen im 15. Strafrechtsänderungsgesetz und im Strafrechtsreform-Ergänzungsgesetz oder aber Teil des nach Herstellung der deutschen Einheit für Gesamtdeutschland neu erlassenen Schwangeren- und Familienhilfegesetzes sind, der verfassungsrechtlichen Pflicht des Staates genügen, das ungeborene menschliche Leben zu schützen.
\end{sourcepara}
\begin{transpara}
\transrn{1}合併言詞辯論及裁判之各規範審查程序，其標的首先在於下列問題：若干關於終止妊娠之刑法、社會保險法及組織法規定，是否符合國家保護未出生人類生命之憲法義務。這些規定一部分屬於德國聯邦憲法法院1975年2月25日判決（BVerfGE 39, 1 ff.）所促成、載於《第十五刑法修正法》及《刑法改革補充法》之新規定；另一部分則屬於德國統一後為全德重新制定之《孕婦及家庭扶助法》。
\end{transpara}

\bisubsubsection{I.}{I.}

\begin{sourcepara}[title={原文理由 A.I.1}]
\sourcern{2}1. Die Frage, ob und auf welche Weise das Problem des Schwangerschaftsabbruchs im Spannungsfeld zwischen dem Schutz des ungeborenen menschlichen Lebens und dem Selbstbestimmungsrecht der schwangeren Frau befriedigender als mit den Mitteln des Strafrechts gelöst werden könne, wird seit vielen Jahren kontrovers diskutiert. Der im Jahr 1974 unternommene Versuch des Gesetzgebers, die ursprünglich generelle Strafbarkeit des Schwangerschaftsabbruchs hauptsächlich durch eine Fristenlösung für die ersten zwölf Schwangerschaftswochen einzuschränken, wurde durch das Bundesverfassungsgericht verworfen. Dessen Erster Senat erklärte in seinem Urteil vom 25. Februar 1975 (\abbrBVerfGE{} 39, 1 \abbrFF{}) § 218a \abbrStGB{} in der Fassung des Fünften Gesetzes zur Reform des Strafrechts (5. \abbrStrRG{}) vom 18. Juni 1974 (\abbrBGBl{} I \abbrS{} 1297) für mit \abbrArt{} 2 \abbrAbs{} 2 Satz 1 in Verbindung mit \abbrArt{} 1 \abbrAbs{} 1 des Grundgesetzes insoweit unvereinbar und nichtig, als er den Schwangerschaftsabbruch auch dann von der Strafbarkeit ausnimmt, wenn keine Gründe vorliegen, die -- im Sinne der Entscheidungsgründe -- vor der Wertordnung des Grundgesetzes Bestand haben.
\end{sourcepara}
\begin{transpara}
\transrn{2}1. 終止妊娠問題處於未出生人類生命之保護與懷孕婦女自我決定權之緊張關係中；是否及如何能比刑法手段更令人滿意地解決此問題，多年來一直爭議不斷。1974年，立法者嘗試主要透過妊娠最初十二週之期限解決方案，限制原本對終止妊娠之一概可罰性；此一嘗試遭聯邦憲法法院否定。第一庭於1975年2月25日判決（BVerfGE 39, 1 ff.）中宣告，1974年6月18日《第五刑法改革法（5. StrRG）》（聯邦法律公報第一部第1297頁）版本之《刑法》第218a條，在其即使於不存在依判決理由所稱能在《基本法》價值秩序前成立之理由時，仍將終止妊娠排除於可罰性之外之範圍內，與《基本法》第2條第2項第1句結合第1條第1項\ggfnArtOneOneTwoTwoOneZH{}不相容並無效。
\end{transpara}

\begin{sourcepara}[title={原文理由 A.I.2}]
\sourcern{3}2. Der Deutsche Bundestag beschloß daraufhin das Fünfzehnte Strafrechtsänderungsgesetz (15. StÄG) vom 18. Mai 1976 (\abbrBGBl{} I \abbrS{} 1213); dadurch erhielten die Vorschriften der §§ 218 \abbrFF{} \abbrStGB{} (im folgenden: §§ 218 \abbrFF{} \abbrStGB{} \abbrAF{}) ihre bis jetzt geltende Fassung, die der Bekanntmachung des Strafgesetzbuches vom 10. März 1987 (\abbrBGBl{} I \abbrS{} 945, 1160) zugrunde liegt.
\end{sourcepara}
\begin{transpara}
\transrn{3}2. 德國聯邦議會其後通過1976年5月18日《第十五刑法修正法（15. StÄG）》（BGBl. I S. 1213）；《刑法》第218條以下規定（以下稱舊法第218條以下）因而取得迄今適用之版本，並成為1987年3月10日《刑法》公告版本（BGBl. I S. 945, 1160）之基礎。
\end{transpara}

\begin{sourcepara}[title={原文理由 A.I.2} Rn 4]
\sourcern{4}Nach diesem Gesetz wird grundsätzlich bestraft, wer eine Schwangerschaft nach Abschluß der Nidation abbricht (§ 218 \abbrAbs{} 1, \abbrAbs{} 3 Satz 1, § 219d \abbrStGB{} \abbrAF{}). Der Abbruch der Schwangerschaft ist jedoch innerhalb bestimmter Fristen nicht strafbar, wenn er durch einen Arzt vorgenommen wird, die Schwangere einwilligt und er nach ärztlicher Erkenntnis mit Rücksicht auf bestimmte schwerwiegende Notlagen der Schwangeren angezeigt ist (Indikationen zum Schwangerschaftsabbruch). Die maßgebende Vorschrift lautet:
\end{sourcepara}
\begin{transpara}
\transrn{4}依該法律，原則上，於著床（\textit{Nidation}）完成後終止妊娠者，應受處罰（舊法第218條第1項、第3項第1句、第219d條）。然而，若終止妊娠係於特定期限內由醫師實施、懷孕婦女同意，且依醫師認識，考量懷孕婦女特定重大困境，終止妊娠具有指標事由（終止妊娠之指標事由，\textit{Indikationen zum Schwangerschaftsabbruch}），則不罰。關鍵規定如下：
\end{transpara}

\begin{sourceboxpara}[title={原文理由 A.I.2 法條}]
\begin{lawboxpart}{first}{lawrn5}
\sourcern{5}\sourcelawtitle{Indikation zum Schwangerschaftsabbruch}
\end{lawboxpart}
\end{sourceboxpara}
\begin{transboxpara}
\begin{lawboxpart}{first}{lawrn5}
\transrn{5}\lawtitle{終止妊娠之指標事由}
\end{lawboxpart}
\end{transboxpara}
\begin{sourceboxpara}
\begin{lawboxpart}{middle}{lawrn6}
\sourcern{6}\sourceabsno{1}{Der Abbruch der Schwangerschaft durch einen Arzt ist nicht nach § 218 strafbar, wenn}
\end{lawboxpart}
\end{sourceboxpara}
\begin{transboxpara}
\begin{lawboxpart}{middle}{lawrn6}
\transrn{6}\absno{1}{由醫師所為之終止妊娠，於下列情形不依第218條處罰：}
\end{lawboxpart}
\end{transboxpara}
\begin{sourceboxpara}
\begin{lawboxpart}{middle}{lawrn7}
\sourcern{7}\sourcenrno{1}{die Schwangere einwilligt und}
\end{lawboxpart}
\end{sourceboxpara}
\begin{transboxpara}
\begin{lawboxpart}{middle}{lawrn7}
\transrn{7}\nrno{1}{懷孕婦女同意；且}
\end{lawboxpart}
\end{transboxpara}
\begin{sourceboxpara}
\begin{lawboxpart}{middle}{lawrn8}
\sourcern{8}\sourcenrno{2}{der Abbruch der Schwangerschaft unter Berücksichtigung der gegenwärtigen und zukünftigen Lebensverhältnisse der Schwangeren nach ärztlicher Erkenntnis angezeigt ist, um eine Gefahr für das Leben oder die Gefahr einer schwerwiegenden Beeinträchtigung des körperlichen oder seelischen Gesundheitszustandes der Schwangeren abzuwenden, und die Gefahr nicht auf eine andere für sie zumutbare Weise abgewendet werden kann.}
\end{lawboxpart}
\end{sourceboxpara}
\begin{transboxpara}
\begin{lawboxpart}{middle}{lawrn8}
\transrn{8}\nrno{2}{依醫師認識，考量懷孕婦女現在及未來生活關係，終止妊娠係為避免其生命危險，或避免其身體或心理健康狀態受到重大損害之危險而具有指標事由，且該危險不能以其他對其可期待之方式排除。}
\end{lawboxpart}
\end{transboxpara}
\begin{sourceboxpara}
\begin{lawboxpart}{middle}{lawrn9}
\sourcern{9}\sourceabsno{2}{Die Voraussetzungen des Absatzes 1 \abbrNr{} 2 gelten auch als erfüllt, wenn nach ärztlicher Erkenntnis}
\end{lawboxpart}
\end{sourceboxpara}
\begin{transboxpara}
\begin{lawboxpart}{middle}{lawrn9}
\transrn{9}\absno{2}{若依醫師認識有下列情形，第1項第2款之要件亦視為成就：}
\end{lawboxpart}
\end{transboxpara}
\begin{sourceboxpara}
\begin{lawboxpart}{middle}{lawrn10}
\sourcern{10}\sourcenrno{1}{dringende Gründe für die Annahme sprechen, daß das Kind infolge einer Erbanlage oder schädlicher Einflüsse vor der Geburt an einer nicht behebbaren Schädigung seines Gesundheitszustandes leiden würde, die so schwer wiegt, daß von der Schwangeren die Fortsetzung der Schwangerschaft nicht verlangt werden kann,}
\end{lawboxpart}
\end{sourceboxpara}
\begin{transboxpara}
\begin{lawboxpart}{middle}{lawrn10}
\transrn{10}\nrno{1}{有迫切理由足以認為，子女因遺傳因素或出生前有害影響，將罹患無法補救之健康損害，其嚴重程度使懷孕婦女不能被要求繼續懷孕；}
\end{lawboxpart}
\end{transboxpara}
\begin{sourceboxpara}
\begin{lawboxpart}{middle}{lawrn11}
\sourcern{11}\sourcenrno{2}{an der Schwangeren eine rechtswidrige Tat nach den §§ 176 bis 179 begangen worden ist und dringende Gründe für die Annahme sprechen, daß die Schwangerschaft auf der Tat beruht, oder}
\end{lawboxpart}
\end{sourceboxpara}
\begin{transboxpara}
\begin{lawboxpart}{middle}{lawrn11}
\transrn{11}\nrno{2}{懷孕婦女曾遭受《刑法》第176條至第179條之違法行為，且有迫切理由足以認為懷孕係基於該行為而生；或}
\end{lawboxpart}
\end{transboxpara}
\begin{sourceboxpara}
\begin{lawboxpart}{middle}{lawrn12}
\sourcern{12}\sourcenrno{3}{der Abbruch der Schwangerschaft sonst angezeigt ist, um von der Schwangeren die Gefahr einer Notlage abzuwenden, die}
\end{lawboxpart}
\end{sourceboxpara}
\begin{transboxpara}
\begin{lawboxpart}{middle}{lawrn12}
\transrn{12}\nrno{3}{終止妊娠另係為避免懷孕婦女陷入困境之危險而具有指標事由，該困境}
\end{lawboxpart}
\end{transboxpara}
\begin{sourceboxpara}
\begin{lawboxpart}{middle}{lawrn13}
\sourcern{13}\sourceletterno{a}{so schwer wiegt, daß von der Schwangeren die Fortsetzung der Schwangerschaft nicht verlangt werden kann, und}
\end{lawboxpart}
\end{sourceboxpara}
\begin{transboxpara}
\begin{lawboxpart}{middle}{lawrn13}
\transrn{13}\letterno{a}{其嚴重程度使懷孕婦女不能被要求繼續懷孕，且}
\end{lawboxpart}
\end{transboxpara}
\begin{sourceboxpara}
\begin{lawboxpart}{middle}{lawrn14}
\sourcern{14}\sourceletterno{b}{nicht auf eine andere für die Schwangere zumutbare Weise abgewendet werden kann.}
\end{lawboxpart}
\end{sourceboxpara}
\begin{transboxpara}
\begin{lawboxpart}{middle}{lawrn14}
\transrn{14}\letterno{b}{不能以其他對懷孕婦女可期待之方式排除。}
\end{lawboxpart}
\end{transboxpara}
\begin{sourceboxpara}
\begin{lawboxpart}{last}{lawrn15}
\sourcern{15}\sourceabsno{3}{In den Fällen des Absatzes 2 \abbrNr{} 1 dürfen seit der Empfängnis nicht mehr als zweiundzwanzig Wochen, in den Fällen des Absatzes 2 \abbrNr{} 2 und 3 nicht mehr als zwölf Wochen verstrichen sein.}
\end{lawboxpart}
\end{sourceboxpara}
\begin{transboxpara}
\begin{lawboxpart}{last}{lawrn15}
\transrn{15}\absno{3}{第2項第1款之情形，自受胎以來不得超過二十二週；第2項第2款及第3款之情形，自受胎以來不得超過十二週。}
\end{lawboxpart}
\end{transboxpara}

\begin{sourcepara}[title={原文理由 A.I.2 續}]
\sourcern{16}Darüber hinaus ist die Schwangere auch beim Fehlen einer Indikation nicht strafbar, wenn ein Arzt die Schwangerschaft innerhalb von 22 Wochen seit der Empfängnis nach einer Beratung gemäß § 218b \abbrStGB{} \abbrAF{} abbricht (§ 218 \abbrAbs{} 3 Satz 2 \abbrStGB{} \abbrAF{}). Liegen diese Voraussetzungen nicht vor, kann das Gericht bei der Frau gleichwohl von Strafe absehen, wenn sie sich zur Zeit des Eingriffs in besonderer Bedrängnis befunden hat (§ 218 \abbrAbs{} 3 Satz 3 \abbrStGB{} \abbrAF{}). In diesen Fällen trifft die Strafe nur den von der Frau herangezogenen Arzt. Für diesen gelten auch sonst strengere Maßstäbe: Wer nämlich eine Schwangerschaft abbricht, ohne daß sich die Schwangere mindestens drei Tage vor dem Eingriff über die zur Verfügung stehenden öffentlichen und privaten Hilfen von einem Berater (vgl. § 218b \abbrAbs{} 2 \abbrStGB{} \abbrAF{}) und über die ärztlich bedeutsamen Gesichtspunkte von einem Arzt hat beraten lassen, macht sich auch bei Vorliegen einer Indikation wegen Abbruchs der Schwangerschaft ohne Beratung der Schwangeren strafbar (§ 218b \abbrAbs{} 1 \abbrStGB{} \abbrAF{}). Ferner macht sich trotz des Vorliegens einer Indikation strafbar, wer eine Schwangerschaft abbricht, ohne daß sich ein anderer Arzt über das Vorliegen einer Indikation schriftlich geäußert hat (§ 219 \abbrStGB{} \abbrAF{}). Die Beratung der Schwangeren über soziale Hilfen und die Beratung über die ärztlich bedeutsamen Gesichtspunkte dürfen auch durch den Arzt erfolgen, der sich zur Indikation äußert; die Beratung über die ärztlich bedeutsamen Gesichtspunkte darf auch der abbrechende Arzt vornehmen. Die Schwangere wird nicht nach den §§ 218b, 219 \abbrStGB{} \abbrAF{} bestraft.
\end{sourcepara}
\begin{transpara}
\transrn{16}此外，即使欠缺指標事由，若醫師於受孕後二十二週內，依舊法第218b條所定諮詢後終止妊娠，懷孕婦女亦不受處罰（舊法第218條第3項第2句）。若上述要件不存在，但婦女於介入時處於特別困迫，法院仍得對其免除刑罰（舊法第218條第3項第3句）。在此等情形中，刑罰僅及於婦女所求助之醫師。對該醫師，其他方面亦適用較嚴格尺度：若終止妊娠時，懷孕婦女未於介入至少三日前，就可利用之公私協助由諮詢者接受諮詢（參舊法第218b條第2項），並就醫學上重要面向由醫師接受諮詢，則即使存在指標事由，行為人仍因未經懷孕婦女諮詢而終止妊娠構成犯罪（舊法第218b條第1項）。此外，即使存在指標事由，若終止妊娠時未有另一名醫師就指標事由存在作成書面表示，亦構成犯罪（舊法第219條）。關於社會協助之懷孕婦女諮詢，及關於醫學上重要面向之諮詢，亦得由表示指標事由意見之醫師為之；關於醫學上重要面向之諮詢，亦得由實施終止妊娠之醫師為之。懷孕婦女不依舊法第218b條、第219條處罰。
\end{transpara}

\begin{sourcepara}[title={原文理由 A.I.3}]
\sourcern{17}3. Das Gesetz über ergänzende Maßnahmen zum Fünften Strafrechtsreformgesetz (Strafrechtsreform-Ergänzungsgesetz -- \abbrStREG{}) vom 28. August 1975 (\abbrBGBl{} I \abbrS{} 2289) verfolgt das Ziel, die Reformbestrebungen des 5. Strafrechtsreformgesetzes durch flankierende sozialpolitische Maßnahmen zu unterstützen (vgl. \abbrBTDrucks{} 7/376 \abbrS{} 1).
\end{sourcepara}
\begin{transpara}
\transrn{17}3. 1975年8月28日《關於第五刑法改革法補充措施之法律（刑法改革補充法，StREG）》（BGBl. I S. 2289），其目的在於以附隨之社會政策措施支持《第五刑法改革法》之改革努力（參 BTDrucks. 7/376 S. 1）。
\end{transpara}

\begin{sourcepara}[title={原文理由 A.I.3} Rn 18]
\sourcern{18}Es beruht auf einem Entwurf der Fraktionen der SPD und FDP (\abbrBTDrucks{} 7/376), der ein Gesetzesvorhaben der Bundesregierung aus dem Jahre 1972 (\abbrBRDrucks{} 104/72) wieder aufgenommen hat. Der Entwurf sieht unter anderem vor, daß Versicherte bei einem Schwangerschaftsabbruch durch einen Arzt Anspruch auf Leistungen der gesetzlichen Krankenversicherung haben.
\end{sourcepara}
\begin{transpara}
\transrn{18}該法以 SPD 與 FDP 黨團草案（BTDrucks. 7/376）為基礎；此草案重新承接聯邦政府1972年之立法計畫（BRDrucks. 104/72）。草案尤其規定，被保險人於由醫師實施終止妊娠時，享有法定醫療保險給付請求權。
\end{transpara}

\begin{sourcepara}[title={原文理由 A.I.3} Rn 19]
\sourcern{19}Nach der Gesetzesbegründung soll der Entwurf die Reformbestrebungen des § 218 \abbrStGB{} ergänzen, der die Notlagen, in denen sich eine Schwangere befinden könne, in zu geringem Maße berücksichtige und damit der Problematik illegaler Abtreibungen nicht gerecht werde. Mit den vorgesehenen Maßnahmen solle auch verhindert werden, daß es zum illegalen Schwangerschaftsabbruch komme; darüber hinaus solle sichergestellt werden, daß Schwangere "in gesetzlich straffrei gestellten Fällen" wegen ihrer wirtschaftlichen Verhältnisse nicht benachteiligt würden. Weiter heißt es in der Begründung (\abbrBTDrucks{} \abbrAAO{} \abbrS{} 5 \abbrF{}):
\end{sourcepara}
\begin{transpara}
\transrn{19}依法律理由，草案旨在補充《刑法》第218條之改革努力；該條對懷孕婦女可能處於之困境考量過少，因而未能適當處理非法終止妊娠問題。所規定措施亦應防止發生非法終止妊娠；此外，並應確保懷孕婦女於「法律上置於不罰之案件」中，不因其經濟關係而受不利對待。理由並進一步記載（BTDrucks. a.a.O. S. 5 f.）：
\end{transpara}

\begin{sourceboxpara}[title={原文理由 A.I.3 引文}]
\begin{lawboxpart}{first}{lawrn20}
\sourcern{20}"Die Beratung und die Behandlung bei einem Schwangerschaftsabbruch sind als Leistungen der gesetzlichen Krankenversicherung ausgestaltet, weil auf diese Weise eine sachgerechte Durchführung gesichert wird.
\end{lawboxpart}
\end{sourceboxpara}
\begin{transboxpara}
\begin{lawboxpart}{first}{lawrn20}
\transrn{20}終止妊娠之諮詢與治療被設計為法定醫療保險給付，因為以此方式可確保其適當實施。
\end{lawboxpart}
\end{transboxpara}
\begin{sourceboxpara}
\begin{lawboxpart}{last}{lawrn21}
\sourcern{21}Die Durchführung dieser Aufgaben liegt auch wesentlich im Interesse der Allgemeinheit. Die Versichertengemeinschaft soll daher nicht allein die Kosten tragen. Aus diesem Grunde ist eine finanzielle Beteiligung des Bundes vorgesehen."
\end{lawboxpart}
\end{sourceboxpara}
\begin{transboxpara}
\begin{lawboxpart}{last}{lawrn21}
\transrn{21}執行此等任務亦具有重要公益。因此，不應由被保險人共同體單獨承擔費用。基於此，規定聯邦財政參與。
\end{lawboxpart}
\end{transboxpara}

\begin{sourcepara}[title={原文理由 A.I.3 續}]
\sourcern{22}Auf Empfehlung des Ausschusses für Arbeit und Sozialordnung nahm der Deutsche Bundestag den Gesetzentwurf mit einigen Änderungen an (Deutscher Bundestag, 7. \abbrWP{}, 88. Sitzung vom 21. März 1974, \abbrStenBer{} \abbrS{} 5763 \abbrFF{}). Diese bestanden im wesentlichen in der Ausdehnung der Leistungen der Krankenversicherung und der Sozialhilfe auch auf die straffrei gestellte Sterilisation durch einen Arzt sowie in der ausdrücklichen Zubilligung des in den Beratungen umstrittenen Anspruchs auf Krankengeld und Lohnfortzahlung bei Arbeitsunfähigkeit infolge Sterilisation oder Schwangerschaftsabbruchs (vgl. \abbrBTDrucks{} 7/1753, \abbrS{} 5 bis 11). Nachdem der Bundesrat seine Zustimmung verweigert hatte, schlug der Vermittlungsausschuß eine Gesetzesfassung vor, die sicherstellen sollte, daß nach einer gesetzlichen Neuregelung der Strafbarkeit des Schwangerschaftsabbruchs im Anschluß an das Urteil des Bundesverfassungsgerichts vom 25. Februar 1975 Leistungen nur bei allen vom Gesetzgeber anerkannten Fällen eines "nicht rechtswidrigen" Abbruchs der Schwangerschaft durch einen Arzt zu gewähren seien (\abbrBTDrucks{} 7/3778 \abbrS{} 2 \abbrF{}). Diese Fassung ist Gesetz geworden.
\end{sourcepara}
\begin{transpara}
\transrn{22}依勞動與社會秩序委員會建議，德國聯邦議會作成若干修正後通過該法律草案（Deutscher Bundestag, 7. WP., 88. Sitzung vom 21. März 1974, StenBer. S. 5763 ff.）。此等修正主要在於，將醫療保險與社會扶助給付擴張至由醫師實施且置於不罰之絕育，並明文承認審議中具爭議之疾病津貼請求權，以及因絕育或終止妊娠而不能工作時之工資續付請求權（參 BTDrucks. 7/1753, S. 5 bis 11）。聯邦參議院拒絕同意後，調解委員會提出一項法律版本，其目的在確保：接續聯邦憲法法院1975年2月25日判決而對終止妊娠可罰性作成法律新規定後，給付僅於立法者所承認之所有由醫師實施「不具違法性」終止妊娠案件中提供（BTDrucks. 7/3778 S. 2 f.）。此一版本成為法律。
\end{transpara}

\begin{sourcepara}[title={原文理由 A.I.3 續} Rn 23]
\sourcern{23}Das Gesetz hat -- soweit dies hier von Interesse ist -- folgenden Wortlaut:
\end{sourcepara}
\begin{transpara}
\transrn{23}該法於本案相關範圍內，文字如下：
\end{transpara}

\begin{sourceboxpara}[title={原文理由 A.I.3 法條}]
\begin{lawboxpart}{first}{lawrn24}
\sourcern{24}\sourcelawtitle{IIIa. Sonstige Hilfen}
\end{lawboxpart}
\end{sourceboxpara}
\begin{transboxpara}
\begin{lawboxpart}{first}{lawrn24}
\transrn{24}\lawtitle{IIIa. 其他協助}
\end{lawboxpart}
\end{transboxpara}
\begin{sourceboxpara}
\begin{lawboxpart}{middle}{lawrn25}
\sourcern{25}\sourcelawtitle{§ 200e}
\end{lawboxpart}
\end{sourceboxpara}
\begin{transboxpara}
\begin{lawboxpart}{middle}{lawrn25}
\transrn{25}\lawtitle{第200e條}
\end{lawboxpart}
\end{transboxpara}
\begin{sourceboxpara}
\begin{lawboxpart}{middle}{lawrn26}
\sourcern{26}\sourcequotepara{Versicherte haben Anspruch auf ärztliche Beratung über Fragen der Empfängnisregelung; zur ärztlichen Beratung gehören auch die erforderliche Untersuchung und die Verordnung von empfängnisregelnden Mitteln.}
\end{lawboxpart}
\end{sourceboxpara}
\begin{transboxpara}
\begin{lawboxpart}{middle}{lawrn26}
\transrn{26}被保險人就受胎調節問題享有醫師諮詢請求權；醫師諮詢亦包括必要檢查及受胎調節手段之處方。
\end{lawboxpart}
\end{transboxpara}
\begin{sourceboxpara}
\begin{lawboxpart}{middle}{lawrn27}
\sourcern{27}\sourcelawtitle{§ 200f}
\end{lawboxpart}
\end{sourceboxpara}
\begin{transboxpara}
\begin{lawboxpart}{middle}{lawrn27}
\transrn{27}\lawtitle{第200f條}
\end{lawboxpart}
\end{transboxpara}
\begin{sourceboxpara}
\begin{lawboxpart}{middle}{lawrn28}
\sourcern{28}\sourcequotepara{Versicherte haben Anspruch auf Leistungen bei einer nicht rechtswidrigen Sterilisation und bei einem nicht rechtswidrigen Abbruch der Schwangerschaft durch einen Arzt. Es werden ärztliche Beratung über die Erhaltung und den Abbruch der Schwangerschaft, ärztliche Untersuchung und Begutachtung zur Feststellung der Voraussetzungen für eine nicht rechtswidrige Sterilisation oder für einen nicht rechtswidrigen Schwangerschaftsabbruch, ärztliche Behandlung, Versorgung mit Arznei-, Verband- und Heilmitteln sowie Krankenhauspflege gewährt. Anspruch auf Krankengeld besteht, wenn Versicherte wegen einer nicht rechtswidrigen Sterilisation oder wegen eines nicht rechtswidrigen Abbruchs der Schwangerschaft durch einen Arzt arbeitsunfähig werden, es sei denn, es besteht Anspruch nach § 182 \abbrAbs{} 1 \abbrNr{} 2.}
\end{lawboxpart}
\end{sourceboxpara}
\begin{transboxpara}
\begin{lawboxpart}{middle}{lawrn28}
\transrn{28}被保險人就不具違法性之絕育及由醫師實施之不具違法性之終止妊娠享有給付請求權。給付包括關於維持懷孕及終止妊娠之醫師諮詢、為確認不具違法性之絕育或不具違法性之終止妊娠要件之醫師檢查與鑑定、醫療、藥品、敷料與療養用品之供給，以及住院照護。若被保險人因不具違法性之絕育或因由醫師實施之不具違法性之終止妊娠而不能工作，享有疾病津貼請求權；但依第182條第1項第2款已有請求權者，不在此限。
\end{lawboxpart}
\end{transboxpara}
\begin{sourceboxpara}
\begin{lawboxpart}{middle}{lawrn29}
\sourcern{29}\sourcelawtitle{§ 200g}
\end{lawboxpart}
\end{sourceboxpara}
\begin{transboxpara}
\begin{lawboxpart}{middle}{lawrn29}
\transrn{29}\lawtitle{第200g條}
\end{lawboxpart}
\end{transboxpara}
\begin{sourceboxpara}
\begin{lawboxpart}{last}{lawrn30}
\sourcern{30}\sourcequotepara{Die für die Krankenhilfe geltenden Vorschriften gelten für die Leistungsgewährung nach den §§ 200e und 200f entsprechend, soweit nichts Abweichendes bestimmt ist. § 192 \abbrAbs{} 1 gilt nicht für die Gewährung von Krankengeld bei einer nicht rechtswidrigen Sterilisation und bei einem nicht rechtswidrigen Abbruch der Schwangerschaft durch einen Arzt."}
\end{lawboxpart}
\end{sourceboxpara}
\begin{transboxpara}
\begin{lawboxpart}{last}{lawrn30}
\transrn{30}關於醫療協助之規定，於第200e條及第200f條給付之提供準用之，但另有不同規定者不在此限。第192條第1項不適用於不具違法性之絕育及由醫師實施之不具違法性之終止妊娠之疾病津貼給付。
\end{lawboxpart}
\end{transboxpara}

\begin{sourcepara}[title={原文理由 A.I.4}]
\sourcern{31}4. Die Indikationenregelung, insbesondere der Tatbestand der allgemeinen Notlagenindikation, sowie die Krankenkassenfinanzierung von Schwangerschaftsabbrüchen blieb auch in der Folgezeit rechtlich und politisch umstritten. Die Bayerische Staatsregierung stellte im März 1990 beim Bundesverfassungsgericht einen Antrag auf verfassungsrechtliche Prüfung von Vorschriften über das Beratungs- und Indikationsfeststellungsverfahren sowie über Krankenversicherungsleistungen bei Schwangerschaftsabbrüchen aufgrund der allgemeinen Notlagenindikation; dieser Antrag (2 \abbrBvF{} 2/90) ist Gegenstand der vorliegenden Entscheidung.
\end{sourcepara}
\begin{transpara}
\transrn{31}4. 指標規則，尤其是一般困境指標事由（\textit{allgemeine Notlagenindikation}）之構成要件，以及醫療保險基金對終止妊娠之財務給付，在其後仍於法律及政治上持續受爭議。巴伐利亞邦政府於1990年3月向聯邦憲法法院提出聲請，請求憲法審查關於諮詢與指標事由確認程序之規定，以及基於一般困境指標事由之終止妊娠醫療保險給付規定；此項聲請（2 BvF 2/90）即為本判決標的。
\end{transpara}

\bisubsubsection{II.}{II.}

\begin{sourcepara}[title={原文理由 A.II}]
\sourcern{32}Die Herstellung der deutschen Einheit am 3. Oktober 1990 und die damit gestellte Aufgabe, das Recht in den vereinigten Teilen Deutschlands zu vereinheitlichen, gab Reformbestrebungen einen neuen Anstoß.
\end{sourcepara}
\begin{transpara}
\transrn{32}1990年10月3日德國統一之實現，以及由此提出之任務，即統一德國合併各部分之法律，為改革努力帶來新的推動。
\end{transpara}

\begin{sourcepara}[title={原文理由 A.II.1}]
\sourcern{33}1. Zunächst blieb es bei unterschiedlichen Regelungen über die Strafbarkeit des Schwangerschaftsabbruchs in den beiden Teilen Deutschlands. Aufgrund des Einigungsvertrages (EV) vom 31. August 1990 in Verbindung mit dem Einigungsvertragsgesetz vom 23. September 1990 (\abbrBGBl{} II \abbrS{} 885; vgl. dort Anlage II, Kapitel III, Sachgebiet C, Abschnitt I \abbrNr{} 1) ist in dem Beitrittsgebiet nach § 153 des Strafgesetzbuches der \abbrDDR{} vom 12. Januar 1968 in der Neufassung vom 14. Dezember 1988 (\abbrGBlDDR{} I 1989 \abbrS{} 33), geändert durch das 6. Strafrechtsänderungsgesetz vom 29. Juni 1990 (\abbrGBlDDR{} I \abbrS{} 526), strafbar, wer "entgegen den gesetzlichen Vorschriften" die Schwangerschaft einer Frau unterbricht. Ebenso wird bestraft, wer eine Frau dazu veranlaßt oder sie dabei unterstützt, ihre Schwangerschaft selbst zu unterbrechen oder eine ungesetzliche Schwangerschaftsunterbrechung vornehmen zu lassen. Die nach dem Einigungsvertrag (\abbrAAO{}, Anlage II, Kapitel III, Sachgebiet C, Abschnitt I \abbrNr{} 4 und 5) fortgeltenden Vorschriften des Gesetzes über die Unterbrechung der Schwangerschaft vom 9. März 1972 (\abbrGBlDDR{} I \abbrS{} 89) und der dazu ergangenen Durchführungsbestimmung vom selben Tage (\abbrGBlDDR{} II \abbrS{} 149) enthalten eine Fristenregelung. Nach § 1 \abbrAbs{} 2 des Gesetzes ist die Schwangere berechtigt, die Schwangerschaft innerhalb von zwölf Wochen nach deren Beginn durch einen ärztlichen Eingriff in einer geburtshilflich-gynäkologischen Einrichtung unterbrechen zu lassen. Nach § 2 darf eine Unterbrechung zu einem späteren Zeitpunkt nur vorgenommen werden, wenn zu erwarten ist, daß die Fortdauer der Schwangerschaft das Leben der Frau gefährdet, oder wenn andere schwerwiegende Gründe vorliegen; darüber entscheidet eine Fachärztekommission. Nach § 3 ist die Unterbrechung der Schwangerschaft grundsätzlich unzulässig, wenn sie zu schweren gesundheitsgefährdenden oder lebensbedrohenden Komplikationen führen kann (\abbrAbs{} 1) oder wenn seit der letzten Unterbrechung weniger als sechs Monate vergangen sind (\abbrAbs{} 2). § 4 \abbrAbs{} 1 stellt die Vorbereitung, Durchführung und Nachbehandlung des zulässigen Schwangerschaftsabbruchs arbeits- und versicherungsrechtlich dem Erkrankungsfall gleich.
\end{sourcepara}
\begin{transpara}
\transrn{33}1. 起初，德國兩部分關於終止妊娠可罰性之規定仍不相同。依1990年8月31日《統一條約》（\textit{Einigungsvertrag, EV}），連同1990年9月23日《統一條約法》（BGBl. II S. 885；參其附件二，第三章，事項領域 C，第一節第1號），在加入地區，依1968年1月12日《德意志民主共和國刑法》第153條，該法為1988年12月14日新版本（GBl. DDR I 1989 S. 33），並經1990年6月29日《第六刑法修正法》（GBl. DDR I S. 526）修正，凡「違反法律規定」對婦女施行終止妊娠者，應受處罰。促使婦女自行終止妊娠，或支持其自行終止妊娠或使其接受違法終止妊娠者，亦同受處罰。依《統一條約》（a.a.O., Anlage II, Kapitel III, Sachgebiet C, Abschnitt I Nr. 4 und 5）繼續適用之1972年3月9日《終止妊娠法》（GBl. DDR I S. 89）及同日制定之施行規定（GBl. DDR II S. 149），包含期限規定（\textit{Fristenregelung}）。依該法第1條第2項，懷孕婦女有權於受孕後十二週內，在產科婦科機構中由醫師介入終止妊娠。依第2條，較後時點之終止妊娠僅於可預期繼續懷孕將危及婦女生命，或存在其他重大理由時，始得實施；是否如此，由專科醫師委員會決定。依第3條，若終止妊娠可能導致嚴重危害健康或危及生命之併發症，原則上不得為之（第1項）；若距上次終止妊娠未滿六個月，亦同（第2項）。第4條第1項在勞動法與保險法上，將合法終止妊娠之準備、實施與後續治療等同於疾病事件。
\end{transpara}

\begin{sourcepara}[title={原文理由 A.II.2.a}]
\sourcern{34}2. \abbrArt{} 31 \abbrAbs{} 4 des Einigungsvertrages vom 31. August 1990 gibt dem gesamtdeutschen Gesetzgeber auf, spätestens bis zum 31. Dezember 1992 eine Regelung zu treffen, die den Schutz des vorgeburtlichen Lebens und die verfassungskonforme Bewältigung von Konfliktsituationen schwangerer Frauen besser gewährleistet, als dies in den beiden Teilen Deutschlands derzeit der Fall ist.
\end{sourcepara}
\begin{transpara}
\transrn{34}2. 1990年8月31日《統一條約》第31條第4項要求全德立法者，最遲於1992年12月31日前制定一項規制，使出生前生命之保護及懷孕婦女衝突情境之合憲處理，較當時德國兩部分之狀態更能獲得保障。
\end{transpara}

\begin{sourcepara}[title={原文理由 A.II.2.a} Rn 35]
\sourcern{35}a) Im Jahre 1991 brachten deshalb die Fraktion der FDP (\abbrBTDrucks{} 12/551), die Abgeordneten Christina Schenk u.a. und die Gruppe Bündnis 90/Die Grünen (\abbrBTDrucks{} 12/696), die Fraktion der SPD (\abbrBTDrucks{} 12/841), die Abgeordneten Petra Bläss u.a. und die Gruppe PDS/Linke Liste (\abbrBTDrucks{} 12/898), die Fraktion der CDU/CSU (\abbrBTDrucks{} 12/1178 [neu]) sowie die Abgeordneten Herbert Werner u.a. (\abbrBTDrucks{} 12/1179) Gesetzentwürfe ein, die eine für das gesamte Deutschland geltende Neuregelung des Schwangerschaftsabbruchs zum Gegenstand haben.
\end{sourcepara}
\begin{transpara}
\transrn{35}a) 因此，1991年 FDP 黨團（BTDrucks. 12/551）、Christina Schenk 等議員與「聯盟90／綠黨」小組（BTDrucks. 12/696）、SPD 黨團（BTDrucks. 12/841）、Petra Bläss 等議員與 PDS／左翼名單小組（BTDrucks. 12/898）、CDU/CSU 黨團（BTDrucks. 12/1178 [neu]），以及 Herbert Werner 等議員（BTDrucks. 12/1179），提出以適用於全德之終止妊娠新規制為標的之法律草案。
\end{transpara}

\begin{sourcepara}[title={原文理由 A.II.2.a} Rn 36]
\sourcern{36}Während der Gesetzesberatung kam der Gesetzentwurf der Abgeordneten Inge Wettig-Danielmeier, Uta Würfel u.a. hinzu (\abbrBTDrucks{} 12/2605, ersetzt durch die \abbrBTDrucks{} 12/2605 [neu]). Kernpunkt des strafrechtlichen Teils dieses später mit Änderungen angenommenen Entwurfes ist ein grundlegend umgestalteter § 218 \abbrStGB{} sowie eine neugestaltete Beratungsregelung (§ 219 \abbrStGB{}). Danach sollen Schwangerschaftsabbrüche, die durch einen Arzt binnen zwölf Wochen seit der Empfängnis mit Einwilligung der schwangeren Frau vorgenommen werden, vom Straftatbestand des § 218 \abbrStGB{} nicht mehr erfaßt werden, sofern die Frau sich mindestens drei Tage vor dem Eingriff durch eine anerkannte Beratungsstelle hat beraten lassen. Die bisherigen Tatbestände der kriminologischen Indikation und der allgemeinen Notlagenindikation sollen entfallen und nur noch die medizinische und die embryopathische Indikation als Rechtfertigungsgründe erhalten bleiben.
\end{sourcepara}
\begin{transpara}
\transrn{36}法律審議期間，Inge Wettig-Danielmeier、Uta Würfel 等議員之法律草案亦加入審議（BTDrucks. 12/2605，後由 BTDrucks. 12/2605 [neu] 取代）。此一後來經修正通過之草案，其刑法部分核心為經根本改造之《刑法》第218條，以及重新設計之諮詢規制（《刑法》第219條）。依其規定，由醫師於受孕後十二週內，經懷孕婦女同意所實施之終止妊娠，若婦女至少於介入前三日已由受承認之諮詢機構接受諮詢，即不再為《刑法》第218條之犯罪構成要件所涵蓋。既有之犯罪學指標事由（\textit{kriminologische Indikation}）及一般困境指標事由（\textit{allgemeine Notlagenindikation}）構成要件應予刪除，僅保留醫學指標事由（\textit{medizinische Indikation}）與胚胎病理指標事由（\textit{embryopathische Indikation}）作為阻卻違法事由。
\end{transpara}

\begin{sourcepara}[title={原文理由 A.II.2.a 理由}]
\sourcern{37}In der Begründung wird hervorgehoben, daß angesichts der Bedeutung des Rechtsguts des werdenden Lebens und dessen verfassungsrechtlicher Gewährleistung ein strafrechtlicher Schutz unverzichtbar sei. Die Erfahrungen mit der 1976 eingeführten Indikationenregelung hätten jedoch gezeigt, daß es nicht möglich sei, hinreichend konkrete, ärztlich und gerichtlich nachprüfbare Kriterien für das Vorliegen einer Notlage zu normieren, die einen Abbruch rechtfertigten. Letztlich könne nur die schwangere Frau selbst die Konfliktlage beurteilen, in der sie sich befinde. Es gelte deshalb eine Regelung zu finden, die sowohl dem hohen Wert des ungeborenen Lebens als auch der Eigenverantwortlichkeit der Frau Rechnung trage. Mit seinem Urteil vom 25. Februar 1975 habe das Bundesverfassungsgericht nicht jede Fristenlösung für verfassungswidrig erklärt. In welchem Maße das Strafrecht zum Schutz des ungeborenen Lebens eingesetzt werden müsse, hänge davon ab, ob andere Regelungen vorhanden seien, durch die ein tatsächlicher Schutz werdenden Lebens gewährleistet werde. Voraussetzung für die verfassungskonforme Ausgestaltung der im Entwurf vorgesehenen Änderungen des Strafrechts sei danach zum einen, daß der Staat in ausreichendem Maße sozialpolitische Mittel zur Verfügung stelle, um auf diese Weise das vorgeburtliche Leben zu schützen. Diesem Erfordernis dienten die vorgeschlagenen sozialpolitischen Maßnahmen. Zum anderen müsse sichergestellt werden, daß die Frau ihre verantwortliche Gewissensentscheidung über einen Schwangerschaftsabbruch nicht losgelöst von der durch die Verfassung vorgegebenen Grundentscheidung für den Schutz des vorgeburtlichen Lebens treffe. Dies werde verfahrensmäßig durch die verpflichtende Beratung gewährleistet, durch die der Frau Rat und Hilfe in ihrem Konflikt angeboten sowie ausreichende Informationen über staatliche Hilfen als Grundlage für eine gründliche Reflexion ihrer Situation zuteil würden. Hierbei werde die Überlegung umgesetzt, daß die Bereitschaft zu einer Entscheidung für das werdende Leben am größten sei, wenn die Frau sich nicht dem Urteil anderer Stellen unterwerfen müsse, sondern nach qualifizierter Beratung und sorgfältiger Überlegung die Entscheidung über die Fortsetzung der Schwangerschaft selbst treffen könne. Die Entscheidungsfreiheit der Frau lasse das werdende Leben nicht schutzlos. Auf diese Weise könne sich die Chance realisieren, daß die Frau -- ohne in der Beratung bevormundet zu werden -- die ihr angebotenen Hilfen in ihrer Konfliktlage annehme und sich für das Kind entscheide. Da der für ein solches Beratungsgespräch vorauszusetzende verantwortungsvolle Kontakt zwischen schwangerer Frau und Beraterin oder Berater nicht erzwingbar sei, würden der Frau keine Darlegungspflicht und kein Rechtfertigungszwang auferlegt. Auf ihren Wunsch erhalte sie jedoch individuelle Lösungsvorschläge zur Bewältigung ihrer Konfliktlage. Die Beratung solle eine vertrauensvolle Atmosphäre zwischen Beraterin oder Berater und schwangerer Frau herstellen, damit die schwangere Frau offen sei, andere Konfliktlösungen als den Schwangerschaftsabbruch zu erwägen.
\end{sourcepara}
\begin{transpara}
\transrn{37}理由中強調，鑑於形成中生命（\textit{werdendes Leben}）此一法益之重要性及其憲法保障，刑法保護不可或缺。然而，1976年引入指標規則（\textit{Indikationenregelung}）之經驗已顯示，不可能將困境存在之標準規範得足夠具體，並可由醫師及法院加以審查，以作為終止妊娠之正當化依據。終究，只有懷孕婦女本人能評估其所處之衝突狀態。因此，必須尋找一種規制，既顧及未出生生命之高度價值，亦顧及婦女之自我責任。聯邦憲法法院於1975年2月25日之判決，並未宣告任何期限規定（\textit{Fristenlösung}）均屬違憲。刑法為保護未出生生命必須運用至何種程度，取決於是否存在其他規制，足以保障形成中生命之實際保護。據此，草案所設刑法修正要合憲地形成，其前提一方面是國家須以充分程度提供社會政策手段，藉此保護出生前生命。草案所建議之社會政策措施即服務於此項要求。另一方面，必須確保婦女就終止妊娠所作之負責任良心決定，不脫離憲法所預設、支持保護出生前生命之基本決定。此點在程序上由強制諮詢加以保障；透過諮詢，婦女於其衝突中獲得建議與協助，並獲得關於國家協助之充分資訊，作為其徹底反思自身處境之基礎。此處所實現之考量是：若婦女無須服從其他機關之判斷，而能在合格諮詢與謹慎思考後，自行決定是否繼續懷孕，則其作成有利於形成中生命之決定的意願最大。婦女之決定自由並不使形成中生命陷於無保護。以此方式，婦女有可能在不於諮詢中被監護式支配之情形下，接受其衝突狀態中所提供之協助，並選擇孩子。由於此種諮詢談話所預設之懷孕婦女與諮詢者之間負責任接觸，不能以強制方式取得，因此不對婦女課以說明義務或正當化壓力。然而，依其願望，她可獲得克服其衝突狀態之個別解決建議。諮詢應在諮詢者與懷孕婦女之間建立信賴氣氛，使懷孕婦女願意考慮終止妊娠以外之其他衝突解決方式。
\end{transpara}

\begin{sourcepara}[title={原文理由 A.II.2.a 引文}]
\sourcern{38}Der Entwurf übernimmt ferner die Vorschriften der Reichsversicherungsordnung über Krankenkassenleistungen bei Schwangerschaftsabbrüchen als §§ 24a, 24b in das Fünfte Buch des Sozialgesetzbuchs. Dazu heißt es in der Begründung (vgl. \abbrBTDrucks{} 12/2605 [neu], \abbrS{} 20):
\end{sourcepara}
\begin{transpara}
\transrn{38}此外，草案將《帝國保險法》中關於終止妊娠時醫療保險基金給付之規定，作為第24a條、第24b條移入《社會法典》第五編。理由對此記載如下（參 BTDrucks. 12/2605 [neu], S. 20）：
\end{transpara}

\begin{sourcepara}[title={原文理由 A.II.2.a 引文 1}]
\begin{lawbox}
\sourcern{39}"§ 24b entspricht im wesentlichen dem bisherigen § 200 \abbrF{} \abbrRVO{}. ... Erfaßt werden auch Schwangerschaftsabbrüche, die nach Beratung innerhalb der ersten 12 Wochen nach der Empfängnis durchgeführt werden, da Artikel 11 in § 218 Absatz 5 \abbrStGB{} den Straftatbestand des Schwangerschaftsabbruchs ausschließt. Es ist daher gesichert, daß hinsichtlich der Kostenübernahme keine Änderung gegenüber der derzeitigen Rechtslage eintritt."
\end{lawbox}
\end{sourcepara}
\begin{transpara}
\begin{lawbox}
\transrn{39}第24b條基本上相當於既有之《帝國保險法》第200f條。……經諮詢後於受孕後最初十二週內實施之終止妊娠亦被涵蓋，因為第11條於《刑法》第218條第5項排除終止妊娠之犯罪構成要件。因此可確保，就費用承擔而言，相較於現行法律狀態不發生變更。
\end{lawbox}
\end{transpara}

\begin{sourcepara}[title={原文理由 A.II.2.a 引文 2 前文}]
\sourcern{40}Die im Entwurf ebenfalls enthaltene Änderung des \abbrArt{} 4 des 5. \abbrStrRG{} (Wegfall der Bundesstatistik über Schwangerschaftsabbrüche, Sicherstellung eines flächendeckenden Netzes von Abbrucheinrichtungen) wird wie folgt begründet (\abbrBTDrucks{} 12/2605 [neu], \abbrS{} 23):
\end{sourcepara}
\begin{transpara}
\transrn{40}草案中同樣包含之《第五刑法改革法》第4條修正（終止妊娠聯邦統計之取消、確保終止妊娠設施之全面網絡），其理由如下（BTDrucks. 12/2605 [neu], S. 23）：
\end{transpara}

\begin{sourcepara}[title={原文理由 A.II.2.a 引文 2}]
\begin{lawbox}
\sourcern{41}"Für den bisherigen Artikel 4 besteht kein Bedarf. Der neue Artikel 4 verpflichtet die Länder, ein ausreichendes Angebot an Einrichtungen zur Vornahme von Schwangerschaftsabbrüchen zur Verfügung zu stellen. Dies gilt sowohl für ambulante als auch für stationäre Einrichtungen. Damit wird ausgeschlossen, daß die Zulassung ambulanter Einrichtungen zum Schwangerschaftsabbruch generell verweigert wird."
\end{lawbox}
\end{sourcepara}
\begin{transpara}
\begin{lawbox}
\transrn{41}既有第4條已無需求。新第4條使各邦負有義務，提供足夠之終止妊娠設施供給。此適用於門診設施及住院設施。藉此排除一般性拒絕准許門診設施實施終止妊娠。
\end{lawbox}
\end{transpara}

\begin{sourcepara}[title={原文理由 A.II.2.b--c}]
\sourcern{42}b) Der Sonderausschuß "Schutz des ungeborenen Lebens" beriet die ersten sechs Gesetzentwürfe in 17 Sitzungen, den Gesetzentwurf in der \abbrBTDrucks{} 12/2605 (neu) in drei Sitzungen.
\end{sourcepara}
\begin{transpara}
\transrn{42}b) 「保護未出生生命」特別委員會就前六項法律草案進行十七次會議審議，並就 BTDrucks. 12/2605 (neu) 所載法律草案進行三次會議審議。
\end{transpara}

\begin{sourcepara}[title={原文理由 A.II.2.b--c} Rn 43]
\sourcern{43}Am 13., 14. und 15. November 1991 sowie am 4. und 6. Dezember 1991 fand vor dem Ausschuß eine öffentliche Anhörung zu den Fragen der Beratung, der Prävention und Sexualerziehung und zu den verfassungs-, straf- und arztrechtlichen Fragen statt (vgl. "Zur Sache, Themen parlamentarischer Beratung", herausgegeben vom Deutschen Bundestag, Band 1/92, \abbrS{} 9 bis 1027).
\end{sourcepara}
\begin{transpara}
\transrn{43}1991年11月13日、14日、15日，以及1991年12月4日、6日，委員會就諮詢、預防與性教育問題，以及憲法、刑法與醫師法問題舉行公開聽證（參德國聯邦議會出版，\textit{Zur Sache, Themen parlamentarischer Beratung}, Band 1/92, S. 9 bis 1027）。
\end{transpara}

\begin{sourcepara}[title={原文理由 A.II.2.b--c} Rn 44]
\sourcern{44}Bei der Überarbeitung der Gesetzentwürfe wich der Ausschuß von der allgemeinen Übung ab: Im Blick darauf, daß die Entwürfe in entscheidenden Punkten miteinander unvereinbare Regelungen enthielten, wurden die Einzelfragen zwar gemeinsam beraten, die Entscheidung über sie und damit über etwaige Änderungen sowie die Endfassung des jeweiligen Gesetzentwurfes ausschließlich seinen im Ausschuß vertretenen Befürwortern und Unterzeichnern überlassen. Eine Schlußabstimmung über die einzelnen Vorlagen fand nicht statt. Der Ausschuß kam einstimmig überein, die Entscheidung über die künftige Regelung der mit ungewollten Schwangerschaften zusammenhängenden Fragen von der Gesamtheit der Mitglieder des Bundestages ohne Vorgabe durch den Ausschuß treffen zu lassen. Dies sei als Empfehlung zu verstehen, der Bundestag möge die Gesetzentwürfe in der Fassung, die sie im Ausschuß erhalten hätten, in zweiter Beratung behandeln und darüber abstimmen (vgl. Empfehlung und Bericht des Sonderausschusses "Schutz des ungeborenen Lebens", \abbrBTDrucks{} 12/2875, \abbrS{} 111).
\end{sourcepara}
\begin{transpara}
\transrn{44}於修訂法律草案時，委員會偏離通常作法：鑑於各草案在關鍵點上包含彼此不相容之規制，雖然各個問題仍共同審議，但關於各問題之決定，並因而關於可能修正及各法律草案最終版本之決定，均專由委員會中支持並簽署該草案者處理。各提案未進行最終表決。委員會一致同意，關於未來如何規制與非意願懷孕相關問題之決定，應由聯邦議會全體成員在不受委員會預設指示之情形下作成。此應理解為一項建議：聯邦議會應以各法律草案在委員會中取得之版本，於第二讀會中審議並表決（參「保護未出生生命」特別委員會之建議與報告，BTDrucks. 12/2875, S. 111）。
\end{transpara}

\begin{sourcepara}[title={原文理由 A.II.2.b--c} Rn 45]
\sourcern{45}c) Bei der namentlichen Abstimmung in der zweiten Beratung des Deutschen Bundestages erhielt der Entwurf der Abgeordneten Inge Wettig-Danielmeier, Uta Würfel u.a. (\abbrBTDrucks{} 12/2605[neu]) in der Ausschußfassung (vgl. dazu Empfehlung und Bericht des Sonderausschusses "Schutz des ungeborenen Lebens", \abbrBTDrucks{} 12/2875, \abbrS{} 85 \abbrFF{}, insbesondere 99 \abbrFF{}) die Mehrheit (Deutscher Bundestag, 12. \abbrWP{}, 99. Sitzung vom 25. Juni 1992, \abbrStenBer{} \abbrS{} 8374). Bei der namentlichen Schlußabstimmung über diesen Entwurf in der dritten Beratung stimmten von 657 Abgeordneten 357 mit Ja und 284 mit Nein. 16 Abgeordnete enthielten sich der Stimme (Deutscher Bundestag, 12. \abbrWP{}, 99. Sitzung vom 25. Juni 1992, \abbrStenBer{} \abbrS{} 8377).
\end{sourcepara}
\begin{transpara}
\transrn{45}c) 在德國聯邦議會第二讀會之記名表決中，Inge Wettig-Danielmeier、Uta Würfel 等議員之草案（BTDrucks. 12/2605[neu]）於委員會版本中（就此參「保護未出生生命」特別委員會之建議與報告，BTDrucks. 12/2875, S. 85 ff., insbesondere 99 ff.）獲得多數（Deutscher Bundestag, 12. WP., 99. Sitzung vom 25. Juni 1992, StenBer. S. 8374）。於第三讀會就此草案進行之記名最終表決中，657名議員中357票贊成、284票反對，16名議員棄權（Deutscher Bundestag, 12. WP., 99. Sitzung vom 25. Juni 1992, StenBer. S. 8377）。
\end{transpara}

\begin{sourcepara}[title={原文理由 A.II.2.b--c} Rn 46]
\sourcern{46}Der Bundesrat stimmte dem Gesetzesbeschluß des Deutschen Bundestages gemäß \abbrArt{} 84 \abbrAbs{} 1 \abbrGG{}\ggfnArtEightyFourOneDE{} gegen die Stimmen Bayerns und bei Enthaltung Baden-Württembergs, Mecklenburg-Vorpommerns und Thüringens zu (Bundesrat, 645. Sitzung vom 10. Juli 1992, \abbrStenBer{} \abbrS{} 375). Der Bundesrat nahm ferner eine Entschließung des Landes Hessen an (vgl. \abbrBRDrucks{} 451/3/92), in der gefordert wird, die Kosten der sozialen Begleitmaßnahmen auf alle Ebenen angemessen zu verteilen, insbesondere im Wege der Erhöhung des Anteils der Länder an der Umsatzsteuer zu Lasten des Bundes.
\end{sourcepara}
\begin{transpara}
\transrn{46}聯邦參議院依《基本法》第84條第1項\ggfnArtEightyFourOneZH{}，在巴伐利亞投反對票，巴登符騰堡、梅克倫堡-前波美拉尼亞及圖林根棄權之情形下，同意德國聯邦議會之法律決議（Bundesrat, 645. Sitzung vom 10. Juli 1992, StenBer. S. 375）。此外，聯邦參議院通過黑森邦之一項決議（參 BRDrucks. 451/3/92），其中要求將社會伴隨措施之費用適當分配於各層級，尤其是透過提高各邦於營業稅中之份額並由聯邦負擔之方式。
\end{transpara}

\begin{sourcepara}[title={原文理由 A.II.3.a}]
\sourcern{47}3. Dieses Gesetz vom 27. Juli 1992 (\abbrBGBl{} I \abbrS{} 1398) trägt die Überschrift "Gesetz zum Schutz des vorgeburtlichen/werdenden Lebens, zur Förderung einer kinderfreundlicheren Gesellschaft, für Hilfen im Schwangerschaftskonflikt und zur Regelung des Schwangerschaftsabbruchs (Schwangeren- und Familienhilfegesetz)" -- im folgenden: \abbrSFHG{} --; es hat im wesentlichen folgenden Inhalt:
\end{sourcepara}
\begin{transpara}
\transrn{47}3. 此項1992年7月27日法律（BGBl. I S. 1398）之標題為《關於保護出生前／形成中生命、促進更友善兒童之社會、懷孕衝突協助及終止妊娠規制之法律（孕婦及家庭扶助法）》（\textit{Schwangeren- und Familienhilfegesetz}；以下稱 SFHG）；其主要內容如下：
\end{transpara}

\begin{sourcepara}[title={原文理由 A.II.3.a} Rn 48]
\sourcern{48}a) \abbrArt{} 1 \abbrSFHG{} ("Gesetz über Aufklärung, Verhütung, Familienplanung und Beratung") verpflichtet die Bundeszentrale für gesundheitliche Aufklärung zur Erstellung von Konzepten und zur Verbreitung von Materialien zur Sexualaufklärung (§ 1) und schafft einen Rechtsanspruch auf Beratung (§ 2) durch anerkannte Beratungsstellen (§ 3). Geschuldet werden Informationen u.a. über Sexualaufklärung, Verhütung und Familienplanung, familienfördernde Leistungen und Hilfen für Kinder und Familien, soziale und wirtschaftliche Hilfen für Schwangere, die Methoden zur Durchführung eines Schwangerschaftsabbruchs und die damit verbundenen Risiken sowie Lösungsmöglichkeiten für psychosoziale Konflikte im Zusammenhang mit einer Schwangerschaft. Die Schwangere ist darüber hinaus bei der Geltendmachung von Ansprüchen sowie bei der Wohnungssuche, bei der Suche nach einer Betreuungsmöglichkeit für das Kind und bei der Fortsetzung ihrer Ausbildung zu unterstützen. Die Länder haben dafür Sorge zu tragen, daß den Beratungsstellen für je 40.000 Einwohner mindestens ein Berater zur Verfügung steht. Die Beratungsstellen haben Anspruch auf eine angemessene öffentliche Förderung der Personal- und Sachkosten (§ 4).
\end{sourcepara}
\begin{transpara}
\transrn{48}a) SFHG 第1條（「關於啟蒙、避孕、家庭計畫及諮詢之法律」）使聯邦健康啟蒙中心負有制定性教育構想並散布性教育資料之義務（第1條），並創設由受承認諮詢機構（第3條）提供諮詢之法律請求權（第2條）。應提供之資訊尤其包括性教育、避孕與家庭計畫，促進家庭之給付及對兒童與家庭之協助，對懷孕婦女之社會與經濟協助，實施終止妊娠之方法及其相關風險，以及與懷孕相關之心理社會衝突的解決可能性。此外，於主張請求權、尋找住居、尋找兒童照顧可能性及繼續教育訓練時，亦應支持懷孕婦女。各邦應確保諮詢機構就每四萬名居民至少有一名諮詢人員可供利用。諮詢機構就人事及物品費用享有適當公共補助之請求權（第4條）。
\end{transpara}

\begin{sourcepara}[title={原文理由 A.II.3.a SGB V}]
\sourcern{49}Die durch \abbrArt{} 2 des Gesetzes neu eingefügten §§ 24a, 24b des Fünften Buches des Sozialgesetzbuches (\abbrSGBV{}) ersetzen die bisherigen §§ 200e, 200f und 200g \abbrRVO{}. Nach § 24a \abbrSGBV{} haben Versicherte Anspruch auf ärztliche Beratung über Fragen der Empfängnisregelung, Versicherte bis zum 20. Lebensjahr darüber hinaus auch auf Versorgung mit empfängnisverhütenden Mitteln, soweit sie ärztlich verordnet werden. In § 24b \abbrSGBV{} wird den Versicherten ein Anspruch auf Leistungen bei einem nicht rechtswidrigen Abbruch der Schwangerschaft durch einen Arzt gewährt, wenn dieser in einer hierfür vorgesehenen Einrichtung vorgenommen wird. Diese Vorschrift lautet:
\end{sourcepara}
\begin{transpara}
\transrn{49}由該法第2條新增插入《社會法典》第五編（SGB V）之第24a條、第24b條，取代既有《帝國保險法》第200e條、第200f條及第200g條。依 SGB V 第24a條，被保險人就受胎調節問題享有醫師諮詢請求權；二十歲以下之被保險人，於醫師開立處方之範圍內，並享有供給避孕手段之請求權。SGB V 第24b條賦予被保險人於由醫師實施不具違法性之終止妊娠時之給付請求權，但以該終止妊娠係於為此設置之機構中實施為限。該規定文字如下：
\end{transpara}

\begin{sourceboxpara}
\begin{lawboxpart}{first}{lawrn50}
\sourcern{50}\sourcelawtitle{Schwangerschaftsabbruch und Sterilisation}
\end{lawboxpart}
\end{sourceboxpara}
\begin{transboxpara}
\begin{lawboxpart}{first}{lawrn50}
\transrn{50}\lawtitle{終止妊娠及絕育}
\end{lawboxpart}
\end{transboxpara}
\begin{sourceboxpara}
\begin{lawboxpart}{middle}{lawrn51}
\sourcern{51}\sourceabsno{1}{Versicherte haben Anspruch auf Leistungen bei einer nicht rechtswidrigen Sterilisation und bei einem nicht rechtswidrigen Abbruch der Schwangerschaft durch einen Arzt. Der Anspruch auf Leistungen bei einem nicht rechtswidrigen Schwangerschaftsabbruch besteht nur, wenn dieser in einem Krankenhaus oder einer sonstigen hierfür vorgesehenen Einrichtung im Sinne des Artikels 3 \abbrAbs{} 1 Satz 1 des Fünften Gesetzes zur Reform des Strafrechts vorgenommen wird.}
\end{lawboxpart}
\end{sourceboxpara}
\begin{transboxpara}
\begin{lawboxpart}{middle}{lawrn51}
\transrn{51}\absno{1}{被保險人就不具違法性之絕育及由醫師實施之不具違法性之終止妊娠，享有給付請求權。就不具違法性之終止妊娠之給付請求權，僅於該終止妊娠係於醫院，或於《第五刑法改革法》第3條第1項第1句意義下其他為此設置之機構中實施時，始存在。}
\end{lawboxpart}
\end{transboxpara}
\begin{sourceboxpara}
\begin{lawboxpart}{last}{lawrn52}
\sourcern{52}\sourceabsno{2}{Es werden ärztliche Beratung über die Erhaltung und den Abbruch der Schwangerschaft, ärztliche Untersuchung und Begutachtung zur Feststellung der Voraussetzungen für eine nicht rechtswidrige Sterilisation oder für einen nicht rechtswidrigen Schwangerschaftsabbruch, ärztliche Behandlung, Versorgung mit Arznei-, Verbands- und Heilmitteln sowie Krankenhauspflege gewährt. Anspruch auf Krankengeld besteht, wenn Versicherte wegen einer nicht rechtswidrigen Sterilisation oder wegen eines nicht rechtswidrigen Abbruchs der Schwangerschaft durch einen Arzt arbeitsunfähig werden, es sei denn, es besteht ein Anspruch nach § 44 \abbrAbs{} 1."}
\end{lawboxpart}
\end{sourceboxpara}
\begin{transboxpara}
\begin{lawboxpart}{last}{lawrn52}
\transrn{52}\absno{2}{給付包括關於維持懷孕及終止妊娠之醫師諮詢、為確認不具違法性之絕育或不具違法性之終止妊娠要件之醫師檢查與鑑定、醫療、藥品、敷料與療養用品之供給，以及住院照護。若被保險人因不具違法性之絕育或因由醫師實施之不具違法性之終止妊娠而不能工作，享有疾病津貼請求權；但依第44條第1項已有請求權者，不在此限。}
\end{lawboxpart}
\end{transboxpara}

\begin{sourcepara}[title={原文理由 A.II.3.a 續}]
\sourcern{53}Das Kinder- und Jugendhilfegesetz wird dahin ergänzt (\abbrArt{} 5 \abbrSFHG{}), daß ein Kind vom vollendeten dritten Lebensjahr an "nach Maßgabe des Landesrechts" Anspruch auf Besuch eines Kindergartens hat; vom 1. Januar 1996 an besteht dieser Anspruch ohne Einschränkung. Ferner sind ab diesem Zeitpunkt für Kinder unter drei Jahren und für Kinder im schulpflichtigen Alter nach Bedarf Plätze in Tageseinrichtungen und Tagespflegeplätze vorzuhalten. Die beschlossenen Änderungen des Bundessozialhilfegesetzes (\abbrArt{} 8 \abbrSFHG{}) betreffen Verbesserungen bei der Anerkennung des Mehrbedarfs von werdenden Müttern und alleinerziehenden Personen sowie eine Ausdehnung des Regreßverbots für Unterhaltsansprüche gegen Verwandte ersten Grades einer Hilfeempfängerin, die schwanger ist oder ihr leibliches Kind bis zur Vollendung seines sechsten Lebensjahres betreut. Weitere Änderungen aus dem Bereich sozialer Hilfen betreffen u.a. das Arbeitsförderungsgesetz (\abbrArt{} 6 \abbrSFHG{}), das Berufsbildungsgesetz (\abbrArt{} 7 \abbrSFHG{}), das Zweite Wohnungsbaugesetz (\abbrArt{} 9 \abbrSFHG{}), das Wohnungsbindungsgesetz (\abbrArt{} 10 \abbrSFHG{}) sowie das Belegungsrechtsgesetz (\abbrArt{} 11 \abbrSFHG{}).
\end{sourcepara}
\begin{transpara}
\transrn{53}《兒童及少年協助法》被補充如下（SFHG 第5條）：兒童自滿三歲起，「依各邦法律之規定」享有就讀幼稚園之請求權；自1996年1月1日起，此項請求權不再受限制。此外，自該時點起，對三歲以下兒童及學齡兒童，應依需求備置日間機構名額及日間照顧名額。已決議之《聯邦社會扶助法》修正（SFHG 第8條），涉及改善對準母親及單親者額外需求之承認，並擴張追償禁止：對於懷孕之受扶助人，或照顧其親生子女至該子女滿六歲之受扶助人，其對一親等親屬之扶養請求權，不得追償。社會協助領域之其他修正，尤其涉及《勞動促進法》（SFHG 第6條）、《職業訓練法》（SFHG 第7條）、《第二住宅建設法》（SFHG 第9條）、《住宅約束法》（SFHG 第10條）及《占用權法》（SFHG 第11條）。
\end{transpara}

\begin{sourcepara}[title={原文理由 A.II.3.b}]
\sourcern{54}b) \abbrArt{} 13 \abbrNr{} 1 des Gesetzes ersetzt die §§ 218 bis 219d des Strafgesetzbuches in der Fassung der Bekanntmachung vom 10. März 1987 (\abbrBGBl{} I \abbrS{} 945, 1160) durch neue §§ 218 bis 219b (im folgenden: §§ 218 \abbrFF{} \abbrStGB{} \abbrNF{}), die -- soweit sie in diesem Verfahren von Interesse sind -- folgenden Wortlaut haben:
\end{sourcepara}
\begin{transpara}
\transrn{54}b) 該法第13條第1款，以新的第218條至第219b條取代1987年3月10日公告版本《刑法》（BGBl. I S. 945, 1160）中第218條至第219d條（以下稱新法第218條以下）。就本程序相關範圍內，其文字如下：
\end{transpara}

\begin{sourceboxpara}
\begin{lawboxpart}{first}{lawrn55}
\sourcern{55}\sourcelawtitle{§ 218\\Schwangerschaftsabbruch}
\end{lawboxpart}
\end{sourceboxpara}
\begin{transboxpara}
\begin{lawboxpart}{first}{lawrn55}
\transrn{55}\lawtitle{第218條\\終止妊娠}
\end{lawboxpart}
\end{transboxpara}
\begin{sourceboxpara}
\begin{lawboxpart}{middle}{lawrn56}
\sourcern{56}\sourceabsno{1}{Wer eine Schwangerschaft abbricht, wird mit Freiheitsstrafe bis zu drei Jahren oder mit Geldstrafe bestraft. Handlungen, deren Wirkung vor Abschluß der Einnistung des befruchteten Eies in der Gebärmutter eintritt, gelten nicht als Schwangerschaftsabbruch im Sinne dieses Gesetzes.}
\end{lawboxpart}
\end{sourceboxpara}
\begin{transboxpara}
\begin{lawboxpart}{middle}{lawrn56}
\transrn{56}\absno{1}{終止妊娠者，處三年以下自由刑或罰金。其效果發生於受精卵在子宮著床完成以前之行為，不視為本法意義下之終止妊娠。}
\end{lawboxpart}
\end{transboxpara}
\begin{sourceboxpara}
\begin{lawboxpart}{middle}{lawrn57}
\sourcern{57}\sourceabsno{2}{In besonders schweren Fällen ist die Strafe Freiheitsstrafe von sechs Monaten bis zu fünf Jahren. Ein besonders schwerer Fall liegt in der Regel vor, wenn der Täter}
\end{lawboxpart}
\end{sourceboxpara}
\begin{transboxpara}
\begin{lawboxpart}{middle}{lawrn57}
\transrn{57}\absno{2}{情節特別重大者，處六月以上五年以下自由刑。通常於行為人有下列情形之一時，為情節特別重大：}
\end{lawboxpart}
\end{transboxpara}
\begin{sourceboxpara}
\begin{lawboxpart}{middle}{lawrn58}
\sourcern{58}\sourcenrno{1}{gegen den Willen der Schwangeren handelt oder}
\end{lawboxpart}
\end{sourceboxpara}
\begin{transboxpara}
\begin{lawboxpart}{middle}{lawrn58}
\transrn{58}\nrno{1}{違反懷孕婦女意願而行為；或}
\end{lawboxpart}
\end{transboxpara}
\begin{sourceboxpara}
\begin{lawboxpart}{middle}{lawrn59}
\sourcern{59}\sourcenrno{2}{leichtfertig die Gefahr des Todes oder einer schweren Gesundheitsschädigung der Schwangeren verursacht.}
\end{lawboxpart}
\end{sourceboxpara}
\begin{transboxpara}
\begin{lawboxpart}{middle}{lawrn59}
\transrn{59}\nrno{2}{輕率造成懷孕婦女死亡或嚴重健康損害之危險。}
\end{lawboxpart}
\end{transboxpara}
\begin{sourceboxpara}
\begin{lawboxpart}{middle}{lawrn60}
\sourcern{60}\sourceabsno{3}{Begeht die Schwangere die Tat, so ist die Strafe Freiheitsstrafe bis zu einem Jahr oder Geldstrafe.}
\end{lawboxpart}
\end{sourceboxpara}
\begin{transboxpara}
\begin{lawboxpart}{middle}{lawrn60}
\transrn{60}\absno{3}{懷孕婦女實施該行為者，處一年以下自由刑或罰金。}
\end{lawboxpart}
\end{transboxpara}
\begin{sourceboxpara}
\begin{lawboxpart}{middle}{lawrn61}
\sourcern{61}\sourceabsno{4}{Der Versuch ist strafbar. Die Schwangere wird nicht wegen Versuchs bestraft.}
\end{lawboxpart}
\end{sourceboxpara}
\begin{transboxpara}
\begin{lawboxpart}{middle}{lawrn61}
\transrn{61}\absno{4}{未遂犯罰之。懷孕婦女不因未遂受處罰。}
\end{lawboxpart}
\end{transboxpara}
\begin{sourceboxpara}
\begin{lawboxpart}{middle}{lawrn62}
\sourcern{62}\sourcelawtitle{§ 218 a\\Straflosigkeit des Schwangerschaftsabbruchs}
\end{lawboxpart}
\end{sourceboxpara}
\begin{transboxpara}
\begin{lawboxpart}{middle}{lawrn62}
\transrn{62}\lawtitle{第218a條\\終止妊娠之不罰}
\end{lawboxpart}
\end{transboxpara}
\begin{sourceboxpara}
\begin{lawboxpart}{middle}{lawrn63}
\sourcern{63}\sourceabsno{1}{Der Schwangerschaftsabbruch ist nicht rechtswidrig, wenn}
\end{lawboxpart}
\end{sourceboxpara}
\begin{transboxpara}
\begin{lawboxpart}{middle}{lawrn63}
\transrn{63}\absno{1}{終止妊娠於下列情形不具違法性：}
\end{lawboxpart}
\end{transboxpara}
\begin{sourceboxpara}
\begin{lawboxpart}{middle}{lawrn64}
\sourcern{64}\sourcenrno{1}{die Schwangere den Schwangerschaftsabbruch verlangt und dem Arzt durch eine Bescheinigung nach § 219 \abbrAbs{} 3 Satz 2 nachgewiesen hat, daß sie sich mindestens drei Tage vor dem Eingriff hat beraten lassen (Beratung der Schwangeren in einer Not- und Konfliktlage),}
\end{lawboxpart}
\end{sourceboxpara}
\begin{transboxpara}
\begin{lawboxpart}{middle}{lawrn64}
\transrn{64}\nrno{1}{懷孕婦女要求終止妊娠，並已透過第219條第3項第2句之證明書向醫師證明，其至少於介入前三日接受諮詢（困境與衝突狀態中之懷孕婦女諮詢）；}
\end{lawboxpart}
\end{transboxpara}
\begin{sourceboxpara}
\begin{lawboxpart}{middle}{lawrn65}
\sourcern{65}\sourcenrno{2}{der Schwangerschaftsabbruch von einem Arzt vorgenommen wird und}
\end{lawboxpart}
\end{sourceboxpara}
\begin{transboxpara}
\begin{lawboxpart}{middle}{lawrn65}
\transrn{65}\nrno{2}{終止妊娠由醫師實施；且}
\end{lawboxpart}
\end{transboxpara}
\begin{sourceboxpara}
\begin{lawboxpart}{middle}{lawrn66}
\sourcern{66}\sourcenrno{3}{seit der Empfängnis nicht mehr als zwölf Wochen vergangen sind.}
\end{lawboxpart}
\end{sourceboxpara}
\begin{transboxpara}
\begin{lawboxpart}{middle}{lawrn66}
\transrn{66}\nrno{3}{自受胎以來未超過十二週。}
\end{lawboxpart}
\end{transboxpara}
\begin{sourceboxpara}
\begin{lawboxpart}{middle}{lawrn67}
\sourcern{67}\sourceabsno{2}{Der mit Einwilligung der Schwangeren von einem Arzt vorgenommene Schwangerschaftsabbruch ist nicht rechtswidrig, wenn nach ärztlicher Erkenntnis der Abbruch notwendig ist, um eine Gefahr für das Leben der Schwangeren oder die Gefahr einer schwerwiegenden Beeinträchtigung ihres körperlichen oder seelischen Gesundheitszustandes abzuwenden, sofern diese Gefahr nicht auf andere für sie zumutbare Weise abgewendet werden kann.}
\end{lawboxpart}
\end{sourceboxpara}
\begin{transboxpara}
\begin{lawboxpart}{middle}{lawrn67}
\transrn{67}\absno{2}{經懷孕婦女同意由醫師實施之終止妊娠，若依醫師認識，該終止妊娠為避免懷孕婦女生命危險，或避免其身體或心理健康狀態遭受重大損害之危險所必要，且該危險不能以其他對其可期待之方式排除，則不具違法性。}
\end{lawboxpart}
\end{transboxpara}
\begin{sourceboxpara}
\begin{lawboxpart}{middle}{lawrn68}
\sourcern{68}\sourceabsno{3}{Die Voraussetzungen des Absatzes 2 gelten auch als erfüllt, wenn nach ärztlicher Erkenntnis dringende Gründe für die Annahme sprechen, daß das Kind infolge einer Erbanlage oder schädlicher Einflüsse vor der Geburt an einer nicht behebbaren Schädigung seines Gesundheitszustandes leiden würde, die so schwer wiegt, daß von der Schwangeren die Fortsetzung der Schwangerschaft nicht verlangt werden kann. Dies gilt nur, wenn die Schwangere dem Arzt durch eine Bescheinigung nach § 219 \abbrAbs{} 3 Satz 2 nachgewiesen hat, daß sie sich mindestens drei Tage vor dem Eingriff hat beraten lassen, und wenn seit der Empfängnis nicht mehr als zweiundzwanzig Wochen verstrichen sind.}
\end{lawboxpart}
\end{sourceboxpara}
\begin{transboxpara}
\begin{lawboxpart}{middle}{lawrn68}
\transrn{68}\absno{3}{若依醫師認識，有迫切理由足以認為，兒童因遺傳因素或出生前之有害影響，將罹患其健康狀態不可消除之損害，且該損害重大到不能要求懷孕婦女繼續懷孕，亦視為符合第2項要件。此僅於懷孕婦女已透過第219條第3項第2句之證明書向醫師證明，其至少於介入前三日接受諮詢，且自受胎以來未超過二十二週時，始適用。}
\end{lawboxpart}
\end{transboxpara}
\begin{sourceboxpara}
\begin{lawboxpart}{middle}{lawrn69}
\sourcern{69}\sourceabsno{4}{Die Schwangere ist nicht nach § 218 strafbar, wenn der Schwangerschaftsabbruch nach Beratung (§ 219) von einem Arzt vorgenommen worden ist und seit der Empfängnis nicht mehr als zweiundzwanzig Wochen verstrichen sind. Das Gericht kann von Strafe nach § 218 absehen, wenn die Schwangere sich zur Zeit des Eingriffs in besonderer Bedrängnis befunden hat.}
\end{lawboxpart}
\end{sourceboxpara}
\begin{transboxpara}
\begin{lawboxpart}{middle}{lawrn69}
\transrn{69}\absno{4}{若終止妊娠係經諮詢（第219條）後由醫師實施，且自受胎以來未超過二十二週，懷孕婦女不依第218條受處罰。若懷孕婦女於介入時處於特別困迫，法院得依第218條免除刑罰。}
\end{lawboxpart}
\end{transboxpara}
\begin{sourceboxpara}
\begin{lawboxpart}{middle}{lawrn70}
\sourcern{70}\sourcelawtitle{§ 218 b\\Schwangerschaftsabbruch ohne ärztliche Feststellung;\\unrichtige ärztliche Feststellung}
\end{lawboxpart}
\end{sourceboxpara}
\begin{transboxpara}
\begin{lawboxpart}{middle}{lawrn70}
\transrn{70}\lawtitle{第218b條\\未經醫師確認之終止妊娠；不正確之醫師確認}
\end{lawboxpart}
\end{transboxpara}
\begin{sourceboxpara}
\begin{lawboxpart}{middle}{lawrn71}
\sourcern{71}\sourceabsno{1}{Wer in Fällen des § 218a \abbrAbs{} 2 oder 3 eine Schwangerschaft abbricht, ohne daß ihm die schriftliche Feststellung eines Arztes, der nicht selbst den Schwangerschaftsabbruch vornimmt, darüber vorgelegen hat, ob die Voraussetzungen des § 218a \abbrAbs{} 2 oder 3 Satz 1 gegeben sind, wird mit Freiheitsstrafe bis zu einem Jahr oder mit Geldstrafe bestraft, wenn die tat nicht in § 218 mit Strafe bedroht ist. Wer als Arzt wider besseres Wissen eine unrichtige Feststellung über die Voraussetzungen des § 218a \abbrAbs{} 2 oder 3 Satz 1 zur Vorlage nach Satz 1 trifft, wird mit Freiheitsstrafe bis zu zwei Jahren oder mit Geldstrafe bestraft, wenn die Tat nicht in § 218 mit Strafe bedroht ist. Die Schwangere ist nicht nach Satz 1 oder 2 strafbar.}
\end{lawboxpart}
\end{sourceboxpara}
\begin{transboxpara}
\begin{lawboxpart}{middle}{lawrn71}
\transrn{71}\absno{1}{於第218a條第2項或第3項之情形，終止妊娠者，若其未取得並非自行實施終止妊娠之醫師，就第218a條第2項或第3項第1句要件是否存在所作之書面確認，且該行為未於第218條受刑罰威嚇，處一年以下自由刑或罰金。醫師明知不實而為關於第218a條第2項或第3項第1句要件之不正確確認，以供依第1句提出者，若該行為未於第218條受刑罰威嚇，處二年以下自由刑或罰金。懷孕婦女不依第1句或第2句受處罰。}
\end{lawboxpart}
\end{transboxpara}
\begin{sourceboxpara}
\begin{lawboxpart}{middle}{lawrn71-abs2}
\sourceabsno{2}{Ein Arzt darf Feststellungen nach § 218a \abbrAbs{} 2 oder 3 Satz 1 nicht treffen, wenn ihm die zuständige Stelle dies untersagt hat, weil er wegen einer rechtswidrigen Tat nach Absatz 1, den §§ 218, 219a oder 219b oder wegen einer anderen rechtswidrigen Tat, die er im Zusammenhang mit einem Schwangerschaftsabbruch begangen hat, rechtskräftig verurteilt worden ist. Die zuständige Stelle kann einem Arzt vorläufig untersagen, Feststellungen nach § 218a \abbrAbs{} 2 und 3 Satz 1 zu treffen, wenn gegen ihn wegen des Verdachts einer der in Satz 1 bezeichneten rechtswidrigen Taten das Hauptverfahren eröffnet worden ist.}
\end{lawboxpart}
\end{sourceboxpara}
\begin{transboxpara}
\begin{lawboxpart}{middle}{lawrn71-abs2}
\absno{2}{若醫師因依第1項、第218條、第219a條或第219b條之違法行為，或因其在終止妊娠關聯中所實施之其他違法行為，而經有罪判決確定，且主管機關因此禁止其為第218a條第2項或第3項第1句之確認，該醫師不得為此等確認。若醫師因涉嫌第1句所稱違法行為之一而已開啟主審程序，主管機關得暫時禁止其為第218a條第2項及第3項第1句之確認。}
\end{lawboxpart}
\end{transboxpara}
\begin{sourceboxpara}
\begin{lawboxpart}{middle}{lawrn72}
\sourcern{72}\sourcelawtitle{§ 219\\Beratung der Schwangeren in einer Not- und Konfliktlage}
\end{lawboxpart}
\end{sourceboxpara}
\begin{transboxpara}
\begin{lawboxpart}{middle}{lawrn72}
\transrn{72}\lawtitle{第219條\\困境與衝突狀態中之懷孕婦女諮詢}
\end{lawboxpart}
\end{transboxpara}
\begin{sourceboxpara}
\begin{lawboxpart}{middle}{lawrn73}
\sourcern{73}\sourceabsno{1}{Die Beratung dient dem Lebensschutz durch Rat und Hilfe für die Schwangere unter Anerkennung des hohen Wertes des vorgeburtlichen Lebens und der Eigenverantwortung der Frau. Die Beratung soll dazu beitragen, die im Zusammenhang mit der Schwangerschaft bestehende Not- und Konfliktlage zu bewältigen. Sie soll die Schwangere in die Lage versetzen, eine verantwortungsbewußte eigene Gewissensentscheidung zu treffen. Aufgabe der Beratung ist die umfassende medizinische, soziale und juristische Information der Schwangeren. Die Beratung umfaßt die Darlegung der Rechtsansprüche von Mutter und Kind und der möglichen praktischen Hilfen, insbesondere solcher, die die Fortsetzung der Schwangerschaft und die Lage von Mutter und Kind erleichtern. Die Beratung trägt auch zur Vermeidung künftiger ungewollter Schwangerschaften bei.}
\end{lawboxpart}
\end{sourceboxpara}
\begin{transboxpara}
\begin{lawboxpart}{middle}{lawrn73}
\transrn{73}\absno{1}{諮詢在承認出生前生命之高度價值及婦女自我責任之下，藉由對懷孕婦女之建議與協助服務於生命保護。諮詢應有助於克服與懷孕相關之困境與衝突狀態。諮詢應使懷孕婦女能夠作成負責任之自身良心決定。諮詢之任務，是對懷孕婦女提供全面之醫學、社會及法律資訊。諮詢包括說明母親與兒童之法律請求權及可能之實際協助，尤其是有助於繼續懷孕並改善母親與兒童處境之協助。諮詢亦有助於避免未來非意願懷孕。}
\end{lawboxpart}
\end{transboxpara}
\begin{sourceboxpara}
\begin{lawboxpart}{middle}{lawrn74}
\sourcern{74}\sourceabsno{2}{Die Beratung hat durch eine auf Grund Gesetzes anerkannte Beratungsstelle zu erfolgen. Der Arzt, der den Schwangerschaftsabbruch vornimmt, ist als Berater ausgeschlossen.}
\end{lawboxpart}
\end{sourceboxpara}
\begin{transboxpara}
\begin{lawboxpart}{middle}{lawrn74}
\transrn{74}\absno{2}{諮詢應由依法受承認之諮詢機構為之。實施終止妊娠之醫師不得擔任諮詢者。}
\end{lawboxpart}
\end{transboxpara}
\begin{sourceboxpara}
\begin{lawboxpart}{last}{lawrn75}
\sourcern{75}\sourceabsno{3}{Die Beratung wird nicht protokolliert und ist auf Wunsch der Schwangeren anonym durchzuführen. Die Beratungsstelle hat über die Tatsache, daß eine Beratung gemäß Absatz 1 stattgefunden hat und die Frau damit die Informationen für ihre Entscheidungsfindung erhalten hat, sofort eine mit Datum versehene Bescheinigung auszustellen."}
\end{lawboxpart}
\end{sourceboxpara}
\begin{transboxpara}
\begin{lawboxpart}{last}{lawrn75}
\transrn{75}\absno{3}{諮詢不作成紀錄，並應依懷孕婦女願望匿名進行。諮詢機構應就已依第1項進行諮詢且婦女因而已取得其作成決定所需資訊之事實，立即出具載明日期之證明書。}
\end{lawboxpart}
\end{transboxpara}

\begin{sourcepara}[title={原文理由 A.II.3.b StPO}]
\sourcern{76}\abbrArt{} 14 \abbrSFHG{} ändert einige Vorschriften der Strafprozeßordnung (StPO) und ergänzt insbesondere den § 108 StPO, der die Beschlagnahme von sogenannten Zufallsfunden betrifft, um ein Verwertungsverbot: Bei einem Arzt gefundene Gegenstände, die den Schwangerschaftsabbruch einer Patientin betreffen, dürfen in einem Strafverfahren gegen die Patientin wegen einer Straftat nach § 218 \abbrStGB{} nicht verwertet werden.
\end{sourcepara}
\begin{transpara}
\transrn{76}SFHG 第14條修正《刑事訴訟法》（StPO）若干規定，並尤其就涉及所謂偶然發現物扣押之 StPO 第108條，增補一項使用禁止：於醫師處發現、涉及病患終止妊娠之物品，不得於因《刑法》第218條犯罪而對該病患進行之刑事程序中使用。
\end{transpara}

\begin{sourcepara}[title={原文理由 A.II.3.c}]
\sourcern{77}c) Zwei weitere Änderungen in \abbrArt{} 15 \abbrSFHG{} betreffen das Fünfte Gesetz zur Reform des Strafrechts (5. \abbrStrRG{}) vom 18. Juni 1974 (\abbrBGBl{} I \abbrS{} 1297): Die Neufassung des \abbrArt{} 3 \abbrAbs{} 1 beseitigt -- jedenfalls dem Wortlaut nach -- das durch \abbrArt{} 3 \abbrAbs{} 1 des Fünfzehnten Strafrechtsänderungsgesetzes eingeführte Erfordernis einer behördlichen Zulassung für Einrichtungen außerhalb von Krankenhäusern, in denen Schwangerschaftsabbrüche vorgenommen werden (Satz 1), und bestimmt, daß der Schwangerschaftsabbruch zum frühest möglichen Zeitpunkt vorgenommen werden soll (Satz 2). \abbrArt{} 4 \abbrNF{} betrifft nunmehr Einrichtungen zur Vornahme von Schwangerschaftsabbrüchen und läßt damit das in \abbrArt{} 4 geregelte Erfordernis einer Bundesstatistik entfallen. Die Vorschrift lautet:
\end{sourcepara}
\begin{transpara}
\transrn{77}c) SFHG 第15條中另有兩項修正涉及1974年6月18日《第五刑法改革法（5. StrRG）》（BGBl. I S. 1297）：第3條第1項之新版本，至少依其文義，排除了由《第十五刑法修正法》第3條第1項所引入、對醫院以外實施終止妊娠機構之行政許可要求（第1句），並規定終止妊娠應於最早可能時點實施（第2句）。新法第4條現在涉及實施終止妊娠之機構，因而使第4條所規定之聯邦統計要求失效。該規定文字如下：
\end{transpara}

\begin{sourceboxpara}[title={原文理由 A.II.3.c 新法第4條}]
\begin{lawboxpart}{first}{lawrn78}
\sourcern{78}\sourcelawtitle{Artikel 4\\Einrichtung zur Vornahme von Schwangerschaftsabbrüchen}
\end{lawboxpart}
\end{sourceboxpara}
\begin{transboxpara}
\begin{lawboxpart}{first}{lawrn78}
\transrn{78}\lawtitle{第4條\\實施終止妊娠之機構}
\end{lawboxpart}
\end{transboxpara}
\begin{sourceboxpara}
\begin{lawboxpart}{last}{lawrn79}
\sourcern{79}\sourcequotepara{Die zuständige oberste Landesbehörde stellt ein ausreichendes und flächendeckendes Angebot sowohl ambulanter als auch stationärer Einrichtungen zur Vornahme von Schwangerschaftsabbrüchen sicher."}
\end{lawboxpart}
\end{sourceboxpara}
\begin{transboxpara}
\begin{lawboxpart}{last}{lawrn79}
\transrn{79}主管之邦最高機關確保提供足夠且遍及全域之供給，包括實施終止妊娠之門診機構及住院機構。
\end{lawboxpart}
\end{transboxpara}

\begin{sourcepara}[title={原文理由 A.II.3.c 舊法第4條前文}]
\sourcern{80}Die frühere Fassung hatte folgenden Wortlaut:
\end{sourcepara}
\begin{transpara}
\transrn{80}先前版本文字如下：
\end{transpara}

\begin{sourceboxpara}[title={原文理由 A.II.3.c 舊法第4條}]
\begin{lawboxpart}{first}{lawrn81}
\sourcern{81}\sourcelawtitle{Artikel 4\\Bundesstatistik}
\end{lawboxpart}
\end{sourceboxpara}
\begin{transboxpara}
\begin{lawboxpart}{first}{lawrn81}
\transrn{81}\lawtitle{第4條\\聯邦統計}
\end{lawboxpart}
\end{transboxpara}
\begin{sourceboxpara}
\begin{lawboxpart}{middle}{lawrn82}
\sourcern{82}\sourcequotepara{Über die unter den Voraussetzungen des § 218a des Strafgesetzbuchs vorgenommenen Schwangerschaftsabbrüche wird beim Statistischen Bundesamt eine Bundesstatistik geführt. Wer als Arzt einen solchen Schwangerschaftsabbruch ausgeführt hat, hat dies bis zum Ende des laufenden Kalendervierteljahres mit Angaben über}
\end{lawboxpart}
\end{sourceboxpara}
\begin{transboxpara}
\begin{lawboxpart}{middle}{lawrn82}
\transrn{82}對於依《刑法》第218a條要件所實施之終止妊娠，由聯邦統計局進行聯邦統計。實施此種終止妊娠之醫師，應於當前日曆季度結束前，向聯邦統計局申報下列資料：
\end{lawboxpart}
\end{transboxpara}
\begin{sourceboxpara}
\begin{lawboxpart}{middle}{lawrn83}
\sourcern{83}\sourcenrno{1}{den Grund des Schwangerschaftsabbruchs,}
\end{lawboxpart}
\end{sourceboxpara}
\begin{transboxpara}
\begin{lawboxpart}{middle}{lawrn83}
\transrn{83}\nrno{1}{終止妊娠之理由；}
\end{lawboxpart}
\end{transboxpara}
\begin{sourceboxpara}
\begin{lawboxpart}{middle}{lawrn84}
\sourcern{84}\sourcenrno{2}{den Familienstand und das Alter der Schwangeren sowie die Zahl der von ihr versorgten Kinder,}
\end{lawboxpart}
\end{sourceboxpara}
\begin{transboxpara}
\begin{lawboxpart}{middle}{lawrn84}
\transrn{84}\nrno{2}{懷孕婦女之婚姻狀態與年齡，以及其所扶養子女之人數；}
\end{lawboxpart}
\end{transboxpara}
\begin{sourceboxpara}
\begin{lawboxpart}{middle}{lawrn85}
\sourcern{85}\sourcenrno{3}{die Zahl der vorangegangenen Schwangerschaften und deren Beendigung,}
\end{lawboxpart}
\end{sourceboxpara}
\begin{transboxpara}
\begin{lawboxpart}{middle}{lawrn85}
\transrn{85}\nrno{3}{先前懷孕之次數及其終結情形；}
\end{lawboxpart}
\end{transboxpara}
\begin{sourceboxpara}
\begin{lawboxpart}{middle}{lawrn86}
\sourcern{86}\sourcenrno{4}{die Dauer der abgebrochenen Schwangerschaft,}
\end{lawboxpart}
\end{sourceboxpara}
\begin{transboxpara}
\begin{lawboxpart}{middle}{lawrn86}
\transrn{86}\nrno{4}{終止妊娠時之懷孕期間；}
\end{lawboxpart}
\end{transboxpara}
\begin{sourceboxpara}
\begin{lawboxpart}{middle}{lawrn87}
\sourcern{87}\sourcenrno{5}{die Art des Eingriffs und beobachtete Komplikationen,}
\end{lawboxpart}
\end{sourceboxpara}
\begin{transboxpara}
\begin{lawboxpart}{middle}{lawrn87}
\transrn{87}\nrno{5}{介入之種類及所觀察到之併發症；}
\end{lawboxpart}
\end{transboxpara}
\begin{sourceboxpara}
\begin{lawboxpart}{middle}{lawrn88}
\sourcern{88}\sourcenrno{6}{den Ort der Vornahme des Eingriffs und im Fall eines Krankenhausaufenthalts dessen Dauer sowie}
\end{lawboxpart}
\end{sourceboxpara}
\begin{transboxpara}
\begin{lawboxpart}{middle}{lawrn88}
\transrn{88}\nrno{6}{實施介入之地點，及於住院情形中住院期間；以及}
\end{lawboxpart}
\end{transboxpara}
\begin{sourceboxpara}
\begin{lawboxpart}{middle}{lawrn89}
\sourcern{89}\sourcenrno{7}{gegebenenfalls den fremden Staat, in dem die Schwangere ihren Wohnsitz oder gewöhnlichen Aufenthalt hat,}
\end{lawboxpart}
\end{sourceboxpara}
\begin{transboxpara}
\begin{lawboxpart}{middle}{lawrn89}
\transrn{89}\nrno{7}{必要時，懷孕婦女有住所或通常居所之外國。}
\end{lawboxpart}
\end{transboxpara}
\begin{sourceboxpara}
\begin{lawboxpart}{last}{lawrn89-tail}
\sourcequotepara{dem Statistischen Bundesamt anzuzeigen; der Name der Schwangeren darf dabei nicht angegeben werden."}
\end{lawboxpart}
\end{sourceboxpara}
\begin{transboxpara}
\begin{lawboxpart}{last}{lawrn89-tail}
申報時不得記載懷孕婦女姓名。
\end{lawboxpart}
\end{transboxpara}

\begin{sourcepara}[title={原文理由 A.II.3.c 續}]
\sourcern{90}Schließlich hebt \abbrArt{} 16 \abbrSFHG{} die aufgrund des Einigungsvertrages fortgeltenden Bestimmungen des Rechts der \abbrDDR{} über den Schwangerschaftsabbruch auf.
\end{sourcepara}
\begin{transpara}
\transrn{90}最後，SFHG 第16條廢止因《統一條約》而繼續適用之 DDR 關於終止妊娠之法律規定。
\end{transpara}

\bisubsubsection{III.}{III.}

\begin{sourcepara}[title={原文理由 A.III}]
\sourcern{91}Mit Urteil vom 4. August 1992 (\abbrBVerfGE{} 86, 390) ordnete das Bundesverfassungsgericht aufgrund von Anträgen der Bayerischen Staatsregierung sowie von 248 Abgeordneten des Deutschen Bundestages gemäß § 32 \abbrBVerfGG{} u.a. an, daß \abbrArt{} 13 \abbrNr{} 1 und \abbrArt{} 16 des Schwangeren- und Familienhilfegesetzes vom 27. Juli 1992 (\abbrBGBl{} I \abbrS{} 1398) einstweilen nicht in Kraft treten und die in \abbrArt{} 4 (Bundesstatistik) des Fünften Strafrechtsreformgesetzes vom 18. Juni 1974 (\abbrBGBl{} I \abbrS{} 1297), geändert durch \abbrArt{} 3 und \abbrArt{} 4 des Gesetzes vom 18. Mai 1976 (\abbrBGBl{} I \abbrS{} 1213) getroffenen Regelungen einstweilen in Kraft bleiben und auch in dem in \abbrArt{} 3 des Einigungsvertrages genannten Gebiet anzuwenden sind (vgl. \abbrBVerfGE{} 86, 390 \abbrFF{}; \abbrBGBl{} 1992 I \abbrS{} 1585). Die einstweilige Anordnung wurde durch Beschluß vom 25. Januar 1993 wiederholt (\abbrBVerfGE{} 88, 83 = \abbrBGBl{} I \abbrS{} 270).
\end{sourcepara}
\begin{transpara}
\transrn{91}聯邦憲法法院依巴伐利亞邦政府及德國聯邦議會248名議員之聲請，依《聯邦憲法法院法》第32條，於1992年8月4日判決（BVerfGE 86, 390）中命令，除其他事項外，1992年7月27日《孕婦及家庭扶助法》（BGBl. I S. 1398）第13條第1款及第16條暫不生效；1974年6月18日《第五刑法改革法》（BGBl. I S. 1297）第4條（聯邦統計）所定規制，該規制經1976年5月18日法律第3條及第4條修正（BGBl. I S. 1213），暫時繼續有效，並亦適用於《統一條約》第3條所稱地區（參 BVerfGE 86, 390 ff.; BGBl. 1992 I S. 1585）。該臨時命令經1993年1月25日裁定重申（BVerfGE 88, 83 = BGBl. I S. 270）。
\end{transpara}

\bisubsection{B.}{B.}

\bisubsubsection{I.}{I.}

\begin{sourcepara}[title={原文理由 B.I 開頭}]
\sourcern{92}Die Bayerische Staatsregierung hat im Verfahren 2 \abbrBvF{} 2/90 gemäß \abbrArt{} 93 \abbrAbs{} 1 \abbrNr{} 2 \abbrGG{}\ggfnArtNinetyThreeOneTwoDE{}, § 13 \abbrNr{} 6 \abbrBVerfGG{} die verfassungsgerichtliche Überprüfung der Vorschriften in § 218b \abbrAbs{} 1 Satz 1 und \abbrAbs{} 2, § 219 \abbrAbs{} 1 Satz 1 \abbrStGB{} in der Fassung des 15. Strafrechtsänderungsgesetzes und der §§ 200f, 200g \abbrRVO{} beantragt, soweit diese Vorschriften Schwangerschaftsabbrüche aufgrund der allgemeinen Notlagenindikation (§ 218a \abbrAbs{} 2 \abbrNr{} 3 \abbrStGB{} in der Fassung des 15. Strafrechtsänderungsgesetzes) betreffen. Sie macht geltend, die Vorschriften der Reichsversicherungsordnung seien im angegebenen Umfang nichtig; für die angegriffenen Bestimmungen des Strafgesetzbuches müsse der Gesetzgeber binnen angemessener Frist eine Neuregelung treffen, die den verfassungsrechtlichen Anforderungen genüge.
\end{sourcepara}
\begin{transpara}
\transrn{92}巴伐利亞邦政府於 2 BvF 2/90 程序中，依《基本法》第93條第1項\ggfnArtNinetyThreeOneTwoZH{}第2款及《聯邦憲法法院法》第13條第6款，聲請憲法法院審查《第十五刑法修正法》版本《刑法》第218b條第1項第1句及第2項、第219條第1項第1句，以及《帝國保險法》第200f條、第200g條之規定；其審查範圍限於此等規定涉及基於一般困境指標事由（《第十五刑法修正法》版本《刑法》第218a條第2項第3款）之終止妊娠。巴伐利亞邦政府主張，《帝國保險法》規定於上述範圍內無效；至於遭攻擊之《刑法》規定，立法者必須於適當期間內制定符合憲法要求之新規制。
\end{transpara}

\begin{sourcepara}[title={原文理由 B.I.1}]
\sourcern{93}1. Die Vorschriften der §§ 218b, 219 \abbrAbs{} 1 \abbrStGB{} \abbrAF{} stellten keinen ausreichenden Ausgleich dafür dar, daß der Schwangerschaftsabbruch bei bestimmten Indikationen nicht mit Strafe bedroht werde. Sie genügten -- wie die Antragstellerin näher darlegt -- weder in ihrer Summe noch im Detail den verfassungsrechtlichen Anforderungen an einen wirksamen Schutz des ungeborenen Lebens, die das Bundesverfassungsgericht in seinem Urteil vom 25. Februar 1975 formuliert habe. Die Vorschriften ermöglichten in Verbindung mit der unscharfen Formulierung der Notlagenindikation bei Einhaltung gewisser Formalitäten Straflosigkeit wie im Falle einer Fristenlösung. Als "allgemeine Notlage" werde auf der Grundlage dieses Konzepts in der Praxis alles anerkannt, was gegen die Bedürfnisse und Lebensperspektiven der Frau gerichtet sei und sie gefährde. Es liege in der Konsequenz dieses Konzepts, daß die Frau letztlich selbst einzuschätzen habe, ob sie sich auf eine Indikation berufen könne. Die Indikationsfeststellung werde auf diese Weise zum Gegenstand eines vermeintlichen Rechtsanspruchs der Frau, der sich in einem Anspruch auf Entlastung vom finanziellen Aufwand des Abbruchs fortsetze. Die Beratung erfolge dabei nur "pro forma" und gelte schwerpunktmäßig der Vermittlung indikationsfeststellender und abbruchbereiter Ärzte.
\end{sourcepara}
\begin{transpara}
\transrn{93}1. 舊法《刑法》第218b條、第219條第1項之規定，並未對於特定指標事由下終止妊娠不受刑罰威嚇一事，提供充分補償。依聲請人進一步說明，此等規定無論整體或細節，均不符合聯邦憲法法院於1975年2月25日判決中所提出、有效保護未出生生命之憲法要求。此等規定結合困境指標事由之模糊文字，使人在遵守若干形式要件時，即可如期限規定之情形般不受處罰。依此構想，在實務上，凡與婦女需求及生活前景相牴觸並危及婦女者，均被承認為「一般困境」。此一構想之結果是，最終由婦女自行評估其是否得主張指標事由。指標事由之確認因而成為婦女一項表面上法律請求權之標的，並延伸為免除終止妊娠財務負擔之請求權。諮詢在此僅「形式上」進行，重點則在媒介願意確認指標事由並願意實施終止妊娠之醫師。
\end{transpara}

\begin{sourcepara}[title={原文理由 B.I.1} Rn 94]
\sourcern{94}Die Vorschriften seien daher mit dem durch \abbrArt{} 2 \abbrAbs{} 2 Satz 1 in Verbindung mit \abbrArt{} 1 \abbrAbs{} 1 \abbrGG{}\ggfnArtOneOneTwoTwoOneDE{} gewährleisteten Schutz des ungeborenen Lebens unvereinbar. Da ihre Nichtigkeit oder Unanwendbarkeit aber selbst die ihnen zukommende begrenzte Schutzwirkung für das werdende Leben beseitigte, seien sie indes ungeachtet ihrer Verfassungswidrigkeit bis zu der gebotenen Neuregelung für weiterhin anwendbar zu erklären. Die Verpflichtung des Gesetzgebers, binnen angemessener Frist eine Nachbesserung vorzunehmen, die den verfassungsrechtlichen Anforderungen genüge, ergebe sich aus der grundgesetzlichen Schutzpflicht für das ungeborene Leben.
\end{sourcepara}
\begin{transpara}
\transrn{94}因此，此等規定與《基本法》第2條第2項第1句結合第1條第1項\ggfnArtOneOneTwoTwoOneZH{}所保障之未出生生命保護不相容。然而，由於宣告其無效或不適用，本身亦會消除其對形成中生命所具有之有限保護效果，因此，儘管其違憲，仍應宣告其至必要新規制制定前繼續適用。立法者負有於適當期間內進行補正、使其符合憲法要求之義務，係源於《基本法》對未出生生命之保護義務。
\end{transpara}

\begin{sourcepara}[title={原文理由 B.I.2.a}]
\sourcern{95}2. Die §§ 200f, 200g \abbrRVO{} seien aus kompetenzrechtlichen und materiellen Gründen verfassungswidrig und nichtig, soweit sie Versicherten einen Anspruch auf Leistungen der gesetzlichen Krankenversicherung bei Schwangerschaftsabbrüchen gewähren, die nach § 218a \abbrAbs{} 2 \abbrNr{} 3 \abbrStGB{} \abbrAF{} (Notlagenindikation) nicht strafbar sind.
\end{sourcepara}
\begin{transpara}
\transrn{95}2. 《帝國保險法》第200f條、第200g條，在其就依舊法《刑法》第218a條第2項第3款（困境指標事由）不受處罰之終止妊娠，賦予被保險人法定醫療保險給付請求權之範圍內，因權限法上及實體法上理由而違憲並無效。
\end{transpara}

\begin{sourcepara}[title={原文理由 B.I.2.a} Rn 96]
\sourcern{96}a) Die Gesetzgebungskompetenz des Bundes lasse sich nicht aus \abbrArt{} 74 \abbrNr{} 12 \abbrGG{}\ggfnArtSeventyFourTwelveDE{} ("Sozialversicherung") herleiten. Die soziale Krankenversicherung diene der Absicherung gegen Krankheit und verwandte Risiken durch Zusammenschluß und Beitragszahlung gleichartig Bedrohter. In den Fällen der sogenannten Notlagenindikation sei der Schutz vor Mutterschaft kein typisches vorsorgefähiges Risiko, das von der Versichertengemeinschaft zu tragen sei. Mutterschaft sei keine Krankheit; ein Schwangerschaftsabbruch, der nach § 218a \abbrAbs{} 2 \abbrNr{} 3 \abbrStGB{} \abbrAF{} indiziert sei, werde auch nicht als Heileingriff bewertet.
\end{sourcepara}
\begin{transpara}
\transrn{96}a) 聯邦之立法權限不能由《基本法》第74條第12款\ggfnArtSeventyFourTwelveZH{}（「社會保險」）導出。社會醫療保險係透過同類受威脅者之結合及繳納保費，保障其對抗疾病及類似風險。於所謂困境指標事由之情形，對母性之保護並非典型可預防、且應由被保險人共同體承擔之風險。母性並非疾病；依舊法《刑法》第218a條第2項第3款具有指標事由之終止妊娠，亦不被評價為治療介入。
\end{transpara}

\begin{sourcepara}[title={原文理由 B.I.2.a} Rn 97]
\sourcern{97}Die Gesetzgebungszuständigkeit ergebe sich auch nicht aus \abbrArt{} 74 \abbrNr{} 7 \abbrGG{}\ggfnArtSeventyFourSevenDE{} ("öffentliche Fürsorge"). Der Erwägung, daß Schwangere bei nicht mit Strafe bedrohten Schwangerschaftsabbrüchen nicht wegen ihrer wirtschaftlichen Verhältnisse benachteiligt werden dürften, könne der Gesetzgeber aufgrund des \abbrArt{} 74 \abbrNr{} 7 \abbrGG{}\ggfnArtSeventyFourSevenDE{} nur im Sozialhilferecht Rechnung tragen (vgl. § 37a BSHG). Er könne aber aufgrund dieser Erwägung nicht allgemeine Leistungsansprüche in der gesetzlichen Krankenversicherung begründen, die die individuelle Leistungsfähigkeit der Versicherten nicht berücksichtigten, und die dadurch entstehenden versicherungsfremden Lasten auf alle Beitragszahler verteilen.
\end{sourcepara}
\begin{transpara}
\transrn{97}立法權限亦非源於《基本法》第74條第7款\ggfnArtSeventyFourSevenZH{}（「公共扶助」）。關於懷孕婦女在不受刑罰威嚇之終止妊娠時，不得因其經濟狀況受不利對待之考量，立法者依《基本法》第74條第7款\ggfnArtSeventyFourSevenZH{}，僅得於社會扶助法中加以顧及（參《聯邦社會扶助法》第37a條）。立法者不能基於此一考量，在法定醫療保險中創設不考量被保險人個別給付能力之一般給付請求權，並將由此產生之保險外負擔分配予全體繳費者。
\end{transpara}

\begin{sourcepara}[title={原文理由 B.I.2.b}]
\sourcern{98}b) Die Vorschriften der §§ 200f, 200g \abbrRVO{} seien auch materiell verfassungswidrig. Die Schutzpflicht für das ungeborene Leben gebiete den Staatsorganen, sich in allen Bereichen der Rechtsordnung schützend und fördernd vor dieses Leben zu stellen. Dagegen verstießen die angegriffenen Bestimmungen: Das Grundgesetz verwehre es dem Gesetzgeber zwar nicht, Schwangerschaftsabbrüche aufgrund einer allgemeinen Notlagenindikation nicht mit Strafe zu bedrohen. Er sei jedoch gehindert, für diesen Fall Leistungen der gesetzlichen Krankenversicherung vorzusehen und dadurch zur Vernichtung des zu schützenden Rechtsgutes Hilfe zu leisten. Der Staat setze damit seine sozialpolitischen Mittel nicht für, sondern gegen das werdende Leben ein. Die Kassenleistungen bildeten zudem einen Anreiz für die ausufernde Anwendung des Tatbestandes der allgemeinen Notlagenindikation. Die Hemmschwelle gegenüber Schwangerschaftsabbrüchen werde ganz allgemein gesenkt, indem der Schwangerschaftsabbruch im "sozialen Netz" aufgefangen werde.
\end{sourcepara}
\begin{transpara}
\transrn{98}b) 《帝國保險法》第200f條、第200g條之規定，在實體上亦屬違憲。對未出生生命之保護義務，要求國家機關在法律秩序所有領域中，站在此一生命之前，對其加以保護並促進。遭攻擊之規定違反此要求：雖然《基本法》並不禁止立法者對基於一般困境指標事由之終止妊娠不以刑罰威嚇；然而，立法者不得於此種情形規定法定醫療保險給付，並藉此協助毀滅應受保護之法益。國家因此並非為形成中生命，而是反其方向運用其社會政策手段。保險基金給付並且形成誘因，使一般困境指標事由構成要件被過度擴張適用。透過將終止妊娠納入「社會安全網」承接，對終止妊娠之抑制門檻被一般性降低。
\end{transpara}

\begin{sourcepara}[title={原文理由 B.I.2.b} Rn 99]
\sourcern{99}Das Bundesverfassungsgericht habe in seiner Entscheidung vom 25. Februar 1975 ausgesprochen, daß die Mißbilligung des Schwangerschaftsabbruchs in der Rechtsordnung klar zum Ausdruck kommen müsse. Es sei der falsche Eindruck zu vermeiden, es handle sich beim Schwangerschaftsabbruch um den gleichen sozialen Vorgang wie etwa einen Gang zum Arzt zwecks Heilung einer Krankheit oder gar um eine rechtlich irrelevante Alternative zur Empfängnisverhütung. Gerade dieses Mißverständnis förderten die angegriffenen Vorschriften. Die Indikationstatbestände nach § 218a \abbrStGB{} \abbrAF{} in Verbindung mit der gesetzlichen Einräumung eines Anspruchs auf Kassenleistungen erweckten den Eindruck, als sei die Abtreibung Gegenstand eines sozialen Anspruchs. Hier liege eine wesentliche Ursache für den vielfach beklagten Verlust des rechtsethischen Bewußtseins für die Schutzbedürftigkeit des ungeborenen Lebens.
\end{sourcepara}
\begin{transpara}
\transrn{99}聯邦憲法法院於1975年2月25日判決中曾表明，法律秩序必須清楚表達其對終止妊娠之否定評價。應避免造成錯誤印象，彷彿終止妊娠與為治療疾病而就醫之社會過程相同，甚至只是避孕之一種法律上無關緊要替代方式。遭攻擊規定正促進此種誤解。舊法《刑法》第218a條之指標事由構成要件，結合法律賦予之保險基金給付請求權，造成終止妊娠彷彿是社會請求權標的之印象。此處正是許多人所抱怨、對未出生生命需受保護性之法律倫理意識喪失的重要原因。
\end{transpara}

\begin{sourcepara}[title={原文理由 B.I.2.b} Rn 100]
\sourcern{100}Schließlich seien die angegriffenen Vorschriften der Reichsversicherungsordnung auch deshalb verfassungswidrig, weil nicht sichergestellt sei, daß die Kassenleistungen nur in den gesetzlich vorgesehenen Fällen in Anspruch genommen würden; sie enthielten keinerlei Ansatz für eine Mißbrauchsabwehr. Die Krankenkassen seien nicht verpflichtet, ihre Leistungen vom Nachweis der Indikationsvoraussetzungen durch ein begründetes ärztliches Votum abhängig zu machen. Das geltende Recht schreibe auch nicht vor, daß der Arzt seine Leistungen bei einem Schwangerschaftsabbruch nur abrechnen dürfe, wenn er seine Meldepflicht gemäß \abbrArt{} 4 des 5. \abbrStrRG{} in Verbindung mit \abbrArt{} 3 \abbrNr{} 2 des 15. StÄG erfüllt habe. Eine solche Koppelung könnte den Arzt zu einer gesetzestreueren Handhabung des § 218a \abbrAbs{} 1 und 2 \abbrStGB{} \abbrAF{} bewegen und die Aufsichtsbehörden der Länder in die Lage versetzen, mit rechtsaufsichtlichen Mitteln auf deren Einhaltung hinzuwirken.
\end{sourcepara}
\begin{transpara}
\transrn{100}最後，遭攻擊之《帝國保險法》規定亦因下列理由違憲：其並未確保保險基金給付僅於法律所規定情形中被利用；其完全欠缺防止濫用之著力點。醫療保險基金並無義務使其給付取決於以有理由之醫師意見證明指標事由要件。現行法亦未規定，醫師僅於其已依《第五刑法改革法》第4條結合《第十五刑法修正法》第3條第2款履行通報義務時，始得就終止妊娠時之服務進行結算。此種連結可能促使醫師更守法地處理舊法《刑法》第218a條第1項及第2項，並使各邦監督機關能夠以法律監督手段促成其遵守。
\end{transpara}

\bisubsubsection{II.}{II.}

\begin{sourcepara}[title={原文理由 B.II.1}]
\sourcern{101}1. Die Bundesregierung und die Landesregierungen von Bremen, Hamburg, Hessen, Niedersachsen, Nordrhein-Westfalen, Saarland und Schleswig-Holstein halten den Antrag bezüglich der Beratung und Indikationsfeststellung für unbegründet; die Landesregierungen sehen ihn überdies als durch den Einigungsvertrag überholt an. Das Grundkonzept des Fünfzehnten Strafrechtsänderungsgesetzes sei vom Bundesverfassungsgericht in seiner Entscheidung vom 25. Februar 1975 gebilligt worden. Im vorliegenden Verfahren -- so führt die Bundesregierung aus -- gehe es allein um die Frage, ob der Gesetzgeber Rechtsänderungen vornehmen müsse, um die Umsetzung seiner verfassungsrechtlich unbedenklichen Bewertung von Schwangerschaftsabbrüchen in der Rechtswirklichkeit zu verbessern. Er könne jedoch von Verfassungs wegen nicht auf bestimmte Details der Beratungsregelung festgelegt sein, wenn eine andere -- auch die bestehende -- Regelung nach seiner Einschätzung ebenfalls geeignet sei, dem Schutz des ungeborenen Lebens zu dienen. Die Wahl eines bestimmten Konzepts sei insoweit verfassungsrechtlich nicht vorgegeben. Die Kritik der Antragstellerin betreffe Detailregelungen, die in den Gestaltungsspielraum des Gesetzgebers fielen. Nach Ansicht der Landesregierungen trifft auch der Vorwurf der Antragstellerin nicht zu, das geltende Recht sei im Blick auf die Generalprävention ineffizient. Es trage im Gegenteil zur Bildung des Rechtsbewußtseins der Bevölkerung bei.
\end{sourcepara}
\begin{transpara}
\transrn{101}1. 聯邦政府及不來梅、漢堡、黑森、下薩克森、北萊茵-西發利亞、薩爾及石勒蘇益格-荷爾斯泰因各邦政府，認為關於諮詢及指標事由確認之聲請無理由；各邦政府並認為，該聲請已因《統一條約》而過時。《第十五刑法修正法》之基本構想，已由聯邦憲法法院於1975年2月25日判決中認可。聯邦政府說明，本程序僅涉及下列問題：立法者是否必須進行法律修正，以改善其對終止妊娠所為、憲法上無疑義之評價在法律現實中之實現。然而，若另一項規制，包括現行規制，依立法者評估同樣適於服務未出生生命之保護，憲法上即不能將立法者固定於諮詢規制之特定細節。就此而言，憲法並未預設選擇某一特定構想。聲請人之批評涉及細部規制，而此屬立法者形成空間。依各邦政府見解，聲請人所稱現行法就一般預防而言無效率，亦不正確。相反地，現行法有助於形成人民之法律意識。
\end{transpara}

\begin{sourcepara}[title={原文理由 B.II.1} Rn 102]
\sourcern{102}Die von der Antragstellerin angegriffenen Vorschriften der Reichsversicherungsordnung seien gleichfalls verfassungsrechtlich nicht zu beanstanden. Dem Bund stehe für die §§ 200f, 200g \abbrRVO{} die Gesetzgebungskompetenz zu; sie unterfielen als "sonstige Hilfen" im Sinne der §§ 200e \abbrFF{} \abbrRVO{} dem Kompetenzbereich der "Sozialversicherung" nach \abbrArt{} 74 \abbrNr{} 12 \abbrGG{}\ggfnArtSeventyFourTwelveDE{}. Der Begriff der Sozialversicherung sei ein weitgefaßter verfassungsrechtlicher Gattungsbegriff, der für neue Lebenssachverhalte und neue Sozialleistungen offen sei, wenn diese in ihren wesentlichen Strukturelementen, insbesondere in der organisatorischen Durchführung und hinsichtlich der abzudeckenden Risiken, dem Bild entsprächen, das durch die klassische Sozialversicherung geprägt sei. Die Gewährung der Leistungen nach §§ 200f, 200g \abbrRVO{} sei nicht nur an einen Träger der klassischen Sozialversicherung angelehnt, sondern in dessen Leistungsgewährung integriert. Die Leistungen wiesen auch inhaltliche Strukturelemente der Sozialversicherung auf. Sie hätten -- so die Bundesregierung -- vor allem auch den gesundheitlichen Schutz versicherter Frauen zum Ziel, die einen von der Rechtsordnung nicht mißbilligten Schwangerschaftsabbruch durchführen ließen. Die umfassenden Kassenleistungen sollten gewährleisten, daß aus medizinischer Sicht alles Erforderliche getan werde, um gesundheitlichen Beeinträchtigungen und Gesundheitsschäden vorzubeugen.
\end{sourcepara}
\begin{transpara}
\transrn{102}聲請人所攻擊之《帝國保險法》規定，在憲法上同樣無可指摘。聯邦對《帝國保險法》第200f條、第200g條具有立法權限；作為《帝國保險法》第200e條以下意義下之「其他協助」，其落入《基本法》第74條第12款\ggfnArtSeventyFourTwelveZH{}「社會保險」之權限領域。社會保險概念是一項廣義之憲法類型概念，對新的生活事實及新的社會給付保持開放，只要其在基本結構要素上，尤其在組織執行及所涵蓋風險方面，符合由典型社會保險所形塑之圖像。依《帝國保險法》第200f條、第200g條提供給付，不僅是附著於典型社會保險之承辦者，並且整合於其給付提供之中。此等給付在內容上亦具有社會保險之結構要素。依聯邦政府見解，其尤其亦以保護被保險婦女之健康為目的，該等婦女實施法律秩序並未予以否定評價之終止妊娠。全面之保險基金給付，應確保從醫學觀點而言，為預防健康不利影響及健康損害，所有必要事項均已完成。
\end{transpara}

\begin{sourcepara}[title={原文理由 B.II.1} Rn 103]
\sourcern{103}Die §§ 200f, 200g \abbrRVO{} verstießen auch nicht gegen das Grundrecht auf Leben und körperliche Unversehrtheit. Für einen kausalen Zusammenhang zwischen dieser Regelung und den Abbruchzahlen gebe es keine tatsächlichen Anhaltspunkte. Im Gegenteil spreche alles dafür, daß keine Frau die schwerwiegende Entscheidung über den Abbruch einer Schwangerschaft von einer Kostenlast von etwa 300 bis 1.500 DM abhängig machen würde. Einen "Anreiz" zum oder eine "Prämierung" des Schwangerschaftsabbruchs stellten die Leistungen der Krankenkasse nicht dar. Die Leistungen prämierten nicht den Schwangerschaftsabbruch, sondern das Sich-Einlassen der Frau auf das Verfahren der Beratung und Indikationsfeststellung und dienten damit dem Schutz des ungeborenen Lebens. Im übrigen dürfe die gesetzliche Krankenversicherung nur dann Leistungen erbringen, wenn der Abbruch nicht rechtswidrig sei. Die sozialversicherungsrechtlichen Normen enthielten keine eigene rechtliche Bewertung von Schwangerschaftsabbrüchen, für die verfassungsrechtlich auch kein Anlaß bestehe; sie stellten sich lediglich als akzessorische Folgeregelungen zu den strafrechtlichen Bestimmungen dar. Der von der Antragstellerin behauptete Mißbrauch der §§ 200f, 200g \abbrRVO{} zur Finanzierung gesetzwidriger Abbrüche stelle dieses Ergebnis nicht in Frage. In einem Normenkontrollverfahren könne es nur um mögliche Versäumnisse des Gesetzgebers, nicht aber um solche von Verwaltung und Justiz gehen.
\end{sourcepara}
\begin{transpara}
\transrn{103}《帝國保險法》第200f條、第200g條亦未違反生命及身體完整性基本權。關於此項規制與終止妊娠數量之間存在因果關係，並無事實根據。相反地，一切均顯示，沒有婦女會使終止妊娠之重大決定取決於約 300 至 1,500 德國馬克之費用負擔。醫療保險基金之給付並非對終止妊娠之「誘因」或「獎勵」。給付所獎勵者並非終止妊娠，而是婦女願意進入諮詢及指標事由確認程序，因而服務於未出生生命之保護。此外，法定醫療保險只有在終止妊娠不具違法性時，始得提供給付。社會保險法規範並未包含對終止妊娠之自身法律評價，憲法上亦無此必要；其僅呈現為刑法規定之附隨後續規制。聲請人主張《帝國保險法》第200f條、第200g條遭濫用以資助違法終止妊娠，並不動搖此一結果。在規範審查程序中，所能涉及者僅為立法者可能之疏失，而非行政與司法之疏失。
\end{transpara}

\begin{sourcepara}[title={原文理由 B.II.1} Rn 104]
\sourcern{104}Zur Beratungspraxis haben die Landesregierungen das Ergebnis einer aufgrund eines gemeinsam erarbeiteten Fragebogens durchgeführten Umfrage bei den Beratungsstellen in den Ländern Niedersachsen, Nordrhein-Westfalen, Saarland und Schleswig-Holstein mitgeteilt sowie einen Erfahrungsbericht über den Vollzug des § 218b \abbrStGB{} \abbrAF{} im Lande Bremen in den Jahren 1987 und 1988 vorgelegt.
\end{sourcepara}
\begin{transpara}
\transrn{104}關於諮詢實務，各邦政府提出依共同制定之問卷，對下薩克森、北萊茵-西發利亞、薩爾及石勒蘇益格-荷爾斯泰因各邦諮詢機構所作調查之結果，並提出不來梅邦於1987年及1988年執行舊法《刑法》第218b條之經驗報告。
\end{transpara}

\begin{sourcepara}[title={原文理由 B.II.2--4}]
\sourcern{105}2. Nach Auffassung der Landesregierung von Baden-Württemberg entspricht die Rechtslage beim Beratungs- und Indikationsfeststellungsverfahren nicht den im Urteil des Bundesverfassungsgerichts vom 25. Februar 1975 gestellten Anforderungen.
\end{sourcepara}
\begin{transpara}
\transrn{105}2. 依巴登符騰堡邦政府見解，關於諮詢及指標事由確認程序之法律狀態，並不符合聯邦憲法法院1975年2月25日判決所提出之要求。
\end{transpara}

\begin{sourcepara}[title={原文理由 B.II.2--4} Rn 106]
\sourcern{106}Der Staat könne seine Verpflichtung, das sich entwickelnde Leben in Schutz zu nehmen, nur dann wirksam erfüllen, wenn er bundeseinheitlich verbindliche Maßstäbe für Beratung und Indikation setze und die bestehende Rechtszersplitterung in den Ländern beseitige. Die §§ 218b, 219 \abbrStGB{} \abbrAF{} regelten das Beratungs- und Indikationsfeststellungsverfahren sowie den Inhalt der Beratung nur lückenhaft und ungenau. So nähmen einige Träger solcher Einrichtungen für sich in Anspruch, die Beratung "wertneutral" führen zu dürfen. Außerdem ließen die geltenden Bestimmungen Vorkehrungen dagegen vermissen, daß die im Interesse eines wirksamen Lebensschutzes unumgängliche Eigenständigkeit von Beratung und Indikation gewahrt bleibe.
\end{sourcepara}
\begin{transpara}
\transrn{106}國家只有在制定全聯邦一致且具拘束力之諮詢與指標事由標準，並排除各邦現存法律分裂時，始能有效履行其保護發展中生命之義務。舊法《刑法》第218b條、第219條對諮詢及指標事由確認程序，以及諮詢內容之規制，僅屬有漏洞且不精確。例如，若干此類機構之承辦者主張自己得以「價值中立」方式進行諮詢。此外，現行規定亦欠缺預防措施，以確保在有效生命保護利益上不可或缺之諮詢與指標事由確認的獨立性得以維持。
\end{transpara}

\begin{sourcepara}[title={原文理由 B.II.2--4} Rn 107]
\sourcern{107}3. Die Landesregierungen von Rheinland-Pfalz und Thüringen beschränken sich in ihren Äußerungen im wesentlichen darauf, zur Beratungspraxis in ihrem Land Stellung zu nehmen.
\end{sourcepara}
\begin{transpara}
\transrn{107}3. 萊茵蘭-普法爾茨及圖林根邦政府之陳述，基本上限於就其各邦諮詢實務表示意見。
\end{transpara}

\begin{sourcepara}[title={原文理由 B.II.2--4} Rn 108]
\sourcern{108}4. Von den obersten Gerichtshöfen des Bundes haben der Bundesgerichtshof, das Bundesverwaltungsgericht, das Bundesarbeitsgericht und das Bundessozialgericht Stellungnahmen einzelner Senate übermittelt. Von den angehörten Beratungsstellen oder deren Trägern, Versicherungsträgern und sonstigen Stellen haben sich geäußert: Sozialdienst katholischer Frauen -- Zentrale e.V., Deutscher Caritasverband e.V., Diakonisches Werk der EKD in Deutschland e.V., Pro Familia Deutsche Gesellschaft für Sexualberatung und Familienplanung e.V., Arbeiterwohlfahrt -- Bundesverband e.V., Deutsche Arbeitsgemeinschaft für Jugend und Eheberatung e.V., Ulmer Beratungsstelle für Problemschwangerschaften e.V., Hannoversche Arbeitsgemeinschaft für Jugend und Eheberatung e.V., Soziale Beratungsstelle der Landeshauptstadt Stuttgart für werdende Mütter, Sozialmedizinische Familienberatung in Düsseldorf, AOK-Bundesverband, Bundesärztekammer und Deutscher Ärztinnenbund e.V.
\end{sourcepara}
\begin{transpara}
\transrn{108}4. 聯邦最高法院中，聯邦普通法院、聯邦行政法院、聯邦勞動法院及聯邦社會法院，已轉交個別庭之意見。受詢問之諮詢機構或其承辦者、保險承辦者及其他機關中，提出意見者包括：天主教婦女社會服務中央協會（Sozialdienst katholischer Frauen -- Zentrale e.V.）、德國明愛協會（Deutscher Caritasverband e.V.）、德國福音教會濟助會（Diakonisches Werk der EKD in Deutschland e.V.）、Pro Familia 德國性諮詢與家庭計畫協會（Pro Familia Deutsche Gesellschaft für Sexualberatung und Familienplanung e.V.）、工人福利聯邦協會（Arbeiterwohlfahrt -- Bundesverband e.V.）、德國青年與婚姻諮詢工作共同體（Deutsche Arbeitsgemeinschaft für Jugend und Eheberatung e.V.）、烏爾姆問題懷孕諮詢所（Ulmer Beratungsstelle für Problemschwangerschaften e.V.）、漢諾威青年與婚姻諮詢工作共同體（Hannoversche Arbeitsgemeinschaft für Jugend und Eheberatung e.V.）、斯圖加特邦首府準母親社會諮詢處（Soziale Beratungsstelle der Landeshauptstadt Stuttgart für werdende Mütter）、杜塞道夫社會醫學家庭諮詢處（Sozialmedizinische Familienberatung in Düsseldorf）、AOK 聯邦總會、聯邦醫師公會及德國女醫師協會（Deutscher Ärztinnenbund e.V.）。
\end{transpara}

\bisubsection{C.}{C.}

\bisubsubsection{I.}{I.}

\begin{sourcepara}[title={原文理由 C.I}]
\sourcern{109}Die Bayerische Staatsregierung (2 \abbrBvF{} 4/92) und 249 Mitglieder des Deutschen Bundestages (2 \abbrBvF{} 5/92) haben gemäß \abbrArt{} 93 \abbrAbs{} 1 \abbrNr{} 2 \abbrGG{}\ggfnArtNinetyThreeOneTwoDE{}, § 13 \abbrNr{} 6 \abbrBVerfGG{} den Antrag auf verfassungsrechtliche Überprüfung der \abbrArt{} 13 \abbrNr{} 1 und 15 \abbrNr{} 2 \abbrSFHG{} gestellt. Sie halten die in \abbrArt{} 13 \abbrNr{} 1 geregelte Neufassung des § 218a \abbrAbs{} 1 \abbrStGB{} und des § 219 \abbrStGB{} (Beratung der Schwangeren in einer Not- und Konfliktlage) sowie die in \abbrArt{} 15 \abbrNr{} 2 vorgesehene Aufhebung des \abbrArt{} 4 des 5. \abbrStrRG{} (Bundesstatistik) für unvereinbar mit dem Grundgesetz, weil diese Vorschriften gegen \abbrArt{} 2 \abbrAbs{} 2 Satz 1 in Verbindung mit \abbrArt{} 1 \abbrAbs{} 1 \abbrGG{}\ggfnArtOneOneTwoTwoOneDE{} verstießen.
\end{sourcepara}
\begin{transpara}
\transrn{109}巴伐利亞邦政府（2 BvF 4/92）及德國聯邦議會249名成員（2 BvF 5/92），依《基本法》第93條第1項\ggfnArtNinetyThreeOneTwoZH{}第2款及《聯邦憲法法院法》第13條第6款，聲請憲法審查 SFHG 第13條第1款及第15條第2款。其認為，第13條第1款所規定之《刑法》第218a條第1項及第219條（困境與衝突狀態中之懷孕婦女諮詢）新版本，以及第15條第2款所規定對《第五刑法改革法》第4條（聯邦統計）之廢止，均與《基本法》不相容，因為此等規定違反《基本法》第2條第2項第1句結合第1條第1項\ggfnArtOneOneTwoTwoOneZH{}。
\end{transpara}

\begin{sourcepara}[title={原文理由 C.I} Rn 110]
\sourcern{110}Die Bayerische Staatsregierung hält darüber hinaus aus denselben Gründen die Sicherstellungsverpflichtung unter \abbrArt{} 15 \abbrNr{} 2 \abbrSFHG{} (Einrichtungen zur Vornahme von Schwangerschaftsabbrüchen) und die Bestimmung des § 24b \abbrSGBV{} in der Fassung des \abbrArt{} 2 \abbrSFHG{} für verfassungswidrig. Sie ist weiter der Meinung, daß es insoweit an der Gesetzgebungszuständigkeit des Bundes fehle. Sie hat zur Stützung ihres Standpunktes auch ein Rechtsgutachten von \abbrProf{} \abbrDr{} Kriele zum Thema der nicht-therapeutischen Abtreibung vor dem Grundgesetz vorgelegt.
\end{sourcepara}
\begin{transpara}
\transrn{110}此外，巴伐利亞邦政府基於相同理由，亦認為 SFHG 第15條第2款中之確保義務（實施終止妊娠之機構），以及 SFHG 第2條版本之 SGB V 第24b條規定違憲。其進一步認為，就此而言，聯邦欠缺立法權限。為支持其立場，巴伐利亞邦政府亦提出 Kriele 教授關於非治療性終止妊娠在《基本法》下地位之法律意見書。
\end{transpara}

\bisubsubsection{II.}{II.}

\begin{sourcepara}[title={原文理由 C.II.1 開頭}]
\sourcern{111}Zur Begründung wird im wesentlichen vorgetragen:
\end{sourcepara}
\begin{transpara}
\transrn{111}理由主要陳述如下：
\end{transpara}

\begin{sourcepara}[title={原文理由 C.II.1 開頭} Rn 112]
\sourcern{112}1. Das Grundgesetz stelle in \abbrArt{} 2 \abbrAbs{} 2 Satz 1 in Verbindung mit \abbrArt{} 1 \abbrAbs{} 1 das sich im Mutterleib entwickelnde Leben unter den Schutz des Staates. Die Schutzpflicht habe nicht das Leben als Abstraktum zum Gegenstand, sondern die je individuelle und einzigartige Existenz jedes einzelnen Menschen. Der so geschützte Mensch existiere als einzigartiges Individuum nicht erst mit der Geburt, sondern auch schon vorher.
\end{sourcepara}
\begin{transpara}
\transrn{112}1. 《基本法》第2條第2項第1句結合第1條第1項\ggfnArtOneOneTwoTwoOneZH{}，將在母體內發展之生命置於國家保護之下。保護義務之標的並非抽象之生命，而是每一個人各自個別且獨特之存在。受此保護之人，作為獨特個體，並非出生後才存在，而是在出生以前即已存在。
\end{transpara}

\begin{sourcepara}[title={原文理由 C.II.1 開頭} Rn 113]
\sourcern{113}Die verschiedenen Regelungskonzepte für den Schwangerschaftsabbruch (generelle Rechtmäßigkeits-Wertung; eingeschränkte Unrechtswertung) könnten nicht nur begriffen werden als je besondere gesetzgeberische "Wege", um den ungeborenen Menschen "möglichst wirksam" zu schützen. Die Verfassung gestatte es dem Gesetzgeber nicht, ein Konzept genereller Legalisierung des Schwangerschaftsabbruchs einzusetzen, um ungeborenes Leben im ganzen besser zu schützen; denn der Verzicht auf das grundgesetzlich gebotene Unrechtsurteil bedeute Verzicht auf den dem einzelnen ungeborenen Menschen durch die Verfassung unmittelbar verliehenen Schutz- und Achtungsanspruch. Sogar für den verfassungsändernden Gesetzgeber seien die Verbürgungen der einzelnen Grundrechte insoweit einer Einschränkung entzogen, als sie zur Aufrechterhaltung einer dem \abbrArt{} 1 \abbrAbs{} 1 und 2 \abbrGG{} entsprechenden Ordnung unverzichtbar seien. Die generelle Ausmerzung des Unrechtsgehalts von Tötungshandlungen greife in diesen Kernbereich ein, weil dadurch der elementarste rechtliche Schutz für das bedrohte Rechtsgut preisgegeben werde.
\end{sourcepara}
\begin{transpara}
\transrn{113}終止妊娠之各種規制構想（一般合法性評價；受限制之不法評價），不能僅被理解為立法者為「盡可能有效」保護未出生人類所採之各種特殊「路徑」。憲法並不允許立法者為整體上更好保護未出生生命，而採用一般合法化終止妊娠之構想；因為放棄《基本法》所要求之不法判斷，即意味放棄憲法直接賦予每一個別未出生人類之保護及尊重請求。即使對修憲立法者而言，各個基本權之保障在其為維持符合《基本法》第1條第1項\ggfnArtOneOneZH{}及第2項之秩序所不可或缺之範圍內，亦不得受限制。一般性消除殺害行為之不法內涵，即侵入此一核心領域，因為此舉放棄了對受威脅法益最基本之法律保護。
\end{transpara}

\begin{sourcepara}[title={原文理由 C.II.1 開頭} Rn 114]
\sourcern{114}Nach wie vor sei das grundsätzliche Unrechtsurteil -- nach Auffassung der Bayerischen Staatsregierung auch eine grundsätzliche und zeitlich nicht eingeschränkte Bedrohung mit Strafe -- neben allen erforderlichen Beratungs- und Hilfsangeboten ein notwendiges und geeignetes Mittel zum Schutz des ungeborenen Lebens. Es habe Einfluß auf die Wertvorstellungen und die Verhaltensweisen der Bevölkerung. Diese rechtsethische Signalwirkung sei auch an anderer Stelle für den Schutz bedeutsamer Rechtsgüter in Anspruch genommen worden (Umweltstrafrecht, Embryonenschutz), ersichtlich ohne Rücksicht darauf, ob in der Praxis eine realistische Chance der strafrechtlichen Verfolgung bestehe.
\end{sourcepara}
\begin{transpara}
\transrn{114}基本之不法判斷，依巴伐利亞邦政府見解並且包括原則上且時間不受限制之刑罰威嚇，迄今仍與所有必要之諮詢及協助提供並列，為保護未出生生命所必要且適當之手段。其對人民之價值觀及行為方式具有影響。此一法律倫理訊號作用，在其他領域亦曾為保護重要法益而被運用（環境刑法、胚胎保護），且顯然不取決於實務上是否存在刑事追訴之現實可能。
\end{transpara}

% --- Servat 原文層（自 C.II.1 後續起） ---
\begin{sourcepara}[title={原文C.1}]
\sourcern{115}Sollte § 218a \abbrAbs{} 1 \abbrStGB{} \abbrNF{} -- so führt die Bayerische Staatsregierung ergänzend aus -- als verfassungsrechtlich unbedenklich beurteilt werden, so werde im Beitrittsgebiet das Konzept der dort seit 1972 geltenden Fristenregelung letztlich bestätigt. Der Gesetzgeber verzichte sonach auf den Versuch, mit seinen Mitteln ein rechtliches Bewußtsein vom Wert und verfassungsrechtlichen Schutz des ungeborenen Lebens in der Bevölkerung der neuen Länder zu schaffen. Spezifische Gefahren drohten dem ungeborenen Leben auch durch die medizinisch- pharmazeutische Entwicklung. Käme zur befristeten strafrechtlichen Freigabe der Abtreibung die Zulassung des Hormonpräparates RU 486 in Deutschland hinzu, so ergäbe sich daraus eine Kombination rechtlicher und medizinisch-organischer Erleichterungen des Schwangerschaftsabbruchs. Die Verbesserung der pränatalen Diagnostik habe zur Folge, daß die Eltern häufig schon jetzt, in jedem Falle aber in absehbarer Zeit innerhalb der ersten zwölf Wochen feststellen lassen könnten, ob das erwartete Kind in jeder Hinsicht gesund sei. Lasse die Frau wegen einer festgestellten Schädigung der Leibesfrucht die Schwangerschaft innerhalb der ersten zwölf Wochen seit Empfängnis abbrechen, so sei dies "nicht rechtswidrig", ohne daß es darauf ankomme, ob die Schädigung des Gesundheitszustandes behebbar sei oder so schwer wiege, daß von der Schwangeren die Fortsetzung der Schwangerschaft nicht verlangt werden könne. Damit werde ein Schwangerschaftsabbruch aus rein eugenischen Gründen möglich. Entsprechend den Erfahrungen aus den Vereinigten Staaten sei auch zu befürchten, daß in Zukunft eine größere Zahl von Frauen einen Schwangerschaftsabbruch verlange, weil das ungeborene Kind das nicht gewünschte Geschlecht habe. Der Arzt könne diesem Wunsch nicht entgegenhalten, der Eingriff sei rechtswidrig; auch derjenige, der einen solchen Abbruch öffentlich empfehle, befinde sich im Einklang mit der Rechtsordnung.
\end{sourcepara}

\begin{transpara}
\transrn{115}巴伐利亞邦政府補充指出，倘若新版本《刑法》第218a條第1項被評價為憲法上無疑義，則加入地區自1972年以來施行之期限規定構想，終究將獲確認。如此一來，立法者即放棄嘗試以其手段，在新邦人民之中形成關於未出生生命之價值及其憲法保護的法律意識。未出生生命亦因醫學及藥學發展而面臨特殊危險。若在有期限之刑法上開放終止妊娠之外，又在德國准許荷爾蒙製劑 RU 486，則將形成終止妊娠在法律上及醫學有機層面上之便利化的結合。產前診斷之改善，其結果是父母往往已經能夠，無論如何在可預見期間內也將能夠，於最初十二週內確認預期出生之子女是否在各方面健康。若婦女因確認胎兒受有損害，而於受孕後最初十二週內終止妊娠，則此即為「不具違法性」，而不問其健康狀態之損害是否可治癒，或是否嚴重到不能要求懷孕婦女繼續懷孕。如此即可能基於純粹優生學理由終止妊娠。依據美國經驗，亦須擔憂未來將有較多婦女因未出生子女之性別不合意而要求終止妊娠。醫師不能以該手術違法為由拒絕此一要求；公開建議為此種終止妊娠者，亦處於與法秩序一致之狀態。
\end{transpara}

\begin{sourcepara}[title={原文C.1} Rn 116]
\sourcern{116}Verfassungsrechtlich komme eine Rechtfertigung des Schwangerschaftsabbruchs nur im Einzelfall aufgrund einer Interessen- und Güterabwägung in Betracht. An dieser fehle es im Fall des § 218a \abbrAbs{} 1 \abbrStGB{} \abbrNF{} Nach dieser Vorschrift würden alle Fälle legalisiert, in denen die Schwangere den Abbruch der Schwangerschaft -- aus welchen Gründen auch immer -- verlange. Eine Beschränkung auf rechtfertigende Ausnahmetatbestände enthalte das Gesetz insoweit nicht. Das Vorliegen einer Not- und Konfliktlage werde an keiner Stelle zur Voraussetzung eines nicht rechtswidrigen Schwangerschaftsabbruchs in den ersten zwölf Wochen gemacht, sondern in der Neufassung der §§ 218a \abbrAbs{} 1 und 219 \abbrAbs{} 1 Satz 2 \abbrStGB{} lediglich generell unterstellt. Auch stelle das regelmäßige Bestehen einer schwierigen Lebenssituation noch keinen zureichenden Rechtfertigungsgrund dar. Es müsse eine außergewöhnliche Belastung im Einzelfall vorliegen, welche das Austragen des Kindes für die Frau als wirklich unzumutbare Härte erscheinen lasse. Das Gesetz fordere jedoch nicht einmal, daß die Frau, die den Schwangerschaftsabbruch verlange, zumindest subjektiv das Austragen des Kindes als unzumutbare außergewöhnliche Belastung empfinde. Die These, Frauen trieben nicht "leichtfertig" ab, erfasse nur einen Teil der Wirklichkeit. Nicht wenige Frauen betrachteten den Schwangerschaftsabbruch als Teil ihrer persönlichen, rechtlich nicht einschränkbaren Freiheit. Die verhältnismäßig hohe Zahl von Mehrfachabbrüchen in Rechtsordnungen mit sogenannter Fristenlösung deute überdies ebenso wie die öffentlichen Kampagnen für die Abtreibung darauf hin, daß der Schwangerschaftsabbruch auch als Mittel der Familienplanung verstanden und praktiziert werde.
\end{sourcepara}
\begin{transpara}
\transrn{116}在憲法上，終止妊娠之正當化僅能於個案中，基於利益與法益衡量而成立。新版本《刑法》第218a條第1項之情形欠缺此種衡量。依該規定，凡懷孕婦女要求終止妊娠者，無論基於何種理由，均被合法化。法律就此並未包含限於正當化例外構成要件之限制。在最初十二週內，困境與衝突狀態之存在，在任何地方都未被作為不具違法性之終止妊娠之前提，而僅在新版本《刑法》第218a條第1項及第219條第1項第2句中被一般性推定。通常存在艱難生活處境，亦尚不足以構成充分之正當化理由。個案中必須存在異常負擔，使為婦女繼續懷孕真正顯得為不可期待之嚴苛。惟法律甚至未要求提出終止妊娠要求之婦女，至少主觀上感受繼續懷孕係不可期待之異常負擔。婦女並非「輕率」終止妊娠此一命題，只掌握了現實之一部分。為數不少之婦女將終止妊娠視為其個人自由之一部分，法律不得對其加以限制。此外，在所謂期限規定之法秩序中，重複終止妊娠之數量相對較高，與支持終止妊娠之公共運動一樣，均顯示終止妊娠也被理解並實行為家庭計畫之手段。
\end{transpara}

\begin{sourcepara}[title={原文C.3}]
\sourcern{117}Das Gesetz versage der schwangeren Frau jeden Maßstab dafür, wann ein Austragen der Schwangerschaft unzumutbar sei. Es lasse damit gerade jene Frauen im Stich, die durch Einwirkungen ihrer Umwelt (Eltern, Kindesvater, Arbeitgeber) zur Abtreibung gedrängt werden, und zwar in einem Zeitraum, in dem die Schwangere solchen Pressionen besonders intensiv ausgesetzt sei. Das Argument, auf Maßstäbe für die Unzumutbarkeit müsse gänzlich verzichtet werden, weil sonst in der Beratung Gewissensnöte vorgetäuscht würden und die "Kommunikation zum Ritual" entarte, überzeuge nicht. Es sei nicht einsichtig, wieso dadurch, daß der Abbruch generell für nicht rechtswidrig erklärt werde, zu einer "offeneren" Beratungsatmosphäre beigetragen werde; denn die Frau sei im Falle der Beratung schon nach geltendem Recht von der Strafdrohung ausgenommen.
\end{sourcepara}
\begin{transpara}
\transrn{117}法律未提供懷孕婦女任何標準，使其得以判斷何時繼續懷孕為不可期待。法律因而正是棄置那些受其環境影響（父母、孩子之父、雇主）而被迫終止妊娠之婦女，而此一時期正是懷孕婦女特別強烈暴露於此等壓力之下的時期。主張必須完全放棄不可期待性標準，否則在諮詢中將會偽裝良心困境，並使「溝通淪為儀式」之論點，並不具說服力。難以理解的是，為何透過將終止妊娠一般性宣告為不具違法性，即能促成一種「更開放」之諮詢氛圍；因為依現行法，婦女在諮詢情形下即已被排除於刑罰威嚇之外。
\end{transpara}

\begin{sourcepara}[title={原文C.4}]
\sourcern{118}Mit der Qualifikation als "nicht rechtswidrig" in § 218a \abbrAbs{} 1 \abbrStGB{} \abbrNF{} gebe der Gesetzgeber der gesamten Rechtsordnung die Grundwertung vor. Es sei klar, daß die genannte Bestimmung den unmittelbaren Anschluß an § 24b \abbrSGBV{} suche und finde. Die Leistungen der Krankenkasse würden damit notwendigerweise auch für jene Schwangerschaftsabbrüche gewährt, die aus Gründen erfolgten, die vor dem Grundgesetz keinen Bestand hätten. Schwere neuartige Konflikte ergäben sich ferner im ärztlichen Berufsrecht und im Organisationsrecht.
\end{sourcepara}
\begin{transpara}
\transrn{118}立法者以新版本《刑法》第218a條第1項中「不具違法性」之定性，為整體法秩序預設基本評價。很清楚的是，該規定尋求並取得與 SGB V 第24b條之直接銜接。由此，醫療保險基金之給付也必然將提供給那些基於在《基本法》面前不能成立之理由所為之終止妊娠。此外，在醫師職業法及組織法上，亦將產生新型且嚴重之衝突。
\end{transpara}

\begin{sourcepara}[title={原文C.5}]
\sourcern{119}Die Zurücknahme des Unrechtsurteils werde nicht durch die im Schwangeren- und Familienhilfegesetz vorgesehenen sozialpolitischen Maßnahmen aufgewogen. Eine rechtliche Mißbilligung des Schwangerschaftsabbruchs komme darin nicht zum Ausdruck. Dies könne im Sozialrecht allenfalls dadurch geschehen, daß der Gesetzgeber nicht in jedem Falle eines Schwangerschaftsabbruchs soziale Vorteile gewähre.
\end{sourcepara}
\begin{transpara}
\transrn{119}不法判斷之撤回，並未被《孕婦及家庭扶助法》所規定之社會政策措施所抵銷。此等措施並未表達對終止妊娠之法律否定評價。在社會法中，這至多只能透過立法者並非在每一終止妊娠案件均給予社會利益而發生。
\end{transpara}

\begin{sourcepara}[title={原文C.6}]
\sourcern{120}Im übrigen stehe die Verwirklichung der sozialpolitischen Maßnahmen in wesentlichen Punkten noch aus und sei -- wie die Verweisung auf den "Finanzausgleich" zeige -- in hohem Maße ungewiß.
\end{sourcepara}
\begin{transpara}
\transrn{120}此外，社會政策措施在重要部分之實現仍尚未完成，且如其指向「財政均衡」所示，具有高度不確定性。
\end{transpara}

\begin{sourcepara}[title={原文C.7}]
\sourcern{121}2. Die Beratung solle die Schutzfunktion übernehmen, die im Indikationenmodell die Feststellung der rechtfertigenden Indikationstatbestände erfülle. An die Stelle objektivierter Kontrolle träten prozedurale Einwirkungen auf den nicht kontrollierbaren Entscheidungsprozeß. Die Beratung bilde damit im gesetzlichen Konzept den "zentralen Punkt".
\end{sourcepara}
\begin{transpara}
\transrn{121}2. 諮詢應承擔在指標模式中由正當化指標事由構成要件之確認所履行之保護功能。取代客觀化控制者，是對不可控制之決定過程施加程序性影響。諮詢因而在法律構想中構成「核心點」。
\end{transpara}

\begin{sourcepara}[title={原文C.8}]
\sourcern{122}Daraus folge, daß die Beratung Pflicht sein müsse. Sie dürfe sich auch nicht auf eine bloße Information über Leistungseinrichtungen und Rechtsansprüche beschränken, sondern müsse darauf gerichtet sein, die Frau zu ermutigen, die Schwangerschaft auszutragen. Dazu müsse die Frau ihre Not- und Konfliktlage darlegen und die Gründe benennen, die sie zu einem Schwangerschaftsabbruch bestimmten. Von einer Beratung könne jedenfalls dann keine Rede sein, wenn sich die Schwangere vollständig ausschweige. Der plausible Gedanke, nur eine Beratung "ohne Druck" könne eine gewisse Aussicht auf lebensschützenden Erfolg haben, könne auch nicht bedeuten, daß die schwangere Frau nicht mit der Bewertung des Schwangerschaftsabbruchs als Unrecht konfrontiert werden dürfe. Erforderlich sei ferner, die beratenden Instanzen und Personen durch normative und institutionelle Vorkehrungen dazu anzuhalten, daß sie die Beratung entsprechend den verfassungsrechtlichen und gesetzlichen Vorgaben durchführen. Dies wiederum setze ein Minimum an Protokollierung des Beratungsvorganges voraus.
\end{sourcepara}
\begin{transpara}
\transrn{122}由此可知，諮詢必須是義務。諮詢亦不得限於僅告知給付機構及法律請求權，而必須以鼓勵婦女繼續懷孕為目標。為此，婦女必須陳述其困境與衝突狀態，並說明促使其終止妊娠之理由。至少在懷孕婦女完全保持沉默時，即無從談論諮詢。唯有「無壓力」之諮詢才可能對生命保護之成功有一定希望，此一看似有理之想法，也不能意味懷孕婦女不得被告知終止妊娠作為不法之評價。此外，亦有必要透過規範及制度性預防安排，促使諮詢機關及人員依憲法與法律要求實施諮詢。這又以對諮詢過程作成最低限度之紀錄為前提。
\end{transpara}

\begin{sourcepara}[title={原文C.9}]
\sourcern{123}Diesen verfassungsrechtlichen Anforderungen genüge die in § 219 \abbrAbs{} 1 \abbrStGB{} \abbrNF{} geregelte Beratung nicht. Sie sei auf das Prinzip der Selbstbestimmung der Frau ausgerichtet. Obwohl in § 219 nicht weniger als neunmal der Begriff der "Beratung" verwendet werde, würden in der Sache nur Informationspflichten begründet. Im Gesetz sei nicht einmal gesagt, daß Gegenstand der Beratung die Not- und Konfliktlage sein solle, in der sich die Schwangere befinde. Die Beratung werde im Gesetz auch nicht als Gespräch definiert. Des weiteren fehle die Vorgabe des Beratungsziels, die Schwangere zur Fortsetzung der Schwangerschaft zu ermutigen. Es werde nur die Erwartung des Gesetzgebers zum Ausdruck gebracht, daß die Beratung dem Lebensschutz diene. Der Kern der Vorschrift liege im Absatz 1 Satz 3, nach dem die Beratung dazu dienen solle, die Schwangere in die Lage zu versetzen, eine "verantwortungsbewußte eigene Gewissensentscheidung" zu treffen. Mit dieser verfehlten Umschreibung werde ein falsches Vorverständnis erzeugt, das zu Denk- und Argumentationsverboten führe und geeignet sei, die unkontrollierbare Entscheidung über den Schwangerschaftsabbruch mit dem Schein einer verfassungsrechtlich geschützten Gewissensentscheidung zu umgeben. Des weiteren fehle eine gesetzliche Verpflichtung der Schwangeren zur Darlegung ihrer persönlichen Notlage ebenso wie eine auch nur minimale Pflicht zur Protokollierung der Beratung.
\end{sourcepara}
\begin{transpara}
\transrn{123}新版本《刑法》第219條第1項所規定之諮詢，並不符合此等憲法要求。該諮詢係以婦女自我決定原則為取向。儘管第219條使用「諮詢」一詞不少於九次，但實質上僅創設資訊義務。法律甚至未說明，諮詢之標的應為懷孕婦女所處之困境與衝突狀態。法律亦未將諮詢定義為對話。再者，法律欠缺諮詢目標之預設，亦即鼓勵懷孕婦女繼續懷孕。法律僅表達立法者之期待，即諮詢服務於生命保護。該規定之核心在於第1項第3句，依其規定，諮詢應使懷孕婦女能夠作成「負責任之自身良心決定」。此一錯誤之描述製造出錯誤之前理解，導致思考與論證禁令，並適於使不可控制之終止妊娠決定披上受憲法保護之良心決定外觀。再者，法律既欠缺懷孕婦女陳述其個人困境之義務，也欠缺哪怕只是最低限度之諮詢紀錄義務。
\end{transpara}

\begin{sourcepara}[title={原文C.10}]
\sourcern{124}3. Die Verfassungswidrigkeit und Nichtigkeit zumindest des § 218a \abbrAbs{} 1 und des § 219 \abbrAbs{} 1 und \abbrAbs{} 3 Satz 1 \abbrStGB{} \abbrNF{} könne nicht isoliert festgestellt werden. Sie führe aus den Gründen des Urteils des Bundesverfassungsgerichts vom 4. August 1992 zur Nichtigkeit des \abbrArt{} 13 \abbrNr{} 1 des \abbrSFHG{} insgesamt.
\end{sourcepara}
\begin{transpara}
\transrn{124}3. 至少新版本《刑法》第218a條第1項及第219條第1項、第3項第1句之違憲與無效，不能孤立確認。依聯邦憲法法院1992年8月4日判決之理由，此將導致 SFHG 第13條第1款整體無效。
\end{transpara}

\begin{sourcepara}[title={原文C.11}]
\sourcern{125}4. Die Beibehaltung der Bundesstatistik über Schwangerschaftsabbrüche (vgl. \abbrArt{} 4 des 5. \abbrStrRG{} \abbrAF{}) sei unter dem Gesichtspunkt der Nachbesserungspflicht des Gesetzgebers verfassungsrechtlich geboten.
\end{sourcepara}
\begin{transpara}
\transrn{125}4. 維持關於終止妊娠之聯邦統計（參見舊版本《第五刑法改革法》第4條），從立法者補正義務之觀點而言，為憲法所要求。
\end{transpara}

\begin{sourcepara}[title={原文C.12}]
\sourcern{126}5. Zur Regelung der in \abbrArt{} 15 \abbrNr{} 2 \abbrSFHG{} enthaltenen Sicherstellungsverpflichtung fehlt dem Bund nach Auffassung der Bayerischen Staatsregierung die Gesetzgebungskompetenz. Die Sicherstellungsverpflichtung des \abbrArt{} 15 \abbrNr{} 2 \abbrSFHG{} verstoße auch inhaltlich gegen das Grundgesetz und sei deshalb nichtig. Sie gehe über eine Verpflichtung der obersten Landesbehörde, im Rahmen der rechtlichen und tatsächlichen Möglichkeiten darauf hinzuwirken, daß ein ausreichendes und flächendeckendes Angebot an Einrichtungen zur Vornahme von Schwangerschaftsabbrüchen zur Verfügung stehe, weit hinaus. Der Gesetzgeber verpflichte eine bestimmte Landesbehörde und mit ihr das Land zu einer Leistung, die rechtlich und tatsächlich nicht möglich oder unzumutbar sei. Dies sei ein Verstoß gegen das Rechtsstaatsprinzip (\abbrArt{} 20 \abbrAbs{} 3, \abbrArt{} 28 \abbrAbs{} 1 Satz 1 \abbrGG{}\ggfnArtTwentyEightOneDE{}) und gegen das Prinzip der Bundestreue.
\end{sourcepara}
\begin{transpara}
\transrn{126}5. 依巴伐利亞邦政府見解，關於 SFHG 第15條第2款所含確保義務之規制，聯邦欠缺立法權限。SFHG 第15條第2款之確保義務，在內容上亦違反《基本法》，因而無效。其遠遠超出最高邦機關在法律及事實可能範圍內促成提供充分且遍及全境之終止妊娠機構之義務。立法者使特定邦機關並連同該邦負擔一項在法律上及事實上不可能或不可期待之給付。此係違反法治國原則（《基本法》第20條第3項\ggfnArtTwentyThreeZH{}、第28條第1項第1句\ggfnArtTwentyEightOneZH{}）及聯邦忠誠原則。
\end{transpara}

\begin{sourcepara}[title={原文C.13}]
\sourcern{127}6. Die Bayerische Staatsregierung hält schließlich § 24b \abbrSGBV{} in der Fassung des \abbrArt{} 2 \abbrSFHG{} aus den im Verfahren 2 \abbrBvF{} 2/90 vorgebrachten kompetenzrechtlichen und materiellen Gründen für verfassungswidrig.
\end{sourcepara}
\begin{transpara}
\transrn{127}6. 巴伐利亞邦政府最後認為，SFHG 第2條版本之 SGB V 第24b條，基於其在 2 BvF 2/90 程序中所提出之權限法及實體法理由，亦屬違憲。
\end{transpara}

\bisubsubsection{III.}{III.}

\begin{sourcepara}[title={原文C.14}]
\sourcern{128}Von den nach § 77 \abbrBVerfGG{} zur Äußerung Berechtigten haben Stellung genommen: der Deutsche Bundestag, der sich ergänzend auf ein von \abbrProf{} \abbrDr{} Eser erstelltes Rechtsgutachten bezieht,sowie -- in einer gemeinsamen Stellungnahme -- die Landesregierungen von Brandenburg, Bremen, Hamburg, Hessen, Niedersachsen, Nordrhein-Westfalen, Rheinland-Pfalz, Saarland und Schleswig-Holstein. Sie halten die Anträge für unbegründet; von einer Äußerung zur Bundesstatistik hat der Deutsche Bundestag abgesehen.
\end{sourcepara}
\begin{transpara}
\transrn{128}依《聯邦憲法法院法》第77條有權表示意見者中，已提出意見者包括：德國聯邦議會，其並補充援引 Eser 教授所作法律意見書；以及以共同意見書提出意見之布蘭登堡、不來梅、漢堡、黑森、下薩克森、北萊茵-西伐利亞、萊茵蘭-普法爾茨、薩爾及石勒蘇益格-荷爾斯泰因各邦政府。其等認為聲請無理由；德國聯邦議會未就聯邦統計表示意見。
\end{transpara}

\begin{sourcepara}[title={原文C.15}]
\sourcern{129}1. Der Gesetzgeber des Schwangeren- und Familienhilfegesetzes sei von einer umfassenden Schutzpflicht des Staates für das Leben, einschließlich des vorgeburtlichen, ausgegangen, habe sich auf die Erkenntnis, daß ein wirksamer Lebensschutz; nicht allein durch Strafdrohungen erreicht werden könne, gestützt und das Prinzip "Hilfe statt Strafe" zugrunde gelegt. Diese Strategie verspreche jedenfalls mittel- und längerfristig einen deutlich besseren Lebensschutz als ein bloßes Abschreckungsstrafrecht, welches die Frau zur unmündigen Normunterworfenen herabstufe; es erreiche auch einen höheren Grad rechtsethischer Integration.
\end{sourcepara}
\begin{transpara}
\transrn{129}1. 《孕婦及家庭扶助法》之立法者係以國家對生命，包括出生前生命，負有全面保護義務為出發點，並基於有效生命保護不能單靠刑罰威嚇達成之認識，以「協助而非刑罰」原則為基礎。此一策略無論如何在中長期而言，比單純之威嚇刑法更可期待明顯較佳之生命保護；後者將婦女貶為無判斷能力之規範服從者。此一策略亦達成較高程度之法律倫理整合。
\end{transpara}

\begin{sourcepara}[title={原文C.16}]
\sourcern{130}Biologisch gesehen sei die Entwicklung menschlichen Lebens ein Kontinuum, das mit der Kernverschmelzung von Ei und Samenzelle beginne und mit dem Tode des Individuums ende. Weder die Nidation noch die Geburt erschienen bei dieser Betrachtung als Einschnitte. Dennoch dürfe der Gesetzgeber an verschiedene Entwicklungsphasen unterschiedliche Rechtsfolgen knüpfen und müsse dies unter Umständen auch tun. Soweit der vorgeburtliche menschliche Entwicklungsprozeß überhaupt Einschnitte aufweise, seien diese im Blick auf die Beziehung zwischen der Mutter und dem Ungeborenen zu ermitteln. Die Schwangere nehme das ungeborene Leben erst allmählich wahr. Erst von dieser Wahrnehmung an könne von der Schwangeren die Entscheidung für die Annahme der Leibesfrucht verlangt werden. Erst mit der Annahme gewinne die Beziehung der Schwangeren zur Leibesfrucht Verbindlichkeit. Dies gelte schon nach bisherigem Recht. Die angegriffene Regelung gehe von denselben Fristen und denselben Wirkungen des Fristablaufs aus. Die bisherige Regelung für den Schwangerschaftsabbruch habe normativ die Unzumutbarkeit der Annahme vorausgesetzt. Die Neufassung des § 218a \abbrStGB{} ändere dies nur insoweit, als sie die kriminologische und die soziale Indikation entfallen und an ihre Stelle die "informierte Entscheidung" der Schwangeren nach entsprechender Beratung treten lasse.
\end{sourcepara}
\begin{transpara}
\transrn{130}從生物學觀點而言，人類生命之發展是一個連續體，始於卵子與精子之核融合，終於個體死亡。在此觀察下，著床與出生均不顯現為斷點。儘管如此，立法者仍得將不同法律效果繫於不同發展階段，並且在特定情形下亦必須如此。出生前人類發展過程若竟具有斷點，亦應從母親與未出生者之關係出發加以確定。懷孕婦女只是逐漸感知未出生生命。只有自此一感知起，始能要求懷孕婦女決定接受胎兒。也只有隨著接受，懷孕婦女與胎兒之關係才取得拘束性。這在既有法律下已然如此。受攻擊之規制以相同期限及期限屆滿之相同效果為出發點。既有終止妊娠規制在規範上預設接受之不可期待性。《刑法》第218a條新版本僅在以下範圍改變此點：其取消犯罪學指標事由及社會指標事由，並以懷孕婦女經相應諮詢後之「知情決定」取而代之。
\end{transpara}

\begin{sourcepara}[title={原文C.17}]
\sourcern{131}Das Bundesverfassungsgericht habe immer wieder betont, daß es Sache des Gesetzgebers sei zu entscheiden, auf welche Weise er seine verfassungsrechtliche Pflicht, das Leben zu schützen, erfüllen wolle. Im Verfahren einer abstrakten Normenkontrolle könne eine "Vertretbarkeitskontrolle" oder sogar eine "intensivierte Inhaltskontrolle" erfolgen. In bezug auf die Abtreibungsproblematik habe das Bundesverfassungsgericht weiter vorgegeben, daß Strafe niemals Selbstzweck sein könne; sie sei die "ultima ratio" im Instrumentarium des Gesetzgebers, deren Einsatz dem Verhältnismäßigkeitsprinzip unterliege. Die Anwendung dieser Maximen erfordere sowohl vom Gesetzgeber als auch vom überprüfenden Gericht ein prognostisches Verhalten. Für die Anforderungen an eine verfassungsrechtlich zu akzeptierende Prognose habe das Bundesverfassungsgericht im Mitbestimmungsurteil Richtlinien entwickelt (vgl. \abbrBVerfGE{} 50, 290 [332 \abbrFF{}]). Diesen Anforderungen sei der Deutsche Bundestag gerecht geworden. Das angegriffene Gesetz realisiere eine neue und bessere Schutztechnik für das werdende Leben und bringe zugleich das Interesse an einer menschlich richtigen und normativ angemessenen rechtlichen Antwort auf die besondere Lage der betroffenen Frau zur Geltung.
\end{sourcepara}
\begin{transpara}
\transrn{131}聯邦憲法法院一再強調，立法者應以何種方式履行其憲法上保護生命之義務，係由立法者決定。在抽象規範審查程序中，可以進行「可支持性審查」，甚至進行「強化內容審查」。關於終止妊娠問題，聯邦憲法法院並進一步預設，刑罰絕不可能是自我目的；刑罰乃立法者工具箱中之「最後手段」，其運用受比例原則拘束。此等準則之適用，要求立法者及審查法院均採取預測性態度。就憲法上可接受之預測所須符合之要求，聯邦憲法法院已在共同決定判決中發展出準則（參見 BVerfGE 50, 290 [332 ff.]）。德國聯邦議會已符合此等要求。受攻擊之法律實現了對形成中生命之新且較佳保護技術，並同時使人性上正確且規範上適當回應受影響婦女特殊處境之利益得以實現。
\end{transpara}

\begin{sourcepara}[title={原文C.18}]
\sourcern{132}Das Schwangeren- und Familienhilfegesetz setze auf Hilfe, Beratung und Überzeugung und nicht auf die Einschüchterung der Frau durch Strafdrohung; es realisiere damit die prinzipielle Vorgabe, das Strafrecht erst als "ultima ratio" einzusetzen. Nach den vorliegenden sozialwissenschaftlichen Untersuchungen hätten die Strafnormen keine abschreckende Wirkung. Die strafrechtliche Aufklärung und Aburteilung von Schwangerschaftsabbrüchen bewege sich statistisch gegen Null. Das Strafrecht sei in diesem Bereich praktisch wirkungslos und wirke überdies selektiv. Begünstigt würden die Informierten, die sozial gut Ausgestatteten, die Wendigen. Diese Verhältnisse seien unerträglich. Der Gesetzgeber habe nicht mehr auf eine Verbesserung des Lebensschutzes in den ersten zwölf Schwangerschaftswochen durch die geltende Indikationenregelung hoffen können.
\end{sourcepara}
\begin{transpara}
\transrn{132}《孕婦及家庭扶助法》依靠協助、諮詢與說服，而非以刑罰威嚇恐嚇婦女；其因而實現了僅將刑法作為「最後手段」使用之原則性要求。依現有社會科學研究，刑罰規範並無嚇阻效果。終止妊娠之刑事偵查及判決，在統計上趨近於零。刑法在此領域實際上無效，且此外具有選擇性。受惠者是資訊充分、社會資源良好、善於迴避者。此等情況難以忍受。立法者已不能再期待藉由現行指標規則，在妊娠最初十二週改善生命保護。
\end{transpara}

\begin{sourcepara}[title={原文C.19}]
\sourcern{133}Die Strafdrohung könne die Chance gefährden, werdendes Leben zu bewahren. Sie könne die Frau, die sich des Vorliegens einer Indikation nicht sicher sei, davon abhalten, sich einer Beratung zu stellen. Bei anderen Frauen könne das Bestreben, eine Indikationsfeststellung um jeden Preis zu erhalten, zu nicht wirklichkeitsgetreuen Darstellungen ihrer Lage führen. Sie erlebten die Beratung als Prüfung, der man sich entziehen müsse, nicht als Hilfsangebot.
\end{sourcepara}
\begin{transpara}
\transrn{133}刑罰威嚇可能危及保存形成中生命之機會。其可能使不確定是否存在指標事由之婦女，不願接受諮詢。對其他婦女而言，不惜一切取得指標事由確認之努力，可能導致其對自身處境作出不符合現實之陳述。她們將諮詢經驗為一種必須逃避之考試，而非協助提供。
\end{transpara}

\begin{sourcepara}[title={原文C.20}]
\sourcern{134}Einer Grundlage entbehre der Vorwurf, § 218a \abbrStGB{} \abbrNF{} gehe in seiner verfassungswidrigen Zielsetzung und Fassung über den früheren § 218a \abbrStGB{} noch hinaus, den das Bundesverfassungsgericht in seinem Urteil vom 25. Februar 1975 für nichtig erklärt habe. In der Sache sei es damals wie heute um den Ausschluß der Rechtswidrigkeit gegangen. Wie das Bundesverfassungsgericht damals richtig erkannt habe (vgl. \abbrBVerfGE{} 39, 1 [54]), ließen sich mit der Rechtswidrigkeit des ärztlichen Eingriffs diejenigen Maßnahmen, welche mit diesem Eingriff zusammenhingen, weder rechtspolitisch noch rechtssystematisch vereinbaren. Deshalb schaffe § 218a \abbrStGB{} in der früheren und in der aktuellen Fassung einen Rechtfertigungsgrund. Die Formulierungen "nicht rechtswidrig" und "nicht strafbar" seien funktional gleichwertig. Verfassungsrechtlich sei die Fassung "nicht rechtswidrig" im Blick auf \abbrArt{} 103 \abbrAbs{} 2 \abbrGG{}\ggfnArtOneZeroThreeTwoDE{} sogar vorzuziehen.
\end{sourcepara}
\begin{transpara}
\transrn{134}關於新版本《刑法》第218a條在其違憲目標與版本上甚至超過聯邦憲法法院於1975年2月25日判決宣告無效之舊《刑法》第218a條的指摘，欠缺根據。實質上，當時與今日所涉及者均為排除違法性。正如聯邦憲法法院當時正確指出（參見 BVerfGE 39, 1 [54]），若醫師介入為違法，則與此介入相關之措施，無論從法律政策或法體系上均無法調和。因此，舊版本及現行版本之《刑法》第218a條均創設一項正當化事由。「不具違法性」與「不罰」兩種表述在功能上等值。從憲法上看，鑑於《基本法》第103條第2項\ggfnArtOneZeroThreeTwoZH{}，「不具違法性」之版本甚至較可取。
\end{transpara}

\begin{sourcepara}[title={原文C.21}]
\sourcern{135}Zu Unrecht werde der Vorwurf erhoben, das Schwangeren- und Familienhilfegesetz mißbillige den Schwangerschaftsabbruch nicht. Es sei abwegig, aus der strafrechtlichen Freistellung einer Handlung kurzerhand zu schließen, daß diese damit zu einem normalen sozialen Vorgang werde. Die rechtliche Mißbilligung ergebe sich zum einen daraus, daß der Gesetzgeber die Strafdrohung unmißverständlich überall dort eingesetzt habe, wo sie den Schutz des werdenden Lebens gewährleisten könne. Auf der anderen Seite böten \abbrArt{} 1 bis 12, \abbrArt{} 14 und 15 \abbrSFHG{} umfängliche Instrumente und Maßnahmen im Dienste des Lebensschutzes. Hier zeige der Gesetzgeber seine Mißbilligung dadurch, daß er alle Kräfte aufwende, um taugliche Gegenmittel zu schaffen und einzusetzen.
\end{sourcepara}
\begin{transpara}
\transrn{135}指摘《孕婦及家庭扶助法》並未否定評價終止妊娠，並不正確。由一項行為在刑法上被排除處罰，即逕行推論該行為因此成為正常社會過程，是荒謬的。法律否定評價一方面來自於，凡刑罰威嚇能夠確保形成中生命之保護處，立法者均明確使用刑罰威嚇。另一方面，SFHG 第1條至第12條、第14條及第15條提供了服務於生命保護之廣泛工具與措施。立法者在此透過投入所有力量以創設並運用適當對策，顯示其否定評價。
\end{transpara}

\begin{sourcepara}[title={原文C.22}]
\sourcern{136}Der Gesetzgeber -- so tragen die Landesregierungen ergänzend vor -- gehe davon aus, daß der Staat bedrohtes Leben nicht in allen Fällen schützen könne, und ziele deshalb darauf, daß er es in möglichst vielen Fällen schütze. Damit werde die Aufgabe, das einzelne Leben zu schützen, nicht in Frage gestellt.
\end{sourcepara}
\begin{transpara}
\transrn{136}各邦政府補充陳述，立法者係以國家不可能在所有情形下保護受威脅生命為出發點，因而以盡可能在更多情形中保護生命為目標。保護個別生命之任務，並未因此受到質疑。
\end{transpara}

\begin{sourcepara}[title={原文C.23}]
\sourcern{137}2. Wer auf Beratung und Hilfe und nicht nur auf Strafdrohung setze, habe zahlreiche und plausible Gründe für die Erwartung, er werde damit das ungeborene Leben besser schützen: Die Beratung sei eine Voraussetzung für eine verantwortliche Entscheidung der schwangeren Frau. Auf sie setze der Gesetzgeber als den wirksamsten Schutzschirm, den der vom Schwangerschaftsabbruch bedrohte Embryo habe.
\end{sourcepara}
\begin{transpara}
\transrn{137}2. 依靠諮詢與協助，而非僅依靠刑罰威嚇者，有諸多且合理之理由期待其將由此更好保護未出生生命：諮詢是懷孕婦女作成負責任決定之前提。立法者將諮詢視為受終止妊娠威脅之胚胎所擁有的最有效保護傘。
\end{transpara}

\begin{sourcepara}[title={原文C.24}]
\sourcern{138}Der Gesetzgeber überlasse das werdende Leben aber nicht einer Entscheidung, in der das Gewissen der Schwangeren sich aus Not, Hilflosigkeit, Bedrängnis und Angst nicht ausreichend Gehör verschaffe. Er habe Maßnahmen zur Verbesserung der sozialen und ökonomischen Situation der Schwangeren und der Mütter und zur Beseitigung oder doch Linderung von Not, Hilflosigkeit, Bedrängnis und Angst getroffen. Er habe auch Sorge dafür getragen, daß diese Hilfen der Frau nicht verborgen blieben. Die Beratung sei das Kernstück auch der strafrechtlichen Regelungen des Schwangeren- und Familienhilfegesetzes und vom Gesetzgeber im Gegensatz zum geltenden Recht allseits mit Strafe bewehrt worden.
\end{sourcepara}
\begin{transpara}
\transrn{138}但立法者並未將形成中生命交由一項決定處置，而在該決定中，懷孕婦女之良心因困境、無助、壓迫與恐懼而不能充分發聲。立法者已採取措施，以改善懷孕婦女與母親之社會及經濟處境，並消除或至少減輕困境、無助、壓迫與恐懼。立法者亦已確保此等協助不會對婦女隱而不見。諮詢同時也是《孕婦及家庭扶助法》刑法規制之核心，且與現行法相反，立法者已在各方面以刑罰加以保障。
\end{transpara}

\begin{sourcepara}[title={原文C.25}]
\sourcern{139}Die Pflicht der beratenden Person zur umfassenden Information der Schwangeren (§ 219 \abbrAbs{} 1 Satz 4, 5 \abbrStGB{} \abbrNF{}) gewährleiste, daß vor allem in bezug auf die sozialen und ökonomischen Bedingungen eine irrtumsfreie Entscheidung der Frau zustandekomme. In der Praxis habe sich gezeigt, daß eine Schwangerschaftskonfliktberatung der einzige Weg sei, die natürliche und unter Umständen auch moralisch fundierte Tendenz der Frau zur Austragung der Schwangerschaft zu verstärken. Selbst wenn Schwangere die Beratungsstellen mit dem als feststehend empfundenen Entschluß aufsuchten, die Schwangerschaft abzubrechen, könne eine solche Beratung den Konflikt aufarbeiten, die Festlegung auf den Schwangerschaftsabbruch lösen und so Lebensschutz bewirken. Die Konfliktberatung könne jedoch nicht vorgeschrieben und erzwungen, sondern nur angeboten werden. Zu diesem Angebot verpflichte § 219 \abbrAbs{} 1 Satz 2, 3 \abbrStGB{} \abbrNF{}
\end{sourcepara}
\begin{transpara}
\transrn{139}諮詢人員全面告知懷孕婦女之義務（新版本《刑法》第219條第1項第4句、第5句），確保婦女尤其就社會及經濟條件作成無錯誤之決定。實務顯示，懷孕衝突諮詢是強化婦女自然且在特定情形下亦有道德基礎之繼續懷孕傾向的唯一道路。即使懷孕婦女是在其感受為既定之終止妊娠決心下前往諮詢機構，此種諮詢仍能處理衝突，鬆動其對終止妊娠之固定決定，並由此產生生命保護效果。然而，衝突諮詢不能被命令及強制，而只能被提供。新版本《刑法》第219條第1項第2句、第3句即課予此一提供義務。
\end{transpara}

\begin{sourcepara}[title={原文C.26}]
\sourcern{140}Daß der Gesetzgeber die unter diesen Voraussetzungen getroffene Entscheidung der Frau als verantwortlich anerkenne, sei die zwingende Konsequenz aus einem Schutzkonzept, das auf Beratung und Hilfe statt auf Strafdrohung setze. Ein Beratungsmodell müsse, um schutzwirksam zu sein, auf der Nicht-Rechtswidrigkeit bestehen. Damit keine über den Abbruch nachdenkende Frau davon abgehalten werde, sich beraten zu lassen, müsse auf jede Form von Repression, Stigmatisierung und Diskriminierung verzichtet werden. Mit der bloßen Straffreiheit sei dies nicht zu erreichen. Der Vorwurf, dadurch werde das Rechtsgut des werdenden Lebens in die freie und unkontrollierbare Entscheidung der Schwangeren gestellt, treffe nicht zu. Die Entscheidung zum Abbruch der Schwangerschaft richte sich im physischen wie im psychischen und regelmäßig auch im sozialen Sinn auf einen existenziellen Eingriff bei der Frau. Sie gehe nicht, wie im Strafrecht sonst, nur auf die Verletzung eines fremden Rechtsguts; sie gehe immer zugleich auch auf Selbstverletzung. Sie werde deshalb auch nicht leichthin getroffen.
\end{sourcepara}
\begin{transpara}
\transrn{140}立法者承認婦女在此等前提下所作決定為負責任之決定，乃出自一項依靠諮詢與協助而非刑罰威嚇之保護構想的必然結果。諮詢模式若要具有保護效果，必須堅持不具違法性。為使任何思考終止妊娠之婦女均不致因而不願接受諮詢，必須放棄任何形式之壓制、污名化與歧視。單純不罰無法達成此一點。指摘此舉使形成中生命之法益被置於懷孕婦女自由且不可控制之決定中，並不正確。終止妊娠之決定，在身體、心理且通常亦在社會意義上，均指向對婦女之一項生存性介入。其並非如刑法中其他情形般，僅指向對他人法益之侵害；它始終同時也指向自我傷害。因此，該決定亦不會被輕易作成。
\end{transpara}

\begin{sourcepara}[title={原文C.27}]
\sourcern{141}Die Landesregierungen führen dazu ergänzend aus, auch das Grundrecht des \abbrArt{} 4 \abbrAbs{} 1 \abbrGG{}\ggfnArtFourOneDE{} verlange, daß die Entscheidung zum Abbruch einer Schwangerschaft gerechtfertigt sei. Das Bundesverfassungsgericht habe bereits ausgesprochen, daß die Entscheidung zum Abbruch einer Schwangerschaft den Rang einer achtenswerten Gewissensentscheidung haben könne (vgl. \abbrBVerfGE{} 39, 1 [48]). Seit 1975 sei die Gesellschaft sensibler dafür geworden, welche Herausforderung und Belastung Schwangerschaft und Mutterschaft für die Frau bedeuteten. Diese träfen Frauen in ihrer sozialen Identität und Lebensführung heute einschneidender als früher, weil Selbstdarstellung und Lebensführung heute für sie schwieriger seien. Das heiße nicht, daß alle Entscheidungen über Annahme oder Abbruch der Schwangerschaft Gewissensentscheidungen wären. Dennoch sei es dem Gesetzgeber nicht verwehrt, die Gewissensentscheidung als die Regel und den Fall, in dem die Gewissensdimension verstellt und verkannt sei, als die Ausnahme anzusehen.
\end{sourcepara}
\begin{transpara}
\transrn{141}各邦政府就此補充指出，《基本法》第4條第1項\ggfnArtFourOneZH{}之基本權亦要求終止妊娠之決定為正當化。聯邦憲法法院已曾表示，終止妊娠之決定得具有值得尊重之良心決定地位（參見 BVerfGE 39, 1 [48]）。自1975年以來，社會對於懷孕與母職對婦女意味何種挑戰與負擔，已變得更為敏感。此等挑戰與負擔今日對婦女之社會身分及生活安排之衝擊較以往更為深刻，因為今日之自我呈現與生活安排對婦女而言更為困難。這並不表示關於接受或終止妊娠之一切決定均為良心決定。儘管如此，立法者並未被禁止將良心決定視為常態，將良心面向被遮蔽或誤認之情形視為例外。
\end{transpara}

\begin{sourcepara}[title={原文C.28}]
\sourcern{142}3. Die Kompetenz des Bundes zum Erlaß des \abbrArt{} 15 \abbrNr{} 2 \abbrSFHG{} (Sicherstellungsverpflichtung) ergebe sich aus \abbrArt{} 74 \abbrNr{} 1 \abbrGG{}\ggfnArtSeventyFourOneDE{} (Strafrecht) unter dem Gesichtspunkt der Annex-Materie. Als Hilfe in einer konkreten Konflikt- und Notsituation schwangerer Frauen, deren Bewältigung der Gesetzgeber zur öffentlichen Aufgabe gemacht habe, folge -- so die Landesregierungen -- diese Gesetzgebungszuständigkeit des Bundes auch aus \abbrArt{} 74 \abbrNr{} 7 \abbrGG{}\ggfnArtSeventyFourSevenDE{} (öffentliche Fürsorge). Die Vorschrift sei auch inhaltlich verfassungsgemäß. Sie verpflichte die Landesbehörden weder zu rechtlich noch zu tatsächlich unmöglichem Verhalten. \abbrArt{} 2 des 5. \abbrStrRG{} bleibe unberührt; die Gewissens- und Glaubensfreiheit werde nicht angetastet. Religiös "neutrale" kommunale und andere Krankenhäuser könnten sich als Institutionen oder als juristische Personen nicht auf das Grundrecht der Gewissensfreiheit (\abbrArt{} 4 \abbrAbs{} 1 \abbrGG{}\ggfnArtFourOneDE{}) berufen. Da die Leistungen bei nicht rechtswidrigem Schwangerschaftsabbruch zu den Pflichtleistungen der gesetzlichen Krankenversicherung zählten, seien deren Leistungsträger auch verpflichtet, diese Leistungen entweder selbst oder durch Dritte anzubieten. \abbrArt{} 15 \abbrNr{} 2 aktualisiere und konkretisiere diese Verpflichtung.
\end{sourcepara}
\begin{transpara}
\transrn{142}3. 聯邦制定 SFHG 第15條第2款（確保義務）之權限，係從附屬事項觀點源於《基本法》第74條第1款\ggfnArtSeventyFourOneZH{}（刑法）。各邦政府並稱，作為對懷孕婦女具體衝突與困境狀態之協助，而其克服已被立法者定為公共任務，聯邦之此一立法權限亦源於《基本法》第74條第7款\ggfnArtSeventyFourSevenZH{}（公共扶助）。該規定在內容上亦合憲。其既未要求邦機關為法律上不可能之行為，亦未要求其為事實上不可能之行為。《第五刑法改革法》第2條不受影響；良心及信仰自由亦未受觸動。宗教上「中立」之市立及其他醫院，作為機構或法人，不能援引良心自由基本權（《基本法》第4條第1項\ggfnArtFourOneZH{}）。由於不具違法性之終止妊娠之給付屬於法定醫療保險之義務給付，其給付主體亦負有自行或透過第三人提供此等給付之義務。第15條第2款使此一義務現實化並具體化。
\end{transpara}

\begin{sourcepara}[title={原文C.29}]
\sourcern{143}4. Die kompetenzrechtlichen und die materiell-rechtlichen Bedenken der Antragstellerin zu 1) gegen § 24b \abbrSGBV{} seien gleichfalls unbegründet. Die Kompetenz des Bundes ergebe sich aus \abbrArt{} 74 \abbrNr{} 12 \abbrGG{}\ggfnArtSeventyFourTwelveDE{}, der die Einbeziehung neuer Lebenssachverhalte in das Gesamtsystem "Sozialversicherung" ermögliche. Auch die materiell-rechtlichen Bedenken griffen nicht durch. Komme die Schwangere trotz der qualifizierten Beratung und der zahlreichen gesetzlich vorgesehenen Hilfen zur Unausweichlichkeit des Abbruchs oder liege einer der Indikationsfälle vor, so sei es konsequent und trage dem Gebot zu sozialer Gerechtigkeit Rechnung, wenn der Gesetzgeber die eigenverantwortliche Entscheidungsfindung von finanziellen Erwägungen nach Möglichkeit freistelle. Die sozialstaatliche Leistungstypisierung erlaube es, dabei auch Einzelfälle mitzuerfassen, die dem normativen Leitgedanken nicht voll entsprächen.
\end{sourcepara}
\begin{transpara}
\transrn{143}4. 第一聲請人針對 SGB V 第24b條所提出之權限法及實體法疑慮，同樣無理由。聯邦權限源於《基本法》第74條第12款\ggfnArtSeventyFourTwelveZH{}；該款使新生活事實得被納入「社會保險」整體制度。實體法疑慮亦不能成立。若懷孕婦女儘管接受合格諮詢及法律所規定之眾多協助，仍認為終止妊娠不可避免，或存在任一指標事由案件，則立法者盡可能使其自主負責之決定形成免於財務考量，乃屬一貫，且符合社會正義要求。社會國給付類型化允許在此同時涵蓋未完全符合規範指導思想之個案。
\end{transpara}

\begin{sourcepara}[title={原文C.30}]
\sourcern{144}5. Die Aufhebung der Vorschrift über die Bundesstatistik so die Landesregierungen -- sei gleichfalls nicht zu beanstanden. Die lebensschützenden Wirkungen der Neuregelung könnten nur durch eine wissenschaftliche Begleitforschung überprüft werden, die auf Methoden der empirischen Sozialforschung zurückgreife. Für derartige Erhebungen sei das Vorhandensein der Statistik keine notwendige Bedingung.
\end{sourcepara}
\begin{transpara}
\transrn{144}5. 各邦政府並稱，廢止關於聯邦統計之規定，同樣無可指摘。新規制之生命保護效果，只能透過使用經驗社會研究方法之科學伴隨研究加以檢驗。就此類調查而言，統計之存在並非必要條件。
\end{transpara}

\bisubsubsection{IV.}{IV.}

\begin{sourcepara}[title={原文C.31}]
\sourcern{145}Zur Vorbereitung der Entscheidung in den Verfahren 2 \abbrBvF{} 4, 5/92 hat der Zweite Senat des Bundesverfassungsgerichts die Professoren Dres. Stürner und Schulin mit der Erstattung eines Rechtsgutachtens zu folgenden Fragen beauftragt:
\end{sourcepara}
\begin{transpara}
\transrn{145}為準備 2 BvF 4、5/92 程序中之裁判，聯邦憲法法院第二庭委託 Stürner 與 Schulin 兩位教授，就下列問題提出法律意見書：
\end{transpara}

\begin{sourcepara}[title={原文C.32}]
\sourcern{146}(1) Welche Auswirkungen ergeben sich de lege lata für die verschiedenen Bereiche der Rechtsordnung (z.B. Arbeits-, Familien-, Sozial-, ärztliches Berufs-, allgemeines Zivilrecht ), wenn die Rechtsordnung Schwangerschaftsabbrüche rechtlich mißbilligt?
\end{sourcepara}
\begin{transpara}
\transrn{146}（1）若法秩序對終止妊娠作出法律否定評價，依現行法，對法秩序不同領域（例如勞動法、家庭法、社會法、醫師職業法、一般民法）將產生何種影響？
\end{transpara}

\begin{sourcepara}[title={原文C.33}]
\sourcern{147}Welche Auswirkungen hat es auf diese Rechtslage, wenn das Strafrecht unter bestimmten Voraussetzungen (derzeit: Indikationenregelung; angegriffenes Gesetz: innerhalb der ersten zwölf Wochen und nach Beratung) für Schwangerschaftsabbrüche einen Rechtfertigungsgrund vorsieht?
\end{sourcepara}
\begin{transpara}
\transrn{147}若刑法在特定前提下（目前：指標規則；受攻擊法律：最初十二週內且經諮詢）對終止妊娠規定正當化事由，對此一法律狀態將有何影響？
\end{transpara}

\begin{sourcepara}[title={原文C.34}]
\sourcern{148}(2) Welche weiteren Möglichkeiten sind denkbar, um die rechtliche Mißbilligung von Schwangerschaftsabbrüchen außerhalb des Strafrechts in einzelnen Bereichen der Rechtsordnung zum Ausdruck zu bringen? Welche rechtlichen Auswirkungen würden sie haben?
\end{sourcepara}
\begin{transpara}
\transrn{148}（2）尚可設想何種其他可能性，以在刑法以外之法秩序個別領域中表達對終止妊娠之法律否定評價？其將具有何種法律效果？
\end{transpara}

\bisubsubsection{V.}{V.}

\begin{sourcepara}[title={原文C.35}]
\sourcern{149}In der mündlichen Verhandlung vom 8. und 9. Dezember 1992, an der auch Abgeordnete aus allen Fraktionen des 12. Deutschen Bundestages teilgenommen haben, haben die Antragsteller, der Deutsche Bundestag und die Landesregierungen von Brandenburg, Bremen, Hamburg, Hessen, Niedersachsen, Nordrhein-Westfalen, Rheinland-Pfalz, Saarland und Schleswig-Holstein ihr schriftsätzliches Vorbringen bekräftigt. Die Sachverständigen \abbrProf{} \abbrDr{} Stürner und \abbrProf{} \abbrDr{} Schulin haben ihre schriftlichen Rechtsgutachten erläutert und ergänzt. Das Gericht hat als Auskunftspersonen zu Fragen des ärztlichen Berufsrechts Vertreter der Bundesärztekammer und ärztlicher Berufsverbände sowie Mitglieder des Landesärztegerichts Baden- Württemberg herangezogen. Zu Fragen der Beratungs- und der Sozialhilfepraxis hat es weitere auf Veranlassung des Senats durch die Antragsteller und durch Äußerungsberechtigte gestellte sachkundige Personen gehört.
\end{sourcepara}
\begin{transpara}
\transrn{149}在1992年12月8日及9日之言詞辯論中，第12屆德國聯邦議會所有黨團之議員亦有參與；聲請人、德國聯邦議會以及布蘭登堡、不來梅、漢堡、黑森、下薩克森、北萊茵-西伐利亞、萊茵蘭-普法爾茨、薩爾及石勒蘇益格-荷爾斯泰因各邦政府，均重申其書面陳述。鑑定人 Stürner 教授及 Schulin 教授說明並補充其書面法律意見書。法院就醫師職業法問題，徵詢聯邦醫師公會、醫師職業團體代表，以及巴登符騰堡邦醫師法院成員作為資訊提供人。就諮詢及社會扶助實務問題，法院另聽取由聲請人及有權表示意見者應審判庭要求所提出之其他專業人士意見。
\end{transpara}

\bisubsection{D.}{D.}

\bisubsubsection{I.}{I.}

\begin{sourcepara}[title={原文D.36}]
\sourcern{150}1. Das Grundgesetz verpflichtet den Staat, menschliches Leben zu schützen. Zum menschlichen Leben gehört auch das ungeborene. Auch ihm gebührt der Schutz des Staates. Die Verfassung untersagt nicht nur unmittelbare staatliche Eingriffe in das ungeborene Leben, sie gebietet dem Staat auch, sich schützend und fördernd vor dieses Leben zu stellen, d.h. vor allem, es auch vor rechtswidrigen Eingriffen von seiten anderer zu bewahren (vgl. \abbrBVerfGE{} 39, 1 [42]). Ihren Grund hat diese Schutzpflicht in \abbrArt{} 1 \abbrAbs{} 1 \abbrGG{}\ggfnArtOneOneDE{}, der den Staat ausdrücklich zur Achtung und zum Schutz der Menschenwürde verpflichtet; ihr Gegenstand und -- von ihm her -- ihr Maß werden durch \abbrArt{} 2 \abbrAbs{} 2 \abbrGG{}\ggfnArtTwoTwoDE{} näher bestimmt.
\end{sourcepara}
\begin{transpara}
\transrn{150}1. 《基本法》課予國家保護人類生命之義務。人類生命亦包括未出生生命。未出生生命同樣應受國家保護。憲法不僅禁止國家直接介入未出生生命，亦命令國家以保護及促進之方式站在此一生命之前，亦即尤其保護其免受他人之違法介入（參見 BVerfGE 39, 1 [42]）。此一保護義務之基礎在於《基本法》第1條第1項\ggfnArtOneOneZH{}；該項明文課予國家尊重及保護人性尊嚴之義務。其標的，以及由此而來之尺度，則由《基本法》第2條第2項\ggfnArtTwoTwoZH{}加以具體確定。
\end{transpara}

\begin{sourcepara}[title={原文D.37}]
\sourcern{151}a) Menschenwürde kommt schon dem ungeborenen menschlichen Leben zu, nicht erst dem menschlichen Leben nach der Geburt oder bei ausgebildeter Personalität (vgl. bereits § 10 I 1 ALR: "Die allgemeinen Rechte der Menschheit gebühren auch den noch ungeborenen Kindern, schon von der Zeit ihrer Empfängnis."). Es bedarf im vorliegenden Verfahren keiner Entscheidung, ob, wie es Erkenntnisse der medizinischen Anthropologie nahelegen, menschliches Leben bereits mit der Verschmelzung von Ei und Samenzelle entsteht. Gegenstand der angegriffenen Vorschriften ist der Schwangerschaftsabbruch, vor allem die strafrechtliche Regelung; entscheidungserheblich ist daher nur der Zeitraum der Schwangerschaft. Dieser reicht nach den -- von den Antragstellern unbeanstandeten und verfassungsrechtlich unbedenklichen -- Bestimmungen des Strafgesetzbuches vom Abschluß der Einnistung des befruchteten Eis in der Gebärmutter (Nidation; vgl. § 218 \abbrAbs{} 1 Satz 2 \abbrStGB{} in der Fassung des \abbrArt{} 13 \abbrNr{} 1 \abbrSFHG{}) bis zum Beginn der Geburt (vgl. § 217 \abbrStGB{} und dazu \abbrBGHSt{} 32, 194 \abbrFF{}). Jedenfalls in der so bestimmten Zeit der Schwangerschaft handelt es sich bei dem Ungeborenen um individuelles, in seiner genetischen Identität und damit in seiner Einmaligkeit und Unverwechselbarkeit bereits festgelegtes, nicht mehr teilbares Leben, das im Prozeß des Wachsens und Sich-Entfaltens sich nicht erst zum Menschen, sondern als Mensch entwickelt (vgl. \abbrBVerfGE{} 39, 1 [37]). Wie immer die verschiedenen Phasen des vorgeburtlichen Lebensprozesses unter biologischen, philosophischen, auch theologischen Gesichtspunkten gedeutet werden mögen und in der Geschichte beurteilt worden sind, es handelt sich jedenfalls um unabdingbare Stufen der Entwicklung eines individuellen Menschseins. Wo menschliches Leben existiert, kommt ihm Menschenwürde zu (vgl. \abbrBVerfGE{} 39, 1 [41]).
\end{sourcepara}
\begin{transpara}
\transrn{151}a) 人性尊嚴已屬於未出生之人類生命，而非僅屬於出生後或已形成位格性之人類生命（參見早已見於 ALR 第10條 I 1：「人類之一般權利，亦歸屬於尚未出生之子女，自其受孕時起即然。」）。在本程序中，無須決定人類生命是否如醫學人類學知見所提示，已於卵子與精子融合時發生。受攻擊規定之標的為終止妊娠，尤其是其刑法規制；因此，與裁判相關者僅為懷孕期間。依《刑法》規定，該期間自受精卵在子宮中完成著床（著床；參見 SFHG 第13條第1款版本之《刑法》第218條第1項第2句）起，至出生開始時（參見《刑法》第217條及 BGHSt 32, 194 ff.）止；聲請人並未指摘此等規定，且其在憲法上無疑義。無論如何，在如此確定之懷孕期間內，未出生者乃個別生命，其基因同一性，並因而其唯一性與不可混淆性，已經確定；其不再可分割，且在成長與自我展開過程中，並非才成為人，而是作為人而發展（參見 BVerfGE 39, 1 [37]）。出生前生命過程之不同階段，無論可從生物學、哲學或神學觀點如何詮釋，並在歷史上曾如何評價，無論如何均是個別人之存在發展中不可或缺之階段。凡有人類生命存在之處，即有人性尊嚴歸屬於其上（參見 BVerfGE 39, 1 [41]）。
\end{transpara}

\begin{sourcepara}[title={原文D.38}]
\sourcern{152}Diese Würde des Menschseins liegt auch für das ungeborene Leben im Dasein um seiner selbst willen. Es zu achten und zu schützen bedingt, daß die Rechtsordnung die rechtlichen Voraussetzungen seiner Entfaltung im Sinne eines eigenen Lebensrechts des Ungeborenen gewährleistet (vgl. auch \abbrBVerfGE{} 39, 1 [37]). Dieses Lebensrecht, das nicht erst durch die Annahme seitens der Mutter begründet wird, sondern dem Ungeborenen schon aufgrund seiner Existenz zusteht, ist das elementare und unveräußerliche Recht, das von der Würde des Menschen ausgeht; es gilt unabhängig von bestimmten religiösen oder philosophischen Überzeugungen, über die der Rechtsordnung eines religiös-weltanschaulich neutralen Staates kein Urteil zusteht.
\end{sourcepara}
\begin{transpara}
\transrn{152}此種人之存在的尊嚴，對未出生生命而言，亦在於其為自身而存在。尊重並保護此一生命，要求法秩序保障其展開之法律前提，亦即保障未出生者自身之生命權（亦參見 BVerfGE 39, 1 [37]）。此一生命權並非始於母親之接受，而是未出生者基於其存在即已享有；它是源於人性尊嚴之基本且不可讓與之權利。其效力不取決於特定宗教或哲學信念；對此，宗教及世界觀中立國家之法秩序無權作成判斷。
\end{transpara}

\begin{sourcepara}[title={原文D.39}]
\sourcern{153}b) Die Schutzpflicht für das ungeborene Leben ist bezogen auf das einzelne Leben, nicht nur auf menschliches Leben allgemein. Ihre Erfüllung ist eine Grundbedingung geordneten Zusammenlebens im Staat. Sie obliegt aller staatlichen Gewalt (\abbrArt{} 1 \abbrAbs{} 1 Satz 2 \abbrGG{}\ggfnArtOneOneTwoDE{}), d.h. dem Staat in allen seinen Funktionen, auch und gerade der gesetzgebenden Gewalt. Die Schutzpflicht bezieht sich zumal auf Gefahren, die von anderen Menschen ausgehen. Sie umfaßt Schutzmaßnahmen mit dem Ziel, Notlagen im Gefolge einer Schwangerschaft zu vermeiden oder ihnen abzuhelfen, ebenso wie rechtliche Verhaltensanforderungen; beides ergänzt sich.
\end{sourcepara}
\begin{transpara}
\transrn{153}b) 對未出生生命之保護義務，係指向個別生命，而非僅指向一般人類生命。履行此一義務，是國家中有秩序共同生活之基本條件。該義務拘束一切國家權力（《基本法》第1條第1項第2句\ggfnArtOneOneTwoZH{}），亦即拘束國家所有功能，尤其也拘束立法權。保護義務特別指向由其他人所生之危險。其包括以避免或救濟懷孕後續困境為目標之保護措施，也包括法律上之行為要求；二者相互補充。
\end{transpara}

\begin{sourcepara}[title={原文D.40}]
\sourcern{154}2. Verhaltensanforderungen zum Schutz des ungeborenen Lebens stellt der Staat, indem er durch Gesetz Gebote und Verbote ausspricht, Handlungs- und Unterlassungspflichten festlegt. Dies gilt auch für den Schutz des nasciturus gegenüber seiner Mutter, ungeachtet der Verbindung, die zwischen beiden besteht und bei Mutter und Kind zu einem Verhältnis der "Zweiheit in Einheit" führt. Ein solcher Schutz des Ungeborenen gegenüber seiner Mutter ist nur möglich, wenn der Gesetzgeber ihr einen Schwangerschaftsabbruch grundsätzlich verbietet und ihr damit die grundsätzliche Rechtspflicht auferlegt, das Kind auszutragen. Das grundsätzliche Verbot des Schwangerschaftsabbruchs und die grundsätzliche Pflicht zum Austragen des Kindes sind zwei untrennbar verbundene Elemente des verfassungsrechtlich gebotenen Schutzes.
\end{sourcepara}
\begin{transpara}
\transrn{154}2. 國家為保護未出生生命而提出行為要求，其方式是透過法律宣示命令與禁止，並確定作為與不作為義務。這同樣適用於保護 \textit{nasciturus} 免受其母親侵害，儘管二者之間存在聯繫，並在母親與子女之間形成「一體中之二元」關係。對未出生者相對於其母親之此種保護，只有在立法者原則上禁止其終止妊娠，並因而課予其繼續懷孕之基本法律義務時，始有可能。終止妊娠之原則禁止，與繼續懷孕之基本義務，是憲法所要求保護之兩個不可分割之要素。
\end{transpara}

\begin{sourcepara}[title={原文D.41}]
\sourcern{155}Nicht weniger ist Schutz gegenüber Einwirkungen gefordert, die von Dritten -- nicht zuletzt aus dem familiären und dem weiteren sozialen Umfeld der schwangeren Frau -- ausgehen; diese können sich unmittelbar gegen den nasciturus richten, aber auch mittelbar, indem der schwangeren Frau geschuldete Hilfe verweigert, sie wegen der Schwangerschaft in Bedrängnis gebracht oder gar Druck auf sie ausgeübt wird, die Schwangerschaft abzubrechen.
\end{sourcepara}
\begin{transpara}
\transrn{155}對於第三人所生影響，同樣要求提供保護；此等影響尤其可能來自懷孕婦女之家庭及較廣泛社會環境。其可能直接指向 \textit{nasciturus}，也可能以間接方式發生，例如拒絕提供懷孕婦女應得之協助，因其懷孕而使其陷入困境，甚至對其施壓以終止妊娠。
\end{transpara}

\begin{sourcepara}[title={原文D.42}]
\sourcern{156}a) Solche Verhaltensgebote können sich nicht darauf beschränken, Anforderungen an die Freiwilligkeit zu sein, sondern sind als Rechtsgebote auszugestalten. Sie müssen, gemäß der Eigenart des Rechts als einer auf tatsächliche Geltung abzielenden und verwiesenen normativen Ordnung, verbindlich und mit Rechtsfolgen versehen sein. Dabei ist die Strafandrohung nicht die einzig denkbare Sanktion, sie kann allerdings den Rechtsunterworfenen in besonders nachhaltiger Weise zur Achtung und Befolgung der rechtlichen Gebote veranlassen.
\end{sourcepara}
\begin{transpara}
\transrn{156}a) 此等行為命令不能限於對自願性之要求，而必須被形塑為法律命令。依法律作為一種以實際效力為目標並以此為依歸之規範秩序的特性，此等命令必須具有拘束力並附有法律效果。在此，刑罰威嚇並非唯一可想像之制裁；然而其能以特別持久之方式促使受法律拘束者尊重並遵循法律命令。
\end{transpara}

\begin{sourcepara}[title={原文D.43}]
\sourcern{157}Rechtliche Verhaltensgebote sollen Schutz in zwei Richtungen bewirken. Zum einen sollen sie präventive und repressive Schutzwirkungen im einzelnen Fall entfalten, wenn die Verletzung des geschützten Rechtsguts droht oder bereits stattgefunden hat. Zum anderen sollen sie im Volke lebendige Wertvorstellungen und Anschauungen über Recht und Unrecht stärken und unterstützen und ihrerseits Rechtsbewußtsein bilden (vgl. \abbrBVerfGE{} 45, 187 [254, 256]), damit auf der Grundlage einer solchen normativen Orientierung des Verhaltens eine Rechtsgutsverletzung schon von vornherein nicht in Betracht gezogen wird.
\end{sourcepara}
\begin{transpara}
\transrn{157}法律上之行為命令應在兩個方向上產生保護。一方面，當受保護法益之侵害即將發生或已經發生時，其應在個案中發揮預防及壓制性保護效果。另一方面，其應強化並支持人民中活生生存在之關於法與不法的價值觀與看法，並且自身形成法律意識（參見 BVerfGE 45, 187 [254, 256]），使在此種行為之規範取向基礎上，法益侵害從一開始即不被納入考量。
\end{transpara}

\begin{sourcepara}[title={原文D.44}]
\sourcern{158}b) Der Schutz des Lebens ist nicht in dem Sinne absolut geboten, daß dieses gegenüber jedem anderen Rechtsgut ausnahmslos Vorrang genösse; das zeigt schon \abbrArt{} 2 \abbrAbs{} 2 Satz 3 \abbrGG{}\ggfnArtTwoTwoThreeDE{}. Der Schutzpflicht ist andererseits nicht dadurch genügt, daß überhaupt Schutzvorkehrungen irgendeiner Art getroffen worden sind. Ihre Reichweite ist vielmehr im Blick auf die Bedeutung und Schutzbedürftigkeit des zu schützenden Rechtsguts -- hier des ungeborenen menschlichen Lebens -- einerseits und mit ihm kollidierender Rechtsgüter andererseits zu bestimmen (vgl. G. Hermes, Das Grundrecht auf Schutz von Leben und Gesundheit, 1987, \abbrS{} 253 \abbrFF{}). Als vom Lebensrecht des Ungeborenen berührte Rechtsgüter kommen dabei -- ausgehend vom Anspruch der schwangeren Frau auf Schutz und Achtung ihrer Menschenwürde (\abbrArt{} 1 \abbrAbs{} 1 \abbrGG{}\ggfnArtOneOneDE{}) -- vor allem ihr Recht auf Leben und körperliche Unversehrtheit (\abbrArt{} 2 \abbrAbs{} 2 \abbrGG{}\ggfnArtTwoTwoDE{}) sowie ihr Persönlichkeitsrecht (\abbrArt{} 2 \abbrAbs{} 1 \abbrGG{}\ggfnArtTwoOneDE{}) in Betracht.
\end{sourcepara}
\begin{transpara}
\transrn{158}b) 生命保護並非在生命相對於其他任何法益均毫無例外享有優先之意義上被絕對要求；《基本法》第2條第2項第3句\ggfnArtTwoTwoThreeZH{}已顯示此點。另一方面，保護義務也並非只要採取某種保護預防措施即已滿足。其範圍毋寧應一方面著眼於受保護法益，亦即此處之未出生人類生命，之重要性與保護需求，另一方面著眼於與之衝突之法益而加以確定（參見 G. Hermes, Das Grundrecht auf Schutz von Leben und Gesundheit, 1987, S. 253 ff.）。作為受未出生者生命權所觸及之法益，從懷孕婦女對其人性尊嚴受保護與尊重之請求（《基本法》第1條第1項\ggfnArtOneOneZH{}）出發，尤其包括其生命權及身體完整權（《基本法》第2條第2項\ggfnArtTwoTwoZH{}），以及其人格權（《基本法》第2條第1項\ggfnArtTwoOneZH{}）。
\end{transpara}

\begin{sourcepara}[title={原文D.45}]
\sourcern{159}Art und Umfang des Schutzes im einzelnen zu bestimmen, ist Aufgabe des Gesetzgebers. Die Verfassung gibt den Schutz als Ziel vor, nicht aber seine Ausgestaltung im einzelnen. Allerdings hat der Gesetzgeber das Untermaßverbot zu beachten (vgl. zum Begriff Isensee in: Handbuch des Staatsrechts, Band V, 1992, § 111 \abbrRdnrn{} 165 \abbrF{}); insofern unterliegt er der verfassungsgerichtlichen Kontrolle. Notwendig ist ein -- unter Berücksichtigung entgegenstehender Rechtsgüter -- angemessener Schutz; entscheidend ist, daß er als solcher wirksam ist. Die Vorkehrungen, die der Gesetzgeber trifft, müssen für einen angemessenen und wirksamen Schutz ausreichend sein und zudem auf sorgfältigen Tatsachenermittlungen und vertretbaren Einschätzungen beruhen (dazu unten I.4.). Das danach verfassungsrechtlich gebotene Maß des Schutzes ist unabhängig vom Alter der Schwangerschaft. Das Grundgesetz enthält für das ungeborene Leben keine vom Ablauf bestimmter Fristen abhängige, dem Entwicklungsprozeß der Schwangerschaft folgende Abstufungen des Lebensrechts und seines Schutzes. Auch in der Frühphase einer Schwangerschaft hat die Rechtsordnung deshalb dieses Maß an Schutz zu gewährleisten.
\end{sourcepara}
\begin{transpara}
\transrn{159}具體確定保護之種類與範圍，是立法者之任務。憲法預設保護作為目標，但並未預設其個別形塑。然而，立法者必須遵守不足禁止（關於此概念，參見 Isensee in: Handbuch des Staatsrechts, Band V, 1992, § 111 Rdnrn. 165 f.）；就此而言，立法者受憲法法院審查。必要者是在考量相反法益下之適當保護；關鍵在於其作為保護本身具有實效。立法者所採預防安排必須足以達成適當且有效之保護，且此外必須建立在審慎事實調查及可支持之評估上（對此見下文 I.4.）。依此在憲法上所要求之保護程度，不取決於懷孕期間之長短。《基本法》對未出生生命並不包含依特定期限經過而定、隨懷孕發展過程而來之生命權及其保護階層化。因此，即使在懷孕早期，法秩序亦必須保障此一程度之保護。
\end{transpara}

\begin{sourcepara}[title={原文D.46}]
\sourcern{160}c) Soll das Untermaßverbot nicht verletzt werden, muß die Ausgestaltung des Schutzes durch die Rechtsordnung Mindestanforderungen entsprechen.
\end{sourcepara}
\begin{transpara}
\transrn{160}c) 若不得違反不足禁止，法秩序對保護之形塑必須符合最低要求。
\end{transpara}

\begin{sourcepara}[title={原文D.47}]
\sourcern{161}aa) Hierzu zählt, daß der Schwangerschaftsabbruch für die ganze Dauer der Schwangerschaft grundsätzlich als Unrecht angesehen wird und demgemäß rechtlich verboten ist (vgl. \abbrBVerfGE{} 39, 1 [44]). Bestünde ein solches Verbot nicht, würde also die Verfügung über das Lebensrecht des nasciturus, wenn auch nur für eine begrenzte Zeit, der freien, rechtlich nicht gebundenen Entscheidung eines Dritten, und sei es selbst der Mutter, überantwortet, wäre rechtlicher Schutz dieses Lebens im Sinne der oben genannten Verhaltensanforderungen nicht mehr gewährleistet. Eine solche Preisgabe des ungeborenen Lebens läßt sich auch unter Hinweis auf die Menschenwürde der Frau und ihre Fähigkeit zu verantwortlicher Entscheidung nicht einfordern. Rechtlicher Schutz bedingt, daß das Recht selbst Umfang und Grenzen zulässigen Einwirkens des einen auf den anderen normativ festlegt und nicht dem Belieben eines der Beteiligten überläßt.
\end{sourcepara}
\begin{transpara}
\transrn{161}aa) 其中包括：終止妊娠在整個懷孕期間原則上均被視為不法，並相應地在法律上被禁止（參見 BVerfGE 39, 1 [44]）。若不存在此種禁止，亦即若 \textit{nasciturus} 之生命權處分，即使僅在有限期間內，被交付予第三人自由且不受法律拘束之決定，即便該第三人為母親自身，則此一生命在上述行為要求意義下之法律保護即不再獲得保障。即使援引婦女之人性尊嚴及其負責任決定之能力，亦不能要求如此交出未出生生命。法律保護要求法律本身以規範方式確定一方對另一方之容許影響的範圍與界限，而非交由參與者之一方任意決定。
\end{transpara}

\begin{sourcepara}[title={原文D.48}]
\sourcern{162}Grundrechte der Frau greifen gegenüber dem grundsätzlichen Verbot des Schwangerschaftsabbruchs nicht durch. Zwar haben diese Rechte auch gegenüber dem Lebensrecht des nasciturus Bestand und sind entsprechend zu schützen. Aber sie tragen nicht so weit, daß die Rechtspflicht zum Austragen des Kindes von Grundrechts wegen -- auch nur für eine bestimmte Zeit -- generell aufgehoben wäre. Die Grundrechtspositionen der Frau führen allerdings dazu, daß es in Ausnahmelagen zulässig, in manchen dieser Fälle womöglich geboten ist, eine solche Rechtspflicht nicht aufzuerlegen.
\end{sourcepara}
\begin{transpara}
\transrn{162}婦女之基本權並不能突破終止妊娠之原則禁止。誠然，此等權利相對於 \textit{nasciturus} 之生命權亦仍存在，並應受到相應保護。但其並未達到使繼續懷孕之法律義務因基本權而一般性取消之程度，即使僅限於特定期間亦然。不過，婦女之基本權地位確實導致在例外狀態中，可以不課予此種法律義務；在其中若干情形中，甚至可能必須如此。
\end{transpara}

\begin{sourcepara}[title={原文D.49}]
\sourcern{163}bb) Solche Ausnahmelagen zu Ausnahmetatbeständen zu fassen, ist Sache des Gesetzgebers. Um dabei das Untermaßverbot nicht zu verletzen, muß er allerdings in Rechnung stellen, daß die miteinander kollidierenden Rechtsgüter hier nicht zu einem verhältnismäßigen Ausgleich gebracht werden können, weil auf der Seite des ungeborenen Lebens in jedem Fall nicht ein Mehr oder Weniger an Rechten, die Hinnahme von Nachteilen oder Einschränkungen, sondern alles, nämlich das Leben selbst, in Frage steht. Ein Ausgleich, der sowohl den Lebensschutz des nasciturus gewährleistet als auch der schwangeren Frau ein Recht zum Schwangerschaftsabbruch einräumt, ist nicht möglich, weil Schwangerschaftsabbruch immer Tötung ungeborenen Lebens ist (vgl. \abbrBVerfGE{} 39, 1 [43]). Ein Ausgleich kann auch nicht -- wie es vertreten wird (vgl. Nelles in: "Zur Sache, Themen parlamentarischer Beratung", herausgegeben vom Deutschen Bundestag, Band 1/92, \abbrS{} 250) -- dadurch erreicht werden, daß für eine gewisse Zeit der Schwangerschaft das Persönlichkeitsrecht der Frau vorgeht und danach erst das Recht des ungeborenen Lebens Vorrang erhält. Das Lebensrecht des Ungeborenen käme dann nur zur Geltung, wenn die Mutter sich nicht in der ersten Phase der Schwangerschaft zu seiner Tötung entschlossen hat.
\end{sourcepara}
\begin{transpara}
\transrn{163}bb) 將此等例外狀態形塑為例外構成要件，是立法者之事。惟為不違反不足禁止，立法者必須考量：在此相互衝突之法益不能被帶入比例性調和，因為在未出生生命一方，無論何種情形中，所涉及者均不是權利之多或少、承受不利益或限制，而是一切，亦即生命本身。既保障 \textit{nasciturus} 之生命保護，又賦予懷孕婦女終止妊娠之權利，二者不可能相互調和，因為終止妊娠始終是殺害未出生生命（參見 BVerfGE 39, 1 [43]）。調和亦不能如有見解所主張（參見 Nelles in: "Zur Sache, Themen parlamentarischer Beratung", herausgegeben vom Deutschen Bundestag, Band 1/92, S. 250），透過在懷孕某一期間使婦女人格權優先，而之後才使未出生生命之權利取得優先而達成。如此一來，未出生者之生命權只有在母親未於懷孕第一階段決定殺害其生命時，才會發生作用。
\end{transpara}

\begin{sourcepara}[title={原文D.50}]
\sourcern{164}Das bedeutet indes nicht, daß eine Ausnahmelage, die es von Verfassungs wegen zuläßt, die Pflicht zum Austragen des Kindes aufzuheben, nur im Falle einer ernsten Gefahr für das Leben der Frau oder einer schwerwiegenden Beeinträchtigung ihrer Gesundheit in Betracht kommt. Ausnahmelagen sind auch darüber hinaus denkbar. Das Kriterium für ihre Anerkennung ist, wie das Bundesverfassungsgericht festgestellt hat, das der Unzumutbarkeit (vgl. \abbrBVerfGE{} 39, 1 [48 \abbrFF{}]). Dieses Kriterium hat -- unbeschadet des Umstandes, daß die Beteiligung der Frau an dem Schwangerschaftsabbruch strafrechtlich nicht als Unterlassungsdelikt einzuordnen ist -- deshalb seine Berechtigung, weil sich das Verbot des Schwangerschaftsabbruchs angesichts der einzigartigen Verbindung von Mutter und Kind nicht in einer Pflicht der Frau erschöpft, den Rechtskreis eines anderen nicht zu verletzen, sondern zugleich eine intensive, die Frau existentiell betreffende Pflicht zum Austragen und Gebären des Kindes enthält und eine darüber hinausgehende Handlungs-, Sorge- und Einstandspflicht nach der Geburt über viele Jahre nach sich zieht (vgl. insoweit auch M. von Renesse, ZRP 1991, \abbrS{} 321 [322 \abbrF{}]). Aus der Vorausschau auf die damit verbundenen Belastungen können in der besonderen seelischen Lage, in der sich werdende Mütter gerade in der Frühphase einer Schwangerschaft vielfach befinden, in Einzelfällen schwere, unter Umständen auch lebensbedrohende Konfliktsituationen entstehen, in denen schutzwürdige Positionen einer schwangeren Frau sich mit solcher Dringlichkeit geltend machen, daß jedenfalls die staatliche Rechtsordnung -- ungeachtet etwa weitergehender moralischer oder religiös begründeter Pflichtauffassungen -- nicht verlangen kann, die Frau müsse hier dem Lebensrecht des Ungeborenen unter allen Umständen den Vorrang geben (vgl. \abbrBVerfGE{} 39, 1 [50]).
\end{sourcepara}
\begin{transpara}
\transrn{164}然而，這並不表示，憲法上允許取消繼續懷孕義務之例外狀態，僅在婦女生命有嚴重危險或其健康受重大損害時始可能存在。除此之外，例外狀態亦可想像。依聯邦憲法法院所確認，承認此等例外狀態之標準是不期待性（參見 BVerfGE 39, 1 [48 ff.]）。此一標準之正當性在於：鑑於母親與子女之獨特聯繫，終止妊娠之禁止並不窮盡於婦女不得侵害他人權利領域之義務，而同時包含懷胎並生下子女之強烈義務，此義務涉及婦女之生存，且在出生後又引發多年之作為、照護及負責義務（就此亦參見 M. von Renesse, ZRP 1991, S. 321 [322 f.]）；此點不受婦女參與終止妊娠在刑法上並非被歸類為不作為犯之影響。預見與此相連之負擔，可能在準母親尤其於懷孕早期常常所處之特殊心理狀態中，於個案產生嚴重、在特定情形下甚至危及生命之衝突狀態；在此等狀態中，懷孕婦女值得保護之地位以如此迫切性提出，以致至少國家法秩序，無論是否存在更進一步之道德或宗教義務觀念，均不能要求婦女在此等情形下必須無條件使未出生者之生命權居於優先（參見 BVerfGE 39, 1 [50]）。
\end{transpara}

\begin{sourcepara}[title={原文D.51}]
\sourcern{165}Eine Unzumutbarkeit kann allerdings nicht aus Umständen herrühren, die im Rahmen der Normalsituation einer Schwangerschaft verbleiben. Vielmehr müssen Belastungen gegeben sein, die ein solches Maß an Aufopferung eigener Lebenswerte verlangen, daß dies von der Frau nicht erwartet werden kann.
\end{sourcepara}
\begin{transpara}
\transrn{165}不過，不期待性不能源於仍停留在懷孕通常狀態範圍內之情事。毋寧必須存在負擔，其要求犧牲自身生命價值之程度，已不能期待婦女承受。
\end{transpara}

\begin{sourcepara}[title={原文D.52}]
\sourcern{166}Für die Pflicht zum Austragen des Kindes folgt daraus, daß neben der hergebrachten medizinischen Indikation auch die kriminologische und -- ihre hinreichend genaue Umgrenzung vorausgesetzt -- die embryopathische Indikation als Ausnahmetatbestände vor der Verfassung Bestand haben können; für andere Notlagen gilt dies nur dann, wenn in ihrer Umschreibung die Schwere des hier vorauszusetzenden sozialen oder psychisch-personalen Konflikts deutlich erkennbar wird, so daß -- unter dem Gesichtspunkt der Unzumutbarkeit betrachtet -- die Kongruenz mit den anderen Indikationsfällen gewahrt bleibt (vgl. auch \abbrBVerfGE{} 39, 1 [50]).
\end{sourcepara}
\begin{transpara}
\transrn{166}由此對繼續懷孕義務可得出：除傳統醫學指標事由外，犯罪學指標事由，以及在其界定足夠精確之前提下之胚胎病理指標事由，均得作為例外構成要件而在憲法前成立；至於其他困境，僅於其描述中清楚顯示此處所預設之社會或心理人格衝突的嚴重性，以致從不期待性觀點觀察時，仍維持與其他指標事由案件之一致性時，始有同樣結論（亦參見 BVerfGE 39, 1 [50]）。
\end{transpara}

\begin{sourcepara}[title={原文D.53}]
\sourcern{167}cc) Soweit die Unzumutbarkeit die Pflicht der Frau, das Kind auszutragen, begrenzt, ist damit die Schutzpflicht des Staates, die gegenüber jedem ungeborenen menschlichen Leben besteht, nicht aufgehoben. Sie wird den Staat insbesondere veranlassen, durch Rat und Hilfe der Frau beizustehen und sie dadurch womöglich doch für das Austragen des Kindes zu gewinnen; davon geht auch die Regelung in § 218a \abbrAbs{} 3 \abbrStGB{} \abbrNF{} aus.
\end{sourcepara}
\begin{transpara}
\transrn{167}cc) 在不期待性限制婦女繼續懷孕義務之範圍內，國家對每一未出生人類生命所負之保護義務並未因此取消。該義務尤其將促使國家透過諮詢與協助支持婦女，並可能藉此仍使其願意繼續懷孕；新版本《刑法》第218a條第3項之規制亦以此為出發點。
\end{transpara}

\begin{sourcepara}[title={原文D.54}]
\sourcern{168}dd) Handelt es sich bei der Aufgabe, das menschliche Leben vor seiner Tötung zu schützen, um eine elementare staatliche Schutzaufgabe, so läßt es das Untermaßverbot nicht zu, auf den Einsatz auch des Strafrechts und die davon ausgehende Schutzwirkung frei zu verzichten.
\end{sourcepara}
\begin{transpara}
\transrn{168}dd) 若保護人類生命免遭殺害之任務是一項基本國家保護任務，則不足禁止並不允許自由放棄運用刑法以及由刑法所生之保護效果。
\end{transpara}

\begin{sourcepara}[title={原文D.55}]
\sourcern{169}Dem Strafrecht kommt seit jeher und auch unter den heutigen Gegebenheiten die Aufgabe zu, die Grundlagen eines geordneten Gemeinschaftslebens zu schützen. Dazu gehört die Achtung und grundsätzliche Unverletzlichkeit menschlichen Lebens. Dementsprechend wird die Tötung anderer umfassend mit Strafe bedroht.
\end{sourcepara}
\begin{transpara}
\transrn{169}刑法自始以來，即使在今日條件下，亦負有保護有秩序共同生活基礎之任務。其中包括尊重人類生命及其原則上不可侵害性。相應地，殺害他人受到全面刑罰威嚇。
\end{transpara}

\begin{sourcepara}[title={原文D.56}]
\sourcern{170}Das Strafrecht ist zwar nicht das primäre Mittel rechtlichen Schutzes, schon wegen seines am stärksten eingreifenden Charakters; seine Verwendung unterliegt daher den Anforderungen der Verhältnismäßigkeit (\abbrBVerfGE{} 6, 389 [433 \abbrF{}]; 39, 1 [47]; 57,250 [270]; 73, 206 [253]). Aber es wird als "ultima ratio" dieses Schutzes eingesetzt, wenn ein bestimmtes Verhalten über sein Verbotensein hinaus in besonderer Weise sozialschädlich und für das geordnete Zusammenleben der Menschen unerträglich, seine Verhinderung daher besonders dringlich ist.
\end{sourcepara}
\begin{transpara}
\transrn{170}刑法固然並非法律保護之首要手段，這已由於其干預性最強之性格；因此，其使用受比例性要求拘束（BVerfGE 6, 389 [433 f.]; 39, 1 [47]; 57,250 [270]; 73, 206 [253]）。但是，當特定行為除其被禁止外，並以特別方式對社會有害且為人類有秩序共同生活所不能忍受，因而防止該行為特別急迫時，刑法即被作為此種保護之「最後手段」使用。
\end{transpara}

\begin{sourcepara}[title={原文D.57}]
\sourcern{171}Danach ist das Strafrecht regelmäßig der Ort, das grundsätzliche Verbot des Schwangerschaftsabbruchs und die darin enthaltene grundsätzliche Rechtspflicht der Frau zum Austragen des Kindes gesetzlich zu verankern. Sofern allerdings wegen verfassungsrechtlich ausreichender Schutzmaßnahmen anderer Art von einer Strafdrohung für nicht gerechtfertigte Schwangerschaftsabbrüche in begrenztem Umfang abgesehen werden darf, kann es auch genügen, das Verbot für diese Fallgruppe auf andere Weise in der Rechtsordnung unterhalb der Verfassung klar zum Ausdruck zu bringen (vgl. \abbrBVerfGE{} 39, 1 [44, 46]).
\end{sourcepara}
\begin{transpara}
\transrn{171}依此，刑法通常是以法律方式錨定終止妊娠之原則禁止，以及其中所含婦女繼續懷孕之基本法律義務的場所。不過，若因其他種類之保護措施已在憲法上足夠，而得在有限範圍內不對未經正當化之終止妊娠施以刑罰威嚇，則以其他方式在憲法以下之法秩序中清楚表達此一案件群之禁止，亦可能已足夠（參見 BVerfGE 39, 1 [44, 46]）。
\end{transpara}

\begin{sourcepara}[title={原文D.58}]
\sourcern{172}3. Der Staat genügt seiner Schutzpflicht gegenüber dem ungeborenen menschlichen Leben nicht allein dadurch, daß er Angriffen wehrt, die diesem von anderen Menschen drohen. Er muß auch denjenigen Gefahren entgegentreten, die für dieses Leben in den gegenwärtigen und absehbaren realen Lebensverhältnissen der Frau und der Familie begründet liegen und der Bereitschaft zum Austragen des Kindes entgegenwirken. Darin berührt sich die Schutzpflicht mit dem Schutzauftrag aus \abbrArt{} 6 \abbrAbs{} 1 und 4 \abbrGG{} (zu \abbrArt{} 6 \abbrAbs{} 1 vgl. \abbrBVerfGE{} 76, 1 [44 \abbrF{}, 49 \abbrF{}]; zu \abbrArt{} 6 \abbrAbs{} 4 vgl. zuletzt \abbrBVerfGE{} 84, 133 [155 \abbrF{}]). Sie verpflichtet die staatliche Gewalt, Problemen und Schwierigkeiten nachzugehen, die der Mutter während und nach der Schwangerschaft erwachsen können. \abbrArt{} 6 \abbrAbs{} 4 \abbrGG{}\ggfnArtSixFourDE{} enthält einen für den gesamten Bereich des privaten und öffentlichen Rechts verbindlichen Schutzauftrag, der sich auch auf die schwangere Frau erstreckt. Diesem Auftrag entspricht es, Mutterschaft und Kinderbetreuung als eine Leistung zu betrachten, die auch im Interesse der Gemeinschaft liegt und deren Anerkennung verlangt.
\end{sourcepara}
\begin{transpara}
\transrn{172}3. 國家對未出生人類生命之保護義務，並非僅透過防禦他人對其所威脅之攻擊即已滿足。國家亦必須對抗那些根基於婦女與家庭之現時及可預見實際生活關係中，並妨礙繼續懷孕意願之危險。保護義務在此與《基本法》第6條第1項\ggfnArtSixOneZH{}及第4項\ggfnArtSixFourZH{}之保護任務相接觸（關於第6條第1項，參見 BVerfGE 76, 1 [44 f., 49 f.]；關於第6條第4項，最近參見 BVerfGE 84, 133 [155 f.]）。其課予國家權力追究母親在懷孕期間及懷孕後可能遭遇之問題與困難的義務。《基本法》第6條第4項\ggfnArtSixFourZH{}包含一項拘束私法與公法全領域之保護任務，並且亦及於懷孕婦女。將母職與子女照顧視為一項同時符合共同體利益並要求承認之給付，正符合此一任務。
\end{transpara}

\begin{sourcepara}[title={原文D.59}]
\sourcern{173}Schon im Ersten Bericht des Sonderausschusses für die Strafrechtsreform (\abbrBTDrucks{} 7/1981 [neu] \abbrS{} 7) wird ausgeführt, daß die ungünstige Wohnungssituation, die Unmöglichkeit, neben einer Ausbildung oder Erwerbstätigkeit ein Kind zu versorgen, sowie wirtschaftliche Not und sonstige materielle Gründe, außerdem bei alleinstehenden Schwangeren die Angst vor gesellschaftlicher Diskriminierung, zu den häufig genannten Ursachen und Motiven für den Wunsch nach einem Abbruch der Schwangerschaft gehören.
\end{sourcepara}
\begin{transpara}
\transrn{173}早在刑法改革特別委員會第一份報告（BTDrucks. 7/1981 [neu] S. 7）中即已指出，不利之居住狀況、在接受教育或從事職業活動之同時無法照顧子女，以及經濟困境與其他物質理由，此外對單身懷孕婦女而言，對社會歧視之恐懼，均屬於終止妊娠願望常被提及之原因與動機。
\end{transpara}

\begin{sourcepara}[title={原文D.60}]
\sourcern{174}a) Die der Mutter geschuldete Fürsorge der Gemeinschaft umfaßt die Verpflichtung des Staates, darauf hinzuwirken, daß eine Schwangerschaft nicht wegen einer bestehenden oder nach der Geburt des Kindes drohenden materiellen Notlage abgebrochen wird. Ebenso sind Nachteile, die einer Frau aus der Schwangerschaft für Ausbildung und Beruf erwachsen können, nach Möglichkeit auszuschließen. Dem Staat obliegt es, in Wahrnehmung seiner Schutzpflicht gegenüber dem ungeborenen menschlichen Leben den Umständen nachzugehen, die die Lage der schwangeren Frau und der Mutter zu erschweren geeignet sind, und sich im Maße des rechtlich und tatsächlich Möglichen und Verantwortbaren um Abhilfe oder Erleichterung zu bemühen. Das betrifft nicht nur den Gesetzgeber, sondern auch Regierung und Verwaltung.
\end{sourcepara}
\begin{transpara}
\transrn{174}a) 共同體對母親所負之照顧，包括國家促成懷孕不因既存或子女出生後迫近之物質困境而被終止的義務。同樣，婦女因懷孕在教育與職業上可能遭受之不利益，亦應盡可能排除。國家在履行其對未出生人類生命之保護義務時，負有追究可能使懷孕婦女及母親處境困難化之情事，並在法律上及事實上可能且可負責之程度內，努力予以救濟或減輕之責。這不僅涉及立法者，亦涉及政府與行政。
\end{transpara}

\begin{sourcepara}[title={原文D.61}]
\sourcern{175}Zwar kann und muß der Staat den Eltern nicht alle Belastungen und Einschränkungen abnehmen, die für sie mit der "Pflege und Erziehung" von Kindern (\abbrArt{} 6 \abbrAbs{} 2 Satz 1 \abbrGG{}\ggfnArtSixTwoOneDE{}) verbunden sind. Indessen gibt es -- über die Regelungen hinaus, die das Schwangeren- und Familienhilfegesetz in seinen Artikeln 5 bis 12 getroffen hat -- weitere Möglichkeiten der Entlastung. So bieten sich schon im staatlichen Bereich viele Möglichkeiten eines wirksameren Schutzes von Mutter und Kind; zu denken ist an das Wohnungswesen, den öffentlichen Dienst und berufs-, insbesondere auch ausbildungsrechtliche Regelungen.
\end{sourcepara}
\begin{transpara}
\transrn{175}誠然，國家不能也無須替父母承擔與子女之「照護與教育」（《基本法》第6條第2項第1句\ggfnArtSixTwoOneZH{}）相連之一切負擔與限制。然而，除《孕婦及家庭扶助法》第5條至第12條所作規制外，尚有進一步減輕負擔之可能。即便在國家領域中，也有許多更有效保護母親與子女之可能；例如可想到住宅制度、公務部門，以及職業法，尤其也包括教育培訓法規制。
\end{transpara}

\begin{sourcepara}[title={原文D.62}]
\sourcern{176}Der Staat kann aber -- und muß gegebenenfalls -- an einem wirksamen Schutz auch Dritte beteiligen. Eltern übernehmen mit der Erziehung ihrer Kinder zugleich Aufgaben, deren Erfüllung sowohl im Interesse der Gemeinschaft als Ganzer als auch jedes einzelnen gelegen ist. Darum ist der Staat gehalten, eine kinderfreundliche Gesellschaft zu fördern, was auch auf den Schutz des ungeborenen Lebens zurückwirkt. Mit Blick darauf hat der Gesetzgeber nicht nur im Bereich des Arbeitsrechts, sondern auch in anderen Bereichen des Privatrechts Regelungen mit besonderer Rücksicht auf Familien mit Kindern zu erwägen; hierher zählen etwa das Verbot einer Kündigung von Mietverträgen über Wohnraum wegen der Aufnahme eines neugeborenen Kindes, aber auch Bestimmungen im Bereich des Kreditwesens über Vertragsgestaltungen oder staatliche Vertragshilfen, welche den Eltern eine Erfüllung von Kreditverpflichtungen nach der Geburt eines Kindes ermöglichen oder erleichtern.
\end{sourcepara}
\begin{transpara}
\transrn{176}但國家也可以，並且在必要時必須，使第三人參與有效保護。父母在教育其子女時，同時承擔了其履行既符合整體共同體利益、亦符合每一個人利益之任務。因此，國家有義務促進一個對子女友善之社會，而此亦反作用於未出生生命之保護。著眼於此，立法者不僅在勞動法領域，也應在私法其他領域考慮特別顧及有子女家庭之規制；例如，禁止因接納新生兒而終止住宅租賃契約，亦包括信貸制度領域中關於契約安排或國家契約協助之規定，使父母在子女出生後得以履行或較易履行信貸義務。
\end{transpara}

\begin{sourcepara}[title={原文D.63}]
\sourcern{177}b) Der Schutz des ungeborenen Lebens, der Schutzauftrag für Ehe und Familie (\abbrArt{} 6 \abbrGG{}\ggfnArtSixDE{}) und die Gleichstellung von Mann und Frau in der Teilhabe am Arbeitsleben (vgl. \abbrArt{} 3 \abbrAbs{} 2 \abbrGG{}\ggfnArtThreeTwoDE{} sowie \abbrArt{} 3, 7 des Internationalen Paktes über wirtschaftliche, soziale und kulturelle Rechte vom 19. Dezember 1966 [\abbrBGBl{} 1973 II, \abbrS{} 1570]) verpflichten den Staat und insbesondere den Gesetzgeber, Grundlagen dafür zu schaffen, daß Familientätigkeit und Erwerbstätigkeit aufeinander abgestimmt werden können und die Wahrnehmung der familiären Erziehungsaufgabe nicht zu beruflichen Nachteilen führt. Dazu zählen auch rechtliche und tatsächliche Maßnahmen, die ein Nebeneinander von Erziehungs- und Erwerbstätigkeit für beide Elternteile ebenso wie eine Rückkehr in eine Berufstätigkeit und einen beruflichen Aufstieg auch nach Zeiten der Kindererziehung ermöglichen. Hierher gehören auch die durch \abbrArt{} 6 und 7 \abbrSFHG{} erfolgten Änderungen des Arbeitsförderungs- und des Berufsbildungsgesetzes; damit beschreitet der Gesetzgeber den richtigen Weg. Gleiches gilt für Regelungen, die auf eine Verbesserung der institutionellen (vgl. \abbrArt{} 5 \abbrSFHG{}) oder familiären Kinderbetreuung zielen (vgl. etwa die Leistungen im Rahmen des sogenannten Familienlastenausgleichs, insbesondere das Erziehungsgeld, sowie die Regelungen über den Erziehungsurlaub und den Unterhaltsvorschuß). Die Bedeutung solcher Leistungen als Maßnahmen präventiven Lebensschutzes hat der Gesetzgeber in Rechnung zu stellen, wenn es erforderlich wird, staatliche Leistungen im Hinblick auf knappe Mittel zu überprüfen.
\end{sourcepara}
\begin{transpara}
\transrn{177}b) 未出生生命之保護、婚姻與家庭之保護任務（《基本法》第6條\ggfnArtSixZH{}），以及男女參與工作生活之平等（參見《基本法》第3條第2項\ggfnArtThreeTwoZH{}，以及1966年12月19日《經濟、社會及文化權利國際公約》第3條、第7條 [BGBl. 1973 II, S. 1570]），均課予國家，尤其是立法者，創設基礎之義務，使家庭活動與職業活動得以相互協調，且履行家庭教育任務不致造成職業上不利益。其中亦包括法律及事實措施，使雙親均得同時從事教育活動與職業活動，並使其在子女教育期間後仍能返回職業活動及取得職業升遷。SFHG 第6條及第7條對《就業促進法》及《職業教育法》所作修改亦屬於此；立法者由此走上正確道路。對於旨在改善機構性（參見 SFHG 第5條）或家庭性子女照顧之規制，亦同（例如所謂家庭負擔均衡範圍內之給付，尤其是育兒津貼，以及育兒假與扶養預付之規制）。當有必要鑑於資金短缺而審查國家給付時，立法者必須將此等給付作為預防性生命保護措施之意義納入考量。
\end{transpara}

\begin{sourcepara}[title={原文D.64}]
\sourcern{178}c) Weiter hat der Staat dafür zu sorgen, daß ein Elternteil, der sich unter Verzicht auf Erwerbseinkommen der Erziehung eines Kindes widmet, für die ihm hieraus erwachsenden versorgungsrechtlichen Nachteile einen angemessenen Ausgleich erhält. Den Ausführungen hierzu im Urteil des Ersten Senats vom 7. Juli 1992 -- 1 \abbrBvL{} 51/86, 50/87 und 1 \abbrBvR{} 873/90, 761/91 (Umdruck \abbrS{} 55 \abbrF{} -- \abbrBVerfGE{} 87, 1 \abbrFF{}) tritt der Senat bei.
\end{sourcepara}
\begin{transpara}
\transrn{178}c) 此外，國家必須確保，因放棄職業收入而投入子女教育之一方父母，就其因此所生之供養法上不利益，獲得適當補償。第二庭贊同第一庭1992年7月7日判決 -- 1 BvL 51/86, 50/87 及 1 BvR 873/90, 761/91（印本 S. 55 f. -- BVerfGE 87, 1 ff.）中就此所作說明。
\end{transpara}

\begin{sourcepara}[title={原文D.65}]
\sourcern{179}d) Der Schutzauftrag verpflichtet den Staat schließlich auch, den rechtlichen Schutzanspruch des ungeborenen Lebens im allgemeinen Bewußtsein zu erhalten und zu beleben. Deshalb müssen die Organe des Staates in Bund und Ländern erkennbar für den Schutz des Lebens eintreten. Das betrifft auch und gerade die Lehrpläne der Schulen. Öffentliche Einrichtungen, die Aufklärung in gesundheitlichen Fragen, Familienberatung oder Sexualaufklärung betreiben, haben allgemein den Willen zum Schutz des ungeborenen Lebens zu stärken; dies gilt insbesondere für die in \abbrArt{} 1 § 1 \abbrSFHG{} vorgesehene Aufklärung. Öffentlich-rechtlicher wie privater Rundfunk sind bei Ausübung ihrer Rundfunkfreiheit (\abbrArt{} 5 \abbrAbs{} 1 \abbrGG{}\ggfnArtFiveOneDE{}) der Würde des Menschen verpflichtet (zum privaten Rundfunk vgl. \abbrArt{} 1 § 23 \abbrAbs{} 1 Satz 1 und 2 des Staatsvertrages über den Rundfunk im vereinten Deutschland vom 31. August 1991); ihr Programm hat daher auch teil an der Schutzaufgabe gegenüber dem ungeborenen Leben.
\end{sourcepara}
\begin{transpara}
\transrn{179}d) 最後，保護任務亦課予國家，在一般意識中維持並活化未出生生命之法律保護請求。因此，聯邦與各邦之國家機關必須可辨識地支持生命保護。這也且尤其涉及學校課程綱要。從事健康問題說明、家庭諮詢或性教育之公共機構，通常必須強化保護未出生生命之意願；這尤其適用於 SFHG 第1條第1項所規定之說明。公法廣播與私人廣播在行使其廣播自由（《基本法》第5條第1項\ggfnArtFiveOneZH{}）時，均受人性尊嚴拘束（關於私人廣播，參見1991年8月31日《統一德國廣播國家條約》第1條第23項第1句及第2句）；因此，其節目亦參與對未出生生命之保護任務。
\end{transpara}

\begin{sourcepara}[title={原文D.66}]
\sourcern{180}4. Nach dem unter 2. und 3. Dargelegten muß der Staat, um seiner Schutzpflicht für das ungeborene Leben zu genügen, ausreichende Maßnahmen normativer und tatsächlicher Art ergreifen, die dazu führen, daß ein -- unter Berücksichtigung entgegenstehender Rechtsgüter -- angemessener und als solcher wirksamer Schutz erreicht wird. Dazu bedarf es eines näher zu entwickelnden, Elemente des präventiven wie des repressiven Schutzes miteinander verbindenden Schutzkonzepts. Die Aufstellung und normative Umsetzung eines solchen Schutzkonzepts ist Sache des Gesetzgebers. Nach den bestehenden verfassungsrechtlichen Anforderungen ist er dabei nicht frei, den Schwangerschaftsabbruch über verfassungsrechtlich unbedenkliche Ausnahmetatbestände hinaus als nicht rechtswidrig, also erlaubt, anzusehen. Allerdings kann der Gesetzgeber nach noch näher auszuführenden Maßstäben entscheiden, wie er das grundsätzliche Verbot des Schwangerschaftsabbruchs im übrigen in den verschiedenen Bereichen der Rechtsordnung zur Wirkung bringt. Insgesamt muß aber das Schutzkonzept so ausgestaltet sein, daß es geeignet ist, den gebotenen Schutz zu entfalten, und nicht in eine -- zeitlich begrenzte -- rechtliche Freigabe des Schwangerschaftsabbruchs übergeht oder als solche wirkt.
\end{sourcepara}
\begin{transpara}
\transrn{180}4. 依前述第2點及第3點，國家為滿足其對未出生生命之保護義務，必須採取足夠之規範性及事實性措施，使在考量相反法益下達成適當且作為保護本身有效之保護。為此，需要一項更進一步發展、將預防性與壓制性保護要素相互結合之保護構想。此一保護構想之建立及規範實現，是立法者之事。依現行憲法要求，立法者在此並無自由將終止妊娠在憲法上無疑義之例外構成要件以外，視為不具違法性，亦即允許。不過，立法者得依下文尚待詳述之標準，決定其如何在其餘法秩序各領域中使終止妊娠之原則禁止發生效力。整體而言，保護構想必須被形塑為適於展開所要求之保護，且不得轉化為，或發生如同，對終止妊娠之有期限法律開放。
\end{transpara}

\begin{sourcepara}[title={原文D.67}]
\sourcern{181}Der Gesetzgeber muß der Wahl und Ausgestaltung seines Schutzkonzepts die verfassungsrechtlich tragfähige Einschätzung zugrunde legen, daß er mit ihm den Schutz des ungeborenen Lebens so gewährleisten werde, wie es das Untermaßverbot verlangt. Soweit seinen Entscheidungen zugleich Prognosen über tatsächliche Entwicklungen, insbesondere die Wirkungen seiner Regelung zugrundeliegen, müssen diese Prognosen verläßlich sein; das Bundesverfassungsgericht prüft, ob sie nach Maßgabe der nachfolgenden Kriterien vertretbar sind.
\end{sourcepara}
\begin{transpara}
\transrn{181}立法者在選擇及形塑其保護構想時，必須以憲法上可承載之評估為基礎，即其將以該構想依不足禁止所要求之方式保障未出生生命之保護。若其決定同時以關於事實發展，尤其其規制效果之預測為基礎，則此等預測必須可靠；聯邦憲法法院審查其依下列標準是否可支持。
\end{transpara}

\begin{sourcepara}[title={原文D.68}]
\sourcern{182}a) Ein Einschätzungs-, Wertungs- und Gestaltungsspielraum kommt dem Gesetzgeber auch dann zu, wenn er -- wie hier -- verfassungsrechtlich verpflichtet ist, wirksame und ausreichende Maßnahmen zum Schutz eines Rechtsguts zu ergreifen. Der Umfang dieses Spielraums hängt von Faktoren verschiedener Art ab, im besonderen von der Eigenart des in Rede stehenden Sachbereichs, den Möglichkeiten, sich -- zumal über künftige Entwicklungen wie die Auswirkungen einer Norm -- ein hinreichend sicheres Urteil zu bilden, und der Bedeutung der auf dem Spiel stehenden Rechtsgüter (vgl. \abbrBVerfGE{} 50, 290 [332 \abbrF{}]; 76, 1 [51 \abbrF{}]; 77, 170 [214 \abbrF{}]). Ob sich hieraus für die verfassungsrechtliche Überprüfung drei voneinander unterscheidbare Kontrollmaßstäbe herleiten lassen (vgl. \abbrBVerfGE{} 50, 290 [333]), bedarf keiner Erörterung; die verfassungsrechtliche Prüfung erstreckt sich in jedem Falle darauf, ob der Gesetzgeber die genannten Faktoren ausreichend berücksichtigt und seinen Einschätzungsspielraum "in vertretbarer Weise" gehandhabt hat. Die im Beschluß des Senats vom 29. Oktober 1987 (vgl. \abbrBVerfGE{} 77, 170 [214 \abbrF{}]) enthaltenen Ausführungen zur Zulässigkeit einer Verfassungsbeschwerde gegen staatliches Unterlassen dürfen nicht dahin verstanden werden, als genügten der Erfüllung der Schutzpflicht des Staates gegenüber menschlichem Leben schon Maßnahmen, "die nicht gänzlich ungeeignet oder völlig unzulänglich sind."
\end{sourcepara}
\begin{transpara}
\transrn{182}a) 即使立法者如同此處在憲法上負有採取有效且足夠措施以保護某一法益之義務，立法者仍享有評估、評價及形成空間。此一空間之範圍取決於不同種類因素，尤其取決於所涉事物領域之特性、形成足夠確定判斷之可能性，特別是關於未來發展例如規範影響之判斷可能性，以及處於風險中之法益的重要性（參見 BVerfGE 50, 290 [332 f.]; 76, 1 [51 f.]; 77, 170 [214 f.]）。是否可由此為憲法審查推導出三種可相互區分之審查標準（參見 BVerfGE 50, 290 [333]），無須討論；憲法審查在任何情形下均及於：立法者是否充分考量上述因素，並以「可支持之方式」運用其評估空間。第二庭1987年10月29日裁定（參見 BVerfGE 77, 170 [214 f.]）中關於針對國家不作為提起憲法訴願許可性之說明，不得被理解為，只要措施「並非全然不適合或完全不足」，即足以履行國家對人類生命之保護義務。
\end{transpara}

\begin{sourcepara}[title={原文D.69}]
\sourcern{183}b) Bei der verfassungsgerichtlichen Beurteilung, ob die gesetzgeberische Einschätzung der Wirksamkeit eines neuen Schutzkonzepts vertretbar ist, darf nicht außer Betracht bleiben, daß der Gesetzgeber in Erfüllung der ihm aus der Verfassung erwachsenden Pflicht zum Schutz des ungeborenen menschlichen Lebens handelt. Die hier in Rede stehenden Rechtsgüter des Ungeborenen und der Frau haben hohen verfassungsrechtlichen Rang. Das ist als Umstand, der die Eigenart des den Regelungsgegenstand bildenden Sachbereichs kennzeichnet, für die Beurteilung einer gesetzlichen Regelung zum Schutz des ungeborenen Lebens ebenso entscheidend wie die Tatsache, daß im Konfliktfall das ungeborene Leben getötet wird, wenn die schwangere Frau sich gegen die Fortsetzung der Schwangerschaft entscheidet. Von wesentlicher Bedeutung ist weiter, daß die Schwangerschaft in ihrer Frühphase oft nur der Mutter bekannt und das Ungeborene ihrem Schutz in jeder Hinsicht anvertraut, von ihm aber auch existenziell abhängig ist. Der Staat sieht sich vor die Aufgabe gestellt, Leben zu schützen, von dessen Vorhandensein er nichts weiß. Dies macht verständlich, warum die Erfahrungen mit allen bisherigen strafrechtlichen Regelungen wenig ermutigend sind. Ins Gewicht fällt schließlich, daß für den Gesetzgeber, wenn er sich zu einer grundlegend neuen Regelung entschließt, die Möglichkeiten, sich ein hinreichend sicheres Urteil über deren Auswirkungen zu bilden, naturgemäß eingeschränkt sind. Hierfür läßt sich auf Erfahrungen aus dem Ausland wegen der nicht gesicherten Vergleichbarkeit der jeweiligen Verhältnisse nur bedingt zurückgreifen. In solcher Lage muß sich der Gesetzgeber des erreichbaren, für die gebotene verläßliche Prognose der Schutzwirkung des Konzepts wesentlichen Materials bedienen und es mit der gebotenen Sorgfalt daraufhin auswerten, ob es seine gesetzgeberische Einschätzung hinreichend zu stützen vermag.
\end{sourcepara}
\begin{transpara}
\transrn{183}b) 憲法法院在評價立法者對新保護構想有效性之評估是否可支持時，不得忽略立法者係在履行其由憲法所生、保護未出生人類生命之義務。此處所涉未出生者及婦女之法益，具有高度憲法位階。這作為標示規制標的所屬事物領域特性之情事，對於評價保護未出生生命之法律規制，與下列事實同樣具有決定性：在衝突案件中，若懷孕婦女決定不繼續懷孕，未出生生命即遭殺害。進一步具有重大意義者是，懷孕在其早期往往只有母親知悉，未出生者在各方面均交付於其保護，並且亦在生存上依賴於其。國家面臨保護其存在並不為國家所知之生命的任務。這使人能理解，為何迄今所有刑法規制之經驗並不令人鼓舞。最後具有分量者是，當立法者決定採取一項根本新規制時，其形成關於該規制影響之足夠確定判斷的可能性，自然受到限制。就此而言，由於各自情況之可比較性並未確定，外國經驗僅能有限援用。在此種情況下，立法者必須使用可取得且對構想保護效果所要求之可靠預測具有重要性之材料，並以所要求之審慎加以評估，以判斷其是否足以支持其立法評估。
\end{transpara}

\bisubsubsection{II.}{II.}

\begin{sourcepara}[title={原文D.70}]
\sourcern{184}Nach dem Dargelegten ist es dem Gesetzgeber verfassungsrechtlich grundsätzlich nicht verwehrt, für den Schutz des ungeborenen Lebens zu einem Schutzkonzept überzugehen, das in der Frühphase der Schwangerschaft in Schwangerschaftskonflikten den Schwerpunkt auf die Beratung der schwangeren Frau legt, um sie für das Austragen des Kindes zu gewinnen, und dabei im Blick auf die notwendige Offenheit und Wirkung der Beratung auf eine indikationsbestimmte Strafdrohung und die Feststellung von Indikationstatbeständen durch einen Dritten verzichtet.
\end{sourcepara}
\begin{transpara}
\transrn{184}依前述，立法者在憲法上原則並未被禁止，為保護未出生生命而轉向一項保護構想；該構想在懷孕早期之懷孕衝突中，將重點置於對懷孕婦女之諮詢，以爭取其繼續懷孕，並且鑑於諮詢必要之開放性與效果，放棄由指標事由決定之刑罰威嚇，以及由第三人確認指標事由構成要件。
\end{transpara}

\begin{sourcepara}[title={原文D.71}]
\sourcern{185}Der Gesetzgeber des Schwangeren- und Familienhilfegesetzes hat den Wechsel im Schutzkonzept mit vertretbaren Einschätzungen vollzogen.
\end{sourcepara}
\begin{transpara}
\transrn{185}《孕婦及家庭扶助法》之立法者，是以可支持之評估完成保護構想之轉換。
\end{transpara}

\begin{sourcepara}[title={原文D.72}]
\sourcern{186}1. Eine Neuregelung der mit dem Schwangerschaftsabbruch zusammenhängenden Fragen ist durch \abbrArt{} 31 \abbrAbs{} 4 EV veranlaßt worden. Die Vorschrift stellt dem gesamtdeutschen Gesetzgeber die Aufgabe, "spätestens bis zum 31. Dezember 1992 eine Regelung zu treffen, die den Schutz vorgeburtlichen Lebens und die verfassungskonforme Bewältigung von Konfliktsituationen schwangerer Frauen vor allem durch rechtlich gesicherte Ansprüche für Frauen, insbesondere auf Beratung und soziale Hilfen, besser gewährleistet, als dies in beiden Teilen Deutschlands derzeit der Fall ist." Im Vordergrund standen dabei zwei Wege. Angesichts der Erfahrungen mit dem Vollzug der nach dem Urteil des Bundesverfassungsgerichts vom 25. Februar 1975 (\abbrBVerfGE{} 39, 1 \abbrFF{}) geschaffenen Indikationenregelung -- deren verfassungsrechtliche Problematik die im Verfahren 2 \abbrBvF{} 2/90 gestellten Anträge aufweisen -- konnte es der Gesetzgeber unternehmen, diese Lösung durch eine deutlichere, damit wohl notwendig auch engere Fassung der Indikationstatbestände zu ersetzen und an die Feststellung der Indikationen strengere Anforderungen zu stellen. Als anderer Weg kam eine Beratungsregelung in Betracht. Für sie kann der gesamtdeutsche Gesetzgeber anführen, daß sie besser geeignet erscheinen kann, die bisher von einer Fristenregelung einerseits und von einer Indikationenregelung andererseits gekennzeichneten deutschen Teilrechtsordnungen und das durch sie in unterschiedlicher Weise geprägte Rechtsbewußtsein der Bevölkerung zusammenzuführen.
\end{sourcepara}
\begin{transpara}
\transrn{186}1. 與終止妊娠相關問題之新規制，係由 EV 第31條第4項所促成。該規定交付全德立法者下列任務：「至遲於1992年12月31日前制定一項規制，其尤其透過婦女受法律保障之請求權，特別是諮詢與社會協助請求權，比德國兩部分目前情形更好地保障出生前生命之保護，以及懷孕婦女衝突狀態之合憲克服。」在此有兩條道路居於前景。鑑於依聯邦憲法法院1975年2月25日判決（BVerfGE 39, 1 ff.）所創設之指標規則的執行經驗，該規制之憲法問題已由 2 BvF 2/90 程序中所提聲請指出，立法者可以嘗試以更明確、因而大概也必然更狹窄之指標事由構成要件版本取代此一方案，並對指標事由之確認提出更嚴格要求。另一道路則是諮詢規制。對於諮詢規制，全德立法者可以主張，其似乎較適於整合德國兩部分法秩序，以及受其以不同方式形塑之人民法律意識；前者迄今一方面以期限規定為特徵，另一方面以指標規則為特徵。
\end{transpara}

\begin{sourcepara}[title={原文D.73}]
\sourcern{187}2. Der Gesetzgeber ist bei der Entwicklung eines neuen Schutzkonzepts nicht nur berechtigt, sondern auch verpflichtet, die Erfahrungen der bisherigen Rechtspraxis zu bewerten und hieran sein Schutzkonzept auszurichten. Auch unterschiedliche Formen eines weitgreifenden strafrechtlichen Schutzes für das ungeborene Leben auf der normativen Ebene -- sei es die strenge Abtreibungsregelung des § 218 \abbrStGB{} 1871, unter der die Rechtsprechung nur eine enge medizinische Indikation anerkannte, sei es die differenzierte Indikationenregelung seit 1976 -- haben nicht zu verhindern vermocht, daß Abtreibung eine Massenerscheinung gewesen und geblieben ist. Es kommt hierfür nicht darauf an, ob die höheren oder niedrigeren Einschätzungen der Abtreibungszahlen in früherer Zeit und in den Jahren seit 1976 die größere Verläßlichkeit und Wirklichkeitsnähe aufweisen; auch die niedrigeren Einschätzungen mußten den Gesetzgeber beunruhigen. Diese hohe Zahl der Abbrüche läßt sich auch nicht zureichend mit Ermittlungs- oder sonstigen Schwierigkeiten bei der Anwendung und Handhabung der jeweils geltenden strafrechtlichen Regelungen erklären oder mit dem Fehlen einer ernsthaften Vollzugsbereitschaft. Soweit solche Umstände als kontinuierliche Erscheinung feststellbar sind, besteht vielmehr gerade Anlaß, nach den Ursachen hierfür zu fragen und sich dem darin zu Tage tretenden Problem zu stellen.
\end{sourcepara}
\begin{transpara}
\transrn{187}2. 立法者在發展新保護構想時，不僅有權，而且也有義務評價既有法律實務經驗，並據此調整其保護構想。即使是在規範層面對未出生生命提供廣泛刑法保護之不同形式，無論是1871年《刑法》第218條之嚴格終止妊娠規制，在該規制下判例僅承認狹窄醫學指標事由，或是自1976年以來之區分化指標規則，均未能防止終止妊娠曾經且仍然是一種大量現象。就此而言，較高或較低之早期及1976年以來終止妊娠數量估計何者具有較大可靠性與現實接近性，並非關鍵；即使較低估計也必須使立法者憂慮。此等高數量之終止妊娠，也不能充分以各該有效刑法規制之適用與操作上的偵查困難或其他困難，或以欠缺認真執行意願來說明。只要此等情事可被確認為持續現象，毋寧正有理由追問其原因，並面對其中顯現之問題。
\end{transpara}

\begin{sourcepara}[title={原文D.74}]
\sourcern{188}Verfassungsrechtlich nicht zu beanstanden ist die aus dieser Analyse gewonnene Einschätzung des Gesetzgebers, die Ausgestaltung der Strafdrohung im geltenden Recht wirke im Falle eines Schwangerschaftskonflikts einer dem Austragen des Kindes aufgeschlossenen Entscheidung der Frau eher entgegen, weil die schwangere Frau ihren Konflikt als höchstpersönlichen erlebt und sich gegen seine Beurteilung und Bewertung durch Dritte wehrt. Die Umstände erheblichen Gewichts, die einer Frau das Austragen eines Kindes bis zur Unzumutbarkeit erschweren können, bestimmen sich nicht nur nach objektiven Komponenten, sondern auch nach ihren physischen und psychischen Befindlichkeiten und Eigenschaften. Eine hinreichend zuverlässige und nachvollziehbare Beurteilung dieser nur in ihrem Zusammenwirken erheblichen und möglicherweise untereinander wieder gegenläufigen Zumutbarkeitskriterien kann daher nur mit erheblichem Aufwand von erfahrenem Sachverstand möglich sein. Je mehr aber Dritte in den innersten Bereich einer Frau eindringen, um so größer wird die Gefahr, daß die Frau sich dem durch Vorschieben von anderen Gründen oder Ausweichen in die Illegalität entzieht. In einem solchen Fall ist aber von vornherein jede Chance vertan, durch verständnisvolle, sachkundige Beratung mit der Frau deren Konflikt zu ergründen und ihr bei seiner Bewältigung so zu helfen, daß sie sich nicht gegen das Kind entscheidet.
\end{sourcepara}
\begin{transpara}
\transrn{188}立法者由此一分析所得之評估，在憲法上無可指摘：現行法中刑罰威嚇之形塑，在懷孕衝突案件中毋寧妨礙婦女作成願意繼續懷孕之決定，因為懷孕婦女將其衝突經驗為極具個人性，並抗拒第三人對其加以判斷及評價。可能使婦女繼續懷孕困難至不可期待程度之重大情事，並非僅由客觀成分決定，也由其身體與心理狀態及特性決定。因此，對此等只有在相互作用中才具有重要性，且可能彼此又相互牴觸之可期待性標準，要作成足夠可靠且可理解之判斷，只能藉由有經驗之專業知識投入相當成本始有可能。然而，第三人越深入婦女最內在領域，婦女即越可能以提出其他理由或逃入非法方式加以迴避。在此種情形下，透過理解且專業之諮詢與婦女共同探究其衝突，並協助其克服衝突以致其不作出反對子女之決定的任何機會，從一開始即已喪失。
\end{transpara}

\begin{sourcepara}[title={原文D.75}]
\sourcern{189}3. Es ist daher eine verfassungsrechtlich unbedenkliche Einschätzung des Gesetzgebers, wenn er sich zur Erfüllung seines Schutzauftrags einem Schutzkonzept zuwendet, das davon ausgeht, jedenfalls in der Frühphase der Schwangerschaft sei ein wirksamer Schutz des ungeborenen menschlichen Lebens nur mit der Mutter, aber nicht gegen sie möglich. Sie allein und nur von ihr selbst ins Vertrauen Gezogene wissen in diesem Stadium der Schwangerschaft um das neue Leben, das noch ganz der Mutter zugehört und von ihr in allem abhängig ist. Diese Unentdecktheit, Hilflosigkeit und Abhängigkeit des auf einzigartige Weise mit der Mutter verbundenen Ungeborenen lassen die Einschätzung berechtigt erscheinen, daß der Staat eine bessere Chance zu seinem Schutz hat, wenn er mit der Mutter zusammenwirkt.
\end{sourcepara}
\begin{transpara}
\transrn{189}3. 因此，立法者為履行其保護任務而轉向一項保護構想，並以至少在懷孕早期，未出生人類生命之有效保護只能與母親一起，而不能對抗母親而實現為出發點，乃屬憲法上無疑義之評估。在此懷孕階段，只有她本人及其本人信任之人知悉此一新生命；該生命尚完全屬於母親，並在一切方面依賴於她。以獨特方式與母親相連之未出生者，其未被發現、無助與依賴，使下列評估顯得正當：若國家與母親合作，國家對其保護即有較佳機會。
\end{transpara}

\begin{sourcepara}[title={原文D.76}]
\sourcern{190}Diese Einschätzung kann sich auch darauf gründen, daß eine Frau, die von einer ungewollten Schwangerschaft erfährt, in besonderer Weise existenziell betroffen sein kann. Sie muß unter Umständen ihre Lebensplanung tiefgreifend ändern und sieht sich weit über die mit einer Schwangerschaft unvermeidlich einhergehenden Beschwerden hinaus im einzelnen unübersehbaren und lange andauernden Handlungs- und Sorgepflichten, gegebenenfalls auch zusätzlichen Lebensrisiken ausgesetzt. Zudem ist eine Frau in der Frühphase einer solchen Schwangerschaft psychisch oft noch nicht auf die Mutterschaft eingestellt und empfindet noch nicht die Bindung an das in ihr wachsende Leben wie in einem späteren Stadium, in dem sie es in sich verspürt. Eine Strafandrohung bewirkt hier wenig, so daß es um so näher liegt, der Frau mit präventiven Mitteln des Rechts bei der Bewältigung ihres Konflikts zu helfen, damit sie ihrer Verantwortung für das Ungeborene gerecht werden kann. Deshalb kann die besondere Lage der Frau und des Ungeborenen in dieser Frühphase der Schwangerschaft zwar Anlaß sein, unter Rücknahme strafrechtlicher Sanktionen besondere Schutzmaßnahmen zu ergreifen; sie darf aber -- wie dargelegt -- nicht dazu führen, die Grundrechtsposition der Frau denen des ungeborenen Lebens überzuordnen. Liegt die Würde des Menschseins auch für das ungeborene Leben im Dasein um seiner selbst willen, so verbieten sich jegliche Differenzierungen der Schutzverpflichtung mit Blick auf Alter und Entwicklungsstand dieses Lebens oder die Bereitschaft der Frau, es weiter in sich leben zu lassen.
\end{sourcepara}
\begin{transpara}
\transrn{190}此一評估亦可基於下列理由：知悉非意願懷孕之婦女，可能以特別方式在生存上受到觸動。其在特定情形下必須深刻改變自身生活規劃，並面臨個別上難以預見且長期持續之作為與照護義務，在必要時亦面臨額外生活風險；此已遠超過懷孕不可避免伴隨之不適。此外，婦女在此種懷孕早期，心理上往往尚未調適為母親，且尚未如在較後階段、當她在自身內感受該生命時一般，感受到與其內部成長生命之連結。刑罰威嚇在此效果有限，因此更接近的是，以法律之預防性手段協助婦女克服其衝突，使其能夠履行對未出生者之責任。因此，婦女與未出生者在懷孕此一早期之特殊處境，固然得成為在撤回刑事制裁下採取特別保護措施之理由；但如前所述，其不得導致將婦女之基本權地位置於未出生生命之地位之上。若人之存在的尊嚴對未出生生命而言亦在於其為自身而存在，則任何著眼於此一生命之年齡與發展狀態，或婦女是否願意讓其繼續在自身內生活，而對保護義務作出差別化，均被禁止。
\end{transpara}

\begin{sourcepara}[title={原文D.77}]
\sourcern{191}4. Es steht mit der einer Frau und werdenden Mutter gebührenden Achtung in Einklang, wenn der Staat Frauen nicht durch generelle Drohung mit Strafe, sondern durch individuelle Beratung und einen Appell an ihre Verantwortung gegenüber dem ungeborenen Leben, durch wirtschaftliche und soziale Förderung und darauf bezogene qualifizierte Information dafür zu gewinnen sucht, sich der Aufgabe als Mutter nicht zu entziehen. Dabei kann der Gesetzgeber davon ausgehen, daß dies eher behindert als gefördert wird, wenn ein Dritter die Gründe, aus denen eine Frau das Austragen ihres Kindes als unzumutbar ansieht, überprüfen und bewerten müßte.
\end{sourcepara}
\begin{transpara}
\transrn{191}4. 國家若不是透過一般性刑罰威嚇，而是透過個別諮詢、訴諸婦女對未出生生命之責任、經濟與社會促進，以及與此相關之合格資訊，試圖爭取婦女不逃避作為母親之任務，這與婦女及準母親所應受之尊重相符。在此，立法者得以為，若須由第三人審查並評價婦女認為繼續懷孕不可期待之理由，則此較可能妨礙而非促進上述目標。
\end{transpara}

\begin{sourcepara}[title={原文D.78}]
\sourcern{192}Damit ist nicht gesagt, daß nach fachkundiger und individueller Beratung nur noch die Frau ihre Schwangerschaft abbricht, die sich in einem Konflikt solchen Ausmaßes befindet, daß ihr das Austragen des Kindes nach den oben dargelegten verfassungsrechtlichen Maßstäben unzumutbar ist. Dies würde die Lebenswirklichkeit verkennen, in der Männer wie Frauen vielfach ihre eigenen Lebensvorstellungen überbewerten und diese auch dann nicht zurückzustellen bereit sind, wenn es bei objektivem Nachvollziehen ihrer individuellen Lebenssituation zumutbar erscheint. Gleichwohl darf der Gesetzgeber aus den genannten Gründen (vgl. oben 2. und 3.) annehmen, auch eine Indikationenüberprüfung im ersten Stadium der Schwangerschaft könnte nicht besser verhindern, daß Frauen sich einseitig an ihren persönlichen Interessen orientieren; demgegenüber könne aber die im Überlassen einer Letztverantwortung zum Ausdruck kommende Achtung vor dem Verantwortungsbewußtsein der Frauen Appellfunktion haben und geeignet sein, allgemein die Verantwortung von Frauen gegenüber dem ungeborenen Leben zu stärken, sofern dies vor dem Hintergrund einer wachgehaltenen Orientierung über die verfassungsrechtlichen Grenzen von Recht und Unrecht geschehe. Der Gesetzgeber darf berücksichtigen, daß Frauen, an die diese Erwartungen gestellt werden, ihre Verantwortung unmittelbarer und stärker empfinden und daher eher Anlaß zu ihrer gewissenhaften Ausübung haben können, als wenn ein Dritter einen ihm genannten Grund -- mehr oder weniger eingehend -- überprüft und bewertet und mit der Feststellung, der Abbruch sei aufgrund eines Indikationstatbestandes erlaubt, der Frau ein Stück Verantwortung abnimmt.
\end{sourcepara}
\begin{transpara}
\transrn{192}這並不是說，經專業且個別之諮詢後，只有處於如此程度衝突中，以致依上述憲法標準繼續懷孕對其不可期待之婦女，才會終止妊娠。如此將誤認生活現實；在該現實中，男性與女性往往過度評價其自身生活想像，且即使在客觀追認其個別生活處境時顯得可以期待，也不願將其退讓。儘管如此，立法者基於上述理由（參見上文 2. 及 3.）得認為，即使在懷孕第一階段進行指標事由審查，也不能更好防止婦女單方面以其個人利益為取向；相反地，在交付最後責任中所表達之對婦女責任意識的尊重，則可能具有呼籲功能，並適於一般性強化婦女對未出生生命之責任，只要這是在持續維持關於法與不法憲法界限之取向背景下發生。立法者得考量，承受此等期待之婦女，相較於由第三人或多或少詳盡審查並評價其所提出之理由，並以確認終止妊娠因指標事由構成要件而被允許之方式替婦女取走一部分責任的情形，可能更直接且更強烈感受其責任，因而更有理由良心謹慎地行使該責任。
\end{transpara}

\begin{sourcepara}[title={原文D.79}]
\sourcern{193}5. Überläßt der Gesetzgeber jenen Frauen, die sich der Beratung unterziehen, die Letztverantwortung für den Schwangerschaftsabbruch und ermöglicht er es ihnen, gegebenenfalls für den Abbruch einen Arzt in Anspruch zu nehmen, so kann er vertretbar daran die Erwartung knüpfen, daß Schwangere im Konfliktfall die Beratung annehmen und ihre Situation offenlegen.
\end{sourcepara}
\begin{transpara}
\transrn{193}5. 若立法者將終止妊娠之最後責任交由接受諮詢之婦女，並使其在必要時得為終止妊娠求助於醫師，則其得可支持地連結下列期待：懷孕婦女在衝突案件中將接受諮詢並揭示其處境。
\end{transpara}

\begin{sourcepara}[title={原文D.80}]
\sourcern{194}a) Der Gesetzgeber kann nicht zuletzt Erfahrungen der Beratungspraxis dafür anführen, daß selbst eine mit der Beratung nicht verbundene, sondern ihr nachfolgende, personell und institutionell von ihr getrennte Indikationsfeststellung eine ungünstige Vorwirkung auf die Beratung ausübt und deren Wirkungschancen wesentlich beeinträchtigt, indem die Frau auf den festzustellenden Indikationstatbestand fixiert und eine offene Erörterung des Schwangerschaftskonflikts verhindert wird.
\end{sourcepara}
\begin{transpara}
\transrn{194}a) 立法者尤其得援引諮詢實務經驗說明，即使指標事由確認並非與諮詢相結合，而是在諮詢之後進行，且在人員及制度上均與諮詢分離，其仍會對諮詢產生不利預先效果，並重大損害諮詢之作用機會，因為婦女會固定於待確認之指標事由構成要件，且妨礙對懷孕衝突作開放討論。
\end{transpara}

\begin{sourcepara}[title={原文D.81}]
\sourcern{195}b) Diese Erwägungen gelten allerdings nur für die allgemeine Notlagenindikation. Bei ihr handelt es sich meist um Notlagen, die sich als Produkt vielfältiger und vielschichtiger, auch dem persönlichen Bereich der Schwangeren zuzurechnender Faktoren ergeben. Sie müßten in gleicher Weise Gegenstand der Drittfeststellung und Bewertung im Verfahren der Indikationsfeststellung sein, wie sie in einer -- zum Lebensschutz gebotenen -- Konfliktberatung aufgedeckt und erörtert werden sollen. Dies macht deutlich, daß die unbefangene Mitarbeit der Frau an der Bewältigung dieses Konflikts durch das Erfordernis seiner Drittbeurteilung behindert sein kann. Anders liegt es bei einer medizinischen, embryopathischen oder auch kriminologischen Indikation. Die Notlage ist hier greifbar, wenn ärztlich festgestellt ist, daß sich die Frau mit der Fortsetzung der Schwangerschaft einer beachtlichen Gefahr für ihre Gesundheit aussetzt oder die erhebliche Gefahr einer schweren Schädigung des Kindes besteht, oder wenn die Frau Opfer einer Straftat ist. Der vorgegebene objektive Sachverhalt verändert auch die Thematik und die Funktion der Beratung so, daß Frauen kaum Anlaß haben, einer solchen Indikationsfeststellung auszuweichen oder die Beratung nicht mit der notwendigen Offenheit anzunehmen.
\end{sourcepara}
\begin{transpara}
\transrn{195}b) 不過，此等考量僅適用於一般困境指標事由。其通常涉及由多樣且多層次因素所形成之困境，且此等因素亦可歸屬於懷孕婦女個人領域。在指標事由確認程序中，此等因素必須以同樣方式成為第三人確認及評價之標的，正如其應在一項為生命保護所要求之衝突諮詢中被揭露並討論。這清楚顯示，婦女在克服此一衝突時之坦然合作，可能因其衝突必須受第三人評價之要求而受妨礙。醫學、胚胎病理或犯罪學指標事由則情形不同。在此，若經醫學確認，婦女繼續懷孕將使其健康承受可觀危險，或存在子女受嚴重損害之重大危險，或者婦女為犯罪被害人，困境即為可掌握。此一既定客觀事實亦改變諮詢之主題與功能，使婦女幾乎沒有理由迴避此種指標事由確認，或不以必要之開放性接受諮詢。
\end{transpara}

\begin{sourcepara}[title={原文D.82}]
\sourcern{196}c) Wissenschaftlich und rechtspolitisch umstritten war und ist allerdings, ob eine Beratungsregelung für Schwangerschaftsabbrüche in der Frühphase der Schwangerschaft eine bessere Schutzwirkung für das ungeborene Leben entfalten kann als die bisherige Regelung. Zwar besteht Übereinstimmung darüber, daß die bisherige Indikationenregelung in ihrer praktischen Handhabung einen ungenügenden Schutz gegen Schwangerschaftsabbrüche geboten habe, jedoch ist die Einschätzung, die Notwendigkeit einer Indikationsfeststellung mindere die Schutzwirkung der Schwangerschaftsberatung, in Wissenschaft und Beratungspraxis unterschiedlich. Angesichts der dargelegten Gründe, die gegen die Beibehaltung der bisherigen Indikationenregelung sprechen, hindern solche Ungewißheiten den Gesetzgeber jedoch nicht grundsätzlich daran, eine Beratungsregelung einzuführen; freilich ist er gehalten, die Auswirkungen seines neuen Schutzkonzepts im Auge zu behalten (Beobachtungs- und Nachbesserungspflicht).
\end{sourcepara}
\begin{transpara}
\transrn{196}c) 不過，懷孕早期終止妊娠之諮詢規制，是否能比既有規制對未出生生命發揮更佳保護效果，在科學及法律政策上曾經且仍有爭議。誠然，關於既有指標規則在其實務操作上對終止妊娠提供不足保護，已有共識；然而，指標事由確認之必要性會削弱懷孕諮詢保護效果此一評估，在學術與諮詢實務中並不一致。鑑於前述反對維持既有指標規則之理由，此等不確定性並不原則上阻礙立法者導入諮詢規制；當然，立法者有義務持續注意其新保護構想之影響（觀察及補正義務）。
\end{transpara}

\bisubsubsection{III.}{III.}

\begin{sourcepara}[title={原文D.83}]
\sourcern{197}Geht der Gesetzgeber in Erfüllung seiner Schutzpflicht zu einem Beratungskonzept über, so bedeutet das, daß die Schutzwirkung für das ungeborene Leben maßgeblich -- präventiv -- durch eine beratende Einflußnahme auf die einen Schwangerschaftsabbruch erwägende Frau erreicht werden soll. Das Beratungskonzept ist darauf angelegt, das Verantwortungsbewußtsein der Frau zu stärken, die -- unbeschadet der Verantwortlichkeiten des familiären und weiteren sozialen Umfeldes sowie des Arztes (vgl. unten V. und VI.) -- letztlich den Abbruch der Schwangerschaft tatsächlich bestimmt und insofern verantworten muß (Letztverantwortung). Das erfordert Rahmenbedingungen, die positive Voraussetzungen für ein Handeln der Frau zugunsten des ungeborenen Lebens schaffen. Nur dann kann trotz des Verzichts auf eine Feststellung von Indikationstatbeständen als Voraussetzung für einen Schwangerschaftsabbruch gleichwohl von einer Schutzwirkung des Beratungskonzepts für das ungeborene Leben ausgegangen werden (1.). Allerdings ist es nicht zulässig, nicht indizierte Schwangerschaftsabbrüche für gerechtfertigt (nicht rechtswidrig) zu erklären, deren Vornahme Frauen nach Beratung in den ersten zwölf Wochen von einem Arzt verlangen (2.). Im übrigen ist der Gesetzgeber nicht gehalten, die sich aus dem grundsätzlichen Verbot des Schwangerschaftsabbruchs an sich aufdrängenden Folgerungen in jeder Hinsicht zu ziehen, wenn das Beratungskonzept um seiner Wirksamkeit willen bestimmte Ausnahmen fordert (3.).
\end{sourcepara}
\begin{transpara}
\transrn{197}若立法者為履行其保護義務而轉向諮詢構想，這表示對未出生生命之保護效果，主要應以預防方式，透過對考慮終止妊娠之婦女施加諮詢影響而達成。諮詢構想旨在強化婦女之責任意識；婦女不妨礙家庭及較廣泛社會環境以及醫師之責任（參見下文 V. 及 VI.），終究在事實上決定懷孕之終止，並就此必須負責（最後責任）。這要求框架條件，使婦女為未出生生命利益而行動之正面前提得以形成。唯有如此，儘管放棄以指標事由構成要件之確認作為終止妊娠之前提，仍可認為諮詢構想對未出生生命具有保護效果（1.）。然而，婦女經諮詢後在最初十二週內要求醫師實施之非指標事由終止妊娠，不得被宣告為正當化（不具違法性）（2.）。此外，若諮詢構想為其有效性而要求特定例外，立法者並非有義務在所有方面均作出終止妊娠原則禁止本身所迫切導出之結論（3.）。
\end{transpara}

\begin{sourcepara}[title={原文D.84}]
\sourcern{198}1. a) Zu den notwendigen Rahmenbedingungen eines Beratungskonzepts gehört an erster Stelle, daß die Beratung für die Frau zur Pflicht gemacht wird und ihrerseits darauf ausgerichtet ist, die Frau zum Austragen des Kindes zu ermutigen. Dabei muß die Beratung nach Inhalt, Durchführung und Organisation geeignet sein, der Frau die Einsichten und Informationen zu vermitteln, deren sie für eine verantwortliche Entscheidung über die Fortsetzung oder den Abbruch der Schwangerschaft bedarf (siehe dazu im einzelnen unten IV.).
\end{sourcepara}
\begin{transpara}
\transrn{198}1. a) 諮詢構想之必要框架條件首先包括：諮詢必須對婦女成為義務，且其本身必須以鼓勵婦女繼續懷孕為取向。在此，諮詢依其內容、實施及組織，必須適於向婦女傳達其為就繼續或終止妊娠作成負責任決定所需之認識與資訊（詳見下文 IV.）。
\end{transpara}

\begin{sourcepara}[title={原文D.85}]
\sourcern{199}b) Darüber hinaus müssen die Personen in das Schutzkonzept einbezogen werden, die -- sei es positiv, sei es negativ -- in einem Schwangerschaftskonflikt auf den Willen der Frau Einfluß nehmen können. Dies gilt insbesondere für den Arzt, den die Schwangere zur Durchführung des Abbruchs aufsucht. Außer ihr und der Person, die sie beraten hat, weiß oft nur er um die Existenz des Ungeborenen; zudem ist der Arzt nach ärztlichem Berufsverständnis ohnehin zu dessen Schutz berufen (siehe dazu unten V.). In das Schutzkonzept einzubeziehen sind aber auch Personen des familiären und des weiteren sozialen Umfeldes einer schwangeren Frau, die diese, wie aus Berichten von Beratern und Ärzten sowie aus wissenschaftlichen Untersuchungen hervorgeht, häufig -- und dies nicht selten in strafwürdiger Weise -- gegen das Kind beeinflussen (siehe dazu unten VI.).
\end{sourcepara}
\begin{transpara}
\transrn{199}b) 此外，必須將可能在懷孕衝突中對婦女意志產生影響之人，不論其影響為正面或負面，納入保護構想。這尤其適用於懷孕婦女為實施終止妊娠而尋求之醫師。除懷孕婦女本人及曾諮詢她之人外，往往只有醫師知悉未出生者之存在；此外，依醫師職業理解，醫師本即受召喚保護未出生者（詳見下文 V.）。但亦應納入保護構想者，還包括懷孕婦女之家庭及較廣泛社會環境中之人；依諮詢人員及醫師報告以及科學研究所示，這些人經常，且不少時候以值得處罰之方式，影響懷孕婦女作出不利於子女之決定（詳見下文 VI.）。
\end{transpara}

\begin{sourcepara}[title={原文D.86}]
\sourcern{200}c) Aus den zu D.II.5.a) und b) genannten Gründen muß die Beratungsregelung davon absehen, den Rechtfertigungsgrund einer allgemeinen Notlagenindikation vorzusehen. Er würde ihrem Konzept zuwiderlaufen. Die Beratungsregelung will wirksamen Schutz dadurch erreichen, daß sie die Frau um ihrer Aufgeschlossenheit gegenüber der Beratung willen davor bewahrt, eine sie rechtfertigende Notlage darlegen und sich deren Feststellung unterziehen zu müssen. Das hat unvermeidlich zur Folge, daß die Beratungsregelung die Möglichkeit einer Rechtfertigung durch die allgemeine Notlagenindikation nicht verheißen kann. Nur wenn die Beratungsregelung generell und ausnahmslos von der Feststellung des Vorliegens einer sozialen Notlage absieht, kann sie erreichen, daß Frauen die Beratung annehmen und sich ihr nicht im Blick auf eine erstrebte Beurteilung ihrer Entscheidung als rechtmäßig -- und daran möglicherweise geknüpfte günstige Rechtsfolgen -- verschließen. Die Beratungsregelung mutet es daher Frauen zu, auf die persönliche Entlastung zu verzichten, die in einer Entscheidung über die Rechtmäßigkeit des von ihnen beabsichtigten Abbruchs liegen kann, auch wenn bei ihnen im Einzelfall eine allgemeine Notlage ohne weiteres nachvollziehbar erscheinen mag.
\end{sourcepara}
\begin{transpara}
\transrn{200}c) 基於 D.II.5.a) 及 b) 所述理由，諮詢規制必須放棄規定一般困境指標事由之正當化事由。該正當化事由將違反其構想。諮詢規制意在透過下列方式達成有效保護：為使婦女對諮詢保持開放，而避免其必須陳述使其正當化之困境並接受該困境之確認。這不可避免地導致，諮詢規制不能允諾透過一般困境指標事由而獲得正當化之可能性。只有當諮詢規制一般且毫無例外地放棄確認社會困境之存在時，才可能達成使婦女接受諮詢，並且不因其期待自身決定被評價為合法以及可能與之相連之有利法律效果，而對諮詢封閉。因此，諮詢規制要求婦女放棄一種個人免責確認；此種確認可能存在於關於其所意圖終止妊娠之合法性決定中，即便在個案中其一般困境可能顯得毫無困難即可理解。
\end{transpara}

\begin{sourcepara}[title={原文D.87}]
\sourcern{201}Die Grundrechtspositionen der Frau fordern einen Rechtfertigungsgrund der allgemeinen Notlagenindikation nicht; ihnen kann auch in anderer Weise entsprochen werden (vgl. dazu D. III.3.). Doch darf die Orientierung über die Rechtspflicht zum Austragen des Kindes und ihre Grenzen auch bei einer Beratungsregelung, die unvermeidlich auf eine Notlagenindikation verzichtet, nicht verlorengehen. Auch im konkreten Schwangerschaftskonflikt kann die normative Orientierung auf den Schutz des ungeborenen Lebens nicht entfallen; der verfassungsrechtliche Rang des Rechtsguts des ungeborenen menschlichen Lebens muß dem allgemeinen Rechtsbewußtsein weiterhin gegenwärtig bleiben (sog. positive Generalprävention). Auch eine Beratungsregelung muß daher in der Rechtsordnung unterhalb der Verfassung zum Ausdruck bringen, daß ein Schwangerschaftsabbruch nur in Ausnahmesituationen rechtmäßig sein kann, wenn der Frau durch das Austragen des Kindes eine Belastung erwächst, die -- vergleichbar den Fällen der medizinischen und embryopathischen Indikation (§ 218a \abbrAbs{} 2 und 3 \abbrStGB{} \abbrNF{}) -- so schwer und außergewöhnlich ist, daß sie die zumutbare Opfergrenze übersteigt. Diese Orientierung gibt der verantwortlich handelnden Frau die Grundlage, ihr Handeln zu beurteilen. Dies gerade ist der Kern der Verantwortung, die der Frau mit einer Beratungsregelung überlassen ist; aus ihrer Wahrnehmung kann freilich eine Rechtfertigung nicht folgen (vgl. D.III.2.b.aa).
\end{sourcepara}
\begin{transpara}
\transrn{201}婦女之基本權地位並不要求一般困境指標事由之正當化事由；其亦可以其他方式獲得相應對待（參見 D.III.3.）。然而，即使在一項不可避免放棄困境指標事由之諮詢規制中，關於繼續懷孕之法律義務及其界限之取向亦不得喪失。即使在具體懷孕衝突中，對未出生生命保護之規範取向亦不能消失；未出生人類生命此一法益之憲法位階，必須繼續呈現於一般法律意識中（所謂積極一般預防）。因此，諮詢規制亦必須在憲法以下之法秩序中表達：終止妊娠只有在例外情境中始可能合法，亦即婦女因繼續懷孕而承受之負擔，類似醫學及胚胎病理指標事由案件（新版本《刑法》第218a條第2項及第3項），如此嚴重且異常，以致超過可期待犧牲界限。此一取向賦予負責任行動之婦女評價其行動之基礎。這正是諮詢規制交由婦女之責任核心；但從其行使中當然不能導出正當化（參見 D.III.2.b.aa）。
\end{transpara}

\begin{sourcepara}[title={原文D.88}]
\sourcern{202}d) Schließlich verlangt eine Beratungsregelung, die vorwiegend auf präventiven Schutz setzt, daß ein Angebot sozialer Hilfen für Mutter und Kind -- auch tatsächlich -- bereitsteht, das durch Beseitigung oder Linderung konkreter Bedrängnisse und sozialer Nöte eine Entscheidung der Eltern für das Kind stützen und die Frau zum Austragen des Kindes ermutigen kann (vgl. oben D.I.3.).
\end{sourcepara}
\begin{transpara}
\transrn{202}d) 最後，一項主要依靠預防性保護之諮詢規制，要求母親與子女之社會協助提供在事實上亦已準備就緒；此種協助能透過消除或減輕具體壓迫與社會困境，支持父母作成有利於子女之決定，並鼓勵婦女繼續懷孕（參見上文 D.I.3.）。
\end{transpara}

\begin{sourcepara}[title={原文D.89}]
\sourcern{203}2. Das mit einem Beratungskonzept verbundene Ziel, Schwangerschaftsabbrüche, die während der ersten zwölf Wochen nach Beratung auf Verlangen der Schwangeren -- ohne Feststellung von Indikationen -- von einem Arzt vorgenommen werden, nicht mit Strafe zu bedrohen, kann der Gesetzgeber nur erreichen, indem er diese Schwangerschaftsabbrüche aus dem Tatbestand des § 218 \abbrStGB{} ausnimmt; sie können nicht für gerechtfertigt (nicht rechtswidrig) erklärt werden.
\end{sourcepara}
\begin{transpara}
\transrn{203}2. 與諮詢構想相連之目標，是對於在最初十二週內、經諮詢後依懷孕婦女要求且未確認指標事由而由醫師實施之終止妊娠，不施以刑罰威嚇；立法者只有透過將此等終止妊娠排除於《刑法》第218條構成要件之外，始能達成此一目標；其不得被宣告為正當化（不具違法性）。
\end{transpara}

\begin{sourcepara}[title={原文D.90}]
\sourcern{204}a) Darf der Schwangerschaftsabbruch nach der Verfassung nur bei Vorliegen bestimmter Ausnahmetatbestände erlaubt werden, so kann er nicht zugleich im Strafrecht unter anderen, weitergehenden Voraussetzungen als erlaubt angesehen werden. Die Rechtsordnung muß das verfassungsrechtliche Verbot des Schwangerschaftsabbruchs bestätigen und verdeutlichen. Dazu dient insbesondere das Strafrecht, das Rechtsgüter von besonderem Rang und in einer besonderen Gefährdungslage schützt und das allgemeine Bewußtsein von Recht und Unrecht am deutlichsten prägt. Wenn das Strafrecht einen Rechtfertigungsgrund vorsieht, muß das im allgemeinen Rechtsbewußtsein so verstanden werden, als sei das im Rechtfertigungstatbestand bezeichnete Verhalten erlaubt. Auch die Rechtsordnung im übrigen würde bei den jeweiligen Regelungen über Recht und Unrecht in ihren Teilbereichen davon ausgehen, daß der Schutz für dieses Leben durch die strafrechtlichen Rechtfertigungsgründe aufgehoben sei. Damit wäre der verfassungsrechtlichen Schutzpflicht nicht genügt. Die Durchschlagskraft, die einem strafrechtlichen Rechtfertigungsgrund für die gesamte Rechtsordnung jedenfalls dann zukommt, wenn es sich um den Schutz elementarer Rechtsgüter handelt, schließt es aus, ihn in seinen Wirkungen allein auf das Strafrecht zu beschränken. Ein Schwangerschaftsabbruch darf deshalb strafrechtlich nur für gerechtfertigt erklärt werden, wenn und soweit die Rechtfertigungsgründe auf die verfassungsrechtlich zugelassenen Ausnahmen vom Verbot des Schwangerschaftsabbruchs tatbestandlich begrenzt sind.
\end{sourcepara}
\begin{transpara}
\transrn{204}a) 若依憲法，終止妊娠只有在特定例外事由成立時方得被許可，則其在刑法中亦不得以其他更寬泛之要件被認定為許可。法秩序必須確認並清楚彰顯終止妊娠之憲法禁令。此目的尤其由刑法達成，因刑法保護特別位階且處於特殊危險情境之法益，並最清楚形塑關於法與不法之一般意識。若刑法設有正當化事由，則在一般法律意識中必須被理解為：正當化事由構成要件所描述之行為係被允許的。其餘法秩序在各分領域關於法與不法之規制中，亦將以對此一生命之保護已因刑法正當化事由而撤除為前提。如此即未滿足憲法上保護義務。刑法正當化事由至少在涉及基本法益之保護時，對整體法秩序所具有之貫穿力，排除了將其效果僅侷限於刑法的可能。因此，終止妊娠在刑法上只有在正當化事由以構成要件方式限於憲法所允許之終止妊娠禁令例外時，方得被認定為正當化。
\end{transpara}

\begin{sourcepara}[title={原文D.91}]
\sourcern{205}Werden hingegen Schwangerschaftsabbrüche unter bestimmten Voraussetzungen aus dem Straftatbestand ausgeklammert, so bedeutet dies lediglich, daß sie nicht mit Strafe bedroht sind. Eine Entscheidung des Gesetzgebers darüber, ob der Schwangerschaftsabbruch in anderen Teilen der Rechtsordnung als rechtmäßig oder rechtswidrig anzusehen und zu behandeln ist, bleibt damit offen (vgl. Lenckner in: Schönke/Schröder, Strafgesetzbuch, 24. \abbrAufl{}, 1991, Vorbem. zu §§ 13 \abbrFF{} \abbrRdnr{} 18; Eser/Burkhardt, Strafrecht I, 4. \abbrAufl{}, 1992, \abbrNr{} 9, \abbrRdnr{} 41). In anderen Bereichen der Rechtsordnung können dann eigenständige Regelungen getroffen werden, die dort den Schwangerschaftsabbruch als rechtswidrig zugrunde legen. Geschieht dies allerdings nicht, so wirkt sich der Ausschluß des Straftatbestandes wie ein Rechtfertigungsgrund aus, womit den Mindestanforderungen der Schutzpflicht nicht mehr genügt wäre.
\end{sourcepara}
\begin{transpara}
\transrn{205}反之，若終止妊娠在特定前提下被排除於刑事構成要件之外，則僅意味著其不受刑罰威脅。立法者關於終止妊娠在法秩序其他部分中應如何就合法或違法加以認定及處理之決定，於此仍懸而未決（參見 Lenckner in: Schönke/Schröder, Strafgesetzbuch, 24. Aufl., 1991, Vorbem. zu §§ 13 ff. Rdnr. 18; Eser/Burkhardt, Strafrecht I, 4. Aufl., 1992, Nr. 9, Rdnr. 41）。在法秩序其他領域，得制定獨立規範，以終止妊娠之違法性為前提加以規制。然而，若未如此為之，構成要件之排除將產生如同正當化事由的效果，從而不再滿足保護義務之最低要求。
\end{transpara}

\begin{sourcepara}[title={原文D.92}]
\sourcern{206}Während der Tatbestandsausschluß so die Möglichkeit offenhält, diesen Mindestanforderungen in anderen Teilen der Rechtsordnung zu entsprechen, gibt ein in das Strafgesetz eingeführter Rechtfertigungsgrund das von der Verfassung geforderte grundsätzliche Verbot des Schwangerschaftsabbruchs von vornherein in einem größeren Umfang preis. Insoweit sind dem Gestaltungsspielraum des Gesetzgebers Grenzen gezogen.
\end{sourcepara}
\begin{transpara}
\transrn{206}構成要件排除因而保持在法秩序其他部分符合此等最低要求之可能；相對地，導入刑法典之正當化事由，則從一開始即在較大範圍內放棄憲法所要求之終止妊娠原則禁止。就此而言，立法者之形成空間受到限制。
\end{transpara}

\begin{sourcepara}[title={原文D.93}]
\sourcern{207}b) Es entspricht unverzichtbaren rechtsstaatlichen Grundsätzen, daß einem Ausnahmetatbestand rechtfertigende Wirkung nur dann zukommen kann, wenn das Vorliegen seiner Voraussetzungen festgestellt werden muß, sei es durch die Gerichte, sei es durch Dritte, denen der Staat kraft ihrer besonderen Pflichtenstellung vertrauen darf und deren Entscheidung nicht jeder staatlichen Überprüfung entzogen ist. Läßt das vom Gesetzgeber gewählte Schutzkonzept der Beratung, soweit es um die überwiegend geltend gemachten allgemeinen Notlagen geht, diese Indikationsregelung nicht zu, weil die Feststellung ihrer Voraussetzungen die Wirksamkeit der Beratung hindern würde, so muß der Gesetzgeber insoweit darauf verzichten, den Abbruch der Schwangerschaft für gerechtfertigt zu erklären.
\end{sourcepara}
\begin{transpara}
\transrn{207}b) 一項例外構成要件只有在其要件之存在必須被確認時，始能具有正當化效果；此一確認無論由法院為之，或由國家基於其特殊義務地位而得信賴且其決定並未完全脫離任何國家審查之第三人為之，均符合不可或缺之法治國原則。若立法者所選擇之諮詢保護構想，在主要被主張之一般困境範圍內，不容許此種指標規則，因為其要件之確認將妨礙諮詢有效性，則立法者就此必須放棄將終止妊娠宣告為正當化。
\end{transpara}

\begin{sourcepara}[title={原文D.94}]
\sourcern{208}aa) Weder der im Schutzkonzept angelegte Verzicht auf die Indikationsregelung noch die einzigartige Verbindung von Mutter und Kind während der Schwangerschaft ermöglichen es, hiervon abzugehen. Es kann nicht eingewandt werden, daß die allgemeine Notlagenindikation als Rechtfertigungsgrund anderen Maßstäben unterliegen müsse, weil ihre Voraussetzungen sich einer Feststellung entzögen. Vielmehr ist eine solche Feststellung ungeachtet der dabei bestehenden Schwierigkeiten (siehe oben II.2.) weder tatsächlich noch rechtlich unmöglich. Auch soweit sie subjektive Umstände betrifft, unterscheidet sie sich nicht von vielen anderen Feststellungen zur Persönlichkeit, Willenskraft und psychischen Befindlichkeit von Personen, die der Rechtsstaat verlangt und trifft, wenn es etwa um die Ahndung einer Verletzung fundamentaler Rechtsgüter oder ihre Bewahrung geht (vgl. etwa §§ 46 \abbrFF{}, 56 \abbrFF{}, 63 \abbrFF{}, 70 \abbrStGB{}).
\end{sourcepara}
\begin{transpara}
\transrn{208}aa) 無論是保護構想中所設定之放棄指標規則，或懷孕期間母親與子女之獨特聯繫，均不能使人偏離此點。不能反駁稱，一般困境指標事由作為正當化事由必須受其他標準支配，因為其要件無法被確認。毋寧，此種確認儘管存在前述困難（見上文 II.2.），在事實上及法律上均非不可能。即使其涉及主觀情事，其亦與許多其他關於人之人格、意志力及心理狀態之確認無異；法治國在例如涉及基本法益侵害之處罰或其維護時，即要求並作成此等確認（例如參見《刑法》第46條以下、第56條以下、第63條以下、第70條）。
\end{transpara}

\begin{sourcepara}[title={原文D.95}]
\sourcern{209}Allerdings würde es dazu in nicht wenigen Fällen erforderlich sein, höchstpersönliche Beziehungen und Verhältnisse der Frau zu beleuchten. Dem kann sie aber, wenn es darum geht, ob die Tötung des auch der Frau gegenüber geschützten Lebens des Ungeborenen erlaubt ist, keine eigenen Rechte entgegenhalten. Ist aber die Feststellung der Indikation nur tatsächlich besonders aufwendig, so ist allein dies noch kein hinreichender Grund, von ihr Abstand zu nehmen und sie etwa durch eine "Selbstindikation" der Frau zu ersetzen. Auch wenn die Beratungsregelung den Frauen Verantwortung bei der Entscheidung über das Austragen ihres Kindes zutraut und mit der Beratung ein Verfahren zur Verfügung stellt, das ihnen die notwendige normative Orientierung und Ermutigung zum Kind geben kann, wäre es gleichwohl mit der rechtsstaatlichen Ordnung des Grundgesetzes unvereinbar, wenn die an dem Konflikt existenziell beteiligten Frauen selbst mit rechtlicher Erheblichkeit feststellten, ob eine Lage gegeben ist, bei der das Austragen des Kindes unzumutbar ist und deshalb der Abbruch der Schwangerschaft auch von Verfassungs wegen erlaubt werden kann. Die Frauen würden dann in eigener Sache über Recht und Unrecht befinden. Das läßt der Rechtsstaat auch und gerade in der besonderen Situation der "Zweiheit in Einheit" nicht zu. Die Verfassung verspricht dem noch in jeder Weise von der Mutter abhängigen Ungeborenen Schutz auch dieser gegenüber. Dann aber darf die besondere Verbindung von Mutter und ungeborenem Leben, die für dieses Schutz bedeutet, aber auch Gefahr bringen kann, nicht zum Anlaß genommen werden, allein der Mutter die Entscheidung über das Vorliegen der Voraussetzungen zu überlassen, die ihr nach der Rechtsordnung die Tötung des Ungeborenen erlauben. Damit wäre gerade ein Mindestmaß an rechtlichem Schutz nicht mehr gewährt.
\end{sourcepara}
\begin{transpara}
\transrn{209}不過，為此在不少案件中將有必要照明婦女極具個人性之關係與狀況。但當問題在於殺害即使相對於婦女亦受保護之未出生生命是否被允許時，婦女不能以自身權利對抗此點。若指標事由之確認僅在事實上特別費力，則單憑此點尚非放棄確認並例如以婦女之「自我指標事由」取代之充分理由。即使諮詢規制信任婦女在是否繼續懷孕之決定上負責，並透過諮詢提供一項程序，能給予其必要之規範取向及對子女之鼓勵；若在衝突中具有生存性參與之婦女自身以具有法律重要性之方式確認，是否存在一種繼續懷孕為不可期待，並因而終止妊娠亦得依憲法被允許之狀態，仍將與《基本法》之法治國秩序不相容。如此婦女將在自身案件中判斷法與不法。法治國即使且正是在「一體中之二元」之特殊情況中，也不允許此點。憲法承諾對尚在一切方面依賴母親之未出生者，亦提供相對於母親之保護。既然如此，母親與未出生生命之特殊聯繫，一方面意味著對此生命之保護，但另一方面也可能帶來危險，即不得成為理由，將是否存在依法秩序允許母親殺害未出生者之前提的決定，單獨交給母親。如此將正是不再給予最低程度之法律保護。
\end{transpara}

\begin{sourcepara}[title={原文D.96}]
\sourcern{210}bb) Der nötige Nachweis einer Notlage kann auch nicht aufgrund von Erfahrungssätzen oder Indizien als allgemein erbracht angesehen werden, so daß auch ohne eine Feststellung im Einzelfall eine Notlagenindikation an die Beratung geknüpft werden könnte.
\end{sourcepara}
\begin{transpara}
\transrn{210}bb) 必要之困境證明，亦不能基於經驗法則或間接證據而被視為一般已經提出，以致即使未於個案中確認，亦可將困境指標事由連結於諮詢。
\end{transpara}

\begin{sourcepara}[title={原文D.97}]
\sourcern{211}Es ist nichts Ungewöhnliches, daß die Rechtsordnung -- insbesondere auch das Strafrecht -- einen Tatsachenkomplex als rechtlich erheblich ansieht, der sich aus äußeren und inneren Tatsachen zusammensetzt. Oft ist deren Feststellung nur dadurch möglich, daß Umstände festgestellt werden, die nach der Lebenserfahrung auf das Vorhandensein der festzustellenden Tatsache schließen lassen (Indizien), etwa indem mit Hilfe von Erfahrungssätzen von einem bestimmten äußeren Geschehen auf innere Tatsachen geschlossen wird. Nur so läßt sich auch die Gesamtheit der Tatsachen feststellen, die Voraussetzung einer rechtfertigenden Notlagenindikation ist.
\end{sourcepara}
\begin{transpara}
\transrn{211}法秩序，尤其刑法，將由外在與內在事實組成之事實複合體視為具有法律重要性，並非異常。此等事實之確認，往往只能透過確認依生活經驗可推論待確認事實存在之情事（間接證據）而可能，例如藉由經驗法則從特定外在事件推論內在事實。只有如此，亦才能確認構成正當化困境指標事由前提之整體事實。
\end{transpara}

\begin{sourcepara}[title={原文D.98}]
\sourcern{212}Die Ausgestaltung der Beratungsregelung ermöglicht es nicht, schon aus ihrer Durchführung unter Zuhilfenahme von Erfahrungssätzen mit hinreichender Sicherheit zu schließen, dem unter den Voraussetzungen der Beratungsregelung geäußerten Abbruchverlangen liege in der Regel tatsächlich eine Konfliktsituation solcher Schwere zugrunde, daß das Austragen des Kindes unzumutbar sei und deshalb der Abbruch der Schwangerschaft auch von Verfassungs wegen erlaubt sein dürfe. Das Beratungskonzept setzt auf die das bedrohte Rechtsgut schützende Wirkung einer Beratung, die bei den Beteiligten die Bereitschaft für das Kind bekräftigen und fördern soll. Es geht auch vertretbar davon aus, daß Frauen sich in der Regel nicht leichten Herzens, ohne eine Konfliktsituation zu empfinden, zu dem Eingriff eines Schwangerschaftsabbruchs entschließen. Das Beratungskonzept mag auch geeignet sein, Frauen, die in der Beratung zum Austragen des Kindes ermutigt werden und denen die von der Rechtsordnung gezogene Grenze zwischen verbotenen und ausnahmsweise erlaubten Schwangerschaftsabbrüchen in einer ihnen verständlichen Weise erläutert worden ist, dazu zu veranlassen, sich um gewissenhafte Beurteilung und Einschätzung ihres Konflikts zu bemühen und in voller Kenntnis des Für und Wider verantwortlich eine Entscheidung zu treffen. Diese Einschätzung ist verfassungsrechtlich nicht zu beanstanden. Sie wäre allerdings nicht mehr nachvollziehbar, nähme der Gesetzgeber darüber hinaus einen Erfahrungssatz des Inhalts an, die Beratung und das ärztliche Gespräch könnten eine Frau in ihrer durch die ungewollte Schwangerschaft hervorgerufenen Konfliktsituation regelmäßig so beeinflussen, daß sie sich dem für sie physisch und psychisch belastenden Schwangerschaftsabbruch nur unterziehen werde, wenn die Voraussetzungen einer unzumutbaren Notlage tatsächlich vorliegen. Eine Frau, die das Kind -- aus welchen Gründen auch immer -- auf keinen Fall haben will, sieht als einzigen Ausweg die Abtreibung. Diese erscheint dann -- wird sie von einem Arzt in den ersten zwölf Wochen der Schwangerschaft ausgeführt -- auch bei Berücksichtigung der nicht zu vernachlässigenden Gefahr psychischer Spätfolgen für die Frau unvergleichlich weniger belastend als das Austragen, Gebären und Aufziehen des Kindes. Es entspricht daher der Erfahrung, daß ungewollt schwanger gewordene Frauen oft ihre Zuflucht in der Abtreibung suchen; das wird durch die sich nach wie vor in überaus großer Zahl ereignenden Abtreibungen bestätigt. Es besteht hingegen kein Erfahrungssatz, daß sie den Abbruch der Schwangerschaft nur dann auf sich nehmen, wenn ihnen in einer Ausnahmelage das Austragen des Kindes aus schwerwiegenden Gründen unzumutbar ist.
\end{sourcepara}
\begin{transpara}
\transrn{212}諮詢規制之形塑，並不能使人僅由其實施並輔以經驗法則，即足夠確定地推論：在諮詢規制前提下所表達之終止妊娠要求，通常實際上係基於如此嚴重之衝突狀態，以致繼續懷孕為不可期待，並因此終止妊娠亦得依憲法被允許。諮詢構想依靠諮詢對受威脅法益之保護效果；此種諮詢應確認並促進參與者對子女之意願。其亦可支持地以婦女通常並非輕率地、未感受衝突狀態即決定接受終止妊娠介入為出發點。諮詢構想也可能適於促使婦女努力良心謹慎地判斷並評估其衝突，並在充分了解利弊下負責任地作成決定；這是指那些在諮詢中被鼓勵繼續懷孕，且法秩序所劃定之被禁止與例外允許之終止妊娠界限已以其可理解方式向其說明之婦女。此一評估在憲法上無可指摘。然而，若立法者進一步採取下列內容之經驗法則，則該評估即不再可理解：諮詢與醫師談話通常能在婦女因非意願懷孕所生之衝突狀態中如此影響她，以致她只有在不可期待困境之前提實際存在時，才會接受對其身體與心理造成負擔之終止妊娠。無論基於何種理由，一名絕不願意要子女之婦女，會將終止妊娠視為唯一出路。若此終止妊娠由醫師於妊娠最初十二週內實施，則即使考量婦女心理晚期後果之不可忽視危險，其相較於懷胎、生產及養育子女，仍顯得負擔不可相比地較低。因此，非意願懷孕之婦女往往求助於終止妊娠，這符合經驗；仍然發生之極大數量終止妊娠亦確認此點。相反地，並不存在一項經驗法則，認為她們只有在例外狀態中因重大理由而不可期待其繼續懷孕時，才會承受終止妊娠。
\end{transpara}

\begin{sourcepara}[title={原文D.99}]
\sourcern{213}Ginge der Gesetzgeber aber davon aus, daß ein Abbruchverlangen nach Beratung und ärztlichem Gespräch sich jedenfalls bei typisierender Betrachtungsweise als verantwortungsvolle und von rechtlichen Maßstäben gesteuerte Entscheidung darstelle und damit die Feststellung einer Indikation durch Dritte im Einzelfall ersetzen könne, so würde er den Mindestanforderungen an den jedem Ungeborenen geschuldeten rechtlichen Schutz nicht genügen. Dazu gehört es -- wie dargelegt --, daß die Tötung des einzelnen Lebens nur erlaubt sein darf, wenn im konkreten Fall für die Frau ausnahmsweise die Voraussetzung der Unzumutbarkeit gegeben ist. Dem wird nicht mit einer Feststellung genügt, es liege eine Situation vor, in der nach Einschätzung des Gesetzgebers Frauen typischerweise nur bei Vorliegen einer rechtfertigenden Ausnahmelage abtreiben.
\end{sourcepara}
\begin{transpara}
\transrn{213}但若立法者以為，經諮詢及醫師談話後之終止妊娠要求，至少在類型化觀察下，呈現為負責任且受法律標準引導之決定，並因而得於個案中取代第三人對指標事由之確認，則其將未滿足對每一未出生者所負法律保護之最低要求。如前所述，其中包括：個別生命之殺害，只有在具體案件中對婦女例外存在不可期待性之前提時，始得被允許。僅確認存在一種依立法者評估婦女典型上只有在正當化例外狀態存在時才會終止妊娠之情境，並不足以滿足此點。
\end{transpara}

\begin{sourcepara}[title={原文D.100}]
\sourcern{214}cc) Die Schutzwirkungen der Beratungsregelung hängen nicht davon ab, daß das nach Beratung der Frau geäußerte Abbruchverlangen den Schwangerschaftsabbruch rechtfertigt. Es gibt keine Anhaltspunkte dafür, daß die Wirksamkeit der Beratung die Gewißheit der Frau voraussetzte, der Abbruch, für den sie sich nach Abschluß des Beratungsverfahrens entscheidet, werde von der Rechtsordnung ausdrücklich gebilligt. Vielmehr würde es einer auf den Schutz des ungeborenen Lebens zielenden Beratung eher abträglich sein, wenn die Rechtsordnung ohnehin jedes Abbruchverlangen nach Beratung erlaubte. Die Beratung könnte nämlich die Verantwortung der Frauen nicht stärken, wenn ihre Entschließung für den Abbruch nach der Beratung in jedem Falle rechtliche Anerkennung erführe.
\end{sourcepara}
\begin{transpara}
\transrn{214}cc) 諮詢規制之保護效果，並不取決於婦女經諮詢後所表達之終止妊娠要求能否正當化終止妊娠。沒有根據顯示，諮詢之有效性以婦女確知其於諮詢程序結束後所決定之終止妊娠將受法秩序明示認可為前提。毋寧，若法秩序反正允許每一經諮詢後之終止妊娠要求，對以保護未出生生命為目標之諮詢反而較為有害。因為若婦女在諮詢後作成終止妊娠之決意在任何情形下均獲法律承認，諮詢即無法強化婦女之責任。
\end{transpara}

\begin{sourcepara}[title={原文D.101}]
\sourcern{215}Die Bewertung eines Schwangerschaftsabbruchs als rechtmäßig auch in Fällen, in denen eine unzumutbare Ausnahmelage nicht festgestellt wird, schwächte zudem den rechtlichen Schutz des ungeborenen menschlichen Lebens, den das prinzipielle Verbot des Schwangerschaftsabbruchs durch die Aufrechterhaltung des Rechtsbewußtseins zu bewirken vermag (positive Generalprävention). Das Rechtsbewußtsein wird durch widersprüchliche rechtliche Bewertungen verunsichert. Ein solcher Widerspruch läge vor, wenn der schwangeren Frau in der Beratung eine rechtliche Orientierung gegeben würde, ihr Schwangerschaftsabbruch sei nur erlaubt, wenn Indikationen vorliegen, andererseits aber ihre Entscheidung für den Abbruch nach Beratung als gerechtfertigt, mithin als erlaubt, angesehen würde, obwohl eine Indikation nicht festgestellt wird.
\end{sourcepara}
\begin{transpara}
\transrn{215}此外，在未確認不可期待例外狀態之案件中亦將終止妊娠評價為合法，將削弱終止妊娠原則禁止透過維持法律意識所能產生之未出生人類生命法律保護（積極一般預防）。法律意識會因矛盾之法律評價而不安定。若在諮詢中給予懷孕婦女一項法律取向，即其終止妊娠只有在存在指標事由時始被允許；另一方面，雖未確認指標事由，卻又將其經諮詢後終止妊娠之決定視為正當化，因而視為允許，即存在此種矛盾。
\end{transpara}

\begin{sourcepara}[title={原文D.102}]
\sourcern{216}Mithin kann der Gesetzgeber im Rahmen einer Beratungsregelung das gewünschte Ergebnis, die Frau nicht mit Strafe zu bedrohen, wenn sie -- nach stattgehabter Beratung -- die Schwangerschaft in der Frühphase durch einen Arzt abbrechen läßt, verfassungsrechtlich unbedenklich nur durch einen Tatbestandsausschluß erreichen. Er muß dann freilich dafür Sorge tragen, daß das grundsätzliche Verbot des Schwangerschaftsabbruchs für die gesamte Dauer der Schwangerschaft, das danach im Umfang des Tatbestandsausschlusses nicht mehr in einer Strafbestimmung enthalten ist, an anderer Stelle der Rechtsordnung in geeigneter Weise zum Ausdruck kommt (vgl. auch oben D.III.1.c).
\end{sourcepara}
\begin{transpara}
\transrn{216}因此，立法者在諮詢規制框架內，若欲達成下列結果：婦女在經過諮詢後，於懷孕早期由醫師終止妊娠時不受刑罰威嚇，則在憲法上無疑義地只能透過構成要件排除達成。不過，立法者隨即必須確保，終止妊娠在整個懷孕期間之原則禁止，既然在構成要件排除範圍內不再包含於刑罰規定中，即應在法秩序其他位置以適當方式表達出來（亦參見上文 D.III.1.c）。
\end{transpara}

\begin{sourcepara}[title={原文D.103}]
\sourcern{217}3. Die Herausnahme des Schwangerschaftsabbruchs aus dem Straftatbestand läßt -- wie dargelegt -- Raum, das grundsätzliche Verbot eines Schwangerschaftsabbruchs, für den rechtfertigende Ausnahmetatbestände nicht festgestellt sind, in den übrigen Bereichen der Rechtsordnung zur Geltung zu bringen. Hierbei verlangen die Besonderheiten des Beratungskonzepts, auch für den Fall eines späteren Schwangerschaftsabbruchs Bedingungen zu schaffen, die nicht im Vorwege der Bereitschaft der Frau entgegenwirken, sich auf die dem Lebensschutz dienende Beratung einzulassen, ihren Konflikt offenzulegen und an seiner Lösung verantwortlich mitzuwirken. Deshalb ist die Rechtslage insgesamt so zu gestalten, daß es sich für die Frau nicht nahelegt, die Beratung gar nicht erst anzunehmen und in die Illegalität auszuweichen. Über die Herausnahme des Schwangerschaftsabbruchs aus der Strafdrohung hinaus muß sichergestellt sein, daß gegen das Handeln der Frau und des Arztes von Dritten Nothilfe zugunsten des Ungeborenen nicht geleistet werden kann. Die Frau muß auch in der Lage sein, den Abbruch durch einen Arzt im Rahmen eines wirksamen privatrechtlichen Vertrags durchführen zu lassen (vgl. dazu unten V.6.). Ebenso muß sie davor geschützt sein, den Abbruch und ihre Gründe hierfür unter Beeinträchtigung ihres Persönlichkeitsrechts weiteren Personen offenbaren zu müssen (vgl. unten E.V.3.b) und 4.b). Um derartige Bedingungen herbeizuführen, muß es möglich sein, in dem jeweils einschlägigen Rechtsbereich davon abzusehen, den nach Beratung vorgenommenen Schwangerschaftsabbruch, obwohl er nicht gerechtfertigt ist, als Unrecht zu behandeln.
\end{sourcepara}
\begin{transpara}
\transrn{217}3. 如前所述，將終止妊娠自犯罪構成要件中取出，留下空間，使未確認正當化例外構成要件之終止妊娠原則禁止，在法秩序其餘領域中發生效力。在此，諮詢構想之特殊性要求，即使對於後續終止妊娠之情形，也必須創設不會事先妨礙婦女願意投入服務於生命保護之諮詢、揭示其衝突並負責任地參與解決之條件。因此，整體法律狀態必須被形塑為，對婦女而言不致接近於一開始即不接受諮詢而逃入非法。除將終止妊娠排除於刑罰威嚇外，亦必須確保第三人不能針對婦女與醫師之行為，為未出生者利益提供緊急救助。婦女也必須能夠在有效私法契約框架內，由醫師實施終止妊娠（就此參見下文 V.6.）。同樣，她必須受到保護，避免在損害其人格權下，不得不向其他人揭露終止妊娠及其理由（參見下文 E.V.3.b) 及 4.b)）。為造成此等條件，在各該相關法律領域中，必須可能放棄將經諮詢後實施之終止妊娠，雖然其未被正當化，作為不法處理。
\end{transpara}

\begin{sourcepara}[title={原文D.104}]
\sourcern{218}Die Schutzpflicht für das ungeborene menschliche Leben läßt dies zu. Der Gesetzgeber ist nicht gehalten, die sich nach dem grundsätzlichen Verbot des Schwangerschaftsabbruchs an sich aufdrängenden Folgerungen in jeder Hinsicht zu ziehen, wenn ein hinreichend auf den Lebensschutz ausgerichtetes Schutzkonzept zu seiner Wirksamkeit bestimmte Ausnahmen fordert. Die Schutzwirkungen, die von dem grundsätzlichen Verbot eines Schwangerschaftsabbruchs ausgehen, indem dieses das allgemeine Rechtsbewußtsein prägt und stützt, gehen dann nicht verloren, wenn Folgewirkungen des Verbots -- mit Rücksicht auf sinnvolle andere Schutzmaßnahmen -- nur in bestimmten Bereichen der Rechtsordnung eingeschränkt werden, in anderen hingegen Geltung haben.
\end{sourcepara}
\begin{transpara}
\transrn{218}對未出生人類生命之保護義務容許此點。若一項足夠以生命保護為取向之保護構想，為其有效性而要求特定例外，立法者並無義務在所有方面均作出終止妊娠原則禁止本身所迫切導出之結論。終止妊娠原則禁止透過形塑並支撐一般法律意識所產生之保護效果，並不會因禁止之後續效果鑑於有意義之其他保護措施而僅在法秩序特定領域受限制、在其他領域仍有效力，即告喪失。
\end{transpara}

\begin{sourcepara}[title={原文D.105}]
\sourcern{219}Diesen Anforderungen kann trotz der Einschränkung der Verbotsfolgen entsprochen werden. Die im Rahmen der Beratungsregelung erfolgenden Abbrüche dürfen in ihren Rechtsfolgen gerechtfertigten Abbrüchen nur insoweit gleichgestellt werden, als dies zur Verwirklichung der Schutzmaßnahme unerläßlich ist. Solche Einschränkungen der Verbotsfolgen sind auch nur in einzelnen Rechtsbeziehungen erforderlich. Im übrigen dürfen Rechtsfolgen, die eine Rechtmäßigkeit des Handelns voraussetzen, nicht an solche Abbrüche geknüpft werden. Die verfassungsrechtlich gebotene Wertung dieser Abbrüche als nicht rechtmäßig prägt des weiteren Maßstäbe des Vertragsrechts. Sie hat etwa auch Richtlinie zu sein für die Ausbildung von Ärzten, ärztlichem Hilfspersonal und Sozialarbeitern. Allgemein ist das aus der Schutzpflicht folgende verfassungsrechtliche Gebot, den Schwangerschaftsabbruch grundsätzlich als Unrecht zu behandeln, bei der Auslegung und Anwendung des Rechts, insbesondere von Generalklauseln, in der Weise zu beachten, daß gerechtfertigte und nicht gerechtfertigte Schwangerschaftsabbrüche nicht gleich behandelt werden dürfen, es sei denn, das Schutzkonzept erfordere es.
\end{sourcepara}
\begin{transpara}
\transrn{219}儘管限制禁止效果，仍可符合此等要求。在諮詢規制框架內所為之終止妊娠，其法律效果僅得在為實現保護措施所不可或缺之範圍內，與正當化終止妊娠等同。此等禁止效果之限制，也僅在個別法律關係中為必要。除此之外，以行為合法性為前提之法律效果，不得連結於此等終止妊娠。將此等終止妊娠評價為非合法，乃憲法所要求；此一評價進一步形塑契約法之標準。其例如亦應成為醫師、醫療輔助人員及社工訓練之準則。一般而言，由保護義務所生、原則上將終止妊娠作為不法處理之憲法命令，在法律解釋及適用，尤其概括條款之解釋及適用時，應以下列方式被遵守：正當化與未正當化之終止妊娠不得被等同處理，除非保護構想有所要求。
\end{transpara}

\begin{sourcepara}[title={原文D.106}]
\sourcern{220}4. Die Beratungsregelung hat nach allem zur Folge, daß die Frau, die ihre Schwangerschaft nach Beratung abbricht, eine von der Rechtsordnung nicht erlaubte Handlung vornimmt. Das Beratungskonzept kann der Frau den Rechtfertigungsgrund der allgemeinen Notlage und deren Feststellung nicht zuteil werden lassen, ohne damit der ihr zugrundeliegenden Schutzkonzeption zu widersprechen. Dies kann die Beratungsregelung der Frau zumuten, ohne sie zum bloßen Objekt eines Schutzkonzepts herabzuwürdigen. Das Beratungskonzept nimmt die Frau als handelnde Person ernst, indem es sie als Verbündete bei dem Schutz des ungeborenen Lebens zu gewinnen sucht und dabei von ihr eine verantwortliche Mitwirkung erwartet. Es schafft des weiteren Bedingungen, die die Rechtsposition der Frau achten (siehe oben 3.) und jene Rechtsnachteile vermeiden, die der Frau Veranlassung geben könnten, sich dem Beratungsverfahren und dem ärztlichen Gespräch zu entziehen. Nur wenn es nicht mehr um solche Bedingungen geht, muß das grundsätzliche Verbot des Schwangerschaftsabbruchs zum Tragen kommen. Dies bedeutet etwa, daß für Schwangerschaftsabbrüche im Rahmen der Beratungsregelung nicht alle die rechtlichen Vorteile gewährt werden können, die nach der Rechtsordnung für rechtmäßige Abbrüche zulässig sind.
\end{sourcepara}
\begin{transpara}
\transrn{220}4. 綜上，諮詢規制之結果是：經諮詢後終止妊娠之婦女，實施一項法秩序所不允許之行為。諮詢構想不能讓婦女享有一般困境之正當化事由及其確認，否則即與其所依據之保護構想相矛盾。諮詢規制可以使婦女承受此點，而不致將其貶抑為保護構想之單純客體。諮詢構想認真看待婦女作為行動之人，因其試圖爭取婦女成為保護未出生生命之同盟者，並在此期待其負責任地參與。此外，該構想創設尊重婦女法律地位之條件（見上文 3.），並避免可能促使婦女逃避諮詢程序與醫師談話之法律不利益。只有當所涉不再是此等條件時，終止妊娠之原則禁止才必須發揮作用。這例如意味著，對諮詢規制框架內之終止妊娠，不能給予所有依法秩序對合法終止妊娠所容許之法律利益。
\end{transpara}

\bisubsubsection{IV.}{IV.}

\begin{sourcepara}[title={原文D.107}]
\sourcern{221}Die Schutzpflicht für das ungeborene menschliche Leben legt dem Gesetzgeber, wenn er sich für ein Beratungskonzept entscheidet, auch Bindungen bei der normativen Ausgestaltung des Beratungsverfahrens auf (siehe oben III.1.a). Dieses erhält mit der Verlagerung des Schwerpunktes der Schutzgewährung auf präventiven Schutz durch Beratung eine zentrale Bedeutung für den Lebensschutz. Der Gesetzgeber muß daher bei der Festlegung des Inhalts einer Beratung (1.), der Regelung ihrer Durchführung (2.) und der Organisation der Beratung einschließlich der Auswahl der an ihr mitwirkenden Personen (3.) unter Bindung an das Untermaßverbot Regelungen treffen, die wirksam und ausreichend sind, um eine Frau, die den Schwangerschaftsabbruch erwägt, für das Austragen des Kindes gewinnen zu können. Nur dann ist die Einschätzung des Gesetzgebers, mit einer Beratung könne wirksamer Lebensschutz erzielt werden, vertretbar.
\end{sourcepara}
\begin{transpara}
\transrn{221}若立法者決定採取諮詢構想，對未出生人類生命之保護義務亦使其在諮詢程序之規範形塑上受到拘束（見上文 III.1.a）。隨著保護提供之重心轉移至透過諮詢之預防性保護，此一程序對生命保護取得核心意義。因此，立法者在確定諮詢內容（1.）、規制其實施（2.）及組織諮詢包括選擇參與諮詢之人員（3.）時，受不足禁止拘束，必須作成有效且足夠之規制，以能夠爭取考慮終止妊娠之婦女繼續懷孕。唯有如此，立法者關於透過諮詢能達成有效生命保護之評估，才是可支持的。
\end{transpara}

\begin{sourcepara}[title={原文D.108}]
\sourcern{222}1. Für die Festlegung des Inhalts der Beratung kann der Gesetzgeber davon ausgehen, daß Beratung nur dann eine Chance hat, das ungeborene menschliche Leben wirklich zu schützen, wenn sie ergebnisoffen geführt wird. Die Beratung muß, um erfolgreich sein zu können, darauf angelegt sein, daß die Frau sich an der Suche nach einer Lösung beteiligt. Dies rechtfertigt es auch, davon abzusehen, die erwartete Gesprächs- und Mitwirkungsbereitschaft der Frau zu erzwingen oder sie zu verpflichten, sich im Beratungsgespräch als Person zu identifizieren.
\end{sourcepara}
\begin{transpara}
\transrn{222}1. 關於諮詢內容之確定，立法者得以為，諮詢只有在以結果開放方式進行時，才有機會真正保護未出生人類生命。諮詢若要成功，必須被設計為使婦女參與尋求解決方案。這也正當化放棄強制婦女具備所期待之談話與合作意願，或強制其在諮詢談話中表明個人身分。
\end{transpara}

\begin{sourcepara}[title={原文D.109}]
\sourcern{223}Die Beratung im Schwangerschaftskonflikt bedarf der Zielorientierung auf den Schutz des ungeborenen Lebens hin. Eine bloß informierende Beratung, die den konkreten Schwangerschaftskonflikt nicht aufnimmt und zum Thema eines persönlich geführten Gesprächs zu machen sucht, sich auch nicht um konkrete Hilfen im Blick auf diesen Konflikt bemüht, ließe die Frau im Stich und verfehlte ihren Auftrag. Die Beraterinnen oder Berater müssen sich von dem Bemühen leiten lassen, die Frau zur Fortsetzung ihrer Schwangerschaft zu ermutigen und ihr Perspektiven für ein Leben mit dem Kind zu eröffnen.
\end{sourcepara}
\begin{transpara}
\transrn{223}懷孕衝突中之諮詢，需要以保護未出生生命為目標取向。若諮詢僅提供資訊，不承接具體懷孕衝突，也不試圖將其作為個人性談話之主題，並且亦不就此一衝突努力提供具體協助，則將棄婦女於不顧，並錯失其任務。諮詢人員必須受下列努力所引導：鼓勵婦女繼續其懷孕，並為其開啟與子女共同生活之前景。
\end{transpara}

\begin{sourcepara}[title={原文D.110}]
\sourcern{224}Eine solche Ermutigung steht solange nicht im Widerspruch zu den Wirkungsbedingungen einer Beratung, wie sie von der personalen Freiheit der Ratsuchenden ausgeht, ihre Verantwortung respektiert und dementsprechend als ergebnisoffene Beratung geführt wird. Als ein Vorgang personaler Kommunikation schließt eine ergebnisoffene Beratung keineswegs aus, daß vom Berater vermittelte normative Vorstellungen und Werthaltungen in sie einfließen; sie bietet vielmehr gerade die Chance dafür, daß normative Orientierungen und Prägungen, die auch bei der Ratsuchenden vorhanden sind, angesprochen werden. Ein Schwangerschaftskonflikt erwächst in aller Regel aus einem elementaren Zwiespalt zwischen der Erkenntnis, schutzbedürftiges menschliches Leben in sich zu tragen, dem Wunsch, das Kind haben zu wollen einerseits und andererseits der Sorge, der damit verbundenen Aufgabe nicht gewachsen, schweren Konflikten im persönlichen Umfeld ausgesetzt zu sein oder eigene Lebensvorstellungen zurückstellen zu müssen. Eine Beratung, die diesen Zwiespalt aufdeckt und sich darum bemüht, die dem Kinderwunsch entgegenstehenden Umstände bewältigen zu helfen und die Frau zum Austragen der Schwangerschaft zu ermutigen, ist deshalb keine von außen kommende Fremdbestimmung. Daß die Frau solcher Hilfe bedarf, belegen auch psychische Folgeprobleme, die nach einem Schwangerschaftsabbruch nicht selten auftreten.
\end{sourcepara}
\begin{transpara}
\transrn{224}只要此種鼓勵以求助者之人格自由為出發點、尊重其責任，並相應地作為結果開放之諮詢進行，即不與諮詢之作用條件相矛盾。作為人格溝通之過程，結果開放之諮詢絕不排除由諮詢人員傳達之規範觀念與價值態度流入其中；毋寧，其正提供機會，使求助者自身亦已具有之規範取向與形塑得以被觸及。懷孕衝突通常源於一種基本矛盾：一方面認知到自身承載需要保護之人類生命，並有想要子女之願望；另一方面則憂慮自己無法勝任與此相連之任務、將暴露於個人環境中之嚴重衝突，或必須推遲自身生活想像。因此，揭示此一矛盾，並努力協助克服與生育願望相對立之情事、鼓勵婦女繼續懷孕之諮詢，並非來自外部之他律。婦女需要此種協助，也由終止妊娠後並不少見之心理後續問題所證明。
\end{transpara}

\begin{sourcepara}[title={原文D.111}]
\sourcern{225}Beratung kann nach alledem nicht nur durch Manipulation und Indoktrination mißglücken, sondern ebenso durch ein unbeteiligtes Anheimstellen, das die Frau mit ihrem Konflikt allein läßt und ihr damit letztlich anteilnehmenden Rat verweigert. Auch eine Beratung, die sich lediglich an der im Beratungsgespräch vorgetragenen Interessenlage der schwangeren Frau orientiert, ohne den vorhandenen Zwiespalt aufzugreifen, wird dem Auftrag der Beratung nicht gerecht. Andererseits würde eine auf die Erzeugung von Schuldgefühlen zielende und in dieser Weise belehrende Einflußnahme die Bereitschaft der Frau, sich der Beratung zu öffnen und sich ihren Zwiespalt bewußt zu machen, behindern.
\end{sourcepara}
\begin{transpara}
\transrn{225}依此，諮詢不僅可能因操縱與灌輸而失敗，同樣也可能因不參與式地任其自行決定而失敗；後者使婦女獨自面對其衝突，並最終拒絕給予其有同理參與之建議。僅以諮詢談話中所提出之懷孕婦女利益狀態為取向，而不承接既存矛盾之諮詢，也不能符合諮詢任務。另一方面，以製造罪惡感為目標、並以此方式進行說教式影響，則會妨礙婦女願意向諮詢敞開並意識到自身矛盾。
\end{transpara}

\begin{sourcepara}[title={原文D.112}]
\sourcern{226}Die Beratung soll ermutigen, nicht einschüchtern; Verständnis wecken, nicht belehren; die Verantwortung der Frau stärken, nicht sie bevormunden. Das stellt hohe Anforderungen an die inhaltliche Ausgestaltung der Beratung und an die Personen, die sie durchführen. Dazu bedarf es entsprechender gesetzlicher Vorgaben.
\end{sourcepara}
\begin{transpara}
\transrn{226}諮詢應鼓勵，而非恐嚇；喚起理解，而非說教；強化婦女之責任，而非監護式地替其作主。這對諮詢之內容形塑及實施諮詢之人提出高度要求。為此，需要相應之法律預設。
\end{transpara}

\begin{sourcepara}[title={原文D.113}]
\sourcern{227}a) Soll die Verantwortung der schwangeren Frau für das ungeborene Leben Grundlage einer gewissenhaften Entscheidung werden, so muß die Frau sich eben dieser Verantwortung bewußt sein, die sie nach dem Beratungskonzept in spezifischer Weise trägt. Dabei muß sie wissen, daß das Ungeborene insbesondere auch ihr gegenüber ein eigenes Recht auf Leben hat, also auch im Frühstadium der Schwangerschaft nach der Rechtsordnung besonderen Schutz genießt. Mithin muß der Frau bewußt sein, daß nur in Ausnahmesituationen nach der Rechtsordnung ein Schwangerschaftsabbruch in Betracht gezogen werden darf, nämlich nur, wenn der Frau eine Belastung erwächst, die so schwer und außergewöhnlich ist, daß sie die zumutbare Opfergrenze übersteigt. Dessen muß sich die beratende Person vergewissern und etwa vorhandene Fehlvorstellungen in für die Ratsuchende verständlicher Weise korrigieren.
\end{sourcepara}
\begin{transpara}
\transrn{227}a) 若懷孕婦女對未出生生命之責任應成為良心謹慎決定之基礎，則婦女必須意識到正是此一責任，亦即其依諮詢構想以特殊方式承擔之責任。在此，她必須知道，未出生者尤其亦相對於她享有自身之生命權，因而即使在懷孕早期，依法秩序亦享有特別保護。因此，婦女必須意識到，依法秩序只有在例外情境中始得考慮終止妊娠，亦即只有當婦女所承受之負擔如此嚴重且異常，以致超過可期待犧牲界限時。諮詢人員必須確定婦女已意識到此點，並以求助者可理解之方式糾正可能存在之錯誤觀念。
\end{transpara}

\begin{sourcepara}[title={原文D.114}]
\sourcern{228}Hiergegen kann nicht eingewandt werden, daß die bloße Mitteilung einer Rechtslage, die für die Frau keine strafrechtlichen Folgen vorsieht, Schutz zu bewirken nicht in der Lage sei. Es gibt keine Anhaltspunkte dafür, daß schwangere Frauen in einer Konfliktsituation allgemein einer solchen Information unzugänglich sind oder sie in ihrer Situation gar nur als belehrende Einflußnahme des Staates empfinden und beiseite schieben. Das Beratungskonzept geht gerade davon aus, daß auch die schwangere Frau, die sich in einer Not- und Konfliktlage sieht, zu einer verantwortlichen, die Belange des ungeborenen Lebens achtenden Entscheidung in der Lage sei; dann ist es nur folgerichtig anzunehmen, daß sie bei ihrer Entscheidung mit in Erwägung ziehen wird, was die Rechtsordnung insoweit als Recht und Unrecht ansieht. Die dazu notwendigen Bewertungsmaßstäbe hat ihr -- soweit erforderlich -- die Beratung in verständlicher Weise zu vermitteln. Danach jedenfalls vermag die Frau in gleicher Weise wie jeder Handelnde im Augenblick seines Handelns ihr Tun rechtlich zu beurteilen.
\end{sourcepara}
\begin{transpara}
\transrn{228}對此不能反駁稱，單純告知一項對婦女未規定刑法後果之法律狀態，不能產生保護。沒有根據顯示，處於衝突狀態之懷孕婦女一般無法接受此種資訊，或在其處境中甚至只將其感受為國家之說教式影響而將之置於一旁。諮詢構想正是以即使是自認處於困境與衝突狀態之懷孕婦女，亦有能力作成一項負責任且尊重未出生生命利益之決定為出發點；既然如此，合乎一貫的假定即是，她在作成決定時，會一併考量法秩序就此何者視為法、何者視為不法。為此所必要之評價標準，諮詢必須在必要範圍內以可理解之方式傳達予她。依此，婦女至少能如同任何行動者在其行動時刻一般，對其作為作法律評價。
\end{transpara}

\begin{sourcepara}[title={原文D.115}]
\sourcern{229}b) Die Wissenschaft hat Methoden dafür entwickelt, wie bei der Bewältigung von Konflikten Hilfe geleistet werden kann; sie werden auch praktiziert. Ein Konzept, das den Schutz des ungeborenen menschlichen Lebens in den ersten zwölf Wochen einer Schwangerschaft vornehmlich durch Beratung gewährleisten will, kann hierauf nicht verzichten. Jede Beratung muß daher darauf angelegt sein, ein Gespräch zu führen und dabei die Methoden einer Konfliktberatung anzuwenden. Dies setzt zum einen voraus, daß die Beratenden über entsprechende Fähigkeiten verfügen und jeder einzelnen Frau hinreichend Zeit widmen können. Zum anderen ist die Aufnahme einer Konfliktberatung von vornherein nur möglich, wenn die Schwangere der beratenden Person die wesentlichen Gründe mitteilt, die sie dazu bewegen, einen Abbruch der Schwangerschaft in Erwägung zu ziehen.
\end{sourcepara}
\begin{transpara}
\transrn{229}b) 科學已發展出如何協助克服衝突之方法；此等方法亦已被實踐。一項主要透過諮詢保障妊娠最初十二週內未出生人類生命保護之構想，不能放棄此等方法。因此，每一諮詢均必須被設計為進行談話，並在其中運用衝突諮詢方法。這一方面以諮詢人員具備相應能力，並能對每一名婦女投入足夠時間為前提。另一方面，衝突諮詢從一開始只有在懷孕婦女向諮詢人員告知促使其考慮終止妊娠之主要理由時，始有可能展開。
\end{transpara}

\begin{sourcepara}[title={原文D.116}]
\sourcern{230}Wenn es auch der Charakter einer Beratung ausschließt, eine Gesprächs- und Mitwirkungsbereitschaft der schwangeren Frau zu erzwingen, ist doch für eine Konfliktberatung, die zugleich die Aufgabe des Lebensschutzes erfüllen soll, die Mitteilung der Gründe unerläßlich, die dazu führen, einen Schwangerschaftsabbruch zu erwägen.
\end{sourcepara}
\begin{transpara}
\transrn{230}即使諮詢之性格排除強制懷孕婦女具備談話與合作意願，對一項同時應履行生命保護任務之衝突諮詢而言，告知導致其考慮終止妊娠之理由仍不可或缺。
\end{transpara}

\begin{sourcepara}[title={原文D.117}]
\sourcern{231}Dies von der Frau zu verlangen, beeinträchtigt weder die Ergebnisoffenheit des Beratungsverfahrens, noch wertet sie die der Frau zukommende Verantwortung ab. Entscheidend ist in diesem Zusammenhang, daß das Konzept, Lebensschutz durch Beratung zu gewähren, darauf verzichtet, die Gründe der Frau für den Abbruch einer Nachprüfung und Bewertung durch Dritte anhand von Indikationstatbeständen zu unterziehen und die nach der Beratung getroffene Entscheidung der Frau gegen das Kind mit Sanktionen zu belegen. Wenn in einem solchen Konzept die Beratung darauf hinwirken soll, daß die Frau ihre Abbruchgründe mitteilt, knüpft dies gerade an die Verantwortung der schwangeren Frau und ihre Fähigkeit an, eine gewissenhafte Entscheidung zu treffen.
\end{sourcepara}
\begin{transpara}
\transrn{231}要求婦女為此告知，既不損害諮詢程序之結果開放性，亦不貶低婦女所享有之責任。在此脈絡中具有決定性的是，透過諮詢提供生命保護之構想，放棄依指標事由構成要件將婦女終止妊娠之理由交由第三人審查及評價，並放棄對婦女經諮詢後作成不利於子女之決定施加制裁。若在此種構想中，諮詢應促使婦女告知其終止妊娠理由，這正是連結於懷孕婦女之責任及其作成良心謹慎決定之能力。
\end{transpara}

\begin{sourcepara}[title={原文D.118}]
\sourcern{232}c) Die Erfüllung der Schutzpflicht gegenüber dem ungeborenen menschlichen Leben erfordert es weiter, daß in der Beratung soziale und sonstige Fördermaßnahmen, die der Staat ergriffen hat, der Frau vorgestellt werden und sie bei der Inanspruchnahme der Leistungen so effektiv wie möglich unterstützt wird. Gleiches gilt, soweit Dritte, wie etwa Kirchen oder Stiftungen, Mittel zum Schutze des Lebens bereithalten. Nur so ist gewährleistet, daß alle diese Hilfen gerade die Frauen, die ihrer oft am meisten bedürfen, auch erreichen und ihren Willen zum Kind fördern.
\end{sourcepara}
\begin{transpara}
\transrn{232}c) 履行對未出生人類生命之保護義務，進一步要求在諮詢中向婦女介紹國家所採取之社會及其他促進措施，並在其請領給付時盡可能有效地予以支持。若第三人，例如教會或基金會，備有保護生命之資源，亦同。唯有如此，方能確保所有此等協助正好到達那些往往最需要協助之婦女，並促進其對子女之意願。
\end{transpara}

\begin{sourcepara}[title={原文D.119}]
\sourcern{233}d) Das Beratungsverfahren löste die ihm zukommende Aufgabe nicht hinreichend, wenn nicht berücksichtigt würde, daß Personen, die der schwangeren Frau nahestehen -- insbesondere der Vater des ungeborenen Kindes und die Eltern einer minderjährigen Schwangeren --, zu ihrer Unterstützung berufen sind und sie in ihrer Entscheidung für oder gegen das Kind beeinflussen können. Jede Beratung muß es deshalb unternehmen zu klären, ob die Hinzuziehung des Vaters, naher Angehöriger oder Vertrauter angezeigt erscheint und die schwangere Frau dafür gewonnen werden kann. In Fällen, in denen die Schwangere von vornherein von einer solchen ihr nahestehenden Person begleitet wird, ist es hingegen geboten zu prüfen, ob von dieser ein für das ungeborene Leben schädlicher Einfluß zu befürchten und die Schwangere daher aufzufordern ist, noch einmal ohne Begleitung zu einem Beratungsgespräch zu kommen.
\end{sourcepara}
\begin{transpara}
\transrn{233}d) 若未考量與懷孕婦女親近之人，尤其未出生子女之父及未成年懷孕婦女之父母，負有支持她之使命，並且可能影響其作成有利或不利於子女之決定，則諮詢程序即未充分完成其任務。因此，每一諮詢均必須嘗試釐清，是否顯得有必要加入父親、近親或信任之人，且是否能爭取懷孕婦女同意。相反地，在懷孕婦女一開始即由此種與其親近之人陪同之案件中，則有必要審查是否須擔心此人對未出生生命有有害影響，並因此要求懷孕婦女再度不受陪同地前來諮詢談話。
\end{transpara}

\begin{sourcepara}[title={原文D.120}]
\sourcern{234}2. Der notwendige Inhalt einer Beratung muß auch die Regelung ihrer Durchführung bestimmen.
\end{sourcepara}
\begin{transpara}
\transrn{234}2. 諮詢之必要內容亦必須決定其實施之規制。
\end{transpara}

\begin{sourcepara}[title={原文D.121}]
\sourcern{235}a) Vorzusehen ist, daß die schwangere Frau nicht notwendig schon nach dem ersten Beratungsgespräch die Ausstellung der Beratungsbescheinigung verlangen kann. Wenn auch ein im Beratungsverfahren ausgeübter Druck der Wirksamkeit einer Beratung im Grundsatz eher abträglich ist, muß doch vorgesehen werden, daß die Beratungsstelle die Bescheinigung erst ausstellt, wenn sie die Beratung als abgeschlossen ansieht. Der besonderen psychischen Situation der Frau kann auch bei einer solchen Regelung hinreichend Rechnung getragen werden. Deshalb werden die Beratungsstellen gehalten sein, etwa erforderliche weitere Beratungstermine unverzüglich zur Verfügung zu stellen. Allerdings darf ein Vorenthalten der Beratungsbescheinigung nicht dazu dienen, die zur Abtreibung entschlossene Frau zu veranlassen, den Abbruch bis zum Ende der Zwölf-Wochen-Frist hinauszuschieben.
\end{sourcepara}
\begin{transpara}
\transrn{235}a) 應規定，懷孕婦女並非必然在第一次諮詢談話後即得要求出具諮詢證明。即使在諮詢程序中施加壓力，原則上較可能不利於諮詢之有效性，仍必須規定，諮詢機構只有在其認為諮詢已結束時始出具證明。在此種規制下，亦能充分顧及婦女之特殊心理狀態。因此，諮詢機構將有義務立即提供可能必要之其他諮詢時段。不過，拒不出具諮詢證明不得用於促使已決定終止妊娠之婦女，將終止妊娠延至十二週期限結束。
\end{transpara}

\begin{sourcepara}[title={原文D.122}]
\sourcern{236}b) Der Schwangerschaftsabbruch darf der Beratung nicht auf dem Fuße folgen. Die Beratung wird vielfach von der schwangeren Frau nicht sofort, sondern erst in einer gewissen Zeit verarbeitet werden können. Die Frau muß noch Gelegenheit haben, mit ihr vertrauten Personen über ihre Entscheidung zu sprechen, wenn diese verantwortungsbewußt getroffen werden soll (vgl. schon \abbrBVerfGE{} 39, 1 [64]).
\end{sourcepara}
\begin{transpara}
\transrn{236}b) 終止妊娠不得緊接著諮詢發生。諮詢往往無法由懷孕婦女立即消化，而是需要一定時間。若其決定應負責任地作成，婦女仍必須有機會與其信任之人談論其決定（早已參見 BVerfGE 39, 1 [64]）。
\end{transpara}

\begin{sourcepara}[title={原文D.123}]
\sourcern{237}3. Will der Staat den von ihm zu gewährenden Schutz des ungeborenen menschlichen Lebens durch ein Beratungsverfahren verwirklichen, so trägt er für dessen Durchführung die volle Verantwortung. Ihm obliegt es, ein angemessenes Beratungsangebot sicherzustellen, und er darf sich seiner Verantwortung nicht etwa dadurch entziehen, daß er die Beratung privaten Organisationen zu unkontrollierter und nach je eigenen religiösen, weltanschaulichen oder politischen Zielvorstellungen ausgerichteter Ausführung überläßt. Wird die Beratung von nichtstaatlichen Organisationen durchgeführt, so ändert dies nichts daran, daß im Rahmen eines Schutzkonzepts, das gerade die Beratung zu einem wesentlichen Instrument des Lebensschutzes macht, die Beratung eine Aufgabe des Staates bleibt; er darf deren Handhabung um der Wirksamkeit des Schutzes willen seiner Aufsicht nicht entgleiten lassen. Daraus folgt:
\end{sourcepara}
\begin{transpara}
\transrn{237}3. 若國家欲透過諮詢程序實現其所應提供之未出生人類生命保護，則國家對該程序之實施負完全責任。國家負有確保適當諮詢提供之義務，且不得藉由將諮詢交由私人組織不受控制地、並依其各自宗教、世界觀或政治目標想像而執行，來逃避其責任。若諮詢由非國家組織實施，這並不改變下列事實：在一項正是將諮詢作為生命保護重要工具之保護構想框架內，諮詢仍是國家任務；為保護之有效性，國家不得使其操作脫離自身監督。由此可得：
\end{transpara}

\begin{sourcepara}[title={原文D.124}]
\sourcern{238}a) Der Staat darf, unabhängig davon, ob es sich um staatliche, kommunale oder private Träger handelt, nur solchen Einrichtungen die Beratung anvertrauen, die nach ihrer Organisation, nach ihrer Grundeinstellung zum Schutz des ungeborenen Lebens, wie sie in ihren verbindlichen Handlungsmaßstäben und öffentlichen Verlautbarungen zum Ausdruck kommt, sowie durch das bei ihnen tätige Personal die Gewähr dafür bieten, daß die Beratung im Sinne der verfassungsrechtlichen und gesetzlichen Vorgaben erfolgt.
\end{sourcepara}
\begin{transpara}
\transrn{238}a) 無論涉及國家、市鎮或私人承辦者，國家均僅得將諮詢委託予下列機構：依其組織、依其對未出生生命保護之基本態度，亦即如其在具拘束力之行動標準及公開聲明中所表達者，以及依其所任用人員，能保證諮詢依憲法及法律預設之意旨進行者。
\end{transpara}

\begin{sourcepara}[title={原文D.125}]
\sourcern{239}b) Eine Beratung, die den dargestellten inhaltlichen Anforderungen genügen soll, muß von persönlich und fachlich qualifiziertem Personal durchgeführt werden, dessen Zahl so bemessen sein muß, daß ein Beratungsgespräch nicht unter Zeitdruck steht. In Fällen, in denen es um die Beratung in einer besonders schwierigen Lage geht, muß es möglich sein, auf Sachverständige zurückzugreifen, die über die jeweils erforderlichen Fachkenntnisse verfügen; so etwa wenn eine embryopathische Indikation in Frage steht, auf den bei den Fachärzten und in den Behindertenverbänden versammelten Sachverstand und deren besondere Erfahrungen.
\end{sourcepara}
\begin{transpara}
\transrn{239}b) 欲符合前述內容要求之諮詢，必須由在人格及專業上合格之人員實施；其人數必須如此衡量，使諮詢談話不處於時間壓力下。在涉及特別困難處境之諮詢案件中，必須可能求助於具備各該必要專業知識之專家；例如當胚胎病理指標事由成為問題時，得求助於專科醫師及身心障礙者團體所匯集之專業知識及其特殊經驗。
\end{transpara}

\begin{sourcepara}[title={原文D.126}]
\sourcern{240}c) Wer den Schwangerschaftsabbruch vornimmt, muß als Berater ausgeschlossen sein. Darüber hinaus ist eine wie immer geartete organisatorische, institutionelle oder wirtschaftliche Integration von Beratungsstellen in Einrichtungen zur Vornahme von Schwangerschaftsabbrüchen unzulässig, wenn hiernach nicht auszuschließen ist, daß die Beratungseinrichtung materiell an der Durchführung von Schwangerschaftsabbrüchen interessiert ist. Das würde sie für eine Beratung, die die Aufgabe hat, auf den Schutz des ungeborenen Lebens hinzuwirken, die schwangere Frau zum Austragen des Kindes zu ermutigen und ihre Notlage einfühlsam mit ihr zu erörtern, untauglich machen.
\end{sourcepara}
\begin{transpara}
\transrn{240}c) 實施終止妊娠者，必須被排除作為諮詢人員。此外，若依此不能排除諮詢機構在物質上對實施終止妊娠具有利益，則諮詢機構與終止妊娠實施機構之任何形式的組織、制度或經濟整合均不被允許。這將使其不適於進行下列諮詢：其任務在於促進未出生生命之保護、鼓勵懷孕婦女繼續懷孕，並以同理心與其討論其困境。
\end{transpara}

\begin{sourcepara}[title={原文D.127}]
\sourcern{241}d) Der Staat genügt seiner Verantwortung für die Organisation und Durchführung der Beratung nicht schon dadurch, daß er die Eignung der Beratungsinstitutionen in dem oben aufgezeigten Sinn bei ihrer Anerkennung prüft, dann aber für die Zukunft eine Rechtsposition verleiht, die nur widerrufen werden kann, wenn eine wesentliche und nachhaltige Verletzung der Pflicht zur Beratung im Sinne der verfassungsrechtlichen Vorgaben und gesetzlichen Vorschriften nachzuweisen ist.
\end{sourcepara}
\begin{transpara}
\transrn{241}d) 國家並非僅因其在承認諮詢機構時，依上述意義審查其適格性，之後卻賦予一項只有在能證明其重大且持續違反依憲法預設及法律規定而為諮詢之義務時始得撤銷之未來法律地位，即已滿足其對諮詢組織及實施所負之責任。
\end{transpara}

\begin{sourcepara}[title={原文D.128}]
\sourcern{242}Die Beratungsstellen trifft im Rahmen des Beratungskonzepts eine besondere Verantwortlichkeit. Der Staat muß daher -- auf gesetzlicher Grundlage -- die Anerkennung dieser Stellen regelmäßig und in nicht zu langen Zeitabständen überprüfen und sich dabei vergewissern, ob die Anforderungen an die Beratung beachtet werden; nur unter diesen Voraussetzungen darf die Anerkennung weiter fortbestehen oder neu bestätigt werden.
\end{sourcepara}
\begin{transpara}
\transrn{242}在諮詢構想框架內，諮詢機構負有特殊責任。因此，國家必須在法律基礎上，定期且以不過長之間隔審查此等機構之承認，並在此確認諮詢要求是否受到遵守；只有在此等前提下，承認始得繼續存在或重新確認。
\end{transpara}

\begin{sourcepara}[title={原文D.129}]
\sourcern{243}e) Diese Überwachungspflicht des Staates setzt voraus, daß das Gesetz auch Möglichkeiten zur wirksamen Überwachung schafft.Es ist Sache des Gesetzgebers zu entscheiden, wie dies zu geschehen hat. Naheliegend ist eine Verpflichtung der beratenden Personen, nach Abschluß einer Beratung deren wesentlichen Inhalt und angebotene Hilfsmaßnahmen in einem Protokoll festzuhalten, das dann allerdings ausschließlich zur späteren Kontrolle der Beratungstätigkeit Verwendung finden darf, nicht hingegen zur Überprüfung und Bewertung einzelner Abbrüche. Bei Beratungsstellen, in denen eine Supervision der Beratungstätigkeit stattfindet -- was grundsätzlich naheliegt --, dürften solche Protokolle ohnehin gebräuchlich sein. Dem Anspruch der an der Beratung Beteiligten auf Wahrung ihres Persönlichkeitsrechts wäre dadurch Rechnung zu tragen, daß das Protokoll Rückschlüsse auf die Identität der Beratenen und eventuell hinzugezogener Dritter nicht erlaubt. Damit würde auch der Gefahr begegnet, durch eine Dokumentation das für eine wirksame Beratung unerläßliche Vertrauensverhältnis der an ihr Beteiligten zu beeinträchtigen.
\end{sourcepara}
\begin{transpara}
\transrn{243}e) 國家之此一監督義務，以法律亦創設有效監督之可能性為前提。如何為之，是立法者之事。較為自然的是課予諮詢人員義務，在諮詢結束後將其重要內容及所提供之協助措施記載於紀錄中；不過，該紀錄僅得用於日後控制諮詢活動，不得用於審查及評價個別終止妊娠。對於進行諮詢活動督導之諮詢機構而言，此等紀錄原本即可能通行；而進行督導原則上是自然可行的。對參與諮詢者維護其人格權之請求，應透過使該紀錄不能推論受諮詢者及可能被加入之第三人身分而加以顧及。如此亦可因應下列危險：透過文件記載損害參與諮詢者之間對有效諮詢不可或缺之信賴關係。
\end{transpara}

\begin{sourcepara}[title={原文D.130}]
\sourcern{244}Für eine wirksame Kontrolle der Beratungsstellen ist es erforderlich, daß diese in Rechenschaftsberichten regelmäßig die in ihrer Beratung gesammelten Erfahrungen niederlegen. Dies könnte auch deshalb notwendig sein, weil solche Berichte eine Grundlage abgeben können für die Überprüfung des Schutzkonzepts und Überlegungen des Staates, es zu verbessern.
\end{sourcepara}
\begin{transpara}
\transrn{244}為有效控制諮詢機構，有必要使其定期在責任報告中記載其於諮詢中所累積之經驗。此亦可能是必要的，因為此等報告得作為審查保護構想及國家思考如何改善該構想之基礎。
\end{transpara}

\bisubsubsection{V.}{V.}

\begin{sourcepara}[title={原文D.131}]
\sourcern{245}Das Schutzkonzept einer Beratungsregelung trifft im Arzt auf einen weiteren Beteiligten, der der Frau -- nunmehr aus ärztlicher Sicht -- Rat und Hilfe schuldet. Der Arzt darf einen verlangten Schwangerschaftsabbruch nicht lediglich vollziehen, sondern hat sein ärztliches Handeln zu verantworten. Er ist Gesundheit und Lebensschutz verpflichtet und darf deshalb nicht unbesehen an einem Schwangerschaftsabbruch mitwirken.
\end{sourcepara}
\begin{transpara}
\transrn{245}諮詢規制之保護構想，在醫師身上遇到另一名參與者；醫師現在從醫療觀點負有向婦女提供建議與協助之義務。醫師不得僅僅執行所要求之終止妊娠，而必須對其醫療行為負責。其受健康與生命保護義務拘束，因此不得不加審視即參與終止妊娠。
\end{transpara}

\begin{sourcepara}[title={原文D.132}]
\sourcern{246}Die staatliche Schutzpflicht erfordert es hier, daß die im Interesse der Frau notwendige Beteiligung des Arztes zugleich Schutz für das ungeborene Leben bewirkt. Der Arzt ist schon durch Berufsethos und Berufsrecht darauf verpflichtet, sich grundsätzlich für die Erhaltung menschlichen Lebens, auch des ungeborenen, einzusetzen. Daß der Arzt dieser Schutzaufgabe bei der ärztlichen Beratung und der Entscheidung über die Mitwirkung am Schwangerschaftsabbruch nachkommen kann, muß der Staat sicherstellen. Er muß ferner gerade dann, wenn die Rechtsordnung darauf verzichtet, das Vorliegen eines Rechtfertigungsgrundes für den Schwangerschaftsabbruch im Einzelfall feststellen zu lassen, den Arzt verpflichten, die ihm zukommende Schutzaufgabe wahrzunehmen.
\end{sourcepara}
\begin{transpara}
\transrn{246}國家保護義務在此要求，醫師為婦女利益所必要之參與，同時也對未出生生命產生保護。醫師本已因職業倫理及職業法而負有原則上致力於保存人類生命，包括未出生生命之義務。國家必須確保，醫師能在醫療諮詢及是否參與終止妊娠之決定中履行此一保護任務。此外，尤其當法秩序放棄在個案中使終止妊娠正當化事由之存在受確認時，國家必須課予醫師履行其所負保護任務之義務。
\end{transpara}

\begin{sourcepara}[title={原文D.133}]
\sourcern{247}1. Im einzelnen kann die rechtliche Regelung der Pflichten des Arztes bei einem Schwangerschaftsabbruch zunächst an die allgemeinen Berufspflichten anknüpfen, die der Arzt bei jedem medizinischen Eingriff zu beachten hat. Seine allgemeinen Pflichten zur Erhebung eines Befundes, zur Aufklärung und Beratung des Patienten sowie zur Dokumentation sind allerdings den Besonderheiten des Schwangerschaftsabbruchs nach dem Beratungskonzept anzupassen.
\end{sourcepara}
\begin{transpara}
\transrn{247}1. 具體而言，關於醫師在終止妊娠時義務之法律規制，首先得連結於醫師在每一醫療介入中均須遵守之一般職業義務。不過，其一般之診察、告知及諮詢病人，以及文件記載義務，必須依諮詢構想下終止妊娠之特殊性加以調整。
\end{transpara}

\begin{sourcepara}[title={原文D.134}]
\sourcern{248}a) Eine ärztlich verantwortbare Entscheidung über die Mitwirkung beim Schwangerschaftsabbruch setzt voraus, daß der Arzt sich über die Voraussetzungen vergewissert, von denen nach dem Schutzkonzept einer Beratungsregelung der Ausschluß der Strafdrohung abhängen muß. Dazu gehört neben der Prüfung, ob die Frau sich hat beraten lassen und ob die Überlegungsfrist zwischen Beratung und Schwangerschaftsabbruch gewahrt ist, die Feststellung des Alters der Schwangerschaft. Hierbei darf der Arzt sich nicht allein auf die Angaben der Frau verlassen, sondern muß sich einer verläßlichen Untersuchungsmethode -- etwa der Ultraschalluntersuchung -- bedienen.
\end{sourcepara}
\begin{transpara}
\transrn{248}a) 醫師對於是否參與終止妊娠作成可負責之醫療決定，預設醫師確認諮詢規制之保護構想中刑罰威嚇排除所必須依附之要件。除審查婦女是否已接受諮詢，以及諮詢與終止妊娠間之思考期間是否已遵守外，尚包括確認懷孕年齡。在此，醫師不得僅依賴婦女之陳述，而必須運用可靠之檢查方法，例如超音波檢查。
\end{transpara}

\begin{sourcepara}[title={原文D.135}]
\sourcern{249}Weiterhin obliegt es dem Arzt, über die rein medizinischen Aspekte des Schwangerschaftsabbruchs hinaus den Schwangerschaftskonflikt, in dem die Frau steht, im Rahmen ärztlicher Erkenntnismöglichkeiten zu erheben. Dazu hat er sich die Gründe, aus denen die Frau den Schwangerschaftsabbruch verlangt, darlegen zu lassen. Soweit diese -- wie etwa der Gesundheitszustand der Frau -- ärztlicher Untersuchung zugänglich sind, obliegt dem Arzt grundsätzlich eine eigene Beurteilung. Bei anderen Gründen darf der Arzt regelmäßig von den Angaben der Frau ausgehen, sofern sie ihm -- gegebenenfalls nach entsprechender Vergewisserung -- glaubhaft erscheinen. Etwaige tieferliegende Ursachen des Schwangerschaftskonflikts soll er in Erfahrung zu bringen suchen. Vor allem hat er sein Augenmerk darauf zu richten, ob die Frau tatsächlich den Schwangerschaftsabbruch innerlich bejaht oder ob sie insbesondere Einflüssen unterlegen ist, die von ihrem familiären oder weiteren sozialen Umfeld -- etwa dem Ehemann, dem Partner, den Eltern oder dem Arbeitgeber -- ausgegangen sind. In solchen Fällen kann vermehrt die Gefahr nachfolgender psychischer Störungen bestehen, die der Arzt bei seiner Beratung und Entscheidung berücksichtigen und über die er die Frau in geeigneter Weise aufklären muß. Der Arzt darf weiterhin Umstände nicht außer acht lassen, die darauf hindeuten, daß der Schwangerschaftsabbruch der Frau in ihrem Konflikt nicht hilft.
\end{sourcepara}
\begin{transpara}
\transrn{249}此外，醫師負有義務，在終止妊娠之純醫學面向之外，於醫療認識可能範圍內掌握婦女所處之懷孕衝突。為此，醫師須讓婦女陳述其要求終止妊娠之理由。若此等理由如婦女健康狀態一般，可受醫學檢查，醫師原則上負有自行判斷義務。對於其他理由，醫師通常得以婦女陳述為出發點，只要該等陳述在必要時經相應確認後對其顯得可信。醫師應嘗試了解懷孕衝突可能較深層之原因。尤其應注意，婦女是否實際上內心肯認終止妊娠，或其是否特別受到來自其家庭或較廣泛社會環境，例如丈夫、伴侶、父母或雇主之影響。在此等案件中，後續心理障礙之危險可能增加；醫師在其諮詢與決定中必須考量此點，並以適當方式告知婦女。醫師亦不得忽略顯示終止妊娠無助於婦女克服其衝突之情事。
\end{transpara}

\begin{sourcepara}[title={原文D.136}]
\sourcern{250}b) Im Beratungsgespräch hat der Arzt der Frau in geeigneter Weise, ohne vorhandene Ängste und seelische Nöte zu verstärken, ein hinreichendes Wissen davon zu vermitteln und zur Sprache zu bringen, daß der Schwangerschaftsabbruch menschliches Leben zerstört. Aus dem einschlägigen Schrifttum sind Schilderungen von Frauen bekannt, daß sie vor dem Schwangerschaftsabbruch falsche Vorstellungen über das tatsächliche Geschehen gehabt hätten und bei ausreichender Kenntnis den Abbruch nicht hätten vornehmen lassen. In solchen Fällen ist der ärztlichen Aufklärungspflicht nicht genügt worden.
\end{sourcepara}
\begin{transpara}
\transrn{250}b) 在諮詢談話中，醫師必須以適當方式，不強化既有恐懼與心理困苦，向婦女傳達並提出足夠認識：終止妊娠毀滅人類生命。相關文獻中可見婦女敘述，其在終止妊娠前對實際發生之情事有錯誤想像，若具備充分知識，即不會讓終止妊娠實施。在此等案件中，醫師之告知義務並未被滿足。
\end{transpara}

\begin{sourcepara}[title={原文D.137}]
\sourcern{251}Der ärztlichen Untersuchung und Information müssen andererseits aufgrund der Schutzpflicht für das ungeborene Leben Grenzen gezogen werden, um der Gefahr von Schwangerschaftsabbrüchen aus Gründen der -- von der Verfassungsordnung mißbilligten -- Geschlechtswahl zu begegnen. Deshalb muß ausgeschlossen sein, daß in der Frühphase der Schwangerschaft anderen als dem Arzt oder seinem Personal das Geschlecht des Kindes bekannt wird, es sei denn, die Mitteilung wäre medizinisch indiziert.
\end{sourcepara}
\begin{transpara}
\transrn{251}另一方面，基於對未出生生命之保護義務，必須對醫師檢查與資訊提供劃定界限，以因應基於性別選擇理由終止妊娠之危險；此種性別選擇受憲法秩序否定評價。因此，必須排除在懷孕早期由醫師或其人員以外之人知悉子女性別，除非該告知具醫學指標事由。
\end{transpara}

\begin{sourcepara}[title={原文D.138}]
\sourcern{252}c) Da verantwortliches ärztliches Handeln normativ am Schutz des ungeborenen Lebens orientiert ist, hat der Arzt zu berücksichtigen, unter welchen Voraussetzungen die Rechtsordnung einen Schwangerschaftsabbruch als nicht rechtswidrig ansieht. Notwendig dafür ist eine -- dem Recht verpflichtete -- ärztliche Beurteilung der Konfliktlage, die als Grundlage des Gesprächs mit der Frau und der eigenen Entscheidung des Arztes dienen kann. Dabei hat der Arzt der Frau die für seine Entscheidung maßgeblichen Gesichtspunkte mitzuteilen.
\end{sourcepara}
\begin{transpara}
\transrn{252}c) 由於負責任之醫療行為在規範上以保護未出生生命為取向，醫師必須考量法秩序在何等前提下將終止妊娠視為不具違法性。為此必要的是一項受法律拘束之醫療判斷，針對衝突狀態而為，並得作為與婦女談話及醫師自身決定之基礎。在此，醫師必須向婦女告知對其決定具有重要性之觀點。
\end{transpara}

\begin{sourcepara}[title={原文D.139}]
\sourcern{253}d) Nach den allgemeinen Grundsätzen des ärztlichen Berufsrechts ist der Arzt gehalten, über die in Ausübung seines Berufes gemachten Feststellungen und getroffenen Maßnahmen hinreichende Aufzeichnungen zu fertigen (vgl. § 11 \abbrAbs{} 1 der Musterberufsordnung). Diese Dokumentationspflicht hat sich auch auf die Erfüllung der erweiterten Aufklärungs- und Beratungspflicht bei einem Schwangerschaftsabbruch, wie sie sich aus den vorstehenden Darlegungen ergibt, zu erstrecken. Dies muß in den berufsrechtlichen Vorschriften klargestellt sein.
\end{sourcepara}
\begin{transpara}
\transrn{253}d) 依醫師職業法之一般原則，醫師有義務就其執行職業時所作確認及所採措施製作足夠紀錄（參見《醫師示範職業規則》第11條第1項）。此一文件記載義務亦必須延伸至終止妊娠時擴張之告知及諮詢義務之履行，如前述說明所得者。此點必須在職業法規定中予以明確化。
\end{transpara}

\begin{sourcepara}[title={原文D.140}]
\sourcern{254}e) Die Erfordernisse eines ärztlichen Gesprächs mit der Frau und einer Entscheidung des Arztes in ärztlicher Verantwortung stehen nicht in Widerspruch zur Konzeption einer Beratungsregelung; denn die Frau wird nicht dem Druck ausgesetzt, die Gründe für ihren Abbruchwunsch durch einen Dritten überprüfen und bewerten zu lassen. Die Feststellung und Beurteilung einer Indikation wird von dem Arzt gerade nicht verlangt, wenn er sich ein Bild darüber machen soll, ob er nach seinem ärztlichen Selbstverständnis seine Mitwirkung bei dem von der Frau gewünschten Abbruch verantworten kann.
\end{sourcepara}
\begin{transpara}
\transrn{254}e) 醫師與婦女談話之要求，以及醫師依醫療責任作成決定之要求，並不與諮詢規制之構想相矛盾；因為婦女並未被置於壓力下，須使其終止妊娠願望之理由受第三人審查及評價。當醫師應形成一項圖像，以判斷其依自身醫師理解能否對參與婦女所希望之終止妊娠負責時，並不要求醫師確認及評價一項指標事由。
\end{transpara}

\begin{sourcepara}[title={原文D.141}]
\sourcern{255}Eine solche ärztliche Entscheidung ist mit dem Beruf des Arztes und seiner ureigenen Aufgabe, Leben zu erhalten, unaufhebbar verbunden. Dementsprechend haben die in der mündlichen Verhandlung geladenen Vertreter der ärztlichen Standesorganisationen übereinstimmend erklärt, auch wenn die strafrechtliche Regelung nicht verlange, daß die Ärzte sich von der Schwangeren die Gründe für den Abbruchwunsch mitteilen lassen, sei diese Mitteilung für ärztlich verantwortliches Handeln ebenso unerläßlich wie die Prüfung, ob das Gespräch mit der Patientin ihnen die Überzeugung vermitteln könne, daß der Abbruchwunsch auf einem eigenen, verantwortlichen Entschluß und achtenswerten Gründen beruhe.
\end{sourcepara}
\begin{transpara}
\transrn{255}此種醫療決定與醫師職業及其保存生命之本有任務不可分離。相應地，在言詞辯論中受傳喚之醫師職業團體代表一致表示，即使刑法規制並未要求醫師讓懷孕婦女告知其終止妊娠願望之理由，此一告知對負責任之醫療行為仍與下列審查同樣不可或缺：與病人之談話是否能使醫師確信，該終止妊娠願望係基於其自身負責任之決意及值得尊重之理由。
\end{transpara}

\begin{sourcepara}[title={原文D.142}]
\sourcern{256}2. a) Verlangt die Frau weiterhin den Schwangerschaftsabbruch, nachdem sie außer der Beratung auch ärztlichen Rat in dem dargestellten Sinne erfahren hat, ist es verfassungsrechtlich unbedenklich, daß in einem Beratungskonzept der Ausschluß der Strafdrohung (Tatbestandsausschluß -- vgl. oben III.2.a) für sie nicht mehr davon abhängt, ob der Arzt den Schwangerschaftsabbruch für zulässig hält. In diesem Konzept soll der Arzt nicht mit verbindlicher Wirkung entscheiden, ob der Schwangerschaftsabbruch im Einzelfall stattfinden darf.
\end{sourcepara}
\begin{transpara}
\transrn{256}2. a) 若婦女除諮詢外亦已依前述意義接受醫療建議後，仍繼續要求終止妊娠，則在諮詢構想中，排除對其之刑罰威嚇（構成要件排除，參見上文 III.2.a）不再取決於醫師是否認為終止妊娠可被允許，在憲法上並無疑義。在此構想中，醫師不應以具拘束力之方式決定終止妊娠在個案中是否得以發生。
\end{transpara}

\begin{sourcepara}[title={原文D.143}]
\sourcern{257}b) Hält der Arzt den Abbruch für ärztlich verantwortbar, so muß er daran mitwirken können, ohne daß ihm Strafe droht. Hält er den Abbruch für ärztlich nicht verantwortbar, so ist er zwar aufgrund seiner allgemeinen Berufspflichten gehalten, seine Mitwirkung abzulehnen. Sanktionen für die Verletzung von Pflichten solcher Art sind in der Rechtspraxis nur schwer durchsetzbar. Die verfassungsrechtliche Schutzpflicht für das ungeborene Leben gebietet es nicht, für Verstöße im Strafrecht eine Ahndung vorzusehen. Zur Erfüllung der Schutzpflicht erscheint es insoweit ausreichend, aber auch erforderlich, daß die genannte Verpflichtung und ihre Durchsetzung in den ärztlichen Berufsordnungen geregelt wird; dies muß unabhängig von einer strafrechtlichen Regelung geschehen.
\end{sourcepara}
\begin{transpara}
\transrn{257}b) 若醫師認為終止妊娠在醫療上可負責，則其必須能夠參與而不受刑罰威脅。若其認為終止妊娠在醫療上不可負責，則依其一般職業義務，固然有義務拒絕參與。此類義務違反之制裁，在法律實務中很難執行。對未出生生命之憲法保護義務，並不命令在刑法中對違反行為規定處罰。就此而言，為履行保護義務，顯得足夠但亦必要的是，在醫師職業規則中規範上述義務及其執行；此必須獨立於刑法規制而為。
\end{transpara}

\begin{sourcepara}[title={原文D.144}]
\sourcern{258}c) Strafrechtlicher Sanktion zugänglich und im Rahmen eines Beratungskonzepts bedürftig ist dagegen, daß der Arzt sich die Gründe der Frau für ihr Abbruchverlangen darlegen läßt, sich der vorausgegangenen Beratung sowie der Überlegungsfrist vergewissert und er seine besondere, dem Lebensschutz dienende Aufklärungs- und Beratungspflicht erfüllt. Außerdem muß die Verpflichtung, das Alter der Schwangerschaft festzustellen und in den ersten zwölf Wochen der Schwangerschaft keine Mitteilungen über das Geschlecht des zu erwartenden Kindes zu machen, strafbewehrt sein. Auf die Sicherung dieses in seinen Berufspflichten begründeten Handelns des Arztes, das nicht allein den Belangen der Frau dient, sondern zugleich eine wichtige Stütze für den Schutz des ungeborenen Lebens darstellt, darf ein auf Schutzwirkung zielendes Beratungskonzept nicht verzichten.
\end{sourcepara}
\begin{transpara}
\transrn{258}c) 相反地，醫師使婦女陳述其終止妊娠要求之理由、確認先前諮詢及思考期間，並履行其特殊且服務於生命保護之告知及諮詢義務，乃可受刑法制裁且在諮詢構想框架內需要刑法制裁之事項。此外，確認懷孕年齡，以及在妊娠最初十二週內不告知預期子女之性別之義務，必須附有刑罰保障。一項以保護效果為目標之諮詢構想，不能放棄保障醫師此種奠基於其職業義務之行為；該行為不僅服務於婦女利益，同時亦構成未出生生命保護之重要支柱。
\end{transpara}

\begin{sourcepara}[title={原文D.145}]
\sourcern{259}Eine nur berufsrechtliche Normierung würde dabei der Schutzpflicht für das ungeborene Leben nicht genügen. Die bisherige Regelung des Schwangerschaftsabbruchs im ärztlichen Berufsrecht verweist auf das Strafrecht. Die Anhörung von Vertretern der ärztlichen Standesorganisationen in der mündlichen Verhandlung hat ergeben, daß diese auch bei einem Übergang zu einem Beratungskonzept keinen Anlaß für ergänzende berufsrechtliche Vorschriften sehen, die höhere Anforderungen an das ärztliche Verhalten normieren als das Strafrecht. Danach kann -- jedenfalls gegenwärtig -- nicht darauf gesetzt werden, daß den dargestellten ärztlichen Verhaltenspflichten durch eine berufsrechtliche Normierung und Sanktionierung in dem Maße zur Beachtung verholfen würde, wie sie für einen wirksamen Schutz des ungeborenen Lebens unabdingbar ist.
\end{sourcepara}
\begin{transpara}
\transrn{259}僅以職業法規範，將不能滿足對未出生生命之保護義務。迄今醫師職業法中關於終止妊娠之規制係指向刑法。言詞辯論中聽取醫師職業團體代表之結果顯示，即使轉向諮詢構想，其等亦不認為有理由制定補充性職業法規定，將高於刑法之要求規範為醫師行為。依此，至少目前不能期待，透過職業法規範與制裁，即能使前述醫師行為義務獲得如有效保護未出生生命所不可或缺之程度的遵守。
\end{transpara}

\begin{sourcepara}[title={原文D.146}]
\sourcern{260}3. Damit der Arzt, an den sich die Frau wegen eines Schwangerschaftsabbruchs wendet, den Anforderungen der erweiterten Aufklärungs- und Beratungspflicht gerecht werden kann, ist es angezeigt, durch Vorschriften über die ärztliche Aus- und Fortbildung in geeigneter Weise dafür zu sorgen, daß dem Arzt über das gynäkologische Fachwissen hinaus die für eine ärztliche Beurteilung von Schwangerschaftskonflikten notwendigen Kenntnisse vermittelt werden.
\end{sourcepara}
\begin{transpara}
\transrn{260}3. 為使婦女因終止妊娠而尋求之醫師，能符合擴張告知及諮詢義務之要求，宜透過關於醫師養成及進修之規定，以適當方式確保醫師除婦科專業知識外，亦被傳授對懷孕衝突作醫療判斷所必要之知識。
\end{transpara}

\begin{sourcepara}[title={原文D.147}]
\sourcern{261}4. Der Arzt kann die ihm im Schutzkonzept einer Beratungsregelung zukommende Funktion nur zuverlässig wahrnehmen, wenn ihm aus seiner Weigerung, Schwangerschaftsabbrüche vorzunehmen, grundsätzlich keine rechtlichen oder tatsächlichen Nachteile erwachsen. Auch fällt das Recht, die Mitwirkung an Schwangerschaftsabbrüchen -- mit Ausnahme medizinisch indizierter -- zu verweigern, in den Schutzbereich seines durch das ärztliche Berufsbild geprägten Persönlichkeitsrechts (\abbrArt{} 2 \abbrAbs{} 1 in Verbindung mit \abbrArt{} 12 \abbrAbs{} 1 \abbrGG{}\ggfnArtTwelveOneDE{}). Dem hat der Gesetzgeber mit \abbrArt{} 2 des 5. \abbrStrRG{} Rechnung getragen. Diese Vorschrift ist bei Berücksichtigung der verfassungsrechtlichen Vorgaben für die Regelung des Schwangerschaftsabbruchs vertraglich nicht abdingbar. Die §§ 627, 628 BGB sind unter Berücksichtigung des Weigerungsrechts anzuwenden. Auch wenn ein Arzt in abhängiger Stellung sich generell weigert, nicht medizinisch indizierte Schwangerschaftsabbrüche vorzunehmen, dürfen sich für ihn daraus regelmäßig keine beruflichen Nachteile ergeben; eine Kündigung des Arbeitsverhältnisses kann auch dann nur ausnahmsweise in Betracht kommen, wenn der Arbeitgeber keine andere Möglichkeit hat, den Arzt zu beschäftigen. Die fachärztliche Ausbildung darf an der Weigerung des Arztes, sich in dem beschriebenen Umfang an Schwangerschaftsabbrüchen zu beteiligen, nicht scheitern.
\end{sourcepara}
\begin{transpara}
\transrn{261}4. 醫師只有在其因拒絕實施終止妊娠原則上不會遭受法律或事實不利益時，才能可靠地履行其在諮詢規制保護構想中所具有之功能。拒絕參與終止妊娠，醫學指標事由者除外，之權利，亦落入由醫師職業圖像所形塑之人格權保護範圍（《基本法》第2條第1項\ggfnArtTwoOneZH{}結合第12條第1項\ggfnArtTwelveOneZH{}）。立法者已以《第五刑法改革法》第2條顧及此點。考量終止妊娠規制之憲法預設，該規定不得以契約排除。《民法》第627條、第628條應在考量拒絕權下適用。即使處於從屬地位之醫師一般性拒絕實施非醫學指標事由之終止妊娠，通常亦不得因此對其產生職業上不利益；勞動關係之終止，即使如此，亦只有在雇主沒有其他可能任用該醫師時，始得例外考慮。專科醫師訓練不得因醫師拒絕在前述範圍內參與終止妊娠而失敗。
\end{transpara}

\begin{sourcepara}[title={原文D.148}]
\sourcern{262}5. \abbrArt{} 3 \abbrAbs{} 1 Satz 1 des 5. \abbrStrRG{} in der Fassung des \abbrArt{} 15 \abbrNr{} 1 \abbrSFHG{} schreibt vor, daß ein Schwangerschaftsabbruch nur in einer Einrichtung vorgenommen werden darf, in der auch die notwendige medizinische Nachbehandlung gewährleistet ist. Es liegt nahe, daß auf der Grundlage dieser Vorschrift Einrichtungen entstehen werden, die sich auf die Durchführung von Schwangerschaftsabbrüchen spezialisieren und ihre Tätigkeit im wesentlichen darauf ausrichten. Die hieraus entstehenden Gefahren für die Erfüllung der dem Arzt, zumal im Rahmen einer Beratungsregelung, zufallenden Aufgabe beim Schutz des ungeborenen menschlichen Lebens liegen auf der Hand. Die verfassungsrechtliche Schutzpflicht gebietet es daher dem Gesetzgeber zu prüfen, in welcher Weise solchen Gefahren wirksam entgegengetreten werden kann, und geeignete Regelungen zu treffen. Für diese Prüfung könnten die Erfahrungen in Betracht kommen, die in Frankreich mit der Begrenzung der Zahl der Abbrüche auf einen bestimmten Anteil der insgesamt vorgenommenen ärztlichen Verrichtungen und der einheitlich festgelegten Höhe der Vergütung für einen Schwangerschaftsabbruch gemacht worden sind (vgl. Eser/Koch, Schwangerschaftsabbruch im internationalen Vergleich, Teil 1 Europa, 1988, \abbrS{} 520 \abbrF{}).
\end{sourcepara}
\begin{transpara}
\transrn{262}5. SFHG 第15條第1款版本之《第五刑法改革法》第3條第1項第1句規定，終止妊娠只能在亦能確保必要醫療後續照護之機構中實施。可預期的是，基於此一規定將產生專門實施終止妊娠並將其活動主要指向此點之機構。由此對醫師在保護未出生人類生命時，尤其在諮詢規制框架內所負任務之履行所生危險，顯而易見。因此，憲法保護義務命令立法者審查，應以何種方式有效對抗此等危險，並制定適當規制。此一審查得考慮法國經驗；法國曾限制終止妊娠數量，使其僅能占全部醫療處置之一定比例，並統一規定終止妊娠報酬數額（參見 Eser/Koch, Schwangerschaftsabbruch im internationalen Vergleich, Teil 1 Europa, 1988, S. 520 f.）。
\end{transpara}

\begin{sourcepara}[title={原文D.149}]
\sourcern{263}6. Die staatliche Schutzpflicht für das ungeborene Leben gebietet nicht, Arzt- und Krankenhausverträge über Schwangerschaftsabbrüche, die nach dem Beratungskonzept nicht mit Strafe bedroht werden, als rechtlich unwirksam anzusehen. Das Konzept erfordert es vielmehr, daß der Leistungsaustausch zwischen Arzt und Frau als Rechtsverhältnis ausgestaltet wird, mithin die Leistungen mit Rechtsgrund gewährt werden. Unbeschadet der Beurteilung von Rechtsfolgen des Vertrags im einzelnen greifen daher die §§ 134, 138 BGB nicht ein. Arzt und Krankenhausträger sollen an dem Schwangerschaftsabbruch nur aufgrund eines wirksamen Vertrages mitwirken, der ihre Rechte, insbesondere den Entgeltanspruch, sichert, ebenso aber auch ihre Pflichten regelt. Vor allem bedarf der vom Arzt zu gewährleistende Schutz des ungeborenen Lebens und der Gesundheit der Frau vertragsrechtlicher Absicherung. Die Schlechterfüllung der Beratungs- und Behandlungspflichten muß deshalb grundsätzlich auch vertrags- und deliktsrechtliche Sanktionen auslösen können.
\end{sourcepara}
\begin{transpara}
\transrn{263}6. 國家對未出生生命之保護義務，並不命令將依諮詢構想不受刑罰威嚇之終止妊娠的醫師契約及醫院契約視為法律上無效。該構想毋寧要求，醫師與婦女間之給付交換被形塑為法律關係，因而給付是基於法律原因而提供。因此，不影響個別評價契約法律效果，《民法》第134條、第138條並不介入。醫師與醫院承辦者應僅基於有效契約參與終止妊娠；該契約保障其權利，尤其是報酬請求權，同時亦規範其義務。尤其，醫師所應保障之未出生生命保護及婦女健康，需要契約法上保障。因此，諮詢及治療義務之不完全履行，原則上亦必須能引發契約法及侵權法制裁。
\end{transpara}

\begin{sourcepara}[title={原文D.150}]
\sourcern{264}Dies bedarf allerdings aus verfassungsrechtlicher Sicht differenzierender Betrachtung. Eine zivilrechtliche Sanktion für die Schlechterfüllung des Vertrages und für eine deliktische Beeinträchtigung der körperlichen Integrität der Frau ist grundsätzlich erforderlich; dies betrifft nicht nur eine Verpflichtung zur Rückzahlung der vergeblich geleisteten Vergütung, sondern auch den Ersatz von Schäden einschließlich -- im Rahmen der §§ 823, 847 BGB -- einer billigen Entschädigung der Frau für die immateriellen Belastungen, die mit dem fehlgeschlagenen Schwangerschaftsabbruch oder mit der Geburt eines behinderten Kindes für sie verbunden waren. Eine rechtliche Qualifikation des Daseins eines Kindes als Schadensquelle kommt hingegen von Verfassungs wegen (\abbrArt{} 1 \abbrAbs{} 1 \abbrGG{}\ggfnArtOneOneDE{}) nicht in Betracht. Die Verpflichtung aller staatlichen Gewalt, jeden Menschen in seinem Dasein um seiner selbst willen zu achten (vgl. oben I.1.a), verbietet es, die Unterhaltspflicht für ein Kind als Schaden zu begreifen. Die Rechtsprechung der Zivilgerichte zur Haftung für ärztliche Beratungsfehler oder für fehlgeschlagene Schwangerschaftsabbrüche (zum Schwangerschaftsabbruch vgl.BGHZ 86, 240 \abbrFF{}; 89, 95 \abbrFF{}; 95, 199 \abbrFF{}; BGH, NJW 1985, \abbrS{} 671 \abbrFF{}; VersR 1985, \abbrS{} 1068 \abbrFF{}; VersR 1986, \abbrS{} 869 \abbrF{}; VersR 1988, \abbrS{} 155 \abbrF{}; NJW 1992, \abbrS{} 1556 \abbrFF{}; zur Sterilisation vgl. BGHZ 76, 249 \abbrFF{}; 76, 259 \abbrFF{}; BGH, NJW 1984, \abbrS{} 2625 \abbrF{}) ist im Blick darauf der Überprüfung bedürftig. Hiervon unberührt bleibt eine Schadensersatzpflicht des Arztes gegenüber dem Kind wegen Schädigungen, die diesem bei einem nicht kunstgerecht ausgeführten, mißlungenen Schwangerschaftsabbruch zugefügt worden sind (vgl. BGHZ 58, 48 [49 \abbrFF{}]; BGH, NJW 1989, \abbrS{} 1538 [1539]).
\end{sourcepara}
\begin{transpara}
\transrn{264}不過，這從憲法觀點需要區分觀察。對契約不完全履行及對婦女身體完整性之侵權損害，原則上需要民事法制裁；這不僅涉及返還徒然支付之報酬的義務，也包括損害賠償，並在《民法》第823條、第847條框架內，包括就失敗之終止妊娠或身心障礙子女出生所伴隨之非財產負擔，給予婦女公平補償。相反地，從憲法上看（《基本法》第1條第1項\ggfnArtOneOneZH{}），不得將子女之存在在法律上定性為損害來源。一切國家權力負有尊重每一人為其自身而存在之義務（參見上文 I.1.a），此禁止將對子女之扶養義務理解為損害。民事法院關於醫師諮詢錯誤或失敗終止妊娠之責任判例（關於終止妊娠，參見 BGHZ 86, 240 ff.; 89, 95 ff.; 95, 199 ff.; BGH, NJW 1985, S. 671 ff.; VersR 1985, S. 1068 ff.; VersR 1986, S. 869 f.; VersR 1988, S. 155 f.; NJW 1992, S. 1556 ff.; 關於絕育，參見 BGHZ 76, 249 ff.; 76, 259 ff.; BGH, NJW 1984, S. 2625 f.），鑑於此點需要重新審查。不受此影響者，是醫師對子女之損害賠償義務，該損害係因未依技術規則實施且失敗之終止妊娠而加諸於子女（參見 BGHZ 58, 48 [49 ff.]; BGH, NJW 1989, S. 1538 [1539]）。
\end{transpara}

\bisubsubsection{VI.}{VI.}

\begin{sourcepara}[title={原文D.151}]
\sourcern{265}Die staatliche Schutzpflicht für das ungeborene Leben umfaßt auch den Schutz vor Gefahren, die von Dritten ausgehen, nicht zuletzt von Personen aus dem familiären oder dem weiteren sozialen Umfeld der Schwangeren (siehe oben I.2. und III.1.b). Eine erhöhte Bedeutung wächst dieser Verpflichtung beim Übergang auf ein Schutzkonzept zu, das in der Frühphase der Schwangerschaft im Schwangerschaftskonflikt primär auf eine Beratung der Frau setzt, um sie für das Kind zu gewinnen.
\end{sourcepara}
\begin{transpara}
\transrn{265}國家對未出生生命之保護義務，亦包括防禦由第三人所生之危險，尤其包括由懷孕婦女家庭或較廣泛社會環境中之人所生之危險（見上文 I.2. 及 III.1.b）。當轉向一項保護構想，而該構想在懷孕早期之懷孕衝突中主要依靠對婦女之諮詢以爭取其接受子女時，此一義務取得更高重要性。
\end{transpara}

\begin{sourcepara}[title={原文D.152}]
\sourcern{266}1. Die Wirksamkeit dieses Schutzkonzepts erfordert in besonderem Maß, die Frau vor Zumutungen zu schützen, die sie wegen der Schwangerschaft in Bedrängnis bringen oder einen Druck auf sie ausüben, die Schwangerschaft abzubrechen. Solche Einwirkungen sind auch geeignet, den Erfolg der Beratung zunichte zu machen, wenn die Personen ihres Umfeldes einer Frau, die sich durch die Beratung zur Fortsetzung der Schwangerschaft ermutigt sieht, nun um so härter zusetzen, um sie unter Hinweis auf die Straflosigkeit und ihre Letztverantwortung zum Schwangerschaftsabbruch zu bringen und sich damit zugleich der eigenen (Mit-)Verantwortung zu entziehen.
\end{sourcepara}
\begin{transpara}
\transrn{266}1. 此一保護構想之有效性，特別要求保護婦女免受因懷孕而使其陷入困境或對其施加終止妊娠壓力之要求。若婦女因諮詢而受到鼓勵繼續懷孕，其環境中之人現在反而更加嚴厲地對待她，並援引不罰及其最後責任而促使其終止妊娠，同時藉此逃避自身之共同責任，則此等影響亦適於使諮詢成果歸於無效。
\end{transpara}

\begin{sourcepara}[title={原文D.153}]
\sourcern{267}Wie Berichte aus der Beratungspraxis, Befragungen und wissenschaftliche Untersuchungen aufweisen, haben Schwangerschaftskonflikte, die schließlich zum Schwangerschaftsabbruch führen, ihre Ursache zu einem erheblichen Teil nicht primär in wirtschaftlich-sozialen Notlagen, sondern in gestörten Partnerschaftsbeziehungen, in der Ablehnung des Kindes durch den Vater oder die Eltern der Frau sowie in einem Druck, der von diesen ausgeübt wird (vgl. Renate Köcher, Schwangerschaftsabbruch -- betroffene Frauen berichten, in: Aus Politik und Zeitgeschichte B 14/90, \abbrS{} 37 \abbrF{}; Deutscher Caritasverband, 13. Erhebung: Werdende Mütter in Not- und Konfliktsituationen, Zeitraum 1989, \abbrS{} 45 \abbrFF{}; Roeder/Sellschopp/Henrich, Die Rolle des Mannes bei Schwangerschaftskonflikten, Abschlußbericht Oktober 1992). Dies zeigt, wie groß die Bedrängnis -- bis hin zur Ausweglosigkeit -- für eine Frau werden kann, wenn die Personen, die ihr nahestehen und von denen sie gerade auch bei einer ungewollten Schwangerschaft mit Grund Beistand und Hilfe erwarten darf, dies verweigern.
\end{sourcepara}
\begin{transpara}
\transrn{267}如諮詢實務報告、調查及科學研究所顯示，最終導致終止妊娠之懷孕衝突，其原因有相當部分並非主要在經濟社會困境，而是在伴侶關係失調、父親或婦女父母拒絕子女，以及其等施加之壓力（參見 Renate Köcher, Schwangerschaftsabbruch -- betroffene Frauen berichten, in: Aus Politik und Zeitgeschichte B 14/90, S. 37 f.; Deutscher Caritasverband, 13. Erhebung: Werdende Mütter in Not- und Konfliktsituationen, Zeitraum 1989, S. 45 ff.; Roeder/Sellschopp/Henrich, Die Rolle des Mannes bei Schwangerschaftskonflikten, Abschlußbericht Oktober 1992）。這顯示，當與婦女親近、且婦女尤其在非意願懷孕時有理由期待其支持與協助之人拒絕此等支持與協助時，婦女所受壓迫可能大到何種程度，甚至達到無路可走。
\end{transpara}

\begin{sourcepara}[title={原文D.154}]
\sourcern{268}2. Dem muß die Rechtsordnung entgegentreten; sie muß der Frau um der Wirkungschancen der Beratung willen einen Raum eigener, nicht durch Druck von außen determinierter Verantwortlichkeit sichern. Daher sind Personen des familiären Umfeldes in das Schutzkonzept einzubeziehen, die für eine Schwangerschaft ebenfalls Verantwortung tragen, wie die Väter, oder für die durch die Schwangerschaft eine besondere Verantwortung entsteht, wie bei den Eltern noch minderjähriger Schwangerer. In Betracht kommen aber auch Personen des weiteren sozialen Umfeldes, wie etwa Vermieter oder Arbeitgeber.
\end{sourcepara}
\begin{transpara}
\transrn{268}2. 法秩序必須對抗此點；為了諮詢之作用機會，法秩序必須為婦女保障一個自身責任空間，該空間不受外部壓力決定。因此，應將家庭環境中亦對懷孕負責之人，例如父親，或因懷孕而產生特殊責任之人，例如未成年懷孕婦女之父母，納入保護構想。但較廣泛社會環境中之人，例如出租人或雇主，亦可列入考量。
\end{transpara}

\begin{sourcepara}[title={原文D.155}]
\sourcern{269}a) Es wäre daher nicht zureichend, in einem Beratungskonzept, das auf den Schutz des ungeborenen Lebens ausgerichtet ist, lediglich an das Verantwortungsbewußtsein dieser Personen zu appellieren. Vielmehr ist darauf hinzuwirken, daß durch gesetzliche Regelungen oder durch die Anwendung bereits bestehender rechtlicher Vorschriften Bedingungen dafür geschaffen werden, daß die familiäre Verantwortung wie auch die im weiteren sozialen Umfeld gebotene Rücksicht eingefordert werden können (vgl. auch oben I.3.a).
\end{sourcepara}
\begin{transpara}
\transrn{269}a) 因此，在一項以保護未出生生命為取向之諮詢構想中，僅訴諸此等人之責任意識並不足夠。毋寧，應促成透過法律規制或適用既有法律規定，創設條件，使家庭責任以及較廣泛社會環境中所要求之顧慮得以被請求（亦參見上文 I.3.a）。
\end{transpara}

\begin{sourcepara}[title={原文D.156}]
\sourcern{270}b) Darüber hinaus sind für Personen des familiären Umfeldes in bestimmtem Umfang strafbewehrte Verhaltensgebote und -verbote unerläßlich. Sie müssen sich zum einen darauf richten, daß die betreffenden Personen der Frau den ihnen zuzumutenden Beistand, dessen sie wegen der Schwangerschaft bedarf, nicht in verwerflicher Weise vorenthalten, zum anderen darauf, daß sie es unterlassen, die Frau zum Schwangerschaftsabbruch zu drängen. Dabei kann die Strafbarkeit davon abhängig gemacht werden, daß es zum Abbruch der Schwangerschaft gekommen ist. Vorschriften dieser Art würden an Überlegungen anknüpfen, wie sie der Strafrechtsreformentwurf 1962 in § 201 erwogen hat (vgl. \abbrBRDrucks{} 200/62 \abbrS{} 45).
\end{sourcepara}
\begin{transpara}
\transrn{270}b) 此外，對家庭環境中之人，在特定範圍內附有刑罰保障之行為命令與禁止不可或缺。其一方面必須指向，使相關人不得以可非難方式不給予婦女因懷孕所需、且可期待其提供之支持；另一方面則指向，使其不得迫使婦女終止妊娠。在此，可將可罰性繫於終止妊娠已經發生。此類規定將連結於1962年刑法改革草案第201條曾考慮之想法（參見 BRDrucks. 200/62 S. 45）。
\end{transpara}

\begin{sourcepara}[title={原文D.157}]
\sourcern{271}Zu prüfen ist, ob vergleichbare Sanktionen auch gegenüber Personen des weiteren sozialen Umfeldes vorzusehen sind, wenn diese die Frau zum Abbruch der Schwangerschaft drängen oder sie in Kenntnis der Schwangerschaft in verwerflicher Weise in eine Notlage bringen, die die Gefahr eines Schwangerschaftsabbruchs herbeiführt.
\end{sourcepara}
\begin{transpara}
\transrn{271}應審查是否亦須對較廣泛社會環境中之人規定可比較之制裁，若其迫使婦女終止妊娠，或明知其懷孕而以可非難方式使其陷入造成終止妊娠危險之困境。
\end{transpara}

\bisubsection{E.}{E.}

\begin{sourcepara}[title={原文E.158}]
\sourcern{272}Prüft man nach diesen Maßstäben die angegriffenen Regelungen des Schwangeren- und Familienhilfegesetzes, so ergibt sich, daß das Gesetz bei dem an sich zulässigen Übergang zu einem Beratungskonzept für die ersten zwölf Schwangerschaftswochen der Verpflichtung aus \abbrArt{} 1 \abbrAbs{} 1 in Verbindung mit \abbrArt{} 2 \abbrAbs{} 2 Satz 1 \abbrGG{}\ggfnArtTwoTwoOneDE{}, das ungeborene Leben wirksam zu schützen,nicht in dem gebotenen Umfang gerecht geworden ist. Die gegen einzelne Vorschriften erhobenen kompetenzrechtlichen Bedenken greifen dagegen nur bezüglich des \abbrArt{} 4 des 5. \abbrStrRG{} \abbrNF{} teilweise durch.
\end{sourcepara}
\begin{transpara}
\transrn{272}依此等標準審查受攻擊之《孕婦及家庭扶助法》規制，可得出：該法雖本得就妊娠最初十二週轉向諮詢構想，卻未在所要求之範圍內履行《基本法》第1條第1項結合第2條第2項第1句\ggfnArtOneOneTwoTwoOneZH{}所生、有效保護未出生生命之義務。相反地，針對個別規定所提出之權限法疑慮，僅就新版本《第五刑法改革法》第4條部分成立。
\end{transpara}

\bisubsubsection{I.}{I.}

\begin{sourcepara}[title={原文E.159}]
\sourcern{273}§ 218a \abbrAbs{} 1 \abbrStGB{} \abbrNF{}, wonach Schwangerschaftsabbrüche, die innerhalb der ersten zwölf Schwangerschaftswochen nach Beratung gemäß § 219 \abbrStGB{} \abbrNF{} auf Verlangen der schwangeren Frau durch einen Arzt vorgenommen werden, "nicht rechtswidrig" sind, ist mit der Schutzpflicht für das ungeborene menschliche Leben (\abbrArt{} 1 \abbrAbs{} 1 in Verbindung mit \abbrArt{} 2 \abbrAbs{} 2 Satz 1 \abbrGG{}\ggfnArtTwoTwoOneDE{}) unvereinbar und nichtig.
\end{sourcepara}
\begin{transpara}
\transrn{273}新版本《刑法》第218a條第1項規定，於妊娠最初十二週內，依新版本《刑法》第219條接受諮詢後，依懷孕婦女要求由醫師實施之終止妊娠為「不具違法性」；該規定與對未出生人類生命之保護義務（《基本法》第1條第1項結合第2條第2項第1句\ggfnArtOneOneTwoTwoOneZH{}）不相容，並且無效。
\end{transpara}

\begin{sourcepara}[title={原文E.160}]
\sourcern{274}1. Diese Vorschrift normiert einen Rechtfertigungsgrund, der in seinen Wirkungen den in § 218a \abbrAbs{} 2 und 3 \abbrStGB{} \abbrNF{} bezeichneten Rechtfertigungsgründen (medizinische und embryopathische Indikation) gleichsteht. Dies war, wie sich aus den Beratungen des Sonderausschusses "Schutz des ungeborenen Lebens" ergibt, auch beabsichtigt (vgl. die Äußerungen der Abgeordneten Baum und Pflüger im Protokoll der 17. Sitzung des Sonderausschusses, \abbrS{} 11, 12). Die Wirkungen dieses Rechtfertigungsgrundes beschränken sich wegen der Durchschlagskraft, die im Bereich des Schutzes elementarer Rechtsgüter einem strafrechtlichen Rechtfertigungsgrund für die gesamte Rechtsordnung zukommt, nicht auf einen Ausschluß der Strafbarkeit (vgl. oben D.III.2.a).
\end{sourcepara}
\begin{transpara}
\transrn{274}1. 該規定規範一項正當化事由，其效果等同於新版本《刑法》第218a條第2項及第3項所稱正當化事由（醫學及胚胎病理指標事由）。如「未出生生命保護」特別委員會之審議所顯示，這亦是其意圖（參見 Baum 及 Pflüger 議員於特別委員會第17次會議紀錄 S. 11, 12 之發言）。由於在基本法益保護領域中，刑法上正當化事由對整體法秩序具有貫穿力，該正當化事由之效果並不限於排除可罰性（參見上文 D.III.2.a）。
\end{transpara}

\begin{sourcepara}[title={原文E.161}]
\sourcern{275}2. Begründet § 218a \abbrAbs{} 1 \abbrStGB{} \abbrNF{} einen allgemeinen Rechtfertigungsgrund für Schwangerschaftsabbrüche, so genügen seine inhaltlichen Voraussetzungen nicht den verfassungsrechtlichen Anforderungen (vgl. oben D.I.2.c)bb). Eine Rechtfertigung des Schwangerschaftsabbruchs kann nur in Betracht kommen, wenn eine -- festzustellende -- tatbestandlich näher umschriebene Notlage gegeben ist, die -- vergleichbar den Fällen der medizinischen, embryopathischen oder kriminologischen Indikation -- so schwer wiegt, daß der schwangeren Frau die Austragung der Schwangerschaft nicht zugemutet werden kann. Eine solche Notlage wird in § 218a \abbrAbs{} 1 \abbrStGB{} \abbrNF{} nicht vorausgesetzt.
\end{sourcepara}
\begin{transpara}
\transrn{275}2. 若新版本《刑法》第218a條第1項為終止妊娠建立一般正當化事由，則其內容前提不符合憲法要求（參見上文 D.I.2.c)bb）。終止妊娠之正當化只有在存在一項須被確認且以構成要件方式更詳盡描述之困境時，始可能被考慮；該困境必須可與醫學、胚胎病理或犯罪學指標事由案件相比，如此嚴重，以致不能期待懷孕婦女繼續懷孕。新版本《刑法》第218a條第1項並未以此種困境為前提。
\end{transpara}

\begin{sourcepara}[title={原文E.162}]
\sourcern{276}Der Tatbestand des § 218a \abbrAbs{} 1 \abbrStGB{} \abbrNF{} handelt in seinen Voraussetzungen nicht von dieser Unzumutbarkeit: Seit der Empfängnis dürfen nicht mehr als zwölf Wochen vergangen sein; der Schwangerschaftsabbruch muß von einem Arzt vorgenommen werden; die schwangere Frau muß ihn verlangen; sie muß dem Arzt durch eine Bescheinigung nachweisen, daß mindestens drei Tage vor dem Eingriff eine Beratung nach § 219 \abbrStGB{} \abbrNF{} stattgefunden hat. Der Hinweis in § 218a \abbrAbs{} 1 \abbrNr{} 1 \abbrStGB{} \abbrNF{} auf die Überschrift des § 219 \abbrStGB{} \abbrNF{} ("Beratung der Schwangeren in einer Not- und Konfliktlage") zeigt zwar, daß das Gesetz vom Vorliegen einer Not- und Konfliktlage ausgeht; diese ist jedoch nicht näher qualifiziert; zudem wird die Rechtfertigung nicht von deren Feststellung abhängig gemacht.
\end{sourcepara}
\begin{transpara}
\transrn{276}新版本《刑法》第218a條第1項之構成要件，其前提並未處理此一不可期待性：自受孕起不得超過十二週；終止妊娠必須由醫師實施；懷孕婦女必須要求之；她必須以證明向醫師證明，至少在介入前三日已依新版本《刑法》第219條接受諮詢。新版本《刑法》第218a條第1項第1款對新版本《刑法》第219條標題（「困境與衝突狀態中之懷孕婦女諮詢」）之指稱，固然顯示法律以困境與衝突狀態之存在為出發點；然而該狀態並未被更詳盡定性；此外，正當化亦未取決於其確認。
\end{transpara}

\begin{sourcepara}[title={原文E.163}]
\sourcern{277}Ein Beratungskonzept, wie es § 218a \abbrAbs{} 1 \abbrStGB{} \abbrNF{} zugrundeliegt, kann nach den oben dargestellten verfassungsrechtlichen Maßstäben nicht zu einer Rechtfertigung des Schwangerschaftsabbruchs führen (vgl. D.I.2.c) und III.2.). Im Rahmen eines Beratungskonzepts ist es nur zulässig, die Rechtswirkungen des grundsätzlichen Verbots des Schwangerschaftsabbruchs unter bestimmten Voraussetzungen in einzelnen Bereichen der Rechtsordnung einzuschränken, wenn das Schutzkonzept zu seiner Wirksamkeit eine solche Ausnahme fordert; danach können insbesondere im Strafrecht die nach dem Beratungskonzept durchgeführten Schwangerschaftsabbrüche durch einen Tatbestandsausschluß von der Strafdrohung ausgenommen werden.
\end{sourcepara}
\begin{transpara}
\transrn{277}依前述憲法標準，如新版本《刑法》第218a條第1項所依據之諮詢構想，不能導致終止妊娠之正當化（參見 D.I.2.c) 及 III.2.）。在諮詢構想框架內，只允許在特定前提下，於法秩序個別領域中限制終止妊娠原則禁止之法律效果，若保護構想為其有效性而要求此種例外；依此，尤其在刑法中，依諮詢構想所實施之終止妊娠得透過構成要件排除而免於刑罰威嚇。
\end{transpara}

\begin{sourcepara}[title={原文E.164}]
\sourcern{278}3. Eine verfassungskonforme Auslegung des § 218a \abbrAbs{} 1 \abbrStGB{} \abbrNF{} im Sinne eines Tatbestandsausschlusses, wie er -- jedenfalls dem Wortlaut nach -- im Gesetzentwurf der Abgeordneten Wettig-Danielmeier, Würfel und weiterer Abgeordneter (\abbrBTDrucks{} 12/2605 [neu]) zunächst als § 218 \abbrAbs{} 5 \abbrStGB{} vorgesehen war, ist nicht möglich. Der Wortlaut, der Normzusammenhang und das vom Gesetzgeber verfolgte Ziel stehen dem entgegen. Der Gesetzgeber hat die Regelung nicht auf die bloße Zurücknahme eines strafrechtlichen Verbots beschränken wollen. Er hat vielmehr die über das Strafrecht hinausreichende Wirkung -- insbesondere im Blick auf § 24b \abbrSGBV{} -- gewollt, wie sich aus den Äußerungen von an der Formulierung des § 218a \abbrAbs{} 1 \abbrStGB{} \abbrNF{} unmittelbar beteiligten Bundestagsabgeordneten in der mündlichen Verhandlung ergibt. Die Fassung des § 218a \abbrAbs{} 1 \abbrStGB{} \abbrNF{} nimmt deshalb entsprechende Formulierungen in außerstrafrechtlichen Vorschriften auf, die rechtliche Vergünstigungen in Fällen eines "nicht rechtswidrigen Abbruchs der Schwangerschaft" vorsehen (vgl. § 616 \abbrAbs{} 2 Satz 3 BGB, § 1 \abbrAbs{} 2 Lohnfortzahlungsgesetz [LFZG], § 133c Satz 4 Gewerbeordnung [GewO], § 63 \abbrAbs{} 1 Satz 2 HGB, § 12 \abbrAbs{} 1 Satz 1 \abbrNr{} 2 Buchst. b Satz 2 Berufsbildungsgesetz [BBiG], § 24b \abbrSGBV{}, § 37a BSHG). Die im beschlossenen Gesetz enthaltenen Worte "nicht rechtswidrig" können danach -- nicht anders als die Indikationstatbestände des § 218a \abbrAbs{} 2 und 3 \abbrStGB{} \abbrNF{} -- nur als allgemeiner Rechtfertigungsgrund ausgelegt werden.
\end{sourcepara}
\begin{transpara}
\transrn{278}3. 將新版本《刑法》第218a條第1項合憲解釋為構成要件排除並不可能；此種排除至少依文字，原先曾在 Wettig-Danielmeier、Würfel 及其他議員之法律草案（BTDrucks. 12/2605 [neu]）中作為《刑法》第218條第5項規定。文字、規範脈絡及立法者所追求之目標均與此相反。立法者並不想將規制限於單純撤回刑法禁止。其毋寧有意使其產生超越刑法之效果，尤其著眼於 SGB V 第24b條；此從直接參與新版本《刑法》第218a條第1項制定之聯邦議會議員於言詞辯論中之發言可知。因此，新版本《刑法》第218a條第1項之版本採納刑法外規定中相應表述；此等規定就「不具違法性之終止妊娠」案件規定法律優惠（參見《民法》第616條第2項第3句、Lohnfortzahlungsgesetz [LFZG] 第1條第2項、Gewerbeordnung [GewO] 第133c條第4句、HGB 第63條第1項第2句、Berufsbildungsgesetz [BBiG] 第12條第1項第1句第2款 b 字第2句、SGB V 第24b條、BSHG 第37a條）。依此，已決議法律中所含「不具違法性」等語，只能與新版本《刑法》第218a條第2項及第3項之指標事由構成要件一樣，被解釋為一般正當化事由。
\end{transpara}

\begin{sourcepara}[title={原文E.165}]
\sourcern{279}4. § 218a \abbrAbs{} 1 \abbrStGB{} \abbrNF{} ist darüber hinaus auch deshalb mit \abbrArt{} 1 \abbrAbs{} 1 in Verbindung mit \abbrArt{} 2 \abbrAbs{} 2 Satz 1 \abbrGG{}\ggfnArtTwoTwoOneDE{} unvereinbar und nichtig, weil die Vorschrift eine Beratung voraussetzt, deren Regelung in § 219 \abbrStGB{} \abbrNF{} ihrerseits nicht den verfassungsrechtlichen Anforderungen genügt (siehe dazu unten II.).
\end{sourcepara}
\begin{transpara}
\transrn{279}4. 新版本《刑法》第218a條第1項此外亦因下列理由與《基本法》第1條第1項結合第2條第2項第1句\ggfnArtOneOneTwoTwoOneZH{}不相容並且無效：該規定以諮詢為前提，而其在新版本《刑法》第219條中之規制本身並不符合憲法要求（對此見下文 II.）。
\end{transpara}

\bisubsubsection{II.}{II.}

\begin{sourcepara}[title={原文E.166}]
\sourcern{280}§ 219 \abbrStGB{} \abbrNF{} genügt in seiner derzeitigen Fassung den verfassungsrechtlichen Anforderungen an den Schutz des ungeborenen menschlichen Lebens (\abbrArt{} 1 \abbrAbs{} 1, \abbrArt{} 2 \abbrAbs{} 2 Satz 1 \abbrGG{}\ggfnArtTwoTwoOneDE{}) nicht; er ist daher mit dem Grundgesetz unvereinbar und nichtig.
\end{sourcepara}
\begin{transpara}
\transrn{280}新版本《刑法》第219條目前版本，不符合保護未出生人類生命之憲法要求（《基本法》第1條第1項\ggfnArtOneOneZH{}、第2條第2項第1句\ggfnArtTwoTwoOneZH{}）；因此，其與《基本法》不相容並且無效。
\end{transpara}

\begin{sourcepara}[title={原文E.167}]
\sourcern{281}Die Beratung hat innerhalb des dem Gesetz zugrundeliegenden Schutzkonzepts eine zentrale Bedeutung. Der Gesetzgeber muß diese Beratung als staatliche Aufgabe so gestalten, daß der Staat seine Verantwortung für ihre Durchführung voll wahrnehmen kann (vgl. oben D.IV.3.). Er muß dabei Ziel und Inhalt der Beratungsaufgabe sowie das zu beachtende Verfahren in Übereinstimmung mit den aus der Verfassung entwickelten Grundsätzen regeln (vgl. oben D.IV.1. und 2.). Bei dem Rang des hier zu schützenden Rechtsguts und dem hohen Grad seiner Gefährdung gerade zu dem Zeitpunkt, zu dem die schwangere Frau einen Abbruch ihrer Schwangerschaft erwägt und deshalb eine Beratungsstelle aufsucht, muß der Gesetzgeber Aufgabe und Durchführung der Beratung so eindeutig und allgemein verständlich regeln, daß das Gesetz auch ohne Zuhilfenahme von Erläuterungen angewandt werden kann.
\end{sourcepara}
\begin{transpara}
\transrn{281}諮詢在法律所依據之保護構想中具有核心意義。立法者必須將此一諮詢作為國家任務加以形塑，使國家能完全履行其對諮詢實施之責任（參見上文 D.IV.3.）。立法者在此必須依與憲法所發展原則一致之方式，規範諮詢任務之目標與內容，以及應遵守之程序（參見上文 D.IV.1. 及 2.）。鑑於此處所欲保護法益之位階，以及其特別在懷孕婦女考慮終止妊娠並因此尋求諮詢機構之時所具有之高度危險，立法者必須如此明確且一般可理解地規範諮詢任務與實施，使法律即使不借助說明亦可被適用。
\end{transpara}

\begin{sourcepara}[title={原文E.168}]
\sourcern{282}1. Die Regelung der Beratung der Schwangeren in einer Not- und Konfliktlage (§ 219 \abbrStGB{} \abbrNF{}) ist schon deshalb verfassungsrechtlich nicht zureichend, weil die nach dem Schutzkonzept erforderliche, auf den Lebensschutz ausgerichtete Beratung nicht durch ausreichende staatliche Befugnisse und Pflichten zur Organisation und Beaufsichtigung der Beratungseinrichtungen sichergestellt ist; der Staat kann auf der Grundlage dieser Regelung nicht seiner Verantwortung für ein Angebot solcher Beratungseinrichtungen genügen, die nach Ziel und Inhalt der Beratung eine dem Schutzkonzept entsprechende Wirksamkeit erwarten lassen.
\end{sourcepara}
\begin{transpara}
\transrn{282}1. 困境與衝突狀態中懷孕婦女諮詢之規制（新版本《刑法》第219條），在憲法上已因下列理由而不足：保護構想所要求且以生命保護為取向之諮詢，並未透過足夠之國家權限及義務以組織並監督諮詢機構而獲確保；國家不能基於此一規制，履行其對提供此類諮詢機構所負之責任，而依諮詢目標與內容，此等機構應可期待具有符合保護構想之有效性。
\end{transpara}

\begin{sourcepara}[title={原文E.169}]
\sourcern{283}a) § 219 \abbrAbs{} 2 \abbrStGB{} \abbrNF{} bestimmt in organisatorischer Hinsicht, daß die Beratung durch eine aufgrund Gesetzes anerkannte Beratungsstelle zu erfolgen habe und ein Arzt, der einen Schwangerschaftsabbruch vornimmt, als Berater ausgeschlossen sei.
\end{sourcepara}
\begin{transpara}
\transrn{283}a) 新版本《刑法》第219條第2項在組織方面規定，諮詢應由依法承認之諮詢機構實施，且實施終止妊娠之醫師不得擔任諮詢人員。
\end{transpara}

\begin{sourcepara}[title={原文E.170}]
\sourcern{284}aa) Damit ist nicht gesichert, daß der Staat nur solchen Beratungseinrichtungen die Aufgabe der Beratung Schwangerer in einer Not- und Konfliktlage anvertraut, die nach ihrer Organisation, nach ihrer Grundeinstellung zum Schutz des ungeborenen Lebens und nach dem bei ihnen tätigen Personal die Gewähr dafür bieten, daß die Beratung im Sinne der verfassungsrechtlichen und gesetzlichen Vorgaben erfolgt (oben D.IV.3.a). Nicht gewährleistet ist auch, daß nur Beratungsstellen mit hinreichender personeller Ausstattung anerkannt werden, bei denen Beratungsgespräche nicht unter Zeitdruck stehen (vgl. oben D.IV.3.b). Schließlich fehlt es an ausreichenden normativen Vorkehrungen dafür, daß nicht Einrichtungen als Beratungsstellen anerkannt werden, die mit Einrichtungen, in denen Schwangerschaftsabbrüche vorgenommen werden, derart organisatorisch oder durch wirtschaftliche Interessen verbunden sind, daß hiernach ein materielles Interesse der Beratungseinrichtung an der Durchführung von Schwangerschaftsabbrüchen nicht auszuschließen ist (vgl. oben D.IV.3.c).
\end{sourcepara}
\begin{transpara}
\transrn{284}aa) 由此並未確保，國家只將困境與衝突狀態中懷孕婦女諮詢之任務，委託予依其組織、依其對未出生生命保護之基本態度，以及依其所任用人員，能保證諮詢依憲法及法律預設之意旨進行之諮詢機構（上文 D.IV.3.a）。亦未保障只有人員配置足夠、諮詢談話不處於時間壓力下之諮詢機構會被承認（參見上文 D.IV.3.b）。最後，亦欠缺足夠之規範預防安排，以避免與實施終止妊娠之機構在組織上或經濟利益上如此連結，以致不能排除諮詢機構對實施終止妊娠具有物質利益之機構，被承認為諮詢機構（參見上文 D.IV.3.c）。
\end{transpara}

\begin{sourcepara}[title={原文E.171}]
\sourcern{285}Die Lücken lassen sich auch nicht durch Rückgriff auf das Gesetz über Aufklärung, Verhütung, Familienplanung und Beratung (\abbrArt{} 1 \abbrSFHG{}) schließen. § 3 \abbrAbs{} 3 dieses Gesetzes sieht als materielle Voraussetzungen der Anerkennung nur vor, daß die Stelle über hinreichend qualifiziertes Personal verfügt, bestimmte Fachkräfte erforderlichenfalls heranziehen kann, mit Stellen zusammenarbeitet, die Hilfe für Mutter und Kind gewähren und zu der in § 2 des Gesetzes vorgesehenen -- vor allem der Information dienenden -- Beratung in der Lage ist. Damit bleiben die materiellen Voraussetzungen, die das Gesetz für die Anerkennung von allgemeinen Beratungsstellen verlangt, jedenfalls für die Gewährleistung einer auf den Lebensschutz ausgerichteten Beratung, die auf die Aufarbeitung des Konflikts im Sinne einer Ermutigung zur Fortsetzung der Schwangerschaft angelegt ist, deutlich hinter den Anforderungen zurück, die von Verfassungs wegen an Einrichtungen der Beratung im Sinne des § 219 \abbrStGB{} \abbrNF{} zu stellen sind. Auch aus § 4 des Gesetzes läßt sich in dieser Hinsicht nichts entnehmen; die Vorschrift enthält allerdings einen kapazitätsbezogenen Sicherstellungsauftrag, regelt aber keine Anerkennungsvoraussetzung für die einzelne Beratungsstelle.
\end{sourcepara}
\begin{transpara}
\transrn{285}此等漏洞亦不能透過援引《說明、避孕、家庭計畫與諮詢法》（SFHG 第1條）而填補。該法第3條第3項作為承認之實體前提，僅規定該機構須具備足夠合格人員、必要時得引入特定專業人員、與提供母親及子女協助之機構合作，並能提供該法第2條所規定，主要服務於資訊提供之諮詢。由此，法律對一般諮詢機構承認所要求之實體前提，至少就保障一項以生命保護為取向、旨在以鼓勵繼續懷孕之方式處理衝突之諮詢而言，明顯低於憲法上對新版本《刑法》第219條意義下諮詢機構所應提出之要求。從該法第4條亦無法就此得出任何內容；該規定固然包含一項與容量相關之確保任務，但並未規範個別諮詢機構之承認前提。
\end{transpara}

\begin{sourcepara}[title={原文E.172}]
\sourcern{286}bb) Lückenhaft ist die gesetzliche Regelung der Organisation der Konfliktberatung auch deshalb, weil § 219 \abbrAbs{} 2 Satz 1 \abbrStGB{} \abbrNF{} nicht zum Ausdruck bringt, daß die dort geforderte Anerkennung "auf Grund Gesetzes" gerade auf die hier gebotene Art der Beratung bezogen sein muß: Der Staat ist es, der jenen Stellen, die die Voraussetzungen für diese Beratung erfüllen, die Beratung als Aufgabe durch staatliche Anerkennung anvertraut, und ihnen durch Rücknahme der Anerkennung die Aufgabe entziehen muß, wenn sie die Voraussetzungen hierfür nicht mehr erfüllen (vgl. oben D.IV.3.d). Auch in dem Gesetz über Aufklärung, Verhütung, Familienplanung und Beratung (\abbrArt{} 1 \abbrSFHG{}) findet sich eine solche Regelung nicht. Die Anerkennung nach § 3 \abbrAbs{} 2 des Gesetzes wird von einer Behörde oder Körperschaft, Anstalt oder Stiftung des öffentlichen Rechts erteilt. Das Gesetz knüpft hiermit an die Bestimmung in § 218b \abbrAbs{} 2 \abbrStGB{} \abbrAF{} an. Wie bisher geht es von einem pluralistischen Angebot anerkannter Beratungsstellen aus (vgl. \abbrBTDrucks{} 12/551 und 12/2605 [neu], jeweils Begründung zu \abbrArt{} 1 § 3); die nähere Regelung der Anerkennung wird damit wie bisher dem Landesrecht überlassen (vgl. zum bisherigen Recht Dreher/Tröndle, 46. \abbrAufl{} \abbrRdnr{} 5 zu § 218b \abbrStGB{} \abbrAF{}). In einzelnen Ländern ist unter der Geltung des § 218b \abbrStGB{} \abbrAF{} auch bestimmt worden, daß Körperschaften, Anstalten oder Stiftungen des öffentlichen Rechts, zu deren gesetzlichen oder satzungsmäßigen Aufgaben die Beratung von Schwangeren gehört, eigene Einrichtungen oder solche ihnen zugehöriger Träger als Beratungsstellen anerkennen können (§ 2 \abbrAbs{} 3 des Hessischen Gesetzes zur Ausführung der §§ 218b und 219 des Strafgesetzbuches und des \abbrArt{} 3 des Fünften Gesetzes zur Reform des Strafrechts vom 2. Mai 1978, \abbrGVBl{} \abbrS{} 273; § 1 \abbrAbs{} 2 des entsprechenden Berliner Gesetzes vom 22. Dezember 1978, \abbrGVBl{} \abbrS{} 2514). In Nordrhein-Westfalen ist vorgesehen, daß für die Anerkennung von in kirchlicher Trägerschaft stehenden Beratungsstellen die Kirchen zuständig sind, wenn sie den Status einer Körperschaft des öffentlichen Rechts haben (§ 1 der Verordnung über Zuständigkeiten bei Schwangerschaftsberatung und -abbruch vom 12. Dezember 1978, \abbrGVBl{} \abbrS{} 632).
\end{sourcepara}
\begin{transpara}
\transrn{286}bb) 關於衝突諮詢組織之法律規制亦有漏洞，因為新版本《刑法》第219條第2項第1句並未表達，該處所要求「依法」承認，必須正是關聯於此處所要求之諮詢種類：由國家透過國家承認，將諮詢作為任務委託予符合此種諮詢前提之機構；若其不再符合此前提，國家即必須透過撤回承認而剝奪其任務（參見上文 D.IV.3.d）。在《說明、避孕、家庭計畫與諮詢法》（SFHG 第1條）中亦無此種規制。依該法第3條第2項之承認，由一機關或公法上社團、機構或基金會作成。法律在此連結於舊版本《刑法》第218b條第2項之規定。其如同過去，以受承認諮詢機構之多元提供為出發點（參見 BTDrucks. 12/551 及 12/2605 [neu]，各該第1條第3項理由）；承認之進一步規制因而如同以往交由邦法處理（關於既有法，參見 Dreher/Tröndle, 46. Aufl. Rdnr. 5 zu § 218b StGB a.F.）。在若干邦，於舊版本《刑法》第218b條施行下亦曾規定，其法律或章程任務包括懷孕婦女諮詢之公法上社團、機構或基金會，得將自身機構或其所屬承辦者之機構承認為諮詢機構（1978年5月2日黑森邦執行《刑法》第218b條、第219條及《第五刑法改革法》第3條之法律第2條第3項，GVBl. S. 273；1978年12月22日相應柏林法律第1條第2項，GVBl. S. 2514）。在北萊茵-西伐利亞，則規定由教會承辦之諮詢機構，若教會具有公法社團地位，其承認由教會主管（1978年12月12日《懷孕諮詢及終止妊娠管轄權命令》第1條，GVBl. S. 632）。
\end{transpara}

\begin{sourcepara}[title={原文E.173}]
\sourcern{287}Für die in \abbrArt{} 1 \abbrSFHG{} geregelte Beratung mag dies nicht zu beanstanden sein. Die Organisation der im Rahmen der Beratungsregelung gebotenen Beratung ist jedoch wesentlicher Teil des vom Bundesgesetzgeber auf dem Boden seiner Zuständigkeit für das Strafrecht (\abbrArt{} 74 \abbrNr{} 1 \abbrGG{}\ggfnArtSeventyFourOneDE{}) zu entwickelnden, insoweit gegebenenfalls unter Zuhilfenahme der Kompetenz aus \abbrArt{} 84 \abbrAbs{} 1 \abbrGG{}\ggfnArtEightyFourOneDE{} normativ zu gestaltenden und insgesamt am Untermaßverbot zu messenden Schutzkonzepts. Würde es der Bundesgesetzgeber den Ländern überlassen, die organisationsrechtlichen Bestimmungen des Schutzkonzepts zu treffen, so müßte er das Inkrafttreten der Gesamtregelung davon abhängig machen, daß in allen Ländern die erforderlichen gesetzlichen Vorschriften ergangen sind. Diesen letztgenannten Weg hat der Gesetzgeber ersichtlich nicht beschritten.
\end{sourcepara}
\begin{transpara}
\transrn{287}對 SFHG 第1條所規定之諮詢而言，這或許無可指摘。然而，諮詢規制框架內所要求諮詢之組織，是聯邦立法者基於其刑法權限（《基本法》第74條第1款\ggfnArtSeventyFourOneZH{}）所應發展之保護構想的重要部分；就此必要時亦須借助《基本法》第84條第1項\ggfnArtEightyFourOneZH{}之權限作規範形塑，且整體須以不足禁止加以衡量。若聯邦立法者將保護構想之組織法規定交由各邦制定，則其必須使整體規制之生效取決於所有邦均已制定必要法律規定。立法者顯然並未採取最後所述道路。
\end{transpara}

\begin{sourcepara}[title={原文E.174}]
\sourcern{288}b) Es fehlt auch an Bestimmungen, die eine ausreichende staatliche Beaufsichtigung der Beratungsstellen gewährleisten. Der Staat muß auf gesetzlicher Grundlage die Anerkennung dieser Beratungsstellen regelmäßig und in nicht zu langen Zeitabständen überprüfen und sich dabei vergewissern, ob die Anforderungen an die Beratung beachtet werden; nur unter diesen Voraussetzungen darf die Anerkennung weiter fortbestehen oder neu bestätigt werden (vgl. D.IV.3.d). Eine solche Überwachung setzt voraus, daß das Gesetz auch Informations- und Kontrollbefugnisse zur Verfügung stellt (vgl. oben D.IV.3.e). Da der Gesetzgeber hierüber nichts bestimmt hat, obwohl dies im Rahmen des Beratungskonzepts nicht ungeregelt bleiben durfte, ist die getroffene Beratungsregelung mangelhaft.
\end{sourcepara}
\begin{transpara}
\transrn{288}b) 亦欠缺確保國家對諮詢機構充分監督之規定。國家必須在法律基礎上，定期且以不過長之間隔審查此等諮詢機構之承認，並在此確認諮詢要求是否受到遵守；只有在此等前提下，承認始得繼續存在或重新確認（參見 D.IV.3.d）。此種監督預設法律亦提供資訊及控制權限（參見上文 D.IV.3.e）。由於立法者對此未作任何規定，儘管此在諮詢構想框架內不得不受規制，故所作諮詢規制有瑕疵。
\end{transpara}

\begin{sourcepara}[title={原文E.175}]
\sourcern{289}c) Die dargelegten Mängel der organisatorischen und verfahrensrechtlichen Einbindung der Beratungsstellen in die dem Staat obliegende Verantwortung für die verfassungsrechtlich gebotene Wirksamkeit des Schutzkonzepts ergreifen die in § 219 \abbrStGB{} \abbrNF{} getroffene Beratungsregelung insgesamt. Aussagen über Ziel und Inhalt der Beratung ermangeln ohne die organisatorischen und aufsichtlichen Vorkehrungen, die ihre Umsetzung sichern, der vom Schutzkonzept gebotenen Wirksamkeit; sie hängen sozusagen in der Luft.
\end{sourcepara}
\begin{transpara}
\transrn{289}c) 前述諮詢機構在組織法及程序法上納入國家對保護構想憲法所要求有效性之責任中的缺陷，影響新版本《刑法》第219條所作諮詢規制整體。關於諮詢目標與內容之陳述，若無確保其實現之組織及監督安排，即欠缺保護構想所要求之有效性；它們可說是懸在空中。
\end{transpara}

\begin{sourcepara}[title={原文E.176}]
\sourcern{290}2. Vor diesem Hintergrund erübrigt sich eine abschließende Entscheidung darüber, ob die Bestimmungen des § 219 \abbrStGB{} \abbrNF{} über Ziel, Inhalt und Durchführung der Beratung verfassungsrechtlicher Prüfung standhalten. Indessen ist diese Vorschrift bei der ohnehin erforderlichen neuen gesetzlichen Regelung jedenfalls so zu fassen, daß sie aus sich heraus eindeutig und allgemein verständlich ist und daher auch ohne Zuhilfenahme von Erläuterungen so angewandt werden kann, daß allen unter D.IV.1. und 2. dargelegten verfassungsrechtlichen Anforderungen aus der staatlichen Schutzpflicht in der Praxis entsprochen wird.
\end{sourcepara}
\begin{transpara}
\transrn{290}2. 在此背景下，已無須就新版本《刑法》第219條關於諮詢目標、內容及實施之規定是否能經受憲法審查作成終局決定。不過，在無論如何均必要之新法律規制中，該規定至少必須如此制定，使其本身明確且一般可理解，因而即使不借助說明亦能如此適用，以致實務上符合 D.IV.1. 及 2. 所述、由國家保護義務所生之一切憲法要求。
\end{transpara}

\begin{sourcepara}[title={原文E.177}]
\sourcern{291}a) Nach § 219 \abbrAbs{} 1 Sätze 1 bis 3 \abbrStGB{} \abbrNF{} dient die Beratung dem Lebensschutz durch Rat und Hilfe für die Schwangere unter Anerkennung des hohen Wertes des vorgeburtlichen Lebens und der Eigenverantwortlichkeit der Frau; sie soll dazu beitragen, die im Zusammenhang mit der Schwangerschaft bestehende Not- und Konfliktlage zu bewältigen und die Schwangere in die Lage versetzen, eine verantwortungsbewußte eigene Gewissensentscheidung zu treffen.
\end{sourcepara}
\begin{transpara}
\transrn{291}a) 依新版本《刑法》第219條第1項第1句至第3句，諮詢在承認出生前生命之高度價值及婦女自我責任下，透過對懷孕婦女之建議與協助服務於生命保護；其應有助於克服與懷孕相關之困境與衝突狀態，並使懷孕婦女能作成負責任之自身良心決定。
\end{transpara}

\begin{sourcepara}[title={原文E.178}]
\sourcern{292}Damit kommen das verfassungsrechtlich vorgegebene Ziel der Beratung und deren wesentlicher Inhalt (vgl. oben D.IV.1.) nicht hinreichend deutlich zum Ausdruck. Die Beratung ist zwar ergebnisoffen zu führen, weil dann am ehesten erwartet werden kann, daß sich die Frau an der Suche nach einer Konfliktlösung beteiligt. Doch darf die Beratung nicht ergebnis- und zieloffen, sondern muß auf den Schutz des ungeborenen Lebens hin orientiert sein. Die Beraterinnen und Berater haben sich von dem Bemühen leiten zu lassen, die Schwangere zur Fortsetzung ihrer Schwangerschaft zu ermutigen und ihr Perspektiven für ein Leben mit dem Kind zu eröffnen. Dazu müssen sie auch vorhandenen Fehlvorstellungen über den grundsätzlichen Vorrang des eigenen Lebensrechts des Ungeborenen und über das notwendige Gewicht der einen Schwangerschaftsabbruch ausnahmsweise erlaubenden Ausnahmelagen entgegentreten.
\end{sourcepara}
\begin{transpara}
\transrn{292}由此，憲法所預設之諮詢目標及其主要內容（參見上文 D.IV.1.）並未足夠清楚地表達出來。諮詢固然應以結果開放方式進行，因為如此最可期待婦女參與尋求衝突解決。但諮詢不得對結果與目標均開放，而必須以保護未出生生命為取向。諮詢人員必須受下列努力所引導：鼓勵懷孕婦女繼續其懷孕，並為其開啟與子女共同生活之前景。為此，其亦必須對抗關於未出生者自身生命權之原則優先，以及關於例外允許終止妊娠之例外狀態所需重量的既有錯誤觀念。
\end{transpara}

\begin{sourcepara}[title={原文E.179}]
\sourcern{293}Die Begründung des Gesetzentwurfs verweist zwar auf den Schutz des ungeborenen menschlichen Lebens als Ziel der Beratung (vgl. \abbrBTDrucks{} 12/2605 [neu], \abbrS{} 22 [Einzelbegründung zu § 219]). Der für die Anwendung der Vorschrift maßgebliche objektive Wortsinn läßt dies aber nicht mit der gebotenen Klarheit und Eindeutigkeit erkennen. § 219 \abbrAbs{} 1 \abbrStGB{} \abbrNF{} enthält auch keinen ausdrücklichen Auftrag, die schwangere Frau zum Austragen des Kindes zu ermutigen. Nicht hinreichend deutlich wird im Gesetz schließlich, daß der Frau bewußt sein muß, dem Schutz des ungeborenen Lebens gebühre grundsätzlich der Vorrang, sofern nicht eine schwerwiegende Notlage im Sinne der eingangs dargelegten verfassungsrechtlichen Maßstäbe gegeben ist.
\end{sourcepara}
\begin{transpara}
\transrn{293}法律草案理由固然指向未出生人類生命之保護作為諮詢目標（參見 BTDrucks. 12/2605 [neu], S. 22 [關於第219條之個別理由]）。但對適用該規定具有決定性之客觀文字意義，並未以所要求之清楚與明確顯示此點。新版本《刑法》第219條第1項亦未包含明示任務，要求鼓勵懷孕婦女繼續懷孕。最後，法律亦未足夠清楚顯示，婦女必須意識到：原則上應賦予未出生生命之保護以優先，除非存在前述憲法標準意義下之重大困境。
\end{transpara}

\begin{sourcepara}[title={原文E.180}]
\sourcern{294}Das Beratungsziel muß deshalb in einer Neufassung des § 219 \abbrStGB{} deutlicheren Niederschlag finden.
\end{sourcepara}
\begin{transpara}
\transrn{294}因此，諮詢目標必須在《刑法》第219條之新版本中更清楚地表現出來。
\end{transpara}

\begin{sourcepara}[title={原文E.181}]
\sourcern{295}b) In § 219 \abbrAbs{} 1 Sätze 4 und 5 \abbrStGB{} \abbrNF{} wird als Aufgabe der Beratung die umfassende medizinische, soziale und juristische Information der Schwangeren genannt; die Beratung umfaßt danach die Darlegung der Rechtsansprüche von Mutter und Kind und der möglichen praktischen Hilfen. Gemäß § 219 \abbrAbs{} 3 Satz 2 \abbrStGB{} \abbrNF{} ist über die Durchführung einer Beratung gemäß Absatz 1 und der der Frau "damit" gegebenen Informationen für ihre Entscheidungsfindung "sofort" eine Bescheinigung auszustellen.
\end{sourcepara}
\begin{transpara}
\transrn{295}b) 新版本《刑法》第219條第1項第4句及第5句將懷孕婦女之全面醫學、社會及法律資訊列為諮詢任務；依此，諮詢包括說明母親及子女之法律請求權及可能實際協助。依新版本《刑法》第219條第3項第2句，關於依第1項實施之諮詢，以及「藉此」給予婦女供其決定形成之資訊，應「立即」出具證明。
\end{transpara}

\begin{sourcepara}[title={原文E.182}]
\sourcern{296}Hiermit wird -- unbeschadet der Tatsache, daß die Gesetzesinitiatoren mehr wollten als bloße Information (vgl. Protokoll der 17. Sitzung des Sonderausschusses "Schutz des ungeborenen Lebens" vom 17. Juni 1992, \abbrS{} 12) -- der Eindruck erweckt, als liege der Schwerpunkt der Beratung nach § 219 \abbrAbs{} 1 \abbrStGB{} \abbrNF{} in der Information. Dem muß der Gesetzgeber entgegentreten, weil dies den verfassungsrechtlichen Anforderungen an die Beratung nach § 219 \abbrStGB{} \abbrNF{} nicht entspräche. Eine bloß informierende Beratung, die den konkreten Schwangerschaftskonflikt nicht aufnimmt und zum Thema eines Gesprächs zu machen sucht, müßte nämlich die der Beratung im Rahmen des Schutzkonzepts zukommende Funktion, durch Ermutigung der Frau wesentlich zum Lebensschutz beizutragen, verfehlen. Für die Aufnahme einer Konfliktberatung ist es zudem erforderlich, daß die Beraterin oder der Berater von der Frau die Mitteilung der Gründe zu erreichen sucht, die sie den Schwangerschaftsabbruch erwägen lassen. Solange der Beraterin oder dem Berater die Möglichkeiten einer Konfliktlösung -- gegebenenfalls unter Einbeziehung dritter Personen -- nicht ausgeschöpft erscheinen, darf die Beratungsbescheinigung, die den Abschluß der Beratung dokumentiert, nicht ausgestellt werden.
\end{sourcepara}
\begin{transpara}
\transrn{296}由此，儘管法律倡議者所欲者多於單純資訊提供（參見1992年6月17日「未出生生命保護」特別委員會第17次會議紀錄 S. 12），仍造成一種印象，彷彿新版本《刑法》第219條第1項所規定諮詢之重點在於資訊提供。立法者必須對抗此點，因為這不符合新版本《刑法》第219條所規定諮詢之憲法要求。單純提供資訊之諮詢若不承接具體懷孕衝突，並試圖將其作為談話主題，即必然錯失諮詢在保護構想框架中所具有之功能，亦即透過鼓勵婦女而對生命保護作出實質貢獻。此外，為展開衝突諮詢，諮詢人員必須嘗試由婦女取得其考慮終止妊娠之理由。只要諮詢人員認為衝突解決之可能性，必要時包括納入第三人，尚未窮盡，證明諮詢結束之諮詢證明即不得出具。
\end{transpara}

\begin{sourcepara}[title={原文E.183}]
\sourcern{297}In einer Neufassung des Gesetzes ist deshalb klar zum Ausdruck zu bringen, daß die Beratung über Information hinaus eine dem Lebensschutz dienende Konfliktberatung jedenfalls anzustreben und die dafür gegebenen Möglichkeiten auszuschöpfen hat
\end{sourcepara}
\begin{transpara}
\transrn{297}因此，在法律新版本中必須清楚表達：諮詢超越資訊提供之外，無論如何必須致力於服務生命保護之衝突諮詢，並窮盡為此所存在之可能性。
\end{transpara}

\begin{sourcepara}[title={原文E.184}]
\sourcern{298}c) Unabhängig von der Regelung des \abbrArt{} l § 2 \abbrAbs{} 2 Satz 2 \abbrSFHG{} bedarf es ferner der Verdeutlichung, daß die Beratungsstellen im Rahmen der Beratung nach § 219 \abbrStGB{} \abbrNF{} die Frau über öffentliche Hilfen nicht nur zu informieren, sondern ihr auch verfügbare Hilfsangebote zu unterbreiten oder sie bei der Erlangung von Hilfen zu unterstützen haben.
\end{sourcepara}
\begin{transpara}
\transrn{298}c) 不論 SFHG 第1條第2項第2句之規制如何，此外亦須明確說明，諮詢機構在新版本《刑法》第219條之諮詢框架中，不僅須告知婦女公共協助，亦須向其提出可取得之協助提供，或支持其取得協助。
\end{transpara}

\begin{sourcepara}[title={原文E.185}]
\sourcern{299}d) § 219 \abbrAbs{} 1 Satz 3 \abbrStGB{} \abbrNF{} bezeichnet die von der Frau nach der Beratung getroffene Entscheidung als "Gewissensentscheidung". Der Gesetzgeber will damit offenbar an eine Formulierung des Bundesverfassungsgerichts (vgl. \abbrBVerfGE{} 39, 1 [48]) anknüpfen, wonach die Entscheidung zum Abbruch der Schwangerschaft den Rang einer achtenswerten Gewissensentscheidung haben kann. Indes kann die Frau, die sich nach Beratung zum Abbruch entschließt, für die damit einhergehende Tötung des Ungeborenen nicht etwa eine grundrechtlich in \abbrArt{} 4 \abbrAbs{} 1 \abbrGG{}\ggfnArtFourOneDE{} geschützte Rechtsposition in Anspruch nehmen. Verfassungsrechtlich zulässig kann das Gesetz nur eine gewissenhaft zustandegekommene und in diesem Sinne achtenswerte Entscheidung meinen. Dies wird in geeigneter Weise klarzustellen sein.
\end{sourcepara}
\begin{transpara}
\transrn{299}d) 新版本《刑法》第219條第1項第3句將婦女經諮詢後所作決定稱為「良心決定」。立法者顯然意在連結聯邦憲法法院之一項表述（參見 BVerfGE 39, 1 [48]），依其所述，終止妊娠之決定得具有值得尊重之良心決定地位。然而，經諮詢後決定終止妊娠之婦女，不能就與此相伴之殺害未出生者主張受《基本法》第4條第1項\ggfnArtFourOneZH{}基本權保護之法律地位。從憲法上可允許者，只能是法律意指一項良心謹慎形成、並在此意義上值得尊重之決定。此點應以適當方式予以明確。
\end{transpara}

\bisubsubsection{III.}{III.}

\begin{sourcepara}[title={原文E.186}]
\sourcern{300}Unabhängig von den zu § 218a \abbrAbs{} 1 und § 219 \abbrStGB{} \abbrNF{} festgestellten verfassungsrechtlichen Mängeln genügt die Regelung des Schwangeren- und Familienhilfegesetzes über Schwangerschaftsabbrüche nach dem Beratungskonzept auch deshalb nicht der Schutzpflicht für das ungeborene menschliche Leben, weil der Gesetzgeber in dem oben näher bezeichneten Umfang die besonderen Pflichten des Arztes, von dem die Frau den Schwangerschaftsabbruch verlangt (vgl. dazu D.V.1. und 2.), sowie der Personen des familiären Umfeldes der schwangeren Frau (vgl. dazu D.VI.2.) nicht festgelegt und bestimmte Pflichtverletzungen nicht mit Strafe bedroht hat. Die fehlenden Bestimmungen hat der Gesetzgeber im Zusammenhang mit der durch die Nichtigerklärung der § 218a \abbrAbs{} 1, 219 \abbrStGB{} \abbrNF{} gebotenen Neuregelung nachzuholen.
\end{sourcepara}
\begin{transpara}
\transrn{300}不論就新版本《刑法》第218a條第1項及第219條所確認之憲法瑕疵如何，《孕婦及家庭扶助法》關於依諮詢構想終止妊娠之規制，亦因下列理由不符合對未出生人類生命之保護義務：立法者未在上文詳述範圍內，規定婦女要求終止妊娠之醫師的特殊義務（對此參見 D.V.1. 及 2.），以及懷孕婦女家庭環境中之人的特殊義務（對此參見 D.VI.2.），亦未以刑罰威嚇特定義務違反。欠缺之規定，立法者必須在因宣告新版本《刑法》第218a條第1項、第219條無效所要求之新規制中補足。
\end{transpara}

\bisubsubsection{IV.}{IV.}

\begin{sourcepara}[title={原文E.187}]
\sourcern{301}\abbrArt{} 15 \abbrNr{} 2 \abbrSFHG{} ist mit \abbrArt{} 1 \abbrAbs{} 1 in Verbindung mit \abbrArt{} 2 \abbrAbs{} 2 Satz 1 \abbrGG{}\ggfnArtTwoTwoOneDE{} unvereinbar und nichtig, soweit dadurch die bisher in \abbrArt{} 4 des 5. \abbrStrRG{} enthaltene Vorschrift betreffend die Bundesstatistik über Schwangerschaftsabbrüche aufgehoben wird.
\end{sourcepara}
\begin{transpara}
\transrn{301}SFHG 第15條第2款，在其廢止迄今包含於《第五刑法改革法》第4條、關於終止妊娠聯邦統計之規定範圍內，與《基本法》第1條第1項結合第2條第2項第1句\ggfnArtOneOneTwoTwoOneZH{}不相容，並且無效。
\end{transpara}

\begin{sourcepara}[title={原文E.188}]
\sourcern{302}1. a) Der Gesetzgeber erfüllt seine Pflicht, das ungeborene menschliche Leben zu schützen, nicht ein für allemal dadurch, daß er ein Gesetz zur Regelung des Schwangerschaftsabbruchs erläßt, welches diesen Schutz bezweckt und nach seiner -- verfassungsrechtlich unbedenklichen -- Einschätzung auch geeignet erscheint, das vom Grundgesetz geforderte Maß an Schutz zu gewährleisten. Er bleibt vielmehr aufgrund seiner Schutzpflicht weiterhin dafür verantwortlich, daß das Gesetz tatsächlich einen -- unter Berücksichtigung entgegenstehender Rechtsgüter -- angemessenen und als solchen wirksamen Schutz vor Schwangerschaftsabbrüchen bewirkt. Stellt sich nach hinreichender Beobachtungszeit heraus, daß das Gesetz das von der Verfassung geforderte Maß an Schutz nicht zu gewährleisten vermag, so ist der Gesetzgeber verpflichtet, durch Änderung oder Ergänzung der bestehenden Vorschriften auf die Beseitigung der Mängel und die Sicherstellung eines dem Untermaßverbot genügenden Schutzes hinzuwirken (Korrektur- oder Nachbesserungspflicht).
\end{sourcepara}
\begin{transpara}
\transrn{302}1. a) 立法者並非一旦制定一部規制終止妊娠之法律，且該法律旨在此種保護，並依其在憲法上無疑義之評估亦顯得適於保障《基本法》所要求之保護程度，即一勞永逸地履行其保護未出生人類生命之義務。立法者毋寧基於其保護義務，仍繼續對該法律是否實際產生一項在考量相反法益下適當且作為保護本身有效之防止終止妊娠保護負責。若經足夠觀察期間後顯示，法律無法保障憲法所要求之保護程度，立法者即有義務透過修改或補充既有規定，促成瑕疵之排除，並確保符合不足禁止之保護（修正或補正義務）。
\end{transpara}

\begin{sourcepara}[title={原文E.189}]
\sourcern{303}Diese Verpflichtung folgt auch daraus, daß der Gesetzgeber von Verfassungs wegen grundsätzlich gehalten ist, die Verfassungswidrigkeit eines Gesetzes sobald als möglich zu beseitigen (vgl. \abbrBVerfGE{} 15, 337 [351]). Sie ist vor allem dann von Bedeutung, wenn ein bei Erlaß verfassungsmäßiges Gesetz nachträglich verfassungswidrig wird, weil sich die tatsächlichen Verhältnisse, auf die es einwirkt, grundlegend gewandelt haben oder sich die beim Erlaß des Gesetzes verfassungsrechtlich unbedenkliche Einschätzung seiner künftigen Wirkungen später als ganz oder teilweise falsch erweist (vgl. \abbrBVerfGE{} 50, 290 [335, 352]; 56, 54 [78 \abbrF{}]; 73, 40 [94]). Die Bindung des Gesetzgebers an die verfassungsmäßige Ordnung (\abbrArt{} 20 \abbrAbs{} 3 \abbrGG{}\ggfnArtTwentyThreeDE{}) erschöpft sich nämlich nicht in der Verpflichtung, bei Erlaß eines Gesetzes die verfassungsrechtlichen Grenzen einzuhalten; sie umfaßt auch die Verantwortung dafür, daß die erlassenen Gesetze in Übereinstimmung mit dem Grundgesetz bleiben (vgl. \abbrBVerfGE{} 15, 337 [350]).
\end{sourcepara}
\begin{transpara}
\transrn{303}此一義務亦源於，立法者依憲法原則上有義務儘速排除法律之違憲性（參見 BVerfGE 15, 337 [351]）。其尤其在一部制定時合憲之法律事後因其作用之事實關係發生根本變化，或因制定法律時在憲法上無疑義之關於其未來效果的評估，日後證明為全部或部分錯誤，而變成違憲時，具有重要性（參見 BVerfGE 50, 290 [335, 352]; 56, 54 [78 f.]; 73, 40 [94]）。因為立法者受合憲秩序拘束（《基本法》第20條第3項\ggfnArtTwentyThreeZH{}），並不窮盡於制定法律時遵守憲法界限之義務；其亦包括負責確保已制定之法律持續與《基本法》一致（參見 BVerfGE 15, 337 [350]）。
\end{transpara}

\begin{sourcepara}[title={原文E.190}]
\sourcern{304}b) Die Nachbesserungspflicht schließt nicht generell eine fortlaufende Kontrolle der Gesetze durch den Gesetzgeber ein. Vielfach aktualisiert sie sich erst dann, wenn die Verfassungswidrigkeit eines Gesetzes erkannt oder doch jedenfalls deutlich erkennbar wird (vgl. \abbrBVerfGE{} 16, 130 [142]). Aus der Schutzpflicht für das Leben, die eine dauernde Verpflichtung für alle Staatsorgane darstellt, ergeben sich jedoch besondere Anforderungen (vgl. \abbrBVerfGE{} 49, 89 [130, 132]). Der hohe Rang des geschützten Rechtsgutes, die Art der Gefährdung ungeborenen Lebens und der in diesem Bereich festzustellende Wandel der gesellschaftlichen Verhältnisse und Anschauungen erfordern es, daß der Gesetzgeber beobachtet, wie sich sein gesetzliches Schutzkonzept in der gesellschaftlichen Wirklichkeit auswirkt (Beobachtungspflicht). Er muß sich in angemessenen zeitlichen Abständen in geeigneter Weise -- etwa durch periodisch zu erstattende Berichte der Regierung -- vergewissern, ob das Gesetz die erwarteten Schutzwirkungen tatsächlich entfaltet oder ob sich Mängel des Konzepts oder seiner praktischen Durchführung offenbaren, die eine Verletzung des Untermaßverbots begründen (vgl. \abbrBVerfGE{} 56, 54 [82 \abbrFF{}]). Diese Beobachtungspflicht besteht auch und gerade nach einem Wechsel des Schutzkonzepts.
\end{sourcepara}
\begin{transpara}
\transrn{304}b) 補正義務並不一般性包含立法者對法律之持續控制。其多半只有在法律違憲性被認識到，或至少明顯可認識時，才現實化（參見 BVerfGE 16, 130 [142]）。然而，由生命保護義務，作為對所有國家機關之持續義務，可得出特殊要求（參見 BVerfGE 49, 89 [130, 132]）。受保護法益之高度位階、未出生生命受危害之種類，以及此領域中可確認之社會關係與觀念變遷，要求立法者觀察其法律保護構想如何在社會現實中發生影響（觀察義務）。立法者必須以適當時間間隔及適當方式，例如透過政府定期提出報告，確認法律是否實際展開所期待之保護效果，或是否顯露構想或其實務執行之瑕疵，足以構成不足禁止之違反（參見 BVerfGE 56, 54 [82 ff.]）。此一觀察義務也且正是在保護構想轉換後存在。
\end{transpara}

\begin{sourcepara}[title={原文E.191}]
\sourcern{305}c) Die Beobachtungspflicht schließt ein, daß der Gesetzgeber im Rahmen seiner Kompetenz dafür sorgt, daß die für die Beurteilung der Wirkungen des Gesetzes notwendigen Daten planmäßig erhoben, gesammelt und ausgewertet werden. Verläßliche Statistiken mit hinreichender Aussagekraft, wie über die absolute Zahl der Schwangerschaftsabbrüche, über die relativen Quoten, die sich aus dem Verhältnis der Abbruchszahl zur Gesamtbevölkerung, zur Zahl der Frauen im gebärfähigen Alter, zur Zahl der Schwangerschaften oder der Lebend- und Totgeburten insgesamt errechnen, sowie über die Verteilung der nicht mit Strafe bedrohten Abbrüche auf die verschiedenen gesetzlichen Grundlagen, sind dazu unerläßlich. Auf welche relevanten Tatsachen der Gesetzgeber im übrigen die statistische Erhebung erstrecken will (etwa Mehrfachabbrüche, Alter der Frau, Familienstand, Kinderzahl) und wie er die Erfassung und Auswertung der Daten im einzelnen regelt, obliegt seiner Entscheidung. Mit der Schutzpflicht nicht mehr zu vereinbaren ist es aber jedenfalls, auf jede staatliche Statistik über Schwangerschaftsabbrüche zu verzichten. Der Gesetzgeber beraubt sich damit des Erfahrungsmaterials, dessen er für die Beobachtung der Auswirkungen seines Schutzkonzepts bedarf. Dagegen kann nicht eingewandt werden, die bisherigen statistischen Unterlagen hätten sich als nicht zuverlässig erwiesen. Soweit dies zutrifft, besteht Veranlassung, die Erhebungen zu verbessern. Vorschläge hierzu wurden im Anhörungsverfahren vor dem Bundestagssonderausschuß "Schutz des ungeborenen Lebens" gemacht.
\end{sourcepara}
\begin{transpara}
\transrn{305}c) 觀察義務包括，立法者在其權限框架內確保有計畫地蒐集、彙整並評估判斷法律效果所必要之資料。為此，具有足夠說明力之可靠統計不可或缺，例如關於終止妊娠之絕對數量、由終止妊娠數與總人口、育齡婦女人數、懷孕數或活產及死產總數之比例所計算之相對比率，以及不受刑罰威嚇之終止妊娠在不同法律基礎間之分布。立法者此外欲將統計調查延伸至何種相關事實（例如重複終止妊娠、婦女年齡、婚姻狀態、子女人數），以及如何具體規範資料之記錄與評估，均由其決定。但無論如何，完全放棄關於終止妊娠之任何國家統計，已不再能與保護義務相容。立法者因此剝奪自己觀察其保護構想影響所需之經驗材料。不能反駁稱，迄今統計資料已證明不可靠。若此點屬實，則有理由改善調查。在聯邦議會「未出生生命保護」特別委員會聽證程序中，已就此提出建議。
\end{transpara}

\begin{sourcepara}[title={原文E.192}]
\sourcern{306}2. Nach diesem verfassungsrechtlichen Maßstab ist die Aufhebung der bisherigen Vorschrift über die Bundesstatistik in \abbrArt{} 4 des 5. \abbrStrRG{} mit der Schutzpflicht für das ungeborene menschliche Leben nicht vereinbar.
\end{sourcepara}
\begin{transpara}
\transrn{306}2. 依此一憲法標準，廢止《第五刑法改革法》第4條關於聯邦統計之既有規定，與對未出生人類生命之保護義務不相容。
\end{transpara}

\begin{sourcepara}[title={原文E.193}]
\sourcern{307}Der mit dem Schwangeren- und Familienhilfegesetz vollzogene Wechsel des gesetzlichen Schutzkonzepts ist -- wie oben ausgeführt -- mit Unsicherheiten über die künftigen Auswirkungen der neuen Regelung verbunden. Der Konzeptwechsel stellt einen Versuch des Gesetzgebers dar, nach den als unbefriedigend angesehenen Erfahrungen mit der Indikationenregelung auf anderem Wege einen besseren Schutz des ungeborenen Lebens zu gewährleisten. Dies verpflichtet den Gesetzgeber zu sorgfältiger Beobachtung der tatsächlichen Auswirkungen des neuen Rechts und zu einer möglichst zuverlässigen Ermittlung der Daten, die für eine empirische Beurteilung der unter dem neuen Recht durchgeführten Schwangerschaftsabbrüche notwendig sind.
\end{sourcepara}
\begin{transpara}
\transrn{307}如前所述，透過《孕婦及家庭扶助法》完成之法律保護構想轉換，伴隨關於新規制未來影響之不確定性。構想轉換是立法者之一項嘗試：在其認為指標規則經驗並不令人滿意後，改以其他道路保障未出生生命之更佳保護。這課予立法者審慎觀察新法實際影響之義務，並要求其盡可能可靠地調查對經驗判斷新法下所實施終止妊娠必要之資料。
\end{transpara}

\begin{sourcepara}[title={原文E.194}]
\sourcern{308}\abbrArt{} 4 des 5. \abbrStrRG{} sieht eine Bundesstatistik über die unter den Voraussetzungen des § 218a \abbrStGB{} vorgenommenen Schwangerschaftsabbrüche vor. Die Vorschrift bezog sich ursprünglich auf die in § 218a \abbrStGB{} \abbrAF{} enthaltene Indikationenregelung. Sie kann aber Aufschluß auch über Abbrüche nach dem Beratungskonzept in den ersten zwölf Schwangerschaftswochen und damit über die wesentlichen Wirkungen des neuen Schutzkonzepts geben; solche Kenntnis ist nicht verzichtbar. Der Gesetzgeber war zwar von Verfassungs wegen nicht gehalten, die bisherige Vorschrift über die Bundesstatistik unverändert beizubehalten. Verwehrt war ihm aber deren ersatzlose Streichung. Diese ist nichtig. Eine Neuregelung ist zur Erfüllung der Schutzpflicht schon deshalb erforderlich, weil sich der Geltungsbereich der Vorschrift nicht auf die neuen Bundesländer erstreckt, in denen sie bisher nur aufgrund der einstweiligen Anordnung des Senats (\abbrBVerfGE{} 86, 390 \abbrFF{}) entsprechend gilt.
\end{sourcepara}
\begin{transpara}
\transrn{308}《第五刑法改革法》第4條規定，對依《刑法》第218a條前提所實施之終止妊娠作聯邦統計。該規定原本涉及舊版本《刑法》第218a條所包含之指標規則。但其亦能提供關於妊娠最初十二週內依諮詢構想所為終止妊娠之資訊，因而提供關於新保護構想主要效果之資訊；此種知識不可放棄。立法者固然並非憲法上必須原封不動維持既有關於聯邦統計之規定。但禁止其無替代地刪除該規定。此一刪除無效。為履行保護義務，已因下列理由有必要作成新規制：該規定之適用範圍並未及於新聯邦各邦；在該等地區，迄今僅基於第二庭暫時處分（BVerfGE 86, 390 ff.）而相應適用。
\end{transpara}

\bisubsubsection{V.}{V.}

\begin{sourcepara}[title={原文E.195}]
\sourcern{309}Die Vorschrift des § 24b \abbrSGBV{} ist mit dem Grundgesetz vereinbar (1. und 2. a). Mit dem Grundgesetz wäre es jedoch nicht vereinbar, für die Vornahme eines Schwangerschaftsabbruchs, dessen Rechtmäßigkeit nicht in einer den zuvor entwickelten verfassungsrechtlichen Maßstäben entsprechenden Weise festgestellt wird, einen Anspruch auf Leistungen der gesetzlichen Krankenversicherung zu gewähren (2. b). Die Gewährung von Sozialhilfe für nicht mit Strafe bedrohte Schwangerschaftsabbrüche nach der Beratungsregelung in Fällen wirtschaftlicher Bedürftigkeit ist demgegenüber ebensowenig verfassungsrechtlich zu beanstanden wie die Fortzahlung des Arbeitsentgelts (3. und 4.). Die verfassungsrechtliche Schutzpflicht für das ungeborene Leben steht Leistungen der sozialen Krankenversicherung für einen Abbruch der Schwangerschaft aufgrund der allgemeinen Notlagenindikation nach bisher geltendem Recht nicht entgegen. Ob es unter anderen Gesichtspunkten mit dem Grundgesetz vereinbar war, Leistungen der Sozialversicherung auch dann zu erbringen, wenn der Abbruch der Schwangerschaft auf der Grundlage des § 218a \abbrAbs{} 2 \abbrNr{} 3 \abbrStGB{} \abbrAF{} erfolgte, bedarf keiner Entscheidung (5.).
\end{sourcepara}
\begin{transpara}
\transrn{309}SGB V 第24b條之規定與《基本法》相容（1. 及 2. a）。然而，若終止妊娠之合法性未以符合前述憲法標準之方式確認，卻就其實施給予法定醫療保險給付請求權，則將與《基本法》不相容（2. b）。相對地，對依諮詢規制不受刑罰威嚇之終止妊娠，在經濟需求案件中給予社會扶助，以及工資繼續支付，均在憲法上無可指摘（3. 及 4.）。對未出生生命之憲法保護義務，並不反對依迄今有效法律，基於一般困境指標事由終止妊娠時給予社會醫療保險給付。至於在其他觀點下，當終止妊娠係基於舊版本《刑法》第218a條第2項第3款而為時，社會保險給付是否亦與《基本法》相容，無須決定（5.）。
\end{transpara}

\begin{sourcepara}[title={原文E.196}]
\sourcern{310}1. Die Vorschriften des § 24b \abbrSGBV{} sind formell verfassungsgemäß zustande gekommen. Die Gesetzgebungskompetenz des Bundes ergibt sich aus \abbrArt{} 74 \abbrNr{} 12 \abbrGG{}\ggfnArtSeventyFourTwelveDE{} ("Sozialversicherung einschließlich der Arbeitslosenversicherung"); das gilt auch, soweit Schwangerschaftsabbrüche aufgrund der allgemeinen Notlagenindikation (§ 218a \abbrAbs{} 2 \abbrNr{} 3 \abbrStGB{} \abbrAF{}) betroffen sind.
\end{sourcepara}
\begin{transpara}
\transrn{310}1. SGB V 第24b條之規定在形式上合憲成立。聯邦之立法權限源於《基本法》第74條第12款\ggfnArtSeventyFourTwelveZH{}（「社會保險，包括失業保險」）；此亦適用於涉及基於一般困境指標事由（舊版本《刑法》第218a條第2項第3款）之終止妊娠的範圍。
\end{transpara}

\begin{sourcepara}[title={原文E.197}]
\sourcern{311}a) Wie das Bundesverfassungsgericht mehrfach (zuletzt \abbrBVerfGE{} 75, 108 [146 \abbrF{}]) entschieden hat, ist der Begriff "Sozialversicherung" in \abbrArt{} 74 \abbrNr{} 12 \abbrGG{}\ggfnArtSeventyFourTwelveDE{} als weitgefaßter "verfassungsrechtlicher Gattungsbegriff" zu verstehen. Er umfaßt alles, was sich der Sache nach als Sozialversicherung darstellt. Neue Lebenssachverhalte können in das Gesamtsystem "Sozialversicherung" einbezogen werden, wenn die neuen Sozialleistungen in ihren wesentlichen Strukturelementen, insbesondere in der organisatorischen Durchführung und hinsichtlich der abzudeckenden Risiken, dem Bild entsprechen, das durch die "klassische" Sozialversicherung geprägt ist. Zur Sozialversicherung gehört jedenfalls die gemeinsame Deckung eines möglichen, in seiner Gesamtheit schätzbaren Bedarfs durch Verteilung auf eine organisierte Vielheit (vgl. BSGE 6, 213 [218, 227 \abbrF{}]). Die Beschränkung auf Arbeitnehmer und auf eine Notlage gehört nicht zum Wesen der Sozialversicherung. Außer dem sozialen Bedürfnis nach Ausgleich besonderer Lasten ist die Art und Weise kennzeichnend, wie die Aufgabe organisatorisch bewältigt wird: Träger der Sozialversicherung sind selbständige Anstalten oder Körperschaften des öffentlichen Rechts, die ihre Mittel durch Beiträge der "Beteiligten" aufbringen.
\end{sourcepara}
\begin{transpara}
\transrn{311}a) 如聯邦憲法法院多次判決（最近 BVerfGE 75, 108 [146 f.]）所示，《基本法》第74條第12款\ggfnArtSeventyFourTwelveZH{}之「社會保險」概念，應理解為廣義之「憲法上類型概念」。其包括一切依事物本質呈現為社會保險者。若新社會給付在其主要結構要素上，尤其在組織實施及所涵蓋風險方面，符合由「古典」社會保險所形塑之圖像，則新生活事實得被納入「社會保險」整體制度。社會保險至少包括：將一項可能且整體上可估計之需求，分配予一個有組織多數而共同負擔（參見 BSGE 6, 213 [218, 227 f.]）。限於勞工及限於困境，並非社會保險之本質。除補償特殊負擔之社會需求外，具有標誌性者是任務在組織上如何被克服：社會保險之承辦者為獨立之公法機構或社團，其資金由「參與者」繳納保費籌措。
\end{transpara}

\begin{sourcepara}[title={原文E.198}]
\sourcern{312}b) In das kompetenzrechtliche Bild der Sozialversicherung läßt sich die durch das Strafrechtsreform-Ergänzungsgesetz erstmals vorgenommene und in das Schwangeren- und Familienhilfegesetz übernommene Erstreckung der Krankenversicherungsleistungen auf den "nicht rechtswidrigen Abbruch der Schwangerschaft" durch einen Arzt einordnen. Dabei kommt es für die kompetenzrechtliche Zuordnung nicht darauf an, ob die Kompetenz für eine inhaltlich rechtmäßige oder rechtswidrige Regelung in Anspruch genommen wird; entscheidend ist, ob die Regelung gegenständlich in den Kompetenzbereich fällt.
\end{sourcepara}
\begin{transpara}
\transrn{312}b) 由《刑法改革補充法》首次實施並由《孕婦及家庭扶助法》承接之醫療保險給付延伸至醫師所為「不具違法性之終止妊娠」，可被歸入社會保險之權限法圖像。在此，對權限歸屬而言，主張該權限是否用於內容上合法或違法之規制並非關鍵；決定性者是該規制就標的而言是否落入權限領域。
\end{transpara}

\begin{sourcepara}[title={原文E.199}]
\sourcern{313}aa) Zu den Aufgaben der gesetzlichen Krankenversicherung gehören seit jeher Maßnahmen der Gesundheitsvorsorge. Der Schwangerschaftsabbruch mit Ausnahme des medizinisch indizierten ist jedoch, auch wenn er von einem Arzt vorgenommen wird, weder eine Maßnahme der Gesundheitsvorsorge noch ein Heileingriff (vgl. \abbrBVerfGE{} 39, 1 [44, 46]; vgl. auch BSGE 39, 167 [169]). Er kann daher der Mutterschaftshilfe, den Vorsorgeuntersuchungen und den vorbeugenden Maßnahmen zur Verhütung von Erbkrankheiten nicht gleichgesetzt werden. Diese Leistungen dienen der Verbesserung der Gesundheit oder der Abwehr gesundheitlicher Gefahren (vgl. Gitter/Wendling in: Eser/Hirsch, Sterilisation und Schwangerschaftsabbruch, 1980, \abbrS{} 215). Hieran fehlt es typischerweise bei den Schwangerschaftsabbrüchen nach § 218a \abbrAbs{} 2 \abbrNr{} 3 \abbrStGB{} \abbrAF{} sowie nach § 218a \abbrAbs{} 1 \abbrStGB{} \abbrNF{} Ein in irgendeiner Weise behandlungsbedürftiger Zustand der Schwangeren entsteht hier in der Regel erst durch den Eingriff, dessen finanzielle Absicherung durch die Solidargemeinschaft in Rede steht.
\end{sourcepara}
\begin{transpara}
\transrn{313}aa) 法定醫療保險之任務自始即包括健康預防措施。然而，除醫學指標事由外之終止妊娠，即使由醫師實施，既非健康預防措施，亦非治療性介入（參見 BVerfGE 39, 1 [44, 46]；亦參見 BSGE 39, 167 [169]）。因此，其不能與母職協助、預防檢查及預防遺傳疾病之措施等同。此等給付服務於改善健康或防禦健康危險（參見 Gitter/Wendling in: Eser/Hirsch, Sterilisation und Schwangerschaftsabbruch, 1980, S. 215）。舊版本《刑法》第218a條第2項第3款及新版本《刑法》第218a條第1項之終止妊娠，典型上欠缺此點。在此，懷孕婦女需以某種方式治療之狀態通常是由介入本身才產生；而其透過團結共同體作財務保障，正是此處所涉問題。
\end{transpara}

\begin{sourcepara}[title={原文E.200}]
\sourcern{314}Gleichwohl fehlt diesen von einem Arzt vorgenommenen Schwangerschaftsabbrüchen nicht ein Bezug zu dem -- kompetenzbegründenden -- Thema der Gesundheitsvorsorge. Auch der nicht medizinisch indizierte Schwangerschaftsabbruch stellt sich notwendigerweise als medizinischer Eingriff an der Frau dar und ist von daher mit gesundheitlichen Risiken für sie verbunden. Um der daraus entstehenden Gesundheitsgefahr vorzubeugen, schreiben sowohl § 218a \abbrAbs{} 1 \abbrStGB{} \abbrAF{} wie § 218a \abbrAbs{} 1 \abbrStGB{} \abbrNF{} vor, daß der Schwangerschaftsabbruch, wenn er denn -- unter den dort festgelegten Voraussetzungen -- geschieht, nur von einem Arzt vorgenommen werden darf, damit er nach den Regeln der ärztlichen Kunst erfolgt.
\end{sourcepara}
\begin{transpara}
\transrn{314}儘管如此，由醫師所實施之此等終止妊娠，並非欠缺與構成權限基礎之健康預防主題的關聯。非醫學指標事由之終止妊娠，也必然呈現為對婦女之醫療介入，並因此與其健康風險相連。為預防由此產生之健康危險，舊版本《刑法》第218a條第1項及新版本《刑法》第218a條第1項均規定，若終止妊娠在該處所定前提下發生，則只能由醫師實施，以使其依醫術規則為之。
\end{transpara}

\begin{sourcepara}[title={原文E.201}]
\sourcern{315}bb) Die Krankenversicherungsleistung wird organisatorisch entsprechend dem Bild erbracht, das durch die "klassische" Sozialversicherung geprägt ist, nämlich durch die Träger der gesetzlichen Krankenversicherung in der gleichen Weise wie in den Fällen krankheitsbedingter Arbeitsunfähigkeit oder Behandlungsbedürftigkeit. Die dazu erforderlichen Mittel werden im Umlageverfahren von den Versicherten und ihren Arbeitgebern über die Krankenversicherungsbeiträge aufgebracht.
\end{sourcepara}
\begin{transpara}
\transrn{315}bb) 醫療保險給付在組織上依「古典」社會保險所形塑之圖像提供，亦即由法定醫療保險承辦者，以與因疾病喪失工作能力或需要治療之案件相同方式提供。為此所需資金，透過分攤程序由被保險人及其雇主以醫療保險費籌措。
\end{transpara}

\begin{sourcepara}[title={原文E.202}]
\sourcern{316}2. a) Die Regelung des § 24b \abbrSGBV{}, die Versicherten "bei einem nicht rechtswidrigen Abbruch der Schwangerschaft" Anspruch auf Leistungen der gesetzlichen Krankenversicherung gewährt, knüpft an die Vorschrift des § 218a \abbrAbs{} 1 bis 3 \abbrStGB{} \abbrNF{} an, in der bestimmt wird, unter welchen Voraussetzungen ein Schwangerschaftsabbruch nicht rechtswidrig ist. Ein Anspruch auf Leistungen soll mithin auch in den von § 218a \abbrAbs{} 1 \abbrStGB{} \abbrNF{} erfaßten Fällen bestehen. Diese Vorschrift hält indessen, wie dargelegt, verfassungsrechtlicher Nachprüfung nicht stand; sie ist verfassungswidrig und nichtig. Ist es aber dem Gesetzgeber verfassungsrechtlich verwehrt, einen Schwangerschaftsabbruch unter den hier geregelten Voraussetzungen als gerechtfertigt zu behandeln, so folgt schon aus dem Wortlaut des § 24b \abbrSGBV{}, daß Leistungen der Krankenversicherung insoweit nicht in Betracht kommen. Mit diesem eingeschränkten Anwendungsbereich ist die Vorschrift mit dem Grundgesetz vereinbar.
\end{sourcepara}
\begin{transpara}
\transrn{316}2. a) SGB V 第24b條之規制，賦予被保險人在「不具違法性之終止妊娠」時享有法定醫療保險給付請求權，其連結於新版本《刑法》第218a條第1項至第3項之規定；後者規定終止妊娠在何等前提下為不具違法性。因此，在新版本《刑法》第218a條第1項所涵蓋案件中，亦應存在給付請求權。然而，如前所述，該規定不能通過憲法審查；其違憲且無效。但若立法者在憲法上被禁止將此處所規定前提下之終止妊娠作為正當化處理，則僅依 SGB V 第24b條文字即可得出：醫療保險給付就此不予考慮。在此受限制之適用範圍內，該規定與《基本法》相容。
\end{transpara}

\begin{sourcepara}[title={原文E.203}]
\sourcern{317}b) Die verfassungsrechtliche Schutzpflicht für das Leben verwehrt eine Auslegung des 24b \abbrSGBV{} dahin, daß Leistungen der Sozialversicherung in gleicher Weise wie bei nicht rechtswidrigen Schwangerschaftsabbrüchen dann gewährt werden können, wenn die Rechtmäßigkeit des Schwangerschaftsabbruchs nicht festgestellt werden kann. Der Rechtsstaat darf eine Tötungshandlung nur zum Gegenstand seiner Finanzierung machen, wenn sie rechtmäßig ist und der Staat sich der Rechtmäßigkeit mit rechtsstaatlicher Verläßlichkeit vergewissert hat.
\end{sourcepara}
\begin{transpara}
\transrn{317}b) 生命保護之憲法保護義務，禁止將 SGB V 第24b條解釋為：當終止妊娠之合法性無法確認時，社會保險給付仍得如同不具違法性之終止妊娠一般給予。法治國只有在殺害行為合法，且國家已以法治國之可靠性確認其合法時，始得將其作為國家資助之標的。
\end{transpara}

\begin{sourcepara}[title={原文E.204}]
\sourcern{318}aa) Ohne den von der Frau verlangten Eingriff, der das Leben des Ungeborenen zerstört und nur deshalb die Gesundheit der Frau gefährdet, bestünde für die Gewährung von Ansprüchen gegen die soziale Krankenversicherung kein Anlaß. Macht der Staat einen solchen medizinischen Vorgang zum Gegenstand eines gegen die gesetzliche Krankenversicherung gerichteten Leistungsanspruchs, so dient dies nicht ausschließlich -- wie alle anderen Leistungen der sozialen Krankenversicherung -- dem Schutz von Leben und Gesundheit. Die Übernahme der Arztkosten und die sozialversicherungsrechtlich vorgesehene Organisation und Gewährleistung der ärztlichen Leistungen für die Vornahme von Schwangerschaftsabbrüchen sind zwar nicht unmittelbar ursächlich für die Tötung des ungeborenen Lebens; sie stellen sich aber als eine Beteiligung des Staates an diesem Vorgang dar. Eine solche Beteiligung ist dem Staat nur gestattet, wenn der Tatbestand einer rechtfertigenden Ausnahme vom prinzipiellen Verbot des Schwangerschaftsabbruchs erfüllt und dies mit rechtsstaatlicher Verläßlichkeit festgestellt worden ist. Der Staat darf sich an der Tötung ungeborenen Lebens nicht beteiligen, solange er von der Rechtmäßigkeit dieses Vorgangs nicht überzeugt sein kann. Die Schutzkonzeption der Beratungsregelung läßt indes keinen Raum, dem Staat eine solche Überzeugung zu vermitteln (vgl. oben D.III.1.c).
\end{sourcepara}
\begin{transpara}
\transrn{318}aa) 若無婦女所要求、毀滅未出生者生命且正因如此危及婦女健康之介入，即無給予對社會醫療保險請求權之理由。國家若將此種醫療過程作為針對法定醫療保險之給付請求權標的，則其並非如社會醫療保險所有其他給付般，專門服務於生命與健康保護。承擔醫師費用，以及社會保險法所規定之醫療給付組織與保障，固然並非殺害未出生生命之直接原因；但其呈現為國家參與此一過程。此種參與僅在符合終止妊娠原則禁止之正當化例外構成要件，且該點已以法治國可靠性確認時，國家始被允許。只要國家不能確信此一過程之合法性，即不得參與殺害未出生生命。然而，諮詢規制之保護構想並未留下空間，使國家得形成此種信念（參見上文 D.III.1.c）。
\end{transpara}

\begin{sourcepara}[title={原文E.205}]
\sourcern{319}Kann bei den unter den Bedingungen einer Beratungsregelung im Frühstadium der Schwangerschaft vorgenommenen Abbrüchen nicht festgestellt werden, ob sie wegen des Vorliegens einer allgemeinen Notlage als erlaubt angesehen werden dürfen, so darf sich der Rechtsstaat mithin an ihnen -- auch finanziell oder durch die Verpflichtung Dritter wie der sozialversicherungsrechtlichen Solidargemeinschaften -- grundsätzlich nicht beteiligen. Durch eine solche Beteiligung übernähme der Staat Mitverantwortung für Vorgänge, deren Rechtmäßigkeit er einerseits schon aus verfassungsrechtlichen Gründen nicht als gegeben ansehen darf und andererseits im Rahmen seines Schutzkonzepts festzustellen gehindert ist.
\end{sourcepara}
\begin{transpara}
\transrn{319}若在諮詢規制條件下於懷孕早期所實施之終止妊娠，無法確認其是否因存在一般困境而得被視為允許，則法治國原則上不得參與此等終止妊娠，包括以財務方式或透過使第三人如社會保險法上之團結共同體負擔義務而參與。透過此種參與，國家將共同承擔對下列過程之責任：其一方面基於憲法理由即不得視其合法性為已存在，另一方面在其保護構想框架內又被阻止確認其合法性。
\end{transpara}

\begin{sourcepara}[title={原文E.206}]
\sourcern{320}bb) Auch die Wirksamkeitsbedingungen des Schutzkonzepts (vgl. oben D.III.3.) fordern keine Ausnahme.
\end{sourcepara}
\begin{transpara}
\transrn{320}bb) 保護構想之有效性條件（參見上文 D.III.3.）亦不要求例外。
\end{transpara}

\begin{sourcepara}[title={原文E.207}]
\sourcern{321}(1) Das Beratungskonzept schließt es ein, daß ein Abbruch der Schwangerschaft unter medizinisch unbedenklichen Bedingungen und unter Umständen erfolgen kann, die das Persönlichkeitsrecht der Frau wahren. Diesen Anforderungen ist nur genügt, wenn keine Frau aus finanziellen Gründen an der Inanspruchnahme eines Arztes gehindert ist. Fehlt es an den dafür erforderlichen Mitteln nicht, entspricht es gesicherter Lebenserfahrung, daß Frauen einen Arzt in Anspruch nehmen, um die bei einem nicht fachgerecht durchgeführten Abbruch drohenden Gefährdungen ihrer eigenen Gesundheit zu vermeiden. Nicht von Ärzten durchgeführte illegale Schwangerschaftsabbrüche sind schon vor Inkrafttreten der §§ 218 \abbrFF{} \abbrStGB{} \abbrAF{} nicht sehr häufig gewesen (vgl. dazu den Ersten Bericht des Sonderausschusses für die Strafrechtsreform, \abbrBTDrucks{} 7/1981 [neu], \abbrS{} 5 \abbrF{}). Soweit die Frauen allerdings nicht über hinreichendes eigenes Einkommen oder Vermögen verfügen, kann der Staat nach den Grundsätzen des Sozialhilferechts auch hier ihren Bedarf decken (vgl. dazu unten 3.). Für die Beurteilung der Bedürftigkeit kommt es dabei auf das zum Zeitpunkt des Abbruchs für die Frau schon verfügbare Einkommen und Vermögen an; weder darf sie insoweit auf etwaige Unterhaltsansprüche gegen die Eltern oder den Ehemann verwiesen noch darf bei diesen Rückgriff genommen werden, sofern die Frau nicht damit einverstanden ist.
\end{sourcepara}
\begin{transpara}
\transrn{321}（1）諮詢構想包含：終止妊娠得在醫學上無疑義之條件下，並在維護婦女人格權之情況下發生。只有在沒有婦女因財務理由而被阻礙尋求醫師時，方符合此等要求。若不欠缺為此所需資金，依確定生活經驗，婦女會尋求醫師，以避免非專業實施終止妊娠時威脅其自身健康之危險。非由醫師實施之非法終止妊娠，在舊版本《刑法》第218條以下生效前即並非很常見（對此參見刑法改革特別委員會第一報告，BTDrucks. 7/1981 [neu], S. 5 f.）。不過，若婦女並無足夠自身收入或財產，國家亦可依社會扶助法原則在此填補其需求（對此見下文 3.）。判斷需求性時，應以終止妊娠時婦女已可支配之收入及財產為準；在此既不得使其指向對父母或丈夫可能享有之扶養請求權，亦不得向其等追償，除非婦女同意。
\end{transpara}

\begin{sourcepara}[title={原文E.208}]
\sourcern{322}(2) Private Krankenversicherungen decken im allgemeinen nur "medizinisch notwendige Behandlung wegen Schwangerschaft" (§ 1 \abbrAbs{} 2 Satz 4 Buchst. a, Allgemeine Versicherungsbedingungen für die Krankheitskosten- und Krankenhaustagegeldversicherung [Musterbedingungen des Verbandes der privaten Krankenversicherung 1976], abgedruckt bei Prölss/Martin, Versicherungsvertragsgesetz, 24. \abbrAufl{} 1988, \abbrS{} 1222). Hierunter werden nur medizinisch indizierte Schwangerschaftsabbrüche zur Beseitigung eines pathologischen Befundes oder zur Abwendung einer physischen oder psychischen Gefahr für die Schwangere verstanden (vgl. Wriede in Bruck/Möller, Kommentar zum Versicherungsvertragsgesetz, 8. \abbrAufl{}, Band VI 2, 1990, G 43; vgl. auch LG Berlin, VersR 1983, \abbrS{} 1180 \abbrF{}; LG Detmold, VersR 1986, \abbrS{} 336). Versicherungsverträge, die darüber hinaus auf die Erbringung von Leistungen für Schwangerschaftsabbrüche gerichtet sind, deren Rechtmäßigkeit nicht festgestellt ist, würden sich in Widerspruch zu dem grundsätzlichen Verbot des Schwangerschaftsabbruchs setzen. Verträge dieses Inhalts wären unwirksam (vgl. D.III.3.).
\end{sourcepara}
\begin{transpara}
\transrn{322}（2）私人醫療保險通常僅涵蓋「因懷孕之醫學必要治療」（疾病費用及住院日額保險一般保險條款第1條第2項第4句 a 字〔私人醫療保險協會1976年示範條款〕，刊於 Prölss/Martin, Versicherungsvertragsgesetz, 24. Aufl. 1988, S. 1222）。此處僅理解為為排除病理性發現或防止懷孕婦女身體或心理危險所為之醫學指標事由終止妊娠（參見 Wriede in Bruck/Möller, Kommentar zum Versicherungsvertragsgesetz, 8. Aufl., Band VI 2, 1990, G 43；亦參見 LG Berlin, VersR 1983, S. 1180 f.; LG Detmold, VersR 1986, S. 336）。除此之外，若保險契約旨在對合法性未經確認之終止妊娠提供給付，將與終止妊娠之原則禁止相矛盾。此種內容之契約將無效（參見 D.III.3.）。
\end{transpara}

\begin{sourcepara}[title={原文E.209}]
\sourcern{323}Nach alledem sind Frauen, die keine oder nur eine private Krankenversicherung haben, ohnehin darauf angewiesen, die ärztlichen Leistungen für den nach der Beratungsregelung nicht mit Strafe bedrohten Schwangerschaftsabbruch aus eigenen Mitteln zu bezahlen. Es ist nicht ersichtlich, daß es die Wirksamkeit des Schutzkonzepts der Beratungsregelung erforderte, demgegenüber Frauen, die in der gesetzlichen Krankenversicherung versichert sind, anders zu behandeln.
\end{sourcepara}
\begin{transpara}
\transrn{323}綜上，沒有醫療保險或僅有私人醫療保險之婦女，無論如何均須依靠自身資金支付依諮詢規制不受刑罰威嚇之終止妊娠的醫療給付。並不能看出，諮詢規制保護構想之有效性要求相對於此，對投保法定醫療保險之婦女作不同處理。
\end{transpara}

\begin{sourcepara}[title={原文E.210}]
\sourcern{324}cc) Nach dem Schutzkonzept der Beratungsregelung hat es die sozial und ärztlich beratene Frau zu verantworten, ob es zu einem -- dann nicht mit Strafe bedrohten -- Schwangerschaftsabbruch kommt oder nicht. Aus dem Beratungskonzept ergibt sich jedoch keine Notwendigkeit, diese Verantwortung sozialrechtlich anzuerkennen. Es erfordert nicht, den Frauen, die einen nicht mit Strafe bedrohten Abbruch unter den Bedingungen der Beratungsregelung durchführen, dieselben sozialrechtlichen Leistungen zuteil werden zu lassen wie denen, bei denen eine medizinische, eine embryopathische oder auch eine kriminologische Indikation, mithin ein Rechtfertigungsgrund, festgestellt ist.
\end{sourcepara}
\begin{transpara}
\transrn{324}cc) 依諮詢規制之保護構想，接受社會及醫療諮詢之婦女，須對是否發生一項不受刑罰威嚇之終止妊娠負責。然而，由諮詢構想並不能得出有必要在社會法上承認此一責任。其並不要求，對於在諮詢規制條件下實施不受刑罰威嚇之終止妊娠之婦女，給予與存在醫學、胚胎病理或犯罪學指標事由，因而確認有正當化事由之婦女相同之社會法給付。
\end{transpara}

\begin{sourcepara}[title={原文E.211}]
\sourcern{325}(1) Die Verantwortung der Frau bei der Entscheidung für einen Schwangerschaftsabbruch unter den Bedingungen der Beratungsregelung ist nicht als "Selbstindikation" zu verstehen, durch die das Vorliegen einer rechtfertigenden Notlage von ihr selbst verbindlich festgestellt wird (siehe oben D.III.2.b)aa). Das Beratungskonzept beruht auf der Einschätzung, daß in der Frühphase einer Schwangerschaft das ungeborene Leben besser mit der Mutter geschützt werden kann. Deshalb sollen, um die in der Beratung liegenden Möglichkeiten einer Ermutigung zum Austragen des Kindes auszuschöpfen, einer Entscheidung, den Abbruch doch durchzuführen, keine der Schutzwirkung entgegenstehenden rechtlichen oder tatsächlichen Hindernisse in den Weg gelegt werden. Einer darüber hinausgehenden Anerkennung der Entscheidung der Frau bedarf es zur Verwirklichung des Schutzes durch beratende Einflußnahme nicht.
\end{sourcepara}
\begin{transpara}
\transrn{325}（1）婦女在諮詢規制條件下決定終止妊娠之責任，不得理解為「自我指標事由」，由其自身以拘束力確認正當化困境之存在（見上文 D.III.2.b)aa）。諮詢構想基於下列評估：在懷孕早期，未出生生命可更好地與母親一同受到保護。因此，為窮盡諮詢中所含鼓勵繼續懷孕之可能性，對於終究仍實施終止妊娠之決定，不應設置與保護效果相牴觸之法律或事實障礙。除此之外，為透過諮詢性影響實現保護，並不需要對婦女之決定作進一步承認。
\end{transpara}

\begin{sourcepara}[title={原文E.212}]
\sourcern{326}(2) Auch das Sozialstaatsprinzip (\abbrArt{} 20 \abbrAbs{} 1 \abbrGG{}\ggfnArtTwentyOneDE{}) gestattet es dem Staat nicht, die im Rahmen der Beratungsregelung nicht mit Strafe bedrohten Abbrüche mit Rücksicht darauf, daß eine Bewertung ihrer Rechtmäßigkeit im Einzelfall nicht stattfindet, so zu behandeln, als seien sie alle erlaubt. Der Sozialstaat kann auf dem Boden des Grundgesetzes nur mit den Mitteln des Rechtsstaats verwirklicht werden. Der Grundsatz der Rechtsstaatlichkeit wäre auch nicht etwa nur am Rande berührt, sondern vielmehr im Kern verletzt, übernähme der Staat allgemein -- also ohne eine dem Gedanken des Sozialstaats Rechnung tragende Differenzierung -- mittelbar oder unmittelbar (Mit-)Verantwortung für Vorgänge, von deren Rechtmäßigkeit er nicht überzeugt sein kann.
\end{sourcepara}
\begin{transpara}
\transrn{326}（2）社會國原則（《基本法》第20條第1項\ggfnArtTwentyOneZH{}）亦不允許國家將諮詢規制框架內不受刑罰威嚇之終止妊娠，基於個案中未評價其合法性，而作為彷彿其全部均被允許般處理。在《基本法》基礎上，社會國只能以法治國之手段實現。若國家一般性地，也就是未依社會國思想作區分地，間接或直接對其無法確信合法性之過程承擔共同責任，法治國原則並非僅邊緣受觸動，而是其核心受到侵害。
\end{transpara}

\begin{sourcepara}[title={原文E.213}]
\sourcern{327}dd) Die Gewährung sozialversicherungsrechtlicher Leistungen für Schwangerschaftsabbrüche, die zwar nicht mit Strafe bedroht werden, deren Rechtmäßigkeit aber nicht feststeht, ist auch deshalb mit der staatlichen Schutzpflicht für das ungeborene menschliche Leben nicht vereinbar, weil dadurch das allgemeine Bewußtsein in der Bevölkerung, daß das Ungeborene auch gegenüber der Mutter ein Recht auf Leben hat und daher der Abbruch der Schwangerschaft grundsätzlich Unrecht ist, erheblich beschädigt würde.
\end{sourcepara}
\begin{transpara}
\transrn{327}dd) 對雖不受刑罰威嚇但其合法性未確定之終止妊娠給予社會保險法給付，亦因下列理由與國家對未出生人類生命之保護義務不相容：此將重大損害人民一般意識中，未出生者亦相對於母親享有生命權，因而終止妊娠原則上是不法之認識。
\end{transpara}

\begin{sourcepara}[title={原文E.214}]
\sourcern{328}Die Nichtgewährung von Sozialleistungen ist zwar nur begrenzt geeignet, ein rechtliches Unwerturteil über einen Sachverhalt zu verdeutlichen, zumal wenn sie sich auf Leistungen -- wie Schwangerschaftsabbrüche -- bezieht, die außerhalb des eigentlichen Aufgabenbereichs der Versicherungen liegen. Jedoch muß umgekehrt, wie auch der Sachverständige \abbrProf{} \abbrDr{} Schulin (Gutachten \abbrS{} 97) ausführt, die Gewährung solcher Leistungen für Sachverhalte, die nicht dem normalen Versicherungsrisiko zuzuordnen sind, in der Bevölkerung den Eindruck erwecken, als handele es sich dabei um einen dem üblichen Risiko gleichkommenden und in diesem Sinne alltäglichen, also der Normalität entsprechenden Vorgang.
\end{sourcepara}
\begin{transpara}
\transrn{328}不給予社會給付，固然僅有限適於清楚顯示對某一事實關係之法律負價值判斷，尤其當其涉及如終止妊娠般位於保險真正任務領域之外之給付時。然而反過來，如鑑定人 Schulin 教授（鑑定書 S. 97）亦指出，對不屬於通常保險風險之事實關係給予此類給付，必然在人民中造成一種印象，彷彿其為等同於通常風險，並在此意義上日常、亦即符合正常性之過程。
\end{transpara}

\begin{sourcepara}[title={原文E.215}]
\sourcern{329}Hierbei kann auch nicht unberücksichtigt bleiben, daß die Leistungstatbestände der Sozialversicherung für etwa 90 vom Hundert der Bevölkerung Bedeutung erlangen (vgl. Schulin, \abbrAAO{}, \abbrS{} 2). Wegen ihrer Auswirkungen auf den persönlichen Lebensbereich werden sie von der Solidargemeinschaft der Versicherten mit besonderem Interesse wahrgenommen und sind damit durchaus geeignet, mit Wertungen, die in ihnen zum Ausdruck kommen, die Vorstellungen der Bevölkerung zu prägen.
\end{sourcepara}
\begin{transpara}
\transrn{329}在此亦不能忽略，社會保險之給付構成要件對約百分之九十之人口具有意義（參見 Schulin, a.a.O., S. 2）。由於其對個人生活領域之影響，被保險人團結共同體以特別興趣感知此等構成要件，因而其所表達之評價完全適於形塑人民觀念。
\end{transpara}

\begin{sourcepara}[title={原文E.216}]
\sourcern{330}Das Konzept der Beratungsregelung kann die Mindestanforderungen an die staatliche Schutzpflicht nur dann erfüllen, wenn es auf die Erhaltung und Stärkung des Rechtsbewußtseins besonderen Bedacht nimmt. Nur wenn das Bewußtsein von dem Recht des Ungeborenen auf Leben wach erhalten wird, kann die unter den Bedingungen der Beratungsregelung von der Frau zu tragende Verantwortung an diesem Recht ausgerichtet und prinzipiell geeignet sein, das Leben des ungeborenen Kindes zu schützen. Diesem Schutz liefe es zuwider, wenn der Staat durch die allgemeine Gewährung von sozialversicherungsrechtlichen Leistungsansprüchen den Abbruch unterstützte und damit unvermeidlich den Eindruck erweckte, der Schwangerschaftsabbruch würde von der Rechtsordnung am Ende doch gutgeheißen. Auch die der schwangeren Frau nahestehenden, für sie und das Ungeborene mitverantwortlichen Personen könnten sich entlastet fühlen, weil sie als normal und rechtmäßig empfinden werden, wofür Leistungen der Sozialversicherung gewährt werden.
\end{sourcepara}
\begin{transpara}
\transrn{330}諮詢規制構想只有在特別注意維持並強化法律意識時，始能滿足國家保護義務之最低要求。只有在未出生者生命權之意識保持清醒時，婦女在諮詢規制條件下所承擔之責任，始能以此權利為取向，並原則上適於保護未出生子女之生命。若國家透過一般給予社會保險法給付請求權支持終止妊娠，並因此不可避免地造成終止妊娠最終仍受法秩序肯認之印象，即與此種保護背道而馳。與懷孕婦女親近、對其及未出生者共同負責之人，亦可能感到責任負擔被解除，因為其會將社會保險給付所涵蓋者感受為正常且合法。
\end{transpara}

\begin{sourcepara}[title={原文E.217}]
\sourcern{331}c) Die Verfassung schließt nach allem die Gewährung sozialversicherungsrechtlicher Leistungen für die ärztliche Vornahme des nicht rechtmäßigen Schwangerschaftsabbruchs und die medizinische Nachsorge bei komplikationslosem Verlauf aus. Insoweit lassen es die dargelegten Grundsätze auch nicht zu, der infolge des Schwangerschaftsabbruchs arbeitsunfähig gewordenen Frau einen Anspruch auf Gewährung von Krankengeld nach § 24b \abbrAbs{} 2 Satz 2 \abbrSGBV{} einzuräumen; auch insoweit kann diese Vorschrift daher nicht über ihre ausdrückliche Voraussetzung (nicht rechtswidrig) hinaus bei nicht mit Strafe bedrohten Schwangerschaftsabbrüchen Grundlage für die Gewährung von Leistungen sein; die Frage, ob sich für diese Fälle ein Anspruch im Rahmen der §§ 44, 52 \abbrSGBV{} ergeben kann, ist damit nicht entschieden.
\end{sourcepara}
\begin{transpara}
\transrn{331}c) 綜上，憲法排除對非合法終止妊娠之醫師實施及無併發症過程中之醫療後續照護給予社會保險法給付。就此，前述原則亦不允許賦予因終止妊娠而喪失工作能力之婦女依 SGB V 第24b條第2項第2句取得病假津貼之請求權；因此，該規定就此亦不能超越其明示前提（不具違法性），成為對不受刑罰威嚇之終止妊娠給予給付之基礎；至於此等案件是否可能在 SGB V 第44條、第52條框架內產生請求權，並未因此決定。
\end{transpara}

\begin{sourcepara}[title={原文E.218}]
\sourcern{332}Die übrigen nach § 24b \abbrAbs{} 2 Satz 1 \abbrSGBV{} zu erbringenden Leistungen sind von dem Verbot nicht betroffen; sie sind geeignet und bestimmt, die Gesundheit der Frau, soweit sie durch eine Schwangerschaft berührt wird, zu erhalten. Die ärztlichen Leistungen, die der versicherten Frau im Vorfeld eines Schwangerschaftsabbruchs erbracht werden, sind ebenso wie Nachbehandlungen, die durch komplikationsbedingte Folgeerscheinungen des Abbruchs veranlaßt sind, Leistungen, die der Gesundheit der Frau unmittelbar dienen und überdies, wenn es nicht zum Schwangerschaftsabbruch kommt, auch dem Leben des Kindes zugute kommen. Nur soweit es sich um den Abbruch selber handelt, steht die Tötung des ungeborenen Lebens so sehr im Vordergrund, daß die Inanspruchnahme der Sozialversicherung aus den zuvor genannten Gründen nicht in Betracht kommt, wenn die Rechtmäßigkeit des Schwangerschaftsabbruchs nicht feststeht.
\end{sourcepara}
\begin{transpara}
\transrn{332}依 SGB V 第24b條第2項第1句應提供之其他給付，不受此禁止影響；其適於且旨在維持婦女健康，只要其受懷孕影響。被保險婦女在終止妊娠前階段所受醫療給付，以及因終止妊娠併發症後續現象而引起之後續治療，同樣是直接服務於婦女健康之給付；此外，若未發生終止妊娠，亦有利於子女生命。只有在涉及終止妊娠本身時，殺害未出生生命才如此居於前景，以致若終止妊娠之合法性未確定，即基於前述理由不得使用社會保險。
\end{transpara}

\begin{sourcepara}[title={原文E.219}]
\sourcern{333}3. a) Nach alledem ist es dem Staat verfassungsrechtlich grundsätzlich untersagt, durch die Gewährung von Leistungen oder durch die normative Begründung von Leistungspflichten Dritter Schwangerschaftsabbrüche zu fördern, deren Rechtmäßigkeit nicht verbindlich festgestellt worden ist. Diesen Grundsatz zu durchbrechen ist dem Staat verfassungsrechtlich, wie dargelegt, nur erlaubt, soweit dies das Konzept zum Schutz des ungeborenen menschlichen Lebens um seiner Wirksamkeit willen gerade erfordert, wenn es also darum geht, daß die Frau einen etwaigen Abbruch nur durch einen Arzt vornehmen läßt. Verlangt die Rechtsordnung nun einerseits zum Schutze der Gesundheit der schwangeren Frau, andererseits aber auch um des Schutzes des Ungeborenen willen die Inanspruchnahme eines Arztes, so darf dies nicht daran scheitern, daß die Frau nicht über die dafür erforderlichen finanziellen Mittel verfügt. In diesem Fall darf es dem Staat nicht verwehrt sein, diese Mittel selbst bereitzustellen.
\end{sourcepara}
\begin{transpara}
\transrn{333}3. a) 綜上，國家原則上在憲法上被禁止透過給予給付或透過規範性建立第三人給付義務，促進其合法性尚未受拘束力確認之終止妊娠。如前所述，國家在憲法上只有在這正是保護未出生人類生命之構想為其有效性所要求時，亦即涉及婦女僅由醫師實施可能終止妊娠時，始得突破此一原則。若法秩序一方面為保護懷孕婦女健康，另一方面也為保護未出生者，而要求求助於醫師，則不得因婦女不具備為此必要之財務資金而使其失敗。在此情形，國家不得被禁止自行提供此等資金。
\end{transpara}

\begin{sourcepara}[title={原文E.220}]
\sourcern{334}b) Sache des Gesetzgebers ist es zu regeln, auf welchem Wege und unter welchen Voraussetzungen in den Fällen, in denen das Schutzkonzept der Beratungsregelung dies erfordert, bei Bedürftigkeit der Frau eine Kostenübernahme durch den Staat erfolgen soll. Es liegt nahe, daß dazu die derzeitige Regelung des § 37a BSHG den sich aus den verfassungsrechtlichen Vorgaben für die Beratungsregelung ergebenden Folgen angepaßt wird. Mit der Gewährung dieser Sozialleistung setzt der Staat sich nicht in Widerspruch zu den Anforderungen seiner Schutzpflicht. Er verhindert damit von vornherein, daß Frauen den Weg in die Illegalität suchen und damit nicht nur sich selbst gesundheitlichen Schaden zufügen, sondern auch dem Ungeborenen die Chance einer Rettung durch ärztliche Beratung nehmen.
\end{sourcepara}
\begin{transpara}
\transrn{334}b) 在諮詢規制保護構想有所要求之案件中，當婦女有需求時，國家應循何種道路及在何等前提下承擔費用，屬立法者之事。較自然的是，將目前 BSHG 第37a條之規制調整為符合諮詢規制憲法預設所生之後果。國家給予此一社會給付，並不與其保護義務之要求相矛盾。國家由此從一開始即防止婦女尋求非法途徑，並因此不僅對自身健康造成損害，也剝奪未出生者透過醫療諮詢獲救之機會。
\end{transpara}

\begin{sourcepara}[title={原文E.221}]
\sourcern{335}Bei der Ausgestaltung des Sozialhilfeanspruchs hat der Gesetzgeber das Persönlichkeitsrecht der Leistungsberechtigten zu schützen. Soweit es der Schutzpflicht für das ungeborene Leben nicht zuwiderläuft, hat er im vorliegenden Zusammenhang insbesondere Regelungen zu treffen, die der Frau möglichst eine wiederholte Darlegung ihrer Lage ersparen. Dies steht einem Rückgriff bei Familienangehörigen gemäß §§ 91 \abbrF{} BSHG entgegen. Es liegt außerdem nahe, alle Verfahren zur staatlichen Gewährung von Schutz und Hilfe möglichst bei einer Behörde, etwa bei der gesetzlichen Krankenversicherung, zusammenzufassen und so zugleich sicherzustellen, daß die Frau nur einmal ihre Situation darlegen muß.
\end{sourcepara}
\begin{transpara}
\transrn{335}在形塑社會扶助請求權時，立法者必須保護給付權利人之人格權。只要不違反對未出生生命之保護義務，立法者在本脈絡中尤其應制定規定，使婦女盡可能免於反覆陳述其處境。這排除依 BSHG 第91條以下向家庭成員追償。此外，較自然的是，將所有國家提供保護與協助之程序盡可能集中於一個機關，例如法定醫療保險，並同時確保婦女僅須一次陳述其處境。
\end{transpara}

\begin{sourcepara}[title={原文E.222}]
\sourcern{336}4. Auch im Lohnfortzahlungsrecht erweist es sich im Blick auf seine besondere arbeitsrechtliche Konzeption und die Erfordernisse des Schutzkonzepts entsprechend den oben (D.III.3.) dargestellten Grundsätzen nicht als geboten, die lediglich vom Straftatbestand des § 218 \abbrStGB{} \abbrNF{} ausgenommenen Schwangerschaftsabbrüche aus der Leistungspflicht auszuschließen.
\end{sourcepara}
\begin{transpara}
\transrn{336}4. 在工資繼續支付法中，鑑於其特殊勞動法構想及保護構想之要求，依上文（D.III.3.）所述原則，也不顯示有必要將僅被排除於新版本《刑法》第218條構成要件之外之終止妊娠，排除於給付義務之外。
\end{transpara}

\begin{sourcepara}[title={原文E.223}]
\sourcern{337}a) aa) Nach § 1 \abbrAbs{} 1 Satz 1 LFZG verliert ein Arbeiter, der nach Beginn der Beschäftigung durch Arbeitsunfähigkeit infolge Krankheit an seiner Arbeitsleistung verhindert ist, ohne daß ihn ein Verschulden trifft, dadurch nicht den Anspruch auf Arbeitsentgelt für die Zeit der Arbeitsunfähigkeit bis zur Dauer von sechs Wochen. Diese Vorschrift gilt entsprechend, wenn die Arbeitsunfähigkeit infolge des Abbruchs einer Schwangerschaft durch einen Arzt eintritt; bei einem nicht rechtswidrigen Abbruch durch einen Arzt gilt die dadurch bedingte Arbeitsunfähigkeit als unverschuldet (§ 1 \abbrAbs{} 2 LFZG). Andere arbeitsrechtliche Vorschriften enthalten ähnliche Bestimmungen (vgl. § 616 BGB, § 63 HGB, § 133c GewO, §§ 52a, 78 Seemannsgesetz, § 12 BBiG). In dem in \abbrArt{} 3 EV genannten Gebiet gilt § 115a des Arbeitsgesetzbuches der \abbrDDR{} fort (\abbrArt{} 9 \abbrAbs{} 2 EV in Verbindung mit Anlage II Kapitel VIII Sachgebiet A Abschnitt III \abbrNr{} 1a), der durch Gesetz vom 22. Juni 1990 (GBl. I \abbrNr{} 35 \abbrS{} 371) in das Arbeitsgesetzbuch eingefügt worden war. Die Vorschrift lehnt sich eng an die Regelungen des Lohnfortzahlungsgesetzes an, enthält indessen keine dem § 1 \abbrAbs{} 2 LFZG entsprechende Regelung und gilt grundsätzlich für alle Arbeitnehmer. Daneben bestimmt § 4 des Gesetzes über die Unterbrechung der Schwangerschaft vom 9. März 1972 (GBl. I \abbrNr{} 5 \abbrS{} 89), daß die Vorbereitung, Durchführung und Nachbehandlung einer nach diesem Gesetz zulässigen Unterbrechung der Schwangerschaft arbeits- und versicherungsrechtlich dem Erkrankungsfall gleichgestellt sind. Diese Vorschrift wird mit dem Wirksamwerden des \abbrArt{} 16 \abbrSFHG{} außer Kraft treten.
\end{sourcepara}
\begin{transpara}
\transrn{337}a) aa) 依 LFZG 第1條第1項第1句，工人在開始受雇後，因疾病所致工作不能而無法履行勞務，且其無過失者，並不因此喪失工作不能期間至六週為止之工資請求權。若工作不能因醫師終止妊娠而發生，該規定相應適用；若由醫師所為者為不具違法性之終止妊娠，因此所生工作不能視為無過失（LFZG 第1條第2項）。其他勞動法規定包含類似規定（參見《民法》第616條、HGB 第63條、GewO 第133c條、Seemannsgesetz 第52a條、第78條、BBiG 第12條）。在 EV 第3條所稱地區，德意志民主共和國《勞動法典》第115a條繼續有效（EV 第9條第2項結合附件 II 第VIII章事項 A 第III節第1a號）；該條係由1990年6月22日法律（GBl. I Nr. 35 S. 371）納入《勞動法典》。該規定緊密依循《工資繼續支付法》規制，但不包含相當於 LFZG 第1條第2項之規制，且原則上適用於所有勞工。此外，1972年3月9日《終止妊娠法》第4條（GBl. I Nr. 5 S. 89）規定，依該法容許之終止妊娠，其準備、實施及後續治療，在勞動法及保險法上與疾病案件等同。該規定將隨 SFHG 第16條生效而失效。
\end{transpara}

\begin{sourcepara}[title={原文E.224}]
\sourcern{338}bb) Das Lohnfortzahlungsrecht sieht vor, daß der Arzt einem Arbeitnehmer zur Vorlage gegenüber dem Arbeitgeber eine Bescheinigung über das Vorliegen einer krankheitsbedingten Arbeitsunfähigkeit erteilt, in der grundsätzlich weder die Art der Erkrankung noch ihre Ursache mitgeteilt werden (vgl. § 3 LFZG). Der Arbeitgeber kann somit nicht ohne weiteres feststellen, ob die Voraussetzungen für die Entgeltfortzahlung möglicherweise wegen eines Verschuldens des Arbeitnehmers nicht gegeben sind. Er kann die Lohnfortzahlung nur verweigern, wenn er konkrete Anhaltspunkte für das Vorliegen eines einen Schuldvorwurf rechtfertigenden Sachverhalts hat, oder der Arbeitnehmer arbeitsvertraglich verpflichtet ist, entsprechende Angaben von sich aus zu machen.
\end{sourcepara}
\begin{transpara}
\transrn{338}bb) 工資繼續支付法規定，醫師為勞工出具關於存在因疾病所致工作不能之證明，以供其向雇主提出；該證明原則上既不告知疾病種類，亦不告知其原因（參見 LFZG 第3條）。因此，雇主不能逕行確定，工資繼續支付之前提是否可能因勞工有過失而不存在。雇主只有在其對足以正當化有責指摘之事實關係存在有具體線索，或勞工依勞動契約負有自行作相應說明之義務時，始得拒絕工資繼續支付。
\end{transpara}

\begin{sourcepara}[title={原文E.225}]
\sourcern{339}b) Die Ursache einer krankheitsbedingten Arbeitsunfähigkeit wird nur dann in den Blick genommen, wenn die Krankheit auf einem gröblichen Verstoß gegen das von einem verständigen Menschen im eigenen Interesse zu erwartende Verhalten beruht und es -- in Ausnahmefällen -- unbillig wäre, die Folgen eines zu einer Krankheit führenden Verhaltens auf den Arbeitgeber abzuwälzen (vgl. BAG, Urteil vom 28. Februar 1979, AP \abbrNr{} 44 zu § 1 LFZG, Urteil vom 7. August 1991, AP \abbrNr{} 94 zu § 1 LFZG; Hueck/Nipperdey, Lehrbuch des Arbeitsrechts, 7. \abbrAufl{}, \abbrBd{} I, 1963, § 44 III 1 a cc; Zöllner/Loritz, Arbeitsrecht, 4. \abbrAufl{}, 1992, § 18 II 2 e; Schmitt, Lohnfortzahlungsgesetz, 1992, § 1 LFZG, \abbrRdnrn{} 60 \abbrFF{} \abbrMwN{}).
\end{sourcepara}
\begin{transpara}
\transrn{339}b) 因疾病所致工作不能之原因，只有在疾病係基於對通情達理之人按自身利益可期待行為之重大違反，且在例外案件中將導致疾病之行為後果轉嫁於雇主為不公平時，始被納入考量（參見 BAG, Urteil vom 28. Februar 1979, AP Nr. 44 zu § 1 LFZG; Urteil vom 7. August 1991, AP Nr. 94 zu § 1 LFZG; Hueck/Nipperdey, Lehrbuch des Arbeitsrechts, 7. Aufl., Bd. I, 1963, § 44 III 1 a cc; Zöllner/Loritz, Arbeitsrecht, 4. Aufl., 1992, § 18 II 2 e; Schmitt, Lohnfortzahlungsgesetz, 1992, § 1 LFZG, Rdnrn. 60 ff. m.w.N.）。
\end{transpara}

\begin{sourcepara}[title={原文E.226}]
\sourcern{340}Der verfassungsrechtlichen Schutzpflicht für das ungeborene menschliche Leben widerspricht es nicht, wenn diese arbeitsrechtlichen Grundsätze dahin ausgelegt und angewendet werden, daß eine Verpflichtung zur Lohnfortzahlung auch dann besteht, wenn die Arbeitsunfähigkeit die Folge eines auf der Grundlage der Beratungsregelung erfolgten Schwangerschaftsabbruchs ist; es ist verfassungsrechtlich nicht zu beanstanden, auch in diesen Fällen entsprechend § 1 \abbrAbs{} 2 LFZG die Arbeitsunfähigkeit als unverschuldet anzusehen.
\end{sourcepara}
\begin{transpara}
\transrn{340}若將此等勞動法原則解釋並適用為，即使工作不能係基於諮詢規制所為終止妊娠之結果，亦存在工資繼續支付義務，並不違反對未出生人類生命之憲法保護義務；在憲法上無可指摘的是，在此等案件中亦相應於 LFZG 第1條第2項，將工作不能視為無過失。
\end{transpara}

\begin{sourcepara}[title={原文E.227}]
\sourcern{341}Ein Ausschluß des Anspruchs der Arbeitnehmerin auf Lohnfortzahlung in dem hier in Rede stehenden Fall erforderte nicht nur eine gesetzliche Regelung, die es ausschließt, eine durch einen nicht rechtmäßigen Abbruch der Schwangerschaft eingetretene Arbeitsunfähigkeit als krankheitsbedingt anzusehen. Es müßte um der Wirksamkeit einer solchen Regelung willen entweder dem behandelnden Arzt untersagt werden, der Arbeitnehmerin Arbeitsunfähigkeit zu bescheinigen, oder es müßte die Arbeitnehmerin selbst verpflichtet werden, dem Arbeitgeber den Grund ihrer Arbeitsunfähigkeit offenzulegen. Dies wiederum ließe sich nur vermeiden, wenn die schwangere Frau für die Vornahme des Abbruchs Urlaub nähme, was nicht immer möglich sein wird. Andere Möglichkeiten, die Offenlegung zu vermeiden, stehen nach geltendem Recht nicht zur Verfügung. Eine Offenlegungspflicht der Frau bestünde auch dann, wenn man annehmen wollte, daß der Schwangerschaftsabbruch nach Beratung ein gröblicher Verstoß gegen das von einem verständigen Menschen im eigenen Interesse zu erwartende Verhalten, also im Sinne der hier einschlägigen arbeitsrechtlichen Regelung verschuldet sein kann. Das Schutzkonzept der Beratungsregelung erfordert es jedoch gerade, daß die schwangere Frau, um dem Anliegen der Beratung, das Leben des Ungeborenen zu erhalten, aufgeschlossen zu begegnen, der Notwendigkeit enthoben wird, die Gründe ihrer Entscheidung für den Abbruch außerhalb der Beratung und des ärztlichen Gesprächs gegenüber Dritten darzulegen. Diese Aufgeschlossenheit würde jedenfalls gefährdet, wenn das Arbeitsrecht auf die eine oder andere Weise eine Verpflichtung der Arbeitnehmerin begründete, dem Arbeitgeber offenzulegen, daß ihre Arbeitsunfähigkeit auf einen nicht gerechtfertigten Schwangerschaftsabbruch zurückzuführen ist. Unter diesen Umständen ist es auch nicht unbillig, die abbruchbedingte Arbeitsunfähigkeit der Arbeitnehmerin dem Risikobereich des Arbeitgebers zuzurechnen.
\end{sourcepara}
\begin{transpara}
\transrn{341}在此處所涉案件中排除女勞工之工資繼續支付請求權，不僅需要一項法律規制，排除將因非合法終止妊娠所生工作不能視為因疾病所致。為使此種規制有效，或須禁止治療醫師向女勞工證明工作不能，或須使女勞工自身負有向雇主揭露其工作不能原因之義務。若要避免此點，則只有懷孕婦女為實施終止妊娠而休假；但這並非總是可能。依現行法，並無其他避免揭露之可能。即使認為經諮詢後之終止妊娠可能是對通情達理之人按自身利益可期待行為之重大違反，亦即在此相關勞動法規制意義下有過失，婦女仍將負有揭露義務。然而，諮詢規制之保護構想正要求，為使懷孕婦女能開放面對諮詢保存未出生者生命之目的，免除其在諮詢及醫師談話以外向第三人陳述其終止妊娠決定理由之必要。若勞動法以某種方式使女勞工負有向雇主揭露其工作不能源於未正當化終止妊娠之義務，此種開放性無論如何將受危及。在此等情況下，將終止妊娠所致之女勞工工作不能歸入雇主風險領域，亦非不公平。
\end{transpara}

\begin{sourcepara}[title={原文E.228}]
\sourcern{342}5. a) Die verfassungsrechtliche Schutzpflicht für das Leben hindert den Gesetzgeber nicht, Leistungen der sozialen Krankenversicherung bei einem nicht rechtswidrigen Abbruch der Schwangerschaft vorzusehen, wie dies in § 24b \abbrSGBV{} im Anschluß an § 200f \abbrRVO{} bestimmt ist. Dies gilt auch mit Bezug auf die Regelungen der §§ 218a \abbrFF{} \abbrStGB{} \abbrAF{}, die bis zum 15. Juni 1993 anzuwenden sind (vgl. II.1. der Entscheidungsformel).
\end{sourcepara}
\begin{transpara}
\transrn{342}5. a) 生命保護之憲法保護義務，並不阻礙立法者如 SGB V 第24b條接續 RVO 第200f條所規定般，就不具違法性之終止妊娠規定社會醫療保險給付。這亦適用於截至1993年6月15日仍應適用之舊版本《刑法》第218a條以下規制（參見判決主文 II.1.）。
\end{transpara}

\begin{sourcepara}[title={原文E.229}]
\sourcern{343}aa) § 24b \abbrSGBV{} begründet in weitgehender Übereinstimmung mit der Vorgängervorschrift des § 200f \abbrRVO{} Ansprüche auf Leistungen der gesetzlichen Krankenversicherung "bei einem nicht rechtswidrigen Abbruch der Schwangerschaft durch einen Arzt", wenn der Abbruch in einem Krankenhaus oder einer sonstigen hierfür vorgesehenen Einrichtung im Sinne des neugefaßten \abbrArt{} 3 \abbrAbs{} 1 Satz 1 des 5. \abbrStrRG{} vorgenommen wird. "Nicht rechtswidrig" in diesem Sinne ist auch ein Schwangerschaftsabbruch bei -- festgestelltem -- Vorliegen der allgemeinen Notlagenindikation. Zwar bezeichnet § 218a \abbrStGB{} \abbrAF{} einen solchen Abbruch lediglich als "nicht nach § 218 strafbar". Der Bundesgerichtshof hat jedoch sowohl in seiner zivilrechtlichen als auch in seiner strafrechtlichen Judikatur den Tatbestand der allgemeinen Notlagenindikation ebenso wie die anderen in § 218a \abbrStGB{} \abbrAF{} normierten Indikationstatbestände als Rechtfertigungsgrund ausgelegt (vgl. BGHZ 86, 240 [245]; 95, 199 [204 \abbrFF{}]; BGHR \abbrStGB{} § 218a \abbrAbs{} 1 Indikation 1). Das Bundessozialgericht (vgl. NJW 1985, \abbrS{} 2215 [2216]) und die sozialrechtliche Praxis sind -- ohne zu der strafrechtsdogmatischen Einordnung der Indikationstatbestände näher Stellung zu nehmen -- davon ausgegangen, daß als "nicht rechtswidrige" Schwangerschaftsabbrüche im Sinne des § 200f \abbrRVO{} alle nach § 218a \abbrStGB{} \abbrAF{} nicht strafbaren Abbrüche anzusehen seien; für die inhaltsgleiche Vorschrift des § 24b \abbrSGBV{}, soweit sie infolge der einstweiligen Anordnung vom 4. August 1992 im Rahmen des § 218a \abbrStGB{} \abbrAF{} Anwendung gefunden hat, wird dieselbe Auslegung zu gelten haben.
\end{sourcepara}
\begin{transpara}
\transrn{343}aa) SGB V 第24b條與前身規定 RVO 第200f條大致一致，於醫師所為「不具違法性之終止妊娠」時建立法定醫療保險給付請求權，前提是終止妊娠於新版本《第五刑法改革法》第3條第1項第1句意義下之醫院或其他為此設置之機構中實施。在此意義下「不具違法性」者，亦包括在一般困境指標事由已被確認存在時之終止妊娠。誠然，舊版本《刑法》第218a條僅將此種終止妊娠稱為「不依第218條處罰」。然而，聯邦普通法院在其民事法及刑事法判例中，均將一般困境指標事由構成要件，如同舊版本《刑法》第218a條所規範之其他指標事由構成要件，解釋為正當化事由（參見 BGHZ 86, 240 [245]; 95, 199 [204 ff.]; BGHR StGB § 218a Abs. 1 Indikation 1）。聯邦社會法院（參見 NJW 1985, S. 2215 [2216]）及社會法實務，在未更詳盡表態指標事由構成要件之刑法教義分類下，均以為 RVO 第200f條意義下之「不具違法性」終止妊娠，係指所有依舊版本《刑法》第218a條不罰之終止妊娠；對內容相同之 SGB V 第24b條，在其因1992年8月4日暫時處分而於舊版本《刑法》第218a條框架內適用之範圍，亦應適用相同解釋。
\end{transpara}

\begin{sourcepara}[title={原文E.230}]
\sourcern{344}bb) Von Verfassungs wegen ist diese fachgerichtliche Auslegung grundsätzlich nicht zu beanstanden, auch soweit der Tatbestand der allgemeinen Notlagenindikation betroffen ist. Bedenken dagegen, daß § 218a \abbrAbs{} 2 \abbrNr{} 3 \abbrStGB{} \abbrAF{} die in Betracht kommenden Notlagen tatbestandlich nicht näher konkretisiere, greifen im Ergebnis nicht durch. Die Notlagen sind so vielfältig, ihr Gewicht unter dem Gesichtspunkt der Unzumutbarkeit einer Fortsetzung der Schwangerschaft ist so sehr von den Umständen des Einzelfalls abhängig, daß eine bestimmtere Fassung des Indikationstatbestandes nur um den Preis erreichbar gewesen wäre, untypische Fallgestaltungen möglicherweise nicht mehr zu erfassen. Im Blick darauf durfte sich der Gesetzgeber mit einer allgemein gehaltenen Umschreibung des Tatbestandes unter Verwendung wertungsbedürftiger Begriffe begnügen und die Konkretisierung der Rechtsprechung überlassen. Auch das Erfordernis, daß die in Betracht kommenden Notlagen unter dem Gesichtspunkt der Unzumutbarkeit einer Fortsetzung der Schwangerschaft vergleichbar schwer wiegen müssen wie die übrigen Indikationsfälle, läßt sich unter Berücksichtigung der Entstehungsgeschichte der Vorschrift, zu der auch das Urteil des Bundesverfassungsgerichts vom 25. Februar 1975 gehört, sowie des Normzusammenhangs im Wege der Interpretation aus der Formulierung erschließen, die Gefahr einer Notlage müsse so schwer wiegen, "daß von der Schwangeren die Fortsetzung der Schwangerschaft nicht verlangt werden kann".
\end{sourcepara}
\begin{transpara}
\transrn{344}bb) 從憲法上看，此一專業法院解釋原則上無可指摘，即使涉及一般困境指標事由構成要件亦然。關於舊版本《刑法》第218a條第2項第3款未以構成要件方式更具體化可考慮困境之疑慮，結果上不能成立。困境如此多樣，其在繼續懷孕不可期待性觀點下之重量又如此取決於個案情況，以致只有付出可能不再涵蓋非典型案件型態之代價，始能使指標事由構成要件有更確定版本。鑑於此，立法者得滿足於使用需要評價之概念對構成要件作一般性描述，並將具體化交由判例。此外，從規定形成史，包括聯邦憲法法院1975年2月25日判決，以及規範脈絡考量，亦可透過解釋從「困境危險必須如此嚴重，以致不能要求懷孕婦女繼續懷孕」之表述中，得出可考慮困境在繼續懷孕不可期待性觀點下必須與其他指標事由案件具可比較重量之要求。
\end{transpara}

\begin{sourcepara}[title={原文E.231}]
\sourcern{345}cc) Allerdings kann als sicher gelten, daß der Tatbestand der allgemeinen Notlagenindikation in der Vergangenheit vielfach zur Rechtfertigung von Schwangerschaftsabbrüchen und damit zur Begründung von Ansprüchen auf Krankenversicherungsleistungen herangezogen worden ist, ohne daß der hier vorauszusetzende soziale Konflikt unter dem Gesichtspunkt der Unzumutbarkeit das Gewicht der anderen Indikationen erreichte. Auch die Begründung des Schwangeren- und Familienhilfegesetzes sieht in der ausufernden praktischen Handhabung einen Grund für die Neuregelung (vgl. \abbrBTDrucks{} 12/2605 [neu], \abbrS{} 3).
\end{sourcepara}
\begin{transpara}
\transrn{345}cc) 不過，可以確定的是，一般困境指標事由構成要件過去經常被用來正當化終止妊娠，並因而建立醫療保險給付請求權，而此處所預設之社會衝突，在不可期待性觀點下並未達到其他指標事由之重量。《孕婦及家庭扶助法》理由亦將實務操作之擴張視為新規制之理由（參見 BTDrucks. 12/2605 [neu], S. 3）。
\end{transpara}

\begin{sourcepara}[title={原文E.232}]
\sourcern{346}Mit der Schutzpflicht für das ungeborene Leben ist es nicht vereinbar, daß die Krankenversicherungsträger auch bei einer derart ausufernden Auslegung und Anwendung des Tatbestandes der allgemeinen Notlage die Rechtmäßigkeit des ärztlichen Eingriffs ohne weiteres unterstellen und diesen als Leistung gewähren. Sie haben sich zu vergewissern, daß die Annahme einer allgemeinen Notlage nach ärztlicher Erkenntnis nicht unvertretbar war -- insoweit können sie sich an den vom Bundesgerichtshof entwickelten Grundsätzen (vgl. BGHR \abbrStGB{} § 218a \abbrAbs{} 2 Erkenntnis 1; BGHZ 95, 199 [204 \abbrFF{}]) orientieren -- und die Vorschriften über die Beratung (§ 218b \abbrStGB{} \abbrAF{}) und das Indikationsfeststellungsverfahren (§ 219 \abbrStGB{} \abbrAF{}) nicht mißachtet wurden. Verfassungsrechtliche Mängel des Vollzugs greifen indes nicht auf die Bestimmungen des Sozialversicherungsrechts über, nach denen Versicherten bei einem nicht rechtswidrigen Schwangerschaftsabbruch Leistungen gewährt werden. Ebenso bleiben diese Bestimmungen unberührt von strukturellen Mängeln, die den Vorschriften der §§ 218b und 219 \abbrStGB{} \abbrAF{} anhaften und sie in ihrer lebensschützenden Wirkung beeinträchtigen; solche Mängel geben dem Gesetzgeber grundsätzlich nur Anlaß zur Nachbesserung.
\end{sourcepara}
\begin{transpara}
\transrn{346}醫療保險承辦者即使在如此擴張解釋及適用一般困境構成要件時，仍毫無保留地推定醫療介入之合法性並將其作為給付提供，與對未出生生命之保護義務不相容。其必須確認，依醫學認識採認一般困境並非不可支持；就此，其得以聯邦普通法院所發展之原則為取向（參見 BGHR StGB § 218a Abs. 2 Erkenntnis 1; BGHZ 95, 199 [204 ff.]），並確認諮詢規定（舊版本《刑法》第218b條）及指標事由確認程序（舊版本《刑法》第219條）未被忽視。然而，執行之憲法瑕疵並不及於社會保險法規定本身；依該等規定，被保險人在不具違法性之終止妊娠時取得給付。同樣地，此等規定亦不受舊版本《刑法》第218b條及第219條本身所附、並削弱其生命保護效果之結構性瑕疵影響；此類瑕疵原則上僅使立法者有補正理由。
\end{transpara}

\begin{sourcepara}[title={原文E.233}]
\sourcern{347}b) Die Antragstellerin zu 1) hat sich allerdings auch unter anderen verfassungsrechtlichen Gesichtspunkten dagegen gewandt, daß -- vor dem Inkrafttreten des Schwangeren- und Familienhilfegesetzes aufgrund des § 200f Satz 1 \abbrRVO{}, danach aufgrund des § 24b \abbrAbs{} 1 \abbrSGBV{} -- Ansprüche auf Leistungen der sozialen Krankenversicherung auch bei Schwangerschaftsabbrüchen aufgrund der allgemeinen Notlagenindikation (§ 218a \abbrAbs{} 2 \abbrNr{} 3 \abbrStGB{} \abbrAF{}) gewährt wurden. Sie macht geltend, diese Aufgabe habe der Sozialversicherung als einem öffentlich-rechtlichen Zwangsverband mit Rücksicht auf die Grundrechte ihrer Mitglieder nicht aufgebürdet werden dürfen; es handele sich um eine versicherungsfremde Last. Das vom Gesetzgeber verfolgte Ziel, Nachteile für schwangere Frauen zu vermeiden, die sich einerseits in einer gesetzlich anerkannten Konfliktlage, andererseits in einer wirtschaftlichen Notlage befinden, hätte mit gleicher Wirkung auch auf andere, die Grundrechte der Versicherungsmitglieder unberührt lassende Weise erreicht werden können.
\end{sourcepara}
\begin{transpara}
\transrn{347}b) 不過，第一聲請人亦從其他憲法觀點反對，在《孕婦及家庭扶助法》生效前基於 RVO 第200f條第1句，生效後基於 SGB V 第24b條第1項，於基於一般困境指標事由（舊版本《刑法》第218a條第2項第3款）之終止妊娠時，亦給予社會醫療保險給付請求權。其主張，鑑於成員之基本權，不得將此任務加諸作為公法上強制團體之社會保險；此乃保險外負擔。立法者所追求之目標，即避免處於法律承認之衝突狀態且同時處於經濟困境之懷孕婦女遭受不利益，本可用同樣效果但不觸及保險成員基本權之其他方式達成。
\end{transpara}

\begin{sourcepara}[title={原文E.234}]
\sourcern{348}Ob diesen Bedenken zu folgen ist, kann offenbleiben. Es bestünde jedenfalls keine Möglichkeit, den aufgrund der früheren Rechtslage erbrachten Leistungen rückwirkend die Rechtsgrundlage zu entziehen; angesichts der Intention des Antrags der Bayerischen Staatsregierung ist es auch nicht erforderlich, für die Zeit der Aufrechterhaltung des bisher geltenden Rechts nach Maßgabe der Nummer 1 der Anordnung gemäß § 35 \abbrBVerfGG{} diese Frage eigens zu entscheiden.
\end{sourcepara}
\begin{transpara}
\transrn{348}此等疑慮是否應被採納，可置而不論。無論如何，均不可能事後剝奪依既有法律狀態所提供給付之法律基礎；鑑於巴伐利亞邦政府聲請之意旨，亦無必要依《聯邦憲法法院法》第35條命令第1款，在維持既有法律有效期間另行決定此一問題。
\end{transpara}

\bisubsubsection{VI.}{VI.}

\begin{sourcepara}[title={原文E.235}]
\sourcern{349}Die in \abbrArt{} 4 des 5. \abbrStrRG{} in der Fassung des \abbrArt{} 15 \abbrNr{} 2 \abbrSFHG{} enthaltene Verpflichtung, ein ausreichendes und flächendeckendes Angebot sowohl ambulanter als auch stationärer Einrichtungen zur Vornahme von Schwangerschaftsabbrüchen sicherzustellen, ist bei restriktiver Auslegung mit der Kompetenzordnung des Grundgesetzes vereinbar. Die Regelung verstößt aber gegen das bundesstaatliche Prinzip und ist nichtig, soweit sie die zuständige oberste Landesbehörde als Träger der Verpflichtung benennt.
\end{sourcepara}
\begin{transpara}
\transrn{349}SFHG 第15條第2款版本之《第五刑法改革法》第4條所含義務，即確保有足夠且遍及全境之門診及住院終止妊娠機構提供，若作限縮解釋，與《基本法》權限秩序相容。但該規制在其將主管最高邦機關指定為義務主體之範圍內，違反聯邦國原則並且無效。
\end{transpara}

\begin{sourcepara}[title={原文E.236}]
\sourcern{350}1. Sieht der Bundesgesetzgeber vor, daß von staatlicher Seite ein ausreichendes und flächendeckendes Angebot sowohl ambulanter als auch stationärer Einrichtungen zur Vornahme von Schwangerschaftsabbrüchen sicherzustellen ist, so begründet er eine Staatsaufgabe.
\end{sourcepara}
\begin{transpara}
\transrn{350}1. 若聯邦立法者規定，國家方面應確保有足夠且遍及全境之門診及住院終止妊娠機構提供，則其建立一項國家任務。
\end{transpara}

\begin{sourcepara}[title={原文E.237}]
\sourcern{351}Wie den Materialien des Gesetzgebungsverfahrens zu entnehmen ist, sollen vor allem Ambulatorien entstehen, in denen beim Schwangerschaftsabbruch die für die Schwangere schonendere Absaugmethode zum Einsatz kommt und den Frauen mehrtägiger stationärer Aufenthalt erspart bleibt (vgl. Begründungen zu den im wesentlichen inhaltsgleichen Vorschriften in den Gesetzentwürfen \abbrBTDrucks{} 12/696 [\abbrS{} 11 \abbrF{}] unter Bezugnahme auf Renate Sadrozinsky, Die ungleiche Praxis des § 218, Köln 1990, \abbrS{} 43, \abbrBTDrucks{} 12/889 [\abbrS{} 12]). Der Schwangerschaftsabbruch ist, wie auch die Neufassung des \abbrArt{} 3 des 5. \abbrStrRG{} durch \abbrArt{} 15 \abbrNr{} 1 \abbrSFHG{} ausweist, nicht mehr grundsätzlich in das Krankenhaus verwiesen. Das Bundesrecht läßt jede Einrichtung zu, die die notwendige Nachbehandlung gewährleistet. Im Zusammenhalt mit dieser Rechtsänderung legt der Wortlaut des neuen \abbrArt{} 4 des 5. \abbrStrRG{} die Auslegung nahe, daß es Aufgabe des Staates sei, für ein ausreichendes Angebot an Abbrucheinrichtungen auch in der Fläche des Landes im Sinne einer Auswahlmöglichkeit zwischen stationären und ambulanten Einrichtungen zu sorgen.
\end{sourcepara}
\begin{transpara}
\transrn{351}如立法程序材料所示，尤其應設立門診機構，在其中於終止妊娠時使用對懷孕婦女較溫和之吸引方法，並使婦女免於多日住院停留（參見法律草案中內容實質相同規定之理由 BTDrucks. 12/696 [S. 11 f.]，援引 Renate Sadrozinsky, Die ungleiche Praxis des § 218, Köln 1990, S. 43；BTDrucks. 12/889 [S. 12]）。如 SFHG 第15條第1款對《第五刑法改革法》第3條之新版本亦顯示，終止妊娠不再原則上被指向醫院。聯邦法允許任何能保障必要後續治療之機構。結合此一法律變更，新版本《第五刑法改革法》第4條之文字，使下列解釋顯得自然：國家之任務在於，也在邦之全境意義下，為充分之終止妊娠機構提供照護，以使住院與門診機構間有選擇可能。
\end{transpara}

\begin{sourcepara}[title={原文E.238}]
\sourcern{352}Eine so verstandene Sicherstellung verlangt ein umfassendes Konzept jeweils für das ganze Land. Gefordert sein können flächenbezogene Erhebungen des voraussichtlichen Bedarfs und der bereits vorhandenen Einrichtungen sowie -- ähnlich wie bei der Krankenhausplanung -- eine landesweite infrastrukturelle Planung, in welche die Einrichtungen privater, frei gemeinnütziger, kommunaler oder staatlicher Träger aufzunehmen und aufeinander abzustimmen sind. Sollen Einrichtungen zum Schwangerschaftsabbruch privaten oder kommunalen Krankenhausträgern zur Pflicht gemacht werden, so bedarf es hierzu gesetzlicher Regelungen, durch die in einer rechtsstaatlichen Anforderungen genügenden Bestimmtheit Maßstäbe und Befugnisse für die erforderlichen behördlichen Anordnungen festgelegt werden.
\end{sourcepara}
\begin{transpara}
\transrn{352}如此理解之確保，要求各該全邦之全面構想。可能被要求者，是就預期需求及既有機構作區域相關調查，以及類似醫院規劃之全邦基礎設施規劃，將私人、自由公益、市鎮或國家承辦者之機構納入並相互協調。若欲使私人或市鎮醫院承辦者負有設置終止妊娠機構之義務，則為此需要法律規制，透過該等規制以符合法治國要求之明確性，確定必要行政命令之標準與權限。
\end{transpara}

\begin{sourcepara}[title={原文E.239}]
\sourcern{353}2. a) Die bundesgesetzliche Begründung einer solchen Staatsaufgabe mit der im Gesetz vorgesehenen weitreichenden Zielsetzung ist durch einen Kompetenztitel des Grundgesetzes nicht gedeckt. \abbrArt{} 74 \abbrNr{} 7 \abbrGG{}\ggfnArtSeventyFourSevenDE{}, der die konkurrierende Gesetzgebung des Bundes für "die öffentliche Fürsorge" eröffnet, bietet eine kompetenzielle Grundlage für die Vorschrift nur bei deren restriktiver Auslegung.
\end{sourcepara}
\begin{transpara}
\transrn{353}2. a) 以法律中所規定之廣泛目標，透過聯邦法律建立此種國家任務，並無《基本法》權限名義涵蓋。《基本法》第74條第7款\ggfnArtSeventyFourSevenZH{}開啟聯邦對「公共扶助」之競合立法，但僅在對該規定作限縮解釋時，方能為受攻擊規制提供權限基礎。
\end{transpara}

\begin{sourcepara}[title={原文E.240}]
\sourcern{354}aa) Der Begriff der öffentlichen Fürsorge im Sinne des Grundgesetzes ist selbst nicht eng auszulegen. Er umfaßt auch präventive Maßnahmen zum Ausgleich von Notlagen und besonderen Belastungen sowie Vorkehrungen gegen die Gefahr der Hilfsbedürftigkeit (Maunz in: Maunz/Dürig, Kommentar zum \abbrGG{}, \abbrArt{} 74 \abbrRdnr{} 106; Rengeling in: Handbuch des Staatsrechts, Band IV, 1990, § 100 \abbrRdnr{} 155). Eingrenzungen ergeben sich insbesondere bei überwiegendem Sachzusammenhang einer Regelung mit anderen Sachkompetenzen; aus der Kompetenz nach \abbrArt{} 74 \abbrNr{} 7 \abbrGG{}\ggfnArtSeventyFourSevenDE{} scheiden vor allem Gesetze aus, die der Krankenversorgung, der Seuchenbekämpfung oder in sonstiger Weise in erster Linie dem Gesundheitswesen dienen. Die Entscheidung der Verfassung (\abbrArt{} 74 \abbrNr{} 19 und 19a \abbrGG{}), dem Bund für das Gesundheitswesen nur in eingeschränktem Maße Gesetzgebungskompetenzen zuzuweisen, darf nicht durch eine erweiternde Auslegung der Gesetzgebungskompetenz für die öffentliche Fürsorge unterlaufen werden.
\end{sourcepara}
\begin{transpara}
\transrn{354}aa) 《基本法》意義下公共扶助之概念本身並不應狹義解釋。其亦包括用以補償困境與特殊負擔之預防措施，以及防止需要協助危險之預防安排（Maunz in: Maunz/Dürig, Kommentar zum GG, Art. 74 Rdnr. 106; Rengeling in: Handbuch des Staatsrechts, Band IV, 1990, § 100 Rdnr. 155）。限制尤其產生於某一規制與其他事項權限具有主要事項關聯時；依《基本法》第74條第7款\ggfnArtSeventyFourSevenZH{}之權限，尤其不包括服務於醫療照護、傳染病防治或以其他方式主要服務於健康制度之法律。憲法決定（《基本法》第74條第19款及第19a款）僅在受限程度內賦予聯邦健康制度之立法權限，不得透過擴張解釋公共扶助立法權限而被規避。
\end{transpara}

\begin{sourcepara}[title={原文E.241}]
\sourcern{355}In der dargestellten weiten Auslegung werden mit \abbrArt{} 4 des 5. \abbrStrRG{} \abbrNF{} Ziele der strukturellen Veränderung des Gesundheitswesens in den Ländern verfolgt. Dafür steht dem Bund keine Gesetzgebungsbefugnis zu.
\end{sourcepara}
\begin{transpara}
\transrn{355}依前述廣義解釋，新版本《第五刑法改革法》第4條追求各邦健康制度結構變革之目標。聯邦就此並無立法權限。
\end{transpara}

\begin{sourcepara}[title={原文E.242}]
\sourcern{356}bb) Unter dem Gesichtspunkt der Hilfe in einer schwangerschaftsbedingten Notlage kann die Kompetenz aus \abbrArt{} 74 \abbrNr{} 7 \abbrGG{}\ggfnArtSeventyFourSevenDE{} für die angegriffene Regelung allerdings eine begrenzte Tragfähigkeit gewinnen.
\end{sourcepara}
\begin{transpara}
\transrn{356}bb) 不過，從懷孕所致困境中提供協助之觀點，《基本法》第74條第7款\ggfnArtSeventyFourSevenZH{}之權限，對受攻擊規制得取得有限承載力。
\end{transpara}

\begin{sourcepara}[title={原文E.243}]
\sourcern{357}Es dient zum einen dem Lebensschutz, wenn sich der Arzt nicht wegen einer weiten Anreise der schwangeren Frau gedrängt sieht, den Schwangerschaftsabbruch an dem Tage, an dem sie sich bei ihm zum ersten Mal einfindet, vorzunehmen. Bei einer unklaren Lage wird sich der Arzt dann eher darauf beschränken, zunächst das ärztliche Gespräch mit der Frau zu führen und sie in Übereinstimmung mit der auf Lebensschutz gerichteten ärztlichen Berufspflicht zu beraten, einen etwaigen Eingriff aber auf einen späteren Tag verschieben. Damit wird noch einmal die Chance für eine Entscheidung der Frau zugunsten des Ungeborenen eröffnet.
\end{sourcepara}
\begin{transpara}
\transrn{357}一方面，若醫師不因懷孕婦女路途遙遠而覺得必須在婦女首次到診當日即實施終止妊娠，即服務於生命保護。在情況不明時，醫師將較可能先限於與婦女進行醫療談話，並依其以生命保護為取向之醫師職業義務加以諮詢，而將可能之介入延至之後日期。如此將再次開啟婦女作成有利於未出生者決定之機會。
\end{transpara}

\begin{sourcepara}[title={原文E.244}]
\sourcern{358}Zum anderen kann es in einer solchen Situation auch der Schwangeren eine Hilfe in der Not sein, wenn sie für einen ersten Arztbesuch die An- und Rückreise -- auch mit öffentlichen Verkehrsmitteln -- an einem Tag bewältigen kann. Es wird ihr leichter, die Betreuung eigener Kinder während ihrer Abwesenheit zu regeln; der Arbeit braucht sie nur für eine relativ kurze Zeit fernzubleiben.
\end{sourcepara}
\begin{transpara}
\transrn{358}另一方面，在此種情況中，若懷孕婦女為首次就醫能在一天內完成往返，即使使用公共交通工具，對其而言亦可能是困境中之協助。她較容易安排自己子女在其外出期間之照顧；工作上也只須離開相對短暫時間。
\end{transpara}

\begin{sourcepara}[title={原文E.245}]
\sourcern{359}In diesen Abmaßen kann es der Bundesgesetzgeber zum Inhalt der öffentlichen Fürsorge im Sinn des \abbrArt{} 74 \abbrNr{} 7 \abbrGG{}\ggfnArtSeventyFourSevenDE{} machen, daß -- über das Land verteilt -- Einrichtungen zum Schwangerschaftsabbruch geschaffen werden. Eine weitergehende Sicherstellung kann er jedoch nicht vorschreiben, ohne die Kompetenzgrenzen des \abbrArt{} 74 \abbrNr{} 7 \abbrGG{}\ggfnArtSeventyFourSevenDE{} zu überschreiten.
\end{sourcepara}
\begin{transpara}
\transrn{359}在此範圍內，聯邦立法者得將在各邦範圍內分散設置終止妊娠機構，作為《基本法》第74條第7款\ggfnArtSeventyFourSevenZH{}意義下公共扶助之內容。然而，超出此範圍之確保，聯邦立法者不得規定，否則即逾越《基本法》第74條第7款\ggfnArtSeventyFourSevenZH{}之權限界限。
\end{transpara}

\begin{sourcepara}[title={原文E.246}]
\sourcern{360}Dafür hat der Bundesgesetzgeber auch keine Annexkompetenz unter dem Gesichtspunkt der Verwirklichung eines von ihm für erforderlich erachteten organisatorischen Folgekonzepts (vgl. \abbrBVerfGE{} 22, 180 [209 \abbrFF{}]; 77, 288 [301]). Die Zuständigkeit zur Gesetzgebung aus \abbrArt{} 74 \abbrNr{} 1 \abbrGG{}\ggfnArtSeventyFourOneDE{} vermag ein bundesgesetzliches, auf Schutz des Ungeborenen und der Frau zugeschnittenes, im Strafrecht wurzelndes Schutzkonzept zu tragen; sie ermöglicht dem Bundesgesetzgeber so, seiner verfassungsrechtlichen Schutzpflicht für das Leben zu genügen. Demgegenüber läßt sich die Sicherstellung eines über den bereits dargestellten Rahmen hinausgehenden Angebots an Einrichtungen, in denen Schwangerschaftsabbrüche vorgenommen werden, nicht als notwendiges Folgekonzept verstehen.
\end{sourcepara}
\begin{transpara}
\transrn{360}聯邦立法者亦無基於實現其認為必要之組織後續構想觀點的附屬權限（參見 BVerfGE 22, 180 [209 ff.]; 77, 288 [301]）。依《基本法》第74條第1款\ggfnArtSeventyFourOneZH{}之立法權限，能承載一項聯邦法律上、根植於刑法且針對未出生者及婦女保護而設計之保護構想；其使聯邦立法者能履行其對生命之憲法保護義務。相對地，確保超越前述框架之終止妊娠實施機構提供，不能被理解為必要後續構想。
\end{transpara}

\begin{sourcepara}[title={原文E.247}]
\sourcern{361}cc) \abbrArt{} 4 des 5. \abbrStrRG{} \abbrNF{} kann, soweit diese Vorschrift eine Aufgabe der öffentlichen Fürsorge begründet, nicht wegen Verfassungswidrigkeit für nichtig erklärt werden, da sie einer verfassungskonformen -- engen -- Auslegung zugänglich ist.
\end{sourcepara}
\begin{transpara}
\transrn{361}cc) 新版本《第五刑法改革法》第4條，在其建立公共扶助任務之範圍內，不能因違憲而被宣告無效，因為該規定可作合憲之限縮解釋。
\end{transpara}

\begin{sourcepara}[title={原文E.248}]
\sourcern{362}Es entspricht ständiger Rechtsprechung des Bundesverfassungsgerichts, daß ein Gesetz nicht verfassungswidrig ist, wenn eine Auslegung möglich ist, die im Einklang mit dem Grundgesetz steht, und das Gesetz bei dieser Auslegung sinnvoll bleibt (vgl. \abbrBVerfGE{} 2, 266 [278]; 69, 1 [55]; st. \abbrRspr{}).
\end{sourcepara}
\begin{transpara}
\transrn{362}依聯邦憲法法院一貫判例，若一部法律有可能作出與《基本法》一致之解釋，且法律在此解釋下仍有意義，則該法律並非違憲（參見 BVerfGE 2, 266 [278]; 69, 1 [55]; st. Rspr.）。
\end{transpara}

\begin{sourcepara}[title={原文E.249}]
\sourcern{363}Mit dem Wortlaut des \abbrArt{} 4 des 5. \abbrStrRG{} \abbrNF{} ist es vereinbar, die Vorschrift dahin auszulegen, daß der Staat zur Verwirklichung des Schutzkonzepts für das Bereitstehen ärztlicher Hilfe zum Abbruch der Schwangerschaft in einer Entfernung zu sorgen hat, die von der Frau nicht die Abwesenheit über einen Tag hinaus verlangt. Bei dieser Auslegung bleibt das Gesetz sinnvoll, weil es dem Lebensschutz dienen kann. Auch die Grenzen, die der klar erkennbare Wille des Gesetzgebers einer verfassungskonformen Auslegung zieht (vgl. \abbrBVerfGE{} 18, 97 [111]; 71, 81 [105]; st. \abbrRspr{}), sind nicht überschritten, weil das Gesetz auch in den Abmaßen der einschränkenden Auslegung diesem Willen jedenfalls entgegenkommt.
\end{sourcepara}
\begin{transpara}
\transrn{363}與新版本《第五刑法改革法》第4條文字相容的是，將該規定解釋為：國家為實現保護構想，必須確保在一段不要求婦女離開超過一天之距離內，備有用於終止妊娠之醫療協助。在此解釋下，法律仍有意義，因為其得服務於生命保護。合憲解釋受立法者明確可辨之意志所劃定之界限（參見 BVerfGE 18, 97 [111]; 71, 81 [105]; st. Rspr.）亦未被逾越，因為即使在限縮解釋之範圍內，法律仍至少符合此一意志。
\end{transpara}

\begin{sourcepara}[title={原文E.250}]
\sourcern{364}3. Beauftragt der Bundesgesetzgeber die zuständige oberste Landesbehörde mit der Ausführung einer von ihm begründeten Staatsaufgabe, so greift er in die Organisationsgewalt der Länder und damit in deren verfassungsmäßige Ordnung ein. Erschließt damit die nach Landesverfassungsrecht zuständigen Organe aus und hindert die Länder damit, die Ausführung der Staatsaufgabe -- etwa im kommunalen Bereich -- nach eigenem, auf sinnvolle Gestaltung der Vollzugsorganisation ausgerichtetem Ermessen zu regeln.
\end{sourcepara}
\begin{transpara}
\transrn{364}3. 若聯邦立法者委託主管最高邦機關執行其所建立之國家任務，即介入各邦組織權，因而介入其憲法秩序。其由此排除依邦憲法有管轄權之機關，並阻礙各邦依其自身、以合理形塑執行組織為取向之裁量，規範國家任務之執行，例如在市鎮領域中。
\end{transpara}

\begin{sourcepara}[title={原文E.251}]
\sourcern{365}a) Dem bundesstaatlichen Prinzip entsprechend ist ein Eingriff der Bundesgewalt in die Verfassungsordnung der Länder nur zulässig, soweit es das Grundgesetz ausdrücklich bestimmt oder zuläßt. Vor allem dort, wo das Grundgesetz die Grenze zwischen den Kompetenzen des Bundes und denen der Länder zieht, verzichtet es grundsätzlich darauf, die Verfassungsorgane der Länder zu bestimmen, die die Landeskompetenz wahrzunehmen haben (vgl. \abbrBVerfGE{} 11, 77 [85 \abbrF{}]).
\end{sourcepara}
\begin{transpara}
\transrn{365}a) 依聯邦國原則，聯邦權力介入各邦憲法秩序，只有在《基本法》明文規定或允許之範圍內始為容許。尤其在《基本法》劃定聯邦權限與各邦權限界限之處，其原則上放棄決定應由各邦何種憲法機關行使邦權限（參見 BVerfGE 11, 77 [85 f.]）。
\end{transpara}

\begin{sourcepara}[title={原文E.252}]
\sourcern{366}So ist es, wenn Bundesgesetze als eigene Angelegenheit ausgeführt werden (\abbrArt{} 83 \abbrGG{}\ggfnArtEightyThreeDE{}): Die Länder -- nicht die Landesregierungen oder einzelne Landesministerien -- sind es, die die Einrichtung der Behörden und das Verfahren zu regeln haben, soweit nicht Bundesgesetze mit Zustimmung des Bundesrates etwas anderes bestimmen (\abbrArt{} 84 \abbrAbs{} 1 \abbrGG{}\ggfnArtEightyFourOneDE{}). Der Grundsatz der Organisationsgewalt der Länder gilt uneingeschränkt, wenn eine bundesgesetzliche Regelung lediglich eine von den Ländern zu erfüllende Staatsaufgabe vorsieht, nicht jedoch Einzelregelungen trifft, die behördlich-administrativ vollzogen werden könnten.
\end{sourcepara}
\begin{transpara}
\transrn{366}聯邦法律作為各邦自身事務執行時即是如此（《基本法》第83條\ggfnArtEightyThreeZH{}）：除非聯邦法律經聯邦參議院同意另有規定，應規範機關設置及程序者，是各邦，而非邦政府或個別邦部會（《基本法》第84條第1項\ggfnArtEightyFourOneZH{}）。若聯邦法律規制僅規定一項由各邦履行之國家任務，而未制定可由行政機關執行之個別規制，則各邦組織權原則毫無限制地適用。
\end{transpara}

\begin{sourcepara}[title={原文E.253}]
\sourcern{367}Auch als Ermächtigung zum Erlaß einer Rechtsverordnung könnte der Bundesgesetzgeber nach der ausdrücklichen Anordnung des \abbrArt{} 80 \abbrAbs{} 1 Satz 1 \abbrGG{}\ggfnArtEightyOneOneDE{} die Aufgabe nur den "Landesregierungen", nicht aber den Landesministern oder den zuständigen obersten Landesbehörden zuweisen (vgl. \abbrBVerfGE{} 11, 77 [86]).
\end{sourcepara}
\begin{transpara}
\transrn{367}即使作為發布法規命令之授權，聯邦立法者依《基本法》第80條第1項第1句\ggfnArtEightyOneOneZH{}之明文命令，亦只能將任務賦予「邦政府」，而非邦部長或主管最高邦機關（參見 BVerfGE 11, 77 [86]）。
\end{transpara}

\begin{sourcepara}[title={原文E.254}]
\sourcern{368}b) \abbrArt{} 4 des 5. \abbrStrRG{} \abbrNF{} wird diesem Maßstab nicht gerecht. Er begründet eine Sicherstellungsaufgabe, die für das Land lediglich ein Handlungsziel vorschreibt, nicht aber behördlich-administrativ vollziehbare Einzelregelungen zur Verwirklichung dieses Ziels enthält. Die bundesgesetzlich bestimmte Sicherstellungsaufgabe ist nicht den Ländern als solchen zugewiesen, sondern den zuständigen obersten Landesbehörden und damit den Landesministern. Eine solche Regelung ist durch das Grundgesetz weder ausdrücklich bestimmt noch zugelassen.
\end{sourcepara}
\begin{transpara}
\transrn{368}b) 新版本《第五刑法改革法》第4條不符合此一標準。其建立一項確保任務，僅對邦規定一項行動目標，但並未包含可由行政機關執行、用以實現此一目標之個別規制。聯邦法律所定之確保任務，並非指派給各邦本身，而是指派給主管最高邦機關，因而指派給邦部長。此種規制並非《基本法》明文規定或允許。
\end{transpara}

\begin{sourcepara}[title={原文E.255}]
\sourcern{369}c) Insoweit ist die Vorschrift gemäß § 78 \abbrBVerfGG{} für nichtig zu erklären. \abbrArt{} 4 des 5. \abbrStrRG{} \abbrNF{} ist einer verfassungskonformen Auslegung nicht zugänglich, soweit dadurch die zuständige oberste Landesbehörde mit der Wahrnehmung der Staatsaufgabe beauftragt wird. Das Gebot verfassungskonformer Auslegung legitimiert nicht dazu, Wortlaut und Sinn des Gesetzes beiseite zu schieben oder zu verändern (vgl. \abbrBVerfGE{} 8, 28 [34]; 72, 278 [295]). Die Bestimmung hat einen präzise formulierten, normativen Gehalt, der es nicht zuläßt, ihr im Wege der Auslegung den Sinn zuzumessen, daß jeweils das Land, nicht die Behörde mit der Wahrnehmung der Aufgabe betraut sei.
\end{sourcepara}
\begin{transpara}
\transrn{369}c) 就此而言，該規定應依《聯邦憲法法院法》第78條宣告無效。新版本《第五刑法改革法》第4條，在其委託主管最高邦機關履行國家任務之範圍內，無法作合憲解釋。合憲解釋命令並不使人有權將法律文字與意義置於一旁或加以改變（參見 BVerfGE 8, 28 [34]; 72, 278 [295]）。該規定具有精確表述之規範內容，不容以解釋方式賦予其下列意義：各該邦，而非機關，受委託履行該任務。
\end{transpara}

\begin{sourcepara}[title={原文E.256}]
\sourcern{370}4. Die Nichtigkeit der Benennung der zuständigen obersten Landesbehörden als Aufgabenträger ergreift nicht den -- verfassungskonform interpretierten -- Sicherstellungsauftrag. Nach dem schon in § 139 BGB enthaltenen Rechtsgedanken kann sich das Bundesverfassungsgericht auf die Feststellung der Nichtigkeit eines Teils der Norm beschränken, wenn es keinem Zweifel unterliegt, daß der Gesetzgeber die sonstige gesetzliche Regelung auch ohne den verfassungswidrigen Teil aufrechterhalten hätte (vgl. \abbrBVerfGE{} 4, 219 [250]). Dies ist angesichts der Bedeutung, die dem -- in ärztlicher Verantwortung beratenden und gegebenenfalls den Abbruch vornehmenden -- Arzt zukommt, der Fall. Die Regelung ist auch vollziehbar, weil nach der Kompetenzverteilung des Grundgesetzes ohnehin die Länder für die Umsetzung und Ausführung der Bundesgesetze zuständig sind (\abbrArt{} 30, 70, 83 \abbrGG{}) und es insofern einer besonderen Zuständigkeitsregelung nicht bedarf.
\end{sourcepara}
\begin{transpara}
\transrn{370}4. 指定主管最高邦機關作為任務主體之無效，並不波及經合憲解釋之確保任務。依已包含於《民法》第139條之法律思想，若毫無疑問可認為立法者即使無違憲部分仍會維持其他法律規制，聯邦憲法法院得限於確認規範一部分無效（參見 BVerfGE 4, 219 [250]）。鑑於醫師，即在醫療責任下進行諮詢並於必要時實施終止妊娠者，所具有之意義，此處即屬此種情形。該規制亦可執行，因為依《基本法》權限分配，各邦本即負責聯邦法律之轉化與執行（《基本法》第30條\ggfnArtThirtyZH{}、第70條、第83條），就此無須特別管轄規制。
\end{transpara}

\begin{sourcepara}[title={原文E.257}]
\sourcern{371}5. Mithin obliegt es den Ländern, im Rahmen der durch das ärztliche Weigerungsrecht (\abbrArt{} 2 des 5. \abbrStrRG{}) eingeschränkten Möglichkeiten und der durch die verfassungskonforme Begrenzung des Sicherstellungsauftrags gesetzten Schranken für die notwendige ärztliche Betreuung der Schwangeren zu sorgen. Andererseits sind sie bei Inanspruchnahme ihrer Kompetenz für das Gesundheitswesen durch die Verpflichtung, ungeborenes menschliches Leben zu schützen, verfassungsrechtlich gebunden; sie haben zusätzliche Maßnahmen zu unterlassen, wenn sich diese als aktive Förderung des Schwangerschaftsabbruchs auswirken.
\end{sourcepara}
\begin{transpara}
\transrn{371}5. 因此，各邦負有義務，在醫師拒絕權（《第五刑法改革法》第2條）所限制之可能範圍內，以及合憲限制確保任務所設定之界限內，照顧懷孕婦女之必要醫療照護。另一方面，各邦在行使其健康制度權限時，亦受保護未出生人類生命之義務的憲法拘束；若額外措施發生積極促進終止妊娠之效果，各邦即應不為該等措施。
\end{transpara}

\bikeepwithnext[6]
\begin{sourcepara}[title={原文E.258}]
F.
\end{sourcepara}
\begin{transpara}
F.
\end{transpara}

\begin{sourcepara}[title={原文E.259}]
\sourcern{372}1. Soweit sich der Normenkontrollantrag der Bayerischen Staatsregierung im Verfahren 2 \abbrBvF{} 2/90 gegen die Vorschriften des § 218b \abbrAbs{} 1 Satz 1 und \abbrAbs{} 2 und des § 219 \abbrAbs{} 1 Satz 1 \abbrStGB{} in der Fassung des Fünfzehnten Strafrechtsänderungsgesetzes richtet, ist er durch die Entscheidung über die Normenkontrollanträge gegen die Vorschriften des Schwangeren- und Familienhilfegesetzes erledigt. Ein Antrag auf abstrakte Normenkontrolle ist nur zulässig bei Vorliegen eines besonderen objektiven Interesses an der Klarstellung der Geltung der Norm (vgl. \abbrBVerfGE{} 6, 104 [110]; 52, 63 [80]). Dieses Interesse ist hier entfallen.
\end{sourcepara}
\begin{transpara}
\transrn{372}1. 巴伐利亞邦政府於 2 BvF 2/90 程序中之規範審查聲請，若其針對第十五刑法修正法版本之《刑法》第218b條第1項第1句及第2項，以及第219條第1項第1句之規定，已因對《孕婦及家庭扶助法》規定之規範審查聲請所作裁判而解決。抽象規範審查聲請，只有在存在澄清規範效力之特殊客觀利益時，始為合法（參見 BVerfGE 6, 104 [110]; 52, 63 [80]）。此一利益在此已不存在。
\end{transpara}

\begin{sourcepara}[title={原文E.260}]
\sourcern{373}Die Antragstellerin hat die bezeichneten Vorschriften über die Schwangerschaftsberatung und über die Indikationsfeststellung nur insoweit angegriffen, als diese die verfahrensrechtlichen Voraussetzungen für einen Schwangerschaftsabbruch aufgrund der allgemeinen Notlagenindikation (§ 218a \abbrAbs{} 2 \abbrNr{} 3 \abbrStGB{} in der Fassung des Fünfzehnten Strafrechtsänderungsgesetzes) regeln; das ergibt sich aus der Antragsbegründung vom 28. Februar 1990. Der Tatbestand der allgemeinen Notlagenindikation wird indes durch die insoweit verfassungsrechtlich nicht beanstandete Vorschrift des \abbrArt{} 13 \abbrNr{} 1 \abbrSFHG{} aufgehoben. Zwar ist die Vorschrift des § 218a \abbrAbs{} 1 \abbrStGB{} \abbrNF{}, die die allgemeine Notlagenindikation ersetzen soll, verfassungswidrig und nichtig. Das ändert aber nichts an der grundsätzlichen verfassungsrechtlichen Zulässigkeit einer auf einem Beratungskonzept beruhenden gesetzlichen Regelung des Schwangerschaftsabbruchs. Der Senat hat für die Zeit bis zu einer verfassungskonformen Neuregelung durch den Gesetzgeber eine Übergangsregelung nach § 35 \abbrBVerfGG{} in der Einschätzung angeordnet, daß der Gesetzgeber aus den Erwägungen, die ihn zu der Gesetzesänderung veranlaßt haben, nicht zu einer Notlagenindikation zurückkehren wird. Die Vorschriften darüber bleiben nur noch für eine kurze Übergangsfrist bis 15. Juni 1993 anwendbar. Damit entfällt für die Zukunft der Gegenstand des verfassungsrechtlichen Angriffs der Antragstellerin.
\end{sourcepara}
\begin{transpara}
\transrn{373}聲請人僅在此等關於懷孕諮詢及指標事由確認之規定，規範基於一般困境指標事由（第十五刑法修正法版本《刑法》第218a條第2項第3款）之終止妊娠的程序法前提範圍內，攻擊該等規定；此由1990年2月28日聲請理由可得知。然而，一般困境指標事由構成要件，由 SFHG 第13條第1款之規定廢止；該規定就此在憲法上未受指摘。誠然，意圖取代一般困境指標事由之新版本《刑法》第218a條第1項規定違憲且無效。但這並不改變以諮詢構想為基礎之終止妊娠法律規制在原則上合憲容許。第二庭依《聯邦憲法法院法》第35條就立法者制定合憲新規制前之期間命令一項過渡規制，其評估為：立法者基於促使其修法之考量，不會回到困境指標事由。關於該等指標事由之規定僅在至1993年6月15日之短暫過渡期間仍可適用。由此，聲請人憲法攻擊之標的就未來而言消失。
\end{transpara}

\begin{sourcepara}[title={原文E.261}]
\sourcern{374}Für die Zeit bis 15. Juni 1993 könnte die beantragte Feststellung der Verfassungswidrigkeit ebenfalls keine rechtlichen Wirkungen mehr äußern. Die Antragstellerin hat zutreffend dargelegt, daß die von ihr erhobenen Beanstandungen -- ihre Begründetheit unterstellt -- nicht dazu führen können, die angegriffenen Vorschriften, soweit sie sich auf Schwangerschaftsabbrüche aufgrund der Notlagenindikation beziehen, rückwirkend für nichtig zu erklären. Denn dies hätte die Unanwendbarkeit der Vorschriften bei Notlagenindikationen und damit einen Rechtszustand zur Folge, der sich in seiner Schutzwirkung noch weiter vom Grundgesetz entfernte als die angegriffene Regelung. Mithin käme nur eine Verpflichtung des Gesetzgebers zur Nachbesserung in Betracht, bis zu der die angegriffenen Vorschriften unverändert fortgelten müßten (vgl. \abbrBVerfGE{} 61, 319 [356]); eine solche in die Zukunft wirkende Nachbesserung aber scheidet wegen des kurzfristig bevorstehenden Auslaufens der alten Notlagenregelung aus.
\end{sourcepara}
\begin{transpara}
\transrn{374}至1993年6月15日止期間，所聲請之違憲確認亦已不能再發生法律效果。聲請人正確說明，其所提出之指摘，即便假定有理由，亦不能導致受攻擊規定在其涉及基於困境指標事由之終止妊娠範圍內，被溯及宣告無效。因為這將導致該等規定於困境指標事由中不可適用，因而造成一種其保護效果比受攻擊規制更遠離《基本法》之法律狀態。因此，僅可能考慮課予立法者補正義務，於補正前受攻擊規定必須原封不動繼續有效（參見 BVerfGE 61, 319 [356]）；但由於舊困境規制即將短期內屆滿，此種向未來發生作用之補正已被排除。
\end{transpara}

\begin{sourcepara}[title={原文E.262}]
\sourcern{375}2. Anders verhält es sich mit dem Antrag, die §§ 200f, 200g \abbrRVO{} für nichtig zu erklären, soweit diese den Versicherten einen Anspruch auf Leistungen der gesetzlichen Krankenversicherung auch bei Schwangerschaftsabbrüchen aufgrund der allgemeinen Notlagenindikation gewähren. Zwar sind die genannten Vorschriften durch \abbrArt{} 3 \abbrSFHG{} mit Wirkung vom 5. August 1992 aufgehoben und durch § 24b \abbrSGBV{} in der Fassung des \abbrArt{} 2 \abbrSFHG{} ersetzt worden. Die Antragstellerin hat indessen dargelegt, daß noch Krankenversicherungsleistungen für Altfälle, in denen Schwangerschaftsabbrüche aufgrund der Notlagenindikation vorgenommen wurden, nach den angegriffenen Vorschriften der Reichsversicherungsordnung abzurechnen und darüber auch Rechtsstreitigkeiten anhängig sind, für deren Entscheidung es auf die Gültigkeit der Vorschriften ankommen kann. Damit besteht das erforderliche objektive Interesse an einer Sachentscheidung fort (vgl. \abbrBVerfGE{} 5, 25 [28]; 79, 311 [327]).
\end{sourcepara}
\begin{transpara}
\transrn{375}2. 關於宣告 RVO 第200f條、第200g條無效之聲請，若該等規定亦於基於一般困境指標事由之終止妊娠時給予被保險人法定醫療保險給付請求權，情形則不同。固然，上述規定已由 SFHG 第3條自1992年8月5日起廢止，並由 SFHG 第2條版本之 SGB V 第24b條取代。然而，聲請人已說明，仍有舊案之醫療保險給付須依受攻擊之《帝國保險條例》規定結算；該等舊案中，終止妊娠係基於困境指標事由而實施，且相關法律爭訟亦仍繫屬，而其裁判可能取決於規定效力。由此，作成實體裁判所需客觀利益仍存在（參見 BVerfGE 5, 25 [28]; 79, 311 [327]）。
\end{transpara}

\begin{sourcepara}[title={原文E.263}]
\sourcern{376}In der Sache gilt hier das gleiche wie für § 24b \abbrSGBV{}, soweit diese Vorschrift Schwangerschaftsabbrüche erfaßt, die seit dem 5. August 1992 aufgrund der allgemeinen Notlagenindikation durchgeführt worden sind (vgl. dazu oben E.V.1. und 5.).
\end{sourcepara}
\begin{transpara}
\transrn{376}實體上，此處適用與 SGB V 第24b條相同之結論，只要該規定涵蓋自1992年8月5日起基於一般困境指標事由所實施之終止妊娠（對此參見上文 E.V.1. 及 5.）。
\end{transpara}

\bikeepwithnext[6]
\begin{sourcepara}[title={原文E.264}]
G.
\end{sourcepara}
\begin{transpara}
G.
\end{transpara}

\begin{sourcepara}[title={原文E.265}]
\sourcern{377}Im Verfahren der abstrakten Normenkontrolle erklärt das Bundesverfassungsgericht nach § 78 \abbrAbs{} 1 Satz 1 \abbrBVerfGG{} das zu prüfende Gesetz, wenn es mit dem Grundgesetz nicht vereinbar ist, für nichtig. Damit wird ausgesprochen, daß das Gesetz nicht seine gewollten Wirkungen haben kann. Demgemäß bewirkt die Nichtigerklärung des § 218a \abbrAbs{} 1 \abbrStGB{} \abbrNF{}, daß die Vorschrift ihre Wirkung als Rechtfertigungsgrund nicht entfaltet. Der für nichtig erklärte § 219 \abbrStGB{} \abbrNF{} kann nicht als Maßstab für Inhalt und Durchführung der Beratung dienen.
\end{sourcepara}
\begin{transpara}
\transrn{377}在抽象規範審查程序中，聯邦憲法法院依《聯邦憲法法院法》第78條第1項第1句，於待審查法律與《基本法》不相容時，宣告其無效。由此即表明，該法律不能具有其所欲之效果。相應地，宣告新版本《刑法》第218a條第1項無效之效果，是該規定不發生作為正當化事由之效果。被宣告無效之新版本《刑法》第219條，不能作為諮詢內容及實施之標準。
\end{transpara}

\begin{sourcepara}[title={原文E.266}]
\sourcern{378}Die Vorschriften der §§ 218a \abbrAbs{} 1 \abbrStGB{} \abbrNF{} und des § 219 \abbrStGB{} \abbrNF{} stehen insofern mit dem Straftatbestand des § 218 \abbrStGB{} \abbrNF{} in engem sachlichen Zusammenhang, als der Gesetzgeber in Ausführung des \abbrArt{} 31 \abbrAbs{} 4 EV für die Zeit der ersten zwölf Wochen einer Schwangerschaft den Lebensschutz auf die Wirksamkeit eines Beratungskonzepts gründen und den Schwangerschaftsabbruch unter den Voraussetzungen des § 218a \abbrAbs{} 1 \abbrStGB{} \abbrNF{} jedenfalls von der Strafbarkeit (\abbrArt{} 103 \abbrAbs{} 2 \abbrGG{}\ggfnArtOneZeroThreeTwoDE{}) ausnehmen wollte. Um zu vermeiden, daß das Schutzkonzept infolge der Nichtigerklärung des § 218a \abbrAbs{} 1 \abbrStGB{} \abbrNF{} und des § 219 \abbrStGB{} \abbrNF{} auch jene vom Gesetzgeber auch für das in \abbrArt{} 3 EV genannte Gebiet gewollten Wirkungen verfehlt, die die verfassungsrechtliche Schutzpflicht zuläßt oder -- im Hinblick auf einen tatsächlich wirksamen Lebensschutz - sogar verlangt, ist es geboten, im Wege einer Anordnung nach § 35 \abbrBVerfGG{} übergangsweise eine verfassungsrechtlich ausreichende Beratungsregelung zu treffen, die unter den übrigen -- vom Gesetzgeber in § 218a \abbrAbs{} 1 \abbrStGB{} \abbrNF{} bestimmten -- Voraussetzungen den Ausschluß des Straftatbestandes des § 218 \abbrStGB{} \abbrNF{} bewirkt. \abbrArt{} 103 \abbrAbs{} 2 und \abbrArt{} 104 \abbrAbs{} 1 \abbrGG{}\ggfnArtOneZeroFourOneDE{} stehen dem nicht entgegen. In den strafrechtlichen Vorschriften des \abbrArt{} 13 \abbrNr{} 1 \abbrSFHG{} sind die mit Strafe bedrohten Fälle des Schwangerschaftsabbruchs tatbestandlich umschrieben und damit Voraussetzungen und Grenzen der Strafbarkeit des Schwangerschaftsabbruchs gesetzlich festgelegt. Die Nichtigerklärung des in § 218a \abbrAbs{} 1 \abbrStGB{} \abbrNF{} normierten Rechtfertigungsgrundes läßt die auf der Anordnung des Gesetzgebers beruhende Strafbarkeit derjenigen Schwangerschaftsabbrüche unberührt, die nicht von § 218a \abbrAbs{} 1 \abbrStGB{} \abbrNF{} oder einer sonstigen die Strafbarkeit ausschließenden Vorschrift tatbestandlich erfaßt werden. Die Entscheidung des Senats dehnt die Strafbarkeit nicht über die vom Gesetzgeber gezogenen Grenzen hinaus aus. Die unter \abbrNr{} II.2. der Urteilsformel getroffene Anordnung nach § 35 \abbrBVerfGG{} stellt im Gegenteil sicher, daß die in § 218a \abbrAbs{} 1 \abbrStGB{} \abbrNF{} tatbestandlich erfaßten Schwangerschaftsabbrüche ungeachtet der Nichtigerklärung der Vorschrift bis zu einer gesetzlichen Neuregelung von der Strafdrohung des § 218 \abbrStGB{} \abbrNF{} ausgenommen bleiben; die Bedeutung dieser gerichtlichen Anordnung erschöpft sich in strafrechtlicher Hinsicht darin, daß der Ausschluß der Strafbarkeit nicht mehr durch einen Rechtfertigungsgrund, sondern mittels eines Tatbestandsausschlusses erfolgt. Für die nicht nach der Beratungsregelung vorgenommenen, gemäß \abbrArt{} 13 \abbrNr{} 1 \abbrSFHG{} mit Strafe bedrohten Schwangerschaftsabbrüche ist Grundlage der künftigen Strafbarkeit das Gesetz, nicht etwa die auf § 35 \abbrBVerfGG{} beruhende Anordnung des Senats. Dem speziellen Gesetzesvorbehalt der \abbrArt{} 103 \abbrAbs{} 2, 104 \abbrAbs{} 1 \abbrGG{} ist damit genügt, auch soweit das nunmehr in seinem strafrechtlichen Teil in Kraft tretende Schwangeren- und Familienhilfegesetz weitergehende Strafvorschriften enthält als das bisher in den neuen Ländern fortgeltende Recht der Deutschen Demokratischen Republik.
\end{sourcepara}
\begin{transpara}
\transrn{378}新版本《刑法》第218a條第1項及新版本《刑法》第219條之規定，與新版本《刑法》第218條之犯罪構成要件，在事物上具有緊密關聯；因為立法者為執行 EV 第31條第4項，意欲在妊娠最初十二週期間將生命保護建立於諮詢構想之有效性，並使新版本《刑法》第218a條第1項前提下之終止妊娠至少免於可罰性（《基本法》第103條第2項\ggfnArtOneZeroThreeTwoZH{}）。為避免因宣告新版本《刑法》第218a條第1項及新版本《刑法》第219條無效，使保護構想亦無法達成立法者對 EV 第3條所稱地區所欲之效果，而此等效果為憲法保護義務所容許，或鑑於實際有效生命保護甚至所要求，故有必要以依《聯邦憲法法院法》第35條之命令，過渡性制定一項憲法上足夠之諮詢規制；其在其他由立法者於新版本《刑法》第218a條第1項所定前提下，產生排除新版本《刑法》第218條犯罪構成要件之效果。《基本法》第103條第2項\ggfnArtOneZeroThreeTwoZH{}及第104條第1項\ggfnArtOneZeroFourOneZH{}並不與此相牴觸。SFHG 第13條第1款之刑法規定，以構成要件方式描述受刑罰威嚇之終止妊娠案件，因而以法律確定終止妊娠可罰性之前提與界限。宣告新版本《刑法》第218a條第1項所規範之正當化事由無效，並不影響基於立法者命令而存在之下列終止妊娠可罰性：其未被新版本《刑法》第218a條第1項或其他排除可罰性之規定以構成要件方式涵蓋。第二庭裁判並未將可罰性擴張至立法者所劃定界限之外。判決主文第 II.2. 款依《聯邦憲法法院法》第35條所作命令，反而確保新版本《刑法》第218a條第1項以構成要件方式涵蓋之終止妊娠，儘管該規定被宣告無效，在法律新規制前仍被排除於新版本《刑法》第218條之刑罰威嚇之外；此一法院命令在刑法上之意義，僅止於排除可罰性不再透過正當化事由，而是透過構成要件排除發生。對於非依諮詢規制實施，且依 SFHG 第13條第1款受刑罰威嚇之終止妊娠，未來可罰性之基礎是法律，而非第二庭基於《聯邦憲法法院法》第35條所作命令。因此，即使現已在刑法部分生效之《孕婦及家庭扶助法》包含較迄今在新邦繼續有效之德意志民主共和國法律更廣泛之刑罰規定，亦已滿足《基本法》第103條第2項\ggfnArtOneZeroThreeTwoZH{}、第104條第1項\ggfnArtOneZeroFourOneZH{}之特殊法律保留。
\end{transpara}

\begin{sourcepara}[title={原文E.267}]
Mahrenholz, Böckenförde, Klein, Graßhof, Kruis, Kirchhof, Winter, Sommer
\end{sourcepara}
\begin{transpara}
Mahrenholz, Böckenförde, Klein, Graßhof, Kruis, Kirchhof, Winter, Sommer
\end{transpara}

\bisection{Abweichende Meinung Mahrenholz/Sommer}{Mahrenholz/Sommer 不同意見}

\begin{sourcepara}[title={原文SondervotumMS.268}]
Abweichende Meinung der Richter Vizepräsident Mahrenholz und Sommer zum Urteil des Zweiten Senats vom 28. Mai 1993 -- 2 \abbrBvF{} 2/90 und 4, 5/92 --
\end{sourcepara}
\begin{transpara}
副院長 Mahrenholz 法官與 Sommer 法官就第二庭1993年5月28日判決 -- 2 BvF 2/90 及 4, 5/92 -- 之不同意見
\end{transpara}

\begin{sourcepara}[title={原文SondervotumMS.269}]
\sourcern{379}Rechtliche Regelungen des Schwangerschaftsabbruchs ergreifen den innersten Bereich menschlichen Lebens und berühren grundlegende Fragen menschlicher Existenz. Zu den spezifischen Grundbedingungen menschlichen Seins gehört, daß Sexualität und Kinderwunsch nicht übereinstimmen. Die Folgen dieser Divergenz haben die Frauen zu tragen. Zu allen Zeiten und in allen, von verschiedenen moralischen und religiösen Wertvorstellungen geprägten Kulturen haben sie Auswege aus der Not ungewollter Schwangerschaften gesucht und gefunden und sich dabei durch Androhung auch schwerster und grausamer Strafen nicht davon abhalten lassen, selbst bei Gefahr für das eigene Leben das ungeborene Leben abzutöten, wenn sie ein Kind nicht wollten. Entsprechend ihrer gewandelten gesellschaftlichen Stellung lösen Frauen heute diesen fundamentalen Lebenskonflikt vornehmlich danach, ob sie nach eigener Einschätzung in ihrer konkreten Lebenssituation die Möglichkeit erkennen, die Aufgaben einer Mutter verantwortlich zu erfüllen.
\end{sourcepara}
\begin{transpara}
\transrn{379}終止妊娠之法律規制觸及人類生命之最內在領域，並涉及人類存在之基本問題。人類存在之特定基本條件之一，是性與生育願望並不一致。此一分歧之後果由婦女承擔。在所有時代及所有受不同道德與宗教價值觀形塑之文化中，婦女均曾尋求並找到擺脫非意願懷孕困境之道路，且即使面對最嚴厲、最殘酷刑罰之威嚇，仍未被阻止；若其不願有子女，即使冒著自身生命危險，亦會殺害未出生生命。依其已變遷之社會地位，今日婦女解決此一根本生活衝突時，主要取決於其依自身評估，在具體生活處境中是否看見負責任履行母親任務之可能。
\end{transpara}

\begin{sourcepara}[title={原文SondervotumMS.270}]
\sourcern{380}Jede Regelung des Schwangerschaftsabbruchs wirft die Frage nach dem Bereich unantastbarer Autonomie des Menschen einerseits und dem Recht des Staates zu Regelungen andererseits auf; der Gesetzgeber befindet sich hier an der Grenze der Regelungsfähigkeit eines Lebensbereichs überhaupt. Er kann sich der Problematik des Schwangerschaftsabbruchs mit einer besseren oder mit einer schlechteren Regelung nähern; "lösen" kann er sie nicht; dem Staat ist hier die Selbstgewißheit zur "richtigen" Gesetzgebung verlorengegangen. Dies zeigt die Länge des Gesetzgebungsprozesses, der wenig mehr als eineinhalb Jahrzehnte nach der letzten grundlegenden Reform begann, und spiegelt sich wider in der Dauer der Beratungen des Senats, der die vom Gesetzgeber aus der Schutzpflicht zu ziehenden Folgerungen anders beurteilt als der Erste Senat im Urteil vom 25. Februar 1975 (\abbrBVerfGE{} 39, 1 \abbrFF{}). Die Maßgaben des Urteils machen es möglich, daß der Gesetzgeber auf dem Weg der Beratungslösung fortschreiten kann, den er mit dem Schwangeren- und Familienhilfegesetz eingeschlagen hat.
\end{sourcepara}
\begin{transpara}
\transrn{380}任何終止妊娠規制均提出一方面人之不可侵犯自主領域，與另一方面國家規制權之問題；立法者在此處於某一生活領域究竟可受規制之界限。立法者能以較佳或較差之規制接近終止妊娠問題；但不能「解決」該問題；國家在此已喪失對「正確」立法之自我確信。這由立法過程之長度可見；該過程在上一次根本改革後僅略多於十五年即開始。這也反映於第二庭審議之期間；第二庭對立法者應由保護義務得出之結論，作出與第一庭1975年2月25日判決（BVerfGE 39, 1 ff.）不同之評價。本判決之準則使立法者能在其透過《孕婦及家庭扶助法》所踏上之諮詢方案道路上繼續前進。
\end{transpara}

\begin{sourcepara}[title={原文SondervotumMS.271}]
\sourcern{381}Wie für den Senat, so steht auch für uns die verfassungsrechtliche Schutzpflicht des Staates für das ungeborene menschliche Leben von dessen Beginn an außer Frage. Wir sind mit dem Senat der Auffassung, daß die Schutzpflicht es dem Gesetzgeber nicht verwehrt, zu einem Schutzkonzept überzugehen, das den Schwerpunkt auf die Beratung der schwangeren Frau legt und auf eine indikationsbestimmte Strafdrohung und die Feststellung von Indikationstatbeständen verzichtet. Wir sind aber der Meinung, daß die vom Senat -- als gesetzgeberische Möglichkeit -- anerkannte faktische Letztverantwortung der Frau für den Abbruch in der Frühphase, nachdem sie sich hat beraten lassen, vom Grundrechtsstatus der Frau geboten und die Schutzpflicht insoweit begrenzt ist (I.). Unseres Erachtens gebietet daneben ein wirksamer Schutz für das ungeborene Leben in der Beratungsregelung gerade eine normative Verdeutlichung über das Erlaubte und Unerlaubte eines Schwangerschaftsabbruchs; dabei läßt das Grundgesetz eine Rechtfertigung des Schwangerschaftsabbruchs nach Beratung jedenfalls zu (II.). Daraus folgt, daß § 218a \abbrAbs{} 1 \abbrStGB{} \abbrNF{} verfassungsmäßig ist und daß Leistungen der gesetzlichen Krankenversicherung für Abbrüche nach Beratung aus § 24b \abbrSGBV{} nicht auszunehmen sind (III.).
\end{sourcepara}
\begin{transpara}
\transrn{381}對我們而言，如同對第二庭而言，國家自未出生人類生命開始時即負有憲法保護義務，並無疑問。我們與第二庭同樣認為，保護義務並不禁止立法者轉向一項保護構想，將重點置於懷孕婦女諮詢，並放棄指標事由決定之刑罰威嚇及指標事由構成要件之確認。但我們認為，第二庭作為立法可能性所承認之婦女在早期經諮詢後對終止妊娠之事實上最後責任，乃婦女基本權地位所要求，且保護義務就此受到限制（I.）。此外，依我們見解，諮詢規制中對未出生生命之有效保護，正要求規範上清楚說明終止妊娠何者被允許、何者不被允許；在此，《基本法》無論如何容許經諮詢後終止妊娠之正當化（II.）。由此可得，新版本《刑法》第218a條第1項合憲，且依 SGB V 第24b條不應排除對經諮詢後終止妊娠之法定醫療保險給付（III.）。
\end{transpara}

\bisubsection{I.}{I.}

\begin{sourcepara}[title={原文SondervotumMS.272}]
\sourcern{382}Nach Auffassung des Senats trifft die Frau während der gesamten Dauer einer Schwangerschaft eine Rechtspflicht zum Austragen des Kindes, die auch nach einer Beratung nur endet, wenn ein gesetzlicher, das Kriterium der Unzumutbarkeit ausformender Ausnahmetatbestand gegeben ist (vgl. Urteil, D.I.2.c). Dem stimmen wir nicht zu. Innerhalb des verfassungsrechtlich vorgegebenen Dreiecks zwischen der Frau, dem ungeborenen Leben und dem Staat nimmt die aus dem Grundgesetz abzuleitende Schutzpflicht für das ungeborene Leben allein den Staat in Anspruch, nicht unmittelbar schon die Frau. Pflichten, die der Staat im Wege der Gesetzgebung der Frau zum Schutz des ungeborenen Lebens auferlegt, müssen zugleich ihre Grundrechtspositionen berücksichtigen.
\end{sourcepara}
\begin{transpara}
\transrn{382}依第二庭見解，婦女在整個懷孕期間負有繼續懷孕之法律義務，即使經諮詢後，亦只有在存在一項法定、形塑不可期待性標準之例外構成要件時，該義務始告終止（參見判決 D.I.2.c）。我們不同意。在由憲法預設之婦女、未出生生命與國家三角關係中，由《基本法》可導出之對未出生生命之保護義務，僅要求國家，而非已直接要求婦女。國家透過立法為保護未出生生命而課予婦女之義務，必須同時考量其基本權地位。
\end{transpara}

\begin{sourcepara}[title={原文SondervotumMS.273}]
\sourcern{383}Jede gesetzliche Regelung des Schwangerschaftsabbruchs muß deshalb nicht nur mit der Schutzpflicht für das ungeborene Leben aus \abbrArt{} 2 \abbrAbs{} 2 \abbrIVm{} \abbrArt{} 1 \abbrAbs{} 1 \abbrGG{}\ggfnArtOneOneTwoTwoDE{}, sondern auch mit dem Anspruch der Frau auf Schutz und Achtung ihrer Menschenwürde (\abbrArt{} 1 \abbrAbs{} 1 \abbrGG{}\ggfnArtOneOneDE{}), ihrem Recht auf körperliche Unversehrtheit (\abbrArt{} 2 \abbrAbs{} 2 \abbrGG{}\ggfnArtTwoTwoDE{}) und ihrem Persönlichkeitsrecht (\abbrArt{} 2 \abbrAbs{} 1 \abbrGG{}\ggfnArtTwoOneDE{}) vereinbar sein (vgl. Urteil, D.I.2.b). Es ist Aufgabe des Gesetzgebers, seine Schutzpflicht gegenüber dem nasciturus und die Grundrechtsposition der Frau einander verhältnismäßig zuzuordnen. Die Verfassung gibt diese Zuordnung nicht vor. Deshalb hat der Gesetzgeber einen Abwägungs- und Gestaltungsspielraum, der auf der einen Seite durch ein Untermaßverbot gegenüber dem ungeborenen Leben, auf der anderen Seite aber durch das Übermaßverbot gegenüber der Frau, letztlich also durch den Verhältnismäßigkeitsgrundsatz beschränkt wird. Wenn der Gesetzgeber der Frau in der Frühphase der Schwangerschaft ermöglicht, nach Beratung, die ihr unter Strafandrohung zur Pflicht gemacht wird, eine eigenverantwortliche Entscheidung über die Fortsetzung der Schwangerschaft zu treffen, und insoweit von dem Verbot des Schwangerschaftsabbruchs und der Rechtspflicht zum Austragen der Schwangerschaft absieht, überschreitet er den ihm zustehenden Abwägungs- und Gestaltungsspielraum nicht. Soweit eine Rechtspflicht der Frau nicht besteht, ist ihr Handeln als Ausübung ihrer Grundrechte erlaubt.
\end{sourcepara}
\begin{transpara}
\transrn{383}因此，任何終止妊娠之法律規制，不僅必須與《基本法》第2條第2項結合第1條第1項\ggfnArtOneOneTwoTwoZH{}所生對未出生生命之保護義務相容，亦必須與婦女受保護及其人性尊嚴受尊重之請求（《基本法》第1條第1項\ggfnArtOneOneZH{}）、其身體完整權（《基本法》第2條第2項\ggfnArtTwoTwoZH{}）及其人格權（《基本法》第2條第1項\ggfnArtTwoOneZH{}）相容（參見判決 D.I.2.b）。立法者之任務，是將其對 \textit{nasciturus} 之保護義務與婦女基本權地位相互作比例性配置。憲法並未預設此一配置。因此，立法者享有衡量及形成空間；其一方面受相對於未出生生命之不足禁止限制，另一方面則受相對於婦女之過度禁止限制，最終亦即受比例原則限制。若立法者使婦女在懷孕早期，經其在刑罰威嚇下被課予義務之諮詢後，能就繼續懷孕作成自主負責決定，並就此放棄終止妊娠之禁止及繼續懷孕之法律義務，並未逾越其所享有之衡量及形成空間。只要婦女並無法律義務，其行為即作為其基本權之行使而被允許。
\end{transpara}

\begin{sourcepara}[title={原文SondervotumMS.274}]
\sourcern{384}1. Der Schwangerschaftskonflikt unterscheidet sich von allen anderen Gefährdungen menschlichen Lebens. Frau und nasciturus stehen sich nicht als mögliche "Täterin" und mögliches "Opfer" gegenüber, sondern bilden in der Person der Schwangeren eine einzigartige Einheit, eine -- wie es im Urteil heißt -- "Zweiheit in Einheit". In den ersten Wochen der Schwangerschaft -- so der Senat -- gehöre das neue Leben noch ganz der Mutter zu und sei von ihr in allem abhängig. Diese Unentdecktheit, Hilflosigkeit und Abhängigkeit des auf einzigartige Weise mit der Mutter verbundenen Ungeborenen ließen die Einschätzung berechtigt erscheinen, daß der Staat eine bessere Chance zu seinem Schutz habe, wenn er mit der Mutter zusammen wirke (vgl. Urteil, D.II.3.).
\end{sourcepara}
\begin{transpara}
\transrn{384}1. 懷孕衝突不同於所有其他對人類生命之危害。婦女與 \textit{nasciturus} 並非作為可能之「行為人」與可能之「被害人」相對立，而是在懷孕婦女人格中形成一種獨特單位，即判決所稱之「一體中之二元」。依第二庭所述，在懷孕最初數週，新生命尚完全屬於母親，並在一切方面依賴於她。以獨特方式與母親相連之未出生者，其未被發現、無助與依賴，使下列評估顯得正當：若國家與母親合作，其保護未出生者之機會較佳（參見判決 D.II.3.）。
\end{transpara}

\begin{sourcepara}[title={原文SondervotumMS.275}]
\sourcern{385}Dies ist der Ausgangspunkt im Tatsächlichen für die Beschreibung der Grundrechtsposition der Frau in der Frühphase der Schwangerschaft. Weitere Annäherungen erbringen die Aussagen im Urteil, wonach die Beratungsregelung auf der Letztverantwortung der schwangeren Frau für die Entscheidung über Fortgang und Abbruch der Schwangerschaft gründet und gründen darf. Verfassungsrechtlich unbedenklich sei die Einschätzung des Gesetzgebers, daß ein wirksamer Schutz des ungeborenen menschlichen Lebens nur mit der Mutter, nicht aber gegen sie möglich ist (vgl. Urteil, D. II. 2., 3. und 4.; so schon der Alternativ-Entwurf zum Besonderen Teil des Strafgesetzbuches von 1970, vgl. \abbrBVerfGE{} 39, 1 [10 \abbrF{}] und das Sondervotum Rupp-von Brünneck/Simon, \abbrAAO{}, \abbrS{} 79).
\end{sourcepara}
\begin{transpara}
\transrn{385}這是描述婦女在懷孕早期之基本權地位的事實出發點。判決中另有若干表述可進一步接近此點：諮詢規制建立於且得建立於懷孕婦女對繼續與終止妊娠決定之最後責任。立法者認為未出生人類生命之有效保護只能與母親一同，而不能對抗母親而實現，此評估在憲法上無疑義（參見判決 D.II.2., 3. 及 4.；1970年《刑法》分則替代草案已如此，參見 BVerfGE 39, 1 [10 f.] 及 Rupp-von Brünneck/Simon 不同意見，a.a.O., S. 79）。
\end{transpara}

\begin{sourcepara}[title={原文SondervotumMS.276}]
\sourcern{386}Hiervon ausgehend ist die Beratungsregelung, wie auch das Urteil deutlich erkennen läßt, nicht das frustrierte Ausweichen vor einem frustrierenden Mißerfolg der Indikationenlösung. Die neue Regelung ist vielmehr die Konsequenz aus einem gewandelten Verständnis von der Personalität und Würde der Frau. Die Aussage des Urteils, die Frau sei zu eigenverantwortlicher Entscheidung über Fortsetzung oder Abbruch der Schwangerschaft fähig, muß aber zu Folgerungen bei der Verfassungsauslegung führen. Sie zwingt nach unserer Auffassung dazu, die Kollision von Menschenwürde des ungeborenen Lebens einerseits und Menschenwürde der schwangeren Frau andererseits, anders als im Urteil geschehen, durch verhältnismäßige Zuordnung aufzulösen: Das einzigartige Zuordnungsproblem der "Zweiheit in Einheit" kann grundrechtlich auch nicht annähernd in einer bloßen Gegenüberstellung von Embryo und Frau eingefangen werden. Ihre eigene grundrechtliche Lage ist vielmehr ihrem Wesen nach durch die Verantwortung für das andere Leben mitbestimmt, weil sie dieses Leben in sich trägt. Damit wird nicht ausgeklammert, daß der Frau dieses andere Leben mit eigener menschlicher Würde auch "gegenübersteht". Doch kommt erst in diesen beiden Aussagen das Spezifische des Ausgleichs zwischen der Grundrechtsposition der schwangeren Frau und der Schutzpflicht zum Ausdruck.
\end{sourcepara}
\begin{transpara}
\transrn{386}由此出發，如判決亦清楚顯示，諮詢規制並非面對令人挫折之指標方案失敗所作之受挫迴避。新規制毋寧是對婦女人格性與尊嚴之理解發生變遷後的結果。判決稱婦女有能力就繼續或終止妊娠作成自主負責決定，則此一陳述必須導致憲法解釋上之結論。依我們見解，其迫使人將未出生生命之人性尊嚴一方與懷孕婦女之人性尊嚴一方之衝突，不同於判決之作法，透過比例性配置加以解決：「一體中之二元」此一獨特配置問題，在基本權上亦不能近似地以胚胎與婦女之單純對立加以掌握。婦女自身之基本權狀態，依其本質毋寧亦由其對另一生命之責任共同決定，因為她在自身內承載此生命。這並不排除該另一具有自身人性尊嚴之生命也「面對」婦女。然而，唯有在此兩項陳述中，懷孕婦女基本權地位與保護義務間調和之特殊性才表達出來。
\end{transpara}

\begin{sourcepara}[title={原文SondervotumMS.277}]
\sourcern{387}Die "Zweiheit in Einheit" unterliegt mit dem Fortschreiten der Schwangerschaft Veränderungen. Während in den ersten Wochen Frau und nasciturus -- wie soeben aus dem Urteil referiert -- noch ganz als Einheit erscheinen, tritt mit dem Wachsen des Embryos die "Zweiheit" stärker hervor. Dieser Entwicklungsprozeß hat auch rechtliche Bedeutung. Zwar bleibt die Eigenverantwortung der Frau, diese verliert aber den Charakter als Letztverantwortung. Die gesetzliche Positivierung des Ausgleichs zwischen Grundrechtsposition der schwangeren Frau und Schutzpflicht für das ungeborene Leben muß ein prozeßhaftes Element im Grundrechtsstatus der schwangeren Frau berücksichtigen, weil die Schwangerschaft selbst einen Prozeß darstellt. Mit dem Fortschreiten der Schwangerschaft und dem Heranwachsen des nasciturus verschieben sich die Gewichte in der Zuordnung; der Inhalt der Grundrechtsposition der Frau und die Rolle des Staates in der Wahrnehmung seiner Schutzpflicht sind in der Frühphase anders zu beurteilen als im fortgeschrittenen Stadium.
\end{sourcepara}
\begin{transpara}
\transrn{387}「一體中之二元」隨懷孕進展而變化。懷孕最初數週中，婦女與 \textit{nasciturus} 如剛才引述判決所示，尚完全呈現為一體；隨胚胎成長，「二元」則更強烈浮現。此一發展過程亦具有法律意義。婦女之自主責任固然仍存在，但其失去最後責任之性格。懷孕婦女基本權地位與對未出生生命之保護義務間調和之法律實證化，必須考量懷孕婦女基本權地位中之過程性要素，因為懷孕本身即是一個過程。隨懷孕進展及 \textit{nasciturus} 成長，配置中之重量移動；婦女基本權地位之內容及國家履行保護義務時之角色，在早期應作不同於後期之評價。
\end{transpara}

\begin{sourcepara}[title={原文SondervotumMS.278}]
\sourcern{388}2. Dieser Prozeß erfordert eine Antwort des Gesetzgebers in der unterschiedlichen Ausgestaltung der staatlichen Schutzpflicht in der Frühphase der Schwangerschaft einerseits und in deren weiterem Verlauf andererseits. Der Staat nimmt die Schutzpflicht in der Frühphase dadurch wahr, daß er die Frau im Konfliktfall unter der Strafandrohung des § 218 \abbrStGB{} zu einer Beratung verpflichtet, in deren Mittelpunkt das ungeborene Leben steht; mit deren Abschluß gibt er aber der Eigenverantwortlichkeit und der Fähigkeit der Frau zu einer letztverantwortlichen Entscheidung Raum. Die Frau ist hier Ansprechpartnerin, nicht Anspruchsgegnerin. Der Staat nimmt in der Ausübung dieser so bestimmten Schutzpflicht die "Zweiheit in Einheit" ernst in der Weise, daß die Frau nicht mehr nur als Gefäß des Embryos erscheint. Im späteren Zeitraum der Schwangerschaft verteidigt die Schutzpflicht nunmehr durch die Androhung von Strafe das Recht des Embryos auf Leben. Sie geht grundsätzlich -- wie bei anderen Kollisionen -- von einander gegenüberstehenden Rechtsgütern aus, wobei es Sache des Gesetzgebers ist, im einzelnen zu bestimmen, von welcher Woche an dies gelten soll.
\end{sourcepara}
\begin{transpara}
\transrn{388}2. 此一過程要求立法者對懷孕早期與其後續進程中之國家保護義務作不同形塑以為回應。國家在早期履行保護義務之方式，是於衝突案件中在《刑法》第218條刑罰威嚇下，使婦女負有諮詢義務；該諮詢以未出生生命為核心。但諮詢完成後，國家為婦女之自主責任及最後負責決定之能力留下空間。婦女在此是對話對象，而非請求相對人。國家在如此界定之保護義務行使中，認真看待「一體中之二元」，使婦女不再僅顯得是胚胎之容器。在懷孕較後期間，保護義務則透過刑罰威嚇保衛胚胎之生命權。其原則上如其他衝突般，以相互對立之法益為出發點；至於自第幾週起應如此，則由立法者具體決定。
\end{transpara}

\begin{sourcepara}[title={原文SondervotumMS.279}]
\sourcern{389}Ob die Ausübung der staatlichen Schutzpflicht im Modus der Beratung stärker wirkt oder nicht, kann nur empirisch beurteilt werden. Maßgebend ist, daß der Staat nicht auf inadäquate Weise schützt; dies ist der Vorwurf gegenüber der Indikationenregelung, der zu der Konzeption des Schwangeren- und Familienhilfegesetzes geführt hat. Der Schutz muß nach der vertretbaren Einschätzung des Gesetzgebers in einem realen Sinne wirksam sein (vgl. Urteil, D.I.2.b). Mit dem Senat sind wir der Meinung, daß die Einschätzung des Gesetzgebers, in der Frühphase besser durch die Beratung, in der darauffolgenden besser durch das Strafrecht schützen zu können, vertretbar ist.
\end{sourcepara}
\begin{transpara}
\transrn{389}國家保護義務以諮詢模式行使是否更具效果，只能以經驗方式判斷。關鍵在於，國家不得以不適當方式保護；這正是對指標規則之指摘，並導致《孕婦及家庭扶助法》之構想。依立法者可支持之評估，保護必須在實際意義上有效（參見判決 D.I.2.b）。我們與第二庭同樣認為，立法者認為在早期較能透過諮詢保護、在其後較能透過刑法保護之評估，是可支持的。
\end{transpara}

\begin{sourcepara}[title={原文SondervotumMS.280}]
\sourcern{390}Allerdings ist mit der Letztverantwortung der Frau nach Beratung ihr Vorrang in der Entscheidung über Abbruch oder Fortsetzung der Schwangerschaft verbunden. Das ist, folgt man dem Urteil, für den Senat im tatsächlichen Sinne nicht anders, während wir dies für grundrechtlich fundiert halten. Das Bild von der "Zweiheit in Einheit" ist für uns nicht nur die Beschreibung eines tatsächlichen Zustandes sondern gibt in Wahrheit den grundrechtlichen Status der Frau wieder. Dabei geht es hier aber nicht einfach um "das Persönlichkeitsrecht der Frau" (vgl. Urteil, D.I.2.c.bb), auch nicht um eine Variante des "Selbstbestimmungsrechts"; damit wäre die Frau erneut nur das "Gegenüber" zum Embryo und dieser nicht auch Teil ihrer selbst. Auch kommt, geht man von einem Vorrang der Rechtsstellung der Frau in dieser Frühphase aus, nicht -- wie der Senat (vgl. Urteil, \abbrAAO{}) meint -- das Lebensrecht des Ungeborenen nur dann zur Geltung, wenn die Mutter die Schwangerschaft nicht abbricht. Hier muß sich der Senat fragen lassen, wie er dem Lebensrecht des Ungeborenen im Urteil auf wirksamere Weise als eben durch die Beratung Geltung verschafft.
\end{sourcepara}
\begin{transpara}
\transrn{390}然而，婦女經諮詢後之最後責任，伴隨其在終止或繼續懷孕決定上之優先地位。若依判決而論，第二庭在事實意義上並非作不同理解；但我們認為此點具有基本權基礎。對我們而言，「一體中之二元」之圖像不只是對事實狀態之描述，而是真正反映婦女之基本權地位。在此所涉並非單純「婦女之人格權」（參見判決 D.I.2.c.bb），亦非「自決權」之一種變體；若如此，婦女將再次只是胚胎之「相對者」，而胚胎並非亦為其自身之一部分。若以婦女在此早期之法律地位優先為出發點，亦不會如第二庭（參見判決，a.a.O.）所認為，未出生者之生命權只有在母親不終止妊娠時才得以實現。第二庭在此必須被追問：其在判決中究竟如何以比諮詢更有效之方式，使未出生者之生命權獲得實現。
\end{transpara}

\begin{sourcepara}[title={原文SondervotumMS.281}]
\sourcern{391}Die Rechtsordnung hat dem beschriebenen Prozeß seit dem Fünften Gesetz zur Reform des Strafrechts (5. \abbrStrRG{}) vom 18. Juni 1974 (\abbrBGBl{} I \abbrS{} 1297), das erstmals die Zwölf-Wochen-Frist in das Strafrecht des Schwangerschaftsabbruchs einführte, Rechnung getragen; die dadurch charakterisierte Frühphase ist durch das Urteil des Ersten Senats über die unterschiedlichen Ausprägungen des Strafrechts hinweg erhalten geblieben -- stets mit wesentlichen Konsequenzen für die Begrenzung der staatlichen Schutzpflicht. Auch alle Gesetzentwürfe, die der Beschlußfassung über das Schwangeren- und Familienhilfegesetz vorangingen und von Fraktionen oder Teilen der Fraktionen des Deutschen Bundestages eingebracht wurden, sind von dieser Frühphase ausgegangen.
\end{sourcepara}
\begin{transpara}
\transrn{391}自1974年6月18日第五次刑法改革法（5. StrRG）（BGBl. I S. 1297）以來，法律秩序已顧及上述過程；該法首次將十二週期限引入終止妊娠刑法。由此標示之早期，經第一庭判決並跨越刑法不同形態而被保留下來，且始終對國家保護義務之限制具有重要後果。所有先於《孕婦及家庭扶助法》決議、由德國聯邦議院各黨團或黨團部分成員提出之法律草案，亦均以此早期為出發點。
\end{transpara}

\begin{sourcepara}[title={原文SondervotumMS.282}]
\sourcern{392}Der Gesetzgeber des Schwangeren- und Familienhilfegesetzes hat den Grundrechtsstatus der schwangeren Frau darin berücksichtigt, daß er der Frau während eines begrenzten Zeitraums zu Beginn der Schwangerschaft die Möglichkeit zur Entscheidung ihres Konflikts einräumt. Dabei nimmt er seine Schutzpflicht wahr, indem er der Frau vor einem Abbruch der Schwangerschaft die Beratung zur Pflicht macht. Erst nach Ablauf eines solchen Zeitraumes begründet er die volle Einstandspflicht der Frau für das ungeborene Leben. Damit schafft er im Grundsatz den je verhältnismäßigen Ausgleich zwischen den im Schwangerschaftskonflikt betroffenen Grundrechtspositionen.
\end{sourcepara}
\begin{transpara}
\transrn{392}《孕婦及家庭扶助法》之立法者，透過在懷孕開始之一段有限期間內賦予婦女決定其衝突之可能性，顧及懷孕婦女之基本權地位。同時，立法者透過使婦女在終止妊娠前負有諮詢義務，履行其保護義務。只有在此一期間屆滿後，立法者才建立婦女對未出生生命之完整擔保義務。藉此，立法者原則上在懷孕衝突所涉及之基本權地位之間，創設各別合比例之調和。
\end{transpara}

\begin{sourcepara}[title={原文SondervotumMS.283}]
\sourcern{393}3. Der Begriff "Zweiheit in Einheit" läßt sich demnach als eine begriffliche Annäherung an die richtige Erfassung einer einzigartigen grundrechtlichen Lage verstehen. Die in ihm mitgedachte natürliche Prozeßhaftigkeit der Schwangerschaft ist dabei im grundrechtsdogmatischen Sinne aufzufassen, nämlich als Entwicklung von einer Letztverantwortung der Frau für das ungeborene Leben, die in der Achtung ihrer Personalität wurzelt (\abbrArt{} 1 \abbrAbs{} 1 \abbrIVm{} \abbrArt{} 2 \abbrAbs{} 1 \abbrGG{}\ggfnArtOneOneTwoOneDE{}), zu einer Übernahme der Letztverantwortung für das ungeborene Kind durch den Staat. Das niederländische Strafgesetzbuch hat folgerichtig den Tatbestand des Schwangerschaftsabbruchs mit (etwa) der 24. Schwangerschaftswoche enden lassen und bestraft von diesem Zeitpunkt an den Abbruch als Totschlag, wenn "redlicherweise erwartet werden darf, daß sie (die Leibesfrucht) imstande ist, außerhalb des Mutterleibes am Leben zu bleiben" (\abbrArt{} 82a \abbrStGB{}, zit. nach Eser/Koch, Schwangerschaftsabbruch im internationalen Vergleich, Teil 1 Europa, 1988, \abbrS{} 1073). Dieses Strafrecht läßt also das ungeborene Leben in der letzten Schwangerschaftsphase zu einer Person im vollen strafrechtlichen Sinne werden.
\end{sourcepara}
\begin{transpara}
\transrn{393}3. 因此，「一體中之二元」此一概念，可理解為對正確掌握一種獨特基本權狀態之概念性接近。其一併思考之懷孕自然過程性，應在基本權教義學意義上理解，亦即理解為一個發展：從根植於對婦女人格性之尊重（《基本法》第1條第1項結合第2條第1項\ggfnArtOneOneTwoOneZH{}）之婦女對未出生生命之最後責任，發展至由國家承擔對未出生子女之最後責任。荷蘭刑法典因而合乎邏輯地使終止妊娠犯罪構成要件約於懷孕第24週終止，並自該時點起，在「合理可期待其（胎兒）有能力在母體外維持生命」時，將終止妊娠作為殺人罪處罰（刑法第82a條，引自 Eser/Koch, Schwangerschaftsabbruch im internationalen Vergleich, Teil 1 Europa, 1988, S. 1073）。此一刑法遂使未出生生命在懷孕最後階段成為完整刑法意義上之人。
\end{transpara}

\begin{sourcepara}[title={原文SondervotumMS.284}]
\sourcern{394}4. Demgegenüber wird der Maßstab der Unzumutbarkeit (Urteil, D.I.2.c.bb) der Einzigartigkeit der Situation nicht gerecht. Der Senat hält an ihm auch im Rahmen der Beratungsregelung fest; auch ohne den Rechtfertigungsgrund einer allgemeinen Notlagenindikation müsse die Rechtsordnung zum Ausdruck bringen, daß ein Schwangerschaftsabbruch nur in Ausnahmesituationen einer Belastung zulässig sei, die für die Frau die zumutbare Opfergrenze übersteige, um so im konkreten Schwangerschaftskonflikt und für das allgemeine Rechtsbewußtsein die notwendige Orientierung über die Rechtspflicht zum Austragen des Kindes und deren Grenzen zu geben (vgl. Urteil, D.III.1.c).
\end{sourcepara}
\begin{transpara}
\transrn{394}4. 相對於此，不可期待性標準（判決 D.I.2.c.bb）並不能適切對應此一情境之獨特性。第二庭即使在諮詢規制架構內仍堅持此標準；即便沒有一般困境指標事由之正當化事由，法律秩序仍必須表明，終止妊娠僅在負擔之例外情境中始得允許，而該負擔對婦女而言超過可期待之犧牲界限，以便在具體懷孕衝突中並對一般法律意識，就繼續懷孕之法律義務及其界限提供必要指引（參見判決 D.III.1.c）。
\end{transpara}

\begin{sourcepara}[title={原文SondervotumMS.285}]
\sourcern{395}Als eine solche orientierende Leitlinie der Rechtsordnung halten wir das Kriterium der Unzumutbarkeit für ungeeignet. Ein eigenes Subsumieren der Konfliktlage unter das Kriterium einer an den beiden gesetzlichen Indikationen orientierten zumutbaren Opfergrenze (vgl. Urteil, D.I.2.c.bb) muß die Frau überfordern. Denn die mit einer ungewollten Schwangerschaft verbundene Konfliktlage stellt sich für jede Frau je nach ihrer physischen und psychischen Situation unterschiedlich dar. In dieser Situation wird sie die "zumutbare Opfergrenze" nur dort zu sehen vermögen, wo sie über das Gebären des Kindes hinaus Chancen und Perspektiven für eine verantwortungsvolle eigene Mutterschaft (vgl. \abbrArt{} 6 \abbrAbs{} 2 \abbrGG{}\ggfnArtSixTwoDE{}) sieht.
\end{sourcepara}
\begin{transpara}
\transrn{395}我們認為，不可期待性標準不適合作為法律秩序之此種指引準則。要求婦女自行將其衝突狀態涵攝於一項以兩項法定指標事由為取向之可期待犧牲界限標準下（參見判決 D.I.2.c.bb），必然使其負擔過重。因為非意願懷孕所伴隨之衝突狀態，依各婦女之身體與心理處境而各不相同。在此情境中，她只有在看見除生下子女之外，尚有負責任地履行自身母職（參見《基本法》第6條第2項\ggfnArtSixTwoZH{}）之機會與前景時，才可能看見「可期待之犧牲界限」所在。
\end{transpara}

\begin{sourcepara}[title={原文SondervotumMS.286}]
\sourcern{396}Rechtlich gesehen führt -- und verführt -- das Kriterium der Unzumutbarkeit zu Subsumierungen, deren Maßstäbe um so vager sind, je mehr die jeweilige Konfliktlage der Subjektivität der Frau verhaftet ist. Diese Subsumierungen sind normativ nur begrenzt zu steuern; unausweichlich sind sie stark soziologisch bestimmt (Konfessionsgefälle, Stadt-Land-Gefälle usw.). Dies hat die allgemeine Notlagenindikation des bisherigen Rechts auch unter dem Gesichtspunkt des Bestimmtheitsgrundsatzes des \abbrArt{} 103 \abbrAbs{} 2 \abbrGG{}\ggfnArtOneZeroThreeTwoDE{} zweifelhaft gemacht (vgl. dazu auch das Sondervotum Rupp-von Brünneck/Simon, \abbrBVerfGE{} 39, 68 [91]). Daß bei einer "Eigensubsumierung" durch die Frau in und nach der Beratung die Schwierigkeiten noch größer werden, bedarf keiner besonderen Begründung.
\end{sourcepara}
\begin{transpara}
\transrn{396}從法律上看，不可期待性標準導致，且誘使人作出涵攝；而各該衝突狀態越是繫於婦女之主觀性，此等涵攝之標準即越加模糊。此等涵攝在規範上僅能有限控制；其不可避免地強烈受社會學因素決定（宗派差異、城鄉差異等）。這也使既有法律之一般困境指標事由，在《基本法》第103條第2項\ggfnArtOneZeroThreeTwoZH{}明確性原則觀點下顯得可疑（另參見 Rupp-von Brünneck/Simon 不同意見，BVerfGE 39, 68 [91]）。若由婦女在諮詢中及諮詢後進行「自行涵攝」，困難將更大，無須特別說明。
\end{transpara}

\begin{sourcepara}[title={原文SondervotumMS.287}]
\sourcern{397}Abwägungen nach dem Maßstab der Zumutbarkeit gehen von sich gegenüberstehenden Rechtsgütern aus, deren eines zerstört wird. Bei diesem Ausgangspunkt dürfte die Schutzpflicht des Staates aber folgerichtig nur durch den rechtfertigenden Notstand des § 34 \abbrStGB{} begrenzt werden. Geschieht dies nicht, läuft die Handhabung des Zumutbarkeitskriteriums auf eben den "Vorrang" für die Frau hinaus, den der Senat ablehnt (vgl. Urteil, D.II.3.). Von den Initiatoren der im Bundestag eingebrachten Gesetzentwürfe war offenbar nur denen des sog. Werner-Entwurfs (\abbrBTDrucks{} 12/1179), der sich auf die medizinische Indikation beschränkte, das Problem bewußt. Die rechtliche Umsetzung des Zumutbarkeitsgedankens in und seit der Entscheidung des Ersten Senats in \abbrBVerfGE{} 39, 1 \abbrFF{} -- insbesondere in der Formulierung der allgemeinen Notlagenindikation in der Entscheidungsformel dieses Urteils -- entsprach deshalb so lange nicht dem Rang des Rechtsguts des ungeborenen menschlichen Lebens, als nicht dessen verhältnismäßige Zuordnung zur Grundrechtsposition der Frau mit in den Blick genommen wurde.
\end{sourcepara}
\begin{transpara}
\transrn{397}依可期待性標準所作之衡量，是以相互對立且其中之一將被毀滅之法益為出發點。但在此出發點下，國家保護義務按理只能受《刑法》第34條正當化緊急避難之限制。若非如此，對可期待性標準之操作，即導向第二庭所拒絕之婦女「優先」（參見判決 D.II.3.）。在提交聯邦議院之法律草案發起者中，顯然只有所謂 Werner 草案（BTDrucks. 12/1179）之提出者意識到此問題；該草案限於醫學指標事由。因此，自第一庭在 BVerfGE 39, 1 ff. 作成判決以來，對可期待性思想之法律落實，尤其是在該判決主文中對一般困境指標事由之表述，只要未同時將其與婦女基本權地位作比例性配置納入視野，即不符合未出生人類生命此一法益之位階。
\end{transpara}

\begin{sourcepara}[title={原文SondervotumMS.288}]
\sourcern{398}Im Rückblick erweisen sich deshalb aus unserer Sicht schon die Begrenzungen der staatlichen Schutzpflicht in den Indikationstatbeständen, insbesondere in dem der allgemeinen Notlage, als Ausdruck der verhältnismäßigen Zuordnung der in einzigartiger Weise miteinander verknüpften Grundrechtspositionen und der Anerkennung einer den Schutz des Embryos grundsätzlich umschließenden Letztverantwortung der Frau in der Frühphase -- wenn auch bislang beschränkt auf die Umstände, die sich nach Auffassung des Bundesverfassungsgerichts und des Gesetzgebers objektivieren ließen.
\end{sourcepara}
\begin{transpara}
\transrn{398}因此，回顧而言，依我們觀點，即使國家保護義務在指標事由構成要件中，尤其在一般困境構成要件中所受之限制，也已表現為對以獨特方式相互連結之基本權地位所作比例性配置，並承認婦女在早期具有原則上包容胚胎保護之最後責任；縱使迄今僅限於依聯邦憲法法院及立法者見解得以客觀化之情況。
\end{transpara}

\begin{sourcepara}[title={原文SondervotumMS.289}]
\sourcern{399}Der Senat verwendet den Zumutbarkeitsgedanken nicht so wie der Erste Senat, also im Sinne einer praktizierten und auf tatsächlich wirksam werdenden Lebensschutz gerichteten Scheidung der gerechtfertigten von den verwerflichen Abbrüchen (vgl. \abbrBVerfGE{} 39, 1 [58]). Um der Offenheit der Beratung willen steht der Unzumutbarkeit nicht mehr der Gegenbegriff zumutbarer Fortsetzung der Schwangerschaft und deshalb strafbaren Abbruchs gegenüber. Dadurch verändert das Zumutbarkeitskriterium seine rechtliche Bedeutung. Indem es in der allgemeinen Notlagenindikation nur noch der -- rechtlich folgenlosen -- Orientierung der schwangeren Frau für ihre Entscheidung dient, der Abbruch aber in jedem Fall rechtswidrig ist, wurzelt das Kriterium letztlich nur noch im sittlichen und kaum noch im rechtlichen Bereich. Maßstäbliche Schutzpflicht und praktische Schutzpflicht lassen sich nicht mehr miteinander verbinden. Insofern steht die Beratungsregelung zum Indikationsmodell des Senats quer.
\end{sourcepara}
\begin{transpara}
\transrn{399}第二庭使用可期待性思想之方式，並非如第一庭那樣，亦即並非意在實務上區分正當化之終止妊娠與應受非難之終止妊娠，並使其實際指向有效生命保護（參見 BVerfGE 39, 1 [58]）。為了諮詢之開放性，不可期待性不再對應於可期待地繼續懷孕並因此終止妊娠可罰之反概念。可期待性標準因而改變其法律意義。當其在一般困境指標事由中僅作為懷孕婦女決定之指引，且該指引在法律上無後果，而終止妊娠在任何情況下均屬違法時，此標準最終只還根植於倫理領域，而幾乎不再根植於法律領域。作為標準之保護義務與實際保護義務不再能相互連結。就此而言，諮詢規制與第二庭之指標模型相互交錯而不相合。
\end{transpara}

\begin{sourcepara}[title={原文SondervotumMS.290}]
\sourcern{400}5. Diese Erwägungen können dem Schwangerschaftskonflikt in tatsächlicher Hinsicht nichts von seiner Schärfe nehmen. Sie zeigen aber auf, warum das Urteil des Bundesverfassungsgerichts vom 25. Februar 1975 (\abbrBVerfGE{} 39, 1 \abbrFF{}) erstmals zu einer differenzierenden Betrachtung gelangte, und vor allem, warum überhaupt der Schritt zur Beratungslösung verantwortet werden konnte. Das vorliegende Urteil selbst veranschaulicht diesen Wandel. Es sieht die Positionen der Frau und des nasciturus je in ihrer Menschenwürde verwurzelt (vgl. Urteil, D.I.2.b), während noch der Erste Senat dies zwar für das ungeborene Leben für zutreffend hielt (vgl. \abbrBVerfGE{} 39, 1 [41]), die Grundrechtsposition der Frau hingegen nur als ihr Recht auf freie Entfaltung ihrer Persönlichkeit (\abbrArt{} 2 \abbrAbs{} 1 \abbrGG{}\ggfnArtTwoOneDE{}) beschrieb, so daß bei der dann folgenden Orientierung am \abbrArt{} 1 \abbrAbs{} 1 \abbrGG{}\ggfnArtOneOneDE{} die Entscheidung für den Vorrang des Schutzes für das ungeborene Leben vor dem Selbstbestimmungsrecht der Schwangeren vorgezeichnet war (vgl. \abbrAAO{}, \abbrS{} 43).
\end{sourcepara}
\begin{transpara}
\transrn{400}5. 此等考量無法在事實上減輕懷孕衝突之尖銳性。但其指出，為何聯邦憲法法院1975年2月25日判決（BVerfGE 39, 1 ff.）首次達成一種區分性觀察，並尤其指出，為何根本能夠負責任地邁向諮詢方案。本判決本身即說明此一變遷。其認為婦女與 \textit{nasciturus} 之地位各自根植於其人性尊嚴（參見判決 D.I.2.b）；而第一庭當時雖認為此點對未出生生命而言正確（參見 BVerfGE 39, 1 [41]），卻僅將婦女之基本權地位描述為其人格自由發展權（《基本法》第2條第1項\ggfnArtTwoOneZH{}），因此在其後以《基本法》第1條第1項\ggfnArtOneOneZH{}為取向時，已預先決定使對未出生生命之保護優先於懷孕婦女自決權（參見 a.a.O., S. 43）。
\end{transpara}

\bisubsection{II.}{II.}

\begin{sourcepara}[title={原文SondervotumMS.291}]
\sourcern{401}Unbeschadet der Ausführungen unter I. gebietet nach unserer Auffassung das Grundgesetz nicht, Schwangerschaftsabbrüchen, die in der Frühphase nach Beratung durch einen Arzt erfolgen und von der Strafdrohung ausgenommen sind, die (strafrechtliche) Rechtfertigung dann zu versagen, wenn nicht ein Dritter festgestellt hat, daß die Fortsetzung der Schwangerschaft unzumutbar ist.
\end{sourcepara}
\begin{transpara}
\transrn{401}不影響上述 I. 之說明，依我們見解，《基本法》並不要求，對於在早期經諮詢後由醫師實施且排除刑罰威嚇之終止妊娠，若未由第三人確認繼續懷孕不可期待，即否定其（刑法上）正當化。
\end{transpara}

\begin{sourcepara}[title={原文SondervotumMS.292}]
\sourcern{402}Die Rechtfertigung der nach Beratung erfolgenden Abbrüche ist der unentbehrliche Schlußstein der Beratungsregelung. Die Annahme einer den Abbruch rechtfertigenden Ausnahmelage ist auch ohne Feststellung ihrer Voraussetzungen durch Dritte mit dem Grundgesetz vereinbar (1.). Ein Rechtswidrigkeitsurteil außerhalb des Strafrechts vermag nach unserer Auffassung zur Erfüllung der staatlichen Schutzpflicht für das ungeborene Leben nichts beizutragen (2.).
\end{sourcepara}
\begin{transpara}
\transrn{402}經諮詢後終止妊娠之正當化，是諮詢規制不可或缺之拱頂石。即使未由第三人確認其前提，承認存在正當化終止妊娠之例外情境，亦與《基本法》相容（1.）。依我們見解，刑法之外之違法判斷，無法對履行國家對未出生生命之保護義務有所貢獻（2.）。
\end{transpara}

\begin{sourcepara}[title={原文SondervotumMS.293}]
\sourcern{403}1. Die Beratungsregelung setzt auf die letztverantwortliche Entscheidung der Frau nach dem Beratungsgespräch. Das Verlangen nach Abbruch der Schwangerschaft, an dem die Frau selbst nach vorangegangener Beratung und ärztlichem Gespräch festhält, stellt sich grundsätzlich als verantwortungsvolle Entscheidung dar. Diese Entscheidung muß von der Rechtsordnung anerkannt sein, wenn die Beratung die ihr zugedachte lebensschützende Wirkung entfalten soll. Die Beratung kann nicht gelingen, wenn die Entscheidung der Frau gegen eine Fortsetzung der Schwangerschaft zwar von Strafe ausgenommen, aber -- außerhalb des Strafrechts -- als nicht gerechtfertigt behandelt wird und hieran rechtliche Nachteile geknüpft werden. Der Gesetzgeber darf bei der für ihn maßgeblichen, vom Regelfall ausgehenden Betrachtungsweise an die Entscheidung der Frau ohne Verfassungsverstoß die Rechtfertigung knüpfen. Auf eine Tatsachenfeststellung durch Dritte, die mit den Wirkungsbedingungen der Beratungsregelung unvereinbar wäre, darf er hierfür verzichten. Ist eine so gestaltete und folgerichtig zu Ende geführte Beratungsregelung geeignet, generell effektiveren Schutz für das ungeborene Leben zu gewährleisten als das bisher geltende Recht, so bedeutet dies zugleich auch wirksameren Schutz für jedes einzelne Ungeborene.
\end{sourcepara}
\begin{transpara}
\transrn{403}1. 諮詢規制仰賴婦女於諮詢談話後之最後負責決定。即使在先前諮詢及醫師談話後，婦女仍堅持之終止妊娠要求，原則上呈現為一項負責任之決定。若諮詢欲發揮其被賦予之生命保護效果，法律秩序即必須承認此一決定。若婦女反對繼續懷孕之決定雖免於刑罰，卻在刑法之外被視為未受正當化，並與法律不利益相連結，諮詢即不可能成功。立法者得在其所採、以通常情形為出發點之觀察方式下，將正當化連結於婦女之決定，而不違反憲法。對於與諮詢規制作用條件不相容之第三人事實確認，立法者得予放棄。若如此設計並合乎邏輯貫徹到底之諮詢規制，適於一般性地保障較既有法律更有效之未出生生命保護，則此同時亦意味著對每一個別未出生者之更有效保護。
\end{transpara}

\begin{sourcepara}[title={原文SondervotumMS.294}]
\sourcern{404}a) Der Staat kann in der Frühphase der Schwangerschaft Schutz für das ungeborene Leben nur dadurch bewirken, daß er die Frau als seine Verbündete bei der Erfüllung seiner Schutzaufgabe gewinnt. Das kann nur gelingen, wenn er sie in ihrer Fähigkeit zu verantwortungsvoller Entscheidung sowie in ihrer besonderen Befindlichkeit am Beginn einer Schwangerschaft ernst nimmt. Die in der mündlichen Verhandlung zur Praxis der Beratung gehörten Beraterinnen haben übereinstimmend bekundet, daß Frauen -- ausgehend von ihrer natürlichen Schutzbereitschaft für das in ihnen heranwachsende ungeborene Leben -- ihre Schwangerschaftskonflikte als Not empfinden und in dieser Situation verantwortlich und gewissenhaft handeln wollen. Sie erleben ihren Konflikt als höchstpersönlichen und wehren sich deshalb gegen dessen Beurteilung nach einem von dritter Seite an sie herangetragenen Maßstab der Zumutbarkeit. Mithin muß die Rechtsordnung, wenn sie dem ungeborenen Leben Schutz geben will, der Frau Raum für eine verantwortliche Entscheidung lassen, ihr also Verantwortung nicht nur auferlegen sondern auch anvertrauen. Aus diesem Grunde sehen wir allerdings -- anders als der Senat (vgl. Urteil, D.IV.1.b) -- auch in der Möglichkeit, daß eine Frau in der Beratung auch schweigen darf, ein wichtiges Element der völligen Offenheit des Beratungsgesprächs. Und wenn sie verantwortlich entscheiden soll, muß ihre Entscheidung -- ohne rechtliche Vorbehalte -- anerkannt werden. Nur dann kann sie in der Beratung wirklich offen sein.
\end{sourcepara}
\begin{transpara}
\transrn{404}a) 國家在懷孕早期，只有透過使婦女成為其履行保護任務之盟友，才能實現對未出生生命之保護。只有在國家認真看待婦女負責任決定之能力，以及其於懷孕開始時之特殊處境時，這才可能成功。在言詞辯論中就諮詢實務所聽取之諮詢員一致陳述，婦女基於其對自身內正在成長之未出生生命之自然保護意願，將其懷孕衝突感受為困境，並欲在此情境中負責任且憑良心行動。她們將其衝突經驗為高度個人性，因此抗拒依第三方帶入之可期待性標準對其衝突加以評價。因此，法律秩序若欲給予未出生生命保護，即必須為婦女留下負責任決定之空間；也就是不僅對其課以責任，也將責任託付於她。基於此理由，我們與第二庭不同（參見判決 D.IV.1.b），認為婦女在諮詢中亦得保持沉默之可能性，是諮詢談話完全開放性之一項重要要素。而若她應負責任地決定，則其決定必須在無法律保留之下被承認。只有如此，她才能在諮詢中真正開放。
\end{transpara}

\begin{sourcepara}[title={原文SondervotumMS.295}]
\sourcern{405}b) Ausgangspunkt für eine gesetzliche Regelung muß sein, daß die schwangere Frau -- selbstverständlich -- zu einer verantwortlichen, für eine Rechtfertigung tragfähigen Entscheidung generell fähig ist. Spätestens nach der obligatorischen Beratung ist ihr auch der mit einem Schwangerschaftsabbruch verbundene Konflikt bewußt. Sie weiß, um welch hohes Gut es sich bei dem in ihr heranwachsenden Ungeborenen handelt. Frauen entschließen sich nicht leichten Herzens und ohne Grund zu einem solchen Eingriff (vgl. Sondervotum Rupp-von Brünneck/Simon, \abbrBVerfGE{} 39, 68 [88]). Die schwangere Frau ist -- was sich ebenfalls von selbst versteht -- bei der Entscheidung über Fortsetzung oder Abbruch einer Schwangerschaft nicht gleichgültig gegenüber Recht und Unrecht. In ihrem Schwangerschaftskonflikt treffen -- gerade für sie erfahrbar -- fundamentale Rechtsgüter aufeinander. Sie will sich hier auch mit einer gegen das Austragen des Kindes getroffenen Entscheidung nicht ausgegrenzt sondern akzeptiert sehen. Sie beansprucht mit dem Abbruch auch nicht etwa nur Vorteile für sich; vielmehr bedeutet für sie die Abwendung von der Erwartung, ein Kind zu bekommen und dem Wunsch, es eigentlich auch haben zu wollen, in jedem Falle eine schwerwiegende Selbstverletzung und einen sie existentiell berührenden Eingriff; das haben auch die in der mündlichen Verhandlung gehörten Beraterinnen aus ihrer Praxis so mitgeteilt.
\end{sourcepara}
\begin{transpara}
\transrn{405}b) 法律規制之出發點必須是：懷孕婦女當然通常有能力作成一項負責任、足以承載正當化之決定。至遲在義務諮詢後，她也已意識到終止妊娠所伴隨之衝突。她知道在其體內成長之未出生者是何等崇高之法益。婦女並非輕率且毫無理由地決定接受此種干預（參見 Rupp-von Brünneck/Simon 不同意見，BVerfGE 39, 68 [88]）。懷孕婦女在決定繼續或終止妊娠時，當然也並非對法與不法漠不關心。在其懷孕衝突中，基本法益彼此碰撞，而這正是她能切身經驗者。即使她作成反對繼續懷孕之決定，她也希望自己不被排除，而是被接受。她藉終止妊娠也並非只是為自己主張利益；毋寧，對她而言，背離將有一個孩子之期待以及其實也想要孩子之願望，在任何情況下都是嚴重之自我傷害，並是一項觸及其存在之干預；在言詞辯論中所聽取之諮詢員亦從其實務如此陳述。
\end{transpara}

\begin{sourcepara}[title={原文SondervotumMS.296}]
\sourcern{406}Aus alledem darf der Gesetzgeber mit dem Gewicht eines tragenden Indizes und ohne die Notwendigkeit weiterer Verifizierung schließen, daß dem Abbruchverlangen eine Konfliktsituation zugrundeliegt, in der sich schutzwürdige Interessen der Frau mit solcher Dringlichkeit geltend machen, daß die staatliche Rechtsordnung nicht verlangen kann, die Schwangere müsse dem Recht des Ungeborenen dennoch den Vorrang einräumen (vgl. \abbrBVerfGE{} 39, 1 [50]). Darin liegt nicht die Anerkennung eines ungebundenen "Selbstbestimmungsrechts" der Frau und auch nicht die Aufgabe des rechtlichen Schutzes des ungeborenen Lebens, wie sich schon aus der Strafdrohung des § 218 \abbrStGB{} \abbrNF{} gegen Abbrüche ohne Beratung ergibt. Die gegenteilige Auffassung des Senats (vgl. Urteil, D.III.2.b), aus dem auch nach Beratung fortbestehenden Abbruchverlangen dürfe der Gesetzgeber nicht auf eine Notlage schließen, führt in ein Dilemma: Entweder traut man der Beratung keinen wirklichen Einfluß oder der schwangeren Frau keine verantwortliche Entscheidung zu. Beides hat der Senat aber zugestanden. Ein dritter Grund bleibt nicht.
\end{sourcepara}
\begin{transpara}
\transrn{406}立法者得由以上一切，以承載性指標之重量且無須進一步驗證，推論終止妊娠要求背後存在一種衝突情境；在該情境中，婦女值得保護之利益以如此急迫性出現，以致國家法律秩序不能要求懷孕婦女仍必須使未出生者之權利居於優先（參見 BVerfGE 39, 1 [50]）。此並非承認婦女不受拘束之「自決權」，亦非放棄對未出生生命之法律保護；單從新版本《刑法》第218條對未經諮詢之終止妊娠所設刑罰威嚇，即可看出。第二庭相反見解（參見判決 D.III.2.b）認為，立法者不得由即使經諮詢後仍持續存在之終止妊娠要求推論困境，這導向一個兩難：不是不相信諮詢有真正影響，就是不相信懷孕婦女能作負責任決定。但第二庭已承認二者。不存在第三個理由。
\end{transpara}

\begin{sourcepara}[title={原文SondervotumMS.297}]
\sourcern{407}Die Gefahr von Mißbräuchen steht dem nicht entgegen. Jede Freiheit zu verantwortlicher Entscheidung, ohne die ein Lebensschutz durch Beratung nicht denkbar ist, schließt die Möglichkeit des Mißbrauchs mit ein; lückenlose Mißbrauchsabwehr müßte diese Freiheit in ihren lebensschützenden Möglichkeiten wieder aufheben. Abgesehen davon haben sich auch Indikationstatbestände mißbrauchen lassen.
\end{sourcepara}
\begin{transpara}
\transrn{407}濫用之危險並不與此相牴觸。任何負責任決定之自由，均包含濫用之可能；而若無此自由，經由諮詢之生命保護不可想像。毫無漏洞之濫用防止，將會重新取消此自由在生命保護上之可能性。除此之外，指標事由構成要件也曾被濫用。
\end{transpara}

\begin{sourcepara}[title={原文SondervotumMS.298}]
\sourcern{408}c) Bei der Regelung des Schwangerschaftsabbruchs steht das Grundgesetz einer besonderen Rechtfertigung auf der Grundlage tragender Indizien auch deshalb nicht entgegen, weil der Schwangerschaftsabbruch in der Frühphase sich von vornherein der herkömmlichen Unterscheidung von Tun oder Unterlassen (vgl. § 13 \abbrStGB{}) entzieht; er ist mit einer Wertung als Rechtsgutverletzung durch aktives Tun allein nicht hinreichend erfaßt, vielmehr lehnt die Frau, in deren Person Rechtsgüter in singulärer Weise einander widerstreiten und sich zugleich vereinen, unter Berufung auf die Grenzen der Aufopferung eigener Rechte vornehmlich die Übernahme einer mit besonders schwerwiegenden Einstands- und Sorgfaltspflichten verbundenen Garantenstellung ab.
\end{sourcepara}
\begin{transpara}
\transrn{408}c) 在終止妊娠規制中，《基本法》之所以亦不反對以承載性指標為基礎之特殊正當化，是因為早期終止妊娠從一開始即逸脫作為或不作為之傳統區分（參見《刑法》第13條）；僅將其評價為以積極作為侵害法益，並不足以掌握其全貌。毋寧，在婦女此一人格中，法益以特殊方式彼此衝突且同時結合，而婦女主張自身權利犧牲之界限，主要是拒絕承擔一項伴隨特別重大擔保與照護義務之保證人地位。
\end{transpara}

\begin{sourcepara}[title={原文SondervotumMS.299}]
\sourcern{409}2. Die verfassungsrechtliche Schutzpflicht verlangt wirksamen Schutz für das ungeborene Leben (vgl. Urteil, D.I.4.). Wird der Schutz auf eine Beratungsregelung gestellt, so kann dadurch, daß der Frau die Rechtfertigung für ihr Handeln versagt bleibt, solcher Schutz nicht bewirkt werden. Das gilt sowohl für den unmittelbaren Schutz des je einzelnen Embryos (a) als auch für einen mittelbaren Schutz des ungeborenen menschlichen Lebens durch Aufrechterhaltung und Stärkung des allgemeinen Rechtsbewußtseins (b).
\end{sourcepara}
\begin{transpara}
\transrn{409}2. 憲法保護義務要求對未出生生命提供有效保護（參見判決 D.I.4.）。若保護建立於諮詢規制之上，則透過否定婦女行為之正當化，無法達成此種保護。這既適用於對每一個別胚胎之直接保護（a），亦適用於透過維持並強化一般法律意識而對未出生人類生命為間接保護（b）。
\end{transpara}

\begin{sourcepara}[title={原文SondervotumMS.300}]
\sourcern{410}a) Die Erfüllung verfassungsrechtlicher Schutzpflichten erfolgt wie die Auflösung von Grundrechtskollisionen durch Gesetzesrecht. Ohne Ausprägung in Gesetzesrecht läuft leer, was sich aus der verfassungsrechtlichen Schutzpflicht ergeben soll. Die Beratungsregelung des Schwangeren- und Familienhilfegesetzes stellt in der Frühphase der Schwangerschaft nur den Abbruch ohne Beratung unter Strafe. Damit gewährleistet der Gesetzgeber rechtlichen Schutz für das ungeborene Leben -- auch mit den Mitteln des Strafrechts --, wenn auch auf andere Weise als mit der Indikationenregelung. Der Senat hält es dennoch auch im Rahmen der Beratungsregelung für erforderlich, daß als orientierende Leitlinie für die Beratung und für das allgemeine Rechtsbewußtsein die Grenze der Rechtspflicht der Frau zum Austragen des Kindes gemäß dem Kriterium der Unzumutbarkeit normativ formuliert wird (vgl. Urteil, D.III.1.c). Das Urteil beruft sich dafür (vgl. Urteil, D.I.2.c.bb) auf die Entscheidung des Ersten Senats (vgl. \abbrBVerfGE{} 39, 1 [48 \abbrFF{}]), läßt dabei aber außer acht, daß dort die Unterscheidung zwischen erlaubten und verbotenen Schwangerschaftsabbrüchen ihre folgerichtige Umsetzung in der Formulierung von Indikationstatbeständen gefunden hatte. Diese Funktion des Unzumutbarkeitsurteils entfällt -- soweit es um die allgemeine Notlage geht -- bei der Beratungsregelung. Insoweit hat für die rechtliche Beurteilung des Abbruchs, den die Frau in der Frühphase der Schwangerschaft nach Beratung verlangt und durchführen läßt, der gesetzliche Maßstab nach der Auffassung des Senats folgenlos zu bleiben; der Frau wird die Rechtfertigung für ihr Handeln -- wegen vermeintlich zwingender Erfordernisse des Schutzkonzepts -- in jedem Falle vorenthalten (vgl. oben I.4.). Der Senat verlagert die von ihm geforderten Rechtsfolgen (vgl. Urteil, D.III.2.a) in die Interpretation des § 24b \abbrSGBV{} (vgl. Urteil, E.V.2.b). Unabhängig davon, daß sich die Versagung von Leistungen aus der gesetzlichen Krankenversicherung als Kennzeichnung der rechtlichen Mißbilligung nicht eignet, kann hierdurch aber -- so auch der Senat -- das einzelne ungeborene Leben nicht unmittelbar geschützt werden.
\end{sourcepara}
\begin{transpara}
\transrn{410}a) 憲法保護義務之履行，如同基本權衝突之解消，是透過法律規範完成。若未在法律規範中成形，所謂由憲法保護義務得出之內容即會落空。《孕婦及家庭扶助法》之諮詢規制，在懷孕早期僅處罰未經諮詢之終止妊娠。藉此，立法者雖以不同於指標規則之方式，仍為未出生生命保障法律保護，且也使用刑法手段。然而第二庭認為，即使在諮詢規制架構內，仍有必要將婦女繼續懷孕之法律義務界限，依不可期待性標準作規範表述，作為諮詢與一般法律意識之指引準則（參見判決 D.III.1.c）。判決為此援引第一庭裁判（參見判決 D.I.2.c.bb；另參見 BVerfGE 39, 1 [48 ff.]），但忽略該處對允許與禁止終止妊娠之區分，已在指標事由構成要件之表述中獲得合乎邏輯之落實。不可期待性判斷之此項功能，在諮詢規制中，就一般困境而言已不復存在。依第二庭見解，就婦女在懷孕早期經諮詢後要求並實施之終止妊娠之法律評價而言，法定標準就此必須不產生後果；婦女行為之正當化，在任何情況下均因據稱為保護構想所強制要求而被拒絕（參見上文 I.4.）。第二庭將其所要求之法律後果（參見判決 D.III.2.a）移轉至對 SGB V 第24b條之解釋（參見判決 E.V.2.b）。且不論拒絕法定醫療保險給付並不適合作為法律否定評價之標示，透過此種方式，如第二庭亦承認，個別未出生生命並不能受到直接保護。
\end{transpara}

\begin{sourcepara}[title={原文SondervotumMS.301}]
\sourcern{411}Eine unmittelbare Schutzwirkung für das ungeborene Leben kann auch von der "normativen Orientierung" der schwangeren Frau auf den Schutz dieses Lebens hin nicht ausgehen. Der Senat sieht es -- wie hervorgehoben -- als Bestandteil des Schutzkonzepts der Beratungslösung an, der einzelnen Frau die Gewißheit zu verweigern, daß der Abbruch, für den sie sich nach Abschluß des Beratungsverfahrens entscheidet, von der Rechtsordnung gebilligt ist (vgl. Urteil, D.III.2.b.cc). Das halten wir für einen Rückschritt gegenüber der bisherigen Rechtslage nach der Indikationenregelung. Die Frau zahlt mit der Versagung der Rechtfertigung gleichsam den Preis für das neue Schutzkonzept. An die Stelle der rechtlichen Klarheit einer Norm tritt für sie eine Lage, in der ihr Handeln in jedem Falle -- trotz einer womöglich gegebenen schweren Notlage -- als nicht erlaubt behandelt wird. Das steht nach unserer Überzeugung mit dem Grundgedanken der Beratungsregelung in einem unvereinbaren Widerspruch, der die Wirksamkeit des Schutzkonzepts insgesamt in Frage stellt (vgl. oben unter 1.). Wir halten es aber auch für nicht angängig, daß der Senat den Staat bei der Erfüllung seiner Schutzpflicht gegenüber dem ungeborenen Leben in verfassungsrechtliche Grenzen verweist, zu deren Wirkungen es gehören soll, der schwangeren Frau die Antwort darauf zu verweigern, ob sie mit einem Abbruch recht handelt. Das stellt sie in die Nähe der Rolle eines unmündigen Objekts im Gefüge des staatlichen Lebensschutzkonzepts. Das Rechtsbewußtsein der Frau kann unseres Erachtens kaum einschneidender geschwächt werden als dadurch, daß sie im Rechtssinne nicht weiß, was sie tut; eben dies ist aber die Folge, wenn ihre verantwortlich getroffene Entscheidung trotz der im Gesetz formulierten normativen Leitlinie, an der sie sich orientieren soll, nicht als gerechtfertigt anerkannt wird. Hieran lassen sich auch rechtsstaatliche Bedenken knüpfen, denen wir aber nicht weiter nachgehen.
\end{sourcepara}
\begin{transpara}
\transrn{411}懷孕婦女朝向保護此一生命之「規範性指引」，也無法對未出生生命產生直接保護效果。如前所強調，第二庭將拒絕使個別婦女確信其在諮詢程序完成後所決定之終止妊娠受到法律秩序肯認，視為諮詢方案保護構想之一部分（參見判決 D.III.2.b.cc）。我們認為，這相較於既有指標規則之法律狀態是一種倒退。婦女彷彿以正當化被拒絕作為新保護構想之代價。對她而言，取代法律規範明確性的是一種狀態：其行為在任何情況下，即使可能存在嚴重困境，仍被視為不被允許。依我們確信，這與諮詢規制之基本思想發生不可相容之矛盾，並使保護構想整體之有效性受到質疑（參見上文 1.）。我們也認為，第二庭不能在國家履行其對未出生生命之保護義務時，將國家指向某些憲法界限，而此等界限之效果竟應包括拒絕回答懷孕婦女其終止妊娠是否依法行事。這使她接近於國家生命保護構想結構中未成年客體之角色。依我們見解，婦女之法律意識幾乎不可能以比此更深刻之方式被削弱：她在法律意義上不知道自己正在做什麼；而這正是當其負責任作成之決定，儘管有法律所表述、她應據以取向之規範準則，卻不被承認為正當化時之後果。此處亦可連結法治國疑慮，但我們不再進一步追究。
\end{transpara}

\begin{sourcepara}[title={原文SondervotumMS.302}]
\sourcern{412}b) Nach der Auffassung des Senats ist die Versagung der Rechtfertigung auch deshalb erforderlich, damit nicht im allgemeinen Rechtsbewußtsein der Eindruck entsteht, der nach Beratung vorgenommene Abbruch sei erlaubt (vgl. Urteil, D.III.2.a). Wir meinen dagegen, daß rechtliche Mißbilligungen außerhalb des Strafrechts im Bereich des Lebensschutzes die Rechtsüberzeugungen der Bevölkerung nicht eigenständig prägen.
\end{sourcepara}
\begin{transpara}
\transrn{412}b) 依第二庭見解，拒絕正當化亦屬必要，以免一般法律意識中產生經諮詢後實施之終止妊娠係被允許之印象（參見判決 D.III.2.a）。我們則認為，在生命保護領域中，刑法之外之法律否定評價並不能獨立形塑人民之法律信念。
\end{transpara}

\begin{sourcepara}[title={原文SondervotumMS.303}]
\sourcern{413}Unseres Erachtens speist sich das Rechtsbewußtsein gerade im Bereich des Schwangerschaftsabbruchs aus den individuellen sittlichen Überzeugungen, die durch Erziehung, eigenes Schicksal und gesellschaftliche Wertvorstellungen mitgeprägt sind. Ein Unwerturteil, wie es dem Senat vorschwebt, bewirkt hier nichts. Schon die Strafverfolgungen und Verurteilungen wegen Abtreibung nach §§ 218 \abbrFF{} \abbrStGB{} bisheriger Fassung tendierten gegen Null. Der Deutsche Bundestag hat vorgetragen, nach Befragungen des Allensbach-Instituts in den Jahren 1983, 1987 und 1988 hielten zwei Drittel der Deutschen den Schwangerschaftsabbruch für erlaubt. Das Strafrecht hat nach dem neuen Konzept einen akzentuierten repressiven und einen entsprechend genau konturierten präventiven Bereich (vgl. § 218 \abbrStGB{} \abbrNF{} einerseits, §§ 218a, 219 \abbrStGB{} \abbrNF{} andererseits). Wo das Recht prägende Kraft entfalten soll, kommt es auf solche Markierungen an und nicht auf Auslegungen im Sozialversicherungsrecht; von diesen geht keine prägende Kraft aus. Im Gegenteil: Wenn der Gesetzgeber im Rahmen eines auf Beratung und Letztverantwortung der Frau setzenden Schutzkonzepts die Strafdrohung zurücknimmt, weil sie sich zum Schutz des ungeborenen Lebens als stumpfes Schwert erwiesen hat, so muß es das Rechtsbewußtsein der Allgemeinheit verunsichern, wenn an die Stelle einer Strafdrohung, die ein Rechtsgut von höchstem Rang schützen soll, nun Rechtsnachteile im Sozialversicherungsrecht als einem bloßen Folgerecht treten, die nach Auffassung des Senats jetzt die Last des verfassungsrechtlichen Unwerturteils tragen sollen.
\end{sourcepara}
\begin{transpara}
\transrn{413}依我們見解，正是在終止妊娠領域中，法律意識由個人倫理信念供給，而此等信念受教育、個人命運及社會價值觀共同形塑。第二庭所構想之負面價值判斷，在此不起作用。即使依迄今版本《刑法》第218條以下對終止妊娠之刑事追訴與定罪，已趨近於零。德國聯邦議院陳述，依 Allensbach 研究所於1983年、1987年及1988年之調查，三分之二德國人認為終止妊娠是被允許的。依新構想，刑法具有一個重點鮮明之壓制領域，及一個相應精確勾勒之預防領域（參見一方面新版本《刑法》第218條，另一方面新版本《刑法》第218a條、第219條）。凡法律應發揮形塑力量之處，關鍵在於此類標記，而非社會保險法中之解釋；後者不具有形塑力量。相反地，若立法者在一項仰賴諮詢與婦女最後責任之保護構想中撤回刑罰威嚇，因其已證明對保護未出生生命而言是鈍劍，則當取代原本應保護最高位階法益之刑罰威嚇者，竟是社會保險法此一單純後續法律中之法律不利益，且依第二庭見解，該不利益現在應承擔憲法負面價值判斷之重量時，必將動搖一般人之法律意識。
\end{transpara}

\begin{sourcepara}[title={原文SondervotumMS.304}]
\sourcern{414}Schon grundrechtsdogmatisch halten wir es für eine Überdehnung der verfassungsrechtlichen Schutzpflicht für das ungeborene Leben, wenn man ihr abverlangt, sie solle das allgemeine Rechtsbewußtsein formen. Es ist zudem eine rechtstatsächliche Frage, ob und bejahendenfalls in welcher Weise praktisch sanktionslose Verfassungsnormen das Rechtsbewußtsein prägen. Gesicherte Erkenntnisse darüber fehlen; bloße Überzeugungen der Interpreten können nicht maßgebend sein. Die Annahme des Senats halten wir nicht einmal für plausibel. Mit dem gleichen Recht ließe sich behaupten, es müsse die allgemeine Rechtsüberzeugung verwirren, wenn Handlungen gegen das Rechtsgut des ungeborenen menschlichen Lebens zwar von Verfassungs wegen in bestimmten Umfang verboten sind, die Grenzen der Reichweite des Verbots für die betroffene Frau aber aus der Natur des Schutzkonzepts heraus nicht zu einer Rechtfertigung ihres Handelns führen dürfen und wegen dieser Unklarheit -- ohne Ansehung einer im Einzelfall gegebenen Notlage -- Leistungen aus der gesetzlichen Krankenversicherung verweigert werden.
\end{sourcepara}
\begin{transpara}
\transrn{414}即使從基本權教義學上看，我們也認為，若要求對未出生生命之憲法保護義務形塑一般法律意識，即是過度擴張該保護義務。此外，實際上無制裁之憲法規範是否以及若是則以何種方式形塑法律意識，乃法律事實問題。對此欠缺可靠認識；解釋者之單純信念不能作為決定性依據。我們甚至認為第二庭之假設並不可信。以同樣理由也可主張：若針對未出生人類生命此一法益之行為，雖依憲法在一定範圍內被禁止，但該禁止射程之界限，對受影響婦女而言依保護構想之性質卻不得導致其行為被正當化，並且因這種不明確性而不顧個案中是否存在困境即拒絕法定醫療保險給付，必將混淆一般法律信念。
\end{transpara}

\begin{sourcepara}[title={原文SondervotumMS.305}]
\sourcern{415}Auch der Erste Senat hat lediglich von der rechtsbewußtseinsbildenden Kraft der Strafnormen gesprochen und zwar im Blick auf den "Versuch", durch eine differenzierte strafrechtliche Regelung einen besseren Lebensschutz zu erreichen (vgl. \abbrBVerfGE{} 39, 1 [65 \abbrF{}]). Er traute auch der bloßen Existenz von Strafandrohungen Einfluß auf die Wertvorstellungen und Verhaltensweisen der Bevölkerung zu (vgl. \abbrBVerfGE{} \abbrAAO{}, \abbrS{} 57;vgl. auch \abbrBVerfGE{} 45, 187 [254 \abbrFF{}]). Darum geht es bei dem hier zu prüfenden Gesetz jedoch nicht; der Gesetzgeber hat diesen "Versuch" mit dem Strafrecht aufgegeben; das billigt auch der Senat.
\end{sourcepara}
\begin{transpara}
\transrn{415}第一庭也僅曾談及刑罰規範形成法律意識之力量，而且是著眼於透過區分性刑法規制達成較佳生命保護之「嘗試」（參見 BVerfGE 39, 1 [65 f.]）。其亦相信單純刑罰威嚇之存在，會影響人民之價值觀與行為方式（參見 BVerfGE a.a.O., S. 57；另參見 BVerfGE 45, 187 [254 ff.]）。然而，本案所審查之法律並非關於此；立法者已放棄以刑法作此「嘗試」；第二庭亦予認可。
\end{transpara}

\begin{sourcepara}[title={原文SondervotumMS.306}]
\sourcern{416}In der mündlichen Verhandlung über die hier zu beurteilenden Anträge hat \abbrProf{} \abbrDr{} Eser darauf hingewiesen, daß im Gegensatz zur Kennzeichnung einer Handlung als "nicht rechtswidrig" die tatbestandliche Ausnahme von einem sonst mit Strafe bedrohten Verhalten als eine Freigabe des Schutzgutes in diesem Bereich verstanden werden müsse (vgl. auch Lenckner in: Schönke/Schröder, \abbrStGB{}, 24. \abbrAufl{}, 1991, Vorbem. zu §§ 13 \abbrFF{} \abbrRdnr{} 17). Der Erste Senat hat es als für das Rechtsbewußtsein unerheblich angesehen, ob der damalige § 218a \abbrStGB{} den Tatbestand des § 218 \abbrStGB{} einengte oder ob er einen Rechtfertigungsgrund setzte oder nur einen Schuld- oder Strafausschließungsgrund zum Inhalt hatte; in jedem Falle müsse der Eindruck entstehen, der Abbruch sei "rechtlich erlaubt" (vgl. \abbrBVerfGE{} 39, 1 [53]). Der Zweite Senat sieht das anders. Auch hieran zeigt sich, wie wenig Sicherheit in der Handhabung der Mittel zur Prägung des Rechtsbewußtseins besteht. Es kann angesichts dessen nicht Sache des Bundesverfassungsgerichts sein, über die Alternative, Handlungen aus dem Tatbestand auszunehmen oder sie für nicht rechtswidrig zu erklären -- und damit über spezifisch strafrechtsdogmatische Fragen -- kategorial im Hinblick auf die Erfüllung der staatlichen Schutzpflicht zu entscheiden. Dies ist Sache des Gesetzgebers und der zuständigen Gerichte.
\end{sourcepara}
\begin{transpara}
\transrn{416}在關於本案聲請之言詞辯論中，Prof. Dr. Eser 指出，與將一行為標示為「不具違法性」相反，將原本受刑罰威嚇之行為排除於構成要件之外，必然被理解為在此領域放開該保護法益（另參見 Lenckner in: Schönke/Schröder, StGB, 24. Aufl., 1991, Vorbem. zu §§ 13 ff. Rdnr. 17）。第一庭認為，當時《刑法》第218a條究竟限縮第218條構成要件，或設置正當化事由，或僅包含阻卻罪責或阻卻刑罰事由，對法律意識並不重要；無論如何都必然產生終止妊娠「法律上被允許」之印象（參見 BVerfGE 39, 1 [53]）。第二庭則另作看法。由此也可見，在操作形塑法律意識之手段時，存在多麼少的確定性。有鑑於此，聯邦憲法法院不應就將行為排除於構成要件之外或宣告其不具違法性此一選擇，亦即就特定刑法教義學問題，從履行國家保護義務角度作範疇性決定。這是立法者及管轄法院之事。
\end{transpara}

\bisubsection{III.}{III.}

\begin{sourcepara}[title={原文SondervotumMS.307}]
\sourcern{417}1. Der Senat ist der Auffassung, daß die Vorschrift des § 218a \abbrAbs{} 1 \abbrStGB{} \abbrNF{} als Rechtfertigungsgrund verfassungswidrig sei, weil sie eine Not- und Konfliktlage, in der der schwangeren Frau das Austragen des Kindes nicht zugemutet werden kann, weder voraussetze noch näher qualifiziere und die Rechtfertigung auch nicht von entsprechenden Feststellungen durch Dritte abhängig mache (vgl. Urteil, E.I.2.). Nach unseren Darlegungen zu I. und II. halten wir insoweit § 218a \abbrAbs{} 1 \abbrStGB{} \abbrNF{} für verfassungsgemäß; soweit die Vorschrift für Schwangerschaftsabbrüche, die unter den dort näher geregelten Voraussetzungen erfolgen, einen allgemeinen Rechtfertigungsgrund schafft, verstößt sie nicht gegen die staatliche Schutzpflicht aus \abbrArt{} 1 \abbrAbs{} 1 \abbrIVm{} \abbrArt{} 2 \abbrAbs{} 2 \abbrGG{}\ggfnArtOneOneTwoTwoDE{}. Soweit der Senat § 218a \abbrAbs{} 1 \abbrStGB{} \abbrNF{} auch wegen des Zusammenhangs mit § 219 \abbrStGB{} \abbrNF{} im Hinblick auf die dem Absatz 2 jener Vorschrift anhaftenden Mängel der Regelung von Organisation und Beaufsichtigung der Beratungseinrichtungen für verfassungswidrig erklärt (vgl. Urteil, E.I.4. und II.1.), tragen wir die Entscheidung mit; auch die Regelung der vom Grundgesetz gebotenen Elemente für Inhalt und Ziel der Beratung mag in § 219 \abbrAbs{} 1 \abbrStGB{} \abbrNF{} nicht mit der Klarheit und Eindeutigkeit zum Ausdruck gekommen sein, die für diese Regelung angezeigt ist (vgl. dazu Urteil, E.II.2.).
\end{sourcepara}
\begin{transpara}
\transrn{417}1. 第二庭認為，新版本《刑法》第218a條第1項作為正當化事由違憲，因為其既未以懷孕婦女不可被期待繼續懷孕之緊急與衝突狀態為前提，亦未進一步限定其性質，且也未使正當化取決於第三人之相應確認（參見判決 E.I.2.）。依我們於 I. 及 II. 之說明，我們就此認為新版本《刑法》第218a條第1項合憲；該規定對於在其中詳定前提下實施之終止妊娠創設一般正當化事由，並不違反《基本法》第1條第1項結合第2條第2項\ggfnArtOneOneTwoTwoZH{}所生之國家保護義務。至於第二庭亦因新版本《刑法》第218a條第1項與新版本《刑法》第219條之關聯，鑑於後者第2項在諮詢機構組織與監督規制上所附著之瑕疵而宣告其違憲（參見判決 E.I.4. 及 II.1.），我們支持此裁判；至於《基本法》所要求之諮詢內容與目標要素之規制，或許也未在新版本《刑法》第219條第1項中以此規制所適合之清楚性與明確性表達出來（就此參見判決 E.II.2.）。
\end{transpara}

\begin{sourcepara}[title={原文SondervotumMS.308}]
\sourcern{418}2. Leistungen der Sozialversicherung für Schwangerschaftsabbrüche, die in den ersten zwölf Wochen nach der Empfängnis nach Beratung von einem Arzt vorgenommen werden, stehen nach unserer Auffassung nicht im Widerspruch zum Grundgesetz. Wir halten es freilich nicht etwa für verfassungsrechtlich geboten, daß der Gesetzgeber über medizinisch indizierte Abbrüche hinaus allgemein die ärztliche Vornahme eines nicht rechtswidrigen Schwangerschaftsabbruchs in den Leistungskatalog der gesetzlichen Krankenversicherung aufnimmt. Tut er es aber, wofür im Blick auf den Schutz der Gesundheit der Frau gute Gründe sprechen, so muß er die gesetzlich festgelegte Leistungspflicht der in den Krankenkassen organisierten Solidargemeinschaft frei von sachlich nicht gerechtfertigten Differenzierungen ausgestalten. Die im Rahmen einer Beratungsregelung erfolgenden Schwangerschaftsabbrüche können dann nicht ausgenommen werden. § 24b \abbrSGBV{} unterliegt daher nicht der einschränkenden Auslegung, die der Senat (vgl. Urteil, E.V.2.b) als geboten erachtet. Das folgt für uns bereits daraus, daß Schwangerschaftsabbrüche nach Beratung insgesamt erlaubt sind und als gerechtfertigt zu behandeln sind oder doch so behandelt werden dürfen.
\end{sourcepara}
\begin{transpara}
\transrn{418}2. 依我們見解，對於在受孕後最初十二週內經諮詢後由醫師實施之終止妊娠，給予社會保險給付，並不牴觸《基本法》。當然，我們並不認為，除醫學上有指標事由之終止妊娠外，立法者在憲法上被要求一般性地將醫師實施之不具違法性之終止妊娠納入法定醫療保險給付目錄。但若立法者如此為之，而鑑於保護婦女健康確有充分理由支持此舉，則其必須使組織於醫療保險基金中之團結共同體依法確定之給付義務，不含事物上無正當理由之差別待遇。在諮詢規制框架內發生之終止妊娠即不得被排除。因此，SGB V 第24b條不受第二庭（參見判決 E.V.2.b）認為必要之限縮解釋拘束。對我們而言，這已可由以下點得出：經諮詢後之終止妊娠整體上被允許，並應被作為正當化處理，或至少得如此處理。
\end{transpara}

\begin{sourcepara}[title={原文SondervotumMS.309}]
\sourcern{419}Wir teilen auch nicht die von der Bayerischen Staatsregierung unter dem Gesichtspunkt, der Sozialversicherung dürften als öffentlich-rechtlichem Zwangsverband mit Rücksicht auf die Grundrechte ihrer Mitglieder keine versicherungsfremden Lasten aufgebürdet werden, vorgebrachten verfassungsrechtlichen Bedenken, über die der Senat in der Sache nicht zu entscheiden brauchte (vgl. Urteil, E.V.5.b). Leistungen der Krankenkassen für von der Rechtsordnung erlaubte Schwangerschaftsabbrüche können Grundrechte der beitragszahlenden Mitglieder -- unbeschadet ihrer abweichenden ethischen oder moralischen Überzeugungen -- nicht verletzen (vgl. dazu schon \abbrBVerfGE{} 78, 320 [331]).
\end{sourcepara}
\begin{transpara}
\transrn{419}我們也不同意巴伐利亞邦政府所提出之憲法疑慮；其觀點是，社會保險作為公法上強制團體，基於其成員基本權，不得被加諸保險外負擔，而第二庭在實體上無須對此作成裁判（參見判決 E.V.5.b）。醫療保險基金對法律秩序所允許之終止妊娠所為給付，不會侵害繳納保費成員之基本權，無論其倫理或道德信念是否不同（就此已參見 BVerfGE 78, 320 [331]）。
\end{transpara}

\begin{sourcepara}[title={原文SondervotumMS.310}]
\sourcern{420}Der Auffassung der Senatsmehrheit, daß § 24b \abbrSGBV{} auf Schwangerschaftsabbrüche nach Beratung nicht angewendet werden dürfe, wäre aber auch dann nicht zu folgen, wenn diese Abbrüche -- entsprechend dem Urteil -- nicht für gerechtfertigt (nicht rechtswidrig) erklärt werden dürften. Insofern schließen wir uns der abweichenden Meinung des Kollegen Böckenförde an.
\end{sourcepara}
\begin{transpara}
\transrn{420}即使此等終止妊娠依判決不得被宣告為正當化（不具違法性），也仍不應追隨第二庭多數認為 SGB V 第24b條不得適用於經諮詢後終止妊娠之見解。就此，我們贊同同僚 Böckenförde 之不同意見。
\end{transpara}

\bisubsection{IV.}{IV.}

\begin{sourcepara}[title={原文SondervotumMS.311}]
\sourcern{421}Gegen die maßstäblichen Ausführungen des Senats zur Einbeziehung des Arztes in das Schutzkonzept der Beratungsregelung (vgl. Urteil, D.V.) haben wir in zwei Punkten Bedenken.
\end{sourcepara}
\begin{transpara}
\transrn{421}對於第二庭關於將醫師納入諮詢規制保護構想之標準性說明（參見判決 D.V.），我們有兩點疑慮。
\end{transpara}

\begin{sourcepara}[title={原文SondervotumMS.312}]
\sourcern{422}1. Die verfassungsrechtliche Schutzpflicht fordert unseres Erachtens -- abweichend von der Auffassung des Senats (vgl. Urteil, D.V.2.c) -- nicht, daß die Verletzung ärztlicher Pflichten im Zusammenhang mit einem Schwangerschaftsabbruch und der ihm vorausgehenden Beratung und Aufklärung der Frau vom Gesetzgeber unter Strafsanktion gestellt wird.
\end{sourcepara}
\begin{transpara}
\transrn{422}1. 依我們見解，憲法保護義務並不如第二庭見解（參見判決 D.V.2.c）所認為，要求立法者將醫師在終止妊娠及其前置之婦女諮詢與告知相關脈絡中違反義務之行為置於刑罰制裁之下。
\end{transpara}

\begin{sourcepara}[title={原文SondervotumMS.313}]
\sourcern{423}2. Die Verfahrensgegenstände gaben keine Veranlassung zu den Ausführungen im Urteil, wonach die Unterhaltspflicht für ein Kind niemals ein Schaden sein könne (vgl. Urteil, D.V.6.). Sie sind ein obiter dictum und entbehren darüber hinaus der erforderlichen Auseinandersetzung mit den eingehenden Ausführungen, mit denen der VI. Zivilsenat des Bundesgerichtshofs begründet hat, unter welchen -- dort eingegrenzten -- Voraussetzungen die Möglichkeit eines Vermögensschadens bestehen kann (BGHZ 76, 249 [253 \abbrFF{}]; BGH, NJW 1984, \abbrS{} 2625 \abbrF{}).
\end{sourcepara}
\begin{transpara}
\transrn{423}2. 本案程序標的並未提供理由，使判決說明對一名子女之扶養義務絕不可能構成損害（參見判決 D.V.6.）。此等說明係 obiter dictum，且除此之外也欠缺對聯邦最高法院第六民事庭詳細論述之必要處理；該庭曾說明在何等於該處加以限定之前提下，可能存在財產損害（BGHZ 76, 249 [253 ff.]；BGH, NJW 1984, S. 2625 f.）。
\end{transpara}

\begin{sourcepara}[title={原文SondervotumMS.314}]
Mahrenholz, Sommer
\end{sourcepara}
\begin{transpara}
Mahrenholz, Sommer
\end{transpara}

\bisection{Abweichende Meinung Böckenförde}{Böckenförde 不同意見}

\begin{sourcepara}[title={原文SondervotumB.315}]
Abweichende Meinung des Richters Böckenförde zum Urteil des Zweiten Senats vom 28. Mai 1993 -- 2 \abbrBvF{} 2/90 und 4, 5/92 --
\end{sourcepara}
\begin{transpara}
Böckenförde 法官就第二庭1993年5月28日判決 -- 2 BvF 2/90 及 4, 5/92 -- 之不同意見
\end{transpara}

\begin{sourcepara}[title={原文SondervotumB.316}]
\sourcern{424}Ich stimme den wesentlichen maßstäblichen Ausführungen des Urteils zu, insbesondere auch der Aussage, daß nicht indizierte Schwangerschaftsabbrüche, die in den ersten zwölf Wochen nach Beratung von einem Arzt vorgenommenen werden, nicht als "nicht rechtswidrig", mithin erlaubt angesehen werden dürfen (vgl. D.III. vor 1.). Ich vermag jedoch den Ausführungen des Urteils (unter E.V.2.b), teilweise antizipiert in D.III.1.c), wonach Sozialversicherungsleistungen für solche Abbrüche (im folgenden: beratene Abbrüche) von Verfassungs wegen ausgeschlossen sind, nicht zuzustimmen. Hierüber zu befinden, ist Sache des Gesetzgebers.
\end{sourcepara}
\begin{transpara}
\transrn{424}我贊同判決主要之標準性說明，尤其也贊同下列命題：未具指標事由、在最初十二週內經諮詢後由醫師實施之終止妊娠，不得被視為「不具違法性」，因而不得被視為允許（參見 D.III. 1. 前）。然而，我無法同意判決之說明（E.V.2.b 下，部分已於 D.III.1.c 預先呈現），即對此等終止妊娠（以下稱：經諮詢終止妊娠）之社會保險給付在憲法上被排除。對此作成決定，是立法者之事。
\end{transpara}

\begin{sourcepara}[title={原文SondervotumB.317}]
\sourcern{425}Dabei geht es nicht darum, daß solche Leistungen etwa verfassungsrechtlich geboten seien -- das sind sie sicher nicht --, sondern nur darum, ob solche Leistungen von vornherein verfassungsrechtlich untersagt sind.
\end{sourcepara}
\begin{transpara}
\transrn{425}此處所涉並非此等給付是否為憲法所要求；其當然不是。問題僅在於，此等給付是否自始即在憲法上被禁止。
\end{transpara}

\begin{sourcepara}[title={原文SondervotumB.318}]
\sourcern{426}1. Mit dem Senat bin ich aus den im Urteil dargelegten Gründen der Auffassung, daß die beratenen Abbrüche zwar von der Strafdrohung ausgenommen, nicht aber generell für gerechtfertigt (nicht rechtswidrig) erklärt werden können. Wie der Senat mit Recht sagt, können nur bestimmte Ausnahmelagen, die eine solche Belastung hervorrufen, daß sie eine rechtliche Pflicht zum Austragen des Kindes für die Frau als unzumutbar erscheinen lassen (vgl. D.I.2.c.bb); konkretisierend aufgenommen in D.III.1.c), dazu führen, Schwangerschaftsabbrüche als gerechtfertigt anzusehen. Daß diese Voraussetzungen bei beratenen Abbrüchen generell vorliegen, läßt sich weder aufgrund eines allgemeinen Erfahrungssatzes annehmen noch typischerweise unterstellen. Das gleiche gilt aber auch umgekehrt: Eine Vermutung, beratene Abbrüche entsprächen diesen Voraussetzungen generell nicht, ist ebenso wenig begründet. Die beratenen Abbrüche stellen mithin eine ungeschiedene Gesamtheit von -- materiell-rechtlich gesehen (unabhängig von der fehlenden Feststellung durch einen Dritten) -- gerechtfertigten und nicht gerechtfertigten Abbrüchen dar.
\end{sourcepara}
\begin{transpara}
\transrn{426}1. 我基於判決所列理由，與第二庭同樣認為，經諮詢終止妊娠雖被排除於刑罰威嚇之外，卻不能一般性地被宣告為正當化（不具違法性）。如第二庭正確指出，只有特定例外情境，造成此種負擔，以致對婦女而言繼續懷孕之法律義務顯得不可期待（參見 D.I.2.c.bb；具體化納入 D.III.1.c），才可能導致終止妊娠被視為正當化。此等前提一般存在於經諮詢終止妊娠，既不能基於一般經驗法則假定，也不能類型化推定。但反之亦然：推定經諮詢終止妊娠一般不符合此等前提，同樣欠缺理由。因此，經諮詢終止妊娠構成一個未區分之整體，其中包含在實體法上（不論是否欠缺第三人確認）正當化與未正當化之終止妊娠。
\end{transpara}

\begin{sourcepara}[title={原文SondervotumB.319}]
\sourcern{427}2. Damit stellt sich für das Sozialversicherungsrecht, das insoweit Folgerecht ist, die Frage, wie mit dieser ungeschiedenen Gesamtheit der beratenen Abbrüche im Blick auf die Gewährung von Versicherungsleistungen zu verfahren ist. Das Schutzkonzept der Beratungsregelung beruht darauf, wie der Senat darlegt (vgl. D.III.1.c), daß generell und ausnahmslos von der Feststellung des Vorliegens der Indikation einer sozialen Notlage abgesehen wird. Die Gesamtheit der beratenen Abbrüche kann also nicht nach rechtswidrigen und nicht rechtswidrigen Abbrüchen aufgelöst werden; sie muß und kann nur einheitlich behandelt werden. Der Senat sieht es als Erfordernis der verfassungsrechtlichen Schutzpflicht für das ungeborene menschliche Leben an, daß alle diese Abbrüche, weil deren Rechtmäßigkeit nicht festgestellt ist und konzeptbedingt nicht festgestellt werden kann, als nicht gerechtfertigte, mithin als Unrecht behandelt werden und demgemäß von Leistungen der Sozialversicherung ausgeschlossen bleiben müssen. Ein solcher Ausschluß mag ein mögliches Ergebnis gesetzgeberischer Regelung sein. Daß er jedoch verfassungsrechtlich zwingend vorgeschrieben ist und jede andere Lösung einen Verfassungsverstoß darstellt, läßt sich weder aus der Schutzpflicht für das ungeborene menschliche Leben noch sonst verfassungsrechtlich herleiten. Indem der Senat dies als notwendige Folge der gebotenen grundsätzlichen rechtlichen Mißbilligung des Schwangerschaftsabbruchs und des abstrakten Prinzips ansieht, daß der Staat sich grundsätzlich nicht an Vorgängen beteiligen darf, deren Rechtmäßigkeit nicht feststeht, negiert er Konsequenzen, die im Beratungskonzept selbst beschlossen liegen.
\end{sourcepara}
\begin{transpara}
\transrn{427}2. 因此，對在此屬後續法律之社會保險法而言，問題在於：鑑於保險給付之給予，應如何處理此一未區分之經諮詢終止妊娠整體。如第二庭所闡述（參見 D.III.1.c），諮詢規制之保護構想建立於一般且無例外地放棄確認社會困境指標事由是否存在。經諮詢終止妊娠之整體因而不能被分解為違法與不具違法性之終止妊娠；其必須且只能被統一處理。第二庭認為，憲法對未出生人類生命之保護義務要求，所有此等終止妊娠，因其合法性未被確認且依構想不能被確認，必須被作為未正當化、因而作為不法處理，並相應地必須被排除於社會保險給付之外。此種排除或許是立法規制之可能結果。但其在憲法上具有強制規定性，且任何其他解決方案均構成違憲，既不能由對未出生人類生命之保護義務，也不能由其他憲法理由推導出來。第二庭將此視為對終止妊娠所要求之原則性法律否定評價，以及國家原則上不得參與合法性未確定事件此一抽象原則之必要後果，因而否定了諮詢構想本身所包含之後果。
\end{transpara}

\begin{sourcepara}[title={原文SondervotumB.320}]
\sourcern{428}a) Die Gesamtheit beratener Abbrüche entzieht sich der bruchlosen Einordnung in die Alternative rechtswidrig -- nicht rechtswidrig; sie stellt gegenüber dieser Alternative als unscheidbare Gesamtheit ein aliud dar. Gleichwohl hält der Senat es für verfassungsrechtlich geboten, daß sie einheitlich als Unrecht -- mit den entsprechenden Folgen im Sozialversicherungsrecht -- behandelt werden.
\end{sourcepara}
\begin{transpara}
\transrn{428}a) 經諮詢終止妊娠之整體，無法無縫納入違法與不具違法性之二分選項；相對於此一選項，其作為不可分離之整體構成一個 \textit{aliud}。然而第二庭仍認為，在憲法上有必要將其統一作為不法，並附隨社會保險法上相應後果來處理。
\end{transpara}

\begin{sourcepara}[title={原文SondervotumB.321}]
\sourcern{429}Eine Notwendigkeit hierfür ergibt sich nicht aus dem vom Senat formulierten -- von mir geteilten -- Grundsatz, daß die Beratungsregelung, weil sie auf ein Indikationsverfahren verzichtet, die Möglichkeit einer Rechtfertigung durch die allgemeine Notlagenindikation nicht verheißen kann. Dieser Grundsatz erklärt und begründet, daß der Frau, wie immer ihr Handeln nach materiell-rechtlichen Kriterien zu beurteilen sein mag, eine förmliche Rechtfertigung nicht zuteil werden kann. Er trägt aber nicht auch die zwingende Konsequenz, daß die beratenen Abbrüche über das Fehlen dieser Rechtfertigung hinaus unterschiedslos und unwiderleglich als materielles Unrecht auch und selbst dann zu qualifizieren sind, wenn sie den vom Senat selbst formulierten materiellen Rechtmäßigkeitskriterien genügen, zu denen -- in Übereinstimmung mit dem Urteil des Ersten Senats von 1975 (vgl. \abbrBVerfGE{} 39, 1 [49 \abbrF{}]) -- nach dem Kriterium der Unzumutbarkeit bestimmte Ausnahmesituationen, einschließlich einer qualifizierten sozialen Notlage, gehören (vgl. D.I.2.c.bb; D.III.1.c).
\end{sourcepara}
\begin{transpara}
\transrn{429}此一必要性並非來自第二庭所表述且我亦同意之原則：諮詢規制既然放棄指標事由程序，即不能許諾透過一般困境指標事由獲得正當化之可能。此一原則說明並證立，無論婦女之行為依實體法標準應如何評價，她均無法獲得形式上之正當化。但其並不支撐一項強制後果：經諮詢終止妊娠除欠缺此一正當化外，尚須一律且不可反駁地被定性為實體不法，即使其符合第二庭本身所表述之實體合法性標準亦然；依不可期待性標準，與第一庭1975年判決（參見 BVerfGE 39, 1 [49 f.]）一致，此等標準包括特定例外情境，含合格之社會困境（參見 D.I.2.c.bb；D.III.1.c）。
\end{transpara}

\begin{sourcepara}[title={原文SondervotumB.322}]
\sourcern{430}Der Senat kommt zu dieser einheitlichen Behandlung der beratenen Abbrüche als Unrecht deshalb, weil er entscheidend nicht auf Konfliktsituation und materiell-rechtliche Lage, sondern allein auf das Fehlen eines formellen Rechtfertigungsverfahrens abstellt. Es kann aber weder ein Erfordernis der Schutzpflicht für das ungeborene Leben noch des Rechtsstaatsprinzips sein, Schwangerschaftsabbrüche, die mit den materiellen Anforderungen der Rechtsordnung in Einklang stehen, für alle Folgefragen des Beratungsmodells unaufhebbar als materielles Unrecht zu qualifizieren. Mit solchen "Konsequenzen" würde der Rechtsstaat sich gegen sich selbst wenden.
\end{sourcepara}
\begin{transpara}
\transrn{430}第二庭之所以將經諮詢終止妊娠統一作為不法處理，是因為其決定性地並非著眼於衝突情境與實體法狀態，而是僅著眼於形式正當化程序之欠缺。然而，將與法律秩序實體要求相符之終止妊娠，在諮詢模式所有後續問題上不可取消地定性為實體不法，既不能是對未出生生命保護義務之要求，也不能是法治國原則之要求。以此等「後果」而言，法治國將轉而反對其自身。
\end{transpara}

\begin{sourcepara}[title={原文SondervotumB.323}]
\sourcern{431}b) Die sozialversicherungsrechtlich einheitliche Behandlung der beratenen Abbrüche als rechtswidrig führt zur generellen Verweigerung von Leistungen der Sozialversicherung. Dies hält der Senat als Folge der erforderlichen grundsätzlichen rechtlichen Mißbilligung des Schwangerschaftsabbruches für geboten. Damit werden notwendig auch diejenigen Abbrüche, die nach den vom Senat dargelegten materiellen Maßstäben die Voraussetzung für eine Rechtfertigung erfüllen, in die rechtliche Mißbilligung mit einbezogen.
\end{sourcepara}
\begin{transpara}
\transrn{431}b) 在社會保險法上將經諮詢終止妊娠統一作為違法處理，會導致一般性拒絕社會保險給付。第二庭認為，這是所需之對終止妊娠原則性法律否定評價之結果，故屬必要。如此一來，依第二庭所闡述之實體標準符合正當化前提之終止妊娠，也必然被納入法律否定評價。
\end{transpara}

\begin{sourcepara}[title={原文SondervotumB.324}]
\sourcern{432}Die rechtliche Mißbilligung auch dieser Abbrüche ist um so fragwürdiger, als der Senat, worin ich ihm zustimme, auch im Blick auf die nähere Ausgestaltung des Beratungskonzepts am Indikationsgedanken festhält: Der Schwangerschaftsabbruch ist, weil Tötung menschlichen Lebens, grundsätzlich Unrecht und verboten und nur besondere, eng zu umgrenzende Ausnahmelagen können ihn -- materiell gesehen -- gerechtfertigt erscheinen lassen (vgl. D.III.1.c). Die Beratung wird dazu verpflichtet, sich zu vergewissern, daß der den Abbruch erwägenden Frau diese rechtlich-normative Orientierung -- die Umfang und Grenzen der Rechtspflicht zum Austragen des Kindes aufzeigen soll -- bewußt ist, und evtl. vorhandene Fehlvorstellungen zu korrigieren.
\end{sourcepara}
\begin{transpara}
\transrn{432}對此等終止妊娠亦作法律否定評價尤其可疑，因為第二庭在諮詢構想之較詳細形塑上也堅持指標事由思想，而我同意此點：終止妊娠因係殺害人類生命，原則上是不法且被禁止，只有特殊且須嚴格界定之例外情境，始能使其在實體上顯得正當化（參見 D.III.1.c）。諮詢因此被賦予義務，須確認考慮終止妊娠之婦女知悉此一法律規範性指引，而該指引應指出繼續懷孕之法律義務範圍與界限，並修正可能存在之錯誤觀念。
\end{transpara}

\begin{sourcepara}[title={原文SondervotumB.325}]
\sourcern{433}Wird die einen Schwangerschaftsabbruch erwägende Frau in dieser Weise in die Verantwortung gerufen und von ihr die Bereitschaft erwartet, im Rahmen der Anforderungen der Rechtsordnung zu handeln, ist es widersprüchlich, zugleich ein Verfassungsgebot zu postulieren, nach dem in allen Rechtsbereichen außerhalb der Strafbarkeit alle Frauen, die einen Abbruch vornehmen lassen, unterschiedslos als unrechtmäßig Handelnde zu behandeln sind, ohne daß sie sich dagegen wehren können. Die Frau soll sich an den Anforderungen der Rechtsordnung zum Schutz des ungeborenen Lebens orientieren, aber in jedem Fall, auch wenn sie das tut, ist und bleibt ihr Handeln -- kraft Verfassungsgebots -- Unrecht. Das ist nicht nur sachwidrig, es trifft auch die Frauen als Person, in ihrer Ehre und Rechtschaffenheit.
\end{sourcepara}
\begin{transpara}
\transrn{433}若以此方式要求考慮終止妊娠之婦女承擔責任，並期待其準備在法律秩序要求之框架內行動，則同時主張一項憲法命令，要求在刑罰之外所有法律領域中，使所有實施終止妊娠之婦女不加區別地被作為違法行為人處理，且無從抗辯，乃是矛盾的。婦女應以法律秩序為保護未出生生命所設要求為取向；但在任何情況下，即使她如此為之，其行為依憲法命令仍然是不法且始終是不法。這不僅事理不當，也打擊婦女作為人格者之身分、其名譽與正直。
\end{transpara}

\begin{sourcepara}[title={原文SondervotumB.326}]
\sourcern{434}c) Ebensowenig ist dem Senat zu folgen, wenn er sich darauf beruft, daß das Beratungskonzept nach seinen Wirkungsbedingungen nicht erfordere, die der Frau überlassene Letztverantwortung für einen Schwangerschaftsabbruch sozialrechtlich anzuerkennen (vgl. E.V.2.b.cc). Entscheidend ist nicht, ob das Beratungskonzept zu seiner Wirksamkeit Versicherungsleistungen bei einem Abbruch als sinnvolle Schutzmaßnahme erfordert -- dies läßt sich mit Grund in Abrede stellen --, sondern daß das Beratungskonzept es, wie der Senat darlegt, mit Notwendigkeit ausschließt, die Rechtmäßigkeit von Schwangerschaftsabbrüchen -- außerhalb der speziellen Indikationen -- feststellen zu lassen. Ist dies aber ein Erfordernis des Beratungskonzepts (vgl. D.III.1.c), so kann daraus nicht zugleich das weitere Erfordernis hergeleitet werden, Frauen, die bei einem Abbruch nach den vom Senat dargelegten materiellen Kriterien nicht rechtswidrig handeln, dennoch ins Unrecht zu setzen, wenn der Senat den eigenen Prämissen (vgl. D.III.3.) treu bleibt. Denn es ist dann eine Situation gegeben, in der davon abgesehen werden darf, Schwangerschaftsabbrüche als Unrecht zu behandeln, obwohl nicht feststeht, daß sie gerechtfertigt sind, weil nämlich andernfalls eine erhebliche Beeinträchtigung von Frauen in ihrem Persönlichkeitsrecht zu besorgen ist. Der Senat anerkennt, daß die Frau davor geschützt sein müsse, den Abbruch und ihre Gründe hierfür außerhalb der Beratung und des ärztlichen Gesprächs weiteren Personen, etwa dem Arbeitgeber, unter Beeinträchtigung ihres Persönlichkeitsrechts offenbaren zu müssen, und rechtfertigt aus diesem Grund die Lohnfortzahlung bei beratenen Abbrüchen (vgl. E.V.4.). Wird aber nicht das Persönlichkeitsrecht der Frau durch die unterschiedslose Behandlung der beratenen Abbrüche als Unrecht, die zu beseitigen oder auch nur infrage zu stellen ihr auch dann nicht möglich ist, wenn der Abbruch materiell mit den Anforderungen der Rechtsordnung in Einklang steht, zumindest in gleichem wenn nicht sogar höherem Maß beeinträchtigt?
\end{sourcepara}
\begin{transpara}
\transrn{434}c) 第二庭若主張，依諮詢構想之作用條件，並不要求在社會法上承認交由婦女承擔之終止妊娠最後責任（參見 E.V.2.b.cc），亦不能追隨。決定性者並非諮詢構想為其有效性是否要求在終止妊娠時提供保險給付作為有意義之保護措施；這點可有理由地否定。決定性者是，如第二庭所闡述，諮詢構想必然排除使人確認特殊指標事由以外終止妊娠之合法性。若這是諮詢構想之要求（參見 D.III.1.c），則不能同時由此推導出另一項要求：若第二庭忠於其自身前提（參見 D.III.3.），仍將依第二庭所闡述實體標準在終止妊娠時並不具違法性行事之婦女置於不法之中。因為此時存在一種情境：雖未確定終止妊娠受到正當化，仍得不將其作為不法處理，因為否則將有嚴重侵害婦女人格權之虞。第二庭承認，婦女必須受到保護，避免其在諮詢與醫師談話之外，還必須向其他人，例如雇主，揭露終止妊娠及其理由而損及其人格權，並基於此理由正當化經諮詢終止妊娠時之工資繼續支付（參見 E.V.4.）。然而，將經諮詢終止妊娠不加區別地作為不法處理，即使終止妊娠在實體上符合法律秩序要求，婦女也無法排除或甚至質疑此一處理，這難道不會至少以相同程度，甚至更高程度侵害婦女人格權嗎？
\end{transpara}

\begin{sourcepara}[title={原文SondervotumB.327}]
\sourcern{435}d) Auch die notwendige Stützung und Stärkung des Rechtsbewußtseins im Blick auf das Lebensrecht der Ungeborenen vermag ein absolutes verfassungsrechtliches Verbot, für beratene Abbrüche Leistungen der Sozialversicherung zu gewähren, nicht zu tragen. Mit dem Senat stimme ich darin überein, daß die Stärkung und Lebendigerhaltung des Rechtsbewußtseins für den Schutz des ungeborenen Lebens von großer Bedeutung ist und verhindert werden muß, daß Schwangerschaftsabbrüche als soziale Normalität erscheinen und als rechtlich "freigegeben" eingestuft werden. Hiernach kommt es darauf an, ein Rechtsbewußtsein zu bilden und zu stützen, das -- gemäß den dargelegten verfassungsrechtlichen Maßgaben -- zwischen gerechtfertigten und unrechtmäßigen Schwangerschaftsabbrüchen zu unterscheiden weiß. Ein solches Rechtsbewußtsein wird aber nicht gefördert, sondern nimmt eher Schaden, wenn einerseits die hierauf bezogene rechtlich-normative Orientierung so stark betont wird wie es der Senat tut -- ich halte dies für notwendig und richtig --, es dann aber konkret, noch dazu von Verfassungs wegen, auf Recht und Unrecht gar nicht ankommt: Auch die Frauen, die sich in einer Notlage befinden, die den Schwangerschaftsabbruch zu rechtfertigen vermag, werden unaufhebbar der rechtlichen Mißbilligung unterworfen und von Leistungen der Sozialversicherung ausgeschlossen. Unabhängig von ihrer konkreten Situation wird Frauen unterschiedslos die Hypothek und Last auferlegt, von Rechts wegen so behandelt zu werden, wie wenn sie unrechtmäßig gehandelt hätten; als Hauptsache erscheint, daß jedenfalls nichts gezahlt wird.
\end{sourcepara}
\begin{transpara}
\transrn{435}d) 鑑於未出生者生命權而有必要支撐並強化法律意識，也不足以支撐一項絕對之憲法禁止，禁止對經諮詢終止妊娠給予社會保險給付。我同意第二庭，強化並維持法律意識之活力，對保護未出生生命具有重大意義，且必須防止終止妊娠顯得是社會常態並被歸類為法律上「放開」。依此，關鍵在於形成並支撐一種法律意識，使其能依所闡述之憲法準則，區分正當化與違法之終止妊娠。然而，若一方面如第二庭所為，如此強烈強調與此相關之法律規範性指引，而我認為這是必要且正確的，但另一方面在具體上，且竟是依憲法要求，根本不問合法與不法，則此種法律意識不會被促進，反而更可能受損。即使婦女處於足以正當化終止妊娠之困境中，仍不可取消地受法律否定評價，並被排除於社會保險給付之外。不論其具體處境如何，婦女均不加區別地被加諸一項負擔與重荷，即在法律上被如同其已違法行事般對待；主要之事彷彿只是無論如何不支付任何款項。
\end{transpara}

\begin{sourcepara}[title={原文SondervotumB.328}]
\sourcern{436}3. Die gegen die Auffassung des Senats erhobenen Einwände besagen allerdings nicht, daß die beratenen Abbrüche etwa von Verfassungs wegen zu finanzieren wären. Dem steht das vom Senat zu Recht gesehene Prinzip entgegen, daß rechtswidrige Schwangerschaftsabbrüche aus Gründen der Schutzwirkung für das ungeborene Leben und Erfordernissen der rechtsstaatlichen Ordnung nicht vom Staat -- oder staatlich veranlaßt -- finanziert werden dürfen. Damit wird allerdings das Dilemma deutlich, das durch die nicht auflösbare Gesamtheit beratener Abbrüche hervorgerufen wird. Dieses Dilemma ist eine notwendige Konsequenz des Beratungskonzepts und kann nicht bruchlos nach der einen oder anderen Seite aufgelöst werden. Läßt sich hinsichtlich der Gesamtheit der beratenen Abbrüche nicht zwischen rechtmäßig und rechtswidrig im Sinne der durch die Verfassung gesetzten Vorgaben unterscheiden, können diese Vorgaben auch nicht durch eine wie immer geartete Finanzierungsregelung verdeutlicht werden. Aus der verfassungsrechtlichen Schutzpflicht ergeben sich Kriterien dafür, wann Schwangerschaftsabbrüche erlaubt sein können und wann sie verboten sind, über die Behandlung einer Gesamtheit, die gerade in Bezug hierauf unscheidbar ist, vermag sie nichts auszusagen. Dann aber kann es nicht ein Erfordernis dieser Schutzpflicht und mithin Verfassungsgebot sein, das dargelegte Dilemma gerade und nur in einer Richtung aufzulösen. Vielmehr ist der Gesetzgeber berufen, eine Regelung und Entscheidung zu finden.
\end{sourcepara}
\begin{transpara}
\transrn{436}3. 然而，對第二庭見解所提出之異議，並不表示經諮詢終止妊娠在憲法上應受資助。第二庭正確看見之原則與此相牴觸：基於對未出生生命之保護效果及法治國秩序之要求，違法終止妊娠不得由國家或由國家促成而資助。這也清楚顯示由經諮詢終止妊娠不可解消之整體所造成之兩難。此一兩難是諮詢構想之必要後果，不能無縫地朝任一方向解消。若就經諮詢終止妊娠之整體而言，無法依憲法所設定之要求區分合法與違法，則此等要求也不能透過任何形式之資助規制被清楚表達。由憲法保護義務可得出何時終止妊娠可能被允許、何時被禁止之標準；但對於一個恰恰就此不可區分之整體，該保護義務無法作出陳述。既然如此，要求所述兩難恰好且僅能朝一個方向解消，即不能成為此一保護義務之要求，因而也不能成為憲法命令。毋寧，立法者負有尋找規制與決定之任務。
\end{transpara}

\begin{sourcepara}[title={原文SondervotumB.329}]
\sourcern{437}4. Damit tritt der Kern der Schwierigkeit, der der Senat nach meiner Auffassung nicht gerecht geworden ist, hervor:
\end{sourcepara}
\begin{transpara}
\transrn{437}4. 由此，第二庭依我見解未能妥善處理之困難核心即浮現出來：
\end{transpara}

\begin{sourcepara}[title={原文SondervotumB.330}]
\sourcern{438}Der entscheidende Schritt des Senats liegt, verfassungsrechtlich gesehen, im Übergang zu einem Schutzkonzept der Beratungsregelung. Dieser Übergang erscheint im Sinne des verfassungsrechtlich gebotenen Schutzes für das ungeborene Leben deshalb grundsätzlich möglich, weil die staatliche Schutzpflicht darauf ausgerichtet ist, durch rechtliche Regelungen und tatsächliche Maßnahmen einen wirksamen Schutz des ungeborenen Lebens zu erreichen und ihr daher nicht genügt ist, wenn zwar auf der normativen Ebene eine in sich geschlossene, auf stringenten Schutz des ungeborenen Lebens ausgerichtete Regelung vorhanden ist, die sich aber -- aus welchen Gründen immer -- nicht in tatsächlich wirksamen Schutz auch für das einzelne menschliche Leben umsetzt oder umzusetzen vermag. Die Gründe dafür sind im Urteil im einzelnen dargelegt.
\end{sourcepara}
\begin{transpara}
\transrn{438}從憲法上看，第二庭之決定性步驟在於轉向諮詢規制之保護構想。此一轉向之所以在憲法所要求之未出生生命保護意義上原則上可能，是因為國家保護義務旨在透過法律規制與實際措施，達成對未出生生命之有效保護；因此，若在規範層面上雖存在一套自身封閉、指向嚴格保護未出生生命之規制，但其無論基於何種原因，並未或無法轉化為對個別人類生命亦實際有效之保護，即不足以滿足該保護義務。判決已詳細說明其理由。
\end{transpara}

\begin{sourcepara}[title={原文SondervotumB.331}]
\sourcern{439}Wird dieser Schritt aber als ein Weg zur Erfüllung der staatlichen Schutzpflicht für das ungeborene Leben für zulässig angesehen -- ungeachtet der dabei notwendigen Anforderungen an seine nähere Ausgestaltung -- so muß er einerseits in seinen Notwendigkeiten und Wirksamkeitsbedingungen akzeptiert werden, zugleich aber mit anderen rechtlichen Positionen und Gesichtspunkten ins Verhältnis gesetzt und "ausgemittelt" werden. Das erfordert an manchen Stellen Kompromisse, die -- gemessen an einem in sich geschlossenen normativen Konzept -- schmerzlich erscheinen, dennoch aber unvermeidlich sind. In diesem Sinne gibt es "Kosten" des Beratungskonzepts, die auch an späterer Stelle, etwa bei der Finanzierung, nicht wieder eingefangen werden können, ohne neue Verwerfungen hervorzurufen, die daraus entstehen, daß andere -- als solche beachtliche -- rechtliche Positionen übergangen werden.
\end{sourcepara}
\begin{transpara}
\transrn{439}然而，若承認此一步驟作為履行國家對未出生生命保護義務之一條道路是可允許的，無論其具體形塑仍須符合必要要求，則一方面必須接受其必要性與作用條件，同時也必須將其與其他法律地位及觀點置於關係之中並作「調節」。這在若干位置要求妥協；以一個自身封閉之規範構想衡量，此等妥協或許顯得令人痛苦，然而仍不可避免。在此意義上，諮詢構想存在「成本」；此等成本在後續位置，例如資助問題上，也不能在不引發新失衡之情況下再被追回，而此等新失衡乃因其他本身值得注意之法律地位被越過而產生。
\end{transpara}

\begin{sourcepara}[title={原文SondervotumB.332}]
\sourcern{440}Eine gesetzliche Regelung des Abtreibungsproblems wird stets ein Stück weit "Notordnung" sein und sein müssen, wenn es ihr nicht nur darum geht, in der normativen Sphäre ein stringent auf den Schutz des ungeborenen Lebens ausgerichtetes Regelwerk zu formulieren, dessen praktische Wirksamkeit dahinstehen mag, sondern auch und gerade darum, soweit möglich wirksamen Schutz für das einzelne, ungeborene menschliche Leben zu erreichen. Eine rechtlich normative Ordnung, die in der normativen Sphäre bei sich bleibt, sich nicht in tatsächlich wirksamen Schutz des ungeborenen Lebens umsetzt, trägt zum Schutz dieses Lebens nicht wirklich bei und erfüllt damit auch nicht die Aufgabe, die dem Recht obliegt. Eine rechtliche Regelung aber, die gemäß der sozialordnenden Aufgabe des Rechts auf tatsächliche Wirksamkeit abzielt und sie herbeiführen will, muß auch die eigenen Wirksamkeitsbedingungen mit in Rechnung stellen. Diese ergeben sich ebenso aus der conditio humana wie aus der konkreten Verfaßtheit einer Gesellschaft. Diese Wirksamkeitsbedingungen gegebenenfalls zu ändern, kann nicht allein oder primär die Aufgabe des Rechts sein; es vermag dazu zwar, aber nur in begrenztem Umfang, beizutragen.
\end{sourcepara}
\begin{transpara}
\transrn{440}對終止妊娠問題之法律規制，若其不僅旨在規範領域中表述一套嚴格指向未出生生命保護之規則體系，而不問其實際有效性如何；而且尤其旨在盡可能達成對個別未出生人類生命之有效保護，則其始終在某種程度上會且必須是一種「緊急秩序」。一項法律規範秩序若停留於規範領域自身，而不轉化為對未出生生命之實際有效保護，則並未真正有助於保護此一生命，也因而未履行法律所負之任務。但一項依法律之社會秩序形成功能而以實際有效性為目標並欲促成之法律規制，也必須將其自身作用條件一併納入計算。此等條件既來自 \textit{conditio humana}，亦來自一個社會之具體構成。必要時改變此等作用條件，不能單獨或主要是法律之任務；法律固然能對此有所貢獻，但僅在有限範圍內。
\end{transpara}

\stoprandnumbers

\begin{sourcepara}[title={原文SondervotumB.333}]
Böckenförde
\end{sourcepara}
\begin{transpara}
Böckenförde
\end{transpara}

\bilingualstop

\phantomsection
\addcontentsline{toc}{section}{Glossar / 譯語表}
\ifbilingualcolumns
  \begin{center}
  \begin{tabular}{@{}p{0.44\textwidth}@{\hspace{0.52cm}}|@{\hspace{0.52cm}}p{0.44\textwidth}@{}}
  \makebox[\linewidth][c]{\begin{minipage}{0.86\linewidth}
  \centering
  {\Large\bfseries Glossar\par}
  \vspace{0.28em}
  {\footnotesize\color{pendingink}Ausgewählte Terminologie des Urteils\par}
  \end{minipage}}
  &
  \makebox[\linewidth][c]{\begin{minipage}{0.86\linewidth}
  \centering\songfont
  {\Large\bfseries 譯語表\par}
  \vspace{0.28em}
  {\footnotesize\color{pendingink}本文件採用之主要法律術語對照\par}
  \end{minipage}}
  \end{tabular}
  \end{center}
  \vspace{0.45em}
\else
  \ifshowtranslation
    \section*{譯語表}
  \else
    \section*{Glossar}
  \fi
\fi

\begin{termbox}
\noindent{\songfont\footnotesize\color{pendingink}
譯語說明：本文正文原則上將 \textit{Schwangerschaftsabbruch} 譯為「終止妊娠」，以保留德文法律與醫療語境中「中止懷孕狀態」的制度性語感，避免中文「墮胎」一詞預先將敘事重心置於胎兒或道德非難。若德文原文使用 \textit{Abtreibung}，或中文學界通稱、歷史語境須凸顯較具價值色彩之用語時，譯為「墮胎」或於語境中另行標示。\par}
\vspace{0.45em}
\begin{termcolumns}
\begin{termitems}
\item[Schwangerschaftsabbruch] 終止妊娠
\item[Abbruch / Abtreibung] 中止、終止；墮胎（依語境）
\item[ungeborenes Leben / ungeborenes menschliches Leben] 未出生生命；未出生人類生命
\item[vorgeburtliches / werdendes Leben] 出生前生命；形成中生命
\item[Embryo / Leibesfrucht] 胚胎；胎兒
\item[nasciturus] \textit{nasciturus}；胎兒；未出生者
\item[Schwangere / schwangere Frau] 懷孕婦女
\item[Austragen des Kindes / Austragung der Schwangerschaft] 繼續懷孕；繼續懷孕至分娩
\item[Schwangerschaftskonflikt] 懷孕衝突
\item[Schutzpflicht] 保護義務
\item[Schutzauftrag] 保護任務
\item[Lebensschutz] 生命保護
\item[Untermaßverbot] 不足禁止
\item[Übermaßverbot] 過度禁止
\item[Verhältnismäßigkeitsgrundsatz] 比例原則
\item[Abwägungs- und Gestaltungsspielraum] 衡量及形成空間
\item[Grundrechtsposition] 基本權地位
\item[Menschenwürde] 人性尊嚴
\item[Persönlichkeitsrecht] 人格權
\item[Recht auf körperliche Unversehrtheit] 身體完整權
\item[Selbstbestimmungsrecht] 自決權
\item[Zweiheit in Einheit] 一體中之二元
\item[Letztverantwortung] 最後責任
\item[Eigenverantwortung / eigenverantwortliche Entscheidung] 自主責任；自主負責決定
\item[Beratungskonzept] 諮詢構想
\item[Beratungsregelung] 諮詢規制
\item[Beratungslösung] 諮詢方案
\item[Beratungsgespräch] 諮詢談話
\item[Beratungsstelle / Beratungseinrichtung] 諮詢所；諮詢機構
\item[Beratungsschein] 諮詢證明
\item[Aufklärung] 告知；說明；啟蒙（依語境）
\item[normative Orientierung] 規範性指引
\item[Rechtsbewußtsein] 法律意識
\item[rechtliche Mißbilligung] 法律否定評價
\item[Indikation / Indikationstatbestand] 指標事由；指標事由構成要件
\item[medizinische Indikation] 醫學指標事由
\item[embryopathische Indikation] 胚胎病理指標事由
\item[kriminologische Indikation] 犯罪學指標事由
\item[allgemeine Notlagenindikation] 一般困境指標事由
\item[Not- und Konfliktlage] 緊急與衝突狀態
\item[soziale Notlage] 社會困境
\item[Unzumutbarkeit / unzumutbar] 不可期待性；不可期待
\item[Zumutbarkeit / zumutbar] 可期待性；可期待
\item[zumutbare Opfergrenze] 可期待之犧牲界限
\item[rechtfertigender Notstand] 正當化緊急避難
\item[Rechtfertigungsgrund] 正當化事由
\item[Schuldausschließungsgrund] 阻卻罪責事由
\item[Strafausschließungsgrund] 阻卻刑罰事由
\item[Rechtswidrigkeit / nicht rechtswidrig] 違法性；不具違法性
\item[Unrecht] 不法
\item[Tatbestand / Tatbestandsausschluß] 構成要件；構成要件排除
\item[Strafbarkeit / Strafdrohung] 可罰性；刑罰威嚇
\item[Bestimmtheitsgrundsatz] 明確性原則
\item[Fristenregelung / Fristenlösung] 期限規定
\item[Zwölf-Wochen-Frist] 十二週期限
\item[seit / nach der Empfängnis] 受孕後（如：受孕後十二週）；5. StrRG 及 DDR 法採此計算基準
\item[Schwangerschaftswochen] 妊娠週數（自末次月經起算，較受孕早約兩週）；SFHG §\,218a n.F. 採此基準；如：妊娠最初十二週
\item[Schwangeren- und Familienhilfegesetz] 孕婦及家庭扶助法
\item[Einigungsvertrag] 統一條約
\item[Beitrittsgebiet / neue Länder] 加入地區；新邦
\item[gesetzliche Krankenversicherung] 法定醫療保險
\item[Sozialversicherung] 社會保險
\item[Krankenkasse] 醫療保險基金
\item[Solidargemeinschaft] 團結共同體
\item[Leistungspflicht / Leistungskatalog] 給付義務；給付目錄
\item[Lohnfortzahlung] 工資繼續支付
\item[Sicherstellungsauftrag] 確保任務
\item[ärztliches Weigerungsrecht] 醫師拒絕權
\item[Beobachtungspflicht] 觀察義務
\item[Nachbesserungspflicht] 改善義務
\item[Übergangsregelung] 過渡規制
\item[Anordnung nach § 35 BVerfGG] 依《聯邦憲法法院法》第35條之命令
\item[Urteilsformel] 判決主文
\item[Antrag（向法院）] 聲請
\item[Antrag（議會動議）] 提案；對應的 Ablehnung 譯「否決」
\item[Antrag（一般行政）] 申請
\item[Antragsteller] 聲請人
\item[abweichende Meinung / Sondervotum] 不同意見
\item[obiter dictum] \textit{obiter dictum}；附帶意見
\item[aliud] \textit{aliud}；另一事物
\item[conditio humana] \textit{conditio humana}；人的處境
\end{termitems}
\end{termcolumns}

\termsectionheading{德文用語簡稱}
\begin{termcolumns}
\begin{abbritems}
\item[GG] \textit{Grundgesetz}；《基本法》
\item[BVerfG] \textit{Bundesverfassungsgericht}；聯邦憲法法院
\item[BVerfGE] \textit{Entscheidungen des Bundesverfassungsgerichts}；《聯邦憲法法院判決集》
\item[BVerfGG] \textit{Bundesverfassungsgerichtsgesetz}；《聯邦憲法法院法》
\item[StGB] \textit{Strafgesetzbuch}；《刑法》
\item[StGB a.F.] \textit{Strafgesetzbuch alter Fassung}；舊版本《刑法》
\item[StGB n.F.] \textit{Strafgesetzbuch neuer Fassung}；新版本《刑法》
\item[StÄG] \textit{Strafrechtsänderungsgesetz}；刑法修正法
\item[5. StrRG] \textit{Fünftes Gesetz zur Reform des Strafrechts}；《第五刑法改革法》
\item[StREG] \textit{Strafrechtsreform-Ergänzungsgesetz}；刑法改革補充法
\item[SFHG] \textit{Schwangeren- und Familienhilfegesetz}；《孕婦及家庭扶助法》
\item[SGB V] \textit{Fünftes Buch Sozialgesetzbuch}；《社會法典》第五編
\item[RVO] \textit{Reichsversicherungsordnung}；《帝國保險法》
\item[EV] \textit{Einigungsvertrag}；《統一條約》
\item[BTDrucks.] \textit{Bundestagsdrucksache}；德國聯邦議會印刷文件
\item[BRDrucks.] \textit{Bundesratsdrucksache}；德國聯邦參議院印刷文件
\item[BGBl.] \textit{Bundesgesetzblatt}；《聯邦法律公報》
\item[GBl. DDR] \textit{Gesetzblatt der Deutschen Demokratischen Republik}；《德意志民主共和國法律公報》
\item[StenBer.] \textit{Stenographischer Bericht}；速記紀錄
\item[WP.] \textit{Wahlperiode}；屆期
\item[a.a.O.] \textit{am angegebenen Ort}；前引處
\item[vgl.] \textit{vergleiche}；參見
\item[Abs.] \textit{Absatz}；項
\item[Art.] \textit{Artikel}；條
\item[Nr.] \textit{Nummer}；款、號
\item[Rdnr. / Rn.] \textit{Randnummer}；邊碼、段碼
\item[S.] \textit{Seite}；頁
\item[ff.] \textit{fortfolgende}；以下數頁、以下各段
\end{abbritems}
\end{termcolumns}
\end{termbox}

\end{document}
